\nonstopmode{}
\documentclass[a4paper]{book}
\usepackage[times,inconsolata,hyper]{Rd}
\usepackage[utf8]{inputenc} % @SET ENCODING@
% \usepackage{graphicx} % @USE GRAPHICX@
\begin{document}
\chapter*{}
\begin{center}
{\textbf{\huge Package `jjstat'}}
\par\bigskip{\large \today}
\end{center}
\inputencoding{utf8}
\ifthenelse{\boolean{Rd@use@hyper}}{\hypersetup{pdftitle = {jjstat: Statistical Analysis Tools for Writing Papers}}}{}
\begin{description}
\raggedright{}
\item[Type]\AsIs{Package}
\item[Title]\AsIs{Statistical Analysis Tools for Writing Papers}
\item[Version]\AsIs{0.1.9}
\item[Author]\AsIs{Park Joonghee}
\item[Maintainer]\AsIs{The package maintainer }\email{shoutjoy@hanmail.net}\AsIs{}
\item[Description]\AsIs{More about what it does (maybe more than one line)
Use four spaces when indenting paragraphs within the Description.}
\item[License]\AsIs{GPL-3}
\item[Encoding]\AsIs{UTF-8}
\item[LazyData]\AsIs{true}
\item[Imports]\AsIs{dplyr ,
knitr ,
ggplot2 (>= 3.4.4),
magrittr (>= 2.0.3),
Hmisc,
ggrepel}
\item[Suggests]\AsIs{broom (>= 1.0.5),
httpuv (>= 1.6.13),
rvest (>= 1.0.3),
tibble (>= 3.2.1),
kableExtra (>= 1.3.4),
testthat (>= 3.0.0),
psych,
gridExtra}
\item[Depends]\AsIs{tidyverse,
lavaan,
broom,
plspm,
R (>= 2.10),}
\item[Config/roxygen2/version]\AsIs{8.1.1}
\end{description}
\Rdcontents{\R{} topics documented:}
\inputencoding{utf8}
\HeaderA{accent\_table}{aaccent\_table: table() analytics as data frame, mosaicplot, usign add\_ratio}{accent.Rul.table}
%
\begin{Description}
aaccent\_table: table() analytics as data frame, mosaicplot, usign add\_ratio
\end{Description}
%
\begin{Usage}
\begin{verbatim}
accent_table(
  data,
  cols = "a1",
  rows = "성조",
  title = "",
  trans = TRUE,
  cex = 1.3,
  color = TRUE,
  ylab = "onset",
  xlab = "accent",
  sub = NULL,
  plot = FALSE,
  type = "res",
  raw = FALSE
)
\end{verbatim}
\end{Usage}
%
\begin{Arguments}
\begin{ldescription}
\item[\code{data}] table data

\item[\code{cols}] column names in data

\item[\code{rows}] row names in data

\item[\code{title}] table main

\item[\code{trans}] transpose

\item[\code{cex}] graph text size

\item[\code{color}] graph color 1:2, mosaicplot

\item[\code{ylab}] ylab

\item[\code{xlab}] xlab

\item[\code{sub}] subtitle

\item[\code{plot}] output option

\item[\code{type}] res, res\_df, ratio, g, all

\item[\code{raw}] raw data, or table data
\end{ldescription}
\end{Arguments}
%
\begin{Value}
matrix type data.frame
\end{Value}
%
\begin{Examples}
\begin{ExampleCode}

## Not run: 
mtcars %>% select(am, vs) %>% table() %>% accent_table("am","vs")

mtcars %>% select(am, vs) %>% table() %>% accent_table("am","vs", plot=TRUE)


mtcars %>% select(am, vs) %>% table() %>% accent_table("am","vs", plot=TRUE, title = "Mtcars")


## End(Not run)



\end{ExampleCode}
\end{Examples}
\inputencoding{utf8}
\HeaderA{add\_bind\_stat}{Functions that item parceling generate multiple pieces of data,}{add.Rul.bind.Rul.stat}
%
\begin{Description}
Functions that item parceling generate multiple pieces of data,
\end{Description}
%
\begin{Usage}
\begin{verbatim}
add_bind_stat(data, ..., fun = mean, type = "res")
\end{verbatim}
\end{Usage}
%
\begin{Arguments}
\begin{ldescription}
\item[\code{data}] data.fram

\item[\code{...}] colname, cols

\item[\code{fun}] fun  mean sd

\item[\code{type}] data, res --data.fram , res, cols\_check
\end{ldescription}
\end{Arguments}
%
\begin{Value}
data.frame
\end{Value}
%
\begin{Examples}
\begin{ExampleCode}
## Not run: 

#Common item parceling methods (group by average )
mtcars %>% mutate(
Mean_col1 = (mpg + cyl)/2 ,
Mean_col2 = (disp + hp)/2 ,
Mean_col3 = (drat + wt )/2 ,
Mean_col4 = (qsec + vs + am)/3 )

#Methods for item parceling using functions
add_bind_stat(mtcars,
c("Mean_col1","mpg", "cyl"),
c("Mean_col2", "disp", "hp"),
c("Mean_col3", "drat", "wt"),
c("Mean_col4", "qsec", "vs", "am") )



add_bind_stat(mtcars,
   c("Mean_col1", 1:2),
   c("Mean_col2", 3:4),
   c("Mean_col3", 5:6),
   c("Mean_col4", 7:9))
#'  ##col_check
#'
add_bind_stat(mtcars,
              c("Mean_col1", "mpg:cyl"),
              c("Mean_col2", "disp:hp"),
              c("Mean_col3", "drat:wt"),
              c("Mean_col4", "qsec:am"))

## End(Not run)

\end{ExampleCode}
\end{Examples}
\inputencoding{utf8}
\HeaderA{add\_col\_ratio}{그룹합에 비율을 만들어 주는 함수}{add.Rul.col.Rul.ratio}
%
\begin{Description}
그룹합에 비율을 만들어 주는 함수
\end{Description}
%
\begin{Usage}
\begin{verbatim}
add_col_ratio(data, col, n_col)
\end{verbatim}
\end{Usage}
%
\begin{Arguments}
\begin{ldescription}
\item[\code{data}] data

\item[\code{col}] 기준이 되는 열

\item[\code{n\_col}] 빈도를 나타내는 열
\end{ldescription}
\end{Arguments}
%
\begin{Value}
df
\end{Value}
%
\begin{Examples}
\begin{ExampleCode}
## Not run: 
#'
# 사용 예제
library(dplyr)

# 주어진 데이터프레임
data <- data.frame(
  Q17 = c("", "고1 남학생", "고1 여학생", "고2 남학생", "고2 여학생",
          "고3 남학생", "고3 여학생", "중1 남학생", "중1 여학생",
          "중2 남학생", "중2 여학생", "중3 남학생", "중3 여학생",
          "초1 남학생", "초1 여학생", "초2 남학생", "초2 여학생",
          "초3 남학생", "초3 여학생", "초4 남학생", "초4 여학생",
          "초5 남학생", "초5 여학생", "초6 남학생", "초6 여학생"),
  Q16 = c("", rep("", 24)),
  n = c(1130, 413, 184, 329, 163, 216, 72, 354, 323, 370, 308,
        195, 169, 7, 4, 25, 23, 60, 53, 174, 162, 237, 216, 307, 293)
)

# 함수 적용: 열 이름으로 지정
result1 <- add_col_ratio(data, col = "Q17", n_col = "n")

# 함수 적용: 숫자 인덱스로 지정
result2 <- add_col_ratio(data, col = 1, n_col = 3)

# 결과 확인
print(result1)
print(result2)



## End(Not run)
\end{ExampleCode}
\end{Examples}
\inputencoding{utf8}
\HeaderA{add\_cor}{correlation analysis}{add.Rul.cor}
%
\begin{Description}
correlation analysis
\end{Description}
%
\begin{Usage}
\begin{verbatim}
add_cor(selectdata, type = "res_all", digits = 3, plot = TRUE, md = FALSE)
\end{verbatim}
\end{Usage}
%
\begin{Arguments}
\begin{ldescription}
\item[\code{selectdata}] seledted daa

\item[\code{type}] res\_all, res, p, r, res\_p, res\_md, res\_lower, g(chart.Correlation)

\item[\code{digits}] default 3

\item[\code{plot}] graph performanceAnalytis::chart.Correlation

\item[\code{md}] markdown table output
\end{ldescription}
\end{Arguments}
%
\begin{Value}
table and picture
\end{Value}
%
\begin{Examples}
\begin{ExampleCode}

## Not run: 

mtcars%>%select(mpg, disp, wt, hp, qsec)%>%add_cor()

## End(Not run)
\end{ExampleCode}
\end{Examples}
\inputencoding{utf8}
\HeaderA{add\_id}{Add row id}{add.Rul.id}
%
\begin{Description}
Add row id
\end{Description}
%
\begin{Usage}
\begin{verbatim}
add_id(data, sel = "names", each = TRUE)
\end{verbatim}
\end{Usage}
%
\begin{Arguments}
\begin{ldescription}
\item[\code{data}] data.frame

\item[\code{sel}] select column

\item[\code{each}] each =TRUE then 111222....FALSE 123123
\end{ldescription}
\end{Arguments}
%
\begin{Value}
add column data
\end{Value}
%
\begin{Examples}
\begin{ExampleCode}

## Not run: 

mtcars %>% mutate(MPG= which(colnames(mtcars) =="mpg"))%>%add_id()
mtcars %>% mutate(MPG= which(colnames(mtcars) =="mpg"))%>%add_id("cyl")
mtcars[1:3, 1:3] %>% add_id()
mtcars[1:3, 1:3] %>% add_id("cyl")
mtcars[1:4, 1:4] %>% to_long() %>% add_id()
mtcars[1:4, 1:4] %>% to_long() %>% add_id(each=FALSE)




## End(Not run)
\end{ExampleCode}
\end{Examples}
\inputencoding{utf8}
\HeaderA{add\_normality}{Add skewness and kurtosis normality checks to descriptive statistics results}{add.Rul.normality}
%
\begin{Description}
Add skewness and kurtosis normality checks to descriptive statistics results
\end{Description}
%
\begin{Usage}
\begin{verbatim}
add_normality(data, skew = "Skew", kurt = "Kurt")
\end{verbatim}
\end{Usage}
%
\begin{Arguments}
\begin{ldescription}
\item[\code{data}] mysummary data

\item[\code{skew}] select skew

\item[\code{kurt}] select kurt
\end{ldescription}
\end{Arguments}
%
\begin{Value}
normalitydata
\end{Value}
%
\begin{Examples}
\begin{ExampleCode}

## Not run: 

mysummary(mtcars)
mysummary(mtcars)%>%add_normality()



## End(Not run)

\end{ExampleCode}
\end{Examples}
\inputencoding{utf8}
\HeaderA{add\_normality\_kline}{Add Normality Kline}{add.Rul.normality.Rul.kline}
%
\begin{Description}
This function evaluates normality for each variable in the dataset based on skewness and kurtosis.
\end{Description}
%
\begin{Usage}
\begin{verbatim}
add_normality_kline(df, skew = 3, kurt = 8)
\end{verbatim}
\end{Usage}
%
\begin{Arguments}
\begin{ldescription}
\item[\code{df}] A data frame containing statistical summary including Skew and Kurt columns.

\item[\code{skew}] Threshold for skewness. Default is 3.0.

\item[\code{kurt}] Threshold for kurtosis. Default is 8.0.
\end{ldescription}
\end{Arguments}
%
\begin{Value}
A data frame with added columns `skew\_sig` and `kurt\_sig` indicating the normality evaluation.
\end{Value}
%
\begin{Examples}
\begin{ExampleCode}
## Not run: 
data <- data.frame(
  var = c("SUX", "FLOW", "PU", "PEOU", "ATT", "USE"),
  N = c(895, 895, 895, 895, 895, 895),
  MEAN = runif(6, 2, 3),
  SD = runif(6, 0.9, 1.1),
  MIN = rep(1, 6),
  MAX = rep(5, 6),
  Skew = c(0.74, 0.23, 0.81, 0.75, 0.69, 0.53),
  Kurt = c(0.16, -0.31, 0.37, 0.06, 0.02, -0.42)
)
add_normality_kline(data)

## End(Not run)
\end{ExampleCode}
\end{Examples}
\inputencoding{utf8}
\HeaderA{add\_partial\_sum}{add\_partial\_sum}{add.Rul.partial.Rul.sum}
%
\begin{Description}
add\_partial\_sum
\end{Description}
%
\begin{Usage}
\begin{verbatim}
add_partial_sum(df, col, freq_col)
\end{verbatim}
\end{Usage}
%
\begin{Arguments}
\begin{ldescription}
\item[\code{df}] df

\item[\code{col}] select col

\item[\code{freq\_col}] select count frequency col
\end{ldescription}
\end{Arguments}
%
\begin{Value}
df
\end{Value}
%
\begin{Examples}
\begin{ExampleCode}
## Not run: 
#'
# 함수 사용 예시
df <- tibble::tribble(
  ~학교급, ~term, ~Freq, ~`prop(%)`,
  "-", "합계", 3110, 100,
  "고등학교", "고1", 341, 11.0,
  "고등학교", "고2", 322, 10.4,
  "고등학교", "고3", 192, 6.17,
  "중학교", "중1", 591, 19.0,
  "중학교", "중2", 527, 16.9,
  "중학교", "중3", 290, 9.32,
  "초등학교", "초1", 8, 0.257,
  "초등학교", "초2", 48, 1.54,
  "초등학교", "초3", 47, 1.51,
  "초등학교", "초4", 191, 6.14,
  "초등학교", "초5", 207, 6.66,
  "초등학교", "초6", 346, 11.1
)

# 함수 실행
result_df <- add_partial_sum(df, col = "학교급", freq_col = "Freq")
print(result_df)

## End(Not run)

\end{ExampleCode}
\end{Examples}
\inputencoding{utf8}
\HeaderA{add\_pca\_scores}{add\_pca\_scores}{add.Rul.pca.Rul.scores}
%
\begin{Description}
add\_pca\_scores
\end{Description}
%
\begin{Usage}
\begin{verbatim}
add_pca_scores(data, ..., method = "nipals", adjust = NULL)
\end{verbatim}
\end{Usage}
%
\begin{Arguments}
\begin{ldescription}
\item[\code{data}] data

\item[\code{...}] select variable

\item[\code{method}] niplas, prcomp

\item[\code{adjust}] Options to change the sign of prcomp c(F, T ...)
\end{ldescription}
\end{Arguments}
%
\begin{Value}
add data
\end{Value}
%
\begin{Examples}
\begin{ExampleCode}

## Not run: 
data(offense)

nfl_path = matrix(
  c(0, 0, 0, 0, 0,
    0, 0, 0, 0, 0,
    0, 0, 0, 0, 0,
    0, 1, 1, 0, 0,
    1, 0, 0, 1, 0),
  nrow = 5, ncol = 5, byrow = TRUE )

rownames(nfl_path) <- c("Speical", "Rushing", "Passing", "Offense", "Scoring")
colnames(nfl_path) <- c("Speical", "Rushing", "Passing", "Offense", "Scoring")
nfl_path
Offense = offense %>% add_pca_scores(rush1 = c(1:3), pass1 = c(4:6))

Offense =  add_pca_scores(offense,
               rush1 = c("YardsRushAtt", "RushYards", "RushFirstDown"),
               pass1 = c("YardsPassComp", "PassYards", "PassFirstDown"))


nfl_blocks2 = list(7:8, 1:3, 4:6,  18:19, 9:11)

plspm_blocks_match(data=Offense, blocks=nfl_blocks2,
                   paths=nfl_path)

plspm_blocks_match(data=Offense, blocks=nfl_blocks2,
                   paths=nfl_path)%>%make_list_text()

nfl_blocks2 = list(
  Speical = c("FieldGoals", "OtherTDs"),
  Rushing = c("YardsRushAtt", "RushYards", "RushFirstDown"),
  Passing = c("YardsPassComp", "PassYards", "PassFirstDown"),
  Offense = c("rush1", "pass1"),
  Scoring = c("PointsGame", "OffensTD", "TDGame")
)
nfl_blocks2
nfl_mode2 =c("B","A","A","A","A")
nfl_pls2 = plspm(Data=Offense, nfl_path, nfl_blocks2, modes = nfl_mode2)
nfl_pls2%>%summary()

nfl_pls_boot2 = plspm_sem(Data=Offense, nfl_path, nfl_blocks2, modes = nfl_mode)
x11()
nfl_pls_boot2 %>%plspm_path_coefs_plot()
nfl_pls_boot2 %>%plspm_semPaths2()

#'
#'

# 함수 테스트
offense <- data.frame(
  YardsRushAtt = rnorm(10),
  RushYards = rnorm(10),
  RushFirstDown = rnorm(10),
  YardsPassComp = rnorm(10),
  PassYards = rnorm(10),
  PassFirstDown = rnorm(10)
)

# 테스트 케이스 1
offense1 <- add_pca_scores(offense,
                           rush1 = c("YardsRushAtt", "RushYards", "RushFirstDown"),
                           pass1 = c("YardsPassComp", "PassYards", "PassFirstDown"))

str(offense1)

# 테스트 케이스 2
offense2 <- add_pca_scores(offense, method = "prcomp",
                           rush1 = c("YardsRushAtt", "RushYards", "RushFirstDown"),
                           pass1 = c("YardsPassComp", "PassYards", "PassFirstDown"))

str(offense2)

# 테스트 케이스 3 (adjust 옵션 포함)
offense3 <- add_pca_scores(offense, method = "prcomp", adjust = c(F, T),
                           rush1 = c("YardsRushAtt", "RushYards", "RushFirstDown"),
                           pass1 = c("YardsPassComp", "PassYards", "PassFirstDown"))

str(offense3)
#'
#'


## End(Not run)

\end{ExampleCode}
\end{Examples}
\inputencoding{utf8}
\HeaderA{add\_prop}{Add Proportions to Frequency Data}{add.Rul.prop}
%
\begin{Description}
This function adds a column to a data frame representing the proportion of each frequency value,
either as a percentage or a simple ratio.
\end{Description}
%
\begin{Usage}
\begin{verbatim}
add_prop(df, freq = "Freq", prop = TRUE)
\end{verbatim}
\end{Usage}
%
\begin{Arguments}
\begin{ldescription}
\item[\code{df}] A data frame containing frequency data.

\item[\code{freq}] A column name or index representing the frequency column. Defaults to "Freq".
If a numeric value is provided, it is converted to the corresponding column name.

\item[\code{prop}] Logical, whether to return proportions as percentages (TRUE, default) or ratios (FALSE).
\end{ldescription}
\end{Arguments}
%
\begin{Value}
A data frame with an additional column for proportions.
If `prop = TRUE`, the column is named `prop(
If `prop = FALSE`, the column is named `prop`.
\end{Value}
%
\begin{Examples}
\begin{ExampleCode}
## Not run: 
# Example data
df <- data.frame(Category = c("A", "B", "C"), Freq = c(30, 50, 20))

# Adding proportions as percentages
add_prop(df, freq = "Freq", prop = TRUE)

# Adding proportions as simple ratios
add_prop(df, freq = "Freq", prop = FALSE)

## End(Not run)
\end{ExampleCode}
\end{Examples}
\inputencoding{utf8}
\HeaderA{add\_ratio}{Convert matrix data to a ratio by using the sum of columns to get a ratio}{add.Rul.ratio}
%
\begin{Description}
Convert matrix data to a ratio by using the sum of columns to get a ratio
\end{Description}
%
\begin{Usage}
\begin{verbatim}
add_ratio(data, type = "res2", digits = 3)
\end{verbatim}
\end{Usage}
%
\begin{Arguments}
\begin{ldescription}
\item[\code{data}] matrix

\item[\code{type}] res2: *100, res1

\item[\code{digits}] defalut 3
\end{ldescription}
\end{Arguments}
%
\begin{Value}
ratio tabl e
\end{Value}
%
\begin{Examples}
\begin{ExampleCode}

## Not run: 


mat = matrix(c(36, 67, 11,
          0, 35, 44,
         41, 108, 1,
          56, 87, 54),
           nrow = 4, byrow = TRUE,
 dimnames = list(c("공명음", "마찰음",
                   "유기음_경음", "평파열음_평파찰음"),
                    c("H", "H(H)", "L")))

## step1
  mat %>% add_ratio()
##conbind
 mat %>% add_ratio_df()
##add margins
 mat %>% add_sum()


## End(Not run)

\end{ExampleCode}
\end{Examples}
\inputencoding{utf8}
\HeaderA{add\_ratio\_df}{Combine matrices and their ratios}{add.Rul.ratio.Rul.df}
%
\begin{Description}
Combine matrices and their ratios
\end{Description}
%
\begin{Usage}
\begin{verbatim}
add_ratio_df(res)
\end{verbatim}
\end{Usage}
%
\begin{Arguments}
\begin{ldescription}
\item[\code{res}] matrix
\end{ldescription}
\end{Arguments}
%
\begin{Value}
combine table /matrix
\end{Value}
%
\begin{Examples}
\begin{ExampleCode}

## Not run: 

#' #'mat = matrix(c(36, 67, 11,
          0, 35, 44,
         41, 108, 1,
          56, 87, 54),
           nrow = 4, byrow = TRUE,
   dimnames = list(c("공명음", "마찰음",
     "유기음_경음", "평파열음_평파찰음"),  c("H", "H(H)", "L")))

## step1
  mat %>% add_ratio()
##conbind
 mat %>% add_ratio_df()
##add margins
 mat %>% add_sum()

## End(Not run)


\end{ExampleCode}
\end{Examples}
\inputencoding{utf8}
\HeaderA{add\_rep\_id}{Create an ordered variable for overlapping groups}{add.Rul.rep.Rul.id}
%
\begin{Description}
Create an ordered variable for overlapping groups
\end{Description}
%
\begin{Usage}
\begin{verbatim}
add_rep_id(data, grp)
\end{verbatim}
\end{Usage}
%
\begin{Arguments}
\begin{ldescription}
\item[\code{data}] data.frame

\item[\code{grp}] select colname or number
\end{ldescription}
\end{Arguments}
%
\begin{Value}
create data
\end{Value}
%
\begin{Examples}
\begin{ExampleCode}

## Not run: 
#'
data <- data.frame(
  weight = c(4.17, 5.58, 5.18, 6.11, 4.50, 4.61, 5.17, 4.53, 5.33, 5.14,
             4.81, 4.17, 4.41, 3.59, 5.87, 3.83, 6.03, 4.89, 4.32, 4.69,
             6.31, 5.12, 5.54, 5.50, 5.37, 5.29, 4.92, 6.15, 5.80, 5.26),
  group = c(rep("ctrl", 10), rep("trt1", 10), rep("trt2", 10))
)
data

#generate id from Colname
add_rep_id(data, "group")

#generate id from Col number
add_rep_id(data, 2)

#data spread wide
data%>%add_rep_id("group")%>%as_trt("id") %>%wide_df("group", "weight")


data%>%
  add_rep_id("group")%>%
  as_trt("id") %>%
  wide_df("group", "weight")


## End(Not run)


\end{ExampleCode}
\end{Examples}
\inputencoding{utf8}
\HeaderA{add\_rows}{Insert data into the required act in the dataframe}{add.Rul.rows}
%
\begin{Description}
Insert data into the required act in the dataframe
\end{Description}
%
\begin{Usage}
\begin{verbatim}
add_rows(df, ..., row = NULL, char = FALSE)
\end{verbatim}
\end{Usage}
%
\begin{Arguments}
\begin{ldescription}
\item[\code{df}] data.frame

\item[\code{...}] data,  c(100, 200), c(200, 300),

\item[\code{row}] row position

\item[\code{char}] transpose character column
\end{ldescription}
\end{Arguments}
%
\begin{Value}
dagta
\end{Value}
%
\begin{Examples}
\begin{ExampleCode}
## Not run: 
#'
#'


df <- data.frame(A = 1:5, B = 6:10)
df
# 2번째 행부터 새 데이터(벡터 c(1,2,3,4)와 c(2,3,4,5))를 삽입

add_rows(df, c(100, 200), c(200, 300), row = 2)

df <- data.frame(
  측정변수 = c("진로목표(L01)", "준비여부(L02)", "직업결정(L03)", "부모추천(L04)", "학업성적(L05)"),
  Class.1 = c(0.2178, 0.3533, 0.1552, 0.2708, 0.1605),
  Class.2 = c(0.9888, 0.9769, 0.9910, 1.0000, 0.9804),
  Class.3 = c(0.9122, 0.8704, 0.9667, 0.0000, 0.6761)
)

# 1행에 새 데이터를 삽입
new_df <- add_rows(df, c("class name", "불명확 진로인식형", "적극적 진로준비형", "조건분 진로지향형"), row = 1)

print(new_df)


df <- data.frame(
  측정변수 = c("진로목표(L01)", "준비여부(L02)", "직업결정(L03)", "부모추천(L04)", "학업성적(L05)"),
  Class.1 = c(0.2178, 0.3533, 0.1552, 0.2708, 0.1605),
  Class.2 = c(0.9888, 0.9769, 0.9910, 1.0000, 0.9804),
  Class.3 = c(0.9122, 0.8704, 0.9667, 0.0000, 0.6761)
)

# 첫 번째 행에 새 데이터를 삽입, char = TRUE로 문자형 변환
new_df_1 <- add_rows(df, c("class name", "불명확 진로인식형", "적극적 진로준비형", "조건분 진로지향형"), row = 1, char = TRUE)

print(new_df_1)

# 마지막 행에 새 데이터를 삽입, char = FALSE로 기존 타입 유지
new_df_2 <- add_rows(df, c("class name", "불명확 진로인식형", "적극적 진로준비형", "조건분 진로지향형"), row = 7, char = FALSE)

print(new_df_2)



## End(Not run)

\end{ExampleCode}
\end{Examples}
\inputencoding{utf8}
\HeaderA{add\_rows\_freq}{Functions that generate rows with a frequency}{add.Rul.rows.Rul.freq}
%
\begin{Description}
Functions that generate rows with a frequency
\end{Description}
%
\begin{Usage}
\begin{verbatim}
add_rows_freq(data, type = "res", sel = ncol(data), fix = NULL)
\end{verbatim}
\end{Usage}
%
\begin{Arguments}
\begin{ldescription}
\item[\code{data}] data

\item[\code{type}] res, all

\item[\code{sel}] frequency column

\item[\code{fix}] final select column
\end{ldescription}
\end{Arguments}
%
\begin{Value}
data
\end{Value}
%
\begin{Examples}
\begin{ExampleCode}

## Not run: 
dff = data.frame(
  trt = c("r","s","z"),
  con = c("f","m","l"),
  num = c(10,0,2),
  fre = c(4,0,5) )
dff
dff %>%add_rows_freq()
dff %>%add_rows_freq(sel=4)
dff %>%add_rows_freq(sel="fre")
dff %>%add_rows_freq(sel=3)
dff %>%add_rows_freq(sel="num")
dff %>%add_rows_freq(fix=1)
dff %>%add_rows_freq(type= "res", fix=1:2)
dff %>%add_rows_freq(type= "res", fix="con")
dff %>%add_rows_freq(type= "all")
dff %>%add_rows_freq(type= "real")
dff %>%add_rows_freq(type= "zero")
dff %>%add_rows_freq(type= "zero", fix="con")


# Making frequency data into real data
dff %>%to_long()%>% add_rows_freq()
 dff_real =  dff %>%to_long()%>% add_rows_freq()
 dff_real


 freq_data =  mtcars %>% select(am, vs) %>% table()
 freq_data%>% data.frame()
 freq_data%>% data.frame() %>% add_rows_freq(type ="zero")
 freq_data%>% data.frame() %>% add_rows_freq(type ="zero") %>%table()



## End(Not run)


\end{ExampleCode}
\end{Examples}
\inputencoding{utf8}
\HeaderA{add\_row\_sum}{Add Row Sum to Data Frame}{add.Rul.row.Rul.sum}
%
\begin{Description}
This function calculates the sum of specified numeric columns and appends a summary row
to the input data frame. The summary row includes totals for `freq\_col` and `prop\_col`.
\end{Description}
%
\begin{Usage}
\begin{verbatim}
add_row_sum(df, freq_col = 2, prop_col = 3)
\end{verbatim}
\end{Usage}
%
\begin{Arguments}
\begin{ldescription}
\item[\code{df}] A data frame containing the data to be summed.

\item[\code{freq\_col}] Column name or index for the frequency column. Default is 2.

\item[\code{prop\_col}] Column name or index for the proportion column. Default is 3.
\end{ldescription}
\end{Arguments}
%
\begin{Value}
A data frame with an added summary row containing the sums.
\end{Value}
%
\begin{Examples}
\begin{ExampleCode}
## Not run: 
data <- data.frame(
  Group = c("A", "B", "C"),
  Freq = c(10, 20, 30),
  `prop(%)` = c(25.0, 50.0, 25.0)
)
add_row_sum(data, freq_col = "Freq", prop_col = "prop(%)")

## End(Not run)
\end{ExampleCode}
\end{Examples}
\inputencoding{utf8}
\HeaderA{add\_sig}{mutate data generate}{add.Rul.sig}
%
\begin{Description}
mutate data generate
\end{Description}
%
\begin{Usage}
\begin{verbatim}
add_sig(var)
\end{verbatim}
\end{Usage}
%
\begin{Arguments}
\begin{ldescription}
\item[\code{p.value}] pvalue
\end{ldescription}
\end{Arguments}
%
\begin{Value}
in data
\end{Value}
%
\begin{Examples}
\begin{ExampleCode}

## Not run: 

aov_table(data = mtcars, dv_var = "mpg",
iv_vars = c("cyl", "gear", "carb"), grp_mean = T,
unite_F=TRUE) %>%
 mutate(STAR= add_sig(p_value))


## End(Not run)
\end{ExampleCode}
\end{Examples}
\inputencoding{utf8}
\HeaderA{add\_significance\_symbols}{Functions that output statistical significance for matrices and dataframes Simple and exact expressions}{add.Rul.significance.Rul.symbols}
%
\begin{Description}
Functions that output statistical significance for matrices and dataframes Simple and exact expressions
\end{Description}
%
\begin{Usage}
\begin{verbatim}
add_significance_symbols(p_value_table, simple = FALSE)
\end{verbatim}
\end{Usage}
%
\begin{Arguments}
\begin{ldescription}
\item[\code{p\_value\_table}] A matrix of p-values, or other numerical matrix

\item[\code{simple}] Whether to represent the * mark simply or precisely with ***.
\end{ldescription}
\end{Arguments}
%
\begin{Value}
Matrices marked with a *.
\end{Value}
%
\begin{Examples}
\begin{ExampleCode}

 ## Not run: 

 #' p_value <- matrix(c(0, 0.00001, 0.00000,
                    0, 0.00000, 0.01497,
                    0, 0.00013, 0.00000,
                    0, 0.00000, 0.00102),
                  nrow = 4, byrow = TRUE,
                  dimnames = list(c("공명음", "마찰음", "유기음_경음", "평파열음_평파찰음"),
                                  c("H", "H(H)", "L")))
     # Add significance symbols
     # p_value

     add_significance_symbols(p_value)

 ##                  H     H(H)  L
 ##공명음            "***" "***" "***"
 ##마찰음            "***" "***" "*"
 ##유기음_경음       "***" "***" "***"
 ##평파열음_평파찰음 "***" "***" "**"




 
## End(Not run)


\end{ExampleCode}
\end{Examples}
\inputencoding{utf8}
\HeaderA{add\_sig\_star}{Functions that output statistical significance for matrices and dataframes Simple and exact expressions}{add.Rul.sig.Rul.star}
%
\begin{Description}
Functions that output statistical significance for matrices and dataframes Simple and exact expressions
\end{Description}
%
\begin{Usage}
\begin{verbatim}
add_sig_star(p_value_table, simple = FALSE)
\end{verbatim}
\end{Usage}
%
\begin{Arguments}
\begin{ldescription}
\item[\code{p\_value\_table}] A matrix of p-values, or other numerical matrix

\item[\code{simple}] Whether to represent the * mark simply or precisely with ***.
\end{ldescription}
\end{Arguments}
%
\begin{Value}
Matrices marked with a *.
\end{Value}
%
\begin{Examples}
\begin{ExampleCode}

 ## Not run: 

 #' p_value <- matrix(c(0, 0.00001, 0.00000,
                    0, 0.00000, 0.01497,
                    0, 0.00013, 0.00000,
                    0, 0.00000, 0.00102),
                  nrow = 4, byrow = TRUE,
                  dimnames = list(c("공명음", "마찰음", "유기음_경음", "평파열음_평파찰음"),
                                  c("H", "H(H)", "L")))
     # Add significance symbols
     # p_value

     add_significance_symbols(p_value)

 ##                  H     H(H)  L
 ##공명음            "***" "***" "***"
 ##마찰음            "***" "***" "*"
 ##유기음_경음       "***" "***" "***"
 ##평파열음_평파찰음 "***" "***" "**"




 
## End(Not run)


\end{ExampleCode}
\end{Examples}
\inputencoding{utf8}
\HeaderA{add\_stat}{데이터에서변수를 택하면 함수에 맞는 새로운 열을 추가하는 함수}{add.Rul.stat}
%
\begin{Description}
데이터에서변수를 택하면 함수에 맞는 새로운 열을 추가하는 함수
\end{Description}
%
\begin{Usage}
\begin{verbatim}
add_stat(
  data,
  ...,
  fun = mean,
  col_name = NULL,
  na.rm = TRUE,
  margin = 1,
  type = "data"
)
\end{verbatim}
\end{Usage}
%
\begin{Arguments}
\begin{ldescription}
\item[\code{data}] data.frame

\item[\code{...}] cols

\item[\code{fun}] mean, ad, ....

\item[\code{col\_name}] New variable

\item[\code{na.rm}] na.rm=TRUE

\item[\code{margin}] margin 1, 2

\item[\code{type}] data, col 2type
\end{ldescription}
\end{Arguments}
%
\begin{Value}
data.fram or vector
\end{Value}
%
\begin{Examples}
\begin{ExampleCode}

## Not run: 

df = data.frame(A = c(1, 2, 3, 4, 5), B = c(6, 7, NA, 9, 10), C = c(11, 12, 13, 14, 15) )

df%>% add_stat( "AC_MEAN","A", "C", fun = mean)
df%>% add_stat( "AC_MEAN","A", "C", fun = mean, type="col")
df %>%add_stat( "A", "C", type="col")
df %>%add_stat( "A", "C")
df %>%add_stat( "A", "C", fun = mean, "AC_MEAN")

df %>%add_stat( "A", "C", "AC_MEAN",fun = sd)

df %>%add_stat(col_name = "AC_MEAN","A", "C", fun = mean)
df %>%add_stat("A", "C", fun = mean, col_name = "AC_MEAN")
df %>%add_stat("A", "C", fun = mean)
df %>%add_stat(c("A", "C"))
df %>%add_stat("ROWs", c("A", "C"))

mtcars %>% add_stat("MEAN", "am","vs")
mtcars %>% add_stat("am","vs")
mtcars %>% add_stat("MEAN", c("am","vs"))
mtcars %>% add_stat("MEAN", c("mpg","qsec"))
mtcars %>% add_stat("MEAN", c("mpg","qsec"), col_name="ACME")
df


## End(Not run)


\end{ExampleCode}
\end{Examples}
\inputencoding{utf8}
\HeaderA{add\_sum}{Functions to sum margins after adding ratios}{add.Rul.sum}
%
\begin{Description}
Functions to sum margins after adding ratios
\end{Description}
%
\begin{Usage}
\begin{verbatim}
add_sum(
  dataset,
  v1 = "a1",
  v2 = "성조",
  raw = FALSE,
  add_ratio = TRUE,
  plot = FALSE,
  color = TRUE,
  ylab = "onset",
  xlab = "speaker",
  cex = 1.2,
  title = NULL,
  type = "df"
)
\end{verbatim}
\end{Usage}
%
\begin{Arguments}
\begin{ldescription}
\item[\code{dataset}] data(rwa=T)or table

\item[\code{v1}] col

\item[\code{v2}] row

\item[\code{raw}] TRUE dataset. not table

\item[\code{add\_ratio}] using add\_ration function

\item[\code{plot}] mosicplot

\item[\code{color}] color 1:4, etc

\item[\code{ylab}] yalb

\item[\code{xlab}] xlab

\item[\code{cex}] size

\item[\code{title}] title

\item[\code{type}] all, df, res
\end{ldescription}
\end{Arguments}
%
\begin{Value}
matrix
\end{Value}
%
\begin{Examples}
\begin{ExampleCode}

#'
## Not run: 

#' #'mat = matrix(c(36, 67, 11,
          0, 35, 44,
         41, 108, 1,
          56, 87, 54),
           nrow = 4, byrow = TRUE.
   dimnames = list(c("공명음", "마찰음",
     "유기음_경음", "평파열음_평파찰음"),  c("H", "H(H)", "L")))
## step1
  mat %>% add_ratio()
##conbind
 mat %>% add_ratio_df()

##add margins
 mat %>% add_sum()

 #                           H H(H)+H(H)_1         L SUM
 #공명음             26(11.5%)   35(19.4%) 36(28.1%)  97
 #마찰음             49(21.6%)   62(34.4%) 35(27.3%) 146
 #유기음_경음       108(47.6%)   82(45.6%)  11(8.6%) 201
 #평파열음_평파찰음  44(19.4%)     1(0.6%) 46(35.9%)  91
 #SUM                      227         180       128 535







## End(Not run)

\end{ExampleCode}
\end{Examples}
\inputencoding{utf8}
\HeaderA{add\_t\_sig}{add\_t\_sig}{add.Rul.t.Rul.sig}
%
\begin{Description}
add\_t\_sig
\end{Description}
%
\begin{Usage}
\begin{verbatim}
add_t_sig(
  data,
  est,
  se,
  col = ncol(data) + 1,
  unite = FALSE,
  digits = 3,
  ns = " ns"
)
\end{verbatim}
\end{Usage}
%
\begin{Arguments}
\begin{ldescription}
\item[\code{data}] data

\item[\code{col}] position t value

\item[\code{unite}] unite sig star +t

\item[\code{digits}] 3

\item[\code{ns}] ns is ns sig
\end{ldescription}
\end{Arguments}
%
\begin{Value}
data
\end{Value}
%
\begin{Examples}
\begin{ExampleCode}

## Not run: 
#'
New_df = data.frame(rownames = c('IMAG -> EXPE', 'IMAG -> QUAL', 'IMAG -> VAL', 'IMAG -> SAT'),
  Original = c(0.578959128452789, 0.491156739303251, 0.393411171344413, 0.490971137879492),
 Mean.Boot = c(0.584251828061525, 0.495822941940513, 0.400070801457516, 0.5022633230577),
 Std.Error = c(0.0518709126012371, 0.0510037451930923, 0.0528454527326974, 0.0681551639743143),
 perc.025 = c(0.477709787839589, 0.391876781940297, 0.308196019466015, 0.369649328752085),
  perc.975 = c(0.664779330969341, 0.58560953201285, 0.495717252789497, 0.614490783702303)
)
#Auto-create a separate T without creating a SIG and combine it with the SIG to move it
New_df %>% add_t_sig(est=3, se=4, col=4, unite=T)


# Move only sigs after creating a t
New_df %>% add_t_sig(est=3, se=4, col=4, unite=F)

# 열을 결합
New_df %>% add_t_sig(est=3, se=4, col=4, unite=F) %>%Unite(2,4)

New_df %>% add_t_sig(est=3, se=4, col=4, unite=F)%>%Round(3) %>%Unite(2,4)


## End(Not run)
\end{ExampleCode}
\end{Examples}
\inputencoding{utf8}
\HeaderA{all\_na\_zero}{all\_na\_zero}{all.Rul.na.Rul.zero}
%
\begin{Description}
all\_na\_zero
\end{Description}
%
\begin{Usage}
\begin{verbatim}
all_na_zero(df, imp = 0)
\end{verbatim}
\end{Usage}
%
\begin{Arguments}
\begin{ldescription}
\item[\code{df}] data.frame

\item[\code{imp}] implement default = 0
\end{ldescription}
\end{Arguments}
%
\begin{Value}
NA to 0
\end{Value}
\inputencoding{utf8}
\HeaderA{all\_na\_zero1}{NA to zero}{all.Rul.na.Rul.zero1}
%
\begin{Description}
NA to zero
\end{Description}
%
\begin{Usage}
\begin{verbatim}
all_na_zero1(data, fun = "extend")
\end{verbatim}
\end{Usage}
%
\begin{Arguments}
\begin{ldescription}
\item[\code{data}] data.frame

\item[\code{fun}] 'extend', 'base' are equal method
\end{ldescription}
\end{Arguments}
%
\begin{Value}
Na data are changed zero
\end{Value}
\inputencoding{utf8}
\HeaderA{alpha\_apa}{alpha\_apa report}{alpha.Rul.apa}
%
\begin{Description}
alpha\_apa report
\end{Description}
%
\begin{Usage}
\begin{verbatim}
alpha_apa(data)
\end{verbatim}
\end{Usage}
%
\begin{Arguments}
\begin{ldescription}
\item[\code{data}] 분석결과
\end{ldescription}
\end{Arguments}
%
\begin{Value}
apa보고서
\end{Value}
%
\begin{Examples}
\begin{ExampleCode}
## Not run: 

#'
# 진로동기 신뢰도 분석
jjstat::bind_alpha_table(jutr2%>% select(A01:A15),
                         c(1:5), c(6:10), c(11:15)) %>%nice_Var()%>% nice_table()%>%
  replace_df("a","진로정체감") %>%
  replace_df("b","진로통찰력") %>%
  replace_df("c","진로탄력성") %>%
  alpha_apa()


## End(Not run)
\end{ExampleCode}
\end{Examples}
\inputencoding{utf8}
\HeaderA{alpha\_table}{Function to tabulate reliability}{alpha.Rul.table}
%
\begin{Description}
Function to tabulate reliability
\end{Description}
%
\begin{Usage}
\begin{verbatim}
alpha_table(
  alpha_data_res,
  title = "",
  show = "data",
  format = "markdown",
  digits = 3,
  check.keys = TRUE,
  variable = "."
)
\end{verbatim}
\end{Usage}
%
\begin{Arguments}
\begin{ldescription}
\item[\code{alpha\_data\_res}] alpha result

\item[\code{title}] your title

\item[\code{show}] 'data', 'markdown'

\item[\code{format}] markdown, html

\item[\code{digits}] default 3

\item[\code{check.keys}] check.keys=TRUE

\item[\code{variable}] . output
\end{ldescription}
\end{Arguments}
%
\begin{Examples}
\begin{ExampleCode}
## Not run: 

stat_onl %>% select(S_Review, S_Add_learn, S_Feedback) %>%
   jjstat::alpha_table(show="data", variable = "A" )


bind_rows(
 stat_onl %>% select(S_Review, S_Add_learn, S_Feedback) %>%
   jjstat::alpha_table(show="data", variable = "A" ),
 stat_onl %>% select(SE_Place, SE_Time) %>%
   jjstat::alpha_table(show="data", variable = "B"),
 stat_onl %>% select(On_Satisfy, On_Joy, On_Easy, On_Engage) %>%
   jjstat::alpha_table(show="data", variable = "C"),
 stat_onl %>% select(upgrade,satisfy,fatigue_inverse) %>%
   jjstat::alpha_table(show="data", variable = "D")
) |> markdown_table(digits=3, font_size = 12,
                   caption= "total Cronbach alpha")

## End(Not run)




\end{ExampleCode}
\end{Examples}
\inputencoding{utf8}
\HeaderA{anova\_apa}{Report aov data (format APA)}{anova.Rul.apa}
%
\begin{Description}
Report aov data (format APA)
\end{Description}
%
\begin{Usage}
\begin{verbatim}
anova_apa(
  aov_post_data,
  ns = 7,
  np = 6,
  ndf = 4,
  nt = 5,
  digits = 3,
  md = FALSE,
  posthoc_type = "The Multi Least Significant Difference(MLSD)"
)
\end{verbatim}
\end{Usage}
%
\begin{Arguments}
\begin{ldescription}
\item[\code{aov\_post\_data}] aov\_post data

\item[\code{ns}] ncol sig column 7

\item[\code{np}] ncol p value column 6

\item[\code{ndf}] ncol df column 4

\item[\code{nt}] ncol t value column 5

\item[\code{digits}] digits 3

\item[\code{md}] md = TRUE is markdown outpput
\end{ldescription}
\end{Arguments}
%
\begin{Value}
report data
\end{Value}
%
\begin{Examples}
\begin{ExampleCode}
## Not run: 
aov(mpg ~ cyl , data=Mtcars)%>% aov_post() %>% anova_apa()
# This is error
aov(mpg ~ cyl , data=Mtcars)%>%aov_post(type="contrast")%>% anova_apa()

# aov(mpg ~ cyl , data=Mtcars)%>%aov_post(type="all")

aov(mpg ~ cyl , data=Mtcars)%>%aov_post(type="all")%>% anova_apa()
aov(mpg ~ cyl , data=Mtcars)%>%aov_post(type="res")%>% anova_apa()


## End(Not run)

\end{ExampleCode}
\end{Examples}
\inputencoding{utf8}
\HeaderA{aov\_apa}{Report aov data (format APA)}{aov.Rul.apa}
%
\begin{Description}
Report aov data (format APA)
\end{Description}
%
\begin{Usage}
\begin{verbatim}
aov_apa(
  aov_post_data,
  ns = 7,
  np = 6,
  ndf = 4,
  nt = 5,
  digits = 3,
  md = FALSE,
  posthoc_type = "The Multi Least Significant Difference(MLSD)"
)
\end{verbatim}
\end{Usage}
%
\begin{Arguments}
\begin{ldescription}
\item[\code{aov\_post\_data}] aov\_post data

\item[\code{ns}] ncol sig column 7

\item[\code{np}] ncol p value column 6

\item[\code{ndf}] ncol df column 4

\item[\code{nt}] ncol t value column 5

\item[\code{digits}] digits 3

\item[\code{md}] md = TRUE is markdown outpput
\end{ldescription}
\end{Arguments}
%
\begin{Value}
report data
\end{Value}
%
\begin{Examples}
\begin{ExampleCode}
## Not run: 
aov(mpg ~ cyl , data=Mtcars)%>%aov_post()%>% aov_apa()
# This is error
aov(mpg ~ cyl , data=Mtcars)%>%aov_post(type="contrast")%>% aov_apa()

# aov(mpg ~ cyl , data=Mtcars)%>%aov_post(type="all")

aov(mpg ~ cyl , data=Mtcars)%>%aov_post(type="all")%>% aov_apa()
aov(mpg ~ cyl , data=Mtcars)%>%aov_post(type="res")%>% aov_apa()


## End(Not run)

\end{ExampleCode}
\end{Examples}
\inputencoding{utf8}
\HeaderA{aov\_df}{\#' Generate an AOV results table for multiple independent variables}{aov.Rul.df}
%
\begin{Description}
\#' Generate an AOV results table for multiple independent variables
\end{Description}
%
\begin{Usage}
\begin{verbatim}
aov_df(
  data,
  dv_var = NULL,
  iv_vars = NULL,
  grp_mean = TRUE,
  unite_p = FALSE,
  unite_F = FALSE,
  digits = 2,
  posthoc = "lsd",
  p.adj = "bonferroni",
  sig = FALSE
)
\end{verbatim}
\end{Usage}
%
\begin{Arguments}
\begin{ldescription}
\item[\code{data}] data.frame

\item[\code{dv\_var}] dv

\item[\code{iv\_vars}] c(iv1, iv2....)

\item[\code{grp\_mean}] 'grooup mean remove

\item[\code{unite\_p}] pvalue unite

\item[\code{unite\_F}] Fvalue unite

\item[\code{posthoc}] 'tukey, scheffe, duncan, lsd

\item[\code{p.adj}] p.adj=c("none","holm","hommel", "hochberg", "bonferroni", "BH", "BY", "fdr") only lsd test

\item[\code{sig}] 'sig = T' is add star
\end{ldescription}
\end{Arguments}
%
\begin{Value}
table
\end{Value}
%
\begin{Examples}
\begin{ExampleCode}

## Not run: 
aov_table(data = mtcars, dv_var = "mpg",
iv_vars = c("cyl", "gear", "carb"))

aov_table(data = mtcars, dv_var = "mpg",
iv_vars = c("cyl", "gear", "carb"),
           mean = TRUE)

aov_table(data = mtcars, dv_var = "mpg",
iv_vars = c("cyl", "gear", "carb"), grp_mean = FALSE)
aov_table(data = mtcars, dv_var = "mpg",
iv_vars = c("cyl", "gear", "carb"), grp_mean = FALSE) %>%


aov_table(data = mtcars, dv_var = "mpg",
iv_vars = c("cyl", "gear", "carb"),
 grp_mean = TRUE)

aov_table(data = mtcars, dv_var = "mpg",
 iv_vars = c("cyl", "gear", "carb"),
 grp_mean = TRUE, unite_F = TRUE)

aov_table(data = mtcars, dv_var = "mpg",
iv_vars = c("cyl", "gear", "carb"),
 grp_mean = TRUE, unite_F = TRUE) %>%
  markdown_table()

# park516d데이터 이용
park516d %>%
  select(1:6,teaching, academyType) %>%  #주요변수 선택
  bind_add_stat(c("RV_Teah",1:6)) %>% # 평균을 이용하여 항목묶음
  aov_df(dv_var="RV_Teah",
         iv_vars =c("teaching", "academyType"), sig = T)



## End(Not run)


\end{ExampleCode}
\end{Examples}
\inputencoding{utf8}
\HeaderA{aov\_df2}{Perform ANOVA with emmeans and extract group letters and p-values}{aov.Rul.df2}
%
\begin{Description}
This function performs ANOVA for each independent variable (IV) and extracts
the estimated marginal means (emmeans), post-hoc comparison group letters, and
p-values, based on the specified multiple comparisons adjustment method.
\end{Description}
%
\begin{Usage}
\begin{verbatim}
aov_df2(
  data,
  dv_var,
  iv_vars,
  adjust = "none",
  alpha = 0.05,
  digits = 2,
  show_group_letter = TRUE,
  adjust_only = TRUE
)
\end{verbatim}
\end{Usage}
%
\begin{Arguments}
\begin{ldescription}
\item[\code{data}] A data.frame containing the variables.

\item[\code{dv\_var}] The name of the dependent variable (character).

\item[\code{iv\_vars}] A character vector of independent variable names.

\item[\code{adjust}] Adjustment method for pairwise comparisons. One of:
"none", "tukey", "bonferroni", "holm", "sidak", "fdr", "scheffe".

\item[\code{alpha}] Significance level (default = 0.05).

\item[\code{digits}] Number of digits to round numeric results (default = 2).

\item[\code{show\_group\_letter}] Logical. Whether to show group letters (default = TRUE).

\item[\code{adjust\_only}] Logical. If TRUE and adjust != "none", only adjusted p-values (adj\_p) are shown (default = TRUE).
\end{ldescription}
\end{Arguments}
%
\begin{Details}
\#\# Adjustment Methods
- `"none"`: No adjustment for multiple comparisons (raw p-values). ⚠️ Can inflate Type I error.
- `"tukey"`: Tukey's Honest Significant Difference method. Simultaneously compares all pairwise mean differences.
- `"bonferroni"`: Conservative method; adjusts p-values by multiplying by the number of comparisons.
- `"holm"`: Sequential Bonferroni procedure; less conservative and more powerful alternative.
- `"sidak"`: Adjustment assuming independent comparisons.
- `"fdr"`: Controls the False Discovery Rate (Benjamini-Hochberg procedure).
- `"scheffe"`: Not supported in `emmeans`; typically used for general contrasts.
- `"duncan"`, `"lsd"`: Not supported in `emmeans`; internally treated as `"none"`.
\end{Details}
%
\begin{Value}
A tibble with ANOVA results per IV level, group letters, and p-values.
\end{Value}
%
\begin{Examples}
\begin{ExampleCode}
## Not run: 
# Basic usage
aov_df_emmeans(data = neulbom_p2_stat2,
               dv_var = "relation",
               iv_vars = c("gender", "age", "child"))

# Using Bonferroni adjustment
aov_df_emmeans(data = neulbom_p2_stat2,
               dv_var = "relation",
               iv_vars = c("gender", "age", "child"),
               adjust = "bonferroni", adjust_only = FALSE)

# Showing only adjusted p-values
aov_df_emmeans(data = neulbom_p2_stat2,
               dv_var = "relation",
               iv_vars = c("gender", "age", "child"),
               adjust = "holm", adjust_only = TRUE)

# Tukey fallback warning
aov_df_emmeans(data = neulbom_p2_stat2,
               dv_var = "relation",
               iv_vars = c("gender", "age", "child"),
               adjust = "tukey")

# Unsupported methods fallback
aov_df_emmeans(data = neulbom_p2_stat2,
               dv_var = "relation",
               iv_vars = c("gender", "age", "child"),
               adjust = "duncan")

## End(Not run)
\end{ExampleCode}
\end{Examples}
\inputencoding{utf8}
\HeaderA{aov\_df\_emmeans}{Perform ANOVA with emmeans and extract group letters and p-values}{aov.Rul.df.Rul.emmeans}
%
\begin{Description}
This function performs ANOVA for each independent variable (IV) and extracts
the estimated marginal means (emmeans), post-hoc comparison group letters, and
p-values, based on the specified multiple comparisons adjustment method.
\end{Description}
%
\begin{Usage}
\begin{verbatim}
aov_df_emmeans(
  data,
  dv_var,
  iv_vars,
  adjust = "none",
  alpha = 0.05,
  digits = 2,
  show_group_letter = TRUE,
  adjust_only = TRUE
)
\end{verbatim}
\end{Usage}
%
\begin{Arguments}
\begin{ldescription}
\item[\code{data}] A data.frame containing the variables.

\item[\code{dv\_var}] The name of the dependent variable (character).

\item[\code{iv\_vars}] A character vector of independent variable names.

\item[\code{adjust}] Adjustment method for pairwise comparisons. One of:
"none", "tukey", "bonferroni", "holm", "sidak", "fdr", "scheffe".

\item[\code{alpha}] Significance level (default = 0.05).

\item[\code{digits}] Number of digits to round numeric results (default = 2).

\item[\code{show\_group\_letter}] Logical. Whether to show group letters (default = TRUE).

\item[\code{adjust\_only}] Logical. If TRUE and adjust != "none", only adjusted p-values (adj\_p) are shown (default = TRUE).
\end{ldescription}
\end{Arguments}
%
\begin{Details}
\#\# Adjustment Methods
- `"none"`: No adjustment for multiple comparisons (raw p-values). ⚠️ Can inflate Type I error.
- `"tukey"`: Tukey's Honest Significant Difference method. Simultaneously compares all pairwise mean differences.
- `"bonferroni"`: Conservative method; adjusts p-values by multiplying by the number of comparisons.
- `"holm"`: Sequential Bonferroni procedure; less conservative and more powerful alternative.
- `"sidak"`: Adjustment assuming independent comparisons.
- `"fdr"`: Controls the False Discovery Rate (Benjamini-Hochberg procedure).
- `"scheffe"`: Not supported in `emmeans`; typically used for general contrasts.
- `"duncan"`, `"lsd"`: Not supported in `emmeans`; internally treated as `"none"`.
\end{Details}
%
\begin{Value}
A tibble with ANOVA results per IV level, group letters, and p-values.
\end{Value}
%
\begin{Examples}
\begin{ExampleCode}
## Not run: 
# Basic usage
aov_df_emmeans(data = neulbom_p2_stat2,
               dv_var = "relation",
               iv_vars = c("gender", "age", "child"))

# Using Bonferroni adjustment
aov_df_emmeans(data = neulbom_p2_stat2,
               dv_var = "relation",
               iv_vars = c("gender", "age", "child"),
               adjust = "bonferroni", adjust_only = FALSE)

# Showing only adjusted p-values
aov_df_emmeans(data = neulbom_p2_stat2,
               dv_var = "relation",
               iv_vars = c("gender", "age", "child"),
               adjust = "holm", adjust_only = TRUE)

# Tukey fallback warning
aov_df_emmeans(data = neulbom_p2_stat2,
               dv_var = "relation",
               iv_vars = c("gender", "age", "child"),
               adjust = "tukey")

# Unsupported methods fallback
aov_df_emmeans(data = neulbom_p2_stat2,
               dv_var = "relation",
               iv_vars = c("gender", "age", "child"),
               adjust = "duncan")

## End(Not run)
\end{ExampleCode}
\end{Examples}
\inputencoding{utf8}
\HeaderA{aov\_post}{ANOVA posthoc and plot, Contrast data}{aov.Rul.post}
%
\begin{Description}
ANOVA posthoc and plot, Contrast data
\end{Description}
%
\begin{Usage}
\begin{verbatim}
aov_post(aov_data, type = "df", adjust = "none", lsd = FALSE)
\end{verbatim}
\end{Usage}
%
\begin{Arguments}
\begin{ldescription}
\item[\code{aov\_data}] aov data

\item[\code{type}] type is all, res, contras, df, raq(plotdata), g(plot), aov, anova

\item[\code{adjust}] adjust ="none", bonferroni, sidak etc

\item[\code{lsd}] plosthoc LSD
\end{ldescription}
\end{Arguments}
%
\begin{Value}
contrast and g
\end{Value}
%
\begin{Examples}
\begin{ExampleCode}

## Not run: 
Mtcars <- mtcars%>%as_trt("cyl")
aov(mpg ~ cyl , data=Mtcars)%>% aov_post()
aov(mpg ~ cyl , data=Mtcars)%>% aov_post(type="all")
aov(mpg ~ cyl , data=Mtcars)%>% aov_post(type="res")
aov(mpg ~ cyl , data=Mtcars)%>% aov_post(type="df")
aov(mpg ~ cyl , data=Mtcars)%>% aov_post(type="call")
aov(mpg ~ cyl , data=Mtcars)%>% aov_post(type="formula")
aov(mpg ~ cyl , data=Mtcars)%>% aov_post(type="g")
aov(mpg ~ cyl , data=Mtcars)%>% aov_post(type="contrast")
aov(mpg ~ cyl , data=Mtcars)%>% aov_post(type="aov")
aov(mpg ~ cyl , data=Mtcars)%>% aov_post(type="anova")
aov(mpg ~ cyl , data=Mtcars)%>% aov_post(type="posthoc", lsd=TRUE)

#report
aov(mpg ~ cyl , data=Mtcars)%>%aov_post()%>% anova_apa()

## End(Not run)


\end{ExampleCode}
\end{Examples}
\inputencoding{utf8}
\HeaderA{aov\_post2}{ANOVA posthoc and plot, Contrast data}{aov.Rul.post2}
%
\begin{Description}
Post-hoc analysis and plotting for ANOVA results. Supports various posthoc methods
including LSD, Duncan, Scheffe, Tukey HSD, SNK, FDR and Bonferroni correction.
\end{Description}
%
\begin{Usage}
\begin{verbatim}
aov_post2(
  aov_data,
  type = "df",
  adjust = "none",
  posthoc = TRUE,
  posthoc_method = "LSD"
)
\end{verbatim}
\end{Usage}
%
\begin{Arguments}
\begin{ldescription}
\item[\code{aov\_data}] An object returned by \code{aov()}; ANOVA model object.

\item[\code{type}] Output type. One of: \code{"all"}, \code{"res"}, \code{"contrast"},
\code{"df"}, \code{"raw"}, \code{"g"}, \code{"aov"}, \code{"anova"}, \code{"call"},
\code{"formula"}, \code{"posthoc"}.

\item[\code{adjust}] Method for p-value adjustment for emmeans contrast. One of \code{"none"},
\code{"bonferroni"}, \code{"fdr"}, etc. Only used in contrast.

\item[\code{posthoc}] Logical; if \code{TRUE}, performs posthoc test using agricolae.

\item[\code{posthoc\_method}] Posthoc method: one of \code{"LSD"}, \code{"BONFERRONI"},
\code{"FDR"}, \code{"SCHEFFE"}, \code{"DUNCAN"}, \code{"TUKEY"}, \code{"SNK"}.
\end{ldescription}
\end{Arguments}
%
\begin{Value}
Depending on \code{type}, returns a list or specific object.
\end{Value}
%
\begin{Examples}
\begin{ExampleCode}
## Not run: 
library(agricolae)
library(dplyr)
library(emmeans)
library(broom)
library(ggplot2)

# Example 1: mtcars dataset (mpg ~ cyl)
mtcars$cyl <- as.factor(mtcars$cyl)
fit1 <- aov(mpg ~ cyl, data = mtcars)
aov_post(fit1, type = "df", posthoc = TRUE, posthoc_method = "LSD")
aov_post(fit1, type = "df", posthoc = TRUE, posthoc_method = "BONFERRONI")
aov_post(fit1, type = "df", posthoc = TRUE, posthoc_method = "FDR")
aov_post(fit1, type = "df", posthoc = TRUE, posthoc_method = "SCHEFFE")
aov_post(fit1, type = "df", posthoc = TRUE, posthoc_method = "DUNCAN")
aov_post(fit1, type = "df", posthoc = TRUE, posthoc_method = "TUKEY")
aov_post(fit1, type = "df", posthoc = TRUE, posthoc_method = "SNK")

# Example 2: ToothGrowth dataset (len ~ dose)
ToothGrowth$dose <- as.factor(ToothGrowth$dose)
fit2 <- aov(len ~ dose, data = ToothGrowth)
aov_post(fit2, type = "df", posthoc = TRUE, posthoc_method = "DUNCAN")

# Example 3: iris dataset (Sepal.Length ~ Species)
fit3 <- aov(Sepal.Length ~ Species, data = iris)
aov_post(fit3, type = "df", posthoc = TRUE, posthoc_method = "SCHEFFE")

## End(Not run)
\end{ExampleCode}
\end{Examples}
\inputencoding{utf8}
\HeaderA{aov\_rename}{aov\_rename}{aov.Rul.rename}
%
\begin{Description}
aov\_rename

aov\_rename
\end{Description}
%
\begin{Usage}
\begin{verbatim}
aov_rename(data)

aov_rename(data)
\end{verbatim}
\end{Usage}
%
\begin{Arguments}
\begin{ldescription}
\item[\code{data}] data
\end{ldescription}
\end{Arguments}
%
\begin{Value}
result

result
\end{Value}
\inputencoding{utf8}
\HeaderA{aov\_table}{Generate an AOV results table for multiple independent variables}{aov.Rul.table}
%
\begin{Description}
Generate an AOV results table for multiple independent variables

Generate an AOV results table for multiple independent variables
\end{Description}
%
\begin{Usage}
\begin{verbatim}
aov_table(
  data,
  dv_var = NULL,
  iv_vars = NULL,
  grp_mean = TRUE,
  unite_p = FALSE,
  unite_F = FALSE,
  digits = 2,
  posthoc = "lsd",
  p.adj = "bonferroni",
  sig = FALSE
)

aov_table(
  data,
  dv_var = NULL,
  iv_vars = NULL,
  grp_mean = TRUE,
  unite_p = FALSE,
  unite_F = FALSE,
  digits = 2,
  posthoc = "lsd",
  p.adj = "bonferroni",
  sig = FALSE
)
\end{verbatim}
\end{Usage}
%
\begin{Arguments}
\begin{ldescription}
\item[\code{data}] data.frame

\item[\code{dv\_var}] dv

\item[\code{iv\_vars}] c(iv1, iv2....)

\item[\code{grp\_mean}] 'grooup mean remove

\item[\code{unite\_p}] pvalue unite

\item[\code{unite\_F}] Fvalue unite

\item[\code{posthoc}] 'tukey, scheffe, duncan, lsd

\item[\code{p.adj}] p.adj=c("none","holm","hommel", "hochberg", "bonferroni", "BH", "BY", "fdr") only lsd test

\item[\code{sig}] 'sig = T' is add star
\end{ldescription}
\end{Arguments}
%
\begin{Value}
table

table
\end{Value}
%
\begin{Examples}
\begin{ExampleCode}

## Not run: 
mtcars <- mtcars%>%as_trt("cyl", "gear", "carb")
aov_table(data = mtcars, dv_var = "mpg", iv_vars = c("cyl", "gear", "carb"))

aov_table(data = mtcars, dv_var = "mpg",
iv_vars = c("cyl", "gear", "carb"), grp_mean = FALSE)
aov_table(data = mtcars, dv_var = "mpg",
iv_vars = c("cyl", "gear", "carb"), grp_mean = FALSE) %>%
aov_table_apaOP

aov_table(data = mtcars, dv_var = "mpg",
 iv_vars = c("cyl", "gear", "carb"), grp_mean = TRUE)

aov_table(data = mtcars, dv_var = "mpg",
 iv_vars = c("cyl", "gear", "carb"), grp_mean = TRUE, unite_F = TRUE)

aov_table(data = mtcars, dv_var = "mpg",
 iv_vars = c("cyl", "gear", "carb"),
  grp_mean = TRUE, unite_F = TRUE) %>% aov_rename()
  markdown_table()

## End(Not run)




## Not run: 
mtcars <- mtcars%>%as_trt("cyl", "gear", "carb")
aov_table(data = mtcars, dv_var = "mpg", iv_vars = c("cyl", "gear", "carb"))

aov_table(data = mtcars, dv_var = "mpg",
iv_vars = c("cyl", "gear", "carb"), grp_mean = FALSE)
aov_table(data = mtcars, dv_var = "mpg",
iv_vars = c("cyl", "gear", "carb"), grp_mean = FALSE) %>%
aov_table_apaOP

aov_table(data = mtcars, dv_var = "mpg",
 iv_vars = c("cyl", "gear", "carb"), grp_mean = TRUE)

aov_table(data = mtcars, dv_var = "mpg",
 iv_vars = c("cyl", "gear", "carb"), grp_mean = TRUE, unite_F = TRUE)

aov_table(data = mtcars, dv_var = "mpg",
 iv_vars = c("cyl", "gear", "carb"),
  grp_mean = TRUE, unite_F = TRUE) %>% aov_rename()
  markdown_table()

## End(Not run)



\end{ExampleCode}
\end{Examples}
\inputencoding{utf8}
\HeaderA{apa\_aov\_df}{total anova report}{apa.Rul.aov.Rul.df}
%
\begin{Description}
total anova report
\end{Description}
%
\begin{Usage}
\begin{verbatim}
apa_aov_df(data, show = TRUE, dv_var = NULL, iv_vars = NULL)
\end{verbatim}
\end{Usage}
%
\begin{Arguments}
\begin{ldescription}
\item[\code{data}] aoa\_df(grp\_mean=F)

\item[\code{show}] show data
\end{ldescription}
\end{Arguments}
%
\begin{Value}
report
\end{Value}
%
\begin{Examples}
\begin{ExampleCode}
## Not run: 
aov_table(data = mtcars, dv_var = "mpg",
iv_vars = c("cyl", "gear", "carb"),
grp_mean = FALSE)%>% apa_aov_df()

apa_aov_df(data = mtcars, dv_var = "mpg",
 iv_vars = c("cyl", "gear", "carb"))

## End(Not run)


\end{ExampleCode}
\end{Examples}
\inputencoding{utf8}
\HeaderA{apa\_aov\_df\_res}{apa\_aov\_df\_res, aov\_df의 분석결과}{apa.Rul.aov.Rul.df.Rul.res}
%
\begin{Description}
apa\_aov\_df\_res, aov\_df의 분석결과
\end{Description}
%
\begin{Usage}
\begin{verbatim}
apa_aov_df_res(data, dv_col = 2, iv_col = 1, show = TRUE, digits = 3)
\end{verbatim}
\end{Usage}
%
\begin{Arguments}
\begin{ldescription}
\item[\code{data}] data

\item[\code{dv\_col}] dv

\item[\code{iv\_col}] iv

\item[\code{show}] TRUE, data result

\item[\code{digits}] 3
\end{ldescription}
\end{Arguments}
%
\begin{Value}
결과해석 text
\end{Value}
%
\begin{Examples}
\begin{ExampleCode}
## Not run: 
#'
# 테스트 데이터로 실행
data <- data.frame(
  독립변수 = c("LCA_class", "", "", "LCA_class",
  "", "", "LCA_class", "", "", "LCA_class", "", ""),
  종속변수 = c("진로동기", "진로동기", "진로동기",
  "진로태도성숙", "진로태도성숙", "진로태도성숙",
  "진로결정자기효능감", "진로결정자기효능감", "진로결정자기효능감",
  "진로준비행동", "진로준비행동", "진로준비행동"),
  df1 = c(2, 2, 2, 2, 2, 2, 2, 2, 2, 2, 2, 2),
  df2 = c(511, 511, 511, 511, 511, 511,
   511, 511, 511, 511, 511, 511),
  F = c(132.64, NA, NA, 107.45, NA, NA,
   138.02, NA, NA, 106.1, NA, NA),
  p = c("4.00665784818538e-47", NA, NA,
   "1.11197564656087e-39", NA, NA,
   "1.18550277068589e-48", NA, NA,
   "2.8829281865653e-39", NA, NA)
)

# show=TRUE로 결과 출력
apa_results <- apa_aov_df_res(data, dv_col = 2, iv_col = 1, show = TRUE)


## End(Not run)
\end{ExampleCode}
\end{Examples}
\inputencoding{utf8}
\HeaderA{aroian\_sobel\_test\_general}{Aroian Approximation Sobel Test generalized function}{aroian.Rul.sobel.Rul.test.Rul.general}
%
\begin{Description}
Aroian Approximation Sobel Test generalized function
\end{Description}
%
\begin{Usage}
\begin{verbatim}
aroian_sobel_test_general(
  a,
  SE,
  roburst = FALSE,
  n = NULL,
  k = 2,
  digits = 5,
  conf = 0.95
)
\end{verbatim}
\end{Usage}
%
\begin{Arguments}
\begin{ldescription}
\item[\code{a}] est

\item[\code{SE}] se

\item[\code{roburst}] roburst T whide correction

\item[\code{n}] sample size

\item[\code{k}] number of regressin coefs, default 2

\item[\code{digits}] 5

\item[\code{conf}] 0.90, 0.95, 0.99, 0.999
\end{ldescription}
\end{Arguments}
%
\begin{Value}
result
\end{Value}
%
\begin{Examples}
\begin{ExampleCode}

## Not run: 
#'
# Example usage
a <- 0.5  # Example coefficient a
b <- 0.3  # Example coefficient b
sa <- 0.1 # Standard error of a
sb <- 0.05 # Standard error of b

Sobel_test(a, b, sa, sb)
Sobel_test(a, b, sa, sb, "all")
Sobel_test(a, b, sa, sb, "Sobel")

Sobel_test(a, b, sa, sb,"Aroian")

Sobel_test(a, b, sa, sb,"Goodman")

#'
# 예제 데이터
a <- c(0.5, 0.3, 0.4)
SE <- c(0.1, 0.1, 0.1)

# 함수 실행 (기본 95% 신뢰구간)
aroian_sobel_test_general(a, SE)
print(result_95)

# 함수 실행 (90% 신뢰구간)
result_90 <- aroian_sobel_test_general(a, SE, conf=0.90)
print(result_90)

# 함수 실행 (99% 신뢰구간)
result_99 <- aroian_sobel_test_general(a, SE, conf=0.99)
print(result_99)

# 함수 실행 (99.9% 신뢰구간)
result_999 <- aroian_sobel_test_general(a, SE, conf=0.999)
print(result_999)


aroian_sobel_test_general(a = c(0.2, 0.3, 0.4, 0.5),
                          SE = c(0.05, 0.04, 0.03, 0.02))%>%unite_ci()


## End(Not run)
\end{ExampleCode}
\end{Examples}
\inputencoding{utf8}
\HeaderA{arrange\_by\_mind}{Arrange dataframe by user-defined priority on rhs column}{arrange.Rul.by.Rul.mind}
%
\begin{Description}
This function arranges a data frame according to a user-defined
priority of values in the \code{rhs} column. If no priority is supplied,
the data frame is sorted alphabetically by \code{rhs}.
\end{Description}
%
\begin{Usage}
\begin{verbatim}
arrange_by_mind(df, priority = NULL)
\end{verbatim}
\end{Usage}
%
\begin{Arguments}
\begin{ldescription}
\item[\code{df}] A data frame containing at least a column named \code{rhs}.

\item[\code{priority}] An optional character vector specifying the preferred
order of values in the \code{rhs} column. If \code{NULL}, rows are
sorted alphabetically by \code{rhs}.
\end{ldescription}
\end{Arguments}
%
\begin{Details}
When \code{priority} is provided, values of \code{rhs} that appear in
\code{priority} are ordered first according to that order. Remaining
values are placed afterward and sorted alphabetically.
If \code{priority} is \code{NULL}, the function simply sorts the data
frame alphabetically by \code{rhs}.
\end{Details}
%
\begin{Value}
A data frame sorted by user-defined priority and/or
alphabetically by \code{rhs}.
\end{Value}
%
\begin{Examples}
\begin{ExampleCode}
## Not run: 
df <- data.frame(
  lhs = c("intention", "access", "efficacy"),
  rhs = c("barrier", "experience", "grade"),
  est = c(0.2, 0.4, -0.1)
)

# Example 1: Alphabetical ordering only
arrange_by_mind(df)

# Example 2: User-defined priority
arrange_by_mind(df,
  priority = c("experience", "grade")
)

# Example 3: Partial priority specification
arrange_by_mind(df,
  priority = "experience"
)

## End(Not run)
\end{ExampleCode}
\end{Examples}
\inputencoding{utf8}
\HeaderA{as\_trt}{Treat selected variables as factors}{as.Rul.trt}
%
\begin{Description}
Treat selected variables as factors
\end{Description}
%
\begin{Usage}
\begin{verbatim}
as_trt(data, ..., fun = factor)
\end{verbatim}
\end{Usage}
%
\begin{Arguments}
\begin{ldescription}
\item[\code{data}] data.frame

\item[\code{...}] variables

\item[\code{fun}] funtion default as.factor
\end{ldescription}
\end{Arguments}
%
\begin{Value}
data treatment
\end{Value}
%
\begin{Examples}
\begin{ExampleCode}
## Not run: 
 str(mrcars)

 mtcars1 = mtcars %>% sel_as_trt("cyl", "am","vs", "gear", "carb")
 str(mtcars1)

 #col number
 mtcars%>%as_trt(c(2,8:11))%>%str()
#'
# 데이터 프레임 생성
df <- data.frame(matrix(1:500, ncol=100))
colnames(df) <- paste0("P", sprintf("%02d", 1:100))

# 열 번호 범위로 변환
as_trt(df, 81:91)%>%str()
as_trt(df, "P81","P91")%>%str()
as_trt(df, 81, 92)%>%str()
#'



## End(Not run)

\end{ExampleCode}
\end{Examples}
\inputencoding{utf8}
\HeaderA{as\_trt2}{Treat selected variables as factors}{as.Rul.trt2}
%
\begin{Description}
Treat selected variables as factors
\end{Description}
%
\begin{Usage}
\begin{verbatim}
as_trt2(data, ..., fun = factor)
\end{verbatim}
\end{Usage}
%
\begin{Arguments}
\begin{ldescription}
\item[\code{data}] data.frame

\item[\code{...}] variables

\item[\code{fun}] funtion default as.factor
\end{ldescription}
\end{Arguments}
%
\begin{Value}
data treatment
\end{Value}
%
\begin{Examples}
\begin{ExampleCode}
## Not run: 
 str(mrcars)

 mtcars1 = mtcars %>% as_trt("cyl", "am","vs", "gear", "carb")
 str(mtcars1)

 #col number
 mtcars%>%as_trt(c(2,8:11))%>%str()
#'
# 데이터 프레임 생성
df <- data.frame(matrix(1:500, ncol=100))
colnames(df) <- paste0("P", sprintf("%02d", 1:100))

# 열 번호 범위로 변환
as_trt(df, 81:91)%>%str()
as_trt(df, "P81:P91")%>%str()
as_trt(df, "P81","P91")%>%str()
as_trt(df, 81, 92)%>%str()
#'



## End(Not run)

\end{ExampleCode}
\end{Examples}
\inputencoding{utf8}
\HeaderA{auto\_coding}{Automated Likert scale conversion functions}{auto.Rul.coding}
%
\begin{Description}
Automated Likert scale conversion functions
\end{Description}
%
\begin{Usage}
\begin{verbatim}
auto_coding(
  data,
  ...,
  range = 1:ncol(data),
  custom = NULL,
  rev = FALSE,
  data.fram = FALSE,
  start_number = 1,
  default = c("전혀 그렇지 않다", "그렇지 않다", "보통이다", "그렇다",
    "매우 그렇다")
)
\end{verbatim}
\end{Usage}
%
\begin{Arguments}
\begin{ldescription}
\item[\code{data}] data.frame, vector

\item[\code{...}] coding levels

\item[\code{range}] data col range\#'

\item[\code{start\_number}] start\_number = 1

\item[\code{defaultc("전혀}] 그렇지 않다", "그렇지 않다","보통이다", "그렇다", "매우 그렇다")
\end{ldescription}
\end{Arguments}
%
\begin{Value}
transforamtion data
\end{Value}
%
\begin{Examples}
\begin{ExampleCode}

## Not run: 
#'

# test
jut1_custom <- c("그렇다", "매우 그렇다", "보통이다",
                 "그렇지 않다", "전혀 그렇지 않다")
jut1_custom
jut1_custom%>%rev()


jut1_custom <- c("그렇다", "매우 그렇다", "보통이다", "그렇지 않다", "전혀 그렇지 않다")

auto_coding(jut1_custom, custom=jut1_custom)
auto_coding(jut1_custom, custom=jut1_custom, rev=TRUE)


# Example data
vec1 <- c("그렇다", "그렇다", "매우 그렇다", "그렇다", "매우 그렇다", "매우 그렇다",
          "매우 그렇다", "그렇다", "그렇다", "그렇지 않다", "그렇다", "그렇다",
          "매우 그렇다", "그렇다", "그렇다", "그렇다", "그렇다", "매우 그렇다",
          "전혀 그렇지 않다", "보통이다", "그렇다", "매우 그렇다", "매우 그렇다",
          "보통이다", "그렇다", "매우 그렇다", "그렇다", "그렇다", "매우 그렇다",
          "매우 그렇다", "그렇다", "그렇다", "그렇다", "그렇다", "보통이다",
          "그렇다", "그렇다", "그렇다", "매우 그렇다", "보통이다", "매우 그렇다",
          "매우 그렇다", "그렇다", "그렇다", "그렇지 않다", "그렇다", "그렇다",
          "그렇다", "보통이다", "그렇다", "그렇다", "보통이다", "그렇지 않다",
          "그렇지 않다", "매우 그렇다", "보통이다", "그렇다", "매우 그렇다",
          "그렇다", "매우 그렇다", "그렇다", "그렇다", "그렇다", "그렇다",
          "그렇다", "매우 그렇다", "그렇다", "보통이다", "그렇다", "그렇다",
          "그렇다", "매우 그렇다", "보통이다", "그렇다", "그렇다", "그렇지 않다",
          "그렇지 않다", "그렇지 않다", "그렇지 않다", "매우 그렇다", "매우 그렇다",
          "그렇지 않다", "그렇다", "전혀 그렇지 않다", "그렇다", "그렇다",
          "그렇다", "보통이다", "그렇다", "그렇지 않다", "그렇다", "보통이다",
          "매우 그렇다")

auto_coding(vec1, custom=jut1_custom, rev=FALSE)
auto_coding(vec1, custom=jut1_custom, rev=TRUE)



# Example data
data1 <- data.frame(
  V1 = c("그렇다", "그렇다", "매우 그렇다", "그렇다", "매우 그렇다",
         "매우 그렇다", "매우 그렇다", "그렇다", "그렇다", "그렇지 않다",
         "그렇다", "그렇다", "매우 그렇다", "그렇다", "그렇다",
         "그렇다", "그렇다", "매우 그렇다",  "전혀 그렇지 않다", "보통이다",
         "그렇다", "매우 그렇다", "매우 그렇다", "보통이다", "그렇다",
         "매우 그렇다", "그렇다", "그렇다", "매우 그렇다",   "매우 그렇다",
         "그렇다", "그렇다", "그렇다", "그렇다", "보통이다",
         "그렇다", "그렇다", "그렇다", "매우 그렇다", "보통이다",
         "매우 그렇다", "매우 그렇다", "그렇다", "그렇다", "그렇지 않다",
         "그렇다", "그렇다",
         "그렇다", "보통이다", "그렇다", "그렇다", "보통이다", "그렇지 않다",
         "그렇지 않다", "매우 그렇다", "보통이다", "그렇다", "매우 그렇다",
         "그렇다", "매우 그렇다", "그렇다", "그렇다", "그렇다", "그렇다",
         "그렇다", "매우 그렇다", "그렇다", "보통이다", "그렇다", "그렇다",
         "그렇다", "매우 그렇다", "보통이다", "그렇다", "그렇다", "그렇지 않다",
         "그렇지 않다", "그렇지 않다", "그렇지 않다", "매우 그렇다", "매우 그렇다",
         "그렇지 않다", "그렇다", "전혀 그렇지 않다", "그렇다", "그렇다",
         "그렇다", "보통이다", "그렇다", "그렇지 않다", "그렇다", "보통이다",
         "매우 그렇다", "보통이다")
)

auto_coding(data1)
data1

# Example data
data2 <- data.frame(
  V1 = c("그렇다", "그렇다", "매우 그렇다", "그렇다", "매우 그렇다",
         "매우 그렇다", "매우 그렇다", "그렇다", "그렇다", "그렇지 않다",
         "그렇다", "그렇다", "매우 그렇다", "그렇다", "그렇다",
         "그렇다", "그렇다", "매우 그렇다",  "전혀 그렇지 않다", "보통이다",
         "그렇다", "매우 그렇다", "매우 그렇다", "보통이다", "그렇다",
         "매우 그렇다", "그렇다", "그렇다", "매우 그렇다",   "매우 그렇다",
         "그렇다", "그렇다", "그렇다", "그렇다", "보통이다",
         "그렇다", "그렇다", "그렇다", "매우 그렇다", "보통이다",
         "매우 그렇다", "매우 그렇다", "그렇다", "그렇다", "그렇지 않다",
         "그렇다", "그렇다"),
  v2=c( "그렇다", "보통이다", "그렇다", "그렇다", "보통이다", "그렇지 않다",
        "그렇지 않다", "매우 그렇다", "보통이다", "그렇다", "매우 그렇다",
        "그렇다", "매우 그렇다", "그렇다", "그렇다", "그렇다", "그렇다",
        "그렇다", "매우 그렇다", "그렇다", "보통이다", "그렇다", "그렇다",
        "그렇다", "매우 그렇다", "보통이다", "그렇다", "그렇다", "그렇지 않다",
        "그렇지 않다", "그렇지 않다", "그렇지 않다", "매우 그렇다", "매우 그렇다",
        "그렇지 않다", "그렇다", "전혀 그렇지 않다", "그렇다", "그렇다",
        "그렇다", "보통이다", "그렇다", "그렇지 않다", "그렇다", "보통이다",
        "매우 그렇다", "보통이다")
)
data2
auto_coding(data2)
# jutr1 %>% auto_coding( range=91:95, custom=c("아니다.","그렇다."),start_number=0)%>%dall("head")
# jutr1 %>% auto_coding( range="L01:L05", custom=c("아니다.","그렇다."),start_number=0)%>% head(30)

## End(Not run)
\end{ExampleCode}
\end{Examples}
\inputencoding{utf8}
\HeaderA{ave\_candidates\_plot}{ave\_candidates\_plot: 문항 제거 전후 AVE 비교 시각화}{ave.Rul.candidates.Rul.plot}
%
\begin{Description}
ave\_candidates\_plot: 문항 제거 전후 AVE 비교 시각화
\end{Description}
%
\begin{Usage}
\begin{verbatim}
ave_candidates_plot(
  df,
  latent = NULL,
  range_focus = TRUE,
  font_size = 4,
  padding = 0.04,
  v_line_show = TRUE
)
\end{verbatim}
\end{Usage}
%
\begin{Arguments}
\begin{ldescription}
\item[\code{df}] lavaan 객체 또는 ave\_raise\_candidates() 결과

\item[\code{latent}] 특정 잠재변수만 시각화 (선택)

\item[\code{range\_focus}] TRUE면 AVE 변화 중심으로 x축 범위 조정

\item[\code{font\_size}] 텍스트 크기

\item[\code{padding}] 범위조절 값 (기본 0.04)

\item[\code{v\_line\_show}] 각 문항의 ave\_current 위치에 수직선 추가 여부
\end{ldescription}
\end{Arguments}
%
\begin{Value}
ggplot2 객체
\end{Value}
%
\begin{Examples}
\begin{ExampleCode}
## Not run: 
library(lavaan)
library(ggplot2)
library(dplyr)

# 예시 데이터: HolzingerSwineford1939
model <- '
  visual  =~ x1 + x2 + x3
  textual =~ x4 + x5 + x6
  speed   =~ x7 + x8 + x9
'
fit <- cfa(model, data = lavaan::HolzingerSwineford1939)

# 1. lavaan 객체 입력 시 자동 처리
ave_candidates_plot(fit)

# 2. 특정 잠재변수만 보기
ave_candidates_plot(fit, latent = "visual")

# 3. AVE 변화량이 두드러진 경우에만 확대해서 보기
ave_candidates_plot(fit, range_focus = TRUE)

# 4. AVE 변화 전후 수직선 제거
ave_candidates_plot(fit, v_line_show = FALSE)

# 5. font 크기 조절
ave_candidates_plot(fit, font_size = 3)

# 6. ave_raise_candidates() 결과를 직접 전달
result <- ave_raise_candidates(fit, print_result = FALSE)
ave_candidates_plot(result)

## End(Not run)
\end{ExampleCode}
\end{Examples}
\inputencoding{utf8}
\HeaderA{ave\_raise\_candidates}{ave\_raise\_candidates: 잠재변수별 AVE 향상 후보 문항 탐색}{ave.Rul.raise.Rul.candidates}
%
\begin{Description}
lavaan 모형 결과를 기반으로 각 잠재변수별로 문항 제거 시
AVE(Average Variance Extracted)의 변화량을 계산합니다.
모든 잠재변수를 자동으로 분석하거나, 특정 잠재변수만 선택할 수 있습니다.
\end{Description}
%
\begin{Usage}
\begin{verbatim}
ave_raise_candidates(fit, latent = NULL, cutoff = -Inf, print_result = TRUE)
\end{verbatim}
\end{Usage}
%
\begin{Arguments}
\begin{ldescription}
\item[\code{fit}] lavaan::cfa 또는 lavaan::sem 결과 객체

\item[\code{latent}] 분석 대상 잠재변수 이름(문자열 벡터).
NULL이면 모델 내 모든 잠재변수를 자동 분석합니다.

\item[\code{cutoff}] AVE 변화량(delta)의 최소 기준. 기본값은 -Inf.

\item[\code{print\_result}] 논리값. TRUE이면 잠재변수별 요약 메시지를 출력합니다.
\end{ldescription}
\end{Arguments}
%
\begin{Value}
tibble
\begin{description}

\item[latent] 잠재변수 이름
\item[Item] 문항 이름
\item[lambda] 표준화 요인적재량
\item[ave\_current] 현재 AVE
\item[ave\_if\_dropped] 문항 제거 시 AVE
\item[delta] AVE 변화량 (제거 후 - 제거 전)

\end{description}

\end{Value}
%
\begin{Examples}
\begin{ExampleCode}
## Not run: 
library(lavaan)
library(dplyr)
library(purrr)
library(magrittr)

#----------------------------------------
# 1. CFA 모델 적합
#----------------------------------------
model <- '
  Visual  =~ x1 + x2 + x3
  Textual =~ x4 + x5 + x6
  Speed   =~ x7 + x8 + x9
'
fit <- cfa(model, data = HolzingerSwineford1939, std.lv = TRUE)

#----------------------------------------
# 2. 모든 잠재변수 자동 분석
#----------------------------------------
ave_raise_candidates(fit)

#----------------------------------------
# 3. 특정 잠재변수만 분석
#----------------------------------------
ave_raise_candidates(fit, latent = "Visual")

#----------------------------------------
# 4. AVE가 0.02 이상 증가하는 경우만 필터링
#----------------------------------------
ave_raise_candidates(fit, cutoff = 0.02)

#----------------------------------------
# 5. 여러 잠재변수 선택 분석
#----------------------------------------
ave_raise_candidates(fit, latent = c("Visual", "Textual"))

## End(Not run)
\end{ExampleCode}
\end{Examples}
\inputencoding{utf8}
\HeaderA{bartlett\_test}{bartlett\_test}{bartlett.Rul.test}
%
\begin{Description}
bartlett\_test
\end{Description}
%
\begin{Usage}
\begin{verbatim}
bartlett_test(
  data,
  formula = NULL,
  Mean = "Mean",
  sd = "SD",
  n = "N",
  var = "variable"
)
\end{verbatim}
\end{Usage}
%
\begin{Arguments}
\begin{ldescription}
\item[\code{data}] data.frame or mean. sd. n

\item[\code{formula}] dv \textasciitilde{} iv

\item[\code{Mean}] mean selection

\item[\code{sd}] sd selection

\item[\code{n}] n selection

\item[\code{var}] variable
\end{ldescription}
\end{Arguments}
%
\begin{Value}
result
\end{Value}
%
\begin{Examples}
\begin{ExampleCode}

## Not run: 

set.seed(20240418)
my_data <- data.frame(
  value = c(group1 <- rnorm(20, mean = 10, sd = 2),
            group2 <- rnorm(20, mean = 12, sd = 2),
            group3 <- rnorm(20, mean = 15, sd = 2)),
  group = rep(c("Group1", "Group2", "Group3"), each = 20)
)
my_data %>%mysummaryBy(value ~ group)%>%select(1:4) %>%bartlett_test()

data = data.frame(row.names = c("1", "2", "3"),
                  var = c('Group1', 'Group2', 'Group3'),
                  Mean = c('10.0782306287367', '11.5824600813951', '14.764623393546'),
                  SD = c('1.88743751774179', '1.9951475292551', '2.07862332265816'),
                  N = c('20', '20', '20'))
data
bartlett_test(data)

# original test
bartlett.test(value ~ group, data=my_data)
#now test
bartlett_test(value ~ group, data=my_data)


#origin
bartlett.test(mpg~cyl, data=mtcars)
#aggregated data
mtcars%>%mysummaryBy(mpg~cyl) %>%bartlett_test()
# Using function
bartlett_test(mpg~cyl, data=mtcars)
#'

## End(Not run)
\end{ExampleCode}
\end{Examples}
\inputencoding{utf8}
\HeaderA{binary\_recode}{binary\_recode}{binary.Rul.recode}
%
\begin{Description}
binary\_recode
\end{Description}
%
\begin{Usage}
\begin{verbatim}
binary_recode(data, variable, recode_list, new_name = NULL, tail = "_grp")
\end{verbatim}
\end{Usage}
%
\begin{Arguments}
\begin{ldescription}
\item[\code{data}] data

\item[\code{variable}] original variable

\item[\code{recode\_list}] list recode variable

\item[\code{new\_name}] new name

\item[\code{tail}] tails \_grp
\end{ldescription}
\end{Arguments}
%
\begin{Value}
data
\end{Value}
%
\begin{Examples}
\begin{ExampleCode}

## Not run: 
#'
# 예제 데이터
data <- data.frame(
  학력 = factor(c("고졸", "대학졸", "석박사졸", "고졸", "석박사졸")),
  grade = 1:5
)
data
binary_recode(data, grade, list(미숙자 = 1:3, 숙달자 = 4:5))
binary_recode(data, grade, list(미숙자 = c(1,2,3), 숙달자 = 4:5))
binary_recode(data, "grade", list(미숙자 = 1:3, 숙달자= 4:5), new_name="수준")

binary_recode(data, 학력, list(저학년= 1:2, 고학년= 3))
binary_recode(data, "학력", list(저학년= 1:2, 고학년= 3))
binary_recode(data, "학력", list(저학년= c("고졸","대학졸"), 고학년= 3))
binary_recode(data, 학력, list(저학년= c("고졸","대학졸"), 고학년= 3))
binary_recode(data, 학력, list(저학년= c("고졸","대학졸"), 고학년= "석박사졸"))


## End(Not run)

\end{ExampleCode}
\end{Examples}
\inputencoding{utf8}
\HeaderA{binary\_recode\_sep}{Multi-recoding binary variables}{binary.Rul.recode.Rul.sep}
%
\begin{Description}
Multi-recoding binary variables
\end{Description}
%
\begin{Usage}
\begin{verbatim}
binary_recode_sep(data, ..., new_names = NULL, tail = "_grp")
\end{verbatim}
\end{Usage}
%
\begin{Arguments}
\begin{ldescription}
\item[\code{data}] data

\item[\code{...}] recode list

\item[\code{new\_names}] newnaem list

\item[\code{tail}] tail \_grp
\end{ldescription}
\end{Arguments}
%
\begin{Value}
data
\end{Value}
%
\begin{Examples}
\begin{ExampleCode}


## Not run: 




# 예제 데이터
data <- data.frame(
  학력 = factor(c("고졸", "대학졸", "석박사졸", "고졸", "석박사졸")),
  grade = 1:5
)
data


# 함수 호출
binary_recode_sep(  data,
                    "학력", list(저학년 = 1:2, 고학년 = 3),
                    "grade", list(저학년 = 1:2, 고학년 = 3:5),
                    new_names = c("학력_grp", "grade_grp")
)





## End(Not run)
\end{ExampleCode}
\end{Examples}
\inputencoding{utf8}
\HeaderA{bind\_alpha\_table}{Alpha Table Combination}{bind.Rul.alpha.Rul.table}
%
\begin{Description}
Alpha Table Combination
\end{Description}
%
\begin{Usage}
\begin{verbatim}
bind_alpha_table(data, ..., check.keys = TRUE)
\end{verbatim}
\end{Usage}
%
\begin{Arguments}
\begin{ldescription}
\item[\code{data}] Data frame containing the variables

\item[\code{...}] Variable names to be included in the alpha table, column name and number. Variable names to be included in the alpha table

\item[\code{check.keys}] check.keys=TRUE
\end{ldescription}
\end{Arguments}
%
\begin{Value}
alpha table
\end{Value}
%
\begin{Examples}
\begin{ExampleCode}
## Not run: 
## select variable
bind_alpha_table(data = stat_onl, c("S_Review", "S_Add_learn", "S_Feedback"))

#error
bind_alpha_table(data = stat_onl, c(S_Review, S_Add_learn, S_Feedback))

#select col number
bind_alpha_table(data = stat_onl, c(9:11))

# mutilple data
bind_alpha_table( stat_onl,
                  c("S_Review", "S_Add_learn", "S_Feedback"),
                  c("SE_Place", "SE_Time"),
                  c("On_Satisfy", "On_Joy", "On_Easy", "On_Engage"),
                  c("upgrade","satisfy","fatigue_inverse") )


#mixed multiple
bind_alpha_table( stat_onl,
                  c("S_Review", "S_Add_learn", "S_Feedback"),
                  c("SE_Place", "SE_Time"),
                  c("On_Satisfy", "On_Joy", "On_Easy", "On_Engage"),
                  c(5:6, 21) )





## End(Not run)

\end{ExampleCode}
\end{Examples}
\inputencoding{utf8}
\HeaderA{bind\_compare\_data}{Table-building functions that pairwise compare two pieces of data}{bind.Rul.compare.Rul.data}
%
\begin{Description}
Table-building functions that pairwise compare two pieces of data
\end{Description}
%
\begin{Usage}
\begin{verbatim}
bind_compare_data(dff1, dff2, digits = 3, pair = TRUE)
\end{verbatim}
\end{Usage}
%
\begin{Arguments}
\begin{ldescription}
\item[\code{dff1}] data1

\item[\code{dff2}] data2

\item[\code{digits}] rounding

\item[\code{pair}] pari option default TRUE
\end{ldescription}
\end{Arguments}
%
\begin{Value}
table
\end{Value}
\inputencoding{utf8}
\HeaderA{bind\_lme\_model}{Tbind\_lme\_model}{bind.Rul.lme.Rul.model}
%
\begin{Description}
Tbind\_lme\_model
\end{Description}
%
\begin{Usage}
\begin{verbatim}
bind_lme_model(
  ...,
  model_name = paste("모형", 1:length(model)),
  title = "모형의 비교"
)
\end{verbatim}
\end{Usage}
%
\begin{Arguments}
\begin{ldescription}
\item[\code{...}] input

\item[\code{model\_name}] model name

\item[\code{title}] model title
\end{ldescription}
\end{Arguments}
%
\begin{Value}
table
\end{Value}
%
\begin{Examples}
\begin{ExampleCode}
## Not run: 
data(birthwt)
str(birthwt)
fitm <- lm(bwt ~ age+lwt+factor(race)+smoke+ptl+ht+ui, data=birthwt)
fitm2 <- lm(bwt ~ lwt+factor(race)+smoke+ht+ui, data=birthwt)
bind_lme_model(fitm, fitm2)


## End(Not run)
\end{ExampleCode}
\end{Examples}
\inputencoding{utf8}
\HeaderA{bind\_mysummaryBy}{bind\_mysummaryBy is group mean of mean}{bind.Rul.mysummaryBy}
%
\begin{Description}
bind\_mysummaryBy is group mean of mean
\end{Description}
%
\begin{Usage}
\begin{verbatim}
bind_mysummaryBy(data, ..., unite = FALSE)
\end{verbatim}
\end{Usage}
%
\begin{Arguments}
\begin{ldescription}
\item[\code{data}] data.frame

\item[\code{...}] fomula data ...

\item[\code{unite}] ... bind persional data and group data  TRUE
\end{ldescription}
\end{Arguments}
%
\begin{Value}
Use group data obtained from group descriptive statistics to obtain group-level descriptive statistics (means, standard deviations, etc.) and to analyze a mixed model
\end{Value}
%
\begin{Examples}
\begin{ExampleCode}
## Not run: 
##
bind_mysummaryBy(mtcars, mpg ~ am)

##compare
mysummaryBy(data=mtcars, mpg ~ am)

group_by(mtcars, am) %>%  summarise(mpg = mean(mpg))

 ## A tibble: 2 × 2
 ##    am   mpg
 ##<dbl> <dbl>
 ## 1     0  17.1
 ##2     1  24.4




bind_mysummaryBy(mtcars, mpg ~ am, mpg ~ vs)


## grp   dv        N  MEAN    SD   MIN   MAX
##  <chr> <chr> <int> <dbl> <dbl> <dbl> <dbl>
##1 am    mpg       2  20.8  5.12  17.1  24.4
##2 vs    mpg       2  20.6  5.61  16.6  24.6
#

bind_mysummaryBy(mtcars, mpg ~ am, mpg ~ vs, mpg ~cyl)

# A tibble: 3 × 9
# grp   dv        N  MEAN    SD   MIN   MAX      Skew  Kurt
# <chr> <chr> <int> <dbl> <dbl> <dbl> <dbl>     <dbl> <dbl>
#   1 am    mpg       2  20.8  5.12  17.1  24.4 -5.28e-16 -2.75
# 2 vs    mpg       2  20.6  5.61  16.6  24.6  0        -2.75
# 3 cyl   mpg       3  20.5  5.82  15.1  26.7  1.28e- 1 -2.33


jjstat::bind_mysummaryBy(mtcars, mpg ~ am, mpg ~ vs, mpg ~ cyl, unite=TRUE)

# A tibble: 3 × 8
# var         N  MEAN    SD   MIN   MAX      Skew  Kurt
# <chr>   <int> <dbl> <dbl> <dbl> <dbl>     <dbl> <dbl>
#   1 am_mpg      2  20.8  5.12  17.1  24.4 -5.28e-16 -2.75
# 2 vs_mpg      2  20.6  5.61  16.6  24.6  0        -2.75
# 3 cyl_mpg     3  20.5  5.82  15.1  26.7  1.28e- 1 -2.33





## End(Not run)
\end{ExampleCode}
\end{Examples}
\inputencoding{utf8}
\HeaderA{box\_cox}{Apply Box-Cox Style Transformations to a Data Frame}{box.Rul.cox}
%
\begin{Description}
This function applies various mathematical transformations to all numeric columns in a data frame.
\end{Description}
%
\begin{Usage}
\begin{verbatim}
box_cox(df, use = "log")
\end{verbatim}
\end{Usage}
%
\begin{Arguments}
\begin{ldescription}
\item[\code{df}] A data frame with numeric columns to be transformed.

\item[\code{use}] A character string specifying the transformation to apply. Options are:
\begin{itemize}

\item{} \code{"log"}: Logarithmic transformation (\eqn{\log(x + 1)}{})
\item{} \code{"sqrt"}: Square root transformation (\eqn{\sqrt{x}}{})
\item{} \code{"inverse"}: Inverse transformation (\eqn{1 / (x + 1)}{})
\item{} \code{"square"}: Square transformation (\eqn{x^2}{})
\item{} \code{"cube\_root"}: Cube root transformation (\eqn{x^{1/3}}{})
\item{} \code{"exp"}: Exponential transformation (\eqn{e^x}{})

\end{itemize}

\end{ldescription}
\end{Arguments}
%
\begin{Value}
A data frame with transformed numeric columns.
\end{Value}
%
\begin{Examples}
\begin{ExampleCode}
## Not run: 
# Example: Apply transformations to a data frame
df <- data.frame(
  col1 = c(1, 2, 3),
  col2 = c(4, 5, 6)
)

# Logarithmic transformation
box_cox(df, use = "log")

# Square root transformation
box_cox(df, use = "sqrt")

# Cube root transformation
box_cox(df, use = "cube_root")

## End(Not run)
\end{ExampleCode}
\end{Examples}
\inputencoding{utf8}
\HeaderA{breslowdata\_test}{Breslow-Day Test for Homogeneity of the Odds Ratios}{breslowdata.Rul.test}
%
\begin{Description}
Breslow-Day Test for Homogeneity of the Odds Ratios
\end{Description}
%
\begin{Usage}
\begin{verbatim}
breslowdata_test(data, OR = NA, correct = FALSE, tidy = TRUE)
\end{verbatim}
\end{Usage}
%
\begin{Arguments}
\begin{ldescription}
\item[\code{data}] 2x2xk table

\item[\code{OR}] the odds ratio to be tested against. If left undefined (default) the Mantel-Haenszel estimate will be used.

\item[\code{correct}] If TRUE, the Breslow-Day test with Tarone's adjustment is computed, which subtracts an adjustment factor to make the resulting statistic asymptotically chi-square.

\item[\code{tidy}] broom::tidy
\end{ldescription}
\end{Arguments}
%
\begin{Details}
For the Breslow-Day test to be valid, the sample size should be relatively large in each stratum, and at least 80
Alternatively, it might be better to cast the entire inference problem into the setting of a logistic regression model. Here, the underlying question of the Breslow-Day test can be answered by investigating whether an interaction term with the strata variable is necessary (e.g. using a likelihood ratio test using the anova function).
\end{Details}
%
\begin{Value}
result
\end{Value}
%
\begin{Examples}
\begin{ExampleCode}
## Not run: 
#'
migraine <- xtabs(freq ~ .,
                  cbind(expand.grid(treatment=c("active", "placebo"),
                                    response =c("better", "same"),
                                    gender   =c("female", "male")),
                        freq=c(16, 5, 11, 20, 12, 7, 16, 19))
)
# get rid of gender
tab <- xtabs(Freq ~ treatment + response, migraine)
DescTools::Desc(tab)
# only the women
female <- migraine[,, 1]
DescTools::Desc(female)
# .. and the men
male <- migraine[,, 2]
DescTools::Desc(male)
breslowdata_test(migraine)
breslowdata_test(migraine, correct = TRUE)


## End(Not run)
\end{ExampleCode}
\end{Examples}
\inputencoding{utf8}
\HeaderA{calculate\_chi\_sig}{Create oserved/expected tapes after chisq calculation and give them individual significance}{calculate.Rul.chi.Rul.sig}
%
\begin{Description}
Create oserved/expected tapes after chisq calculation and give them individual significance
\end{Description}
%
\begin{Usage}
\begin{verbatim}
calculate_chi_sig(observed, type = "data", digits = 3, exact = TRUE)
\end{verbatim}
\end{Usage}
%
\begin{Arguments}
\begin{ldescription}
\item[\code{observed}] contigency table

\item[\code{type}] data, Res(all), chisq\_test, cramers v, observed\_over\_expected\_ratios,cell\_p\_values, p\_sig, data, data2

\item[\code{digits}] digits= 3

\item[\code{exact}] yYates's continuity correction
\end{ldescription}
\end{Arguments}
%
\begin{Value}
matrix, data.frame
\end{Value}
%
\begin{Examples}
\begin{ExampleCode}

## Not run: 
# # calculate_chi_sig(data)
# Example data
matrix(c(36, 67, 11,
          0, 35, 44,
         41, 108, 1,
          56, 87, 54),
           nrow = 4, byrow = TRUE,
   dimnames = list(c("공명음", "마찰음",
     "유기음_경음", "평파열음_평파찰음"),  c("H", "H(H)", "L"))) %>%
     calculate_chi_sig("p_sig")


     matrix(c(36, 67, 11,
              40, 35, 44,
              41, 108, 1,
              56, 87, 54), nrow = 4, byrow = TRUE,
                     dimnames = list(c("공명음", "마찰음",
                                    "유기음_경음", "평파열음_평파찰음"),
                                 c("H", "H(H)", "L"))) %>%
                calculate_chi_sig(simple = TRUE)



  calculate_chi_sig(matrix(c(36, 67, 11,
                      40, 35, 44,
                   41, 108, 1,
              56, 87, 54), nrow = 4, byrow = TRUE,
               dimnames = list(c("공명음", "마찰음",
              "유기음_경음", "평파열음_평파찰음"),
              c("H", "H(H)", "L")))   )


 # calculate_chi_sig(data)

## End(Not run)

\end{ExampleCode}
\end{Examples}
\inputencoding{utf8}
\HeaderA{cal\_vif}{calculation vif}{cal.Rul.vif}
%
\begin{Description}
calculation vif
\end{Description}
%
\begin{Usage}
\begin{verbatim}
cal_vif(model)
\end{verbatim}
\end{Usage}
%
\begin{Arguments}
\begin{ldescription}
\item[\code{model}] lm model data
\end{ldescription}
\end{Arguments}
%
\begin{Value}
vif
\end{Value}
%
\begin{Examples}
\begin{ExampleCode}

## Not run: 
#'
# Create the linear model
model <- lm(hp ~ wt + disp + drat, data = mtcars)

# Calculate VIF using the custom function
cal_vif(model)

lm(mpg ~ wt + disp + drat + hp , data = mtcars) %>%cal_vif()



## End(Not run)
\end{ExampleCode}
\end{Examples}
\inputencoding{utf8}
\HeaderA{centering}{centering each vector}{centering}
%
\begin{Description}
centering each vector
\end{Description}
%
\begin{Usage}
\begin{verbatim}
centering(variable, na.rm = TRUE)
\end{verbatim}
\end{Usage}
%
\begin{Arguments}
\begin{ldescription}
\item[\code{variable}] vector variable in data

\item[\code{na.rm}] missing data chedk
\end{ldescription}
\end{Arguments}
\inputencoding{utf8}
\HeaderA{cfa2}{Create CFA result summary by lavaan sem estimated result By Park Joonghee Ph.D}{cfa2}
%
\begin{Description}
Create CFA result summary by lavaan sem estimated result By Park Joonghee Ph.D
\end{Description}
%
\begin{Usage}
\begin{verbatim}
cfa2(
  x,
  format = "markdown",
  htmt_cut = 0.9,
  htmt2 = FALSE,
  cut = 0.7,
  angle = 90,
  cex = 11,
  hjust = 0.9,
  val.size = 4,
  dis.sort = TRUE,
  rename = FALSE,
  var_name = NA,
  digits = 3,
  type = "all"
)
\end{verbatim}
\end{Usage}
%
\begin{Arguments}
\begin{ldescription}
\item[\code{x}] lavaan result

\item[\code{format}] knitr markdown decision

\item[\code{htmt\_cut}] htmt\_cut is default =1.
if you want roburst option htmt\_cut =0.9

\item[\code{htmt2}] TRUE using htmt2

\item[\code{cut}] loading criteran

\item[\code{angle}] tesxt angle

\item[\code{cex}] text size

\item[\code{hjust}] on text value gap

\item[\code{val.size}] text size

\item[\code{dis.sort}] item sort

\item[\code{rename}] item validiation new name apply

\item[\code{var\_name}] new text name default NULL

\item[\code{digits}] value rounding

\item[\code{dataset}] this is import name that dataset column name

\item[\code{model}] lavaan model need indices validation name

\item[\code{res}] result type
\end{ldescription}
\end{Arguments}
%
\begin{Examples}
\begin{ExampleCode}
## Not run: 
library(lavaan)
data(PoliticalDemocracy)

model1 <- "
  ind60 =~ x1 + x2 + x3
  dem60 =~ y1 + y2 + y3 + y4
  dem65 =~ y5 + y6 + y7 + y8
  dem60 ~ ind60
  dem65 ~ ind60 + dem60
"

fit1 <- sem(model1, data = PoliticalDemocracy)

# 기본 사용법: 전체 결과 출력
cfa2(fit1)

# 특정 결과(res)만 출력 예시
cfa2(fit1, res = "modelfit")        # 모델 적합도
cfa2(fit1, res = "loadings")        # 요인적재량 테이블
cfa2(fit1, res = "alpha")           # 신뢰도(Cronbach alpha)
cfa2(fit1, res = "CR_AVE")          # 수렴타당도(CR, AVE)
cfa2(fit1, res = "Discriminant")    # 판별타당도 (Fornell-Larcker 기준)
cfa2(fit1, res = "htmt", dataset = PoliticalDemocracy, model = model1)  # HTMT 지표
cfa2(fit1, res = "str_cor")         # 잠재변수 간 경로 유의성

# 시각화 결과: 요인적재량 바그래프 출력
cfa2(fit1, res = "loadings_bar")

# 사용자 정의 변수명 적용(rename + var_name)
new_names <- c("Indicator1", "Indicator2", "Indicator3",
               "Democracy1", "Democracy2", "Democracy3", "Democracy4",
               "Democracy5", "Democracy6", "Democracy7", "Democracy8")

# 그래프의 색을 바꾸고 싶은 경우
cfa2(fit, type = "loadings_bar") +
scale_fill_grey()
# rename = TRUE일 경우에만 var_name이 반영됨
cfa2(fit1, rename = TRUE, var_name = new_names)

# rename = FALSE일 경우 var_name은 무시됨
cfa2(fit1, rename = FALSE, var_name = new_names)

# 다양한 출력 포맷 지정
cfa2(fit1, format = "markdown")
cfa2(fit1, format = "html") %>% kableExtra::kable_minimal(full_width = FALSE)

# 모델 간 비교
model2 <- "
  ind60 =~ x1 + x2 + x3
  dem60 =~ y1 + y2 + y3 + y4
  dem65 =~ y5 + y6 + y7 + y8
  dem60 ~ ind60
  dem65 ~ ind60 + dem60
"
fit2 <- sem(model2, data = PoliticalDemocracy)
CompareFit_diff(fit1, fit2)

# AVE 직접 확인
AVE(fit1)
AVE(fit2)

# 통합 요약 테이블 출력 (CR, AVE 포함)
library(dplyr)
full_join(
  cfa2(fit1, type = "item") %>% dplyr::select(-Accept, -p),
  cfa2(fit1, type = "Convergent") %>% row2col("Latent") %>%
    dplyr::select(-Cronbach) %>% Round(3),
  by = "Latent"
) %>% nice_table()

## End(Not run)
\end{ExampleCode}
\end{Examples}
\inputencoding{utf8}
\HeaderA{cfa\_ml}{Multilevel CFA summary wrapper}{cfa.Rul.ml}
%
\begin{Description}
다층 구조방정식(CFA/SEM) 모형의 within / between 수준 결과를
한 번에 요약하여 제공하는 래퍼 함수이다.
\end{Description}
%
\begin{Usage}
\begin{verbatim}
cfa_ml(
  fit,
  level = c("within", "between"),
  type = c("all", "model_fit", "indicator_validity", "CR_AVE", "item_CR_AVE", "fornell",
    "htmt"),
  digits = 3
)
\end{verbatim}
\end{Usage}
%
\begin{Arguments}
\begin{ldescription}
\item[\code{fit}] lavaan multilevel SEM object

\item[\code{level}] character. "within" 또는 "between"

\item[\code{type}] character. 출력 제어 옵션
\begin{itemize}

\item{} "all"모든 결과를 list로 반환
\item{} "model\_fit"모형 적합도
\item{} "indicator\_validity"요인부하량 및 지표 타당도
\item{} "CR\_AVE"잠재변수 수준 CR / AVE
\item{} "item\_CR\_AVE"문항 수준 CR / AVE
\item{} "fornell"Fornell–Larcker 판별타당도
\item{} "htmt"HTMT 판별타당도

\end{itemize}


\item[\code{digits}] numeric. 소수점 자리수
\end{ldescription}
\end{Arguments}
%
\begin{Value}
data.frame 또는 list
\end{Value}
%
\begin{Examples}
\begin{ExampleCode}
## Not run: 
res_all <- cfa_ml(fit_ml, level = "within", type = "all")
res_fit <- cfa_ml(fit_ml, level = "between", type = "model_fit")

## End(Not run)
\end{ExampleCode}
\end{Examples}
\inputencoding{utf8}
\HeaderA{change\_node\_name}{Change node names in semPlot / qgraph object}{change.Rul.node.Rul.name}
%
\begin{Description}
semPlot::semPaths()로 생성된 구조방정식 시각화 객체에서
노드 이름을 사용자가 지정한 (원래이름, 변경라벨) 쌍에 따라
일괄적으로 변경한다.
본 함수는 pipe(
\end{Description}
%
\begin{Usage}
\begin{verbatim}
change_node_name(plot_obj, var_named)
\end{verbatim}
\end{Usage}
%
\begin{Arguments}
\begin{ldescription}
\item[\code{plot\_obj}] semPlot 또는 qgraph 객체

\item[\code{var\_named}] character vector.
c(old1, new1, old2, new2, ...)
\end{ldescription}
\end{Arguments}
%
\begin{Value}
노드 명칭이 변경된 semPlot 객체
\end{Value}
%
\begin{Examples}
\begin{ExampleCode}
## Not run: 

#################################
# 1. SCCT 구조방정식 모형 정의
#################################

model_scct_indirect <- '

REFF =~ efficacy2 + efficacy3 + efficacy4
OEXP =~ Outcome1 + Outcome2 + Outcome3 + Outcome4 + Outcome5 + Outcome8
INTR =~ Interest1 + Interest2 + Interest3 + Interest4

REFF ~ e1*grade + e2*experience0
OEXP ~ a1*REFF + d1*grade + d2*experience0
INTR ~ b1*REFF + b2*OEXP

intention ~ c1*INTR + c2*OEXP + c3*REFF

grade ~~ experience0

# 간접효과
ind_grade_REFF        := e1 * c3
ind_grade_REFF_INTR   := e1 * b1 * c1
ind_grade_REFF_OEXP   := e1 * a1 * c2
ind_grade_REFF_OEXP_INTR := e1 * a1 * b2 * c1

ind_exp_REFF          := e2 * c3
ind_exp_REFF_INTR     := e2 * b1 * c1
ind_exp_REFF_OEXP     := e2 * a1 * c2
ind_exp_REFF_OEXP_INTR := e2 * a1 * b2 * c1
'

#################################
# 2. 모형 추정
#################################

library(lavaan)
fit_scct_indirect <- sem(
  model = model_scct_indirect,
  data  = shindata,
  se = "bootstrap",
  bootstrap = 2000
)

#################################
# 3. 노드 명칭 매핑 벡터
#################################

var_named1 <- c(
  "REFF", "연구효능감",
  "efficacy2","연구효능2",
  "efficacy3","연구효능3",
  "efficacy4","연구효능4",

  "INTR", "진로흥미",
  "Interest1","진로흥미1",
  "Interest2","진로흥미2",
  "Interest3","진로흥미3",
  "Interest4","진로흥미4",

  "OEXP", "결과기대",
  "Outcome1","결과기대1",
  "Outcome2","결과기대2",
  "Outcome3","결과기대3",
  "Outcome4","결과기대4",
  "Outcome5","결과기대5",
  "Outcome8","결과기대8",

  "grade","학년",
  "experience0","연구경험",
  "intention","진로의지"
)

#################################
# 4. SCCT 전용 layout 행렬
#################################

layout_scct <- matrix(
  c(
    NA, "efficacy2", "efficacy3", "efficacy4", NA,
        "Interest1", "Interest2", "Interest3", "Interest4", NA,

    "experience0", NA, "REFF", NA, NA,
        NA, "INTR", NA, NA, NA,

    "grade", NA, "OEXP", NA, NA,
        NA, NA, "intention", NA, NA,

    NA, "Outcome1", "Outcome2", "Outcome3", "Outcome4",
        "Outcome5", "Outcome8", NA, NA, NA
  ),
  nrow = 4,
  byrow = TRUE
)

#################################
# 5. 구조방정식 그림 + 노드명 변경
#################################

library(semPlot)
library(semptools)

semPlot::semPaths(
  fit_scct_indirect,
  layout        = layout_scct,
  whatLabels    = "std",
  style         = "ram",
  residuals     = FALSE,
  intercepts    = FALSE,
  nCharNodes    = 0,
  sizeLat       = 9,
  sizeMan       = 9,
  edge.label.cex = 0.85,
  edge.label.position = 0.65
) %>%
  semptools::mark_sig(fit_scct_indirect) %>%
  change_node_name(var_named1) %>%  ##### 이것을 이용하여 이름 변경
  plot()


## End(Not run)
\end{ExampleCode}
\end{Examples}
\inputencoding{utf8}
\HeaderA{Char2num}{Functions to convert character to numeric}{Char2num}
%
\begin{Description}
Functions to convert character to numeric
\end{Description}
%
\begin{Usage}
\begin{verbatim}
Char2num(data, iter = NULL, char_col = 1, last_col = FALSE)
\end{verbatim}
\end{Usage}
%
\begin{Arguments}
\begin{ldescription}
\item[\code{data}] data.frame

\item[\code{iter}] Select the range of variables to convert

\item[\code{char\_col}] Specify which columns should be kept as char col
\end{ldescription}
\end{Arguments}
%
\begin{Value}
data.frame
\end{Value}
%
\begin{Examples}
\begin{ExampleCode}

## Not run: 

## Replacing a character variable with numeric
mysummary(sats, "sux1") %>%
    mutate(across(everything(), ~ format(., digits = 4))) %>%
    tibble() %>% Char2num()


## Set remove=TRUE if you only want to process the selected variable and output only that.
  mysummary(sats, "sux1") %>%
  mutate(across(everything(), ~ format(., digits = 5))) %>%
    tibble() %>%Char2num(auto=FALSE, iter= 4:ncol(.), remove= TRUE)

##
mysummary(sats, "sux1") %>%
  mutate(across(everything(), ~ format(., digits = 5))) %>%
    tibble() %>%Char2num(auto=FALSE, iter= 4:ncol(.), remove= FALSE)
## same
  mysummary(sats, "sux1") %>%
     mutate(across(everything(), ~ format(., digits = 5))) %>%
     tibble() %>%Char2num(auto=FALSE, iter= 4:ncol(.))


#새로운 데이터
reg_text <- "var B se b t
성취감 -.03 .09 -.03 -.33
인정감 .24 .09 .20 2.50
책임감 .30 .08 .23 3.56***
발전성 -.08 .06 -.10 -1.29
감독 .06 .05 .08 1.35
보수 -.14 .07 -.18 -2.01*
대인관계 .06 .06 .06 .89
작업환경 .33 .09 .34 3.82***
안정성 .17 .05 .26 3.42**
성별 -.05 .09 -.03 -.52
연령 -.20 .07 -.22 -3.01**
학력 .05 .07 .04 .70
소득 .00 .04 .01 .12
종사기간 .17 .06 .23 2.71**
기관규모 -.10 .03 -.20 -3.36**"

# 함수 호출
text2df(reg_text) %>% tibble()
text2df(reg_text)%>% Char2num() %>% tibble()
text2df(reg_text)%>% Char2num(iter=2:4) %>% tibble()

## End(Not run)

\end{ExampleCode}
\end{Examples}
\inputencoding{utf8}
\HeaderA{Char2num\_fct}{fact varaible to num  in data.frame}{Char2num.Rul.fct}
%
\begin{Description}
fact varaible to num  in data.frame
\end{Description}
%
\begin{Usage}
\begin{verbatim}
Char2num_fct(df, col = 1)
\end{verbatim}
\end{Usage}
%
\begin{Arguments}
\begin{ldescription}
\item[\code{df}] df

\item[\code{col}] select col
\end{ldescription}
\end{Arguments}
%
\begin{Value}
df
\end{Value}
\inputencoding{utf8}
\HeaderA{chisq.multcomp}{RVAideMemoire::chisq.multcomp, Pairwise comparisons after a chi-squared goodness-of-fit test}{chisq.multcomp}
%
\begin{Description}
RVAideMemoire::chisq.multcomp, Pairwise comparisons after a chi-squared goodness-of-fit test
\end{Description}
%
\begin{Usage}
\begin{verbatim}
chisq.multcomp(x, p.method = "fdr")
\end{verbatim}
\end{Usage}
%
\begin{Arguments}
\begin{ldescription}
\item[\code{x}] numeric vector (counts).

\item[\code{p.method}] method for p-values correction. See help of p.adjust.
\end{ldescription}
\end{Arguments}
%
\begin{Value}
result
\end{Value}
%
\begin{Examples}
\begin{ExampleCode}

## Not run: 
#'
counts <- c(49,30,63,59)
chisq.test(counts)
chisq.multcomp(counts)
#comapre
chisq_multcomp(counts)


## End(Not run)
\end{ExampleCode}
\end{Examples}
\inputencoding{utf8}
\HeaderA{chisqTest}{chisq\_test 2way}{chisqTest}
%
\begin{Description}
chisq\_test 2way
\end{Description}
%
\begin{Usage}
\begin{verbatim}
chisqTest(Obs, correct = FALSE, digits = 3, cat = TRUE)
\end{verbatim}
\end{Usage}
%
\begin{Arguments}
\begin{ldescription}
\item[\code{Obs}] observation

\item[\code{correct}] Yates' continuity correction

\item[\code{digits}] 3

\item[\code{cat}] chisq.test form
\end{ldescription}
\end{Arguments}
%
\begin{Value}
chisq
\end{Value}
%
\begin{Examples}
\begin{ExampleCode}
## Not run: 

table(mtcars$am, mtcars$vs) %>% chisq_test()


table(mtcars$am, mtcars$vs) %>% chisq_test(correct=F)

## compare original chisq.test
table(mtcars$am, mtcars$vs) %>% chisq.test()


##Pearson's Chi-squared test with Yates' continuity correction
table(mtcars$am, mtcars$vs) %>% chisq_test(correct=T)

## End(Not run)


\end{ExampleCode}
\end{Examples}
\inputencoding{utf8}
\HeaderA{chisq\_apa}{카이제곱 apa보고}{chisq.Rul.apa}
%
\begin{Description}
카이제곱 apa보고
\end{Description}
%
\begin{Usage}
\begin{verbatim}
chisq_apa(data, chi_type = NULL, show = TRUE)
\end{verbatim}
\end{Usage}
%
\begin{Arguments}
\begin{ldescription}
\item[\code{data}] chisq.test result data

\item[\code{chi\_type}] "적합성 검정", "독립성 검정", "동질성 검정"

\item[\code{show}] TRUE cat output
\end{ldescription}
\end{Arguments}
%
\begin{Value}
data result
\end{Value}
%
\begin{Examples}
\begin{ExampleCode}

## Not run: 
#'
dfa= matrix(c(22, 26, 55,13, 31, 33,13, 16, 38), ncol=3, byrow=TRUE)
colnames(dfa)=c("가","나","다")
rownames(dfa)=c("A","B","C")
dfa
dfat= t(dfa)
dfat%>%chisq.test()%>%chisq_apa()

## End(Not run)
\end{ExampleCode}
\end{Examples}
\inputencoding{utf8}
\HeaderA{chisq\_contrast}{Chisq test Contrast test}{chisq.Rul.contrast}
%
\begin{Description}
Chisq test Contrast test
\end{Description}
%
\begin{Usage}
\begin{verbatim}
chisq_contrast(
  data,
  ...,
  test = chisq.test,
  p.adjust = "bonferroni",
  adj = FALSE,
  name = FALSE
)
\end{verbatim}
\end{Usage}
%
\begin{Arguments}
\begin{ldescription}
\item[\code{data}] data

\item[\code{...}] contrast

\item[\code{test}] chisq.test, fisher.test

\item[\code{p.adjust}] c("bonferroni","fdr","BH","BY","holm","hochberg","hommel")

\item[\code{adj}] cadjust p value TRUE

\item[\code{name}] rowname, when  pairwise
\end{ldescription}
\end{Arguments}
%
\begin{Value}
data
\end{Value}
%
\begin{Examples}
\begin{ExampleCode}

## Not run: 
#'
d <- matrix(c(584, 24, 1721, 56, 2400, 90, 8164, 289),
            nrow = 4, byrow = TRUE,
            dimnames = list(c("kangnam", "Moksdong", "nowon", "dongtan"), c("Yes", "No")))

d%>%chisq.test()
chisq_contrast(d)
chisq_contrast(d, c(1,2,3,4))


chisq_contrast(d, c(1,1,2,3),c(2,1,2,2),c(2,2,1,2),c(2,2,2,1))
# fisher.test
chisq_contrast(d, c(1,2,2,2),c(2,1,2,2),c(2,2,1,2),c(2,2,2,1), test=fisher.test)
# fisher.exact test
chisq_contrast(d, c(1,2,2,3),c(2,1,2,2),c(2,2,1,2),c(2,2,2,1), test=fisher.test)

chisq_contrast(d, c(1,2,2,2),c(2,1,1,2),c(2,2,1,2),c(2,2,2,1), test=fisher.test)
chisq_contrast(d, c(1,2,2,2),c(2,3,1,2),c(2,2,1,2),c(2,2,2,1), test=fisher.test)
chisq_contrast(d, c(1,2,2,2),c(2,3,1,2),c(2,2,1,3),c(2,2,2,1), test=fisher.test)
chisq_contrast(d, c(1,2,2,2),c(2,3,1,2),c(2,2,1,3),c(3,4,2,1), test=fisher.test)


#other data
data_text = "Tobacco_Dependence (0,5] (5,10] (10,15] (15,20] (20,98]
No_Nicotine_Dependence 130 210 43 114 20
Nicotine_Dependence 119 267 91 254 67"

#method2
data_text%>% text2df(header=T) %>% make_df_text()
nt
nt = data_text%>% text2df(header=T)%>% jjstat::col2row()
nt %>%tibble()
nt=Char2num(nt,1:5)
nt%>%t() %>%as.data.frame()%>%
  chisq_contrast(c(1,2,2,2,2),
                 c(2,1,2,2,2),
                 c(2,2,1,2,2),
                 c(2,3,4,1,2))

## End(Not run)


\end{ExampleCode}
\end{Examples}
\inputencoding{utf8}
\HeaderA{chisq\_gof\_posthoc}{chisq\_gof\_posthoc}{chisq.Rul.gof.Rul.posthoc}
%
\begin{Description}
chisq\_gof\_posthoc
\end{Description}
%
\begin{Usage}
\begin{verbatim}
chisq_gof_posthoc(vector, type = "res", method = "bonferroni", names = TRUE)
\end{verbatim}
\end{Usage}
%
\begin{Arguments}
\begin{ldescription}
\item[\code{vector}] data

\item[\code{type}] all, res, p, chisq

\item[\code{method}] p.adust fdr, bonferroni, ...

\item[\code{names}] TRUE using names
\end{ldescription}
\end{Arguments}
%
\begin{Value}
data
\end{Value}
%
\begin{Examples}
\begin{ExampleCode}

## Not run: 
#'

X <- c(49, 30, 63, 59)

chisq_gof_posthoc(X, "res")
chisq_gof_posthoc(X,"all")$gof_test
chisq_gof_posthoc(X,"all")$p
chisq_gof_posthoc(X,"all")$chisq
chisq_gof_posthoc(X, typ="all")
chisq_gof_posthoc(X, typ="res")
chisq_gof_posthoc(X, type="p")
chisq_gof_posthoc(X, type="chisq")

chisq_gof_posthoc(X, method="bonferroni")

X1 <- c(4, 3, 6, 19)
chisq_gof_posthoc(X1)
chisq_gof_posthoc(X1, typ="all")
chisq_gof_posthoc(X1, type="p")
chisq_gof_posthoc(X1, type="chisq")
chisq_gof_posthoc(X1, method="bonferroni")


chisq_gof_posthoc(X, method="bonferroni")

chisq_gof_posthoc(c(49, 30, 63, 59, 40, 60))

#' #'

## End(Not run)



\end{ExampleCode}
\end{Examples}
\inputencoding{utf8}
\HeaderA{chisq\_gof\_posthoc\_apa}{chisq\_gof\_posthoc\_apa}{chisq.Rul.gof.Rul.posthoc.Rul.apa}
%
\begin{Description}
chisq\_gof\_posthoc\_apa
\end{Description}
%
\begin{Usage}
\begin{verbatim}
chisq_gof_posthoc_apa(df, adj = FALSE)
\end{verbatim}
\end{Usage}
%
\begin{Arguments}
\begin{ldescription}
\item[\code{df}] df

\item[\code{adj}] adj
\end{ldescription}
\end{Arguments}
%
\begin{Value}
df
\end{Value}
\inputencoding{utf8}
\HeaderA{chisq\_multcomp2}{Pairwise comparisons after a chi-squared goodness-of-fit test}{chisq.Rul.multcomp2}
%
\begin{Description}
Pairwise comparisons after a chi-squared goodness-of-fit test
\end{Description}
%
\begin{Usage}
\begin{verbatim}
chisq_multcomp2(x, p.method = "fdr", ko = TRUE, sep = " / ")
\end{verbatim}
\end{Usage}
%
\begin{Arguments}
\begin{ldescription}
\item[\code{x}] numeric vector (counts).

\item[\code{p.method}] method for p-values correction. See help of p.adjust

\item[\code{ko}] TRUE variable names is changed to korean language

\item[\code{sep}] sep =" / " default
\end{ldescription}
\end{Arguments}
%
\begin{Value}
result
\end{Value}
%
\begin{Section}{Functions}
\begin{itemize}

\item{} \code{chisq\_multcomp2()}: pairwise comparisons after a global chi-squared goodness-of-fit test.


\end{itemize}
\end{Section}
%
\begin{Examples}
\begin{ExampleCode}
## Not run: 
#'
counts <- c(49,30,63,59)
chisq.test(counts)
chisq_multcomp2(counts)
# compare
chisq.multcomp2(counts)


## End(Not run)
\end{ExampleCode}
\end{Examples}
\inputencoding{utf8}
\HeaderA{chisq\_pivot\_table}{pivot}{chisq.Rul.pivot.Rul.table}
%
\begin{Description}
pivot
\end{Description}
%
\begin{Usage}
\begin{verbatim}
chisq_pivot_table(
  data,
  formula,
  names_from = NULL,
  unite = FALSE,
  col = 1:2,
  remove = TRUE,
  change_col = TRUE
)
\end{verbatim}
\end{Usage}
%
\begin{Arguments}
\begin{ldescription}
\item[\code{data}] data

\item[\code{formula}] formul

\item[\code{names\_from}] column

\item[\code{unite}] 1,2 colum unite

\item[\code{col}] 합쳐질 컬럼

\item[\code{remove}] 기존것 지우기

\item[\code{change\_col}] 열의 위치 바꾸기
\end{ldescription}
\end{Arguments}
%
\begin{Value}
데이터
\end{Value}
%
\begin{Examples}
\begin{ExampleCode}

## Not run: 

#' kge_bind3a %>%
filter(type == "고유어1")%>%
  pivot_table(~ area + age + 성조, unite=TRUE)
#'
kge_bind3a %>%
  filter(type == "고유어1")%>%
  pivot_table(~  age + area+ 성조, unite=FALSE)

#고유어1 공명음  지역별 연령
kge_bind3a %>%
  filter(type == "고유어1" & a1=="공명음")%>%
  pivot_table(~ age + area + 성조, unite=T)  #%>% md()



kge_bind3a %>%
  filter(type == "고유어1")%>%
  pivot_table(~  age +  성조, unite=FALSE)

#고유어1 공명음  지역별 연령
kge_bind3a %>%
  filter(type == "고유어1" & a1=="공명음")%>%
  pivot_table(~  area + 성조, unite=F)  #%>% md()


  
## End(Not run)
#'

\end{ExampleCode}
\end{Examples}
\inputencoding{utf8}
\HeaderA{chisq\_posthoc}{chisq post hoc}{chisq.Rul.posthoc}
%
\begin{Description}
chisq post hoc
\end{Description}
%
\begin{Usage}
\begin{verbatim}
chisq_posthoc(
  tbl,
  test = fisher.test,
  Rows = TRUE,
  control = c("bonferroni", "fdr", "BH", "BY", "holm", "hochberg", "hommel"),
  digits = 4,
  ...
)
\end{verbatim}
\end{Usage}
%
\begin{Arguments}
\begin{ldescription}
\item[\code{tbl}] tavke

\item[\code{test}] fusger.test, chisq.test

\item[\code{Rows}] rows , FLASE = cols

\item[\code{control}] p adjust"bonferroni","fdr","BH","BY","holm","hochberg","hommel"

\item[\code{digits}] 4

\item[\code{...}] ...
\end{ldescription}
\end{Arguments}
%
\begin{Value}
result
\end{Value}
%
\begin{Examples}
\begin{ExampleCode}

## Not run: 
as.table(cbind("T"=c(30,40,60,120),"S"= c(50,70,60,80)))%>%
  chisq_posthoc(Rows  = TRUE)

as.table(cbind("T"=c(30,40,60,120),"S"= c(50,70,60,80)))%>%
  chisq_posthoc(Rows  = FALSE)

as.table(cbind("T"=c(30,40,60,120),"S"= c(50,70,60,80)))%>%
  chisq_posthoc(test="chisq.test", Rows  = TRUE)

## End(Not run)


\end{ExampleCode}
\end{Examples}
\inputencoding{utf8}
\HeaderA{chisq\_retest}{Functions to check the chi-square test of a paper}{chisq.Rul.retest}
%
\begin{Description}
Functions to check the chi-square test of a paper
\end{Description}
%
\begin{Usage}
\begin{verbatim}
chisq_retest(
  input,
  ncol = 2,
  colname = NULL,
  rowname = NULL,
  type = "all",
  simulate.p.value = FALSE,
  digits = 3,
  correct = TRUE,
  B = 2000
)
\end{verbatim}
\end{Usage}
%
\begin{Arguments}
\begin{ldescription}
\item[\code{input}] vectro

\item[\code{ncol}] matrix col

\item[\code{colname}] input

\item[\code{rowname}] input

\item[\code{type}] all, chisq

\item[\code{simulate.p.value}] default FaLSE fisher exact test

\item[\code{digits}] round 3

\item[\code{correct}] default is Pearson’s Chi-squared test with Yates’ continuity correction (If you have something with a small sample count, you should use it. ), FALSE Pearson's Chi-squared test

\item[\code{B}] B is bootstrap default 2000
\end{ldescription}
\end{Arguments}
%
\begin{Value}
chisquare data
\end{Value}
%
\begin{Examples}
\begin{ExampleCode}

## Not run: 
## Mina Choi. (2007). Analysis of learning effects
# according to online and offline learning environments
#  in the same class by the same instructor.
#  Educational Information Media Research, 13(4>  5-23.
## jjpaper
## https://blog.naver.com/shoutjoy/223005233887

chisq_retest(c(111, 33, 64, 58), ncol=2 )
## 논문에 나온 값이 잘못됨 :chisq = 17.793
## this chisq = 16.7

chisq_retest(c(111, 33, 64, 58), ncol=2, type="chisq" )

chisq_retest(c(111, 33, 64, 58), ncol=2,
  colname = c("experience", "Non-experience"),
  rowname =c("E-Leanring","FaceToFace"))

chisq_retest(c(4,10, 78, 46, 6, 0,5,79,36,2),
  ncol=2, simulate.p.value=T )

chisq_retest(c(4,10, 78, 46, 6, 0,5,90,36,2),
  ncol=2 ,simulate.p.value=T)

chisq_retest(c(4,10, 78, 46, 6, 0,5,79,36,2), ncol=2,
  colname = c("experience", "Non-experience"),
  rowname = c("maxlow","low","mid","high","maxhigh") )


##  Comparison of gender learning methods
# --> An invalid value was calculated.
## chisq = 5.238, p =0.015
# But recalculate chisq = 4.69, p = 0.0303
chisq_retest(c(84, 60, 54, 68), ncol=2 ,
             colname = c("experience", "Non-experience"),
                          rowname = c("male", "female"))

## Comparing Learner Characteristics by Grade Level, Mina Choi (2007)

chisq_retest(c(18,38, 40, 48, 52,35,19,16), ncol=2 ,
colname = c("experience", "Non-experience"),
rowname = c("1학년","2학년","3학년","4학년")


##If the value entered is not a vector value, but a value calculated with table(),
# it is automatically calculated with the ....
table(mtcars$am, mtcars$vs) %>% Chisq_retest()


## End(Not run)


\end{ExampleCode}
\end{Examples}
\inputencoding{utf8}
\HeaderA{chisq\_table\_sum}{chisq table resize}{chisq.Rul.table.Rul.sum}
%
\begin{Description}
chisq table resize
\end{Description}
%
\begin{Usage}
\begin{verbatim}
chisq_table_sum(data, ..., trans = FALSE, sep = "_")
\end{verbatim}
\end{Usage}
%
\begin{Arguments}
\begin{ldescription}
\item[\code{data}] mat

\item[\code{...}] con 1,1,2,2

\item[\code{trans}] matr transpose

\item[\code{sep}] sep rowname combine  "-"
\end{ldescription}
\end{Arguments}
%
\begin{Value}
mat
\end{Value}
%
\begin{Examples}
\begin{ExampleCode}

## Not run: 
#'
# Test data
data <- data.frame(T = c(30, 40, 60, 120),
                   S = c(50, 70, 60, 80))
rownames(data) <- LETTERS[1:4]
data <- as.matrix(data)

# Example usage
print(chisq_table_sum(data, 1, 1, 2, 2))

chisq_table_sum(data, 1, 1, 2, 2)
chisq_table_sum(data, con = c(1, 1, 2, 2))
chisq_table_sum(data, la = c(1, 1, 2, 2))


## End(Not run)



\end{ExampleCode}
\end{Examples}
\inputencoding{utf8}
\HeaderA{chisq\_test}{chisq\_test}{chisq.Rul.test}
%
\begin{Description}
chisq\_test
\end{Description}
%
\begin{Usage}
\begin{verbatim}
chisq_test(
  x,
  y = NULL,
  correct = TRUE,
  p = rep(1/length(x), length(x)),
  rescale.p = FALSE,
  simulate.p.value = FALSE,
  B = 2000
)
\end{verbatim}
\end{Usage}
%
\begin{Arguments}
\begin{ldescription}
\item[\code{x}] a numeric vector or matrix. x and y can also both be factors.

\item[\code{y}] a numeric vector; ignored if x is a matrix. If x is a factor,
y should be a factor of the same length.

\item[\code{correct}] a logical indicating whether to apply continuity correction when computing
the test statistic for 2 by 2 tables: one half is subtracted from all ∣O−E∣ differences
; however, the correction will not be bigger than the differences themselves.
No correction is done if simulate.p.value = TRUE.

\item[\code{p}] a vector of probabilities of the same length as x.
An error is given if any entry of p is negative.

\item[\code{rescale.p}] a logical scalar; if TRUE then p is rescaled (if necessary) to sum to 1.
If rescale.p is FALSE, and p does not sum to 1, an error is given.

\item[\code{simulate.p.value}] a logical indicating whether to compute p-values by Monte Carlo simulation

\item[\code{B}] an integer specifying the number of replicates used in the Monte Carlo test.
\end{ldescription}
\end{Arguments}
%
\begin{Value}
A list with class "htest" containing the following components:
\end{Value}
%
\begin{Examples}
\begin{ExampleCode}
## Not run: 

## From Agresti(2007) p.39
M <- as.table(rbind(c(762, 327, 468), c(484, 239, 477)))
dimnames(M) <- list(gender = c("F", "M"),
party = c("Democrat","Independent", "Republican"))
(Xsq <- chisq.test(M))  # Prints test summary
Xsq$observed   # observed counts (same as M)
Xsq$expected   # expected counts under the nul
Xsq$residuals  # Pearson residuals
Xsq$stdres     # standardized residuals


x <- matrix(c(12, 5, 7, 7), ncol = 2)
chisq_test(x)$p.value           # 0.4233
chisq_test(x, simulate.p.value = TRUE, B = 10000)$p.value

## End(Not run)
\end{ExampleCode}
\end{Examples}
\inputencoding{utf8}
\HeaderA{chisq\_test\_By1}{one sample chisq test by group}{chisq.Rul.test.Rul.By1}
%
\begin{Description}
one sample chisq test by group
\end{Description}
%
\begin{Usage}
\begin{verbatim}
chisq_test_By1(dataset, sel = NULL, correct = FALSE, simulate.p.value = FALSE)
\end{verbatim}
\end{Usage}
%
\begin{Arguments}
\begin{ldescription}
\item[\code{dataset}] data.frame

\item[\code{sel}] select column group

\item[\code{correct}] chisq correct

\item[\code{simulate.p.value}] simulate.p.value
\end{ldescription}
\end{Arguments}
%
\begin{Value}
Chi-square results by group
\end{Value}
%
\begin{Examples}
\begin{ExampleCode}
## Not run: 
mtcars%>%chisq_test_By1("cyl")

mtcars%>%chisq_test_By1("am")
iris%>%chisq_test_By1("Species")

## End(Not run)
\end{ExampleCode}
\end{Examples}
\inputencoding{utf8}
\HeaderA{chisq\_test\_By2}{twosample chisq.test by group}{chisq.Rul.test.Rul.By2}
%
\begin{Description}
twosample chisq.test by group
\end{Description}
%
\begin{Usage}
\begin{verbatim}
chisq_test_By2(
  data,
  v1,
  v2,
  sel,
  correct = TRUE,
  simulate.p.value = FALSE,
  warning = FALSE,
  type = "all"
)
\end{verbatim}
\end{Usage}
%
\begin{Arguments}
\begin{ldescription}
\item[\code{data}] data,frame

\item[\code{v1}] first variable

\item[\code{v2}] second variable

\item[\code{sel}] group

\item[\code{correct}] correct

\item[\code{simulate.p.value}] simulate.p.value

\item[\code{warning}] warning

\item[\code{type}] "all", "res"
\end{ldescription}
\end{Arguments}
%
\begin{Value}
chisq result
\end{Value}
%
\begin{Examples}
\begin{ExampleCode}
## Not run: 

mtcars%>%  chisq_test_By2("am", "vs", "cyl")

mtcars%>% chisq_test_By2("am", "vs", sel= "cyl")

Mtcars <- mtcars  # 데이터를 원본과 겹치지 않도록 복사
Mtcars$am = factor(Mtcars$am, levels=c(0,1), labels= c("automatic","manual" ))
Mtcars$vs = factor(Mtcars$vs, levels=c(0,1), labels= c("V-shaped","straight" ))
Mtcars$cyl = factor(Mtcars$cyl, levels=c(4,6,8),
             labels= c("cyl-4","cyl-6","cyl-8"))

chisq_test_By2(data = Mtcars, v1 = "am", v2 = "vs", sel = "cyl")


## End(Not run)




\end{ExampleCode}
\end{Examples}
\inputencoding{utf8}
\HeaderA{chisq\_test\_By3}{그룹별 카이제곱 분석}{chisq.Rul.test.Rul.By3}
%
\begin{Description}
그룹별 카이제곱 분석
\end{Description}
%
\begin{Usage}
\begin{verbatim}
chisq_test_By3(data, grp, formula, raw = FALSE)
\end{verbatim}
\end{Usage}
%
\begin{Arguments}
\begin{ldescription}
\item[\code{data}] data

\item[\code{grp}] group variable

\item[\code{formula}] \textasciitilde{} vq +v2

\item[\code{raw}] warning
\end{ldescription}
\end{Arguments}
%
\begin{Value}
table
\end{Value}
%
\begin{Examples}
\begin{ExampleCode}

## Not run: 
#'
chisq_test_By3(mtcars, cyl, ~ am+vs )
chisq_test_By3(mtcars, cyl, ~ am+vs , raw=T)


## End(Not run)
\end{ExampleCode}
\end{Examples}
\inputencoding{utf8}
\HeaderA{chisq\_test\_cell}{Chi-square analysis by cell}{chisq.Rul.test.Rul.cell}
%
\begin{Description}
Chi-square analysis by cell
\end{Description}
%
\begin{Usage}
\begin{verbatim}
chisq_test_cell(data, type = "res", trans = FALSE, plot = TRUE, digits = 3)
\end{verbatim}
\end{Usage}
%
\begin{Arguments}
\begin{ldescription}
\item[\code{data}] table

\item[\code{type}] res, alll()default), chisq\_result , observed   expected,   residuals ,  chisq\_cell ,  stdres, chisq\_cell\_sig ,  cramersv ,  OE\_ratio,    OE\_sig ,   graph

\item[\code{trans}] plot transpose

\item[\code{plot}] TRUE

\item[\code{digits}] digits= 3
\end{ldescription}
\end{Arguments}
%
\begin{Value}
chisq
\end{Value}
%
\begin{Examples}
\begin{ExampleCode}

## Not run: 
#'


M <- as.table(rbind(c(762, 327, 468), c(484, 239, 477)))
dimnames(M) <- list(gender = c("F", "M"),
                    party = c("Democrat","Independent", "Republican"))
M %>% chisq_test_cell()

#notsig
examtest  = matrix(c(7,14,13,6), nrow=2,
                   dimnames= list(c("사전검사","사후검사"),
                                  c("합격","불합격")))
examtest
fisher.test(examtest)

examtest %>% chisq_test_cell()

## End(Not run)


\end{ExampleCode}
\end{Examples}
\inputencoding{utf8}
\HeaderA{chisq\_test\_cor}{kge\_chisq\_table chisq table observed/Expected table-}{chisq.Rul.test.Rul.cor}
%
\begin{Description}
kge\_chisq\_table chisq table observed/Expected table-
\end{Description}
%
\begin{Usage}
\begin{verbatim}
chisq_test_cor(
  dataset,
  v1 = "var1",
  v2 = "var2",
  title = "Table",
  type = "result",
  digits = 3,
  yadd = 0.4,
  Ncol = NULL,
  trans = FALSE,
  simple = FALSE,
  ko = TRUE,
  simulate.p.value = FALSE,
  correct = FALSE,
  size_bartext = 5,
  strip_size = 16,
  axis_size = 16,
  text_size = 13,
  xlab = "관측기대비율",
  ylab = "var2",
  cramer = "adjust"
)
\end{verbatim}
\end{Usage}
%
\begin{Arguments}
\begin{ldescription}
\item[\code{dataset}] data.frame

\item[\code{v1}] col

\item[\code{v2}] row

\item[\code{title}] title

\item[\code{type}] res2 default

\item[\code{digits}] 3

\item[\code{yadd}] up yadd

\item[\code{Ncol}] grid

\item[\code{trans}] transpose

\item[\code{simple}] * style

\item[\code{ko}] graph korean

\item[\code{simulate.p.value}] exact

\item[\code{correct}] exactt

\item[\code{size\_bartext}] graph text

\item[\code{strip\_size}] strip

\item[\code{axis\_size}] axis\_size

\item[\code{text\_size}] text\_size

\item[\code{xlab}] xlab

\item[\code{ylab}] ylab

\item[\code{cramer}] cramer cor
\end{ldescription}
\end{Arguments}
%
\begin{Value}
multiple data
\end{Value}
%
\begin{Examples}
\begin{ExampleCode}

## Not run: 
kge_chisq_table(mtcars,"am","cyl")


#chisq test result
vote_chi <- Chisq_retest(c(183, 35, 824, 50), ncol=2,
                        colname = c("Did not vote", "Voted"),
                         rowname =c("All other ages","Youngest"))

vote_chi


# generate matrix
voteMat = matrix(c(183, 35, 824, 50), ncol=2)
colnames(voteMat) = c("Did not vote", "Voted")
rownames(voteMat) = c("All other ages","Youngest")

voteMat%>% data.frame() %>%
  long_df("voted", rowname="youth")%>% #long format
  add_rows_freq() %>%     # add rows
  kge_chisq_table("voted","youth")    #chisq test


#신진영(2013)
shin2013_mat =matrix(c(14,1,23,3,4,0,9,0,13,18,89,8,11,40,46,12,3,4,22,0), ncol= 5)
# colnames(shin2013_mat)=c("지능지수가높다","학업성취도가높다","창의적사고를한다","특정분야에서뛰어난능력을보인다","모든분야에서뛰어난능력을보인다")
colnames(shin2013_mat)=c("지능지수","학업성취도","창의적사고","특정분야","모든분야")
rownames(shin2013_mat)=c("고졸","초대졸","대졸","대학원이상")
shin2013_mat

shin2013_mat %>% data.frame() %>%
  rownames_to_column("학력") %>%
  to_long(names_to = "영재성",cols=2:6) %>%
  unCount() %>%
  chisq_test_cor(v1="영재성", v2="학력")


shin2013_mat%>%  data.frame() %>%
  long_df("학력","영재성") %>%
  unCount() %>%
  chisq_test_cor("학력","영재성")


## End(Not run)
\end{ExampleCode}
\end{Examples}
\inputencoding{utf8}
\HeaderA{chisq\_test\_kge}{kge\_chisq\_table chisq table observed/Expected table-}{chisq.Rul.test.Rul.kge}
%
\begin{Description}
kge\_chisq\_table chisq table observed/Expected table-
\end{Description}
%
\begin{Usage}
\begin{verbatim}
chisq_test_kge(
  dataset,
  v1 = "var1",
  v2 = "var2",
  title = "Table",
  type = "result",
  digits = 3,
  yadd = 0.4,
  Ncol = NULL,
  trans = FALSE,
  simple = FALSE,
  ko = TRUE,
  simulate.p.value = FALSE,
  correct = FALSE,
  size_bartext = 5,
  strip_size = 16,
  axis_size = 16,
  text_size = 13,
  xlab = "관측기대비율",
  ylab = "var2",
  cramer = "adjust"
)
\end{verbatim}
\end{Usage}
%
\begin{Arguments}
\begin{ldescription}
\item[\code{dataset}] data.frame

\item[\code{v1}] col

\item[\code{v2}] row

\item[\code{title}] title

\item[\code{type}] res2 default

\item[\code{digits}] 3

\item[\code{yadd}] up yadd

\item[\code{Ncol}] grid

\item[\code{trans}] transpose

\item[\code{simple}] * style

\item[\code{ko}] graph korean

\item[\code{simulate.p.value}] exact

\item[\code{correct}] exactt

\item[\code{size\_bartext}] graph text

\item[\code{strip\_size}] strip

\item[\code{axis\_size}] axis\_size

\item[\code{text\_size}] text\_size

\item[\code{xlab}] xlab

\item[\code{ylab}] ylab

\item[\code{cramer}] cramer cor
\end{ldescription}
\end{Arguments}
%
\begin{Value}
multiple data
\end{Value}
%
\begin{Examples}
\begin{ExampleCode}

## Not run: 
kge_chisq_table(mtcars,"am","cyl")


#chisq test result
vote_chi <- Chisq_retest(c(183, 35, 824, 50), ncol=2,
                        colname = c("Did not vote", "Voted"),
                         rowname =c("All other ages","Youngest"))

vote_chi


# generate matrix
voteMat = matrix(c(183, 35, 824, 50), ncol=2)
colnames(voteMat) = c("Did not vote", "Voted")
rownames(voteMat) = c("All other ages","Youngest")

voteMat%>% data.frame() %>%
  long_df("voted", rowname="youth")%>% #long format
  add_rows_freq() %>%     # add rows
  kge_chisq_table("voted","youth")    #chisq test


#신진영(2013)
shin2013_mat =matrix(c(14,1,23,3,4,0,9,0,13,18,89,8,11,40,46,12,3,4,22,0), ncol= 5)
# colnames(shin2013_mat)=c("지능지수가높다","학업성취도가높다","창의적사고를한다","특정분야에서뛰어난능력을보인다","모든분야에서뛰어난능력을보인다")
colnames(shin2013_mat)=c("지능지수","학업성취도","창의적사고","특정분야","모든분야")
rownames(shin2013_mat)=c("고졸","초대졸","대졸","대학원이상")
shin2013_mat

shin2013_mat %>% data.frame() %>%
  rownames_to_column("학력") %>%
  to_long(names_to = "영재성",cols=2:6) %>%
  unCount() %>%
  chisq_test_kge(v1="영재성", v2="학력")


shin2013_mat%>%  data.frame() %>%
  long_df("학력","영재성") %>%
  unCount() %>%
  chisq_test_kge("학력","영재성")


## End(Not run)
\end{ExampleCode}
\end{Examples}
\inputencoding{utf8}
\HeaderA{cmh\_test}{cochran\_mantel\_haenszel\_test}{cmh.Rul.test}
%
\begin{Description}
cochran\_mantel\_haenszel\_test
\end{Description}
%
\begin{Usage}
\begin{verbatim}
cmh_test(
  x,
  y = NULL,
  z = NULL,
  alternative = c("two.sided", "less", "greater"),
  correct = TRUE,
  exact = FALSE,
  conf.level = 0.95
)
\end{verbatim}
\end{Usage}
%
\begin{Arguments}
\begin{ldescription}
\item[\code{x}] array

\item[\code{y}] NULL

\item[\code{z}] NULL

\item[\code{correct}] TRUE

\item[\code{exact}] FALSE

\item[\code{conf.level}] 0.95

\item[\code{alternativec("two.sided", }] "less", "greater")
\end{ldescription}
\end{Arguments}
%
\begin{Value}
result
\end{Value}
%
\begin{Examples}
\begin{ExampleCode}

## Not run: 

CAD <-array(c(8, 4, 10, 11, 21, 6, 9, 9), dim = c(2, 2, 2),
            dimnames = list(
              EKG = c(">=0.1 ST Dep", "< 0.1 ST Dep"),
              Response = c("Case", "Control"),
              Penicillin.Level = c("Female", "Male")))
CAD

#Example: Coronary Artery Disease
mantelhaen.test(CAD,correct=FALSE)

cmh_test(CAD,correct=FALSE)%>%tidy()


## End(Not run)
\end{ExampleCode}
\end{Examples}
\inputencoding{utf8}
\HeaderA{col2row}{column to rownames col2row}{col2row}
%
\begin{Description}
column to rownames col2row
\end{Description}
%
\begin{Usage}
\begin{verbatim}
col2row(data, col = 1)
\end{verbatim}
\end{Usage}
%
\begin{Arguments}
\begin{ldescription}
\item[\code{data}] data.frame

\item[\code{col}] Columns that should be rows
\end{ldescription}
\end{Arguments}
%
\begin{Value}
data
\end{Value}
%
\begin{Examples}
\begin{ExampleCode}

## Not run: 
#'
data_text = "Tobacco_Dependence (0,5] (5,10] (10,15] (15,20] (20,98]
No_Nicotine_Dependence 130 210 43 114 20
Nicotine_Dependence 119 267 91 254 67"
#method1
# nt = data_text%>% text2df(header=T)%>%
#         tibble::column_to_rownames("Tobacco_Dependence")
#method2
nt = data_text%>% text2df(header=T)%>% jjstat::col2row()
nt=Char2num(nt,1:5)
nt
nt %>%tibble()

#NEW data
New = data.frame(
  rownames= c('No_Nicotine_Dependence', 'Nicotine_Dependence'),

  "A" = c(130, 119),
  "B" = c(210, 267),
  "C" = c(43, 91),
  "D" = c(114, 254),
  "E" = c(20, 67)
);New
New%>%col2row()
New%>%col2row(col="B")
#mtcars
data_mtcars=mtcars%>%tibble::rownames_to_column()
data_mtcars%>%col2row()




## End(Not run)


\end{ExampleCode}
\end{Examples}
\inputencoding{utf8}
\HeaderA{collapse\_path}{Extract only the corresponding paths from the data with the given material and the given paths}{collapse.Rul.path}
%
\begin{Description}
Extract only the corresponding paths from the data with the given material and the given paths
\end{Description}
%
\begin{Usage}
\begin{verbatim}
collapse_path(data, path_string, est = "Original")
\end{verbatim}
\end{Usage}
%
\begin{Arguments}
\begin{ldescription}
\item[\code{data}] boot

\item[\code{path\_string}] path IMAG -> EXPE -> VAL

\item[\code{est}] Original, Mean.Boot
\end{ldescription}
\end{Arguments}
%
\begin{Value}
paths
\end{Value}
%
\begin{Examples}
\begin{ExampleCode}


## Not run: 
#'

# Example usage
sdata <- data.frame(
  paths = c("IMAG -> EXPE", "IMAG -> SAT", "IMAG -> LOY",
   "EXPE -> QUAL", "EXPE -> VAL", "EXPE -> SAT",
            "QUAL -> VAL", "QUAL -> SAT",
            "VAL -> SAT", "SAT -> LOY"),
  Original = c(0.578959128, 0.200724200, 0.275149576,
   0.848344408, 0.105477650, -0.002753995,
               0.676655529, 0.122144566, 0.589330650, 0.495479322),
  Mean.Boot = c(0.58232765, 0.21209094, 0.27486564,
  0.84793899, 0.10992051, -0.01008222,
                0.67780527, 0.15132543, 0.56312717, 0.50757755),
  Std.Error = c(0.04358293, 0.06056988, 0.07729000,
  0.01864573, 0.06968826, 0.06507903,
                0.07760527, 0.08440236, 0.08723293, 0.07856694),
  perc.025 = c(0.497644719, 0.107561558, 0.140641502,
   0.815794494, -0.025566471, -0.147788188,
               0.519959683, -0.008629516, 0.392664443, 0.351774130),
  perc.975 = c(0.6646198, 0.3275312, 0.4161131,
   0.8828875, 0.2545785, 0.1086856,
               0.8166486, 0.2991045, 0.7012494, 0.6556333)
)
collapse_path(sdata, path_string= "IMAG -> EXPE -> VAL -> SAT -> LOY")
## paths  Original  Std.Error
## 1 IMAG -> EXPE 0.5789591 0.04358293
## 2  EXPE -> VAL 0.1054776 0.06968826
## 3   VAL -> SAT 0.5893307 0.08723293
## 4   SAT -> LOY 0.4954793 0.07856694


#활용방법
collapse_path(sdata,
path_string= "IMAG -> EXPE -> VAL -> SAT -> LOY")%>%
find_paths()%>%
find_paths_cal()

#' 
## End(Not run)

\end{ExampleCode}
\end{Examples}
\inputencoding{utf8}
\HeaderA{collapse\_path2}{소벡테스트에 입력될 값으로 산정하여 list로 출력}{collapse.Rul.path2}
%
\begin{Description}
소벡테스트에 입력될 값으로 산정하여 list로 출력
\end{Description}
%
\begin{Usage}
\begin{verbatim}
collapse_path2(data, path_string, show = FALSE, est = "Original")
\end{verbatim}
\end{Usage}
%
\begin{Arguments}
\begin{ldescription}
\item[\code{data}] boot

\item[\code{path\_string}] paths

\item[\code{show}] message

\item[\code{est}] Original, Mean.Boot
\end{ldescription}
\end{Arguments}
%
\begin{Value}
data
\end{Value}
%
\begin{Examples}
\begin{ExampleCode}

## Not run: 
#'

# Example usage
sdata <- data.frame(
  paths = c("IMAG -> EXPE", "IMAG -> SAT", "IMAG -> LOY",
   "EXPE -> QUAL", "EXPE -> VAL", "EXPE -> SAT",
            "QUAL -> VAL", "QUAL -> SAT",
            "VAL -> SAT", "SAT -> LOY"),
  Original = c(0.578959128, 0.200724200, 0.275149576,
   0.848344408, 0.105477650, -0.002753995,
               0.676655529, 0.122144566, 0.589330650, 0.495479322),
  Mean.Boot = c(0.58232765, 0.21209094, 0.27486564,
  0.84793899, 0.10992051, -0.01008222,
                0.67780527, 0.15132543, 0.56312717, 0.50757755),
  Std.Error = c(0.04358293, 0.06056988, 0.07729000,
  0.01864573, 0.06968826, 0.06507903,
                0.07760527, 0.08440236, 0.08723293, 0.07856694),
  perc.025 = c(0.497644719, 0.107561558, 0.140641502,
   0.815794494, -0.025566471, -0.147788188,
               0.519959683, -0.008629516, 0.392664443, 0.351774130),
  perc.975 = c(0.6646198, 0.3275312, 0.4161131,
   0.8828875, 0.2545785, 0.1086856,
               0.8166486, 0.2991045, 0.7012494, 0.6556333)
)
#check 1
collapse_path2(sdata, path_string= "IMAG -> EXPE -> VAL -> SAT -> LOY")

#check 2
collapse_path2(sdata, path_string= "IMAG -> EXPE -> VAL -> SAT -> LOY", show=TRUE)
#'
# TEST

collapse_path2(sdata, path_string= "IMAG -> EXPE -> VAL -> SAT -> LOY")
collapse_path2(sdata, path_string= "IMAG -> EXPE -> VAL -> SAT -> LOY", show=TRUE)
collapse_path2(satpls_boot, path_string= "IMAG -> EXPE -> VAL -> SAT -> LOY", show=TRUE)

# 소벨테스트
sobel_test_extend(coefficients= c("IMAG -> EXPE" = 0.578959128,
                                  "EXPE -> SAT" = -0.002753995,
                                  "SAT -> LOY" = 0.495479322),
                  se_values = c(se1 = 0.04358293,
                                se2 = 0.06507903,
                                se3 = 0.0785669))


Resdata  =   collapse_path2(sdata, path_string= "IMAG -> EXPE -> VAL -> SAT -> LOY")

sobel_test_extend(coefficients= Resdata$coefficients, Resdata$se_values)


sobel_test_extend(
  coefficients= collapse_path2(sdata,
              path_string= "IMAG -> EXPE -> VAL -> SAT -> LOY")$coefficients,
  se_values =  collapse_path2(sdata,
        path_string= "IMAG -> EXPE -> VAL -> SAT -> LOY")$se_values)


sobel_test_extend(coefficients= collapse_path2(sdata,
         path_string=  "IMAG -> EXPE -> SAT -> LOY")$coefficients,
         se_values =  collapse_path2(sdata, path_string= "IMAG -> EXPE -> SAT -> LOY")$
                    se_values, show=FALSE)


## End(Not run)


\end{ExampleCode}
\end{Examples}
\inputencoding{utf8}
\HeaderA{colnames\_by\_dict}{colnames\_by\_dict: Rename only column names using a dictionary}{colnames.Rul.by.Rul.dict}
%
\begin{Description}
This function renames the column names of a data frame using a dictionary vector.
\end{Description}
%
\begin{Usage}
\begin{verbatim}
colnames_by_dict(df, dict)
\end{verbatim}
\end{Usage}
%
\begin{Arguments}
\begin{ldescription}
\item[\code{df}] A data.frame or tibble whose column names you want to rename.

\item[\code{dict}] A character vector of key-value pairs.
It can be either:
- Named vector: c("old1" = "new1", "old2" = "new2")
- Unnamed pairwise vector: c("old1", "new1", "old2", "new2")
\end{ldescription}
\end{Arguments}
%
\begin{Value}
A data.frame with renamed column names.
\end{Value}
%
\begin{Examples}
\begin{ExampleCode}
var_named <- c(
  "access","연구접근성",
  "barrier","진로장벽",
  "efficacy","연구효능감",
  "interest","진로흥미",
  "outcome","결과기대"
)

df <- data.frame(access = 1:2, barrier = 3:4, outcome = 5:6)
colnames_by_dict(df, var_named)
\end{ExampleCode}
\end{Examples}
\inputencoding{utf8}
\HeaderA{colnames\_change}{colnames\_change: Change column names partially or entirely, with optional preservation of original names}{colnames.Rul.change}
%
\begin{Description}
colnames\_change: Change column names partially or entirely, with optional preservation of original names
\end{Description}
%
\begin{Usage}
\begin{verbatim}
colnames_change(df, ...)
\end{verbatim}
\end{Usage}
%
\begin{Arguments}
\begin{ldescription}
\item[\code{df}] A data frame. The target data frame whose column names need to be changed.

\item[\code{...}] Character vector of new column names. Use an empty string (`""`) to keep the original column name for specific positions.
\end{ldescription}
\end{Arguments}
%
\begin{Value}
A data frame with updated column names. Columns corresponding to `""` in the input retain their original names.
\end{Value}
%
\begin{Examples}
\begin{ExampleCode}
## Not run: 
# Example with partial column renaming using mtcars
data(mtcars)
head(mtcars)

# Rename selected columns
mtcars_updated <- mtcars %>%
  colnames_change("Miles_per_Gallon", "Cylinders", "Displacement", "", "Weight")

head(mtcars_updated)

## End(Not run)

\end{ExampleCode}
\end{Examples}
\inputencoding{utf8}
\HeaderA{colnames\_change\_df}{colnames\_change\_df: Change column names of a data frame within a specified range}{colnames.Rul.change.Rul.df}
%
\begin{Description}
colnames\_change\_df: Change column names of a data frame within a specified range
\end{Description}
%
\begin{Usage}
\begin{verbatim}
colnames_change_df(df, range = 1:ncol(df), impute = NULL)
\end{verbatim}
\end{Usage}
%
\begin{Arguments}
\begin{ldescription}
\item[\code{df}] A data frame whose column names need to be changed.

\item[\code{range}] A vector or range indicating the columns to change. It can be numeric (e.g., 1:27) or character specifying column names (e.g., "disp:vs").

\item[\code{impute}] A character vector of new column names to assign to the specified range.
\end{ldescription}
\end{Arguments}
%
\begin{Value}
The data frame with updated column names.
\end{Value}
%
\begin{Examples}
\begin{ExampleCode}
# Example 1: Change the first 11 columns of mtcars
new_colnames <- c(
  paste0("feature_", 1:6),
  paste0("metric_", 1:5)
)
updated_mtcars <- colnames_change_df(
  df = mtcars,
  range = 1:11,
  impute = new_colnames
)

# Example 2: Change a subset of columns in mtcars by name range
new_colnames <- c(
  paste0("feature_", 1:6)   # Replace columns "disp" to "vs"
)
updated_mtcars <- colnames_change_df(
  df = mtcars,
  range = "disp:vs",
  impute = new_colnames
)

\end{ExampleCode}
\end{Examples}
\inputencoding{utf8}
\HeaderA{colnames\_to\_row}{colnames\_to\_row: Insert column names as a row in a data frame}{colnames.Rul.to.Rul.row}
%
\begin{Description}
colnames\_to\_row: Insert column names as a row in a data frame
\end{Description}
%
\begin{Usage}
\begin{verbatim}
colnames_to_row(df, row = 1)
\end{verbatim}
\end{Usage}
%
\begin{Arguments}
\begin{ldescription}
\item[\code{df}] A data frame whose column names will be inserted as a row.

\item[\code{row}] An integer specifying the row number where column names should be inserted. Default is 1.
\end{ldescription}
\end{Arguments}
%
\begin{Value}
A data frame with column names inserted as a row at the specified position.
\end{Value}
%
\begin{Examples}
\begin{ExampleCode}
## Not run: 
df <- data.frame(a = 1:3, b = 4:6, c = 7:9)
colnames_to_row(df)
colnames_to_row(df, row = 2)
mtcars
mtcars %>% nrow()
mtcars %>% colnames_to_row(1)
mtcars %>% colnames_to_row(3)
mtcars %>% colnames_to_row(32)

## End(Not run)
\end{ExampleCode}
\end{Examples}
\inputencoding{utf8}
\HeaderA{col\_capital}{col\_capital: Convert letters in a selected column to upper or lower case}{col.Rul.capital}
%
\begin{Description}
col\_capital: Convert letters in a selected column to upper or lower case
\end{Description}
%
\begin{Usage}
\begin{verbatim}
col_capital(df, col, direction = TRUE)
\end{verbatim}
\end{Usage}
%
\begin{Arguments}
\begin{ldescription}
\item[\code{df}] A data.frame. The input dataset.

\item[\code{col}] A column name (character) or index (numeric) indicating which column to transform.

\item[\code{direction}] Logical. If TRUE (default), convert to upper case. If FALSE, convert to lower case.
\end{ldescription}
\end{Arguments}
%
\begin{Value}
A data.frame with the selected column transformed to upper or lower case.
\end{Value}
%
\begin{Examples}
\begin{ExampleCode}
## Not run: 
df <- data.frame(
  구분 = c("chisq", "df", "pvalue"),
  criterian = c("Chisq", "df", "p > .05"),
  Value = c("1732.243", "269", "0")
)

# Use column name
col_capital(df, "criterian", TRUE)

# Use column index
col_capital(df, 2, FALSE)

## End(Not run)
\end{ExampleCode}
\end{Examples}
\inputencoding{utf8}
\HeaderA{col\_to\_colnames}{one select column to colnames}{col.Rul.to.Rul.colnames}
%
\begin{Description}
one select column to colnames
\end{Description}
%
\begin{Usage}
\begin{verbatim}
col_to_colnames(df, col = 1, remove = TRUE)
\end{verbatim}
\end{Usage}
%
\begin{Arguments}
\begin{ldescription}
\item[\code{df}] data frame

\item[\code{col}] column to use as new column names (default = 1)

\item[\code{remove}] logical. If TRUE, removes the selected column (default = TRUE).
\end{ldescription}
\end{Arguments}
%
\begin{Value}
data frame with updated column names
\end{Value}
%
\begin{Examples}
\begin{ExampleCode}
## Not run: 
# Example usage
df <- data.frame(
  요인 = c("에듀테크", "학습몰입", "지각된편리성", "지각된유용성", "태도", "지속사용"),
  SUX = c(NA, "0.948*", "0.933*", "0.947*", "0.902*", "0.865*"),
  FLOW = c(NA, NA, "0.951*", "0.967*", "0.929*", "0.876*"),
  PEOU = c(NA, NA, NA, "0.992*", "0.972*", "0.940*"),
  PU = c(NA, NA, NA, NA, "0.956*", "0.927*"),
  ATT = c(NA, NA, NA, NA, NA, "0.972*")
)

# col_to_colnames 실행
new_df <- col_to_colnames(df, col = 1, remove = FALSE)
print(new_df)

new_df2 <- col_to_colnames(df, col = 1, remove = TRUE)
print(new_df2)

## End(Not run)

\end{ExampleCode}
\end{Examples}
\inputencoding{utf8}
\HeaderA{col\_to\_rownames}{column to rownames , col\_to\_rownames}{col.Rul.to.Rul.rownames}
%
\begin{Description}
column to rownames , col\_to\_rownames
\end{Description}
%
\begin{Usage}
\begin{verbatim}
col_to_rownames(data, col = 1)
\end{verbatim}
\end{Usage}
%
\begin{Arguments}
\begin{ldescription}
\item[\code{data}] data.frame

\item[\code{col}] Columns that should be rows
\end{ldescription}
\end{Arguments}
%
\begin{Value}
data
\end{Value}
%
\begin{Examples}
\begin{ExampleCode}

## Not run: 
#'
data_text = "Tobacco_Dependence (0,5] (5,10] (10,15] (15,20] (20,98]
No_Nicotine_Dependence 130 210 43 114 20
Nicotine_Dependence 119 267 91 254 67"
#method1
# nt = data_text%>% text2df(header=T)%>%
#         tibble::column_to_rownames("Tobacco_Dependence")
#method2
nt = data_text%>% text2df(header=T)%>% jjstat::col2row()
nt=Char2num(nt,1:5)
nt
nt %>%tibble()

#NEW data
New = data.frame(
  rownames= c('No_Nicotine_Dependence', 'Nicotine_Dependence'),

  "A" = c(130, 119),
  "B" = c(210, 267),
  "C" = c(43, 91),
  "D" = c(114, 254),
  "E" = c(20, 67)
);New
New%>%col_to_rownames()

#mtcars
data_mtcars=mtcars%>%tibble::rownames_to_column()
data_mtcars%>%col_to_rownames()




## End(Not run)


\end{ExampleCode}
\end{Examples}
\inputencoding{utf8}
\HeaderA{combine\_data}{A function that combines two matrices of data, for example, to combine observation/expectation tables and p-values.}{combine.Rul.data}
%
\begin{Description}
A function that combines two matrices of data, for example, to combine observation/expectation tables and p-values.
\end{Description}
%
\begin{Usage}
\begin{verbatim}
combine_data(mat_observed_expected, mat_p_value, left = " ", right = "")
\end{verbatim}
\end{Usage}
%
\begin{Arguments}
\begin{ldescription}
\item[\code{mat\_observed\_expected}] matrix  observed/expected

\item[\code{mat\_p\_value}] matirx pvalue

\item[\code{left}] "(" etc input

\item[\code{right}] ")"
\end{ldescription}
\end{Arguments}
%
\begin{Value}
matrix
\end{Value}
%
\begin{Examples}
\begin{ExampleCode}
 ## Not run: 

# Example data
observed_expected <- matrix(c(1.0587162, 1.1477347, 0.50877193,
                              1.1269248, 0.5743712, 1.94957983,
                              0.9163776, 1.4060606, 0.03515152,
                              0.9530237, 0.8624314, 1.44531611),
                            nrow = 4, byrow = TRUE,
         dimnames = list(c("공명음", "마찰음", "유기음_경음", "평파열음_평파찰음"),
                           c("H", "H(H)", "L")))

p_value <- matrix(c(0, 0.00001, 0.00000,
                    0, 0.00000, 0.01497,
                    0, 0.00013, 0.00000,
                    0, 0.00000, 0.00102),
                  nrow = 4, byrow = TRUE,
       dimnames = list(c("공명음", "마찰음", "유기음_경음", "평파열음_평파찰음"),
                      c("H", "H(H)", "L")))

sig_sym = add_significance_symbols(p_value)

# Combine the data
combined_data1 <- combine_data(observed_expected, p_value)
 print(combined_data1)
combined_data2 <- combine_data(observed_expected, sig_sym)
 print(combined_data2)

#iris cor
iris_cor = iris %>% select(-Species) %>%psych::corr.test()
iris_cor$r
iris_cor$p

combine_data(iris_cor$r%>%round(2),
             iris_cor$p%>%round(5),
             left="(", right=")")
 sig = add_significance_symbols(iris_cor$p)
 sig

 combine_data( round(iris_cor$r, 3) , sig)

 combine_data( round(iris_cor$r, 3) , sig) %>%
 lowerMat(diag = 1, fill="")

## End(Not run)


\end{ExampleCode}
\end{Examples}
\inputencoding{utf8}
\HeaderA{combine\_rows}{중복된 데이터를 합쳐주는 함수}{combine.Rul.rows}
%
\begin{Description}
중복된 데이터를 합쳐주는 함수
\end{Description}
%
\begin{Usage}
\begin{verbatim}
combine_rows(df, rows = NULL, name_col = 1, value_col = 2, prop_col = 3)
\end{verbatim}
\end{Usage}
%
\begin{Arguments}
\begin{ldescription}
\item[\code{df}] df

\item[\code{rows}] rows =c(1,2,3) 지정된 행을 합침

\item[\code{name\_col}] 이름 열 숫자로 입력 가능

\item[\code{value\_col}] 빈도 열

\item[\code{prop\_col}] 비율 열
\end{ldescription}
\end{Arguments}
%
\begin{Value}
df
\end{Value}
%
\begin{Examples}
\begin{ExampleCode}
## Not run: 
#'
# 예제 데이터
library(dplyr)

df <- data.frame(
  term = c("지평선중", "한국게임과학고", "가천초", "감곡초", "감곡초",
           "강호항공고교", "강호항공고", "개야도초", "고", "고부초",
           "고산고", "고산고"),
  Freq = c(1, 1, 1, 2, 1, 1, 1, 1, 4, 1, 1, 2),
  prop.. = c(0.1237624, 0.1237624, 0.1237624, 0.2475248, 0.1237624,
             0.1237624, 0.1237624, 0.1237624, 0.4950495, 0.1237624,
             0.1237624, 0.2475248)
)

# 함수 호출: 디폴트 값 사용 (열 번호)
result <- combine_rows(df)

# 결과 출력
print(result)

# 함수 호출: 열 이름 사용
result_named <- combine_rows(df, name_col = "term", value_col = "Freq", score_col = "prop..")
print(result_named)

#'

## End(Not run)
\end{ExampleCode}
\end{Examples}
\inputencoding{utf8}
\HeaderA{comma}{number comma}{comma}
%
\begin{Description}
number comma
\end{Description}
%
\begin{Usage}
\begin{verbatim}
comma(x, digits = 2, big_mark = ",")
\end{verbatim}
\end{Usage}
%
\begin{Arguments}
\begin{ldescription}
\item[\code{x}] vector

\item[\code{digits}] digits 2

\item[\code{big\_mark}] "," is default
\end{ldescription}
\end{Arguments}
%
\begin{Value}
include comma
\end{Value}
%
\begin{Examples}
\begin{ExampleCode}

## Not run: 
comma(3452345)
comma(.12358124331)

## End(Not run)
\end{ExampleCode}
\end{Examples}
\inputencoding{utf8}
\HeaderA{compare\_bar}{plspm\_compare\_bar}{compare.Rul.bar}
%
\begin{Description}
plspm\_compare\_bar
\end{Description}
%
\begin{Usage}
\begin{verbatim}
compare_bar(
  data,
  x = NULL,
  y = NULL,
  fill = NULL,
  size = 3,
  size_x = 12,
  angle = 90,
  position = "dodge",
  xlab = "",
  ylab = "Values",
  text = TRUE,
  digits = 3
)
\end{verbatim}
\end{Usage}
%
\begin{Arguments}
\begin{ldescription}
\item[\code{data}] arragne data

\item[\code{x}] paths

\item[\code{y}] values

\item[\code{fill}] group

\item[\code{size}] bargrahp text size

\item[\code{size\_x}] textsize

\item[\code{angle}] xangle

\item[\code{position}] dodge, stack

\item[\code{xlab}] default ""

\item[\code{ylab}] values

\item[\code{text}] hide bartext
\end{ldescription}
\end{Arguments}
%
\begin{Value}
graph
\end{Value}
%
\begin{Examples}
\begin{ExampleCode}

## Not run: 
#'
# 예시 데이터
REs <- tibble::tribble(
  ~paths, ~grp, ~values,
  "자기효능감->진로동기", "group.남자", 0.902,
  "자기효능감->진로동기", "group.여자", 0.866,
  "자기효능감->진로태도", "group.남자", -0.639,
  "자기효능감->진로태도", "group.여자", -0.672,
  "자기효능감->진로준비", "group.남자", 0.149,
  "자기효능감->진로준비", "group.여자", 0.487,
  "진로동기->진로준비", "group.남자", 0.63,
  "진로동기->진로준비", "group.여자", 0.432,
  "진로태도->진로준비", "group.남자", -0.153,
  "진로태도->진로준비", "group.여자", 0.059
)

# # 함수 호출 예시
REs %>% compare_bar()
REs %>% compare_bar(text=FALSE)





## End(Not run)
\end{ExampleCode}
\end{Examples}
\inputencoding{utf8}
\HeaderA{compare\_bind\_data}{Table-building functions that pairwise compare two pieces of data}{compare.Rul.bind.Rul.data}
%
\begin{Description}
Table-building functions that pairwise compare two pieces of data
\end{Description}
%
\begin{Usage}
\begin{verbatim}
compare_bind_data(dff1, dff2, digits = 3, pair = TRUE)
\end{verbatim}
\end{Usage}
%
\begin{Arguments}
\begin{ldescription}
\item[\code{dff1}] data1

\item[\code{dff2}] data2

\item[\code{digits}] rounding

\item[\code{pair}] pari option default TRUE
\end{ldescription}
\end{Arguments}
%
\begin{Value}
table
\end{Value}
\inputencoding{utf8}
\HeaderA{confint\_mark\_sig}{Confint mark significant}{confint.Rul.mark.Rul.sig}
%
\begin{Description}
Confint mark significant
\end{Description}
%
\begin{Usage}
\begin{verbatim}
confint_mark_sig(data, unite = FALSE, digits = 3)
\end{verbatim}
\end{Usage}
%
\begin{Arguments}
\begin{ldescription}
\item[\code{data}] data.frame Confidence interval significance test: When the confidence interval is in the last column.  Significance test when the last two columns of data are confidence intervals

\item[\code{unite}] TRUE is combine CI and sig

\item[\code{digits}] default 3
\end{ldescription}
\end{Arguments}
%
\begin{Examples}
\begin{ExampleCode}

## Not run: 

lm(mpg ~ hp + wt, data=mtcars)%>%
 confint()%>%as.data.frame()%>%
  rownames_to_column()%>%
  confint_mark_sig()

## combind data unite=TRUE
  lm(mpg ~ hp + wt, data=mtcars)%>%
     confint()%>%
     as.data.frame()%>%
     rownames_to_column()%>%
     confint_mark_sig(unite=TRUE, digits=2)



  
## End(Not run)

\end{ExampleCode}
\end{Examples}
\inputencoding{utf8}
\HeaderA{confint\_plot\_data}{Plot Confidence Intervals from a Summary Data Frame}{confint.Rul.plot.Rul.data}
%
\begin{Description}
Draws horizontal error bars using data with estimate and confidence interval columns.
\end{Description}
%
\begin{Usage}
\begin{verbatim}
confint_plot_data(
  data,
  var_named = NULL,
  size_axis = 14,
  size_text = 4,
  linewidth = 2
)
\end{verbatim}
\end{Usage}
%
\begin{Arguments}
\begin{ldescription}
\item[\code{data}] A data.frame or tibble containing path analysis results with columns: `Path`, `est`, `ci.lower`, `ci.upper`. Optionally `sig`.

\item[\code{var\_named}] Optional named vector to rename `Path` labels.

\item[\code{size\_axis}] Font size for axis text. Default is 14.

\item[\code{size\_text}] Font size for estimate labels. Default is 4.

\item[\code{linewidth}] Line width of confidence interval bars. Default is 2.
\end{ldescription}
\end{Arguments}
%
\begin{Value}
A ggplot object showing the confidence intervals.
\end{Value}
%
\begin{Examples}
\begin{ExampleCode}
## Not run: 
df <- sem_effect_ie(fit_le)
confint_plot_data(df)

# With renaming
named_vec <- c("IE.경험_효능감_진로의지" = "경험→효능감→진로의지")
confint_plot_data(df, var_named = named_vec)

## End(Not run)
\end{ExampleCode}
\end{Examples}
\inputencoding{utf8}
\HeaderA{confint\_plot\_lm}{Plot Confidence Intervals for Linear Model or Data Frame}{confint.Rul.plot.Rul.lm}
%
\begin{Description}
Plot Confidence Intervals for Linear Model or Data Frame
\end{Description}
%
\begin{Usage}
\begin{verbatim}
confint_plot_lm(
  lm_data,
  input = c("lm", "data.frame"),
  size_text = 12,
  size_title = 14,
  color = "steelblue",
  linewidth = 1,
  alpha = 1,
  intercept = FALSE,
  errorbar_height = 0.3,
  x.title = "Estimate",
  y.title = "Item",
  type = "g"
)
\end{verbatim}
\end{Usage}
%
\begin{Arguments}
\begin{ldescription}
\item[\code{lm\_data}] An object of class 'lm' (linear model) or a data.frame containing results with estimates and confidence intervals.

\item[\code{input}] Character, one of 'lm' or 'data.frame'. Determines how lm\_data is handled.

\item[\code{size\_text}] Numeric, size of text in the plot (default = 12).

\item[\code{size\_title}] Numeric, size of axis titles in the plot (default = 14).

\item[\code{color}] Character, color of error bars (default = "steelblue").

\item[\code{linewidth}] Numeric, thickness of error bars (default = 1).

\item[\code{alpha}] Numeric, transparency of error bars (default = 1).

\item[\code{intercept}] Logical, whether to include the intercept in the plot (default = FALSE).

\item[\code{errorbar\_height}] Numeric, height of error bar ends (default = 0.3).

\item[\code{x.title}] Character, label for the x-axis (default = "Estimate").

\item[\code{y.title}] Character, label for the y-axis (default = "Item").

\item[\code{type}] Character, output type ('g' for plot, 'res' for results, 'df' for tidy dataframe).
\end{ldescription}
\end{Arguments}
%
\begin{Value}
A ggplot object, results list, or tidy dataframe based on `type`.
\end{Value}
%
\begin{Examples}
\begin{ExampleCode}
## Not run: 
# Example with lm object
model <- lm(mpg ~ wt + hp, data = mtcars)
confint_plot_lm(model, input = "lm", errorbar_height = 0.1, size_title = 16)

# Example with data.frame
df <- data.frame(term = c("wt", "hp"), estimate = c(-5.344, -0.018),
                 conf.low = c(-7.344, -0.032), conf.high = c(-3.344, -0.004))
confint_plot_lm(df, input = "data.frame", errorbar_height = 0.3,
                x.title = "Beta Coefficients", y.title = "Predictors")

## End(Not run)
\end{ExampleCode}
\end{Examples}
\inputencoding{utf8}
\HeaderA{confint\_plot\_sem}{Plot Confidence Intervals for SEM Paths using lavaan}{confint.Rul.plot.Rul.sem}
%
\begin{Description}
Create a confidence interval plot of structural paths from a lavaan SEM object.
\end{Description}
%
\begin{Usage}
\begin{verbatim}
confint_plot_sem(
  lavaan_obj,
  type = "g",
  effect = "~",
  var_named = NULL,
  size_axis = 14,
  size_text = 4,
  linewidth = 2
)
\end{verbatim}
\end{Usage}
%
\begin{Arguments}
\begin{ldescription}
\item[\code{lavaan\_obj}] A fitted `lavaan` object (e.g., from `sem()` or `cfa()`).

\item[\code{type}] Type of output: `"g"` for the ggplot object, `"res"` for the result table, `"all"` for both.

\item[\code{effect}] Operator type to filter by, default is `"\textasciitilde{}"` for regression paths. Use `"=\textasciitilde{}"` for latent loadings.

\item[\code{var\_named}] Named vector for renaming variables using `replace\_df\_rep()`. Default is `NULL`.

\item[\code{size\_axis}] Numeric. Font size for axis labels in the plot. Default is 14.

\item[\code{size\_text}] Numeric. Size of text showing estimates in the plot. Default is 4.

\item[\code{linewidth}] Numeric. Width of confidence interval error bars. Default is 2.
\end{ldescription}
\end{Arguments}
%
\begin{Value}
A ggplot object, a data frame, or a list depending on the `type` argument.
\end{Value}
%
\begin{Examples}
\begin{ExampleCode}
## Not run: 
library(lavaan)
model <- '
  visual  =~ x1 + x2 + x3
  textual =~ x4 + x5 + x6
  speed   =~ x7 + x8 + x9
  speed ~ visual + textual
'
fit <- sem(model, data = HolzingerSwineford1939)
named_vec <- c("visual" = "시각", "textual" = "언어", "speed" = "속도")
confint_plot_sem(fit, type = "g", effect = "~", var_named = named_vec)

## End(Not run)
\end{ExampleCode}
\end{Examples}
\inputencoding{utf8}
\HeaderA{contrast\_msdn}{Contrast test with mean, standard deviation, and number of samples}{contrast.Rul.msdn}
%
\begin{Description}
Contrast test with mean, standard deviation, and number of samples
\end{Description}
%
\begin{Usage}
\begin{verbatim}
contrast_msdn(M, sd, n, ..., digits = 3, type = "res", meanadd = TRUE)
\end{verbatim}
\end{Usage}
%
\begin{Arguments}
\begin{ldescription}
\item[\code{M}] means M = c(26.66, 19.74, 15.10)

\item[\code{sd}] sd  sd = c(4.51, 1.45, 2.56)

\item[\code{n}] sample  n = c(11, 7, 14)

\item[\code{...}] contrast c(1, -1,  0), c(1,  0, -1), c(0,  1, -1) ....

\item[\code{digits}] default 3

\item[\code{type}] res, res\_ci

\item[\code{meanadd}] include means
\end{ldescription}
\end{Arguments}
%
\begin{Value}
conrast result
\end{Value}
%
\begin{Examples}
\begin{ExampleCode}
## Not run: 


contrast_msdn(M = c(26.66, 19.74, 15.10),
              sd = c(4.51, 1.45, 2.56),
              n = c(11, 7, 14),
              c(1, -1, 0) )

contrast_msdn(M = c(26.66, 19.74, 15.10),
              sd = c(4.51, 1.45, 2.56),
              n = c(11, 7, 14),
              c(1, -1, 0) ,
              meanadd = FALSE)

contrast_msdn(M = c(26.66, 19.74, 15.10),
              sd = c(4.51, 1.45, 2.56),
              n = c(11, 7, 14),
              c(1, -1, 0) ,
             meanadd = FALSE,
              type="res_ci")


contrast_msdn(M = c(26.66, 19.74, 15.10),
              sd = c(4.51, 1.45, 2.56),
             n = c(11, 7, 14),
              c(1, -1,  0),
              c(1,  0, -1),
              c(0,  1, -1)
)
contrast_msdn(M = c(26.66, 19.74, 15.10),
              sd = c(4.51, 1.45, 2.56),
              n = c(11, 7, 14),
              contr1 = c(1, -1,  0),
              contr2= c(1,  0, -1),
              contr3= c(0,  1, -1),
              meanadd=F
)
contrast_msdn(M = c(26.66, 19.74, 15.10),
              sd = c(4.51, 1.45, 2.56),
              n = c(11, 7, 14),
             contr1 = c(1, -1,  0),
             contr2= c(1,  0, -1),
              contr3= c(0,  1, -1),
              meanadd=FALSE, type="res_ci"
)
#error
contrast_msdn(M = c(26.66, 19.74, 15.10),
              sd = c(4.51, 1.45, 2.56),
             n = c(11, 7, 14),
              contr = list(c(1, -1,  0),
                            c(1,  0, -1),
                           c(0,  1, -1)
              ))



## End(Not run)
\end{ExampleCode}
\end{Examples}
\inputencoding{utf8}
\HeaderA{convert\_min\_max\_scaler}{비율 척도로 변환하는 함수 정의}{convert.Rul.min.Rul.max.Rul.scaler}
%
\begin{Description}
비율 척도로 변환하는 함수 정의
\end{Description}
%
\begin{Usage}
\begin{verbatim}
convert_min_max_scaler(data, columns = NULL)
\end{verbatim}
\end{Usage}
%
\begin{Arguments}
\begin{ldescription}
\item[\code{data}] data

\item[\code{columns}] colum all
\end{ldescription}
\end{Arguments}
%
\begin{Value}
df
\end{Value}
%
\begin{Examples}
\begin{ExampleCode}
## Not run: 
#'
 사용 예
data_frame: 데이터프레임 객체
변환할 열 이름을 명시적으로 지정하거나 NULL로 모든 열을 변환
scaled_data <- min_max_scaler(data_frame, columns = c("V1_1", "V2_1", "V3_1"))
print(scaled_data)

## End(Not run)
\end{ExampleCode}
\end{Examples}
\inputencoding{utf8}
\HeaderA{convert\_scale\_df}{범위를 변환하는 데이터프레임 전용 함수 정의}{convert.Rul.scale.Rul.df}
%
\begin{Description}
범위를 변환하는 데이터프레임 전용 함수 정의
\end{Description}
%
\begin{Usage}
\begin{verbatim}
convert_scale_df(
  df,
  from_min = -4,
  from_max = 4,
  to_min = 1,
  to_max = 9,
  select_col = everything(),
  remove = FALSE,
  post_fix = "_re"
)
\end{verbatim}
\end{Usage}
%
\begin{Arguments}
\begin{ldescription}
\item[\code{df}] df

\item[\code{from\_min}] from\_min

\item[\code{from\_max}] from\_max

\item[\code{to\_min}] tomin

\item[\code{to\_max}] to max

\item[\code{select\_col}] col

\item[\code{remove}] FALSE 유지

\item[\code{post\_fix}] \_re
\end{ldescription}
\end{Arguments}
%
\begin{Value}
data
\end{Value}
%
\begin{Examples}
\begin{ExampleCode}
## Not run: 

convert_scale_df(jen01, from_min = -4, from_max = 4, to_min = 1, to_max = 9,
                 select_col = contains("Q6_"), remove = FALSE, post_fix = "_re") %>%
                 str()

# 변환된 데이터 구조 확인
str(jen01_transformed)

## End(Not run)
\end{ExampleCode}
\end{Examples}
\inputencoding{utf8}
\HeaderA{Cor2cov}{cor2cov}{Cor2cov}
%
\begin{Description}
cor2cov
\end{Description}
%
\begin{Usage}
\begin{verbatim}
Cor2cov(cor_matrix, sds = rep(1, ncol(cor_matrix)))
\end{verbatim}
\end{Usage}
%
\begin{Arguments}
\begin{ldescription}
\item[\code{cor\_matrix}] cor

\item[\code{sds}] sds =c(...), default sds 1
\end{ldescription}
\end{Arguments}
%
\begin{Value}
cov
\end{Value}
%
\begin{Examples}
\begin{ExampleCode}

## Not run: 
# correlation matrix
mtcars[1:6] %>%cor()

#change cor to cov
mtcars[1:6] %>%cor() %>% cor2cov(sds = mtcars[1:6] %>% sapply(sd))

#real cov matrix
mtcars[1:6] %>%cov()


## End(Not run)


\end{ExampleCode}
\end{Examples}
\inputencoding{utf8}
\HeaderA{count\_ne}{0이 아닌 요소의 개수를 세는 함수 count\_nonzero\_elements}{count.Rul.ne}
%
\begin{Description}
0이 아닌 요소의 개수를 세는 함수 count\_nonzero\_elements
\end{Description}
%
\begin{Usage}
\begin{verbatim}
count_ne(matrix)
\end{verbatim}
\end{Usage}
%
\begin{Arguments}
\begin{ldescription}
\item[\code{matrix}] matrix
\end{ldescription}
\end{Arguments}
%
\begin{Value}
count
\end{Value}
%
\begin{Examples}
\begin{ExampleCode}

## Not run: 
#'

#  행렬
matrix2 <- matrix(c(0, 1, 2, 0, 0, 0,
                    0, 0, 0, 0, 3, 4), nrow = 6, byrow = TRUE)
matrix2
count_ne(matrix2)

# 주어진 theta 행렬

theta <- matrix(c(5, 0, 0, 0, 0, 0,
                  0, 6, 0, 0, 0, 0,
                  0, 0, 7, 0, 0, 0,
                  0, 0, 0, 8, 0, 0,
                  0, 0, 0, 0, 9, 0,
                  0, 0, 0, 0, 0, 10), nrow = 6, byrow = TRUE)
theta
count_ne(theta)

## End(Not run)

\end{ExampleCode}
\end{Examples}
\inputencoding{utf8}
\HeaderA{Cov2cor}{cov to cor}{Cov2cor}
%
\begin{Description}
cov to cor
\end{Description}
%
\begin{Usage}
\begin{verbatim}
Cov2cor(cov_matrix)
\end{verbatim}
\end{Usage}
%
\begin{Arguments}
\begin{ldescription}
\item[\code{cov\_matrix}] cov matrix
\end{ldescription}
\end{Arguments}
%
\begin{Value}
cor mat
\end{Value}
%
\begin{Examples}
\begin{ExampleCode}

## Not run: 
cov(mtcars[,c("mpg","hp")])%>%cov2cor()

## End(Not run)
\end{ExampleCode}
\end{Examples}
\inputencoding{utf8}
\HeaderA{cov\_lm\_sem}{Regression analysis with covariance matrix, for paper checking}{cov.Rul.lm.Rul.sem}
%
\begin{Description}
Regression analysis with covariance matrix, for paper checking
\end{Description}
%
\begin{Usage}
\begin{verbatim}
cov_lm_sem(
  formula = NULL,
  cov,
  n,
  sds = NULL,
  dv,
  ivs,
  method = "none",
  md = FALSE,
  lm = TRUE
)
\end{verbatim}
\end{Usage}
%
\begin{Arguments}
\begin{ldescription}
\item[\code{formula}] regression formula

\item[\code{cov}] cov matrix if your data cor, then you must sds

\item[\code{n}] sample

\item[\code{dv}] dv

\item[\code{ivs}] iv variales

\item[\code{method}] c("holm", "hochberg", "hommel", "bonferroni", "BH", "BY")

\item[\code{md}] markdown default FALSE

\item[\code{lm}] lm option TREW, FALSE sem
\end{ldescription}
\end{Arguments}
%
\begin{Value}
result table, plot, confint sig plot
\end{Value}
%
\begin{Examples}
\begin{ExampleCode}

## Not run: 
#' #method1
cov_lm_sem(formula="mpg ~ cyl+disp+hp+drat+wt" ,
n = 32, cov = cov(mtcars[1:6]))
# method 2
cov_lm_sem(dv = "mpg", ivs = c("cyl", "disp", "hp", "drat", "wt"),
n = 32, cov = cov(mtcars[1:6]))
# method 2 cor matrix
cov_lm_sem(formula="mpg ~ cyl+disp+hp+drat+wt" ,n = 32,
cov = cor(mtcars[1:6]),
           sds = mtcars[1:6]%>%sapply(sd))
#est
cov_lm(formula="mpg ~ cyl+disp+hp+drat+wt; qsec ~ mpg +vs;" ,
       n = 32, cov = cov(mtcars[1:8]),lm=FALSE)
#std
cov_lm(formula="mpg ~ cyl+disp+hp+drat+wt; qsec ~ mpg +vs;" ,
       n = 32, cov = cor(mtcars[1:8]), sds = sapply(mtcars[1:8], sd),
       lm=FALSE)

###################
#'
Songkj2024 = lav_matrix_lower2full(
  c(
    1,
    .38,   1,
    .49, .54, 1,
    .58, .33, .45, 1,
    .35, .17, .35, .60, 1,
    .22, .11, .20, .30, .28, 1,
    .46, .25, .46, .61, .67, .31, 1,
    .52, .29, .37, .57, .47, .31, .55, 1,
    -.09, -.03,-.02, -.20,-.05, .03, .00, .01, 1,
    .23, .10, .13, .07, -.03, .06, -.06, .40,-.18, 1   )
)

Songkj2024

colnames(Songkj2024)=c("학습애정열정","학습기회개방성",
                       "솔선수범독립성","효율학습자자아개념",
                       "학업적효능감","사회적효능감",
                       "자기조절효능감","경력개발역량",
                       "성별","연령")
rownames(Songkj2024)=c("학습애정열정","학습기회개방성",
                       "솔선수범독립성","효율학습자자아개념",
                       "학업적효능감","사회적효능감",
                       "자기조절효능감","경력개발역량",
                       "성별","연령")

#기술통계 데이터
songki2024_stat=  data.frame(
  var = c("학습애정열정","학습기회개방성",
          "솔선수범독립성","효율학습자자아개념",
          "학업적효능감","사회적효능감","자기조절효능감","경력개발역량",
          "성별","연령"),
  mean = c(4.14, 3.86, 4.14, 3.86, 3.61, 3.44, 3.95, 3.89, 0.65, 1.70),
  sd =   c(.52, .72, .70, .59, .64, .82, .54, .60, 0.48, 0.82)

);songki2024_stat

Songkj2024_cov = cor2cov(Songkj2024, sds= songki2024_stat$sd)
# Songkj2024_cov
# original
Songkj2024_model_1="
학업적효능감 ~ 학습애정열정 + 학습기회개방성 +
솔선수범독립성 + 효율학습자자아개념+ 성별 + 연령
"
Songkj2024_sem1 = sem( Songkj2024_model_1,
                       sample.cov = Songkj2024 ,
                       sample.nobs =  389  )

#cor, std
cov_lm(formula="학업적효능감 ~ 학습애정열정 +
학습기회개방성 + 솔선수범독립성 + 효율학습자자아개념 +
성별 + 연령" ,n = 389, cov = Songkj2024)
# cov
cov_lm(formula="학업적효능감 ~ 학습애정열정 +
학습기회개방성 + 솔선수범독립성 + 효율학습자자아개념 +
성별 + 연령" ,n = 389, cov = Songkj2024, sds = songki2024_stat$sd)
# cov
cov_lm(formula="학업적효능감 ~ 학습애정열정 +
학습기회개방성 + 솔선수범독립성 +
효율학습자자아개념 + 성별 + 연령" ,n = 389, cov = Songkj2024_cov)

## End(Not run)

\end{ExampleCode}
\end{Examples}
\inputencoding{utf8}
\HeaderA{cov\_simdata}{Regression with covariance, generate structural equation data restoration}{cov.Rul.simdata}
%
\begin{Description}
Regression with covariance, generate structural equation data restoration
\end{Description}
%
\begin{Usage}
\begin{verbatim}
cov_simdata(
  covmat = NULL,
  n = 300,
  dv_mean = 0,
  iv_means = rep(0, ncol(covmat) - 1),
  means = NULL,
  cormat = NULL,
  sds = rep(1, ncol(cormat)),
  seed = NULL
)
\end{verbatim}
\end{Usage}
%
\begin{Arguments}
\begin{ldescription}
\item[\code{covmat}] cov matrix

\item[\code{n}] sample size default = 300

\item[\code{dv\_mean}] dv mean

\item[\code{iv\_means}] iv means

\item[\code{means}] total means default null

\item[\code{cormat}] correlation matrix

\item[\code{sds}] sds is need cor to cov

\item[\code{seed}] set seed today Sys.Data
\end{ldescription}
\end{Arguments}
%
\begin{Value}
sim data
\end{Value}
%
\begin{Examples}
\begin{ExampleCode}

## Not run: 
#'
#simdata mean setting 0
#set.seed(todaty)
cov(mtcars[,c("mpg","hp","wt")]) %>%cov_simdata(n=32)
cov(mtcars[,c("mpg","hp","wt")]) %>%cov_simdata(n=32, means=c(0,0,0))

cov(mtcars[,c("mpg","hp","wt")]) %>%
  cov_simdata(n=32, means=c(20.09, 146.69 ,3.22))

cov(mtcars[,c("mpg","hp","wt")]) %>%
  cov_simdata(n=32, dv_mean=20.09, iv_means=c(146.69 ,3.22))

#simdata
cov(mtcars[,c("mpg","hp","wt")]) %>%
  cov_simdata(n=32) %>%
  lm(formula = mpg ~ hp + wt)%>%summary()

#revised mean data Closer to real-world data
cov(mtcars[,c("mpg","hp","wt")]) %>%
  cov_simdata(n=32, dv_mean=20.09, iv_means=c(146.69 ,3.22)) %>%
  lm(formula = mpg ~ hp + wt)%>%summary()


#real data
lm(formula = mpg ~ hp +wt,data=mtcars)%>%summary()



# define sample statistics
x <- MASS::Cars93[,c("EngineSize", "Horsepower", "RPM")]
covx <- cov(x)
covx <- matrix(c(1.076122, 39.776999, -339.163745,
                 39.776999, 2743.079, 1146.634,
                 -339.163745, 1146.634, 356088.7097), nrow = 3,
               dimnames = list(c("EngineSize", "Horsepower", "RPM"),
                               c("EngineSize", "Horsepower", "RPM")))
covx
# 함수 실행
sim_data <- cov_simdata(covmat = covx, n = 100)
head(sim_data)
# check
sim_data <- cov_simdata(covmat = covx, n = 100)
lm( EngineSize~ Horsepower+RPM, data = sim_data) %>% summary()
#real check
lm( EngineSize ~ Horsepower +RPM , data=x)%>% summary()



#'

## End(Not run)
\end{ExampleCode}
\end{Examples}
\inputencoding{utf8}
\HeaderA{cramers\_v}{Tcramers\_v}{cramers.Rul.v}
%
\begin{Description}
Tcramers\_v
\end{Description}
%
\begin{Usage}
\begin{verbatim}
cramers_v(
  data,
  type = "cramer",
  digits = 3,
  correct = FALSE,
  simulate.p.value = FALSE
)
\end{verbatim}
\end{Usage}
%
\begin{Arguments}
\begin{ldescription}
\item[\code{data}] data table

\item[\code{type}] cramer

\item[\code{digits}] 3
\end{ldescription}
\end{Arguments}
%
\begin{Value}
cramer correlation
\end{Value}
\inputencoding{utf8}
\HeaderA{cronbach\_alpha}{Model-based Cronbach's Alpha from Correlation Matrix}{cronbach.Rul.alpha}
%
\begin{Description}
This function computes Cronbach's alpha using a correlation matrix.
The correlation matrix can be directly provided or extracted from a
fitted lavaan object (e.g., CFA or SEM results).
\end{Description}
%
\begin{Usage}
\begin{verbatim}
cronbach_alpha(x, type = c("ov", "lv"), digits = 3)
\end{verbatim}
\end{Usage}
%
\begin{Arguments}
\begin{ldescription}
\item[\code{x}] A correlation matrix or a fitted lavaan object.

\item[\code{type}] Character. If \code{x} is a lavaan object, choose which
correlation matrix to use:
\begin{description}

\item["ov"] Observed variable correlation matrix (default).
\item["lv"] Latent variable correlation matrix.

\end{description}


\item[\code{digits}] Integer. Number of digits to round the result.
\end{ldescription}
\end{Arguments}
%
\begin{Details}
The formula used is:
alpha = k over (k - 1) ( 1 - trace(C) over sum(C) )

where k is the number of items and C is a correlation matrix.

This alpha should be interpreted as a model-implied internal consistency
rather than a raw data-based reliability coefficient.
\end{Details}
%
\begin{Value}
A numeric value of Cronbach's alpha.
\end{Value}
%
\begin{Examples}
\begin{ExampleCode}
## Example 1: Using a correlation matrix
R <- matrix(c(
  1.00, 0.45, 0.38,
  0.45, 1.00, 0.41,
  0.38, 0.41, 1.00
), nrow = 3, byrow = TRUE)
cronbach_alpha_model(R)

## Example 2: Using a CFA lavaan object
## Not run: 
library(lavaan)

model <- '
  F =~ x1 + x2 + x3
'
fit <- cfa(model, data = mydata)

cronbach_alpha_model(fit)

## End(Not run)

\end{ExampleCode}
\end{Examples}
\inputencoding{utf8}
\HeaderA{cronbach\_from\_cor}{Cronbach's Alpha from Model-Implied Correlation Matrix}{cronbach.Rul.from.Rul.cor}
%
\begin{Description}
This function computes Cronbach's alpha using a correlation matrix.
It is intended for use with model-implied correlation matrices
(e.g., from CFA or SEM), not raw data.
\end{Description}
%
\begin{Usage}
\begin{verbatim}
cronbach_from_cor(C)
\end{verbatim}
\end{Usage}
%
\begin{Arguments}
\begin{ldescription}
\item[\code{C}] A correlation matrix.
\end{ldescription}
\end{Arguments}
%
\begin{Details}
The formula used is:
alpha = k over (k - 1) ( 1 - sum(diag(C)) over sum(C) )

where k is the number of items and C is a correlation matrix.

This coefficient represents a model-based internal consistency estimate.
\end{Details}
%
\begin{Value}
A numeric value of Cronbach's alpha.
\end{Value}
%
\begin{Examples}
\begin{ExampleCode}
## Example 1: Simple correlation matrix
R <- matrix(c(
  1.00, 0.40, 0.35,
  0.40, 1.00, 0.42,
  0.35, 0.42, 1.00
), nrow = 3, byrow = TRUE)

cronbach_from_cor(R)

## Example 2: Model-implied correlation matrix
## Not run: 
library(lavaan)

fit <- cfa(model, data = dat)
C <- lavInspect(fit, "cor.ov")
cronbach_from_cor(C)

## End(Not run)

\end{ExampleCode}
\end{Examples}
\inputencoding{utf8}
\HeaderA{cr\_ave\_from\_lambda}{Compute Composite Reliability (CR) and AVE from standardized loadings}{cr.Rul.ave.Rul.from.Rul.lambda}
%
\begin{Description}
표준화 요인부하량(lambda)을 이용하여
합성신뢰도(CR)와 평균분산추출지수(AVE)를 계산한다.
\end{Description}
%
\begin{Usage}
\begin{verbatim}
cr_ave_from_lambda(lambda, na.rm = TRUE, check = TRUE)
\end{verbatim}
\end{Usage}
%
\begin{Arguments}
\begin{ldescription}
\item[\code{lambda}] Numeric vector.
표준화 요인부하량(standardized factor loadings).

\item[\code{na.rm}] Logical.
결측값 제거 여부. 기본값은 TRUE.

\item[\code{check}] Logical.
TRUE일 경우 입력값의 범위(|lambda| ≤ 1)를 점검하고 경고를 출력한다.
\end{ldescription}
\end{Arguments}
%
\begin{Value}
A named numeric vector with elements:
\begin{itemize}

\item{} CR: Composite Reliability
\item{} AVE: Average Variance Extracted

\end{itemize}

\end{Value}
%
\begin{References}
Fornell, C., \& Larcker, D. F. (1981).
Evaluating structural equation models with unobservable variables
and measurement error. Journal of Marketing Research, 18(1), 39–50.

Hair, J. F., Hult, G. T. M., Ringle, C. M., \& Sarstedt, M. (2019).
A Primer on Partial Least Squares Structural Equation Modeling (PLS-SEM).
\end{References}
%
\begin{Examples}
\begin{ExampleCode}
lambda <- c(0.72, 0.81, 0.76)
cr_ave_from_lambda(lambda)

lambda_bad <- c(1.05, 0.80)
cr_ave_from_lambda(lambda_bad, check = TRUE)

\end{ExampleCode}
\end{Examples}
\inputencoding{utf8}
\HeaderA{cr\_candidates\_plot}{cr\_candidates\_plot: 문항 제거 전후 CR 비교 시각화}{cr.Rul.candidates.Rul.plot}
%
\begin{Description}
cr\_candidates\_plot: 문항 제거 전후 CR 비교 시각화
\end{Description}
%
\begin{Usage}
\begin{verbatim}
cr_candidates_plot(
  df,
  latent = NULL,
  range_focus = TRUE,
  font_size = 4,
  padding = 0.04,
  v_line_show = TRUE
)
\end{verbatim}
\end{Usage}
%
\begin{Arguments}
\begin{ldescription}
\item[\code{df}] lavaan 객체 또는 cr\_raise\_candidates() 결과

\item[\code{latent}] 특정 잠재변수만 시각화 (선택)

\item[\code{range\_focus}] TRUE면 CR 변화 중심으로 x축 범위 조정

\item[\code{font\_size}] 텍스트 크기

\item[\code{padding}] 범위조절 값 (기본 0.04)

\item[\code{v\_line\_show}] 각 문항의 cr\_current 위치에 수직선 추가 여부
\end{ldescription}
\end{Arguments}
%
\begin{Value}
ggplot2 객체
\end{Value}
%
\begin{Examples}
\begin{ExampleCode}
## Not run: 
library(lavaan)

model <- '
  Visual  =~ x1 + x2 + x3
  Textual =~ x4 + x5 + x6
  Speed   =~ x7 + x8 + x9
'

fit <- cfa(model, data = HolzingerSwineford1939, std.lv = TRUE)

# CR 비교 시각화 (전체)
cr_candidates_plot(fit)

# 특정 잠재변수만 시각화
cr_candidates_plot(fit, latent = "Speed")

## End(Not run)
\end{ExampleCode}
\end{Examples}
\inputencoding{utf8}
\HeaderA{cr\_raise\_candidates}{cr\_raise\_candidates: 잠재변수별 CR 향상 후보 문항 탐색}{cr.Rul.raise.Rul.candidates}
%
\begin{Description}
lavaan 모형 결과를 기반으로 각 잠재변수별로 문항 제거 시 CR 변화량을 계산합니다.
모든 잠재변수를 자동 탐색하거나, 특정 잠재변수를 지정할 수 있습니다.
\end{Description}
%
\begin{Usage}
\begin{verbatim}
cr_raise_candidates(fit, latent = NULL, cutoff = -Inf, print_result = TRUE)
\end{verbatim}
\end{Usage}
%
\begin{Arguments}
\begin{ldescription}
\item[\code{fit}] lavaan::cfa 또는 sem 결과 객체

\item[\code{latent}] 분석 대상 잠재변수 이름(문자열). NULL이면 전체 잠재변수 자동 분석.

\item[\code{cutoff}] CR 변화량(delta)의 최소 기준. 기본값은 -Inf.

\item[\code{print\_result}] 논리값. TRUE면 변수별 요약 메시지 출력.
\end{ldescription}
\end{Arguments}
%
\begin{Value}
tibble
\end{Value}
%
\begin{Examples}
\begin{ExampleCode}
## Not run: 
## 예제: lavaan 내장 데이터와 CFA 모델 사용
library(lavaan)

# 모델 정의
model <- '
  Visual  =~ x1 + x2 + x3
  Textual =~ x4 + x5 + x6
  Speed   =~ x7 + x8 + x9
'

# CFA 적합
fit <- cfa(model, data = HolzingerSwineford1939, std.lv = TRUE)

# CR 향상 후보 탐색 (전체 변수)
cr_raise_candidates(fit)

# 특정 잠재변수만 분석
cr_raise_candidates(fit, latent = "Visual", cutoff = 0.01)

## End(Not run)
\end{ExampleCode}
\end{Examples}
\inputencoding{utf8}
\HeaderA{cut\_print}{Functions that cutoff a given value}{cut.Rul.print}
%
\begin{Description}
Functions that cutoff a given value
\end{Description}
%
\begin{Usage}
\begin{verbatim}
cut_print(
  data,
  cut = 0,
  digits = 3,
  rm_row = FALSE,
  rm_cut = FALSE,
  star = FALSE,
  each = NULL
)
\end{verbatim}
\end{Usage}
%
\begin{Arguments}
\begin{ldescription}
\item[\code{data}] data.frame

\item[\code{cut}] cutoff  default 0

\item[\code{digits}] round 3

\item[\code{rm\_row}] remove row

\item[\code{rm\_cut}] remove cutoff rows

\item[\code{star}] over cut marked star pasted

\item[\code{each}] each col using cut
\end{ldescription}
\end{Arguments}
%
\begin{Value}
data
\end{Value}
%
\begin{Examples}
\begin{ExampleCode}

## Not run: 
#'
mtcars%>% cut_print(cut = 280)
mtcars%>% cut_print(cut = 280, rm_row = T)
mtcars%>% cut_print(cut = 280, rm_row = T,rm_cut=T)

#'
# 예제 데이터
data <- data.frame(
  name = c("C_S1", "C_S2", "C_S3", "C_S4", "C_S5", "A_M1", "A_M2",
           "A_M3", "B_A2", "B_A3", "D_P1", "D_P2", "D_P3"),
  block = c("자기효능감", "자기효능감", "자기효능감", "자기효능감",
            "자기효능감", "진로동기", "진로동기", "진로동기",
            "진로태도", "진로태도", "진로준비", "진로준비", "진로준비"),
  자기효능감 = c(0.88334202, 0.82571836, 0.79357537, 0.78069220,
            0.63356092, 0.79936329, 0.75998517, 0.66930167, 0.08384875,
            0.67950207, 0.75093254, 0.70621491, 0.33677926),
  진로동기 = c(0.819634450, 0.756922172, 0.732994108, 0.617655739,
           0.507011139, 0.859130054, 0.851836536, 0.815780379, 0.006420011,
           0.699152545, 0.769048149, 0.688524060, 0.336711208),
  진로태도 = c(0.6497353, 0.6010873, 0.5346407, 0.5128086,
           0.3223576, 0.6769730, 0.6030949, 0.4615263, 0.1463340,
           0.9964532, 0.5030203, 0.6625537, 0.4041093),
  진로준비 = c(0.76921767, 0.67147592, 0.58862568, 0.63738845,
           0.50527519, 0.69871013, 0.68835534, 0.66810321, 0.03263257,
           0.69001534, 0.86708694, 0.84986139, 0.55547251)
)

# 각 열이 numeric 변수인지 확인
numeric_cols <- sapply(data, is.numeric)
print(numeric_cols)

# 함수 실행 예시
result <- cut_print(data, cut = 0.01, digits = 2, rm_row = FALSE,
                    rm_cut = TRUE, star = FALSE, each = c(0.78, 0.8, 0.8, 0.8))
print(result)

#'

## End(Not run)
\end{ExampleCode}
\end{Examples}
\inputencoding{utf8}
\HeaderA{cut\_print\_all}{overall cut}{cut.Rul.print.Rul.all}
%
\begin{Description}
overall cut
\end{Description}
%
\begin{Usage}
\begin{verbatim}
cut_print_all(data, cut = 0, imp = "", digits = 3)
\end{verbatim}
\end{Usage}
%
\begin{Arguments}
\begin{ldescription}
\item[\code{data}] data

\item[\code{cut}] cut =0.5

\item[\code{imp}] "

\item[\code{digits}] 3
\end{ldescription}
\end{Arguments}
%
\begin{Value}
data
\end{Value}
\inputencoding{utf8}
\HeaderA{dall}{Quick View Utility for data.frame or tibble}{dall}
%
\begin{Description}
This function provides a convenient way to preview or fully display
a data.frame or tibble object in a pipe-friendly workflow.
By default, it returns the first \code{n} rows of the data.
\end{Description}
%
\begin{Usage}
\begin{verbatim}
dall(data, n = 10, type = NULL)
\end{verbatim}
\end{Usage}
%
\begin{Arguments}
\begin{ldescription}
\item[\code{data}] A data.frame or tibble object.

\item[\code{n}] Integer. Number of rows to display when previewing. Default is 10.

\item[\code{type}] Character. One of \code{NULL}, \code{"head"}, or \code{"all"}.
If NULL (default), the first \code{n} rows are returned.
If \code{"all"}, the full data.frame is returned.
If \code{"head"}, the first \code{n} rows are returned explicitly.
\end{ldescription}
\end{Arguments}
%
\begin{Details}
This function is designed as a lightweight alternative to \code{View()},
optimized for console-based inspection and pipe  workflows.
The default behavior (type = NULL) prioritizes safe data previewing.
\end{Details}
%
\begin{Value}
A data.frame.
\end{Value}
%
\begin{Examples}
\begin{ExampleCode}
## Example 1: Basic preview (default n = 10)
dall(mtcars)

## Example 2: Preview first 5 rows
dall(mtcars, n = 5)

## Example 3: Explicit head option
dall(mtcars, type = "head", n = 3)

## Example 4: Return full data
dall(mtcars, type = "all")

## Example 5: Pipe-friendly usage
library(dplyr)
mtcars %>%
  dplyr::select(mpg, cyl, hp) %>%
  dall(n = 6)

\end{ExampleCode}
\end{Examples}
\inputencoding{utf8}
\HeaderA{datapaste}{exceldata to dataframe}{datapaste}
%
\begin{Description}
exceldata to dataframe
\end{Description}
%
\begin{Usage}
\begin{verbatim}
datapaste(type = "matrix", x = "clipboard", header = TRUE)
\end{verbatim}
\end{Usage}
%
\begin{Arguments}
\begin{ldescription}
\item[\code{type}] "clipboard" to table data. dataframe, tibble

\item[\code{x}] "clipboard"

\item[\code{header}] header=FALSE
\end{ldescription}
\end{Arguments}
%
\begin{Value}
paste data.frame
\end{Value}
%
\begin{Examples}
\begin{ExampleCode}
#first you must copy data (Ctrl+C)
## then excute below ctrl+Enter
df = datapaste()
df = datapaste("tibble")
df = datapaste("table")
df = datapaste("data.frame")


datapaste("datapasta")
New = data.frame(row.names = c("1", "2", "3", "4", "5", "6"),
                   응답자 = c("응답1", "응답2", "응답3", "응답4", "응답5", "응답6"),
                   role_v_inst1 = c(5, 5, 5, 5, 4, 5),
                   role_v_inst2 = c(5, 5, 5, 5, 5, 5),
                   role_v_inst3 = c(4, 4, 5, 5, 4, 4),
                   role_v_inst4 = c(5, 4, 5, 5, 3, 4))


\end{ExampleCode}
\end{Examples}
\inputencoding{utf8}
\HeaderA{datapaste\_text}{data.frame to input text}{datapaste.Rul.text}
%
\begin{Description}
data.frame to input text
\end{Description}
%
\begin{Usage}
\begin{verbatim}
datapaste_text(data, rownames = TRUE)
\end{verbatim}
\end{Usage}
%
\begin{Arguments}
\begin{ldescription}
\item[\code{data}] data.frame

\item[\code{rownames}] rownames=TRUE, FALSE -> nothing rownames
\end{ldescription}
\end{Arguments}
%
\begin{Value}
text
\end{Value}
%
\begin{Examples}
\begin{ExampleCode}

## Not run: 

# 데이터 프레임 생성
cho2013_cor <- data.frame(
  필요성 = c(1.000, 0.320, 0.266, 0.506, 0.429, 0.401),
  수업 = c(0.320, 1.000, 0.647, 0.463, 0.548, 0.554),
  교사 = c(0.266, 0.647, 1.000, 0.383, 0.491, 0.545),
  계획 = c(0.506, 0.463, 0.383, 1.000, 0.531, 0.491),
  영향 = c(0.429, 0.548, 0.491, 0.531, 1.000, 0.743),
  만족도 = c(0.401, 0.554, 0.545, 0.491, 0.743, 1.000)
)
cho2013_cor %>% datapaste_text()
#'#정성령(2014)
Jungsy2014 = "
수강료 1
강사 .165** 1
수업 .303** .144** 1
개인관리 .494** .069 .284** 1
인지도 .425** .176** .407** .424** 1
이상적영향력 .210** -.056 .212** .286** .222** 1
영감적동기화 .190** -.133** .211** .290** .230** .659** 1
지적자극 .300** -.136** .159** .402** .238** .533** .517** 1
개별적 배려 .202** -.153** .150** .250** .291** .684** .618** .530** 1
"
lav_matrix(Jungsy2014)

Jungsy2014_cor = lav_matrix(Jungsy2014)

Jungsy2014_cor %>% make_df_text()

## End(Not run)

\end{ExampleCode}
\end{Examples}
\inputencoding{utf8}
\HeaderA{DescribeBy}{DescribeBy  statistics by group}{DescribeBy}
%
\begin{Description}
DescribeBy  statistics by group
\end{Description}
%
\begin{Usage}
\begin{verbatim}
DescribeBy(
  data,
  formula,
  add_var = NULL,
  stat = FALSE,
  agg = FALSE,
  gm = FALSE,
  digits = 2
)
\end{verbatim}
\end{Usage}
%
\begin{Arguments}
\begin{ldescription}
\item[\code{data}] data.frame

\item[\code{formula}] general formula

\item[\code{add\_var}] one sample name, Used to obtain statistics for one variable

\item[\code{stat}] TRUE t.test and aov(), thes stat ="t.test", or 'aov'

\item[\code{agg}] TRUE, mnay to many variable is showed. This functuion that makes aggregate() work full

\item[\code{gm}] default FALSE, TRUE Use group data obtained from group descriptive statistics to obtain group-level descriptive statistics (means, standard deviations, etc.) and to analyze a mixed model

\item[\code{digits}] Troundings
\end{ldescription}
\end{Arguments}
%
\begin{Value}
Mea, SD, N, min, max, skew, kurt
\end{Value}
%
\begin{Examples}
\begin{ExampleCode}
## Not run: 
## Used to obtain statistics for one variable
DescribeBy(mpg ~ 1, data = mtcars)
DescribeBy(mpg ~ 1, data = mtcars, "mpg")
DescribeBy(mpg ~ 1, data = mtcars, add_var = "mpg")

##The group variable is 2 or greater
DescribeBy(mpg ~ vs, data = mtcars)
DescribeBy(mpg ~ vs+am, data = mtcars)
DescribeBy(mpg ~ vs+am+cyl, data = mtcars)

##statistic data output
DescribeBy(mpg ~ 1, data = mtcars, "mpg", stat="t.test")
DescribeBy(mpg ~ vs, data = mtcars, stat="t.test")
DescribeBy(mpg ~ vs+am, data = mtcars, stat="aov")
DescribeBy(mpg ~ vs*cyl, data = mtcars, stat="aov")

##new group mean by group
#' group_by(mtcars, am) %>%
 summarise(mpg = mean(mpg)) %>%
  summarise(mean(mpg),sd(mpg),length(mpg), min(mpg),max(mpg))

### A tibble: 1 × 5
###`mean(mpg)` `sd(mpg)` `length(mpg)` `min(mpg)` `max(mpg)`
###<dbl>     <dbl>         <int>      <dbl>      <dbl>
###  1        20.8      5.12             2       17.1       24.4

DescribeBy(mpg ~ am , data= mtcars) %>%
    rename(mpg = Mean) %>%   #change name
      DescribeBy(formula = mpg ~ 1)  # mean of mpg

## End(Not run)



\end{ExampleCode}
\end{Examples}
\inputencoding{utf8}
\HeaderA{DE\_effect}{SEM direct effect result}{DE.Rul.effect}
%
\begin{Description}
SEM direct effect result
\end{Description}
%
\begin{Usage}
\begin{verbatim}
DE_effect(model_sem, effect = "~", effect2 = NULL)
\end{verbatim}
\end{Usage}
%
\begin{Arguments}
\begin{ldescription}
\item[\code{model\_sem}] lavaan object

\item[\code{effect}] select op

\item[\code{effect2}] select op 2nd option
\end{ldescription}
\end{Arguments}
%
\begin{Value}
repot
\end{Value}
%
\begin{Examples}
\begin{ExampleCode}
## Not run: 

 #library(lavaan) # only needed once per session

models1 <- '
  # measurement model
    ind60 =~ x1 + x2 + x3
    dem60 =~ y1 + y2 + y3 + y4
    dem65 =~ y5 + y6 + y7 + y8
    # regressions
    dem60 ~ ind60
    dem65 ~ ind60 + dem60
      # residual correlations
    y1 ~~ y5
    y2 ~~ y4 + y6
    y3 ~~ y7
    y4 ~~ y8
    y6 ~~ y8
'
models2 ='
  # measurement model
    ind60 =~ x1 + x2 + x3
    dem60 =~ y1 + y2 + y3 + y4
    dem65 =~ y5 + y6 + y7 + y8
     # regressions
    dem60 ~ a1*ind60
    dem65 ~ a2*ind60 + a3*dem60
 #NEw
   ie := a1*a3
   te := a2 +a1*a3
'
fits1 <- sem(models1, data=PoliticalDemocracy)
fits2 <- sem(models2, data=PoliticalDemocracy)
summary(fits1, standardized=TRUE)
summary(fits2, standardized=TRUE)

parameterEstimates(fits1, ci=F, stand=T)
parameterEstimates(fits2, ci=F, stand=T)
DE_effect(fits1)
DE_effect(fits1, effect="=~")
DE_effect(fits1, effect="=~", effect2="~")
DE_effect(fits1, effect="~~")
DE_effect(fits1, effect=":=")

DE_effect(fits2)
DE_effect(fits2, effect="=~")
DE_effect(fits2, effect="~~")
DE_effect(fits2, effect=":=")
DE_effect(fits2, effect=":=", effect2="~")
 DE_effect(fits)%>% interpretation_de()
 # same method
 DE_effect(fits)%>% report_sem_de()


## End(Not run)


\end{ExampleCode}
\end{Examples}
\inputencoding{utf8}
\HeaderA{diagram}{diagram sem plot}{diagram}
%
\begin{Description}
diagram sem plot
\end{Description}
%
\begin{Usage}
\begin{verbatim}
diagram(
  lav_obj,
  what = "est",
  rotation = 4,
  nDigits = 2,
  residScale = 18,
  sig = TRUE,
  edge.label.cex = 1.5
)
\end{verbatim}
\end{Usage}
%
\begin{Arguments}
\begin{ldescription}
\item[\code{lav\_obj}] lavaan

\item[\code{what}] est, std, model

\item[\code{rotation}] 4

\item[\code{nDigits}] 3

\item[\code{residScale}] 18

\item[\code{sig}] FALSE

\item[\code{edge.label.cex}] 1
\end{ldescription}
\end{Arguments}
%
\begin{Value}
plot
\end{Value}
%
\begin{Examples}
\begin{ExampleCode}

## Not run: 
example(cfa)
fit %>% diagram()


## End(Not run)

\end{ExampleCode}
\end{Examples}
\inputencoding{utf8}
\HeaderA{diagram2}{diagram sem plot 2}{diagram2}
%
\begin{Description}
diagram sem plot 2
\end{Description}
%
\begin{Usage}
\begin{verbatim}
diagram2(
  lav_obj,
  whatLabels = "est",
  rotation = 2,
  nDigits = 2,
  residScale = 18,
  mar = c(2, 8, 3, 8),
  edge.label.cex = 1,
  edge.label.position = 0.5,
  sizeMan = 6,
  sizeMan2 = 4,
  sizeLat = 8,
  sizeLat2 = 6,
  border.width = 2,
  edge.width = 1.5,
  style = "lisrel",
  sig = TRUE,
  asize = 1.5,
  shapeLat = "circle",
  residuals = FALSE,
  curve = 1,
  layout = "tree",
  exoVar = TRUE,
  exoCov = TRUE,
  structual = FALSE,
  nodeLabels = NULL,
  groups = FALSE
)
\end{verbatim}
\end{Usage}
%
\begin{Arguments}
\begin{ldescription}
\item[\code{lav\_obj}] lavaan

\item[\code{whatLabels}] est, std, model(sig=FALSE)

\item[\code{rotation}] rotation= 4

\item[\code{nDigits}] 3

\item[\code{residScale}] residScale=18

\item[\code{mar}] c(2,8,3,10)

\item[\code{edge.label.cex}] edge.label.cex= 1

\item[\code{edge.label.position}] edge.label.position= 0.5

\item[\code{sizeMan}] sizeMan = 8,

\item[\code{sizeMan2}] sizeMan2 = 4,

\item[\code{sizeLat}] sizeLat = 10,

\item[\code{sizeLat2}] sizeLat2 = 6

\item[\code{border.width}] border.width 2

\item[\code{edge.width}] edge.width 1.5

\item[\code{style}] style="lisrel"

\item[\code{sig}] sig=TRUE

\item[\code{asize}] asize = 1.5

\item[\code{shapeLat}] shapeLat ="circle",ellipse

\item[\code{residuals}] residuals=TRUE

\item[\code{curve}] curve=1

\item[\code{layout}] "tree"

\item[\code{exoVar}] exoVar=TRUE,

\item[\code{exoCov}] exoCov=TRUE,

\item[\code{structual}] structual=FALSE

\item[\code{nodeLabels}] nodeLabels=NULL

\item[\code{groups}] true color
\end{ldescription}
\end{Arguments}
%
\begin{Value}
plot
\end{Value}
%
\begin{Examples}
\begin{ExampleCode}

## Not run: 

example(sem)
fit%>% diagram2()
#'
fit %>%diagram2(mar=c(3,8,3,6), sizeLat = 6, sizeMan = 6,
                edge.label.cex = 0.7)

#model
fit %>%diagram2("model", sig=F)

fit %>%diagram2("model", sig=F, residuals=F, curve=2)

fit %>%diagram2("model", sig=F, residuals=F, curve=2, rotation = 4)

#making method layout
make_layout(nc=10)
make_layout(nc=10, nr=6)
make_layout(3,6)

lay1 = make_layout(6)
lay1

edit_layout(lay1)
lay1 %>%edit_layout()
lay2 = lay1 %>%edit_layout()
lay2


## End(Not run)


\end{ExampleCode}
\end{Examples}
\inputencoding{utf8}
\HeaderA{diagram\_model}{Drawing Structural Equation Models Using the lavaan Syntax (park JH PhD)}{diagram.Rul.model}
%
\begin{Description}
Drawing Structural Equation Models Using the lavaan Syntax (park JH PhD)
\end{Description}
%
\begin{Usage}
\begin{verbatim}
diagram_model(
  model,
  residuals = FALSE,
  whatLabels = "model",
  rotation = 2,
  residScale = 12,
  layout = "tree2",
  edge.label.cex = 1,
  edge.label.position = 0.5,
  edge.color = "black",
  sizeMan = 6,
  sizeMan2 = 3,
  sizeLat = 8,
  sizeLat2 = 4,
  style = "lisrel",
  sig = FALSE,
  exoVar = FALSE,
  exoCov = TRUE,
  curve = 1.5,
  asize = 2,
  mar = c(3, 3, 3, 3),
  nDigits = 3,
  shapeLat = "circle",
  shapeMan = "rectangle",
  shapeInt = "triangle",
  sample.nobs = 100,
  border.width = 2,
  edge.width = 1.5,
  Groups = FALSE,
  group_match = "lat",
  growth = FALSE,
  structural = FALSE,
  edgeLabels = NULL,
  nodeLabels = NULL,
  fixedStyle = 1,
  color = list(lat = "white", man = "white"),
  type = "plot"
)
\end{verbatim}
\end{Usage}
%
\begin{Arguments}
\begin{ldescription}
\item[\code{model}] lavaan syntax

\item[\code{residuals}] residauls TRUE

\item[\code{whatLabels}] "model

\item[\code{rotation}] rotation=2

\item[\code{residScale}] residScale=12

\item[\code{layout}] layout="tree"

\item[\code{edge.label.cex}] 1

\item[\code{edge.label.position}] 0.5

\item[\code{edge.color}] edge.color= black

\item[\code{sizeMan}] 8

\item[\code{sizeMan2}] 4

\item[\code{sizeLat}] 10

\item[\code{sizeLat2}] 6

\item[\code{style}] lisrel

\item[\code{sig}] FALSE

\item[\code{exoVar}] FALSE

\item[\code{exoCov}] TRUE

\item[\code{curve}] 1.5

\item[\code{asize}] 1.5

\item[\code{mar}] c(2,8,3,10)

\item[\code{nDigits}] 2

\item[\code{shapeLat}] "circle, ellipse

\item[\code{shapeMan}] retangle

\item[\code{shapeInt}] triangle ,

\item[\code{sample.nobs}] sample.nobs = 100

\item[\code{border.width}] border.width=1

\item[\code{edge.width}] edge.width=1.5

\item[\code{Groups}] FALSE

\item[\code{group\_match}] "FALSE"lat\#'

\item[\code{growth}] growth= FALSE, TRUE  model.type ="grouwth"

\item[\code{structural}] structural = FALSE

\item[\code{edgeLabels}] edgeLabels=NULL

\item[\code{nodeLabels}] nodeLabels=NULL

\item[\code{fixedStyle}] fixedStyle=1 default rbase linetype

\item[\code{color}] color

\item[\code{type}] type ="plot, and res is lavaan simdata result
\end{ldescription}
\end{Arguments}
%
\begin{Value}
plot
\end{Value}
%
\begin{Examples}
\begin{ExampleCode}
## Not run: 
#'
model_test <- '
  # latent variable definitions
     ind60 =~ x1 + x2 + x3
     dem60 =~ y1 + a*y2 + b*y3 + c*y4
     dem65 =~ y5 + a*y6 + b*y7 + c*y8

  # regressions
    dem60 ~ ind60
    dem65 ~ ind60 + dem60

  # residual correlations
    # y1 ~~ y5
    # y2 ~~ y4 + y6
    # y3 ~~ y7
    # y4 ~~ y8
    # y6 ~~ y8
'

HSmodel <- ' visual  =~ x1 + x2 + x3
             textual =~ x4 + x5 + x6
             speed   =~ x7 + x8 + x9 '


m1a <- "
y1 =~ x1 + x2 +x5
y2 =~ x3 + x4 +x6
y1 ~~ y2"

diagram_model(model= HSmodel, Goups=F,  residuals=F)

diagram_model(model= model_test, style="ram", groups=F)
diagram_model(model= model_test, residuals=T)
diagram_model(model= m1a, residuals=T)


#plspm model
#model 1
model1 = "
EXPE ~ IMAG
QUAL ~ EXPE
VAL ~  QUAL+ EXPE
SAT ~ IMAG + EXPE + QUAL + VAL
LOY ~ SAT + IMAG
"
sat_path = plspm_lav2path(model1,
                          fct_order=c("IMAG", "EXPE","QUAL", "VAL", "SAT", "LOY" ))
sat_path

blocks_model ="
IMAG =~ imag1 + imag2 + imag3 + imag4 + imag5
EXPE =~ expe1 + expe2 + expe3 + expe4 + expe5
QUAL =~ qual1 + qual2 + qual3 + qual4 + qual5
VAL =~ val1 + val2 + val3 + val4
SAT =~ sat1 + sat2 + sat3 + sat4
LOY =~ loy1 + loy2 + loy3 + loy4
"
blocks_model
plspm_lav2blocks(blocks_model)

#model 그리기
paste(model1, blocks_model)%>%

  diagram_model(rotation = 1, sizeLat = 6, sizeMan = 4, sizeMan2=3,
                mar=c(2,2,2,2))

#'
#'
#'

## End(Not run)
\end{ExampleCode}
\end{Examples}
\inputencoding{utf8}
\HeaderA{Diff\_effect}{Diff effect}{Diff.Rul.effect}
%
\begin{Description}
Diff effect
\end{Description}
%
\begin{Usage}
\begin{verbatim}
Diff_effect(
  model_sem,
  effect1 = "Diff",
  effect2 = "diff",
  label = FALSE,
  ci = FALSE
)
\end{verbatim}
\end{Usage}
%
\begin{Arguments}
\begin{ldescription}
\item[\code{model\_sem}] lavaan objgect

\item[\code{effect1}] "Diff"

\item[\code{effect2}] "diff"

\item[\code{label}] label

\item[\code{ci}] confidence interval
\end{ldescription}
\end{Arguments}
%
\begin{Value}
parameterEstimates()
\end{Value}
%
\begin{Examples}
\begin{ExampleCode}

## Not run: 
#library(lavaan) # only needed once per session

models1 <- '
  # measurement model
    ind60 =~ x1 + x2 + x3
    dem60 =~ y1 + y2 + y3 + y4
    dem65 =~ y5 + y6 + y7 + y8
    # regressions
    dem60 ~ ind60
    dem65 ~ ind60 + dem60
      # residual correlations
    y1 ~~ y5
    y2 ~~ y4 + y6
    y3 ~~ y7
    y4 ~~ y8
    y6 ~~ y8

models2 ='
  # measurement model
    ind60 =~ x1 + x2 + x3
    dem60 =~ y1 + y2 + y3 + y4
    dem65 =~ y5 + y6 + y7 + y8
     # regressions
    dem60 ~ a1*ind60
    dem65 ~ a2*ind60 + a3*dem60
 #NEw
   ie := a1*a3
   te := a2 +a1*a3
   Diff := a1 - a3




fits1 <- sem(models1, data=PoliticalDemocracy)
fits2 <- sem(models2, data=PoliticalDemocracy)
summary(fits1, standardized=TRUE)
summary(fits2, standardized=TRUE)
fits2 %>% Diff_effect()

## End(Not run)
\end{ExampleCode}
\end{Examples}
\inputencoding{utf8}
\HeaderA{e}{Google translaste english  to koean}{e}
%
\begin{Description}
Google translaste english  to koean
\end{Description}
%
\begin{Usage}
\begin{verbatim}
e(
  text = "hello let's start translate english to korean.",
  slang = "en",
  elang = "ko",
  show = "translate"
)
\end{verbatim}
\end{Usage}
%
\begin{Arguments}
\begin{ldescription}
\item[\code{text}] Input text

\item[\code{slang}] english

\item[\code{elang}] korean

\item[\code{show}] 'translate', 'browser', 'all' select 3 opttions
\end{ldescription}
\end{Arguments}
%
\begin{Examples}
\begin{ExampleCode}
#eexamples code
## Not run: 

e("Translating English to Korean is the default option.")

## End(Not run)



\end{ExampleCode}
\end{Examples}
\inputencoding{utf8}
\HeaderA{EAI\_index}{EAI\_index: Compute and visualize the Eduteck Acceptance Index (EAI) for given profiles}{EAI.Rul.index}
%
\begin{Description}
EAI\_index: Compute and visualize the Eduteck Acceptance Index (EAI) for given profiles
\end{Description}
%
\begin{Usage}
\begin{verbatim}
EAI_index(
  df,
  profile = NULL,
  caption = "",
  sel_col = "Class",
  level_start = 9,
  jump = 5,
  adj = 0,
  size_label = 5,
  size_text = 5,
  color_line = "gray20",
  color_label = "gray20",
  lpos = 2.35,
  rpos = 2.65,
  print = TRUE,
  web = FALSE
)
\end{verbatim}
\end{Usage}
%
\begin{Arguments}
\begin{ldescription}
\item[\code{df}] Data frame containing the necessary columns for EAI calculation.

\item[\code{profile}] A specific profile (e.g., a category from a column) to filter the data. Default is NULL, meaning all data will be used.

\item[\code{caption}] Optional caption for the plot title. Default is an empty string.

\item[\code{sel\_col}] The column name in the data frame used for selecting profiles. Default is "Class".

\item[\code{level\_start}] Starting level for categorization in the plot. Default is 9.

\item[\code{jump}] Increment to calculate subsequent levels. Default is 5.

\item[\code{adj}] Adjustment factor for level calculations. Default is 0.

\item[\code{size\_label}] Font size for level annotations. Default is 4.

\item[\code{size\_text}] Font size for score labels in the plot. Default is 5.

\item[\code{color\_line}] Color for the dashed lines in the plot. Default is "gray20".

\item[\code{color\_label}] Color for the level annotations. Default is "gray20".

\item[\code{lpos}] Horizontal position for left-side annotations. Default is 2.35.

\item[\code{rpos}] Horizontal position for right-side annotations. Default is 2.65.

\item[\code{print}] Logical indicating whether to print the summary table. Default is TRUE.

\item[\code{web}] Logical indicating whether to print web the summary table. Default is FALSE.
\end{ldescription}
\end{Arguments}
%
\begin{Value}
A ggplot object visualizing the EAI index along with calculated levels and annotations.
\end{Value}
%
\begin{Examples}
\begin{ExampleCode}
## Not run: 
set.seed(123)
df <- data.frame(
  Class = rep(c("A", "B"), each = 10),
  EAI = rnorm(20, 15, 5),
  PU_sum = rnorm(20, 20, 4),
  PEOU_sum = rnorm(20, 18, 3),
  ATT_sum = rnorm(20, 17, 4),
  USE_sum = rnorm(20, 19, 5),
  EAI_check = rnorm(20, 16, 2)
)

g <- EAI_index(df, profile = "A", caption = "EAI Analysis", sel_col = "Class")
print(g)

## End(Not run)
\end{ExampleCode}
\end{Examples}
\inputencoding{utf8}
\HeaderA{edit\_layout}{Edit an Existing Layout Matrix}{edit.Rul.layout}
%
\begin{Description}
Edit an Existing Layout Matrix
\end{Description}
%
\begin{Usage}
\begin{verbatim}
edit_layout(data)
\end{verbatim}
\end{Usage}
%
\begin{Arguments}
\begin{ldescription}
\item[\code{data}] A matrix or data frame to be edited interactively.
\end{ldescription}
\end{Arguments}
%
\begin{Value}
A matrix or data frame representing the edited layout.
\end{Value}
%
\begin{Examples}
\begin{ExampleCode}
## Not run: 
layout <- matrix(1:9, ncol = 3)
edited_layout <- edit_layout(layout)
print(edited_layout)

## End(Not run)
\end{ExampleCode}
\end{Examples}
\inputencoding{utf8}
\HeaderA{effect1}{effec1 is the path coefficient of the structural equation.}{effect1}
%
\begin{Description}
effec1 is the path coefficient of the structural equation.
\end{Description}
%
\begin{Usage}
\begin{verbatim}
effect1(
  x,
  format = "markdown",
  op1 = "~",
  op2 = ":=",
  what = "std",
  digits = 3,
  layout = "tree2",
  curve = 1,
  mar = c(8, 8, 8, 8),
  residScale = 15,
  sizeLat = 12,
  sizeLat2 = 6,
  sizeMan = 8,
  sizeMan2 = 4,
  edge.label.cex = 0.9,
  label.cex = 1.5,
  asize = 1.4,
  name_change = FALSE,
  label_list = NULL,
  type = "res"
)
\end{verbatim}
\end{Usage}
%
\begin{Arguments}
\begin{ldescription}
\item[\code{x}] lavaan result

\item[\code{format}] 'markdown', 'html'

\item[\code{op1}] first'\textasciitilde{}'

\item[\code{op2}] first ':='  newparameters

\item[\code{what}] default 'std', Can be changeed to 'est'

\item[\code{digits}] round 3

\item[\code{layout}] 'tree2'

\item[\code{curve}] curve =1

\item[\code{mar}] margins default c(8,8,8,8)

\item[\code{residScale}] length of residuals. default 15

\item[\code{sizeLat}] Horizontal length of latent variable

\item[\code{sizeLat2}] Vertical length of latent variable

\item[\code{sizeMan}] horizontal length of the item

\item[\code{sizeMan2}] Vertical length of the item

\item[\code{edge.label.cex}] edge size

\item[\code{label.cex}] label size

\item[\code{asize}] arrow size

\item[\code{name\_change}] When set to default, the variable name of the data is used, and when set to TRUE, a new variable name can be entered. Just enter them all in label\_list using the c() function.

\item[\code{label\_list}] label list

\item[\code{type}] type output
\end{ldescription}
\end{Arguments}
%
\begin{Examples}
\begin{ExampleCode}

## Not run: 
## change label_list
abel_listnode.name =list(
  list(node="f인정욕구", to="인정신념"),
  list(node="f파국화",  to="불안신념1"),
  list(node="f과잉불안", to="불안신념2"),
  list(node="f무력감",  to="의존신념"),
  list(node="f개인완벽", to="해결신념"),
  list(node="IRB",  to="비합리적신념"),
  list(node="PREACT",  to="진로준비행동"),
  list(node="DECISION",  to="진로결정수준"),
  list(node="f정보수집",  to="정보수집"),
  list(node="f직업체험",  to="직업체험"),
  list(node="f진로탐색",  to="정보탐색"),
  list(node="f구체성",  to="구체성"),
  list(node="f확신성",  to="확신성"),
  list(node="f적합성",  to="적합성"))

## End(Not run)


\end{ExampleCode}
\end{Examples}
\inputencoding{utf8}
\HeaderA{factor\_layout}{Extracting the layout matrix of a plspm model}{factor.Rul.layout}
%
\begin{Description}
Extracting the layout matrix of a plspm model
\end{Description}
%
\begin{Usage}
\begin{verbatim}
factor_layout(plspm_model)
\end{verbatim}
\end{Usage}
%
\begin{Arguments}
\begin{ldescription}
\item[\code{plspm\_model}] sempaths
\end{ldescription}
\end{Arguments}
%
\begin{Value}
mat
\end{Value}
%
\begin{Examples}
\begin{ExampleCode}

## Not run: 
#'
data(offense)
offense %>%str()

nfl_path = matrix(
  c(0, 0, 0, 0, 0,
    0, 0, 0, 0, 0,
    0, 0, 0, 0, 0,
    0, 1, 1, 0, 0,
    1, 0, 0, 1, 0),
  nrow = 5, ncol = 5, byrow = TRUE )

rownames(nfl_path) <- c("Speical", "Rushing", "Passing", "Offense", "Scoring")
colnames(nfl_path) <- c("Speical", "Rushing", "Passing", "Offense", "Scoring")
nfl_path

nfl_blocks = list(7:8, 1:3, 4:6,1:6, 9:11)

plspm_blocks_match(data=offense, blocks=nfl_blocks,
                   paths=nfl_path)

nfl_mode =c("B","A","A","A","A")

nfl_pls = plspm(Data=offense, nfl_path, nfl_blocks, modes = nfl_mode)
nfl_pls%>%summary()

#########
nfl_pls_boot = plspm_sem(Data=offense, nfl_path, nfl_blocks, modes = nfl_mode)

#Layouts extracted from mockups
plspm_semPaths2(nfl_pls_boot,sizeLat = 6)%>%
  plspm_factor_layout()

#Make it input material
plspm_semPaths2(nfl_pls_boot,sizeLat =6)%>%
  plspm_factor_layout() %>%layout_mat(text=T)

#Organizing materials for input
plspm_semPaths2(nfl_pls_boot,sizeLat =6)%>%
  plspm_factor_layout() %>%layout_mat()%>%
  mat_text_arrange()

## End(Not run)



\end{ExampleCode}
\end{Examples}
\inputencoding{utf8}
\HeaderA{factor\_scores}{Add factor score calculation functions and data}{factor.Rul.scores}
%
\begin{Description}
Add factor score calculation functions and data
\end{Description}
%
\begin{Usage}
\begin{verbatim}
factor_scores(
  data,
  range = 1:ncol(data),
  nfactor = NULL,
  prefix = "",
  rotate = "varimax"
)
\end{verbatim}
\end{Usage}
%
\begin{Arguments}
\begin{ldescription}
\item[\code{data}] data

\item[\code{range}] range default 1:ncol(data)

\item[\code{nfactor}] number of factor

\item[\code{prefix}] data prefic

\item[\code{rotate}] varimax
\end{ldescription}
\end{Arguments}
%
\begin{Value}
data
\end{Value}
%
\begin{Examples}
\begin{ExampleCode}


## Not run: 
#'
# 활용방법 1
factor_scores(mtcars,
              range = c("disp", "hp", "drat", "wt", "qsec"),
              nfactor = 2, prefix="a")

mtcars %>% select("disp", "hp", "drat", "wt", "qsec") %>% factor_scores(nfactor=2)

# 활용방법 2
factor_scores(mtcars, range = c(3:7), nfactor = 2)

factor_scores(mtcars, c("disp", "hp", "drat", "wt", "qsec"), nfactor = 2)
factor_scores(mtcars, c(3, 4, 5, 6, 7), nfactor = 2)

factor_scores(mtcars, range = "disp:qsec", nfactor = 2)


mtcars %>% select("disp", "hp", "drat", "wt", "qsec") %>%
  factor_scores(c("disp","hp"),nfactor=2, prefix="a") %>%
  factor_scores(c("drat", "wt", "qsec"),nfactor=3, prefix="b") %>%head()


## End(Not run)
\end{ExampleCode}
\end{Examples}
\inputencoding{utf8}
\HeaderA{fa\_print}{fa\_print: A custom function to format and display Factor Analysis (FA) results}{fa.Rul.print}
%
\begin{Description}
fa\_print: A custom function to format and display Factor Analysis (FA) results
\end{Description}
%
\begin{Usage}
\begin{verbatim}
fa_print(
  data,
  cut = 0.4,
  sort = TRUE,
  abs = TRUE,
  desc = TRUE,
  star = FALSE,
  justify = "right"
)
\end{verbatim}
\end{Usage}
%
\begin{Arguments}
\begin{ldescription}
\item[\code{data}] A data frame containing the FA results to format.

\item[\code{cut}] Numeric. The threshold for suppressing small factor loadings. Defaults to 0.4.

\item[\code{sort}] Logical. If TRUE, sorts the factor loadings. Defaults to TRUE.

\item[\code{abs}] Logical. If TRUE, applies the absolute value to the factor loadings for threshold comparison. Defaults to TRUE.

\item[\code{desc}] Logical. If TRUE, sorts in descending order. Defaults to TRUE.

\item[\code{star}] Logical. If TRUE, marks factor loadings exceeding the cut-off with "*". Defaults to FALSE.

\item[\code{justify}] Character. Specifies justification ("left", "center", or "right") for formatted values. Defaults to "right".
\end{ldescription}
\end{Arguments}
%
\begin{Value}
A data frame with formatted and optionally sorted FA results, keeping "Eigenvalues",
"Proportion Var", and "Cumulative Var" unaffected by filtering and sorting.
\end{Value}
%
\begin{Examples}
\begin{ExampleCode}
## Not run: 
# Example usage with a FA result
library(psych)
results <- psych::fa(mtcars, nfactors = 2, rotate = "promax")
summary <- summary_fa(results)
formatted_results <- fa_print(summary, cut = 0.3, star = TRUE)
print(formatted_results)

## End(Not run)
\end{ExampleCode}
\end{Examples}
\inputencoding{utf8}
\HeaderA{fei\_Chisq\_test}{fei분석을위한 카이제곱함수}{fei.Rul.Chisq.Rul.test}
%
\begin{Description}
fei분석을위한 카이제곱함수
\end{Description}
%
\begin{Usage}
\begin{verbatim}
fei_Chisq_test(tabledf, col = "빈도", name = "구분", ns = FALSE, sep = " / ")
\end{verbatim}
\end{Usage}
%
\begin{Arguments}
\begin{ldescription}
\item[\code{tabledf}] 빈도분석 Freq\_table 데이터

\item[\code{col}] "빈도"

\item[\code{name}] 구분

\item[\code{ns}] FALSE 생성한 것사용 ,TRUE RVAideMemoire로 행렬로 나타남

\item[\code{sep}] sep="/"
\end{ldescription}
\end{Arguments}
%
\begin{Value}
df
\end{Value}
%
\begin{Examples}
\begin{ExampleCode}
## Not run: 
#'
eduteck_s$Q31%>% na_remover() %>% Freq_table(prop=TRUE) %>%
  fei_Chisq_test()

## End(Not run)
\end{ExampleCode}
\end{Examples}
\inputencoding{utf8}
\HeaderA{fei\_dv\_res}{T종속변수의 빈도를 세는 함수 (에듀테크 연구보고용 )}{fei.Rul.dv.Rul.res}
%
\begin{Description}
T종속변수의 빈도를 세는 함수 (에듀테크 연구보고용 )
\end{Description}
%
\begin{Usage}
\begin{verbatim}
fei_dv_res(
  df,
  title = "",
  x_size = 20,
  label_text = 5,
  x_range = 100,
  type = "res"
)
\end{verbatim}
\end{Usage}
%
\begin{Arguments}
\begin{ldescription}
\item[\code{df}] df

\item[\code{title}] 변수명

\item[\code{x\_size}] 20 사이즈로 조절

\item[\code{label\_text}] 막대그래프 끄트

\item[\code{x\_range}] 범위조절

\item[\code{type}] res all, freq, console, md, gg, graph na\_Count, cleandata
\end{ldescription}
\end{Arguments}
%
\begin{Value}
df
\end{Value}
%
\begin{Examples}
\begin{ExampleCode}
## Not run: 

eduteck_s$Q14_5 %>% fei_dv_res("자기주도 공부력 향상")



## End(Not run)
\end{ExampleCode}
\end{Examples}
\inputencoding{utf8}
\HeaderA{filter\_rows}{filter and,  or function}{filter.Rul.rows}
%
\begin{Description}
filter and,  or function
\end{Description}
%
\begin{Usage}
\begin{verbatim}
filter_rows(data, op = "or", ...)
\end{verbatim}
\end{Usage}
%
\begin{Arguments}
\begin{ldescription}
\item[\code{data}] data

\item[\code{op}] operatror or, and

\item[\code{...}] condition
\end{ldescription}
\end{Arguments}
%
\begin{Value}
data
\end{Value}
%
\begin{Examples}
\begin{ExampleCode}

## Not run: 
#'
mtcars %>% filter_rows(cyl = 6)
mtcars %>% filter_rows(vs = 1)
mtcars %>% filter_rows(cyl = 6, vs = 1, op ="and")
mtcars %>% filter_rows(op="and", cyl = 6, vs = 1)
mtcars %>% filter_rows("and", cyl = 6, vs = 1 )

mtcars %>% filter_rows("or", cyl = 6, vs = 1 )
mtcars %>% filter_rows(op= "or", cyl = 6, vs = 1 )
mtcars %>% filter_rows(cyl = 6, vs = 1, op ="or")

mtcars %>% filter_rows(cyl = 6, vs = 1, gear = 4, op = "and")
mtcars %>% filter_rows(cyl = 6, vs = 1, gear = 4, op = "or")
mtcars %>% filter_rows(cyl = 6, vs = 1, gear = 4, op = "or")

mtcars %>% filter_rows(mpg > 25) #error
mtcars %>% filter(mpg > 25)

## End(Not run)

\end{ExampleCode}
\end{Examples}
\inputencoding{utf8}
\HeaderA{find\_paths}{Resolver knows the path}{find.Rul.paths}
%
\begin{Description}
Resolver knows the path
\end{Description}
%
\begin{Usage}
\begin{verbatim}
find_paths(data, type = "paths", est = "Original", allpaths = FALSE)
\end{verbatim}
\end{Usage}
%
\begin{Arguments}
\begin{ldescription}
\item[\code{data}] boot data

\item[\code{type}] es, all, est, ind, paths(default)

\item[\code{est}] Original

\item[\code{allpaths}] allpaths TRUE all, FALSE is ind

\item[\code{start}] IMAG
\end{ldescription}
\end{Arguments}
%
\begin{Value}
result text
\end{Value}
%
\begin{Examples}
\begin{ExampleCode}

## Not run: 
#'
# 예제 데이터
satpls_boot <- list(boot = list(paths = data.frame(
  Original = c(0.578959128, 0.200724200, 0.275149576,
  0.848344408, 0.105477650, -0.002753995,
  0.676655529, 0.122144566, 0.589330650, 0.495479322),
  Mean.Boot = c(0.57466537, 0.20926883,
  0.28748175, 0.84907549, 0.09629728,
  -0.02304841, 0.68789406, 0.13300281,
  0.59112195, 0.48611617),
  Std.Error = c(0.04941293, 0.05440409,
   0.07095248, 0.02022514, 0.05731729,
   0.06209752, 0.06286102, 0.08889983,
   0.08591013, 0.07920673),
  perc.025 = c(0.46120153, 0.09494506,
  0.14255079, 0.80126496, -0.02555946,
  -0.15182035, 0.57655200, -0.01836564,
   0.41297885, 0.34267859),
  perc.975 = c(0.65528114, 0.30696127,
   0.41537966, 0.88210331, 0.19678402,
   0.09426639, 0.81986383, 0.32270662,
    0.75274086, 0.63147051),
  row.names = c("IMAG -> EXPE",
  "IMAG -> SAT",
   "IMAG -> LOY", "EXPE -> QUAL",
   "EXPE -> VAL", "EXPE -> SAT",
   "QUAL -> VAL", "QUAL -> SAT",
   "VAL -> SAT", "SAT -> LOY")
)))

# 함수 실행
find_paths_all(satpls_boot$boot$paths)
find_paths_all(satpls_boot$boot$paths,"all")

## End(Not run)

\end{ExampleCode}
\end{Examples}
\inputencoding{utf8}
\HeaderA{find\_paths\_cal}{real paths calculation est}{find.Rul.paths.Rul.cal}
%
\begin{Description}
real paths calculation est
\end{Description}
%
\begin{Usage}
\begin{verbatim}
find_paths_cal(paste_path_result, type = "all", digits = 3, est = "Original")
\end{verbatim}
\end{Usage}
%
\begin{Arguments}
\begin{ldescription}
\item[\code{paste\_path\_result}] find\_paths data

\item[\code{type}] all=res, ind,  est

\item[\code{digits}] round 3

\item[\code{est}] Original  or Mean.Boot
\end{ldescription}
\end{Arguments}
%
\begin{Value}
data
\end{Value}
%
\begin{Examples}
\begin{ExampleCode}


## Not run: 
#'

# Example usage
sdata <- data.frame(
  paths = c("IMAG -> EXPE", "IMAG -> SAT", "IMAG -> LOY",
   "EXPE -> QUAL", "EXPE -> VAL", "EXPE -> SAT",
            "QUAL -> VAL", "QUAL -> SAT",
            "VAL -> SAT", "SAT -> LOY"),
  Original = c(0.578959128, 0.200724200, 0.275149576,
   0.848344408, 0.105477650, -0.002753995,
               0.676655529, 0.122144566, 0.589330650, 0.495479322),
  Mean.Boot = c(0.58232765, 0.21209094, 0.27486564,
  0.84793899, 0.10992051, -0.01008222,
                0.67780527, 0.15132543, 0.56312717, 0.50757755),
  Std.Error = c(0.04358293, 0.06056988, 0.07729000,
  0.01864573, 0.06968826, 0.06507903,
                0.07760527, 0.08440236, 0.08723293, 0.07856694),
  perc.025 = c(0.497644719, 0.107561558, 0.140641502,
   0.815794494, -0.025566471, -0.147788188,
               0.519959683, -0.008629516, 0.392664443, 0.351774130),
  perc.975 = c(0.6646198, 0.3275312, 0.4161131,
   0.8828875, 0.2545785, 0.1086856,
               0.8166486, 0.2991045, 0.7012494, 0.6556333)
)

#'

##realdata analysis
data(satisfaction)
sat_path = plspm_paths(
  row_names = c("IMAG","EXPE","QUAL","VAL","SAT","LOY"),
  relationship = list(
    path(from="IMAG", to=c("EXPE","SAT","LOY")),
    path("EXPE", c("QUAL","VAL","SAT")),
    path("QUAL", c("VAL","SAT")),
    path("VAL",c("SAT")),
    path("SAT","LOY"))
)
sat_path

# blokcs
sat_blocks1 <- plspm_blocks(
  IMAG = item(paste0("imag",1:5)),
  EXPE = item(paste0("expe", 1:5)),
  QUAL = item( paste0("qual", 1:5)),
  VAL = item(paste0("val", 1:4)),
  SAT = item(paste0("sat", 1:4)),
  LOY = item(paste0("loy", 1:4))
)

sat_blocks1

satpls_boot = plspm_sem(satisfaction, sat_path, sat_blocks1, scaled = FALSE,
                        boot.val =TRUE, br=100)
#' find_paths(satpls_boot)
find_paths_cal(find_paths(satpls_boot,"all") )

#값을 계산하여 보기
find_paths(jutpls_boot,"all")%>%find_paths_cal()
#Mean.Boot자료 보기
find_paths(jutpls_boot,"all", est="Mean.Boot")%>%find_paths_cal()




## End(Not run)


\end{ExampleCode}
\end{Examples}
\inputencoding{utf8}
\HeaderA{format\_col}{select\_format}{format.Rul.col}
%
\begin{Description}
select\_format
\end{Description}
%
\begin{Usage}
\begin{verbatim}
format_col(data, ..., nsmall = 3, digits = 3)
\end{verbatim}
\end{Usage}
%
\begin{Arguments}
\begin{ldescription}
\item[\code{data}] data

\item[\code{...}] col selection

\item[\code{nsmall}] default 3

\item[\code{digits}] round 3
\end{ldescription}
\end{Arguments}
%
\begin{Value}
data
\end{Value}
%
\begin{Examples}
\begin{ExampleCode}


## Not run: 


# Usage example
data2 <- data.frame(a = rnorm(5), b = rnorm(5), c = rnorm(5), d = rnorm(5), e = rnorm(5))
format_col(data2, 1, 2)
format_col(data2, 1:5)
format_col(data2, a)
format_col(data2, a:e)
format_col(data2, "a")

## End(Not run)
\end{ExampleCode}
\end{Examples}
\inputencoding{utf8}
\HeaderA{format\_number}{Define a function: Convert numbers to scientific notation or common number}{format.Rul.number}
%
\begin{Description}
Define a function: Convert numbers to scientific notation or common number
\end{Description}
%
\begin{Usage}
\begin{verbatim}
format_number(number, n1 = 8, n2 = 3, n3 = 5)
\end{verbatim}
\end{Usage}
%
\begin{Arguments}
\begin{ldescription}
\item[\code{number}] number

\item[\code{n1}] length

\item[\code{n2}] scientfic

\item[\code{n3}] scientfic
\end{ldescription}
\end{Arguments}
%
\begin{Value}
formatting
\end{Value}
%
\begin{Examples}
\begin{ExampleCode}
## Not run: 
## number of data
nchar(sub("\\d+\\.", "", as.character(0.123456)))

numbers <- c(0.123456, 123.456789, 0.0000123456)
sapply(numbers, format_number)

for (i in seq_along(numbers)) {
   cat("Original:", numbers[i], "Formatted:", formatted_numbers[i], "\n")}

format_number(0.123456)


## End(Not run)

\end{ExampleCode}
\end{Examples}
\inputencoding{utf8}
\HeaderA{Freq\_rowcounter}{Freq\_rowcounter}{Freq.Rul.rowcounter}
%
\begin{Description}
Freq\_rowcounter
\end{Description}
%
\begin{Usage}
\begin{verbatim}
Freq_rowcounter(df, type = "counter")
\end{verbatim}
\end{Usage}
%
\begin{Arguments}
\begin{ldescription}
\item[\code{df}] df

\item[\code{type}] counter, term, freq\_counter, freq\_term
\end{ldescription}
\end{Arguments}
%
\begin{Value}
df
\end{Value}
%
\begin{Examples}
\begin{ExampleCode}
## Not run: 

 Freq_rowcounter(df = eduteck2, type = "counter")
Freq_rowcounter(df = eduteck2, type = "freq_counter")

## End(Not run)
\end{ExampleCode}
\end{Examples}
\inputencoding{utf8}
\HeaderA{Freq\_table}{Frequency analysis table}{Freq.Rul.table}
%
\begin{Description}
Frequency analysis table
\end{Description}
%
\begin{Usage}
\begin{verbatim}
Freq_table(
  data,
  ...,
  prop = FALSE,
  plot = FALSE,
  angle = 0,
  size_text = 4,
  size_axis = 10,
  vjust = -0.05,
  hjust = 0.5,
  yadd = 100,
  legend.position = "",
  reorder = FALSE,
  wt = NULL,
  na.rm = TRUE,
  na_treatment = "무응답",
  remove_rows = NULL,
  add_row_sum = FALSE,
  xlab = "Category",
  flip = FALSE,
  size_axis_title = 12,
  sort_bar = TRUE,
  type = "res"
)
\end{verbatim}
\end{Usage}
%
\begin{Arguments}
\begin{ldescription}
\item[\code{data}] data.frame

\item[\code{...}] data.frame col names

\item[\code{prop}] proportion value, default TRUE

\item[\code{plot}] plot defualt false

\item[\code{angle}] angle = 0

\item[\code{size\_text}] default 3, bargraph text

\item[\code{size\_axis}] default 12, axis test

\item[\code{vjust}] default -0.5

\item[\code{hjust}] default -0.5

\item[\code{yadd}] default 0.8

\item[\code{legend.position}] legend.position="" none. top, bottom, left, right

\item[\code{reorder}] reorder Sorting Graph Arrays

\item[\code{wt}] weight

\item[\code{na.rm}] TRUE

\item[\code{na\_treatment}] "무응답", \# 결측값 대체값

\item[\code{remove\_rows}] Enter the number of rows to hide and they will be removed

\item[\code{add\_row\_sum}] default FALSE, Sum excluding removed rows

\item[\code{xlab}] xalb = "Category",

\item[\code{flip}] FALSE

\item[\code{size\_axis\_title}] default 12,

\item[\code{sort\_bar}] TRUE(default)

\item[\code{type}] default  "res",  all, res=data, g=plot=graph
\end{ldescription}
\end{Arguments}
%
\begin{Value}
Freqency table
\end{Value}
%
\begin{Examples}
\begin{ExampleCode}
## Not run: 
#사용방법
eduteck_s%>% Freq_table("Q9", prop=TRUE, add_row_sum = TRUE)

eduteck_s$Q9 %>% Freq_table(prop=TRUE, add_row_sum = TRUE, na.rm= TRUE)
eduteck_s$Q9 %>% Freq_table(prop=TRUE, add_row_sum = TRUE, na.rm=FALSE)

Mtcars = mtcars
Mtcars$am = factor(Mtcars$am, levels=c(0,1), labels= c("automatic","manual" ))
Mtcars$vs  = factor(Mtcars$vs, levels=c(0,1), labels= c("V-shaped","straight" ))
Mtcars$cyl  = factor(Mtcars$cyl, levels=c(4,6,8), labels= c("cyl-4","cyl-6","cyl-8"))

#vector variables
Freq_table(c("a","b","c","a","b","a","c","c","c"))
Freq_table(c("a","b","c","a","b","a","c","c","c"), prop=TRUE)
Freq_table(c("a","b","c","a","b","a","c","c","c"), plot=T)


Mtcars %>% Freq_table("am")
Mtcars %>% Freq_table("am", prop=F)
Mtcars %>% Freq_table("vs", plot=FALSE)

Mtcars %>% Freq_table("vs")
Mtcars %>% Freq_table("vs","am")

Mtcars %>% Freq_table("vs","am","cyl", type="all")

#graph
Mtcars %>% Freq_table("vs","am","cyl", type="g")

##only graph
Mtcars %>% Freq_table("vs","am","cyl", angle=90, type="g")
Mtcars %>% Freq_table("vs","am","cyl", angle=90, reorder = TRUE, type="g")

Mtcars %>% Freq_table("vs","am","cyl", angle=90, type="g")+ylim(0,14)
Mtcars %>% Freq_table("vs","am","cyl", angle=90, reorder = TRUE, type="g")+ylim(0,14)

Mtcars %>% Freq_table("vs","am") %>% arrange(am)
Mtcars %>% Freq_table("vs","am") %>% arrange(vs)
Mtcars %>% Freq_table("vs","am") %>% arrange(Freq)

##when the data has already been aggregated once
df_sample <- tribble(
  ~name,    ~gender,   ~ sum,
  "Max",    "male",       10,
  "Sandra", "female",      1,
  "Susan",  "female",      4)
df_sample

#weight frequency
df_sample %>% unCount(sel = "sum") %>% Freq_table("gender")

# this function
df_sample %>% Freq_table("gender", wt="sum")
#'
# 데이터 생성
df <- data.frame(
  term = c("", "10-20분", "20-30분", "30분 이상", "5-10분", "5분 이내"),
  Freq = c(728, 222, 77, 120, 319, 135)
)

# 기본 그래프
Freq_table(df, term, prop = TRUE, remove_rows = 1, add_row_sum = FALSE, plot = TRUE)

# xlab과 flip 옵션 사용
Freq_table(df, term, prop = TRUE, remove_rows = 1, add_row_sum = TRUE, xlab = "응답 시간", flip = TRUE, plot = TRUE)
#'

#빈도분석 및 그래프
Freq_table(eduteck2$Q9, prop=TRUE, remove_rows=1, add_row_sum=TRUE, vjust=-0.5,
           plot=T, size_axis = 14, flip=FALSE, yadd=130, xlab="Eduteck Time", size_text = 6) #'

## End(Not run)
\end{ExampleCode}
\end{Examples}
\inputencoding{utf8}
\HeaderA{Freq\_table\_apa}{Interpreting frequency analysis results from Freq\_table}{Freq.Rul.table.Rul.apa}
%
\begin{Description}
Interpreting frequency analysis results from Freq\_table
\end{Description}
%
\begin{Usage}
\begin{verbatim}
Freq_table_apa(df, title = "", print = TRUE, digits = 2)
\end{verbatim}
\end{Usage}
%
\begin{Arguments}
\begin{ldescription}
\item[\code{df}] data

\item[\code{title}] title =""

\item[\code{print}] df output

\item[\code{digits}] 2
\end{ldescription}
\end{Arguments}
%
\begin{Value}
apa
\end{Value}
%
\begin{Examples}
\begin{ExampleCode}

## Not run: 
#'
Mtcars = mtcars

Mtcars$am = factor(Mtcars$am, levels=c(0,1), labels= c("automatic","manual" ))
Mtcars$vs  = factor(Mtcars$vs, levels=c(0,1), labels= c("V-shaped","straight" ))
Mtcars$cyl  = factor(Mtcars$cyl, levels=c(4,6,8), labels= c("cyl-4","cyl-6","cyl-8"))

Mtcars %>% Freq_table("am","vs", prop=T) %>% Freq_table_apa()
Mtcars %>% table_freq("vs","am","cyl") %>% Freq_table_apa()
Mtcars %>% table_freq("vs","am","cyl", angle=90)%>% Freq_table_apa()

## End(Not run)

\end{ExampleCode}
\end{Examples}
\inputencoding{utf8}
\HeaderA{Freq\_table\_colnames}{Freq\_table\_colnames caange korean names}{Freq.Rul.table.Rul.colnames}
%
\begin{Description}
Freq\_table\_colnames caange korean names

table\_df\_colnames
\end{Description}
%
\begin{Usage}
\begin{verbatim}
Freq_table_colnames(df, input = "구분")

Freq_table_colnames(df, input = "구분")
\end{verbatim}
\end{Usage}
%
\begin{Arguments}
\begin{ldescription}
\item[\code{df}] df

\item[\code{input}] "구분
\end{ldescription}
\end{Arguments}
%
\begin{Value}
df

df
\end{Value}
%
\begin{Examples}
\begin{ExampleCode}
## Not run: 
#' table_df_multit(eduteck2 %>% select(contains("Q5_"), -Q5_4_TEXT)) %>%
arrange(desc(Freq)) %>%Freq_table_colnames()

## End(Not run)

## Not run: 



## End(Not run)
\end{ExampleCode}
\end{Examples}
\inputencoding{utf8}
\HeaderA{Freq\_table\_multi}{Create a Frequency Table from Delimited Multiple Responses}{Freq.Rul.table.Rul.multi}
%
\begin{Description}
This function creates a frequency table from a column containing delimited multiple responses.
It accepts either a data frame with a specified column or a simple character vector.
\end{Description}
%
\begin{Usage}
\begin{verbatim}
Freq_table_multi(df, col = 1, sep = ",\\s*")
\end{verbatim}
\end{Usage}
%
\begin{Arguments}
\begin{ldescription}
\item[\code{df}] A data frame or a character vector. If a vector is provided, it is converted to a one-column data frame named \code{response}.

\item[\code{col}] Column name or index to be processed when \code{df} is a data frame. Defaults to \code{1}.

\item[\code{sep}] A regular expression string for splitting responses. Defaults to \code{",\bsl{}s*"} (comma followed by optional spaces).
\end{ldescription}
\end{Arguments}
%
\begin{Value}
A data frame with two columns: \code{item} and \code{n}, where \code{item} is the response item and \code{n} is the count.
\end{Value}
%
\begin{Examples}
\begin{ExampleCode}
# Example 1: Using a character vector
responses <- c("apple, banana", "banana, apple", "banana", "orange, apple")
Freq_table_multi(responses)

# Example 2: Using a data frame column with different separator
df <- data.frame(Q1 = c("cat / dog", "dog", "cat / bird", "bird / dog"))
Freq_table_multi(df, col = "Q1", sep = "\\s*/\\s*")
\end{ExampleCode}
\end{Examples}
\inputencoding{utf8}
\HeaderA{fs}{Add factor score calculation functions and data}{fs}
%
\begin{Description}
Add factor score calculation functions and data
\end{Description}
%
\begin{Usage}
\begin{verbatim}
fs(data, range = 1:ncol(data), nfactor = NULL, prefix = "", rotate = "varimax")
\end{verbatim}
\end{Usage}
%
\begin{Arguments}
\begin{ldescription}
\item[\code{data}] data

\item[\code{range}] range default 1:ncol(data)

\item[\code{nfactor}] number of factor

\item[\code{prefix}] data prefic

\item[\code{rotate}] varimax
\end{ldescription}
\end{Arguments}
%
\begin{Value}
data
\end{Value}
%
\begin{Examples}
\begin{ExampleCode}


## Not run: 
#'
# 활용방법 1
factor_scores(mtcars,
              range = c("disp", "hp", "drat", "wt", "qsec"),
              nfactor = 2, prefix="a")

mtcars %>% select("disp", "hp", "drat", "wt", "qsec") %>% factor_scores(nfactor=2)

# 활용방법 2
factor_scores(mtcars, range = c(3:7), nfactor = 2)

factor_scores(mtcars, c("disp", "hp", "drat", "wt", "qsec"), nfactor = 2)
factor_scores(mtcars, c(3, 4, 5, 6, 7), nfactor = 2)

factor_scores(mtcars, range = "disp:qsec", nfactor = 2)


mtcars %>% select("disp", "hp", "drat", "wt", "qsec") %>%
  factor_scores(c("disp","hp"),nfactor=2, prefix="a") %>%
  factor_scores(c("drat", "wt", "qsec"),nfactor=3, prefix="b") %>%head()


## End(Not run)
\end{ExampleCode}
\end{Examples}
\inputencoding{utf8}
\HeaderA{g}{Googole Translate korean to english}{g}
%
\begin{Description}
Googole Translate korean to english
\end{Description}
%
\begin{Usage}
\begin{verbatim}
g(text = "hi.", slang = "ko", elang = "en", show = "translate")
\end{verbatim}
\end{Usage}
%
\begin{Arguments}
\begin{ldescription}
\item[\code{text}] input your text

\item[\code{slang}] start language. default Korean

\item[\code{elang}] translated language. default English

\item[\code{show}] show Methode. 'translate' is console, 'browser'  is Print to your computer's browser, 'all' are There are two ways to output it. 'data' is char data
\end{ldescription}
\end{Arguments}
%
\begin{Examples}
\begin{ExampleCode}
# example code
## Not run: 
#korean to english : default
g("안녕하세요 번역을 시작합니다!(by 박중희).", slang ="ko", elang = "en")
g("안녕하세요 번역을 시작합니다!(by 박중희).", "ko", "en")
# english to korean
g("Hello, let’s start translating! (by Park Joong-hee", slang ="ko", elang = "en")
g("Hello, let’s start translating! (by Park Joong-hee","ko","en")
#japan to korean
g("こんにちは翻訳を始めます！（byパク・ジュンヒ", slang ="ja", elang = "ko")

# ko: korean , en: English , ja: Japan, china : zh
# https://cloud.google.com/translate/docs/languages?hl=ko

## End(Not run)
\end{ExampleCode}
\end{Examples}
\inputencoding{utf8}
\HeaderA{general\_Sobel\_test}{Sobeltest toal}{general.Rul.Sobel.Rul.test}
%
\begin{Description}
Sobeltest toal
\end{Description}
%
\begin{Usage}
\begin{verbatim}
general_Sobel_test(a, SE, roburst = FALSE, n = NULL, k = 2, digits = 5)
\end{verbatim}
\end{Usage}
%
\begin{Arguments}
\begin{ldescription}
\item[\code{a}] est

\item[\code{SE}] se

\item[\code{roburst}] roburst T whide correction

\item[\code{n}] sample size

\item[\code{k}] number of regressin coefs, default 2

\item[\code{digits}] 5

\item[\code{conf}] 0.90, 0.95, 0.99, 0.999

\item[\code{type}] all, res, sobel, aroian, goodman
\end{ldescription}
\end{Arguments}
%
\begin{Value}
result
\end{Value}
%
\begin{Examples}
\begin{ExampleCode}

## Not run: 
#'
# Example usage
a <- 0.5  # Example coefficient a
b <- 0.3  # Example coefficient b
sa <- 0.1 # Standard error of a
sb <- 0.05 # Standard error of b

Sobel_test(a, b, sa, sb)
Sobel_test(a, b, sa, sb, "all")
Sobel_test(a, b, sa, sb, "Sobel")

Sobel_test(a, b, sa, sb,"aroian")

Sobel_test(a, b, sa, sb,"Goodman")
#'
a <- c(0.5, 0.3, 0.4)
SE <- c(0.1, 0.1, 0.1)
general_Sobel_test(a, SE)
general_Sobel_test(a, SE)%>%unite_ci(digits=4)



## End(Not run)
\end{ExampleCode}
\end{Examples}
\inputencoding{utf8}
\HeaderA{GetCov}{text to get covariance matrix}{GetCov}
%
\begin{Description}
text to get covariance matrix
\end{Description}
%
\begin{Usage}
\begin{verbatim}
GetCov(
  x,
  lower = TRUE,
  diagonal = TRUE,
  sds = NULL,
  names = paste("V", 1:nvar, sep = "")
)
\end{verbatim}
\end{Usage}
%
\begin{Arguments}
\begin{ldescription}
\item[\code{x}] text

\item[\code{lower}] True

\item[\code{diagonal}] TRUE

\item[\code{sds}] NULL if cor is need

\item[\code{names}] mat colnames,
\end{ldescription}
\end{Arguments}
%
\begin{Value}
matrix
\end{Value}
%
\begin{Examples}
\begin{ExampleCode}

## Not run: 
#'
library(jjstat)

lower <- '
 11.834
  6.947   9.364
  6.819   5.091  12.532
  4.783   5.028   7.495   9.986
 -3.839  -3.889  -3.841  -3.625  9.610
-21.899 -18.831 -21.748 -18.775 35.522 450.288 '

GetCov(lower)
#' GetCov(lower, names = c("anomia67", "powerless67", "anomia71", "powerless71","education", "sei"))

wheaton.cov <-
  GetCov(lower, names = c("anomia67", "powerless67",
                          "anomia71", "powerless71",
                          "education", "sei"))
wheaton.cov
wheaton.model <- '
  # latent variables
    ses     =~ education + sei
    alien67 =~ anomia67 + powerless67
    alien71 =~ anomia71 + powerless71
  # regressions
    alien71 ~ alien67 + ses
    alien67 ~ ses
  # correlated residuals
    anomia67 ~~ anomia71
    powerless67 ~~ powerless71
'
fit <- sem(wheaton.model,
           sample.cov = wheaton.cov,
           sample.nobs = 932)
summary(fit, standardized = TRUE)
diagram2(fit, curve=2)

## End(Not run)


\end{ExampleCode}
\end{Examples}
\inputencoding{utf8}
\HeaderA{goodman\_sobel\_test\_general}{Goodman Sobel Test generalized function}{goodman.Rul.sobel.Rul.test.Rul.general}
%
\begin{Description}
Goodman Sobel Test generalized function
\end{Description}
%
\begin{Usage}
\begin{verbatim}
goodman_sobel_test_general(
  a,
  SE,
  roburst = FALSE,
  n = NULL,
  k = 2,
  digits = 5,
  conf = 0.95
)
\end{verbatim}
\end{Usage}
%
\begin{Arguments}
\begin{ldescription}
\item[\code{a}] est

\item[\code{SE}] se

\item[\code{roburst}] roburst T whide correction

\item[\code{n}] sample size

\item[\code{k}] number of regressin coefs, default 2

\item[\code{digits}] 5

\item[\code{conf}] 0.90, 0.95, 0.99, 0.999
\end{ldescription}
\end{Arguments}
%
\begin{Value}
result
\end{Value}
%
\begin{Examples}
\begin{ExampleCode}

## Not run: 
#'
# Example usage
a <- 0.5  # Example coefficient a
b <- 0.3  # Example coefficient b
sa <- 0.1 # Standard error of a
sb <- 0.05 # Standard error of b

Sobel_test(a, b, sa, sb)
Sobel_test(a, b, sa, sb, "all")
Sobel_test(a, b, sa, sb, "Sobel")

Sobel_test(a, b, sa, sb,"Aroian")

Sobel_test(a, b, sa, sb,"Goodman")
#'
# 예제 데이터
a <- c(0.5, 0.3, 0.4)
SE <- c(0.1, 0.1, 0.1)

# 함수 실행 (기본 95% 신뢰구간)
goodman_sobel_test_general(a, SE)

# 함수 실행 (90% 신뢰구간)
result_90 <- goodman_sobel_test_general(a, SE, conf=0.90)
print(result_90)

# 함수 실행 (99% 신뢰구간)
result_99 <- goodman_sobel_test_general(a, SE, conf=0.99)
print(result_99)

# 함수 실행 (99.9% 신뢰구간)
result_999 <- goodman_sobel_test_general(a, SE, conf=0.999)
print(result_999)
#'


## End(Not run)
\end{ExampleCode}
\end{Examples}
\inputencoding{utf8}
\HeaderA{gtl}{Googole Translate korean to english}{gtl}
%
\begin{Description}
Googole Translate korean to english
\end{Description}
%
\begin{Usage}
\begin{verbatim}
gtl(text = "hi.", slang = "ko", elang = "en", show = "data")
\end{verbatim}
\end{Usage}
%
\begin{Arguments}
\begin{ldescription}
\item[\code{text}] input your text

\item[\code{slang}] start language. default Korean

\item[\code{elang}] translated language. default English

\item[\code{show}] show Methode. 'translate' is console, 'browser'  is Print to your computer's browser, 'all' are There are two ways to output it. 'data' is char data
\end{ldescription}
\end{Arguments}
%
\begin{Examples}
\begin{ExampleCode}
# example code
## Not run: 
#korean to english : default
g("안녕하세요 번역을 시작합니다!(by 박중희).", slang ="ko", elang = "en")
g("안녕하세요 번역을 시작합니다!(by 박중희).", "ko", "en")
# english to korean
g("Hello, let’s start translating! (by Park Joong-hee", slang ="ko", elang = "en")
g("Hello, let’s start translating! (by Park Joong-hee","ko","en")
#japan to korean
g("こんにちは翻訳を始めます！（byパク・ジュンヒ", slang ="ja", elang = "ko")

# ko: korean , en: English , ja: Japan, china : zh
# https://cloud.google.com/translate/docs/languages?hl=ko

## End(Not run)
\end{ExampleCode}
\end{Examples}
\inputencoding{utf8}
\HeaderA{gtl\_df}{Google translate data.frame \#gtl\_df() dataframe 번역 ---------- \#영어논문을 만들 때, 표에 들어가는 단어를 번역하여 넣기 \#너무 많은 수는 오래걸림--> 웹페이지를 갔다와야 하므로... Title}{gtl.Rul.df}
%
\begin{Description}
Google translate data.frame
\#gtl\_df() dataframe 번역 ----------
\#영어논문을 만들 때, 표에 들어가는 단어를 번역하여 넣기
\#너무 많은 수는 오래걸림--> 웹페이지를 갔다와야 하므로...
Title
\end{Description}
%
\begin{Usage}
\begin{verbatim}
gtl_df(df, col = 1)
\end{verbatim}
\end{Usage}
%
\begin{Arguments}
\begin{ldescription}
\item[\code{df}] data.frame

\item[\code{col}] Select the column you want to translate
\end{ldescription}
\end{Arguments}
%
\begin{Examples}
\begin{ExampleCode}
## Not run: 
 df0 <-  data.frame(word=c("온라인","진행","교육"),
                   n=c(100, 200, 300))
 df0 %>% gtl_df()

## End(Not run)


\end{ExampleCode}
\end{Examples}
\inputencoding{utf8}
\HeaderA{hist2}{histogram}{hist2}
%
\begin{Description}
histogram
\end{Description}
%
\begin{Usage}
\begin{verbatim}
hist2(data, color = "lightblue", linecolor = "red")
\end{verbatim}
\end{Usage}
%
\begin{Arguments}
\begin{ldescription}
\item[\code{data}] vector

\item[\code{color}] fill color

\item[\code{linecolor}] line color
\end{ldescription}
\end{Arguments}
%
\begin{Value}
graph histogram
\end{Value}
%
\begin{Examples}
\begin{ExampleCode}
 ## Not run: 

set.seed(123)

data <- rnorm(100)

hist2(data)

 
## End(Not run)

\end{ExampleCode}
\end{Examples}
\inputencoding{utf8}
\HeaderA{histogram}{histogram}{histogram}
%
\begin{Description}
histogram
\end{Description}
%
\begin{Usage}
\begin{verbatim}
histogram(data, color = "gray60")
\end{verbatim}
\end{Usage}
%
\begin{Arguments}
\begin{ldescription}
\item[\code{data}] vector

\item[\code{color}] color
\end{ldescription}
\end{Arguments}
%
\begin{Value}
plot
\end{Value}
%
\begin{Examples}
\begin{ExampleCode}
## Not run: 
set.seed(123)
data <- rnorm(100)
histogram(data)


## End(Not run)

\end{ExampleCode}
\end{Examples}
\inputencoding{utf8}
\HeaderA{HTMT}{HTMT for Single-level SEM/CFA (lavaan)}{HTMT}
%
\begin{Description}
This function computes the Heterotrait–Monotrait ratio (HTMT)
for single-level CFA/SEM models fitted with \code{lavaan}.
The output is a lower-triangular HTMT matrix (upper triangle and diagonal are NA).
\end{Description}
%
\begin{Usage}
\begin{verbatim}
HTMT(fit, digits = 3)
\end{verbatim}
\end{Usage}
%
\begin{Arguments}
\begin{ldescription}
\item[\code{fit}] A \code{lavaan} cfa/sem fitted object.

\item[\code{digits}] Integer. Number of digits to round the HTMT values. Default is 3.
\end{ldescription}
\end{Arguments}
%
\begin{Details}
HTMT is computed using observed-variable correlations from
\code{lavaan::lavInspect(fit, "cor.ov")}.
Only latent variables with at least two indicators are included.

The HTMT for latent variables i and j is computed as:
HTMT\_ij = mean( |r\_xy| ) /  mean( |r\_xx| ) * mean( |r\_yy| )  over 2
where r\_xy are heterotrait correlations (indicators across constructs),
and r\_xx, r\_yy are monotrait correlations (within-construct indicators).
\end{Details}
%
\begin{Value}
A data.frame containing the HTMT matrix (lower triangle only).
\end{Value}
%
\begin{References}
Henseler, J., Ringle, C. M., \& Sarstedt, M. (2015).
A new criterion for assessing discriminant validity in variance-based
structural equation modeling. Journal of the Academy of Marketing Science,
43(1), 115–135.
\end{References}
%
\begin{Examples}
\begin{ExampleCode}
## Not run: 
HS.model <- '
  visual  =~ x1 + x2 + x3
  textual =~ x4 + x5 + x6
  speed   =~ x7 + x8 + x9
'
fit <- lavaan::cfa(HS.model, data = lavaan::HolzingerSwineford1939)

HTMT(fit)
HTMT(fit, digits = 4)

## End(Not run)
\end{ExampleCode}
\end{Examples}
\inputencoding{utf8}
\HeaderA{htmt\_mlsem}{HTMT for Multilevel SEM}{htmt.Rul.mlsem}
%
\begin{Description}
This function computes the Heterotrait–Monotrait ratio (HTMT)
for multilevel CFA/SEM models estimated using lavaan.
HTMT is calculated separately for within-level or between-level
measurement models based on observed-variable correlations.
\end{Description}
%
\begin{Usage}
\begin{verbatim}
htmt_mlsem(fit, level = c("within", "between"), digits = 3)
\end{verbatim}
\end{Usage}
%
\begin{Arguments}
\begin{ldescription}
\item[\code{fit}] A lavaan object fitted with a multilevel CFA/SEM model.

\item[\code{level}] Character. One of \code{"within"} or \code{"between"}.

\item[\code{digits}] Integer. Number of digits to round the HTMT values.
\end{ldescription}
\end{Arguments}
%
\begin{Details}
HTMT is defined as the ratio of the mean of heterotrait–heteromethod
correlations to the geometric mean of monotrait–heteromethod
correlations.

Only latent variables with at least two indicators are included.
The function uses level-specific observed-variable correlation
matrices extracted via \code{lavInspect(fit, "cor.ov")}.
\end{Details}
%
\begin{Value}
A data.frame containing the HTMT matrix
(lower triangle only; diagonal and upper triangle are NA).
\end{Value}
%
\begin{References}
Henseler, J., Ringle, C. M., \& Sarstedt, M. (2015).
A new criterion for assessing discriminant validity in
variance-based structural equation modeling.
Journal of the Academy of Marketing Science, 43(1), 115–135.
\end{References}
%
\begin{Examples}
\begin{ExampleCode}
## Not run: 
fit <- lavaan::sem(model, data = Demo.twolevel, cluster = "cluster")

mlsem_htmt(fit, level = "within")
mlsem_htmt(fit, level = "between")

## End(Not run)
\end{ExampleCode}
\end{Examples}
\inputencoding{utf8}
\HeaderA{htmt\_sem}{HTMT for Single-level SEM/CFA (lavaan)}{htmt.Rul.sem}
%
\begin{Description}
This function computes the Heterotrait–Monotrait ratio (HTMT)
for single-level CFA/SEM models fitted with \code{lavaan}.
The output is a lower-triangular HTMT matrix (upper triangle and diagonal are NA).
\end{Description}
%
\begin{Usage}
\begin{verbatim}
htmt_sem(fit, digits = 3)
\end{verbatim}
\end{Usage}
%
\begin{Arguments}
\begin{ldescription}
\item[\code{fit}] A \code{lavaan} cfa/sem fitted object.

\item[\code{digits}] Integer. Number of digits to round the HTMT values. Default is 3.
\end{ldescription}
\end{Arguments}
%
\begin{Details}
HTMT is computed using observed-variable correlations from
\code{lavaan::lavInspect(fit, "cor.ov")}.
Only latent variables with at least two indicators are included.

The HTMT for latent variables i and j is computed as:
HTMT\_ij = mean( |r\_xy| ) /  mean( |r\_xx| ) * mean( |r\_yy| )  over 2
where r\_xy are heterotrait correlations (indicators across constructs),
and r\_xx, r\_yy are monotrait correlations (within-construct indicators).
\end{Details}
%
\begin{Value}
A data.frame containing the HTMT matrix (lower triangle only).
\end{Value}
%
\begin{References}
Henseler, J., Ringle, C. M., \& Sarstedt, M. (2015).
A new criterion for assessing discriminant validity in variance-based
structural equation modeling. Journal of the Academy of Marketing Science,
43(1), 115–135.
\end{References}
%
\begin{Examples}
\begin{ExampleCode}
## Not run: 
HS.model <- '
  visual  =~ x1 + x2 + x3
  textual =~ x4 + x5 + x6
  speed   =~ x7 + x8 + x9
'
fit <- lavaan::cfa(HS.model, data = lavaan::HolzingerSwineford1939)

htmt_sem(fit)
htmt_sem(fit, digits = 4)

## End(Not run)
\end{ExampleCode}
\end{Examples}
\inputencoding{utf8}
\HeaderA{ICC}{ICC}{ICC}
%
\begin{Description}
ICC
\end{Description}
%
\begin{Usage}
\begin{verbatim}
ICC(lmedata, type = "mixed", digits = 4)
\end{verbatim}
\end{Usage}
%
\begin{Arguments}
\begin{ldescription}
\item[\code{lmedata}] lmer result

\item[\code{type}] 'mixed', 'logit', 'pois'

\item[\code{digits}] digits = 3 default
\end{ldescription}
\end{Arguments}
%
\begin{Examples}
\begin{ExampleCode}
## Not run: 
data(sleepstudy)
sleepstudy |> str()
lmer(Reaction ~ Days + (Days | Subject), sleepstudy) |> ICC()



## End(Not run)
\end{ExampleCode}
\end{Examples}
\inputencoding{utf8}
\HeaderA{icc\_batch}{Batch calculation of ICCs for multiple variables}{icc.Rul.batch}
%
\begin{Description}
This function applies \code{icc\_lme4()} to multiple outcome variables
and returns a combined table of ICC results.
\end{Description}
%
\begin{Usage}
\begin{verbatim}
icc_batch(data, vars, cluster)
\end{verbatim}
\end{Usage}
%
\begin{Arguments}
\begin{ldescription}
\item[\code{data}] data.frame. Dataset containing outcome variables and cluster variable.

\item[\code{vars}] character vector. Names of numeric outcome variables.

\item[\code{cluster}] character string. Name of the clustering variable.
\end{ldescription}
\end{Arguments}
%
\begin{Value}
A tibble with ICC results for each variable.
\end{Value}
%
\begin{Examples}
\begin{ExampleCode}
## Not run: 
# Built-in dataset
data(mtcars)

# Single ICC calculation
icc_lme4(
  y = mtcars$hp,
  cluster = mtcars$cyl
)

# Batch ICC calculation
icc_batch(
  data = mtcars,
  vars = c("mpg", "hp", "wt"),
  cluster = "cyl"
)

## End(Not run)
\end{ExampleCode}
\end{Examples}
\inputencoding{utf8}
\HeaderA{icc\_lme4}{Calculate Intraclass Correlation Coefficient (ICC) using lme4}{icc.Rul.lme4}
%
\begin{Description}
This function computes the Intraclass Correlation Coefficient (ICC)
based on a random-intercept-only linear mixed-effects model.
\end{Description}
%
\begin{Usage}
\begin{verbatim}
icc_lme4(y, cluster)
\end{verbatim}
\end{Usage}
%
\begin{Arguments}
\begin{ldescription}
\item[\code{y}] numeric vector. Outcome variable.

\item[\code{cluster}] vector or factor. Grouping (cluster) variable.
\end{ldescription}
\end{Arguments}
%
\begin{Value}
A data.frame with the following components:
\begin{description}

\item[ICC] Intraclass correlation coefficient
\item[between\_var] Between-cluster variance (tau\textasciicircum{}2)
\item[within\_var] Within-cluster variance (sigma\textasciicircum{}2)
\item[J] Number of clusters
\item[kbar] Average cluster size

\end{description}

\end{Value}
%
\begin{Examples}
\begin{ExampleCode}
## Not run: 

# Example using a built-in R dataset
data(ToothGrowth)

# ICC of tooth length clustered by supplement type
icc_lme4(
  y = ToothGrowth$len,
  cluster = ToothGrowth$supp
)

# ICC of tooth length clustered by dose
icc_lme4(
  y = ToothGrowth$len,
  cluster = ToothGrowth$dose


# Built-in dataset
data(mtcars)

# Single ICC calculation
icc_lme4(
  y = mtcars$mpg,
  cluster = mtcars$cyl
)

# Batch ICC calculation using icc_batch()
icc_batch(
  data = mtcars,
  vars = c("mpg", "hp", "wt"),
  cluster = "cyl"
)

## End(Not run)
\end{ExampleCode}
\end{Examples}
\inputencoding{utf8}
\HeaderA{IE\_effect}{SEM indirect effect}{IE.Rul.effect}
%
\begin{Description}
SEM indirect effect
\end{Description}
%
\begin{Usage}
\begin{verbatim}
IE_effect(model_sem, effect1 = "IE", effect2 = "IE", label = FALSE, ci = TRUE)
\end{verbatim}
\end{Usage}
%
\begin{Arguments}
\begin{ldescription}
\item[\code{model\_sem}] lavaan objgect

\item[\code{effect1}] "IE"

\item[\code{effect2}] "IE"

\item[\code{label}] label

\item[\code{ci}] confidence interval
\end{ldescription}
\end{Arguments}
%
\begin{Value}
parameterEstimates()
\end{Value}
%
\begin{Examples}
\begin{ExampleCode}

## Not run: 
#library(lavaan) # only needed once per session

models1 <- '
  # measurement model
    ind60 =~ x1 + x2 + x3
    dem60 =~ y1 + y2 + y3 + y4
    dem65 =~ y5 + y6 + y7 + y8
    # regressions
    dem60 ~ ind60
    dem65 ~ ind60 + dem60
      # residual correlations
    y1 ~~ y5
    y2 ~~ y4 + y6
    y3 ~~ y7
    y4 ~~ y8
    y6 ~~ y8

models2 ='
  # measurement model
    ind60 =~ x1 + x2 + x3
    dem60 =~ y1 + y2 + y3 + y4
    dem65 =~ y5 + y6 + y7 + y8
     # regressions
    dem60 ~ a1*ind60
    dem65 ~ a2*ind60 + a3*dem60
 #NEw
   ie := a1*a3
   te := a2 +a1*a3
   diff := a1 - a3




fits1 <- sem(models1, data=PoliticalDemocracy)
fits2 <- sem(models2, data=PoliticalDemocracy)
summary(fits1, standardized=TRUE)
summary(fits2, standardized=TRUE)
fits2 %>% IE_effect()

## End(Not run)

\end{ExampleCode}
\end{Examples}
\inputencoding{utf8}
\HeaderA{indicator\_validity\_ml}{Indicator validity table for (Multi-level) CFA}{indicator.Rul.validity.Rul.ml}
%
\begin{Description}
Indicator validity table for (Multi-level) CFA
\end{Description}
%
\begin{Usage}
\begin{verbatim}
indicator_validity_ml(fit, level = c("within", "between"), digits = 3)
\end{verbatim}
\end{Usage}
%
\begin{Arguments}
\begin{ldescription}
\item[\code{fit}] lavaan SEM / ML-SEM object

\item[\code{level}] "within" or "between"

\item[\code{digits}] numeric
\end{ldescription}
\end{Arguments}
%
\begin{Value}
data.frame
\end{Value}
\inputencoding{utf8}
\HeaderA{interlist}{Functions that compensate for the method of input}{interlist}
%
\begin{Description}
Functions that compensate for the method of input
\end{Description}
%
\begin{Usage}
\begin{verbatim}
interlist(...)
\end{verbatim}
\end{Usage}
%
\begin{Arguments}
\begin{ldescription}
\item[\code{...}] LONG LIST
\end{ldescription}
\end{Arguments}
%
\begin{Value}
list
\end{Value}
%
\begin{Examples}
\begin{ExampleCode}

## Not run: 


# want list
intbasic = list(c("IMAG", 1, 2), c("SAT", 1, 2))
intbasic

inta = interlist(list(c("IMAG", 1, 2, 3), c("SAT", 1, 2)));inta
intc = interlist(list("IMAG", 1, 2, 3, "SAT", 1, 2));intc

intb = interlist(c("IMAG", 1, 2, 3), c("SAT", 1, 2));intb
intd = interlist("IMAG", 1, 2, 3, "SAT", 1, 2);intd
#'

# Partially select to create an interaction term
satisfaction%>% plspm_prod(sat_blocks1,
                           list(c("IMAG", 1, 2, 3), c("SAT", 1, 2, 3))) %>%str()

satisfaction%>% plspm_prod(sat_blocks1, inta) %>%str()
satisfaction%>% plspm_prod(sat_blocks1, intb) %>%str()
satisfaction%>% plspm_prod(sat_blocks1, intc) %>%str()



satisfaction%>%
plspm_prod(sat_blocks1, interlist(list("IMAG", 1, 2, 3, "SAT", 1, 2))) %>%str()
satisfaction%>%
plspm_prod(sat_blocks1, interlist(list("IMAG", 1, "SAT", 1, 2))) %>%str()
satisfaction%>%
plspm_prod(sat_blocks1, interlist(list("IMAG", 1,3,  "SAT", 2, 3))) %>%str()





## End(Not run)



\end{ExampleCode}
\end{Examples}
\inputencoding{utf8}
\HeaderA{interplot}{interplot is using data}{interplot}
%
\begin{Description}
interplot is using data
\end{Description}
%
\begin{Usage}
\begin{verbatim}
interplot(
  data_long,
  type = "plot",
  size_point = 4,
  size_text = 5,
  linewidth = 1,
  vjust = -0.4,
  hjust = -2,
  title = NULL,
  xlab = "x_variable",
  ylab = "y_variable",
  legend_title = "v2",
  yadd = 0.5,
  size_element = 12,
  t = FALSE,
  xy = FALSE,
  digits = 3
)
\end{verbatim}
\end{Usage}
%
\begin{Arguments}
\begin{ldescription}
\item[\code{data\_long}] long foramt data

\item[\code{size\_point}] point size default 4

\item[\code{size\_text}] text size  default 5

\item[\code{linewidth}] line width default 1

\item[\code{vjust}] vjust

\item[\code{hjust}] hjust

\item[\code{title}] title null

\item[\code{xlab}] xlab

\item[\code{ylab}] ylab

\item[\code{legend\_title}] Change legend\_title

\item[\code{yadd}] fig size add value

\item[\code{size\_element}] axis element size

\item[\code{t}] matrix transpose efault FALSE, use when matrix data is wide format

\item[\code{xy}] plot ranspose default FALSE

\item[\code{digits}] default 3
\end{ldescription}
\end{Arguments}
%
\begin{Value}
plot
\end{Value}
%
\begin{Examples}
\begin{ExampleCode}
## Not run: 

## interplot_data use!!
interplot_data( 6, 1, 4, 5, type = "long")%>%
interplot(size_element = 10)

interplot_data( 2, 1, 4, 5, type="long")%>%interplot()

## 2 *3 matriix
interplot_data( 4, 1, 2, 5, 6, 3, type="long", row=3, col=2)%>%interplot()

##3*2
interplot_data( 4, 1, 2, 5, 6, 3, type="long", row=2, col=3)%>%interplot()

##all data output
interplot_data( 4, 1, 2, 5, 6, 3, type="long", row=2, col=3)%>%interplot("all")

#generated direct data matrix and data frame
matrix(c(4, 1, 2, 5, 6, 3), nrow=2,ncol=3)  #marix

#long format
matrix(c(4, 1, 2, 5, 6, 3), nrow=2,ncol=3) %>%
data.frame() %>% rownames_to_column("v1") %>%
pivot_longer(cols=2:4)

##plot
matrix(c(4, 1, 2, 5, 6, 3), nrow=2,ncol=3) %>%
 data.frame() %>%
rownames_to_column("v1") %>% pivot_longer(cols=2:4)%>%
interplot()

## onestep method when use this function
matrix(c(4, 1, 2, 5, 6, 3), nrow=2, ncol=3) %>% interplot(t=TRUE)


#  3 * 2
matrix(c(4, 1, 2, 5, 6, 3), nrow=3, ncol=2)

matrix(c(4, 1, 2, 5, 6, 3), nrow=3, ncol=2) %>%
data.frame() %>%
rownames_to_column("v1") %>%
pivot_longer(cols=2:3)

matrix(c(4, 1, 2, 5, 6, 3), nrow=3, ncol=2) %>%
 data.frame() %>%
 rownames_to_column("v1") %>% p
 ivot_longer(cols=2:3)%>%
 interplot()


##one step output
matrix(c(4, 1, 2, 5, 6, 3), nrow=3, ncol=2) %>%
interplot(t=TRUE)


#### all data output
matrix(c(4, 1, 2, 5, 6, 3), nrow=3, ncol=2) %>% interplot(t=TRUE, type="all")


## End(Not run)

\end{ExampleCode}
\end{Examples}
\inputencoding{utf8}
\HeaderA{interplot\_data}{interplot\_data is matrix, long\_data, plotting function}{interplot.Rul.data}
%
\begin{Description}
interplot\_data is matrix, long\_data, plotting function
\end{Description}
%
\begin{Usage}
\begin{verbatim}
interplot_data(
  ...,
  col = 2,
  row = 2,
  type = "long",
  title = NULL,
  xlab = "x",
  ylab = "y",
  vjust = -0.4,
  hjust = 0.5,
  size_element = 12,
  byrow = TRUE,
  xy = FALSE,
  yadd = 0.3
)
\end{verbatim}
\end{Usage}
%
\begin{Arguments}
\begin{ldescription}
\item[\code{...}] input data c() or as select()

\item[\code{col}] matrix ncol setting

\item[\code{row}] matrix nrow

\item[\code{type}] type is mat, long, all, plot

\item[\code{title}] plot title

\item[\code{xlab}] plot xlab

\item[\code{ylab}] plot ylab

\item[\code{vjust}] vjust text in plot postion

\item[\code{hjust}] hjust text in plot position

\item[\code{size\_element}] plot element text size

\item[\code{byrow}] byrow=TRUE is default

\item[\code{xy}] plot transpose default FALSE

\item[\code{yadd}] plot range add
\end{ldescription}
\end{Arguments}
%
\begin{Value}
data and plot
\end{Value}
%
\begin{Examples}
\begin{ExampleCode}
## Not run: 
 ## data output
interplot_data( 6, 1, 4, 5, type = "mat")
#' interplot_data( 6, 1, 4, 5, type = "all")
interplot_data( 6, 1, 4, 5, type = "long")
interplot_data( 6, 1, 4, 5, type = "plot")

# 2*3 data
interplot_data( c(6, 1, 4, 5,6,7), type = "mat",  col = 3,row = 2)
interplot_data( c(6, 1, 4, 5,6,7), type = "long",  col = 3,row = 2)
interplot_data( 6, 1, 4, 5,6,7, type = "all",  col = 3,row = 2)
interplot_data( 6, 1, 4, 5,6,7, type = "plot",  col = 3,row = 2)

## 3*2 data
interplot_data( 6, 1, 4, 5,6,7, type = "mat",  col = 2,row = 3)
interplot_data( 6, 1, 4, 5,6,7, type = "all",  col = 2,row = 3)
interplot_data( 6, 1, 4, 5,6,7, type = "long",  col = 2, row = 3)
interplot_data( c(6, 1, 4, 5,6,7), type = "plot",  col = 2,row = 3)

### ii is very hard example
interplot_data( 6, 1, 4, 5) %>%
ggplot(aes(x = v1, y = value))+
geom_point(size=4, aes(color=v2))+
geom_line(aes(group = v2, linetype = v2 ), linewidth = 1)+
geom_text(aes(label = value), vjust= -0.4, hjust= 2, size= 5)+
theme_bw()


## End(Not run)
\end{ExampleCode}
\end{Examples}
\inputencoding{utf8}
\HeaderA{interpretation\_de}{SEM direct effect interpretation}{interpretation.Rul.de}
%
\begin{Description}
SEM direct effect interpretation
\end{Description}
%
\begin{Usage}
\begin{verbatim}
interpretation_de(data, md = FALSE)
\end{verbatim}
\end{Usage}
%
\begin{Arguments}
\begin{ldescription}
\item[\code{data}] DE\_effect obj

\item[\code{md}] markdown\_table output
\end{ldescription}
\end{Arguments}
%
\begin{Value}
report apa
\end{Value}
%
\begin{Examples}
\begin{ExampleCode}

## Not run: 

 #library(lavaan) # only needed once per session
 models <- '
   # measurement model
     ind60 =~ x1 + x2 + x3
     dem60 =~ y1 + y2 + y3 + y4
     dem65 =~ y5 + y6 + y7 + y8
 # regressions
     dem60 ~ ind60
     dem65 ~ ind60 + dem60
   # residual correlations
     y1 ~~ y5
     y2 ~~ y4 + y6
     y3 ~~ y7
     y4 ~~ y8
    y6 ~~ y8
 '
 fits <- sem(models, data=PoliticalDemocracy)
 summary(fits, standardized=TRUE)

 ##Direct effect
 DE_effect(fits)
 DE_effect(fits)%>% interpretation_de()


## End(Not run)

\end{ExampleCode}
\end{Examples}
\inputencoding{utf8}
\HeaderA{interpretation\_ie}{interpretation\_ie}{interpretation.Rul.ie}
%
\begin{Description}
interpretation\_ie
\end{Description}
%
\begin{Usage}
\begin{verbatim}
interpretation_ie(
  data,
  msg = "IE",
  md = FALSE,
  est = 3,
  std = 4,
  sig = 7,
  p = 6
)
\end{verbatim}
\end{Usage}
%
\begin{Arguments}
\begin{ldescription}
\item[\code{data}] data

\item[\code{msg}] msg="IE"

\item[\code{md}] md = FALSE

\item[\code{est}] est=3

\item[\code{std}] std=4

\item[\code{sig}] sig = 7

\item[\code{p}] p=6
\end{ldescription}
\end{Arguments}
%
\begin{Value}
ie report
\end{Value}
%
\begin{Examples}
\begin{ExampleCode}

## Not run: 
example(sem)
IE_effect(fit)
IE_effect(fit) %>% interpretation_ie()

## End(Not run)
\end{ExampleCode}
\end{Examples}
\inputencoding{utf8}
\HeaderA{interpretation\_ie\_CI}{sem\_apa\_ie\_CI}{interpretation.Rul.ie.Rul.CI}
%
\begin{Description}
sem\_apa\_ie\_CI
\end{Description}
%
\begin{Usage}
\begin{verbatim}
interpretation_ie_CI(
  data,
  caption = "매개효과 분석 결과 ",
  md = FALSE,
  est = 2,
  std = 3,
  sig = 6,
  CI = 7,
  Z = FALSE,
  tot_filter = "총",
  ind_filter = "IE"
)
\end{verbatim}
\end{Usage}
%
\begin{Arguments}
\begin{ldescription}
\item[\code{data}] lavaan object

\item[\code{caption}] "매개효과 분석 결과",

\item[\code{md}] markdown

\item[\code{est}] est sel 2

\item[\code{std}] std sel 3

\item[\code{sig}] sig sel 6

\item[\code{CI}] CI seledct 7

\item[\code{Z}] Z=FALSE

\item[\code{tot\_filter}] tot\_filter="총",

\item[\code{ind\_filter}] ind\_filter="IE"
\end{ldescription}
\end{Arguments}
%
\begin{Value}
table
\end{Value}
\inputencoding{utf8}
\HeaderA{item}{item}{item}
%
\begin{Description}
item
\end{Description}
%
\begin{Usage}
\begin{verbatim}
item(...)
\end{verbatim}
\end{Usage}
%
\begin{Arguments}
\begin{ldescription}
\item[\code{...}] input
\end{ldescription}
\end{Arguments}
%
\begin{Value}
list
\end{Value}
\inputencoding{utf8}
\HeaderA{item\_CR\_AVE}{Item-level CR \& AVE table}{item.Rul.CR.Rul.AVE}
%
\begin{Description}
Item-level CR \& AVE table
\end{Description}
%
\begin{Usage}
\begin{verbatim}
item_CR_AVE(fit, level = c("within", "between"), digits = 3)
\end{verbatim}
\end{Usage}
%
\begin{Arguments}
\begin{ldescription}
\item[\code{fit}] lavaan (ML-)SEM object

\item[\code{level}] "within" or "between"

\item[\code{digits}] numeric
\end{ldescription}
\end{Arguments}
%
\begin{Value}
data.frame
\end{Value}
\inputencoding{utf8}
\HeaderA{item\_parceling}{Item Parceling Utility Function}{item.Rul.parceling}
%
\begin{Description}
This function creates item parcels from observed item-level data using
aggregation methods (mean, sum, weighted mean) or
factor-score-based methods (EFA / CFA).
\end{Description}
%
\begin{Usage}
\begin{verbatim}
item_parceling(
  data,
  ...,
  fun = "mean",
  type = c("new", "data", "all"),
  weights = NULL,
  standardize_items = FALSE,
  min_valid = 1,
  factor_score = c("regression", "bartlett"),
  lavaan_estimator = "ML",
  std.lv = FALSE,
  name_in_vector = FALSE,
  verbose = FALSE
)
\end{verbatim}
\end{Usage}
%
\begin{Arguments}
\begin{ldescription}
\item[\code{data}] A data.frame containing observed item variables.

\item[\code{...}] Parcel definitions.

\item[\code{fun}] Scoring method. Options include "mean", "sum", "median", "wmean",
"factor\_efa", "fs\_efa", "factor\_cfa", "fs\_cfa".
Abbreviated forms "efa" and "cfa" are allowed with a warning.

\item[\code{type}] Output type:
"new" (parcel data only),
"data" (append parcels to original data),
"all" (data + metadata).

\item[\code{weights}] Numeric vector or a named list of numeric vectors.
For fun = "wmean", parcel-specific weights can be supplied as a list.
Parcels without weights will be aggregated using mean().

\item[\code{standardize\_items}] Logical. If TRUE, items are standardized before scoring.

\item[\code{min\_valid}] Minimum number of valid items required per row.

\item[\code{factor\_score}] Factor score method for EFA ("regression", "bartlett").

\item[\code{lavaan\_estimator}] Estimator for CFA.

\item[\code{std.lv}] Logical. If TRUE, latent variance is fixed to 1 in CFA.

\item[\code{name\_in\_vector}] Logical or "auto".
- TRUE: c("ParcelName","item1","item2",...)
- FALSE: ParcelName = c("item1","item2",...) (default)
- "auto": if the first element is not a column name, treat it as parcel name.

\item[\code{verbose}] Logical. If TRUE, prints reliability summaries.
\end{ldescription}
\end{Arguments}
%
\begin{Value}
Depending on type, a data.frame or a list.
\end{Value}
%
\begin{Examples}
\begin{ExampleCode}
## -------------------------------------------------
## Example data
## -------------------------------------------------
set.seed(123)
shindata_example <- data.frame(
  efficacy2 = rnorm(100, 3.5, 0.8),
  efficacy3 = rnorm(100, 3.6, 0.7),
  efficacy4 = rnorm(100, 3.7, 0.6),

  Accessibility1 = rnorm(100, 4.0, 0.9),
  Accessibility2 = rnorm(100, 3.9, 0.8),
  Accessibility3 = rnorm(100, 4.1, 0.7),
  Accessibility4 = rnorm(100, 4.0, 0.6),
  Accessibility5 = rnorm(100, 3.8, 0.9),

  Interest1 = rnorm(100, 3.2, 0.7),
  Interest2 = rnorm(100, 3.3, 0.8),
  Interest3 = rnorm(100, 3.4, 0.6),
  Interest4 = rnorm(100, 3.5, 0.7)
)

## -------------------------------------------------
## Basic example: sum
## -------------------------------------------------
parcel_sum <- item_parceling(
  data = shindata_example,
  효능감 = c("efficacy2","efficacy3","efficacy4"),
  연구접근성 = c("Accessibility1","Accessibility2",
                  "Accessibility3","Accessibility4","Accessibility5"),
  fun = "sum",
  type = "new"
)
head(parcel_sum)

## -------------------------------------------------
## Weighted mean: single parcel
## -------------------------------------------------
item_parceling(
  data = shindata_example,
  효능감 = c("efficacy2","efficacy3","efficacy4"),
  fun = "wmean",
  weights = c(0.40, 0.35, 0.25),
  type = "new"
)

## -------------------------------------------------
## Weighted mean: parcel-specific weights (recommended)
## -------------------------------------------------
parcel_wmean <- item_parceling(
  data = shindata_example,
  효능감 = c("efficacy2","efficacy3","efficacy4"),
  연구접근성 = c("Accessibility1","Accessibility2",
                  "Accessibility3","Accessibility4","Accessibility5"),
  fun = "wmean",
  weights = list(
    효능감 = c(0.40, 0.35, 0.25)
  ),
  type = "new"
)
head(parcel_wmean)

## Not run: 
## -------------------------------------------------
## Error example: invalid weights usage
## -------------------------------------------------
item_parceling(
  data = shindata_example,
  효능감 = c("efficacy2","efficacy3","efficacy4"),
  연구접근성 = c("Accessibility1","Accessibility2",
                  "Accessibility3","Accessibility4","Accessibility5"),
  fun = "wmean",
  weights = c(0.40, 0.35, 0.25)
)

## End(Not run)

## -------------------------------------------------
## Factor score (EFA)
## -------------------------------------------------
parcel_fs_efa <- item_parceling(
  data = shindata_example,
  효능감 = c("efficacy2","efficacy3","efficacy4"),
  연구접근성 = c("Accessibility1","Accessibility2",
                  "Accessibility3","Accessibility4","Accessibility5"),
  fun = "factor_efa",
  factor_score = "regression",
  type = "all"
)
parcel_fs_efa$score_info$효능감

## -------------------------------------------------
## EFA alias (warning example)
## -------------------------------------------------
item_parceling(
  data = shindata_example,
  효능감 = c("efficacy2","efficacy3","efficacy4"),
  fun = "efa",
  type = "new"
)

## -------------------------------------------------
## CFA factor score
## -------------------------------------------------
parcel_fs_cfa <- item_parceling(
  data = shindata_example,
  효능감 = c("efficacy2","efficacy3","efficacy4"),
  fun = "factor_cfa",
  type = "all"
)
parcel_fs_cfa$score_info$효능감$CR

## -------------------------------------------------
## name_in_vector usage
## -------------------------------------------------
item_parceling(
  data = shindata_example,
  c("효능감","efficacy2","efficacy3","efficacy4"),
  c("진로흥미","Interest1","Interest2","Interest3","Interest4"),
  fun = "factor_efa",
  name_in_vector = TRUE,
  type = "new"
)

item_parceling(
  data = shindata_example,
  c("efficacy2","efficacy3","efficacy4"),
  진로흥미 = c("Interest1","Interest2","Interest3","Interest4"),
  fun = "factor_efa",
  name_in_vector = "auto",
  type = "new"
)
\end{ExampleCode}
\end{Examples}
\inputencoding{utf8}
\HeaderA{jjstat\_plot\_path}{SEM Path Diagram (Base R)}{jjstat.Rul.plot.Rul.path}
%
\begin{Description}
Draws a simple SEM path diagram using base R graphics.
Only latent-latent edges (B, Phi) are weighted by coefficient size and sign.
Measurement paths are drawn with fixed style for readability.
\end{Description}
%
\begin{Usage}
\begin{verbatim}
jjstat_plot_path(
  object,
  main = "SEM Path Diagram",
  layout = c("tree", "tree2", "tree3", "circle"),
  rotate = 1,
  show_phi = TRUE,
  show_theta = FALSE,
  std.all = TRUE,
  digits = 3,
  sizeLat = 8,
  sizeLat2 = 4,
  sizeMan = 5,
  sizeMan2 = 3,
  edge.width.range = c(1, 6),
  edge.color.pos = "black",
  edge.color.neg = "red",
  curve = 0.6,
  curve.max = 2,
  zigzag = TRUE,
  zigzag.amp = 0.035
)
\end{verbatim}
\end{Usage}
%
\begin{Arguments}
\begin{ldescription}
\item[\code{object}] A \code{jjstat\_sem} object or a list containing matrices:
\code{Lambda}, \code{B}, \code{Phi} (and optionally \code{Theta}).

\item[\code{main}] Title.

\item[\code{layout}] Layout for latent variables: tree/tree2/tree3/circle.

\item[\code{rotate}] Integer 1\textasciitilde{}4. Rotate coordinates by 90-degree steps (clockwise).

\item[\code{show\_phi}] Logical. If TRUE, draw latent correlations (Phi).

\item[\code{show\_theta}] Logical. If TRUE, draw residual loops for observed variables (Theta diagonal).

\item[\code{digits}] Numeric. Rounding digits for coefficients.

\item[\code{sizeLat, sizeLat2}] Latent ellipse radii (x/y) scaling.

\item[\code{sizeMan, sizeMan2}] Manifest rectangle half-width/half-height scaling.

\item[\code{edge.width.range}] Numeric length-2. Width range for latent-latent edges.

\item[\code{edge.color.pos}] Color for positive coefficients.

\item[\code{edge.color.neg}] Color for negative coefficients.

\item[\code{curve}] A single numeric alias to control Phi curvature intensity.
Internally mapped to \code{curve.base}. (Recommended to use this.)

\item[\code{curve.max}] Maximum curvature for Phi.

\item[\code{zigzag}] Logical. If TRUE, apply small zigzag offsets to latent positions.

\item[\code{zigzag.amp}] Numeric. Zigzag amplitude.

\item[\code{std}] Logical. If TRUE, uses correlation for Phi (standardized visualization).

\item[\code{curve.base}] Base curvature parameter for Phi (distance-adjusted).
\end{ldescription}
\end{Arguments}
%
\begin{Examples}
\begin{ExampleCode}
#'
library(jjstat)
# 검증용 python data
df <- data.frame(
  x1 = c( 1.730516,  1.159965, -0.648269, -0.593655,  1.213306,  0.895361, -2.729817, -1.785706,  0.354516,  0.294027),
  x2 = c( 2.208429,  1.111993, -0.671010, -0.309758,  0.893759,  1.354703, -2.800578, -1.492191,  0.510968,  0.176220),
  x3 = c( 1.791936,  0.665374, -0.004232, -0.253564,  0.835220,  0.945070, -2.115333, -1.176772,  0.018319, -0.232088),
  y1 = c( 1.581619,  1.372783, -0.346918, -0.078805,  1.076466,  1.521254, -1.682830, -1.047084,  1.098595, -0.667011),
  y2 = c( 2.393086,  1.230622, -0.902185, -0.714770,  1.501231,  1.572681, -2.962813, -1.159968,  0.709137, -0.200424),
  y3 = c( 0.896092,  1.639311, -0.756823, -0.286587,  0.717714,  1.296822, -1.155036, -1.229061,  0.482557,  0.030295)
)

mydata<-df

mydata


latent2 <- list(
  n1 = c("x1","x2","x3"),
  n2 = c("y1","y2","y3")
)

res2 <- jjstat_SEM(
  data = mydata,
  latent_list = latent2,
  structural = list(
    path(from="n1", to="n2")
  ),
  residual_cov = list(c("y1","y2")),
  nboot = 300,
  rmseaCI = 0.90
)

res2

summary(res2)

res2$parameters

jjstat_plot_path(res2)
# jjstat_plot_path(res, show_phi = TRUE, show_theta = TRUE, curve = 2, layout="circle")
\end{ExampleCode}
\end{Examples}
\inputencoding{utf8}
\HeaderA{jjstat\_SEM}{Educational Structural Equation Modeling (Closed-form Engine)}{jjstat.Rul.SEM}
%
\begin{Description}
행렬 기반 수기계산 방식으로 구조방정식 모형을 추정하는
교육용 SEM 엔진이다. 잠재변수 점수(unit-weight),
OLS 회귀, 모형암시 공분산 행렬을 이용하여
ML 적합도를 근사하며, bootstrap 기반 표준오차와
RMSEA 신뢰구간을 제공한다.
\end{Description}
%
\begin{Usage}
\begin{verbatim}
jjstat_SEM(
  data,
  latent_list,
  structural,
  residual_cov = NULL,
  variance = NULL,
  data_format = c("data", "cov", "cor"),
  n_cov = 200,
  nboot = 200,
  seed = 2025,
  rmseaCI = 0.9
)
\end{verbatim}
\end{Usage}
%
\begin{Arguments}
\begin{ldescription}
\item[\code{data}] data.frame 또는 matrix.
\begin{itemize}

\item{} \code{data\_format="data"}: 관측자료
\item{} \code{data\_format="cov"}: 공분산 행렬
\item{} \code{data\_format="cor"}: 상관행렬

\end{itemize}


\item[\code{latent\_list}] named list.
잠재변수와 관측지표의 매핑 정보.

\item[\code{structural}] list of \code{path()}.
잠재변수 간 구조경로 정의.

\item[\code{residual\_cov}] list or NULL.
관측변수 잔차공분산 설정.

\item[\code{variance}] list or NULL.
관측·잠재변수 분산 고정 설정.

\item[\code{data\_format}] character.
"data", "cov", "cor" 중 하나.

\item[\code{n\_cov}] integer.
공분산/상관행렬 입력 시 생성할 표본 크기.

\item[\code{nboot}] integer.
bootstrap 반복 횟수.

\item[\code{seed}] integer.
난수 시드.

\item[\code{rmseaCI}] numeric.
RMSEA bootstrap 신뢰수준 (예: 0.90).
\end{ldescription}
\end{Arguments}
%
\begin{Details}
본 함수는 lavaan, OpenMx와 달리
\strong{모든 계산 단계를 명시적으로 드러내는 교육용 구현}이다.
수기계산 검증, SEM 이론 강의, 알고리즘 이해를 목적으로 한다.
\end{Details}
%
\begin{Value}
\code{jjstat\_sem} 클래스 객체.
\begin{itemize}

\item{} call: 함수 호출 정보
\item{} n: 표본 크기
\item{} Lambda, B, Phi, Theta, Sigma\_hat
\item{} fit: 적합도 지수
\item{} parameters: lavaan 스타일 파라미터 테이블
\item{} bootstrap: bootstrap 결과

\end{itemize}

\end{Value}
%
\begin{Examples}
\begin{ExampleCode}
latent <- list(
  F1 = c("y1","y2","y3"),
  F2 = c("y4","y5","y6")
)

res <- jjstat_SEM(
  data = HolzingerSwineford1939[,1:6],
  latent_list = latent,
  structural = list(path("F1","F2")),
  nboot = 50
)

summary(res)

\end{ExampleCode}
\end{Examples}
\inputencoding{utf8}
\HeaderA{jtl}{Haehyoung translate}{jtl}
%
\begin{Description}
Haehyoung translate
\end{Description}
%
\begin{Usage}
\begin{verbatim}
jtl(text = "", slang = "ko", elang = "en", translator = "deepl", print_url = F)
\end{verbatim}
\end{Usage}
%
\begin{Arguments}
\begin{ldescription}
\item[\code{text}] text to translate

\item[\code{slang}] starting language

\item[\code{elang}] ending language

\item[\code{translator}] translator type. lists are above."alibaba", "apertium", "argos", "baidu", "bing","caiyun", "cloudTranslation", "deepl", "elia", "google","iciba", "iflytek", "iflyrec", "itranslate", "judic", "languageWire", "lingvanex", "niutrans", "mglip", "mirai", "modernMt", "myMemory", "papago", "qqFanyi", "qqTranSmart", "reverso", "sogou", "sysTran", "tilde", "translateCom","translateMe", "utibet", "volcEngine", "yandex", "yeekit", "youdao"

\item[\code{print\_url}] Debugging Setting
\end{ldescription}
\end{Arguments}
\inputencoding{utf8}
\HeaderA{k}{Kakao translaste}{k}
%
\begin{Description}
Kakao translaste
\end{Description}
%
\begin{Usage}
\begin{verbatim}
k(text = "", slang = "en", elang = "ko", show = "data")
\end{verbatim}
\end{Usage}
%
\begin{Arguments}
\begin{ldescription}
\item[\code{text}] input text

\item[\code{slang}] "ko"

\item[\code{elang}] "en"

\item[\code{show}] nomal, data is each sentence
\end{ldescription}
\end{Arguments}
%
\begin{Value}
translate result
\end{Value}
\inputencoding{utf8}
\HeaderA{kakaoi}{Kakao translaste}{kakaoi}
%
\begin{Description}
Kakao translaste
\end{Description}
%
\begin{Usage}
\begin{verbatim}
kakaoi(text = "", slang = "en", elang = "ko", show = "data")
\end{verbatim}
\end{Usage}
%
\begin{Arguments}
\begin{ldescription}
\item[\code{text}] input text

\item[\code{slang}] "ko"

\item[\code{elang}] "en"

\item[\code{show}] nomal, data is each sentence
\end{ldescription}
\end{Arguments}
%
\begin{Value}
translate result
\end{Value}
\inputencoding{utf8}
\HeaderA{kge\_ca\_analysis}{Correspondence analysis visualization functions}{kge.Rul.ca.Rul.analysis}
%
\begin{Description}
Correspondence analysis visualization functions
\end{Description}
%
\begin{Usage}
\begin{verbatim}
kge_ca_analysis(
  dataset,
  typekey,
  selcol_3 = "w1f",
  type = "all",
  arrows = c(F, T),
  selcol_1 = "성조",
  selcol_2 = "onset",
  area = "부산",
  addtext = "",
  xlim = NULL,
  ylim = NULL,
  size_text = 5,
  opt = 1
)
\end{verbatim}
\end{Usage}
%
\begin{Arguments}
\begin{ldescription}
\item[\code{dataset}] data

\item[\code{typekey}] filter

\item[\code{selcol\_3}] col3

\item[\code{type}] all

\item[\code{arrows}] c(F,T)

\item[\code{selcol\_1}] col1

\item[\code{selcol\_2}] col2

\item[\code{area}] 지역

\item[\code{addtext}] add text

\item[\code{xlim}] range c(0, 0.3)

\item[\code{ylim}] rangec(-0.3, 0.4)

\item[\code{size\_text}] 5

\item[\code{opt}] opt 1\textasciitilde{} 6
\end{ldescription}
\end{Arguments}
\inputencoding{utf8}
\HeaderA{kge\_chisq\_table}{kge\_chisq\_table chisq table observed/Expected table-}{kge.Rul.chisq.Rul.table}
%
\begin{Description}
kge\_chisq\_table chisq table observed/Expected table-
\end{Description}
%
\begin{Usage}
\begin{verbatim}
kge_chisq_table(
  dataset,
  v1 = "a1",
  v2 = "성조",
  title = "Table",
  type = "res2",
  digits = 3,
  yadd = 0.4,
  Ncol = NULL,
  trans = FALSE,
  simple = FALSE,
  ko = TRUE,
  simulate.p.value = FALSE,
  correct = FALSE,
  size_bartext = 5,
  strip_size = 16,
  axis_size = 16,
  text_size = 13,
  xlab = "관측기대비율",
  ylab = "성조형",
  cramer = "adjust",
  raw = TRUE
)
\end{verbatim}
\end{Usage}
%
\begin{Arguments}
\begin{ldescription}
\item[\code{dataset}] data.frame

\item[\code{v1}] col

\item[\code{v2}] row

\item[\code{title}] title

\item[\code{type}] res2 default

\item[\code{digits}] 3

\item[\code{yadd}] up yadd

\item[\code{Ncol}] grid

\item[\code{trans}] transpose

\item[\code{simple}] * style

\item[\code{ko}] graph korean

\item[\code{simulate.p.value}] exact

\item[\code{correct}] exactt

\item[\code{size\_bartext}] graph text

\item[\code{strip\_size}] strip

\item[\code{axis\_size}] axis\_size

\item[\code{text\_size}] text\_size

\item[\code{xlab}] xlab

\item[\code{ylab}] ylab

\item[\code{cramer}] cramer cor

\item[\code{raw}] TRUE data.grame FALSE table data
\end{ldescription}
\end{Arguments}
%
\begin{Value}
multiple data
\end{Value}
%
\begin{Examples}
\begin{ExampleCode}

## Not run: 
kge_chisq_table(mtcars,"am","cyl")

# generate matrix
voteMat = matrix(c(183, 35, 824, 50), ncol=2)
colnames(voteMat) = c("Did not vote", "Voted")
rownames(voteMat) = c("All other ages","Youngest")

#chisq test result
vote_chi <- Chisq_retest(c(183, 35, 824, 50), ncol=2,
                        colname = c("Did not vote", "Voted"),
                         rowname =c("All other ages","Youngest"))

vote_chi

voteMat%>% data.frame() %>%
  long_df("voted", rowname="youth")%>% #long format
  add_rows_freq() %>%     # add rows
  kge_chisq_table("voted","youth")    #chisq test


## End(Not run)
\end{ExampleCode}
\end{Examples}
\inputencoding{utf8}
\HeaderA{kge\_weigth}{generative weight variable, kge(2024)}{kge.Rul.weigth}
%
\begin{Description}
generative weight variable, kge(2024)
\end{Description}
%
\begin{Usage}
\begin{verbatim}
kge_weigth(df, type = "res", pattern = "", remove = FALSE)
\end{verbatim}
\end{Usage}
%
\begin{Arguments}
\begin{ldescription}
\item[\code{df}] data.frame

\item[\code{type}] result res, res1, res2, res3, res4

\item[\code{pattern}] pattern = "", If you want something different, you can specify a different pattern.

\item[\code{remove}] ubite option
\end{ldescription}
\end{Arguments}
%
\begin{Value}
add weight
\end{Value}
%
\begin{Examples}
\begin{ExampleCode}


## Not run: 
 # first
 kge_bind0 %>% split_kw_df_match()
#데이터검증: 입력시 " "와 같은 빈칸들이 존재함.
#' kge_bind11 <- kge_bind0 %>% split_kw_df_match()
kge_bind11$A3 %>% unique()
kge_bind11$B3 %>% unique()
kge_bind11$C3 %>% unique()
kge_bind11$D3 %>% unique()
# kge_bind10$d3 %>% unique()
#데이터 전처리: 빈칸 규칙 체크하여 정리
kge_bind11$A3 <- kge_bind11$A3 %>% replace_df(imp = "", pattern = " ")
kge_bind11$B3 <- kge_bind11$B3 %>% replace_df(imp = "", pattern = " ")
kge_bind11$C3 <- kge_bind11$C3 %>% replace_df(imp = "", pattern = " ")
kge_bind11$D3 <- kge_bind11$D3 %>% replace_df(imp = "", pattern = " ")

# result
kge_bind11 %>% data.frame()

# same data
kge_bind10 %>% data.frame()

# apply kge_weight
kge_bind10%>% kge_weigth() %>%  data.frame()

kge_bind1<- kge_bind10%>% kge_weigth()

kge_bind1 %>% data.frame()
 kge_bind1 %>% data.frame() %>% nrow()  #4924

##
 kge_bind1 %>% filter(성조 != "") %>% nrow() #4360

kge_bind2 %>% data.frame() %>% nrow()#4360


## End(Not run)
\end{ExampleCode}
\end{Examples}
\inputencoding{utf8}
\HeaderA{kge\_weigth\_add\_ca}{k generative weight variable korean, kge(2024)}{kge.Rul.weigth.Rul.add.Rul.ca}
%
\begin{Description}
k generative weight variable korean, kge(2024)
\end{Description}
%
\begin{Usage}
\begin{verbatim}
kge_weigth_add_ca(df001, type = "res", pattern = "", remove = FALSE)
\end{verbatim}
\end{Usage}
%
\begin{Arguments}
\begin{ldescription}
\item[\code{type}] result res, res1, res2, res3, res4

\item[\code{pattern}] pattern = "", If you want something different, you can specify a different pattern.

\item[\code{remove}] ubite option

\item[\code{df}] data.frame
\end{ldescription}
\end{Arguments}
%
\begin{Value}
add weight
\end{Value}
%
\begin{Examples}
\begin{ExampleCode}


## Not run: 
 # first
 kge_bind0 %>% split_kw_df_match()
#데이터검증: 입력시 " "와 같은 빈칸들이 존재함.
#' kge_bind11 <- kge_bind0 %>% split_kw_df_match()
kge_bind11$A3 %>% unique()
kge_bind11$B3 %>% unique()
kge_bind11$C3 %>% unique()
kge_bind11$D3 %>% unique()
# kge_bind10$d3 %>% unique()
#데이터 전처리: 빈칸 규칙 체크하여 정리
kge_bind11$A3 <- kge_bind11$A3 %>% replace_df(imp = "", pattern = " ")
kge_bind11$B3 <- kge_bind11$B3 %>% replace_df(imp = "", pattern = " ")
kge_bind11$C3 <- kge_bind11$C3 %>% replace_df(imp = "", pattern = " ")
kge_bind11$D3 <- kge_bind11$D3 %>% replace_df(imp = "", pattern = " ")

# result
kge_bind11 %>% data.frame()

# same data
kge_bind10 %>% data.frame()

# apply kge_weight
kge_bind10%>% kge_weigth() %>%  data.frame()

kge_bind1<- kge_bind10%>% kge_weigth()

kge_bind1 %>% data.frame()
 kge_bind1 %>% data.frame() %>% nrow()  #4924

##
 kge_bind1 %>% filter(성조 != "") %>% nrow() #4360

kge_bind2 %>% data.frame() %>% nrow()#4360


## End(Not run)
\end{ExampleCode}
\end{Examples}
\inputencoding{utf8}
\HeaderA{kge\_weigth\_ko}{generative weight variable korean, kge(2024)}{kge.Rul.weigth.Rul.ko}
%
\begin{Description}
generative weight variable korean, kge(2024)
\end{Description}
%
\begin{Usage}
\begin{verbatim}
kge_weigth_ko(df, type = "res", pattern = "", remove = FALSE)
\end{verbatim}
\end{Usage}
%
\begin{Arguments}
\begin{ldescription}
\item[\code{df}] data.frame

\item[\code{type}] result res, res1, res2, res3, res4

\item[\code{pattern}] pattern = "", If you want something different, you can specify a different pattern.

\item[\code{remove}] ubite option
\end{ldescription}
\end{Arguments}
%
\begin{Value}
add weight
\end{Value}
%
\begin{Examples}
\begin{ExampleCode}


## Not run: 
 # first
 kge_bind0 %>% split_kw_df_match()
#데이터검증: 입력시 " "와 같은 빈칸들이 존재함.
#' kge_bind11 <- kge_bind0 %>% split_kw_df_match()
kge_bind11$A3 %>% unique()
kge_bind11$B3 %>% unique()
kge_bind11$C3 %>% unique()
kge_bind11$D3 %>% unique()
# kge_bind10$d3 %>% unique()
#데이터 전처리: 빈칸 규칙 체크하여 정리
kge_bind11$A3 <- kge_bind11$A3 %>% replace_df(imp = "", pattern = " ")
kge_bind11$B3 <- kge_bind11$B3 %>% replace_df(imp = "", pattern = " ")
kge_bind11$C3 <- kge_bind11$C3 %>% replace_df(imp = "", pattern = " ")
kge_bind11$D3 <- kge_bind11$D3 %>% replace_df(imp = "", pattern = " ")

# result
kge_bind11 %>% data.frame()

# same data
kge_bind10 %>% data.frame()

# apply kge_weight
kge_bind10%>% kge_weigth() %>%  data.frame()

kge_bind1<- kge_bind10%>% kge_weigth()

kge_bind1 %>% data.frame()
 kge_bind1 %>% data.frame() %>% nrow()  #4924

##
 kge_bind1 %>% filter(성조 != "") %>% nrow() #4360

kge_bind2 %>% data.frame() %>% nrow()#4360


## End(Not run)


\end{ExampleCode}
\end{Examples}
\inputencoding{utf8}
\HeaderA{KURT}{kurtosis 계산 함수}{KURT}
%
\begin{Description}
kurtosis 계산 함수
\end{Description}
%
\begin{Usage}
\begin{verbatim}
KURT(x)
\end{verbatim}
\end{Usage}
%
\begin{Arguments}
\begin{ldescription}
\item[\code{x}] vector
\end{ldescription}
\end{Arguments}
%
\begin{Value}
kurtosis
\end{Value}
\inputencoding{utf8}
\HeaderA{k\_df}{k df rep}{k.Rul.df}
%
\begin{Description}
k df rep
\end{Description}
%
\begin{Usage}
\begin{verbatim}
k_df(df, range = 1:ncol(df), slang = "en", elang = "ko")
\end{verbatim}
\end{Usage}
%
\begin{Arguments}
\begin{ldescription}
\item[\code{df}] data.frame

\item[\code{range}] range

\item[\code{slang}] en

\item[\code{elang}] ko
\end{ldescription}
\end{Arguments}
%
\begin{Value}
text
\end{Value}
%
\begin{Examples}
\begin{ExampleCode}

## Not run: 

#'
Authentic_Leadership <- data.frame(
  id = paste0("Q",1:16),
  factor= c(
    rep("self Awareness", 4),
    rep("Relational Transparency", 4),
    rep("Internalized Moral perspective", 4),
    rep("Balanced Processing", 4)
  ),
  Question = c(
    "I can list my three greatest weaknesses.",
    "My actions reflect my core values.",
    "I seek others' opinions before making up my own mind.",
    "I openly share my feelings with others.",
    "I can list my three greatest strengths.",
    "I do not allow group pressure to control me.",
    "I listen closely to the ideas of those who disagree with me.",
    "I let others know who I truly am as a person.",
    "I seek feedback as a way of understanding who I really am as a person.",
    "Other people know where I stand on controversial issues.",
    "I do not emphasize my own point of view at the expense of others.",
    "I rarely present a “false” front to others.",
    "I accept the feelings I have about myself.",
    "My morals guide what I do as a leader.",
    "I listen very carefully to the ideas of others before making decisions.",
    "I admit my mistakes to others."
  ))%>%tibble()
Authentic_Leadership
# 데이터프레임 번역
Alp_ko = k_df(Authentic_Leadership,
              range = 2:3,
              slang = "en",
              elang = "ko") %>%select(3,2)
Alp_ko
Alp_ko %>%replace_df_rep("받아들입니다.","받아들인다.", "압니다","안다")
#'






## End(Not run)




\end{ExampleCode}
\end{Examples}
\inputencoding{utf8}
\HeaderA{LavaanPlot}{Lavaan plot draw}{LavaanPlot}
%
\begin{Description}
Lavaan plot draw
\end{Description}
%
\begin{Usage}
\begin{verbatim}
LavaanPlot(
  x,
  rotation = "LR",
  coefs = TRUE,
  col = "gray60",
  stand = FALSE,
  label = NULL,
  covs = FALSE,
  stars = c("covs", "latent", "regress"),
  overlap = FALSE,
  layout = "dot",
  digits = 3
)
\end{verbatim}
\end{Usage}
%
\begin{Arguments}
\begin{ldescription}
\item[\code{x}] lavaan result

\item[\code{rotation}] rotation 'LR','RL', 'TB','BT'

\item[\code{coefs}] TRUE

\item[\code{col}] default 'gray60'

\item[\code{stand}] std coeff

\item[\code{label}] input name labels=c(f1\_1="exp", f1\_2="want",f1\_3="problem") etc

\item[\code{covs}] TRUE

\item[\code{stars}] stars = c("covs", "latent","regress"),

\item[\code{overlap}] default FALSE

\item[\code{layout}] 'dot', 'neato', 'fdp', 'sfdp', 'circo', 'twopi', 'nop','nop2', 'osage', 'patchwork'

\item[\code{digits}] digits = 3\#'
\end{ldescription}
\end{Arguments}
%
\begin{Value}
plot
\end{Value}
%
\begin{Examples}
\begin{ExampleCode}
# lavaan code

## Bollen (1989), page 332
model <- '
  # latent variable definitions
     ind60 =~ x1 + x2 + x3
     dem60 =~ y1 + a*y2 + b*y3 + c*y4
     dem65 =~ y5 + a*y6 + b*y7 + c*y8

  # regressions
    dem60 ~ ind60
  dem65 ~ ind60 + dem60

  # residual correlations
    y1 ~~ y5
    y2 ~~ y4 + y6
    y3 ~~ y7
    y4 ~~ y8
    y6 ~~ y8
'

fit <- lavaan::sem(model, data = PoliticalDemocracy)
 ## summary
summary(fit, fit.measures = TRUE, stand=T)
##
LavaanPlot(fit, cov=F)
LavaanPlot(fit, cov=F, layout="circo")

LavaanPlot(fit, cov=F, layout="dot",
           label=c(x1="ind1",ind60="IND60"))


\end{ExampleCode}
\end{Examples}
\inputencoding{utf8}
\HeaderA{lav\_change\_node\_name}{Change node names in semPlot / qgraph object}{lav.Rul.change.Rul.node.Rul.name}
%
\begin{Description}
semPlot::semPaths()로 생성된 구조방정식 시각화 객체에서
노드 이름을 사용자가 지정한 (원래이름, 변경라벨) 쌍에 따라
일괄적으로 변경한다.
본 함수는 pipe(
\end{Description}
%
\begin{Usage}
\begin{verbatim}
lav_change_node_name(plot_obj, var_named)
\end{verbatim}
\end{Usage}
%
\begin{Arguments}
\begin{ldescription}
\item[\code{plot\_obj}] semPlot 또는 qgraph 객체

\item[\code{var\_named}] character vector.
c(old1, new1, old2, new2, ...)
\end{ldescription}
\end{Arguments}
%
\begin{Value}
노드 명칭이 변경된 semPlot 객체
\end{Value}
%
\begin{Examples}
\begin{ExampleCode}
## Not run: 

#################################
# 1. SCCT 구조방정식 모형 정의
#################################

model_scct_indirect <- '

REFF =~ efficacy2 + efficacy3 + efficacy4
OEXP =~ Outcome1 + Outcome2 + Outcome3 + Outcome4 + Outcome5 + Outcome8
INTR =~ Interest1 + Interest2 + Interest3 + Interest4

REFF ~ e1*grade + e2*experience0
OEXP ~ a1*REFF + d1*grade + d2*experience0
INTR ~ b1*REFF + b2*OEXP

intention ~ c1*INTR + c2*OEXP + c3*REFF

grade ~~ experience0

# 간접효과
ind_grade_REFF        := e1 * c3
ind_grade_REFF_INTR   := e1 * b1 * c1
ind_grade_REFF_OEXP   := e1 * a1 * c2
ind_grade_REFF_OEXP_INTR := e1 * a1 * b2 * c1

ind_exp_REFF          := e2 * c3
ind_exp_REFF_INTR     := e2 * b1 * c1
ind_exp_REFF_OEXP     := e2 * a1 * c2
ind_exp_REFF_OEXP_INTR := e2 * a1 * b2 * c1
'

#################################
# 2. 모형 추정
#################################

library(lavaan)
fit_scct_indirect <- sem(
  model = model_scct_indirect,
  data  = shindata,
  se = "bootstrap",
  bootstrap = 2000
)

#################################
# 3. 노드 명칭 매핑 벡터
#################################

var_named1 <- c(
  "REFF", "연구효능감",
  "efficacy2","연구효능2",
  "efficacy3","연구효능3",
  "efficacy4","연구효능4",

  "INTR", "진로흥미",
  "Interest1","진로흥미1",
  "Interest2","진로흥미2",
  "Interest3","진로흥미3",
  "Interest4","진로흥미4",

  "OEXP", "결과기대",
  "Outcome1","결과기대1",
  "Outcome2","결과기대2",
  "Outcome3","결과기대3",
  "Outcome4","결과기대4",
  "Outcome5","결과기대5",
  "Outcome8","결과기대8",

  "grade","학년",
  "experience0","연구경험",
  "intention","진로의지"
)

#################################
# 4. SCCT 전용 layout 행렬
#################################

layout_scct <- matrix(
  c(
    NA, "efficacy2", "efficacy3", "efficacy4", NA,
        "Interest1", "Interest2", "Interest3", "Interest4", NA,

    "experience0", NA, "REFF", NA, NA,
        NA, "INTR", NA, NA, NA,

    "grade", NA, "OEXP", NA, NA,
        NA, NA, "intention", NA, NA,

    NA, "Outcome1", "Outcome2", "Outcome3", "Outcome4",
        "Outcome5", "Outcome8", NA, NA, NA
  ),
  nrow = 4,
  byrow = TRUE
)

#################################
# 5. 구조방정식 그림 + 노드명 변경
#################################

library(semPlot)
library(semptools)

semPlot::semPaths(
  fit_scct_indirect,
  layout        = layout_scct,
  whatLabels    = "std",
  style         = "ram",
  residuals     = FALSE,
  intercepts    = FALSE,
  nCharNodes    = 0,
  sizeLat       = 9,
  sizeMan       = 9,
  edge.label.cex = 0.85,
  edge.label.position = 0.65
) %>%
  semptools::mark_sig(fit_scct_indirect) %>%
  change_node_name(var_named1) %>%  ##### 이것을 이용하여 이름 변경
  plot()


## End(Not run)
\end{ExampleCode}
\end{Examples}
\inputencoding{utf8}
\HeaderA{lav\_de\_hypo}{Hypothesis generation functions}{lav.Rul.de.Rul.hypo}
%
\begin{Description}
Hypothesis generation functions
\end{Description}
%
\begin{Usage}
\begin{verbatim}
lav_de_hypo(model)
\end{verbatim}
\end{Usage}
%
\begin{Arguments}
\begin{ldescription}
\item[\code{model}] lavaan syntax
\end{ldescription}
\end{Arguments}
%
\begin{Value}
text
\end{Value}
%
\begin{Examples}
\begin{ExampleCode}

## Not run: 
#'
# 예제 모델
model1a <- "
EXPE ~ H1*IMAG
QUAL ~ H2*EXPE
VAL ~  H3*QUAL + H4*EXPE
SAT ~ H5*IMAG + H6*EXPE + H7*QUAL + H8*VAL
LOY ~ H9*SAT + H10*IMAG
"

# 함수 실행
lav_de_hypo(model1a)

#'

## End(Not run)

\end{ExampleCode}
\end{Examples}
\inputencoding{utf8}
\HeaderA{lav\_edge\_value}{Extract edge labels for semPlot from lavaan / MLSEM object}{lav.Rul.edge.Rul.value}
%
\begin{Description}
lavaan(일반 SEM/CFA) 또는 다층 SEM(ML-SEM) 적합 객체에서
semPlot::semPaths()의 edgeLabels 인수에 바로 넣을 수 있는
"간선 라벨(edge label) 벡터"를 자동 생성하는 함수이다.

본 함수는 lavaan::parameterEstimates()를 기반으로
측정경로(=\textasciitilde{}), 구조경로(\textasciitilde{}), 공분산(\textasciitilde{}\textasciitilde{})을 모두 인식한다.
다층모형의 경우 within/between 수준(block)을 선택할 수 있다.

핵심 설계 목표는 다음과 같다.
(1) 기준지표(fixed loading=1)에서 pvalue가 NA로 반환되는 경우에도
라벨에 "NA"가 붙지 않도록 처리한다.
(2) semPlot이 요구하는 형식(unnamed character vector)에 맞게 안전 반환한다.
(3) 측정경로/공분산 값을 숨기거나, 유의성 별표를 제거하는 옵션을 제공한다.
(4) 구조경로(\textasciitilde{})만 뽑아 구조모형 라벨에만 쓰는 옵션(path\_only)을 제공한다.
\end{Description}
%
\begin{Usage}
\begin{verbatim}
lav_edge_value(
  obj,
  level = c("within", "between"),
  type = c("est", "std"),
  sig = TRUE,
  digits = 3,
  include_cov = TRUE,
  fillin = FALSE,
  hide_mea = "",
  hide_cov = "",
  hide_sig = FALSE,
  hide_sig_mea = "",
  hide_sig_cov = "",
  path_only = FALSE
)
\end{verbatim}
\end{Usage}
%
\begin{Arguments}
\begin{ldescription}
\item[\code{obj}] lavaan fitted object. (cfa/sem 또는 다층 sem 객체)

\item[\code{level}] character. "within" 또는 "between".
다층모형에서만 의미가 있으며, within=block 1, between=block 2로 처리한다.
일반 lavaan 단일수준 모형에서는 block이 0 또는 1로 나올 수 있어,
이 경우 level을 "within"으로 두는 것이 보통 안전하다.

\item[\code{type}] character. "est"(비표준화 추정치) 또는 "std"(표준화 계수 std.all).

\item[\code{sig}] logical. 유의성 별표(*, **, ***)를 붙일지 여부. 기본 TRUE.
단, lavaan에서 pvalue가 NA인 경우(예: 기준지표 fixed loading=1)는
별표를 붙이지 않는다(NA-safe).

\item[\code{digits}] numeric. 반올림 소수점 자리수. 기본 3.

\item[\code{include\_cov}] logical. 공분산(\textasciitilde{}\textasciitilde{})을 포함할지 여부. 기본 TRUE.
FALSE이면 \textasciitilde{}\textasciitilde{}는 제외된다.

\item[\code{fillin}] logical. TRUE이면 측정경로(=\textasciitilde{}) 및 공분산(\textasciitilde{}\textasciitilde{}) 라벨을
사용자가 지정한 문자로 "대체"한다. (예: 측정경로는 '지표'로 숨기고,
구조경로(\textasciitilde{})만 수치 라벨로 보이게)

\item[\code{hide\_mea}] character. fillin=TRUE일 때 =\textasciitilde{} 경로 라벨을 이 문자로 대체.
기본 ""(빈 문자열).

\item[\code{hide\_cov}] character. fillin=TRUE일 때 \textasciitilde{}\textasciitilde{} 경로 라벨을 이 문자로 대체.
기본 ""(빈 문자열).

\item[\code{hide\_sig}] logical. TRUE이면 모든 유의성 별표를 제거한다.
주의: hide\_sig=TRUE이면 sig 인수는 사실상 무시된다(별표 제거가 우선).

\item[\code{hide\_sig\_mea}] character. hide\_sig=TRUE일 때,
=\textasciitilde{} 경로에 대해 특정 문자열로 대체하고 싶을 때 사용한다.
기본 ""(대체 안 함).

\item[\code{hide\_sig\_cov}] character. hide\_sig=TRUE일 때,
\textasciitilde{}\textasciitilde{} 경로에 대해 특정 문자열로 대체하고 싶을 때 사용한다.
기본 ""(대체 안 함).

\item[\code{path\_only}] logical. TRUE이면 구조경로(\textasciitilde{})만 추출한다.
구조모형만 그리고 싶을 때 사용한다.
이 경우 fillin 옵션은 적용되지 않는다(측정/공분산을 뽑지 않기 때문).
\end{ldescription}
\end{Arguments}
%
\begin{Value}
unnamed character vector.
semPlot::semPaths(edgeLabels= )에 그대로 투입 가능한 형식이다.
\end{Value}
%
\begin{Examples}
\begin{ExampleCode}
## Not run: 
# ============================================================
# [0] 전제: fit_mlsem_expB, model_within_plot, diagram_model()이 존재한다고 가정
# ============================================================

# ------------------------------------------------------------
# (1) 기본 사용: within 수준, 표준화 계수(std) + 유의성 별표
# ------------------------------------------------------------
edge_std <- lav_edge_value(
  obj   = fit_mlsem_expB,
  level = "within",
  type  = "std"
)
edge_std

# diagram_model에 그대로 주입
diagram_model(
  model      = model_within_plot,
  layout     = layout_scct1,
  whatLabels = "none",
  sizeMan    = 8,
  sizeLat    = 10,
  edge.label.position = 0.7,
  edge.label.cex      = 0.9,
  edgeLabels = edge_std,
  color = list(lat="black", man="White")
)

# ------------------------------------------------------------
# (2) 비표준화 계수(est) 라벨
# ------------------------------------------------------------
edge_est <- lav_edge_value(
  obj   = fit_mlsem_expB,
  level = "within",
  type  = "est"
)
edge_est

diagram_model(
  model      = model_within_plot,
  layout     = layout_scct1,
  whatLabels = "none",
  edge.label.position = 0.7,
  edge.label.cex      = 0.9,
  edgeLabels = edge_est
)

# ------------------------------------------------------------
# (3) 측정경로(=~)를 숨기고 구조경로(~)만 숫자로 보여주기
#     - fillin=TRUE + hide_mea="지표"
#     - 공분산(~~)은 빈문자("")로 숨김
# ------------------------------------------------------------
edge_std_path <- lav_edge_value(
  obj      = fit_mlsem_expB,
  level    = "within",
  type     = "std",
  fillin   = TRUE,
  hide_mea = "지표",
  hide_cov = ""
)

diagram_model(
  model      = model_within_plot,
  layout     = layout_scct1,
  whatLabels = "none",
  edge.label.position = 0.6,
  edge.label.cex      = 0.9,
  edgeLabels = edge_std_path
)

# ------------------------------------------------------------
# (4) 유의성 별표(*,**,***)를 모두 숨기고 숫자만 표시
#     - hide_sig=TRUE가 최우선
# ------------------------------------------------------------
edge_std_nostar <- lav_edge_value(
  obj      = fit_mlsem_expB,
  level    = "within",
  type     = "std",
  hide_sig = TRUE
)

diagram_model(
  model      = model_within_plot,
  layout     = layout_scct1,
  whatLabels = "none",
  edgeLabels = edge_std_nostar
)

# ------------------------------------------------------------
# (5) 구조경로(~)만 추출: 구조모형 라벨만 뽑기
#     - 측정경로/공분산 라벨이 전혀 포함되지 않는다.
# ------------------------------------------------------------
edge_struct_only <- lav_edge_value(
  obj       = fit_mlsem_expB,
  level     = "within",
  type      = "std",
  path_only = TRUE
)
edge_struct_only

# (주의) diagram_model()에서 structural=TRUE를 쓰면
# semPlot이 일부 변수를 자동 제외할 수 있어 레이아웃과 충돌 가능성이 있다.
# 이 경우에는 whatLabels="none" + fillin을 쓰는 방식이 더 안정적일 수 있다.

# ------------------------------------------------------------
# (6) between 수준 라벨 추출 (ML-SEM)
# ------------------------------------------------------------
edge_between <- lav_edge_value(
  obj   = fit_mlsem_expB,
  level = "between",
  type  = "std",
  fillin   = TRUE,
  hide_mea = "지표",
  hide_cov = ""
)
edge_between

## End(Not run)
\end{ExampleCode}
\end{Examples}
\inputencoding{utf8}
\HeaderA{lav\_extract\_imlist}{Function to extract interaction terms from the lavaan syntax, excluding commented lines.}{lav.Rul.extract.Rul.imlist}
%
\begin{Description}
Function to extract interaction terms from the lavaan syntax,
excluding commented lines.
\end{Description}
%
\begin{Usage}
\begin{verbatim}
lav_extract_imlist(lav_syn, prefix = ".")
\end{verbatim}
\end{Usage}
%
\begin{Arguments}
\begin{ldescription}
\item[\code{lav\_syn}] lavaan syntax

\item[\code{prefix}] . dataname prefix
\end{ldescription}
\end{Arguments}
%
\begin{Value}
list
\end{Value}
%
\begin{Examples}
\begin{ExampleCode}
## Not run: 
#'

# # Example usage:
data(MarshWenHau)
str(MarshWenHau)

mwh = MarshWenHau

model2 <- "
f1 =~ x1 + x2 + x3
f1Xf2 =~ f1:f2
f2 =~ x4 + x5 + x6
f3 =~ y1 + y2 + y3

f3 ~ f1 + f2
f3 ~ f1Xf2

f1Xf2 ~~ 0*f1
f1Xf2 ~~ 0*f2
"

# interaction model
intsem =  SEM(model2, data= MarshWenHau)
summary(intsem)
intsem %>%lav_semPaths2(color=list(lat="darkred",man="gold"))


##2nd
#base model
lav_extract_mm(model2) %>%cat()
#interaction item list
lav_extract_imlist(model3)
#data
lav_latentProd(mwh, model3)%>%head()

#extracted interaction term
lav_extract_int(model2, lav_extract_imlist(model3))

# or interaction items ; imlist
imlist <- list(
  .f1Xf2 = c("f1", "f2"),
  .i_f1_f4 = c("f1", "f4")
)
# names(imlist)
lav_extract_int(model2, imlist)
#direct
lav_extract_int(model2, imlist= list(.f1Xf2 = c("f1", "f2"),
                                     .i_f1_f4 = c("f1", "f4")))
#estimates model
intsem =  SEM(lav_new_model(model2, lav_extract_imlist(model2)),
              data= lav_latentProd(mwh, model2))


## End(Not run)

\end{ExampleCode}
\end{Examples}
\inputencoding{utf8}
\HeaderA{lav\_extract\_ind}{lav\_extract\_ind, indirect effect}{lav.Rul.extract.Rul.ind}
%
\begin{Description}
lav\_extract\_ind, indirect effect
\end{Description}
%
\begin{Usage}
\begin{verbatim}
lav_extract_ind(
  model,
  start_node = NULL,
  end_node = NULL,
  prefix = "a",
  cat = FALSE,
  type = "ind",
  auto = TRUE,
  paths_name = FALSE,
  n_name = 1
)
\end{verbatim}
\end{Usage}
%
\begin{Arguments}
\begin{ldescription}
\item[\code{model}] lavaan

\item[\code{start\_node}] path start node

\item[\code{end\_node}] path end node

\item[\code{prefix}] hypothesis "a" if U want change "H"

\item[\code{cat}] test show cat TRUR

\item[\code{type}] type "ind","new\_parameters","combine"

\item[\code{auto}] TRUE path auto, FLASE all ind path

\item[\code{paths\_name}] TRUE change paths name

\item[\code{n\_name}] Number of letters in the name being extracted, default =1
\end{ldescription}
\end{Arguments}
%
\begin{Value}
text
\end{Value}
%
\begin{Examples}
\begin{ExampleCode}


## Not run: 
#'

# Example usage
model1 <- "
IMAG =~ imag1 + imag2 + imag3 + imag4 + imag5
EXPE =~ expe1 + expe2 + expe3 + expe4 + expe5
QUAL =~ qual1 + qual2 + qual3 + qual4 + qual5
VAL =~ val1 + val2 + val3 + val4
SAT =~ sat1 + sat2 + sat3 + sat4
LOY =~ loy1 + loy2 + loy3 + loy4

EXPE ~ IMAG
QUAL ~ EXPE
VAL ~  QUAL+ EXPE
SAT ~ IMAG + EXPE + QUAL + VAL
LOY ~ SAT + IMAG
"


# Extracting paths from the model with default start_node = NULL and end_node = NULL
lav_extract_mm(model1) %>%cat("\n")
lav_extract_sm(model1) %>%cat("\n")

lav_extract_ind(model1)%>%cat("\n")
lav_extract_ind(model1, auto=TRUE)%>%cat("\n")

lav_extract_ind(model1,start_node="EXPE", auto=TRUE)%>%cat("\n")
lav_extract_ind(model1,start_node="EXPE")%>%cat("\n")
lav_extract_ind(model1,start_node="QUAL")%>%cat("\n")
lav_extract_ind(model1,start_node="IMAG", auto=TRUE)%>%cat("\n")

lav_extract_ind(lav_extract_sm(model1), de=TRUE) %>%cat("\n")

lav_extract_ind(lav_extract_sm(model1), de=TRUE, auto=TRUE)%>%cat("\n")
lav_extract_ind(lav_extract_sm(model1), start_node="IMAG", de=TRUE) %>%cat('\n')

lav_extract_ind(lav_extract_sm(model1),"IMAG", prefix="H", de=TRUE) %>%cat('\n')

# Example usage
de_paths <- "
EXPE  ~ a1* IMAG
QUAL  ~ a2* EXPE
VAL  ~ a3* QUAL
VAL  ~ a4* EXPE
SAT  ~ a5* IMAG
SAT  ~ a6* EXPE
SAT  ~ a7* QUAL
SAT  ~ a8* VAL
LOY  ~ a9* SAT
LOY  ~ a10* IMAG
"

lav_extract_ind(de_paths)
#출력
#  lav_extract_ind(model1, cat=TRUE)

lav_extract_ind(de_paths, cat=TRUE)

lav_extract_ind(de_paths)%>%make_c_text()

## End(Not run)

\end{ExampleCode}
\end{Examples}
\inputencoding{utf8}
\HeaderA{lav\_extract\_ind\_paths}{lav\_extract\_ind\_paths, Create an indirect effect hypothesis}{lav.Rul.extract.Rul.ind.Rul.paths}
%
\begin{Description}
lav\_extract\_ind\_paths, Create an indirect effect hypothesis
\end{Description}
%
\begin{Usage}
\begin{verbatim}
lav_extract_ind_paths(
  paths_data,
  first_node = NULL,
  paths_name = FALSE,
  allpaths = FALSE
)
\end{verbatim}
\end{Usage}
%
\begin{Arguments}
\begin{ldescription}
\item[\code{paths\_data}] direct path

\item[\code{first\_node}] start node name

\item[\code{paths\_name}] TRUE hypothesis output

\item[\code{allpaths}] all and ind(FALSE)
\end{ldescription}
\end{Arguments}
%
\begin{Value}
text
\end{Value}
%
\begin{Examples}
\begin{ExampleCode}

## Not run: 
#'
# 예시 데이터
paths_data <- data.frame(
  hypo = c("H1", "H2", "H3", "H4", "H5",
  "H6", "H7", "H8", "H9", "H10"),
  paths = c("IMAG->EXPE", "EXPE->QUAL",
  "QUAL->VAL", "EXPE->VAL", "IMAG->SAT",
   "EXPE->SAT", "QUAL->SAT", "VAL->SAT", "SAT->LOY", "IMAG->LOY")
)
# 함수 실행
lav_extract_ind_paths(paths_data, allpaths = TRUE)
lav_extract_ind_paths(paths_data, allpaths = FALSE)
#'
lav_extract_ind_paths(paths_data, allpaths = FALSE, paths_name = TRUE)
lav_extract_ind_paths(paths_data, allpaths = FALSE, paths_name = FALSE)
#'
lav_extract_path(model1a)%>%lav_extract_ind_paths()

lav_extract_path(model1)%>%lav_extract_ind_paths()
#'
lav_extract_path(model1a) %>%
  lav_extract_ind_paths() %>%
  sem_lav_ind_hypo()
#'
## 구조모형에서 나타나는 간접효과 가설은 다음과 같다.
## 간접효과[H1]: 간접효과 IMAG -> SAT ->
LOY 는 통계적으로 유의할 것이다.
## 간접효과[H2]: 간접효과 IMAG -> EXPE -
> SAT -> LOY 는 통계적으로 유의할 것이다.
## 간접효과[H3]: 간접효과 IMAG -> EXPE -
> VAL -> SAT -> LOY 는 통계적으로 유의할 것이다.
## 간접효과[H4]: 간접효과 IMAG -> EXPE
-> QUAL -> SAT -> LOY 는 통계적으로 유의할 것이다.
## 간접효과[H5]: 간접효과 IMAG -> EXPE
-> QUAL -> VAL -> SAT -> LOY 는 통계적으로 유의할 것이다.
#' #'
#'

# Example usage with the provided model
model10 <- "
IMAG =~ imag1 + imag2 + imag3 + imag4 + imag5
EXPE =~ expe1 + expe2 + expe3 + expe4 + expe5
QUAL =~ qual1 + qual2 + qual3 + qual4 + qual5
VAL =~ val1 + val2 + val3 + val4
SAT =~ sat1 + sat2 + sat3 + sat4
LOY =~ loy1 + loy2 + loy3 + loy4

EXPE ~ IMAG
QUAL ~ EXPE
VAL ~  QUAL+ EXPE
SAT ~ IMAG + EXPE + QUAL + VAL
LOY ~ SAT + IMAG
"


# Extract paths from the model
lav_extract_path(model10)
lav_extract_path(model10)%>%lav_extract_ind_paths()

lav_extract_sm(model10)%>%lav_extract_path()%>%lav_extract_ind_paths()
lav_extract_sm(model10)%>%lav_extract_ind_paths()
lav_extract_sm(model10)%>%lav_extract_ind_paths("EXPE")


## End(Not run)

\end{ExampleCode}
\end{Examples}
\inputencoding{utf8}
\HeaderA{lav\_extract\_int}{lav\_extract\_int}{lav.Rul.extract.Rul.int}
%
\begin{Description}
lav\_extract\_int
\end{Description}
%
\begin{Usage}
\begin{verbatim}
lav_extract_int(lav_syn, imlist = NULL)
\end{verbatim}
\end{Usage}
%
\begin{Arguments}
\begin{ldescription}
\item[\code{lav\_syn}] lavaan

\item[\code{imlist}] item list
\end{ldescription}
\end{Arguments}
%
\begin{Value}
interation data
\end{Value}
%
\begin{Examples}
\begin{ExampleCode}

## Not run: 
#'

# # Example usage:
data(MarshWenHau)
str(MarshWenHau)

mwh = MarshWenHau

model2 <- "
f1 =~ x1 + x2 + x3
f1Xf2 =~ f1:f2
f2 =~ x4 + x5 + x6
f3 =~ y1 + y2 + y3

f3 ~ f1 + f2
f3 ~ f1Xf2

f1Xf2 ~~ 0*f1
f1Xf2 ~~ 0*f2
"

# interaction model
intsem =  SEM(model2, data= MarshWenHau)
summary(intsem)
intsem %>%lav_semPaths2(color=list(lat="darkred",man="gold"))


##2nd
#base model
lav_extract_mm(model2) %>%cat()
#interaction item list
lav_extract_imlist(model3)
#data
lav_latentProd(mwh, model3)%>%head()

#extracted interaction term
lav_extract_int(model2, lav_extract_imlist(model3))

# or interaction items ; imlist
imlist <- list(
  .f1Xf2 = c("f1", "f2"),
  .i_f1_f4 = c("f1", "f4")
)
# names(imlist)
lav_extract_int(model2, imlist)
#direct
lav_extract_int(model2, imlist= list(.f1Xf2 = c("f1", "f2"),
                                     .i_f1_f4 = c("f1", "f4")))
#estimates model
intsem =  SEM(lav_new_model(model2, lav_extract_imlist(model2)),
              data= lav_latentProd(mwh, model2))


## End(Not run)

\end{ExampleCode}
\end{Examples}
\inputencoding{utf8}
\HeaderA{lav\_extract\_mm}{extranct measurement}{lav.Rul.extract.Rul.mm}
%
\begin{Description}
extranct measurement
\end{Description}
%
\begin{Usage}
\begin{verbatim}
lav_extract_mm(lav_syn, text = FALSE)
\end{verbatim}
\end{Usage}
%
\begin{Arguments}
\begin{ldescription}
\item[\code{lav\_syn}] model lavaan syntax

\item[\code{text}] text output
\end{ldescription}
\end{Arguments}
%
\begin{Value}
text model
\end{Value}
%
\begin{Examples}
\begin{ExampleCode}

## Not run: 
#'

# # Example usage:
data(MarshWenHau)
str(MarshWenHau)

mwh = MarshWenHau

model2 <- "
f1 =~ x1 + x2 + x3
f1Xf2 =~ f1:f2
f2 =~ x4 + x5 + x6
f3 =~ y1 + y2 + y3

f3 ~ f1 + f2
f3 ~ f1Xf2

f1Xf2 ~~ 0*f1
f1Xf2 ~~ 0*f2
"

# interaction model
intsem =  SEM(model2, data= MarshWenHau)
summary(intsem)
intsem %>%lav_semPaths2(color=list(lat="darkred",man="gold"))


##2nd
#base model
lav_extract_mm(model2) %>%cat()
#interaction item list
lav_extract_imlist(model3)
#data
lav_latentProd(mwh, model3)%>%head()

#extracted interaction term
lav_extract_int(model2, lav_extract_imlist(model3))

# or interaction items ; imlist
imlist <- list(
  .f1Xf2 = c("f1", "f2"),
  .i_f1_f4 = c("f1", "f4")
)
# names(imlist)
lav_extract_int(model2, imlist)
#direct
lav_extract_int(model2, imlist= list(.f1Xf2 = c("f1", "f2"),
                                     .i_f1_f4 = c("f1", "f4")))
#estimates model
intsem =  SEM(lav_new_model(model2, lav_extract_imlist(model2)),
              data= lav_latentProd(mwh, model2))


## End(Not run)

\end{ExampleCode}
\end{Examples}
\inputencoding{utf8}
\HeaderA{lav\_extract\_mm2}{extranct measurement, Extract all structural models}{lav.Rul.extract.Rul.mm2}
%
\begin{Description}
extranct measurement, Extract all structural models
\end{Description}
%
\begin{Usage}
\begin{verbatim}
lav_extract_mm2(lav_syn, text = FALSE, interact = TRUE)
\end{verbatim}
\end{Usage}
%
\begin{Arguments}
\begin{ldescription}
\item[\code{lav\_syn}] model lavaan syntax

\item[\code{text}] text output

\item[\code{interact}] FALSE exclude :, TRue include :
\end{ldescription}
\end{Arguments}
%
\begin{Value}
text model
\end{Value}
%
\begin{Examples}
\begin{ExampleCode}

## Not run: 
#'

# # Example usage:
data(MarshWenHau)
str(MarshWenHau)

mwh = MarshWenHau

model2 <- "
f1 =~ x1 + x2 + x3
f1Xf2 =~ f1:f2
f2 =~ x4 + x5 + x6
f3 =~ y1 + y2 + y3

f3 ~ f1 + f2
f3 ~ f1Xf2

f1Xf2 ~~ 0*f1
f1Xf2 ~~ 0*f2
"

# interaction model
intsem =  SEM(model2, data= MarshWenHau)
summary(intsem)
intsem %>%lav_semPaths2(color=list(lat="darkred",man="gold"))


##2nd
#base model
lav_extract_mm(model2) %>%cat()
#interaction item list
lav_extract_imlist(model3)
#data
lav_latentProd(mwh, model3)%>%head()

#extracted interaction term
lav_extract_int(model2, lav_extract_imlist(model3))

# or interaction items ; imlist
imlist <- list(
  .f1Xf2 = c("f1", "f2"),
  .i_f1_f4 = c("f1", "f4")
)
# names(imlist)
lav_extract_int(model2, imlist)
#direct
lav_extract_int(model2, imlist= list(.f1Xf2 = c("f1", "f2"),
                                     .i_f1_f4 = c("f1", "f4")))
#estimates model
intsem =  SEM(lav_new_model(model2, lav_extract_imlist(model2)),
              data= lav_latentProd(mwh, model2))

#'
model2 <- "
 f1 =~ x1 + x2 + x3
 f1Xf2 =~ f1:f2
 f2 =~ x4 + x5 + x6
 f3 =~ y1 + y2 + y3

 f3 ~ f1 + f2
 f3 ~ f1Xf2

 f1Xf2 ~~ 0*f1
 f1Xf2 ~~ 0*f2
 "

lav_extract_mm2(model2)
lav_extract_mm2(model2, interact=F)
lav_extract_mm2(model2, interact=T)
#'
#'

## End(Not run)

\end{ExampleCode}
\end{Examples}
\inputencoding{utf8}
\HeaderA{lav\_extract\_path}{lav\_extract\_path}{lav.Rul.extract.Rul.path}
%
\begin{Description}
lav\_extract\_path
\end{Description}
%
\begin{Usage}
\begin{verbatim}
lav_extract_path(model)
\end{verbatim}
\end{Usage}
%
\begin{Arguments}
\begin{ldescription}
\item[\code{model}] model
\end{ldescription}
\end{Arguments}
%
\begin{Value}
text
\end{Value}
%
\begin{Examples}
\begin{ExampleCode}

## Not run: 
#'
# 모델 문자열
model1a <- "
EXPE ~ H1*IMAG
QUAL ~ H2*EXPE
VAL ~ H3*QUAL + H4*EXPE
SAT ~ H5*IMAG + H6*EXPE + H7*QUAL + H8*VAL
LOY ~ H9*SAT + H10*IMAG
"

model1 <- "
EXPE ~ IMAG
QUAL ~ EXPE
VAL ~ QUAL + EXPE
SAT ~ IMAG + EXPE + QUAL + VAL
LOY ~ SAT + IMAG
"


# 함수 호출
lav_extract_path(model1a)

lav_extract_path(model1)


## End(Not run)


\end{ExampleCode}
\end{Examples}
\inputencoding{utf8}
\HeaderA{lav\_extract\_sm}{lav\_extract\_sm, Extracting a Structural Path Model from a Model}{lav.Rul.extract.Rul.sm}
%
\begin{Description}
lav\_extract\_sm, Extracting a Structural Path Model from a Model
\end{Description}
%
\begin{Usage}
\begin{verbatim}
lav_extract_sm(model, hypo = TRUE, prefix = "a", all = FALSE)
\end{verbatim}
\end{Usage}
%
\begin{Arguments}
\begin{ldescription}
\item[\code{model}] lavaan model

\item[\code{hypo}] hypothesis, default TRUE

\item[\code{prefix}] "a"

\item[\code{all}] all paths
\end{ldescription}
\end{Arguments}
%
\begin{Value}
text
\end{Value}
%
\begin{Examples}
\begin{ExampleCode}

## Not run: 
#'
# Example usage
model1 <- "
IMAG =~ imag1 + imag2 + imag3 + imag4 + imag5
EXPE =~ expe1 + expe2 + expe3 + expe4 + expe5
QUAL =~ qual1 + qual2 + qual3 + qual4 + qual5
VAL =~ val1 + val2 + val3 + val4
SAT =~ sat1 + sat2 + sat3 + sat4
LOY =~ loy1 + loy2 + loy3 + loy4

EXPE ~ IMAG
QUAL ~ EXPE
VAL ~  QUAL+ EXPE
SAT ~ IMAG + EXPE + QUAL + VAL
LOY ~ SAT + IMAG
"


lav_extract_sm(model1)%>%cat("\n")
lav_extract_sm(model1, hypo=FALSE)%>%cat("\n")
lav_extract_sm(model1, prefix="lamda")%>%cat("\n")
lav_extract_sm(model1, prefix="H")%>%cat("\n")
lav_extract_sm(model1, all=FALSE) %>%cat("\n")
lav_extract_sm(model1, all=TRUE) %>%cat("\n")
lav_extract_sm(model1,FALSE) %>%cat("\n")


lav_extract_sm(model1, all=TRUE) %>%
  diagram_model(sizeLat=6)
#'


## End(Not run)


\end{ExampleCode}
\end{Examples}
\inputencoding{utf8}
\HeaderA{lav\_factor\_layout}{자동화된 semPaths모델 layout}{lav.Rul.factor.Rul.layout}
%
\begin{Description}
자동화된 semPaths모델 layout
\end{Description}
%
\begin{Usage}
\begin{verbatim}
lav_factor_layout(
  sem_res,
  positions = NULL,
  nrow = 4,
  ncol = 5,
  new_position = NULL,
  variables = NULL,
  ...
)
\end{verbatim}
\end{Usage}
%
\begin{Arguments}
\begin{ldescription}
\item[\code{sem\_res}] lavaan object

\item[\code{positions}] 각 변수의 위치를 변경 position = list("motiv" =c(1, 3))

\item[\code{nrow}] 4

\item[\code{ncol}] 5

\item[\code{new\_position}] list(c(1, 1), c(1, 2), c(3, 1), c(2, 3),c(2, 4))

\item[\code{variables}] 직업입력하는 변수명

\item[\code{...}] ETC
\end{ldescription}
\end{Arguments}
%
\begin{Value}
MAT
\end{Value}
%
\begin{Examples}
\begin{ExampleCode}

## Not run: 
#'
data(motivation)
str(motivation)
# Test the function with debug prints
motive_model <- "
achieve ~ motiv + ses + iq
motiv ~ ses + health + iq
"

# Assuming 'motivation' is your data frame
motiv_sem <- SEM(motive_model, data = motivation)

motiv_sem%>%summary(stand=TRUE)

# 적용 예시
motiv_sem %>%
  lav_semPaths2(rotation = 2, se=TRUE,
  edge.label.position = 0.4,
                layout =
                lav_factor_layout(motiv_sem,
       new_position = list(c(1, 1), c(1, 2),
       c(3, 1), c(2, 3),c(2, 4)) ),
            rotate_resid = c("motiv" =160, "achieve"=160))


# 적용 예시
motiv_sem %>%
  lav_semPaths2(rotation = 2, se=TRUE,
  edge.label.position = 0.4,
                layout = lav_factor_layout(motiv_sem,
                                           position = list("motiv" =c(1, 3))),
                rotate_resid = c("motiv" =160, "achieve"=160))

#'

## End(Not run)

\end{ExampleCode}
\end{Examples}
\inputencoding{utf8}
\HeaderA{lav\_htmt}{lavaan htmt}{lav.Rul.htmt}
%
\begin{Description}
lavaan htmt
\end{Description}
%
\begin{Usage}
\begin{verbatim}
lav_htmt(
  sem_res = NULL,
  data = NULL,
  blocks = NULL,
  htmt2 = TRUE,
  sig = TRUE,
  cut = 0.9,
  digits = 3,
  imp = ""
)
\end{verbatim}
\end{Usage}
%
\begin{Arguments}
\begin{ldescription}
\item[\code{sem\_res}] lavaan result

\item[\code{data}] sem data

\item[\code{blocks}] sem blocks list data

\item[\code{htmt2}] true using htmt2

\item[\code{sig}] true *

\item[\code{cut}] cut default 0.9, roburst 0.85, general 1

\item[\code{digits}] 3

\item[\code{imp}] "" upper triangle matrix
\end{ldescription}
\end{Arguments}
%
\begin{Value}
data
\end{Value}
%
\begin{Examples}
\begin{ExampleCode}

## Not run: 
#'
data(satisfaction)
model1 <- "
IMAG =~ imag1 + imag2 + imag3 + imag4 + imag5
EXPE =~ expe1 + expe2 + expe3 + expe4 + expe5
QUAL =~ qual1 + qual2 + qual3 + qual4 + qual5
VAL =~ val1 + val2 + val3 + val4
SAT =~ sat1 + sat2 + sat3 + sat4
LOY =~ loy1 + loy2 + loy3 + loy4

EXPE ~ IMAG
QUAL ~ EXPE
VAL ~  QUAL + EXPE
SAT ~ IMAG + EXPE + QUAL + VAL
LOY ~ SAT + IMAG
"

sat_sem0 <- sem(model1, data = satisfaction)

#check
lav_htmt(sat_sem0)
# using data
lav_htmt(data= lav_return_data(sat_sem0), blocks= lav_return_blocks(sat_sem0))
#model data
lav_htmt(data= satisfaction, blocks =  lav_extract_mm(model1)%>% plspm_lav2blocks())
# using data
lav_htmt(data= satisfaction, blocks = model1)

#option
lav_htmt(sat_sem0, htmt2=T, cut=0.85)



## End(Not run)

\end{ExampleCode}
\end{Examples}
\inputencoding{utf8}
\HeaderA{lav\_ind\_hypo}{Function that takes path data as input and outputs an indirect effect hypothesis.}{lav.Rul.ind.Rul.hypo}
%
\begin{Description}
Function that takes path data as input and outputs an indirect effect hypothesis.
\end{Description}
%
\begin{Usage}
\begin{verbatim}
lav_ind_hypo(paths)
\end{verbatim}
\end{Usage}
%
\begin{Arguments}
\begin{ldescription}
\item[\code{paths}] indpath
\end{ldescription}
\end{Arguments}
%
\begin{Value}
text
\end{Value}
%
\begin{Examples}
\begin{ExampleCode}

## Not run: 
#'
# 모델 문자열
model1a <- "
EXPE ~ H1*IMAG
QUAL ~ H2*EXPE
VAL ~ H3*QUAL + H4*EXPE
SAT ~ H5*IMAG + H6*EXPE + H7*QUAL + H8*VAL
LOY ~ H9*SAT + H10*IMAG
"

model1 <- "
EXPE ~ IMAG
QUAL ~ EXPE
VAL ~ QUAL + EXPE
SAT ~ IMAG + EXPE + QUAL + VAL
LOY ~ SAT + IMAG
"
lav_extract_path(model1a)%>%
  lav_generate_ind_paths() %>%
  lav_ind_hypo()


## End(Not run)

\end{ExampleCode}
\end{Examples}
\inputencoding{utf8}
\HeaderA{lav\_interact\_model}{lav\_new\_model: Generate interaction new model}{lav.Rul.interact.Rul.model}
%
\begin{Description}
lav\_new\_model: Generate interaction new model
\end{Description}
%
\begin{Usage}
\begin{verbatim}
lav_interact_model(lav_syn, imlist = NULL)
\end{verbatim}
\end{Usage}
%
\begin{Arguments}
\begin{ldescription}
\item[\code{lav\_syn}] original lavaan syntax model

\item[\code{imlist}] interaction data lav\_extract\_imlist(model)
\end{ldescription}
\end{Arguments}
%
\begin{Value}
model text
\end{Value}
%
\begin{Examples}
\begin{ExampleCode}

## Not run: 
#'

# # Example usage:
data(MarshWenHau)
str(MarshWenHau)

mwh = MarshWenHau

model2 <- "
f1 =~ x1 + x2 + x3
f1Xf2 =~ f1:f2
f2 =~ x4 + x5 + x6
f3 =~ y1 + y2 + y3

f3 ~ f1 + f2
f3 ~ f1Xf2

f1Xf2 ~~ 0*f1
f1Xf2 ~~ 0*f2
"

# interaction model
intsem =  SEM(model2, data= MarshWenHau)
summary(intsem)
intsem %>%lav_semPaths2(color=list(lat="darkred",man="gold"))


##2nd
#base model
lav_extract_mm(model2) %>%cat()
#interaction item list
lav_extract_imlist(model3)
#data
lav_latentProd(mwh, model3)%>%head()

#extracted interaction term
lav_extract_int(model2, lav_extract_imlist(model3))

# or interaction items ; imlist
imlist <- list(
  .f1Xf2 = c("f1", "f2"),
  .i_f1_f4 = c("f1", "f4")
)
# names(imlist)
lav_extract_int(model2, imlist)
#direct
lav_extract_int(model2, imlist= list(.f1Xf2 = c("f1", "f2"),
                                     .i_f1_f4 = c("f1", "f4")))
#estimates model
intsem =  sem(lav_interact_model(model2, lav_extract_imlist(model2)),
              data= lav_latentProd(mwh, model2))





## End(Not run)

\end{ExampleCode}
\end{Examples}
\inputencoding{utf8}
\HeaderA{lav\_latentProd}{Functions to generate data for interaction terms in a dataset}{lav.Rul.latentProd}
%
\begin{Description}
Functions to generate data for interaction terms in a dataset
\end{Description}
%
\begin{Usage}
\begin{verbatim}
lav_latentProd(data, lav_syn, prefix = ".")
\end{verbatim}
\end{Usage}
%
\begin{Arguments}
\begin{ldescription}
\item[\code{data}] original data

\item[\code{lav\_syn}] lavaan syntax

\item[\code{prefix}] . prefix
\end{ldescription}
\end{Arguments}
%
\begin{Value}
data.frame
\end{Value}
%
\begin{Examples}
\begin{ExampleCode}

## Not run: 
#'

# # Example usage:
data(MarshWenHau)
str(MarshWenHau)

mwh = MarshWenHau

model2 <- "
f1 =~ x1 + x2 + x3
f1Xf2 =~ f1:f2
f2 =~ x4 + x5 + x6
f3 =~ y1 + y2 + y3

f3 ~ f1 + f2
f3 ~ f1Xf2

f1Xf2 ~~ 0*f1
f1Xf2 ~~ 0*f2
"

# interaction model
intsem =  SEM(model2, data= MarshWenHau)
summary(intsem)
intsem %>%lav_semPaths2(color=list(lat="darkred",man="gold"))


##2nd
#base model
lav_extract_mm(model2) %>%cat()
#interaction item list
lav_extract_imlist(model3)
#data
lav_latentProd(mwh, model3)%>%head()

#extracted interaction term
lav_extract_int(model2, lav_extract_imlist(model3))

# or interaction items ; imlist
imlist <- list(
  .f1Xf2 = c("f1", "f2"),
  .i_f1_f4 = c("f1", "f4")
)
# names(imlist)
lav_extract_int(model2, imlist)
#direct
lav_extract_int(model2, imlist= list(.f1Xf2 = c("f1", "f2"),
                                     .i_f1_f4 = c("f1", "f4")))
#estimates model
intsem =  SEM(lav_new_model(model2, lav_extract_imlist(model2)),
              data= lav_latentProd(mwh, model2))


## End(Not run)

\end{ExampleCode}
\end{Examples}
\inputencoding{utf8}
\HeaderA{lav\_matrix}{paper table to matrix}{lav.Rul.matrix}
%
\begin{Description}
paper table to matrix
\end{Description}
%
\begin{Usage}
\begin{verbatim}
lav_matrix(text, type = "res")
\end{verbatim}
\end{Usage}
%
\begin{Arguments}
\begin{ldescription}
\item[\code{text}] input copy data
\end{ldescription}
\end{Arguments}
%
\begin{Value}
matrix
\end{Value}
%
\begin{Examples}
\begin{ExampleCode}
## Not run: 

Input data Cho(2014)
text <- "필 요 성 1
         수 업 0.320** 1
         교 사 0.266** 0.647** 1
         계 획 0.506** 0.463** 0.383** 1
         영 향 0.429** 0.548** 0.491** 0.531** 1
         만 족 도 0.401** 0.554** 0.545** 0.491** 0.743** 1"

lav_matrix(text)

#Kim(2014)
kim2014text = "
자기주도성  1
컴퓨터효능감 .284** 1
학습설계  .369** .234** 1
상호작용  .124** .179** .361** 1
강사선호도  .267** .130** .423** .368** 1
학업성취도  .281** .197** .548** .306** .513** 1
동기변화  .246** .108* .453** .248** .405** .445** 1
자신감  .262** .092 .471** .149** .353** .482** .576** 1
흥미  .317** .149** .535** .272** .470** .571** .596** .628** 1"
#
lav_matrix(kim2014text)

# 마이너스가 포함되어 있는 경우
#정성령(2014)
Jungsy2014 = "
수강료 1
강사 .165** 1
수업 .303** .144** 1
개인관리 .494** .069 .284** 1
인지도 .425** .176** .407** .424** 1
이상적영향력 .210** -.056 .212** .286** .222** 1
영감적동기화 .190** -.133** .211** .290** .230** .659** 1
지적자극 .300** -.136** .159** .402** .238** .533** .517** 1
개별적 배려 .202** -.153** .150** .250** .291** .684** .618** .530** 1
"
lav_matrix(Jungsy2014)

Jungsy2014_cor = lav_matrix(Jungsy2014)
Jungsy2014_cor
#'
지경선(2016)
# 개인배경: 수학성적, 기대수준
# 선행학습실태: 학습시간, 학습진도, 반복횟수, 이해정도, 비용, 향후계획
# 선행학습유발요인:입시준비, 학벌사회,학교수업
# 수학학습태도: 흥미와동기, 학습자세, 수학효능감


Jiks2016="
수학성적 1
기대수준 .256*** 1
학습시간 .128** .079 1
학습진도 .304*** .146** .381*** 1
반복횟수 .115* .073 .155*** .173*** 1
이해정도 .150*** .283*** .142** .070 .236*** 1
비용 .132** .214*** .295*** .212*** .042 .086 1
향후계획 .228*** .368*** .134** .183*** .061 .302*** .226*** 1
입시준비 .201*** .300*** .155*** .249*** .121** .128*** .224*** .624*** 1
학벌사회 .097* .196*** .106* .176*** .063 -.032 .152*** .287*** .497*** 1
학교수업 .072 .258*** .120** .140** .087 .242*** .049 .292*** .235*** .098* 1
흥미동기 .357*** .358*** .113* .182*** .062 .350*** .175*** .514*** .359*** .182*** .373*** 1
학습자세 .257*** .441*** .177*** .193*** .101* .320*** .186*** .478*** .448*** .268*** .391*** .636*** 1
수학효능감 .419*** .385*** .249*** .322*** .121** .373*** .182***
           .505*** .423*** .261*** .459*** .703*** .614*** 1
"
lav_matrix(Jiks2016, "check")
lav_matrix(Jiks2016)
#'
## data
textdf =  "motiv 400
harm      77  300
stabi     32   25   100
ppsych   -25  -25   -40    100
ses       25   26    40    -42 100
verbal    59   58    27    -16  18    100
read      53   42    56    -39  43     36  100
arith     60   44    49    -24  37     38   73   100
spell     59   45    48    -31  33     38   87    72   100"

lav_matrix(textdf)
lav_matrix(textdf,type= "all")

lav_matrix(textdf, "check", init=3)



## End(Not run)

\end{ExampleCode}
\end{Examples}
\inputencoding{utf8}
\HeaderA{lav\_modelfit}{lavaan model fit}{lav.Rul.modelfit}
%
\begin{Description}
lavaan model fit
\end{Description}
%
\begin{Usage}
\begin{verbatim}
lav_modelfit(..., nice_table = FALSE)
\end{verbatim}
\end{Usage}
%
\begin{Arguments}
\begin{ldescription}
\item[\code{...}] sem models

\item[\code{nice\_table}] output model fit web
\end{ldescription}
\end{Arguments}
%
\begin{Value}
data
\end{Value}
%
\begin{Examples}
\begin{ExampleCode}

## Not run: 
#'
example(sem)
lav_modelfit(fit)
lav_modelfit(fit, nice_table = T)

lav_modelfit("model1"=fit,"model2"= fit)

#'

## End(Not run)


\end{ExampleCode}
\end{Examples}
\inputencoding{utf8}
\HeaderA{lav\_new\_model}{lav\_new\_model: Generate interaction new model}{lav.Rul.new.Rul.model}
%
\begin{Description}
lav\_new\_model: Generate interaction new model
\end{Description}
%
\begin{Usage}
\begin{verbatim}
lav_new_model(lav_syn, imlist = NULL)
\end{verbatim}
\end{Usage}
%
\begin{Arguments}
\begin{ldescription}
\item[\code{lav\_syn}] original lavaan syntax model

\item[\code{imlist}] interaction data lav\_extract\_imlist(model)
\end{ldescription}
\end{Arguments}
%
\begin{Value}
model text
\end{Value}
%
\begin{Examples}
\begin{ExampleCode}

## Not run: 
#'

# # Example usage:
data(MarshWenHau)
str(MarshWenHau)

mwh = MarshWenHau

model2 <- "
f1 =~ x1 + x2 + x3
f1Xf2 =~ f1:f2
f2 =~ x4 + x5 + x6
f3 =~ y1 + y2 + y3

f3 ~ f1 + f2
f3 ~ f1Xf2

f1Xf2 ~~ 0*f1
f1Xf2 ~~ 0*f2
"

# interaction model
intsem =  SEM(model2, data= MarshWenHau)
summary(intsem)
intsem %>%lav_semPaths2(color=list(lat="darkred",man="gold"))


##2nd
#base model
lav_extract_mm(model2) %>%cat()
#interaction item list
lav_extract_imlist(model3)
#data
lav_latentProd(mwh, model3)%>%head()

#extracted interaction term
lav_extract_int(model2, lav_extract_imlist(model3))

# or interaction items ; imlist
imlist <- list(
  .f1Xf2 = c("f1", "f2"),
  .i_f1_f4 = c("f1", "f4")
)
# names(imlist)
lav_extract_int(model2, imlist)
#direct
lav_extract_int(model2, imlist= list(.f1Xf2 = c("f1", "f2"),
                                     .i_f1_f4 = c("f1", "f4")))
#estimates model
intsem =  SEM(lav_new_model(model2, lav_extract_imlist(model2)),
              data= lav_latentProd(mwh, model2))


## End(Not run)

\end{ExampleCode}
\end{Examples}
\inputencoding{utf8}
\HeaderA{lav\_Pattern}{응답패턴 분석}{lav.Rul.Pattern}
%
\begin{Description}
응답패턴 분석
\end{Description}
%
\begin{Usage}
\begin{verbatim}
lav_Pattern(data, n = 2, dim = 0L, top = 20)
\end{verbatim}
\end{Usage}
%
\begin{Arguments}
\begin{ldescription}
\item[\code{data}] data

\item[\code{n}] 구분되는 개수

\item[\code{dim}] 0L

\item[\code{top}] 20
\end{ldescription}
\end{Arguments}
%
\begin{Value}
data
\end{Value}
%
\begin{Examples}
\begin{ExampleCode}

## Not run: 
#'
latentPattern(MarshWenHau, n=3)

HS9 <- HolzingerSwineford1939[,c("x1","x2","x3","x4","x5",
                                 "x6","x7","x8","x9")]
latentPattern(HS9)

#explain
HSbinary <- as.data.frame( lapply(HS9, cut, 2, labels=FALSE) )
HSbinary %>%head()
lavTables(HSbinary, dim = 0L, categorical = names(HSbinary)) %>%head(20)

#'
#'

x11()
latentPattern(park516_kor[,1:6], n=5) %>%
  ggplot(aes(x= fct_reorder(pattern, desc(obs.freq)), y= obs.freq))+
  geom_bar(stat="identity")+theme_bw()+
  theme(axis.text.x= element_text(size=14, angle=90))

## End(Not run)


\end{ExampleCode}
\end{Examples}
\inputencoding{utf8}
\HeaderA{lav\_remodel}{lav\_remodel, Models that create indirect effects}{lav.Rul.remodel}
%
\begin{Description}
lav\_remodel, Models that create indirect effects
\end{Description}
%
\begin{Usage}
\begin{verbatim}
lav_remodel(
  model,
  prefix = "a",
  auto = TRUE,
  start_node = NULL,
  end_node = NULL,
  cat = FALSE,
  text = FALSE,
  interact = TRUE,
  paths_name = TRUE,
  n_name = 1,
  add_model = NULL
)
\end{verbatim}
\end{Usage}
%
\begin{Arguments}
\begin{ldescription}
\item[\code{model}] lavaan(original model)

\item[\code{prefix}] model hypothesis "a", or "H"

\item[\code{auto}] TRUe indirect effec auto, FALSE all path

\item[\code{start\_node}] NULL If you type it, it starts with

\item[\code{end\_node}] NULL and that's it

\item[\code{cat}] text puput console

\item[\code{text}] Output text to save and type

\item[\code{interact}] TRUe :(use) using SESM, FALSE (:) using sem()

\item[\code{paths\_name}] paths name

\item[\code{n\_name}] paste name number

\item[\code{add\_model}] add model syntax
\end{ldescription}
\end{Arguments}
%
\begin{Value}
text
\end{Value}
%
\begin{Examples}
\begin{ExampleCode}

## Not run: 
#'


data(satisfaction)

# Example usage
model1 <- "
IMAG =~ imag1 + imag2 + imag3 + imag4 + imag5
EXPE =~ expe1 + expe2 + expe3 + expe4 + expe5
QUAL =~ qual1 + qual2 + qual3 + qual4 + qual5
VAL =~ val1 + val2 + val3 + val4
SAT =~ sat1 + sat2 + sat3 + sat4
LOY =~ loy1 + loy2 + loy3 + loy4

EXPE ~ IMAG
QUAL ~ EXPE
VAL ~  QUAL+ EXPE
SAT ~ IMAG + EXPE + QUAL + VAL
LOY ~ SAT + IMAG
"



# lav_remodel(model1)
lav_remodel(model1)%>%cat("\n")
lav_remodel(model1, interact=FALSE)%>%cat("\n")

lav_remodel(model1, auto=FALSE)%>%cat("\n")

lav_remodel(model1, start="EXPE", prefix="H", cat=TRUE)
lav_remodel(model1, cat=TRUE)

lav_remodel(model1, cat=TRUE, text=TRUE)
lav_remodel(model1, cat=TRUE, text=FALSE)

lav_remodel(model1, cat=FALSE, text=TRUE)
lav_remodel(model1, cat=FALSE, text=FALSE)


satis_sem = sem(lav_remodel(model1), satisfaction)
summary(satis_sem)
model1 %>%class()



## End(Not run)


\end{ExampleCode}
\end{Examples}
\inputencoding{utf8}
\HeaderA{lav\_return\_blocks}{lavaan:sem에서 blocks을 생성}{lav.Rul.return.Rul.blocks}
%
\begin{Description}
lavaan:sem에서 blocks을 생성
\end{Description}
%
\begin{Usage}
\begin{verbatim}
lav_return_blocks(sem_res)
\end{verbatim}
\end{Usage}
%
\begin{Arguments}
\begin{ldescription}
\item[\code{sem\_res}] sem result
\end{ldescription}
\end{Arguments}
%
\begin{Value}
list
\end{Value}
%
\begin{Examples}
\begin{ExampleCode}

## Not run: 
#'
data(satisfaction)
model1 <- "
IMAG =~ imag1 + imag2 + imag3 + imag4 + imag5
EXPE =~ expe1 + expe2 + expe3 + expe4 + expe5
QUAL =~ qual1 + qual2 + qual3 + qual4 + qual5
VAL =~ val1 + val2 + val3 + val4
SAT =~ sat1 + sat2 + sat3 + sat4
LOY =~ loy1 + loy2 + loy3 + loy4

EXPE ~ IMAG
QUAL ~ EXPE
VAL ~  QUAL + EXPE
SAT ~ IMAG + EXPE + QUAL + VAL
LOY ~ SAT + IMAG
"

sat_sem0 <- sem(model1, data = satisfaction)

# 함수 결과 확인
lav_return_blocks(sat_sem0)

## End(Not run)

\end{ExampleCode}
\end{Examples}
\inputencoding{utf8}
\HeaderA{lav\_return\_data}{lavaan결과에서 데이터 추출하는 함수}{lav.Rul.return.Rul.data}
%
\begin{Description}
lavaan결과에서 데이터 추출하는 함수
\end{Description}
%
\begin{Usage}
\begin{verbatim}
lav_return_data(sem_res)
\end{verbatim}
\end{Usage}
%
\begin{Arguments}
\begin{ldescription}
\item[\code{sem\_res}] sem result
\end{ldescription}
\end{Arguments}
%
\begin{Value}
data
\end{Value}
%
\begin{Examples}
\begin{ExampleCode}

## Not run: 
#'
data(satisfaction)
model1 <- "
IMAG =~ imag1 + imag2 + imag3 + imag4 + imag5
EXPE =~ expe1 + expe2 + expe3 + expe4 + expe5
QUAL =~ qual1 + qual2 + qual3 + qual4 + qual5
VAL =~ val1 + val2 + val3 + val4
SAT =~ sat1 + sat2 + sat3 + sat4
LOY =~ loy1 + loy2 + loy3 + loy4

EXPE ~ IMAG
QUAL ~ EXPE
VAL ~  QUAL + EXPE
SAT ~ IMAG + EXPE + QUAL + VAL
LOY ~ SAT + IMAG
"

sat_sem0 <- sem(model1, data = satisfaction)

# 함수 결과 확인
lav_return_data(sat_sem0)
lav_return_data(sat_sem0)%>%head()

## End(Not run)


\end{ExampleCode}
\end{Examples}
\inputencoding{utf8}
\HeaderA{lav\_return\_model}{lavaan::sem으로 분석한 함수에서 model추출}{lav.Rul.return.Rul.model}
%
\begin{Description}
lavaan::sem으로 분석한 함수에서 model추출
\end{Description}
%
\begin{Usage}
\begin{verbatim}
lav_return_model(sem_res, sep = "\n")
\end{verbatim}
\end{Usage}
%
\begin{Arguments}
\begin{ldescription}
\item[\code{sem\_res}] sem data

\item[\code{sep}] sep =
\end{ldescription}
\end{Arguments}
%
\begin{Value}
text
\end{Value}
%
\begin{Examples}
\begin{ExampleCode}

## Not run: 
#'
# 예시 사용법
data(satisfaction)
model1 <- "
IMAG =~ imag1 + imag2 + imag3 + imag4 + imag5
EXPE =~ expe1 + expe2 + expe3 + expe4 + expe5
QUAL =~ qual1 + qual2 + qual3 + qual4 + qual5
VAL =~ val1 + val2 + val3 + val4
SAT =~ sat1 + sat2 + sat3 + sat4
LOY =~ loy1 + loy2 + loy3 + loy4

EXPE ~ IMAG
QUAL ~ EXPE
VAL ~  QUAL + EXPE
SAT ~ IMAG + EXPE + QUAL + VAL
LOY ~ SAT + IMAG
"

sat_sem0 <- sem(model1, data = satisfaction)

# 함수 결과 확인
lav_return_model(sat_sem0)
lav_return_model(sat_sem0)%>%cat("\n")
#data
lav_return_data(sat_sem0)

#lav sem reuslt  model to blocks
lav_return_model(sat_sem0)%>%
  lav_extract_mm()%>%
  plspm_lav2blocks()



## End(Not run)


\end{ExampleCode}
\end{Examples}
\inputencoding{utf8}
\HeaderA{lav\_rotate\_resid}{plspm\_rotate\_resid}{lav.Rul.rotate.Rul.resid}
%
\begin{Description}
plspm\_rotate\_resid
\end{Description}
%
\begin{Usage}
\begin{verbatim}
lav_rotate_resid(plotdata, resid_list = NULL, plot = TRUE)
\end{verbatim}
\end{Usage}
%
\begin{Arguments}
\begin{ldescription}
\item[\code{plotdata}] plotdata

\item[\code{resid\_list}] resid list

\item[\code{plot}] TRUE
\end{ldescription}
\end{Arguments}
%
\begin{Value}
plot data
\end{Value}
%
\begin{Examples}
\begin{ExampleCode}

## Not run: 
#'

pathmodel = "
자기효능감 =~ C_S1 +  C_S2 + C_S3 + C_S4 + C_S5
진로태도 =~ B02 + B04 + B10 + B14
진로동기 =~ A_M1 + A_M2 + A_M3
진로준비 =~ D_P1 + D_P2 + D_P3

진로동기 ~ H1*자기효능감
진로태도 ~  H2*자기효능감
진로준비 ~ H5 *진로태도
진로준비 ~  H4 * 진로동기
진로준비 ~  H3*자기효능감
"

plspm_diagram_model(pathmodel, layout="tree2",
                    rotation = 2,residuals = T, residScale = 6,
                    edge.label.position = 0.6)%>%
  plspm_rotate_resid(c(진로동기= 160, 진로태도= 160))

#
plspm_diagram_model(pathmodel, layout="tree2",
                    rotation = 2,residuals = T, residScale = 6,
                    edge.label.position = 0.6)%>%
  plspm_curve(c("진로준비~~ 진로동기" = -2))%>%
  lav_rotate_resid(c(진로동기= 160, 진로태도= 30))

## End(Not run)


\end{ExampleCode}
\end{Examples}
\inputencoding{utf8}
\HeaderA{lav\_semPaths}{Drawing Structural Equation Models Using the lavaan Syntax (park JH PhD)}{lav.Rul.semPaths}
%
\begin{Description}
Drawing Structural Equation Models Using the lavaan Syntax (park JH PhD)
\end{Description}
%
\begin{Usage}
\begin{verbatim}
lav_semPaths(
  lav_obj,
  whatLabels = "est",
  rotation = 2,
  residuals = TRUE,
  residScale = 12,
  layout = "tree2",
  edge.label.cex = 0.7,
  edge.label.position = 0.55,
  edge.color = "black",
  sizeMan = 6,
  sizeMan2 = 4,
  sizeLat = 9,
  sizeLat2 = 6,
  style = "lisrel",
  exoVar = FALSE,
  exoCov = TRUE,
  curve = 1.5,
  asize = 2,
  mar = c(3, 8, 3, 8),
  nDigits = 3,
  shapeLat = "circle",
  shapeMan = "rectangle",
  shapeInt = "triangle",
  border.width = 2,
  edge.width = 1.5,
  Groups = FALSE,
  group_match = "lat",
  structural = FALSE,
  edgeLabels = NULL,
  nodeLabels = NULL,
  fixedStyle = 1,
  intercept = FALSE,
  color = list(lat = "gray10", man = "white"),
  label.cex = 1,
  sig = TRUE,
  rotate_resid = NULL,
  se = FALSE,
  type = "plot"
)
\end{verbatim}
\end{Usage}
%
\begin{Arguments}
\begin{ldescription}
\item[\code{whatLabels}] "model

\item[\code{rotation}] rotation=2

\item[\code{residuals}] residauls TRUE

\item[\code{residScale}] residScale=12

\item[\code{layout}] layout="tree"

\item[\code{edge.label.cex}] 1

\item[\code{edge.label.position}] 0.5

\item[\code{edge.color}] edge.color= black

\item[\code{sizeMan}] 8

\item[\code{sizeMan2}] 4

\item[\code{sizeLat}] 10

\item[\code{sizeLat2}] 6

\item[\code{style}] lisrel

\item[\code{exoVar}] FALSE

\item[\code{exoCov}] TRUE

\item[\code{curve}] 1.5

\item[\code{asize}] 1.5

\item[\code{mar}] c(2,8,3,10)

\item[\code{nDigits}] 2

\item[\code{shapeLat}] "circle, ellipse

\item[\code{shapeMan}] retangle

\item[\code{shapeInt}] triangle ,

\item[\code{border.width}] border.width=1

\item[\code{edge.width}] edge.width=1.5

\item[\code{Groups}] FALSE

\item[\code{group\_match}] "FALSE"lat\#'

\item[\code{structural}] structural = FALSE

\item[\code{edgeLabels}] edgeLabels=NULL

\item[\code{nodeLabels}] nodeLabels=NULL

\item[\code{fixedStyle}] fixedStyle=1 default rbase linetype

\item[\code{color}] list(lat="gray10", man="white")

\item[\code{label.cex}] label.cex

\item[\code{sig}] FALSE

\item[\code{rotate\_resid}] c(lv= angle), c(f1 =90, f2=180)

\item[\code{se}] se=FALSE

\item[\code{type}] type ="plot, and res is lavaan simdata result

\item[\code{plsres\_boot}] plsres\_boot to model syntax
\end{ldescription}
\end{Arguments}
%
\begin{Value}
plot jutpls\_boot
\end{Value}
%
\begin{Examples}
\begin{ExampleCode}
## Not run: 
#'

#data
data(MarshWenHau) # Marsh, Wen & Hau (2000)
# colnames(MarshWenHau) = c( paste0("x",1:6), paste0("y",1:3))
MarshWenHau %>%str()


model1 = "

f1 =~ x1 + x2 + x3
f2 =~ x4 + x5 + x6
f3 =~ y1 + y2 + y3
"

mwh_sem1 = sem(model1, MarshWenHau)
mwh_sem1%>%summary(fit=T)


re_mwh_data = lavPredict(mwh_sem1)%>%data.frame()
re_mwh_data



# 이방법인 LMS방법임

Intf1f2  = re_mwh_data[,"f1"] *  re_mwh_data[,"f2"]

#fitst
MarshWenHau$f1_f2 = Intf1f2
MarshWenHau %>% str()

model2 = "

f1 =~ x1 + x2 + x3
f2 =~ x4 + x5 + x6
f3 =~ y1 + y2 + y3

f1Xf2 =~ f1_f2

f3 ~ f1 + f2 + f1Xf2

f1Xf2 ~~ 0*f1
f1Xf2 ~~ 0*f2

"

mwh_sem2 = sem(model2, MarshWenHau)
mwh_sem2%>%summary(fit=T)
mwh_sem2%>% sem_effect()
x11()
mwh_sem2%>% lav_semPaths2(exoCov = FALSE, sig=F,
                          rotate_resid = c(f3 = 160))

mwh_sem2%>% lav_semPaths2(exoCov = FALSE, sig=TRUE, Groups=TRUE,
                          rotate_resid = c(f3 = 160))

mwh_sem2%>% lav_semPaths2(exoCov = FALSE, sig=TRUE,se=TRUE,
                          rotate_resid = c(f3 = 160))


#'


## End(Not run)

\end{ExampleCode}
\end{Examples}
\inputencoding{utf8}
\HeaderA{lav\_semPaths2}{Drawing Structural Equation Models Using the lavaan Syntax (park JH PhD)}{lav.Rul.semPaths2}
%
\begin{Description}
Drawing Structural Equation Models Using the lavaan Syntax (park JH PhD)
\end{Description}
%
\begin{Usage}
\begin{verbatim}
lav_semPaths2(
  lav_obj,
  whatLabels = "est",
  rotation = 2,
  residuals = TRUE,
  residScale = 12,
  layout = "tree2",
  edge.label.cex = 0.7,
  edge.label.position = 0.55,
  edge.color = "black",
  sizeMan = 6,
  sizeMan2 = 4,
  sizeLat = 9,
  sizeLat2 = 6,
  style = "lisrel",
  exoVar = FALSE,
  exoCov = TRUE,
  curve = 1.5,
  asize = 2,
  mar = c(3, 8, 3, 8),
  nDigits = 3,
  shapeLat = "circle",
  shapeMan = "rectangle",
  shapeInt = "triangle",
  border.width = 2,
  edge.width = 1.5,
  Groups = FALSE,
  group_match = "lat",
  structural = FALSE,
  edgeLabels = NULL,
  nodeLabels = NULL,
  fixedStyle = 1,
  intercept = TRUE,
  color = list(lat = "gray10", man = "white"),
  label.cex = 1,
  sig = TRUE,
  rotate_resid = NULL,
  se = FALSE,
  type = "plot"
)
\end{verbatim}
\end{Usage}
%
\begin{Arguments}
\begin{ldescription}
\item[\code{whatLabels}] "model

\item[\code{rotation}] rotation=2

\item[\code{residuals}] residauls TRUE

\item[\code{residScale}] residScale=12

\item[\code{layout}] layout="tree"

\item[\code{edge.label.cex}] 1

\item[\code{edge.label.position}] 0.5

\item[\code{edge.color}] edge.color= black

\item[\code{sizeMan}] 8

\item[\code{sizeMan2}] 4

\item[\code{sizeLat}] 10

\item[\code{sizeLat2}] 6

\item[\code{style}] lisrel

\item[\code{exoVar}] FALSE

\item[\code{exoCov}] TRUE

\item[\code{curve}] 1.5

\item[\code{asize}] 1.5

\item[\code{mar}] c(2,8,3,10)

\item[\code{nDigits}] 2

\item[\code{shapeLat}] "circle, ellipse

\item[\code{shapeMan}] retangle

\item[\code{shapeInt}] triangle ,

\item[\code{border.width}] border.width=1

\item[\code{edge.width}] edge.width=1.5

\item[\code{Groups}] FALSE

\item[\code{group\_match}] "FALSE"lat\#'

\item[\code{structural}] structural = FALSE

\item[\code{edgeLabels}] edgeLabels=NULL

\item[\code{nodeLabels}] nodeLabels=NULL

\item[\code{fixedStyle}] fixedStyle=1 default rbase linetype

\item[\code{intercept}] intercept

\item[\code{color}] list(lat="gray10", man="white")

\item[\code{label.cex}] label.cex

\item[\code{sig}] FALSE

\item[\code{rotate\_resid}] c(lv= angle), c(f1 =90, f2=180)

\item[\code{se}] se=FALSE

\item[\code{type}] type ="plot, and res is lavaan simdata result

\item[\code{plsres\_boot}] plsres\_boot to model syntax
\end{ldescription}
\end{Arguments}
%
\begin{Value}
plot jutpls\_boot
\end{Value}
%
\begin{Examples}
\begin{ExampleCode}
## Not run: 
#'

#data
data(MarshWenHau) # Marsh, Wen & Hau (2000)
# colnames(MarshWenHau) = c( paste0("x",1:6), paste0("y",1:3))
MarshWenHau %>%str()


model1 = "

f1 =~ x1 + x2 + x3
f2 =~ x4 + x5 + x6
f3 =~ y1 + y2 + y3
"

mwh_sem1 = sem(model1, MarshWenHau)
mwh_sem1%>%summary(fit=T)


re_mwh_data = lavPredict(mwh_sem1)%>%data.frame()
re_mwh_data



# 이방법인 LMS방법임

Intf1f2  = re_mwh_data[,"f1"] *  re_mwh_data[,"f2"]

#fitst
MarshWenHau$f1_f2 = Intf1f2
MarshWenHau %>% str()

model2 = "

f1 =~ x1 + x2 + x3
f2 =~ x4 + x5 + x6
f3 =~ y1 + y2 + y3

f1Xf2 =~ f1_f2

f3 ~ f1 + f2 + f1Xf2

f1Xf2 ~~ 0*f1
f1Xf2 ~~ 0*f2

"

mwh_sem2 = sem(model2, MarshWenHau)
mwh_sem2%>%summary(fit=T)
mwh_sem2%>% sem_effect()
x11()
mwh_sem2%>% lav_semPaths2(exoCov = FALSE, sig=F,
                          rotate_resid = c(f3 = 160))

mwh_sem2%>% lav_semPaths2(exoCov = FALSE, sig=TRUE, Groups=TRUE,
                          rotate_resid = c(f3 = 160))

mwh_sem2%>% lav_semPaths2(exoCov = FALSE, sig=TRUE,se=TRUE,
                          rotate_resid = c(f3 = 160))


#'


## End(Not run)

\end{ExampleCode}
\end{Examples}
\inputencoding{utf8}
\HeaderA{layout\_mat}{layout\_mat}{layout.Rul.mat}
%
\begin{Description}
layout\_mat
\end{Description}
%
\begin{Usage}
\begin{verbatim}
layout_mat(semplot, text = TRUE, byrow = TRUE, variable = TRUE)
\end{verbatim}
\end{Usage}
%
\begin{Arguments}
\begin{ldescription}
\item[\code{semplot}] semPaths2 data

\item[\code{text}] text True output text matrix

\item[\code{byrow}] byrow TRUE

\item[\code{variable}] variable is rownames adapt
\end{ldescription}
\end{Arguments}
%
\begin{Value}
data
\end{Value}
%
\begin{Examples}
\begin{ExampleCode}


## Not run: 
#'
data(offense)
offense %>%str()

nfl_path = matrix(
  c(0, 0, 0, 0, 0,
    0, 0, 0, 0, 0,
    0, 0, 0, 0, 0,
    0, 1, 1, 0, 0,
    1, 0, 0, 1, 0),
  nrow = 5, ncol = 5, byrow = TRUE )

rownames(nfl_path) <- c("Speical", "Rushing", "Passing", "Offense", "Scoring")
colnames(nfl_path) <- c("Speical", "Rushing", "Passing", "Offense", "Scoring")
nfl_path

nfl_blocks = list(7:8, 1:3, 4:6,1:6, 9:11)

plspm_blocks_match(data=offense, blocks=nfl_blocks,
                   paths=nfl_path)

nfl_mode =c("B","A","A","A","A")

nfl_pls = plspm(Data=offense, nfl_path, nfl_blocks, modes = nfl_mode)
nfl_pls%>%summary()

#########
nfl_pls_boot = plspm_sem(Data=offense, nfl_path, nfl_blocks, modes = nfl_mode)

#Layouts extracted from mockups
plspm_semPaths2(nfl_pls_boot,sizeLat = 6)%>%
  plspm_factor_layout()

#Make it input material
plspm_semPaths2(nfl_pls_boot,sizeLat =6)%>%
  plspm_factor_layout() %>%layout_mat(text=T)

#Organizing materials for input
plspm_semPaths2(nfl_pls_boot,sizeLat =6)%>%
  plspm_factor_layout() %>%layout_mat()%>%
  mat_text_arrange()

## End(Not run)

\end{ExampleCode}
\end{Examples}
\inputencoding{utf8}
\HeaderA{lcaGraphData}{lcaGraphData, lcadata to graph data}{lcaGraphData}
%
\begin{Description}
lcaGraphData, lcadata to graph data
\end{Description}
%
\begin{Usage}
\begin{verbatim}
lcaGraphData(lcaResult, ..., yes_col = 3, yes = `Y = 2`)
\end{verbatim}
\end{Usage}
%
\begin{Arguments}
\begin{ldescription}
\item[\code{lcaResult}] lcadata

\item[\code{...}] iteam names, nothng auto

\item[\code{yes\_col}] yes col , default 3

\item[\code{yes}] yes name default `Y = 2`
\end{ldescription}
\end{Arguments}
%
\begin{Value}
data
\end{Value}
%
\begin{Examples}
\begin{ExampleCode}
## Not run: 
#'
#함수를 이용하여 그래프 데이터 생성
jutLca3 %>%lcaGraphData()

jutLca3 %>%lcaGraphData( "학업성적"="L05",
                         "부모지지"="L04",
                         "진로결정"="L03",
                         "준비미흡"="L02",
                         "목표확실"="L01")


## End(Not run)

\end{ExampleCode}
\end{Examples}
\inputencoding{utf8}
\HeaderA{lca\_class\_add}{lca\_class\_add,기존데이터에 잠재클래스 추가}{lca.Rul.class.Rul.add}
%
\begin{Description}
lca\_class\_add,기존데이터에 잠재클래스 추가
\end{Description}
%
\begin{Usage}
\begin{verbatim}
lca_class_add(data, lcadata)
\end{verbatim}
\end{Usage}
%
\begin{Arguments}
\begin{ldescription}
\item[\code{data}] df

\item[\code{lcadata}] lcadata
\end{ldescription}
\end{Arguments}
%
\begin{Value}
data
\end{Value}
%
\begin{Examples}
\begin{ExampleCode}
## Not run: 

jut6 %>% lca_class_add(jutLca3) %>% head()

## End(Not run)
\end{ExampleCode}
\end{Examples}
\inputencoding{utf8}
\HeaderA{lca\_class\_freq}{lca\_class\_freq cal}{lca.Rul.class.Rul.freq}
%
\begin{Description}
lca\_class\_freq cal
\end{Description}
%
\begin{Usage}
\begin{verbatim}
lca_class_freq(lacdata, show = "")
\end{verbatim}
\end{Usage}
%
\begin{Arguments}
\begin{ldescription}
\item[\code{lacdata}] lca result

\item[\code{show}] df, ""
\end{ldescription}
\end{Arguments}
%
\begin{Value}
df
\end{Value}
%
\begin{Examples}
\begin{ExampleCode}
## Not run: 

 jutLca3 %>%lca_class_freq()%>% data.frame()

## End(Not run)
\end{ExampleCode}
\end{Examples}
\inputencoding{utf8}
\HeaderA{lca\_class\_plot}{lca\_class\_plot,잠재클래스 패턴 그래프}{lca.Rul.class.Rul.plot}
%
\begin{Description}
lca\_class\_plot,잠재클래스 패턴 그래프
\end{Description}
%
\begin{Usage}
\begin{verbatim}
lca_class_plot(lcadata_graph, type = "g")
\end{verbatim}
\end{Usage}
%
\begin{Arguments}
\begin{ldescription}
\item[\code{lcadata\_graph}] lcaGraphData 데이터처리된 것

\item[\code{type}] g, full, line
\end{ldescription}
\end{Arguments}
%
\begin{Value}
plot
\end{Value}
%
\begin{Examples}
\begin{ExampleCode}
## Not run: 
#'
## using
jutLca3 %>%lcaGraphData() %>% lca_class_plot()
jutLca3 %>%lcaGraphData() %>% lca_class_plot("line")

jutLca3 %>%lcaGraphData( "학업성적"="L05",
                         "부모지지"="L04",
                         "진로결정"="L03",
                         "준비미흡"="L02",
                         "목표확실"="L01") %>%
  lca_class_plot()
#full graph
jutLca3 %>%lcaGraphData( "학업성적"="L05",
                         "부모지지"="L04",
                         "진로결정"="L03",
                         "준비미흡"="L02",
                         "목표확실"="L01") %>%
  replace_df_rep(
    "L01", "L01:목표확실",
    "L02", "L02:준비미흡",
    "L03", "L03:진로결정",
    "L04", "L04:부모지지",
    "L05", "L05:학업성적")%>%
  replace_df_rep(
    "Class1", "Class1:적극적 진로준비형",
    "Class2", "Class2:불명확 진로인식형",
    "Class3", "Class3:조건부 진로지향형"
  )%>%lca_class_plot("line")


## End(Not run)
\end{ExampleCode}
\end{Examples}
\inputencoding{utf8}
\HeaderA{lca\_freq\_plot}{lca\_freq\_plot LCA 빈도 그래프}{lca.Rul.freq.Rul.plot}
%
\begin{Description}
lca\_freq\_plot LCA 빈도 그래프
\end{Description}
%
\begin{Usage}
\begin{verbatim}
lca_freq_plot(data, lcadata)
\end{verbatim}
\end{Usage}
%
\begin{Arguments}
\begin{ldescription}
\item[\code{data}] df

\item[\code{lcadata}] lcadata
\end{ldescription}
\end{Arguments}
%
\begin{Value}
plot
\end{Value}
%
\begin{Examples}
\begin{ExampleCode}
## Not run: 

jut7c %>% lca_freq_plot(jutLca3)+ ylim(0, 380)

## End(Not run)
\end{ExampleCode}
\end{Examples}
\inputencoding{utf8}
\HeaderA{lca\_gamma}{lca\_gamma 사전확률}{lca.Rul.gamma}
%
\begin{Description}
lca\_gamma 사전확률
\end{Description}
%
\begin{Usage}
\begin{verbatim}
lca_gamma(lcadata)
\end{verbatim}
\end{Usage}
%
\begin{Arguments}
\begin{ldescription}
\item[\code{lcadata}] lcadata
\end{ldescription}
\end{Arguments}
%
\begin{Value}
vedtor
\end{Value}
%
\begin{Examples}
\begin{ExampleCode}
## Not run: 
jutLca3%>% lca_gamma()

## End(Not run)
\end{ExampleCode}
\end{Examples}
\inputencoding{utf8}
\HeaderA{lca\_gof}{lca\_gof}{lca.Rul.gof}
%
\begin{Description}
lca\_gof
\end{Description}
%
\begin{Usage}
\begin{verbatim}
lca_gof(...)
\end{verbatim}
\end{Usage}
%
\begin{Arguments}
\begin{ldescription}
\item[\code{...}] models
\end{ldescription}
\end{Arguments}
%
\begin{Value}
data
\end{Value}
%
\begin{Examples}
\begin{ExampleCode}

## Not run: 
#'
# 함수 사용 예시
lca_gof(jutLca2, jutLca3, jutLca4)

lca_gof(jutLca3)%>%row2col()

## End(Not run)
\end{ExampleCode}
\end{Examples}
\inputencoding{utf8}
\HeaderA{lca\_GraphData}{lcaGraphData, lcadata to graph data}{lca.Rul.GraphData}
%
\begin{Description}
lcaGraphData, lcadata to graph data
\end{Description}
%
\begin{Usage}
\begin{verbatim}
lca_GraphData(lcaResult, ..., yes_col = 3, yes = `Y = 2`)
\end{verbatim}
\end{Usage}
%
\begin{Arguments}
\begin{ldescription}
\item[\code{lcaResult}] lcadata

\item[\code{...}] iteam names, nothng auto

\item[\code{yes\_col}] yes col , default 3

\item[\code{yes}] yes name default `Y = 2`
\end{ldescription}
\end{Arguments}
%
\begin{Value}
data
\end{Value}
%
\begin{Examples}
\begin{ExampleCode}
## Not run: 
#'
#함수를 이용하여 그래프 데이터 생성
jutLca3 %>%lcaGraphData()

jutLca3 %>%lcaGraphData( "학업성적"="L05",
                         "부모지지"="L04",
                         "진로결정"="L03",
                         "준비미흡"="L02",
                         "목표확실"="L01")


## End(Not run)

\end{ExampleCode}
\end{Examples}
\inputencoding{utf8}
\HeaderA{lca\_result}{lca\_result lca분석후 결과 정리}{lca.Rul.result}
%
\begin{Description}
lca\_result lca분석후 결과 정리
\end{Description}
%
\begin{Usage}
\begin{verbatim}
lca_result(lcadata, class_name = NULL)
\end{verbatim}
\end{Usage}
%
\begin{Arguments}
\begin{ldescription}
\item[\code{lcadata}] lcadata

\item[\code{class\_name}] input name
\end{ldescription}
\end{Arguments}
%
\begin{Value}
data
\end{Value}
%
\begin{Examples}
\begin{ExampleCode}
## Not run: 
#'
jutLca3 %>%lca_result() %>%
  replace_df_rep( "L01", "진로목표(L01)",
                  "L02", "준비여부(L02)",
                  "L03", "직업결정(L03)",
                  "L04", "부모추천(L04)",
                  "L05", "학업성적(L05)" ) #변수명을 바꾸기

jutLca3 %>%lca_result(class_name=c(1,2,3))

jutLca3 %>%lca_result(class_name=c("조건부_진로지향형\n(Conditional-Oriented)",
                                   "불명확_진로인식형\n(Unclear-Awareness)",
                                   "적극적_진로준비형\n(Active-Preparedness)")


                      jutLca3 %>%lca_result(class_name=c("조건부_진로지향형",
                                                         "불명확_진로인식형",
                                                         "적극적_진로준비형"))


## End(Not run)
\end{ExampleCode}
\end{Examples}
\inputencoding{utf8}
\HeaderA{lme\_report}{linear mixed model report}{lme.Rul.report}
%
\begin{Description}
Mixed model summary
\end{Description}
%
\begin{Usage}
\begin{verbatim}
lme_report(
  lmedata,
  type = "basic",
  form = "lmer",
  apa = FALSE,
  fit_more = FALSE,
  ranef_sig = FALSE,
  show.effect = FALSE,
  show.ci = FALSE
)
\end{verbatim}
\end{Usage}
%
\begin{Arguments}
\begin{ldescription}
\item[\code{lmedata}] lmedata is lmer() function result

\item[\code{type}] = 'all' is res, 'Fixed\_effect','Random\_effect','ICC', 'ConfidenceInterval\_95','Satterthwaite\_method','FIT','APA'

\item[\code{form}] 'lmer','glmer' = 'logit'='dich', 'poisson' ='pois'

\item[\code{apa}] = default FALSE, if you  REML is FALSE

\item[\code{fit\_more}] = default FALSE, TRUE detailed report

\item[\code{ranef\_sig}] ranef\_sig =TRUE  random effect test

\item[\code{show.effect}] if you show.effect = TRUE -> show ranef, fixef, coef

\item[\code{show.ci}] TRUE calculation 95percent CI
\end{ldescription}
\end{Arguments}
%
\begin{Examples}
\begin{ExampleCode}
## Not run: 
data(sleepstudy)

sleepstudy |> str()

lmer(Reaction ~ Days + (Days | Subject), sleepstudy) |> lme_report()

lmer(Reaction ~ Days + (Days | Subject), sleepstudy) |> lme_report(apa=TRUE)


book.url <- "https://stat.ethz.ch/~meier/teaching/book-anova"
quality <- readRDS(url(file.path(book.url, "data/quality.rds")))
str(quality)
fit.quality <- lmer(score ~ (1 | employee) + (1 | batch) +
                     (1 | employee:batch), data = quality)
summary(fit.quality)
fit.quality|> lme_report()
fit.quality|> lme_report(apa= TRUE) #error the app is FALSE
fit.quality|> lme_report(apa=FALSE) #error the app is FALSE
fit.quality|> ICC()

lmer(Reaction ~ Days + (Days | Subject), sleepstudy) |>
SEM212::lme_report(apa=TRUE, fit_more = T)
lmer(Reaction ~ Days + (Days | Subject), sleepstudy) |>
SEM212::lme_report(apa=TRUE, fit_more = F)  #default FALSE, it is show that AIC and BIC


## End(Not run)

\end{ExampleCode}
\end{Examples}
\inputencoding{utf8}
\HeaderA{long\_df}{longdata transfomation}{long.Rul.df}
%
\begin{Description}
longdata transfomation
\end{Description}
%
\begin{Usage}
\begin{verbatim}
long_df(
  data,
  rowname = "rownames",
  names_to = "names",
  values_to = "values",
  cols = NULL,
  fix = NULL,
  rownames_to_column = TRUE
)
\end{verbatim}
\end{Usage}
%
\begin{Arguments}
\begin{ldescription}
\item[\code{data}] widedata

\item[\code{rowname}] defailt accent

\item[\code{names\_to}] name, speaker(kge)

\item[\code{values\_to}] value, freq

\item[\code{cols}] colrange

\item[\code{fix}] column fixed (thereofore not using cols )
\end{ldescription}
\end{Arguments}
%
\begin{Value}
longdata
\end{Value}
%
\begin{Examples}
\begin{ExampleCode}
## Not run: 

mtcars[,1:4] %>%long_df("car", "value")
mtcars[1:3, 1:7] %>% long_df(fix= 1:3)
mtcars[1:3, 1:7] %>% long_df(fix= 2:7)
mtcars[1:3, 1:7] %>% long_df(cols= 3:7, rownames_to_column=F)
mtcars[1:3, 1:7] %>% long_df(cols= 3:7)
mtcars[1:3, 1:7] %>% long_df(cols= 5:7)


 mtcars[,1:6] %>%long_df("car", "value", rownames_to_column = F, cols = c(1,3:6))




## End(Not run)
\end{ExampleCode}
\end{Examples}
\inputencoding{utf8}
\HeaderA{lowerdf}{df or mat lower tri}{lowerdf}
%
\begin{Description}
df or mat lower tri
\end{Description}
%
\begin{Usage}
\begin{verbatim}
lowerdf(data)
\end{verbatim}
\end{Usage}
%
\begin{Arguments}
\begin{ldescription}
\item[\code{data}] data
\end{ldescription}
\end{Arguments}
%
\begin{Value}
matrix
\end{Value}
\inputencoding{utf8}
\HeaderA{lowerMat}{lower pull matrix}{lowerMat}
%
\begin{Description}
lower pull matrix
\end{Description}
%
\begin{Usage}
\begin{verbatim}
lowerMat(mat, fill = 0, diag = 0, digits = 3)
\end{verbatim}
\end{Usage}
%
\begin{Arguments}
\begin{ldescription}
\item[\code{mat}] matrix data

\item[\code{fill}] matrix data fill upper 0 or if that's what you want , enter it

\item[\code{diag}] matrix data diag enter 0
\end{ldescription}
\end{Arguments}
%
\begin{Value}
lowermatrix, upper 0
\end{Value}
\inputencoding{utf8}
\HeaderA{lowerMat\_na}{lower pull matrix fill NA}{lowerMat.Rul.na}
%
\begin{Description}
lower pull matrix fill NA
\end{Description}
%
\begin{Usage}
\begin{verbatim}
lowerMat_na(
  mat,
  fill = "",
  diag = 1,
  numeric = FALSE,
  digits = 3,
  justify = "right"
)
\end{verbatim}
\end{Usage}
%
\begin{Arguments}
\begin{ldescription}
\item[\code{mat}] matrix data

\item[\code{fill}] matrix data NA or .

\item[\code{diag}] matrix  diagonal data NA or .
\end{ldescription}
\end{Arguments}
%
\begin{Value}
lowermatrix, upper 0
\end{Value}
\inputencoding{utf8}
\HeaderA{lpa}{Latent Profile Analysis (LPA) Function}{lpa}
%
\begin{Description}
This function performs latent profile analysis (LPA) using the `mclust` package, creates visualizations,
and returns summary statistics and plots.
\end{Description}
%
\begin{Usage}
\begin{verbatim}
lpa(
  df,
  profile = 4,
  model = c("EII", "VII", "EEI", "VEI", "EVI", "VVI", "EEE", "EVE", "VEE", "VVE", "EEV",
    "VEV", "EVV", "VVV"),
  seed = 123,
  size_axis = 14,
  size_strip = 14,
  name_strip = NULL,
  md = FALSE
)
\end{verbatim}
\end{Usage}
%
\begin{Arguments}
\begin{ldescription}
\item[\code{df}] Data frame containing the numeric variables for LPA.

\item[\code{profile}] Integer specifying the number of profiles (clusters). Default is 4.

\item[\code{model}] Character vector specifying the covariance model to use in the analysis.
Options include "EII", "VII", "EEI", "VEI", "EVI", "VVI", "EEE", "EVE", "VEE", "VVE",
"EEV", "VEV", "EVV", "VVV". Default is "EEE".

\item[\code{seed}] Integer for setting the random seed for reproducibility. Default is 123.

\item[\code{size\_axis}] Numeric value specifying the size of the axis text in plots. Default is 14.

\item[\code{size\_strip}] Numeric value specifying the size of the strip text in plots. Default is 14.

\item[\code{name\_strip}] Optional vector of names for customizing profile strip labels.

\item[\code{md}] Logical indicating whether to convert the combined data to markdown format. Default is FALSE.
\end{ldescription}
\end{Arguments}
%
\begin{Value}
A list containing:
\begin{ldescription}
\item[\code{summary\_table}] Data frame with raw means, frequencies, and proportions for each profile.
\item[\code{scaled\_summary\_table}] Data frame with standardized means, frequencies, and proportions for each profile.
\item[\code{profile\_frequencies}] Data frame with profile frequencies and proportions.
\item[\code{combined\_data}] Combined data in either raw or markdown format.
\item[\code{mclust\_object}] Summary of the Mclust object from the analysis.
\item[\code{raw\_plot}] Plot of raw scores by profile.
\item[\code{scaled\_plot}] Plot of standardized scores by profile.
\item[\code{combined\_plot}] Combined plot of raw and standardized scores.
\end{ldescription}
\end{Value}
%
\begin{Examples}
\begin{ExampleCode}
## Not run: 
data <- mtcars[, c("mpg", "disp", "hp", "drat", "wt")]
result <- lpa(data, profile = 3, model = "EEE", seed = 123)
result$raw_plot
result$scaled_plot
result$combined_data

## End(Not run)
\end{ExampleCode}
\end{Examples}
\inputencoding{utf8}
\HeaderA{lpa\_BIC\_plot}{lpa\_BIC\_plot: Visualize BIC Values for Latent Profile Analysis}{lpa.Rul.BIC.Rul.plot}
%
\begin{Description}
This function generates a plot to visualize BIC values from latent profile analysis.
\end{Description}
%
\begin{Usage}
\begin{verbatim}
lpa_BIC_plot(
  data,
  n_profiles_range = 1:6,
  type = "all",
  legend_position = "top",
  covar_model = c("EII", "EEE", "VII", "EEV", "EEI", "VVI", "VVV"),
  set_n_profile = NULL,
  xintercept = 4,
  size.text = 4,
  flip = TRUE,
  na2zero = FALSE,
  new.window = TRUE,
  basic = TRUE
)
\end{verbatim}
\end{Usage}
%
\begin{Arguments}
\begin{ldescription}
\item[\code{data}] A data frame containing BIC values for different models and profiles.

\item[\code{n\_profiles\_range}] A numeric vector specifying the range of profiles to analyze.

\item[\code{type}] Type of output to return. Options include "all", "res", "data", "long\_data", "wide\_data", or "graph".

\item[\code{legend\_position}] The position of the legend in the plot (e.g., "top", "bottom").

\item[\code{covar\_model}] A character vector specifying covariance models to include in the analysis.

\item[\code{set\_n\_profile}] A fixed number of profiles to use in the analysis (optional).

\item[\code{xintercept}] The x-axis intercept line for visualization.

\item[\code{size.text}] The size of text labels in the plot.

\item[\code{flip}] Logical. If TRUE, flips the y-axis values for visualization. Default is TRUE.

\item[\code{na2zero}] Logical. If TRUE, replaces NA values with zero. Default is FALSE.

\item[\code{new.window}] Logical. If TRUE, opens a new graphics window for the plot. Default is TRUE.

\item[\code{basic}] Logical. If TRUE, uses only basic covariance models (e.g., "EII", "EEE"). Default is TRUE.
\end{ldescription}
\end{Arguments}
%
\begin{Value}
A list containing:
\begin{ldescription}
\item[\code{data}] The wide-format data used for plotting.
\item[\code{graph}] The ggplot object for visualization.
\item[\code{to\_plot}] The long-format data used for plotting.
\end{ldescription}
\end{Value}
%
\begin{Examples}
\begin{ExampleCode}
## Not run: 
# Example usage of lpa_BIC_plot:
lpa_BIC_plot(data, n_profiles_range = 1:6, basic = FALSE, set_n_profile = 4)

## End(Not run)
\end{ExampleCode}
\end{Examples}
\inputencoding{utf8}
\HeaderA{lpa\_bind\_classdata}{lpa\_bind\_classdata}{lpa.Rul.bind.Rul.classdata}
%
\begin{Description}
lpa\_bind\_classdata
\end{Description}
%
\begin{Usage}
\begin{verbatim}
lpa_bind_classdata(
  data,
  n_profiles = 3,
  models = NULL,
  package = c("mclust"),
  type = "all",
  add_line = TRUE,
  var_model = "EII",
  Class_colors = 1:4,
  linewidth = 0.9,
  sd = TRUE,
  set_seed = 123
)
\end{verbatim}
\end{Usage}
%
\begin{Arguments}
\begin{ldescription}
\item[\code{data}] data

\item[\code{n\_profiles}] default 3

\item[\code{package}] mclucs , MplusAutomation

\item[\code{type}] all, res, fit, est, data, all, plot, group\_plot

\item[\code{add\_line}] add\_lineTRUE

\item[\code{var\_model}] "EEE","EEI,"VVV","EII"

\item[\code{Class\_colors}] 1:4

\item[\code{linewidth}] linewidth 0.9

\item[\code{sd}] TRUE graph
\end{ldescription}
\end{Arguments}
%
\begin{Value}
data
\end{Value}
%
\begin{Examples}
\begin{ExampleCode}

## Not run: 

#'
lpa_bind_classdata(select(iris, -Species))

lpa_bind_classdata(select(iris, -Species), type="res")
lpa_bind_classdata(select(iris, -Species), type="fit")
lpa_bind_classdata(select(iris, -Species), type="est")
lpa_bind_classdata(select(iris, -Species), type="data")%>%Round(3)%>%dall
lpa_bind_classdata(select(iris, -Species), type="all")
#'
# "EEE": variances = "equal", covariances = "equal"
# 프로파일 간 분산과 공분산이 모두 동일하다고 가정.
# "EEV": variances = "equal", covariances = "varying"
# 분산은 동일하지만, 공분산은 프로파일마다 다르게 추정.
# "EEI": variances = "equal", covariances = "zero"
# 분산은 동일하지만, 공분산은 추정하지 않음 (공변량 간 독립성 가정).
# "EII": variances = "equal", covariances = "zero"
# 동일한 가정으로 분산이 동일하고, 공분산은 추정되지 않음.

## End(Not run)



\end{ExampleCode}
\end{Examples}
\inputencoding{utf8}
\HeaderA{lpa\_bind\_classdata\_tidyLPA}{lpa\_bind\_classdata\_tidyLPA}{lpa.Rul.bind.Rul.classdata.Rul.tidyLPA}
%
\begin{Description}
lpa\_bind\_classdata\_tidyLPA
\end{Description}
%
\begin{Usage}
\begin{verbatim}
lpa_bind_classdata_tidyLPA(
  data,
  n_profiles = 3,
  models = NULL,
  package = c("mclust", "MplusAutomation"),
  type = "all",
  add_line = TRUE,
  var_model = "EEE",
  Class_colors = 1:4,
  linewidth = 0.9,
  sd = TRUE
)
\end{verbatim}
\end{Usage}
%
\begin{Arguments}
\begin{ldescription}
\item[\code{data}] data

\item[\code{n\_profiles}] default 3

\item[\code{package}] mclucs , MplusAutomation

\item[\code{type}] all, res, fit, est, data, all, plot, group\_plot

\item[\code{add\_line}] add\_lineTRUE

\item[\code{var\_model}] "EEE","EEI,"VVV","EII"

\item[\code{Class\_colors}] 1:4

\item[\code{linewidth}] linewidth 0.9

\item[\code{sd}] TRUE graph
\end{ldescription}
\end{Arguments}
%
\begin{Value}
data
\end{Value}
%
\begin{Examples}
\begin{ExampleCode}

## Not run: 

#'
lpa_bind_classdata(select(iris, -Species))

lpa_bind_classdata(select(iris, -Species), type="res")
lpa_bind_classdata(select(iris, -Species), type="fit")
lpa_bind_classdata(select(iris, -Species), type="est")
lpa_bind_classdata(select(iris, -Species), type="data")%>%Round(3)%>%dall
lpa_bind_classdata(select(iris, -Species), type="all")
#'
# "EEE": variances = "equal", covariances = "equal"
# 프로파일 간 분산과 공분산이 모두 동일하다고 가정.
# "EEV": variances = "equal", covariances = "varying"
# 분산은 동일하지만, 공분산은 프로파일마다 다르게 추정.
# "EEI": variances = "equal", covariances = "zero"
# 분산은 동일하지만, 공분산은 추정하지 않음 (공변량 간 독립성 가정).
# "EII": variances = "equal", covariances = "zero"
# 동일한 가정으로 분산이 동일하고, 공분산은 추정되지 않음.

## End(Not run)



\end{ExampleCode}
\end{Examples}
\inputencoding{utf8}
\HeaderA{lpa\_create\_profiles}{lpa\_create\_profiles}{lpa.Rul.create.Rul.profiles}
%
\begin{Description}
lpa\_create\_profiles
\end{Description}
%
\begin{Usage}
\begin{verbatim}
lpa_create_profiles(
  df,
  n_profiles = 3,
  variance_structure = "freed",
  covariance_structure = "freed"
)
\end{verbatim}
\end{Usage}
%
\begin{Arguments}
\begin{ldescription}
\item[\code{df}] df

\item[\code{n\_profiles}] default = 3

\item[\code{variance\_structure}] freed

\item[\code{covariance\_structure}] freed
\end{ldescription}
\end{Arguments}
%
\begin{Value}
data
\end{Value}
%
\begin{Examples}
\begin{ExampleCode}

## Not run: 
#'
iris[,-5]%>% lpa_create_profiles(n_profiles =  4,
                                    variance_structure = "constrained",
                                    covariance_structure = "constrained")

## End(Not run)


\end{ExampleCode}
\end{Examples}
\inputencoding{utf8}
\HeaderA{lpa\_create\_profiles2}{lpa\_create\_profiles2}{lpa.Rul.create.Rul.profiles2}
%
\begin{Description}
lpa\_create\_profiles2
\end{Description}
%
\begin{Usage}
\begin{verbatim}
lpa_create_profiles2(
  df,
  n_profiles = 4,
  model_name = NULL,
  type = "mean",
  seed = 123
)
\end{verbatim}
\end{Usage}
%
\begin{Arguments}
\begin{ldescription}
\item[\code{df}] data

\item[\code{n\_profiles}] 3, ...

\item[\code{model\_name}] EEI, VVI, EEE, VVV
\end{ldescription}
\end{Arguments}
%
\begin{Value}
data
\end{Value}
%
\begin{Examples}
\begin{ExampleCode}

## Not run: 
#'

iris[,-5] %>% lpa_create_profiles2(n_profiles = 4)

#v:constrained, Co: fixed
iris[,-5] %>% lpa_create_profiles2(n_profiles = 4, model_name="EEE")
#variance_structure == "fixed"
iris[,-5] %>% lpa_create_profiles2(n_profiles = 4, model_name="VVV")


# variance_structure == "freed" &covariance_structure == "fixed"
iris[,-5] %>% lpa_create_profiles2(n_profiles = 4, model_name="VVI")


#variance_structure == "constrained" & covariance_structure == "fixed"
iris[,-5] %>% lpa_create_profiles2(n_profiles = 4, model_name="EEI")
#'

## End(Not run)

\end{ExampleCode}
\end{Examples}
\inputencoding{utf8}
\HeaderA{lpa\_explore\_modelfit}{lpa\_explore\_modelfit}{lpa.Rul.explore.Rul.modelfit}
%
\begin{Description}
lpa\_explore\_modelfit
\end{Description}
%
\begin{Usage}
\begin{verbatim}
lpa_explore_modelfit(
  df,
  n_profiles_range = 1:9,
  modelNames = c("EEV", "EEI", "EEE", "EII", "VVI", "VII", "VVV"),
  chk = FALSE,
  basic = FALSE,
  set_n_profile = NULL,
  na2zero = TRUE
)
\end{verbatim}
\end{Usage}
%
\begin{Arguments}
\begin{ldescription}
\item[\code{df}] data

\item[\code{n\_profiles\_range}] 1:9

\item[\code{modelNames}] "EEV", "EEI", "EEE", "EII", "VVI", "VII", "VVV"

\item[\code{chk}] default FALSE, TRUE show

\item[\code{basic}] select EII, EEE, VII, EEI

\item[\code{set\_n\_profile}] default select row number then result show

\item[\code{na2zero}] na2zero is NA data to 0
\end{ldescription}
\end{Arguments}
%
\begin{Value}
data
\end{Value}
%
\begin{Examples}
\begin{ExampleCode}

## Not run: 

df <- select(iris, -Species)
lpa_explore_modelfit(df, 1:4)
#' lpa_explore_modelfit(df, 1:9)
lpa_explore_modelfit(pwjlpa,1:4, basic=TRUE)
lpa_explore_modelfit(pwjlpa,1:4, basic=TRUE, set_n_profile=4)
lpa_explore_modelfit(pwjlpa,1:6, basic=TRUE)

lpa_explore_modelfit(pwjlpa,1:6, basic=FALSE)
lpa_explore_modelfit(pwjlpa,1:10, basic=TRUE, set_n_profile=4)
lpa_explore_modelfit(pwjlpa,1:10, basic=TRUE, set_n_profile=4, chk=TRUE)


## End(Not run)

\end{ExampleCode}
\end{Examples}
\inputencoding{utf8}
\HeaderA{lpa\_gmm}{lpa\_gmm: Latent Profile Analysis using Gaussian Mixture Models}{lpa.Rul.gmm}
%
\begin{Description}
This function performs latent profile analysis (LPA) using Gaussian Mixture Models (GMM).
It provides model parameters, profile classifications, and visualizations for the analysis.
\end{Description}
%
\begin{Usage}
\begin{verbatim}
lpa_gmm(
  df,
  profiles = 4,
  modelNames = "EII",
  g_type = c("uncertainty", "classification")
)
\end{verbatim}
\end{Usage}
%
\begin{Arguments}
\begin{ldescription}
\item[\code{df}] A data frame containing the variables to be analyzed.

\item[\code{profiles}] An integer specifying the number of profiles to extract. Default is 4.

\item[\code{modelNames}] A character string specifying the GMM covariance model. Default is "EII".

\item[\code{g\_type}] A character string specifying the type of GMM plot to generate.
Choices are "uncertainty" or "classification". Default is "uncertainty".
\end{ldescription}
\end{Arguments}
%
\begin{Value}
A list containing:
\begin{description}

\item[result] The GMM parameters for the fitted model.
\item[lpa\_reuslt] A summary table of profile means and frequencies.
\item[graph] The GMM plot based on the specified type.
\item[profile\_plot] A bar plot visualizing profile frequencies.
\item[table] A formatted markdown table summarizing profile results.

\end{description}

\end{Value}
%
\begin{Examples}
\begin{ExampleCode}
## Not run: 
# Example usage of lpa_gmm
library(mclust)
library(dplyr)
library(ggplot2)

# Example data
set.seed(123)
df <- data.frame(x = rnorm(100), y = rnorm(100))

# Perform LPA with 3 profiles
lpa_results <- lpa_gmm(df, profiles = 3)

# View results
print(lpa_results$lpa_reuslt)
plot(lpa_results$graph)

## End(Not run)
\end{ExampleCode}
\end{Examples}
\inputencoding{utf8}
\HeaderA{lpa\_plot}{Latent Profile Analysis Profile Plot}{lpa.Rul.plot}
%
\begin{Description}
Latent Profile Analysis Profile Plot
\end{Description}
%
\begin{Usage}
\begin{verbatim}
lpa_plot(
  data,
  n_profiles = 3,
  model_name = "EEE",
  view = "pair",
  alpha = 0.6,
  gtype = "normal",
  angle = 90,
  size.x = 13,
  size.p = 5,
  ncol = 2,
  fct_reorder = NULL,
  show.legend = FALSE,
  legend.text.size = 16,
  legend_position = "top",
  linewidth = 1,
  title_size = 18,
  seed = 123
)
\end{verbatim}
\end{Usage}
%
\begin{Arguments}
\begin{ldescription}
\item[\code{data}] A data frame containing numeric variables for latent profile analysis.

\item[\code{n\_profiles}] Integer. Number of latent profiles to identify. Default is 3.

\item[\code{model\_name}] Character. Model name to use in Mclust for Gaussian mixture modeling (e.g., "EEE", "EEI", "VVI", "VVV"). Default is "EEE".

\item[\code{view}] Character. Determines the view type of the plots ("pair", "each"). Default is "pair".

\item[\code{alpha}] Numeric. Transparency level for bar plots. Default is 0.6.

\item[\code{gtype}] Character. Graph type to generate ("normal", "trans", "each"). Default is "normal".

\item[\code{angle}] Numeric. Angle for the x-axis text labels. Default is 90.

\item[\code{size.x}] Numeric. Font size for the x-axis text labels. Default is 13.

\item[\code{size.p}] Numeric. Size of points in the plot. Default is 5.

\item[\code{ncol}] Integer. Number of columns for the facet plots in "each" view. Default is 2.

\item[\code{fct\_reorder}] Character vector or NULL. Specifies the order of factor levels for plotting. Default is NULL.

\item[\code{show.legend}] Logical. Whether to show legends for line plots. Default is FALSE.

\item[\code{legend.text.size}] Numeric. Font size for legend text. Default is 16.

\item[\code{legend\_position}] Character. Position of the legend ("top", "bottom", "left", "right"). Default is "top".

\item[\code{linewidth}] Numeric. Line width for plots. Default is 1.

\item[\code{title\_size}] Numeric. Font size for plot titles. Default is 18.

\item[\code{seed}] Integer or NULL. Seed value for reproducibility. Default is 123.
\end{ldescription}
\end{Arguments}
%
\begin{Value}
A list containing:
\begin{ldescription}
\item[\code{std}] Standardized mean data for each profile.
\item[\code{est}] Estimated raw mean data for each profile.
\item[\code{graph}] Combined plots for the profiles.
\item[\code{plot\_data}] Data used for plotting.
\item[\code{Class\_freq}] Frequency table of profile classifications.
\end{ldescription}
\end{Value}
%
\begin{Examples}
\begin{ExampleCode}
## Not run: 
# Example usage with iris dataset
result <- lpa_plot(iris[, -5], n_profiles = 3, model_name = "EEE")
print(result$graph)

## End(Not run)
\end{ExampleCode}
\end{Examples}
\inputencoding{utf8}
\HeaderA{lpa\_profile\_plot}{Latent Profile Analysis (LPA) Profile Plot Function}{lpa.Rul.profile.Rul.plot}
%
\begin{Description}
This function generates plots to visualize latent profile analysis results, including raw scores and standardized scores for profiles.
\end{Description}
%
\begin{Usage}
\begin{verbatim}
lpa_profile_plot(
  data,
  n_profiles = 3,
  model_name = "EEE",
  view = "pair",
  alpha = 0.6,
  gtype = "normal",
  angle = 90,
  size.x = 13,
  size.p = 5,
  ncol = 2,
  fct_reorder = NULL,
  show.legend = FALSE,
  legend.text.size = 16,
  legend_position = "top",
  linewidth = 1,
  title_size = 18,
  seed = 123
)
\end{verbatim}
\end{Usage}
%
\begin{Arguments}
\begin{ldescription}
\item[\code{data}] A data frame containing numeric variables for latent profile analysis.

\item[\code{n\_profiles}] Integer. Number of latent profiles to identify. Default is 3.

\item[\code{model\_name}] Character. Model name for Gaussian mixture modeling in Mclust (e.g., "EEE", "EEI", "VVI", "VVV"). Default is "EEE".

\item[\code{view}] Character. Determines the view type of the plots ("pair", "each"). Default is "pair".

\item[\code{alpha}] Numeric. Transparency level for bar plots. Default is 0.6.

\item[\code{gtype}] Character. Type of graph to generate ("normal", "trans"). Default is "normal".

\item[\code{angle}] Numeric. Angle for the x-axis text labels. Default is 90.

\item[\code{size.x}] Numeric. Font size for the x-axis text labels. Default is 13.

\item[\code{size.p}] Numeric. Size of points in the plot. Default is 5.

\item[\code{ncol}] Integer. Number of columns for grid-arranged plots. Default is 2.

\item[\code{fct\_reorder}] Character vector or NULL. Specifies the order of factor levels for plotting. Default is NULL.

\item[\code{show.legend}] Logical. Whether to show legends for line plots. Default is FALSE.

\item[\code{legend.text.size}] Numeric. Font size for legend text. Default is 16.

\item[\code{legend\_position}] Character. Position of the legend ("top", "bottom", "left", "right"). Default is "top".

\item[\code{linewidth}] Numeric. Line width for line plots. Default is 1.

\item[\code{title\_size}] Numeric. Font size for plot titles. Default is 18.

\item[\code{seed}] Integer or NULL. Seed value for reproducibility. Default is 123.
\end{ldescription}
\end{Arguments}
%
\begin{Value}
A list containing:
\begin{ldescription}
\item[\code{std}] Standardized mean data for each profile.
\item[\code{est}] Estimated raw mean data for each profile.
\item[\code{graph}] Combined plots for the profiles.
\item[\code{plot\_data}] Data used for plotting.
\item[\code{Class\_freq}] Frequency table of profile classifications.
\end{ldescription}
\end{Value}
%
\begin{Examples}
\begin{ExampleCode}
## Not run: 
# Example usage with the iris dataset
result <- lpa_profile_plot(iris[, -5], n_profiles = 3, model_name = "EEE", gtype = "normal")
print(result$graph)

## End(Not run)
\end{ExampleCode}
\end{Examples}
\inputencoding{utf8}
\HeaderA{lpa\_profile\_plot2}{profile\_plot2 Brargrahp}{lpa.Rul.profile.Rul.plot2}
%
\begin{Description}
profile\_plot2 Brargrahp
\end{Description}
%
\begin{Usage}
\begin{verbatim}
lpa_profile_plot2(data, n_profiles = 3, model_name = NULL, view = "pair")
\end{verbatim}
\end{Usage}
%
\begin{Arguments}
\begin{ldescription}
\item[\code{data}] data

\item[\code{n\_profiles}] profiles number

\item[\code{model\_name}] "EEE","EEI", "VVV","VVI",

\item[\code{view}] each, pair
\end{ldescription}
\end{Arguments}
%
\begin{Value}
plot
\end{Value}
%
\begin{Examples}
\begin{ExampleCode}

## Not run: 

iris[,-5] %>%profile_plot2(n_profiles = 4, model_name = "EEE")


## End(Not run)


\end{ExampleCode}
\end{Examples}
\inputencoding{utf8}
\HeaderA{lpa\_profile\_plot\_summary}{lpa\_profile\_plot summary data}{lpa.Rul.profile.Rul.plot.Rul.summary}
%
\begin{Description}
lpa\_profile\_plot summary data
\end{Description}
%
\begin{Usage}
\begin{verbatim}
lpa_profile_plot_summary(resl, digits = 3)
\end{verbatim}
\end{Usage}
%
\begin{Arguments}
\begin{ldescription}
\item[\code{resl}] result

\item[\code{digits}] default 3
\end{ldescription}
\end{Arguments}
%
\begin{Value}
df
\end{Value}
%
\begin{Examples}
\begin{ExampleCode}
## Not run: 
#'
lpa_profile_plot(edlcarename, n_profiles = 4, model_name = "EII", gtype="trans",
                 fct_reorder= f_order,
                 size.x=16, alpha = 0.2)%>%
  lpa_profile_plot_summary() %>%
  nice_table()
#'

## End(Not run)

\end{ExampleCode}
\end{Examples}
\inputencoding{utf8}
\HeaderA{LRT\_apa}{LRT\_apa}{LRT.Rul.apa}
%
\begin{Description}
LRT\_apa
\end{Description}
%
\begin{Usage}
\begin{verbatim}
LRT_apa(data, md = FALSE)
\end{verbatim}
\end{Usage}
%
\begin{Arguments}
\begin{ldescription}
\item[\code{data}] anova data, LRT data

\item[\code{md}] md=TRUE
\end{ldescription}
\end{Arguments}
%
\begin{Value}
table explain
\end{Value}
%
\begin{Examples}
\begin{ExampleCode}

## Not run: 
#'

library(jjstat)
HS.model <- '
    visual  =~ x1 + b1*x2 + x3
    textual =~ x4 + b2*x5 + x6
    speed   =~ x7 + b3*x8 + x9
'
fit1 <- cfa(HS.model, data = HolzingerSwineford1939)
fit0 <- cfa(HS.model, data = HolzingerSwineford1939,
            orthogonal = TRUE)

fit1 %>% sem_effect(effect ="=~")
fit0 %>% sem_effect(effect ="=~")
lavTestLRT(fit1, fit0)
lavTestLRT(fit1, fit0) %>% LRT_apa()

## End(Not run)


\end{ExampleCode}
\end{Examples}
\inputencoding{utf8}
\HeaderA{LRT\_apa2}{LRT\_apa}{LRT.Rul.apa2}
%
\begin{Description}
LRT\_apa
\end{Description}
%
\begin{Usage}
\begin{verbatim}
LRT_apa2(data, md = FALSE)
\end{verbatim}
\end{Usage}
%
\begin{Arguments}
\begin{ldescription}
\item[\code{data}] anova data, LRT data

\item[\code{md}] md=TRUE
\end{ldescription}
\end{Arguments}
%
\begin{Value}
table explain
\end{Value}
%
\begin{Examples}
\begin{ExampleCode}

## Not run: 
#'

library(jjstat)
HS.model <- '
    visual  =~ x1 + b1*x2 + x3
    textual =~ x4 + b2*x5 + x6
    speed   =~ x7 + b3*x8 + x9
'
fit1 <- cfa(HS.model, data = HolzingerSwineford1939)
fit0 <- cfa(HS.model, data = HolzingerSwineford1939,
            orthogonal = TRUE)

fit1 %>% sem_effect(effect ="=~")
fit0 %>% sem_effect(effect ="=~")
lavTestLRT(fit1, fit0)
lavTestLRT(fit1, fit0) %>% LRT_apa()

## End(Not run)


\end{ExampleCode}
\end{Examples}
\inputencoding{utf8}
\HeaderA{make\_c\_text}{make\_c\_text}{make.Rul.c.Rul.text}
%
\begin{Description}
make\_c\_text
\end{Description}
%
\begin{Usage}
\begin{verbatim}
make_c_text(data)
\end{verbatim}
\end{Usage}
%
\begin{Arguments}
\begin{ldescription}
\item[\code{data}] vector data
\end{ldescription}
\end{Arguments}
%
\begin{Value}
text
\end{Value}
%
\begin{Examples}
\begin{ExampleCode}


## Not run: 

# Example usage
data <- c("FieldGoals", "OtherTDs", "YardsRushAtt", "RushYards",
          "RushFirstDown", "YardsPassComp", "PassYards", "PassFirstDown",
          "YardsRushAtt", "RushYards", "RushFirstDown", "YardsPassComp",
          "PassYards", "PassFirstDown", "PointsGame", "OffensTD",
          "TDGame", "Speical", "Rushing", "Passing",
          "Offense", "Scoring", 123, 456.78)

make_c_text(data)


## End(Not run)
\end{ExampleCode}
\end{Examples}
\inputencoding{utf8}
\HeaderA{make\_df\_text}{data.frame to input text}{make.Rul.df.Rul.text}
%
\begin{Description}
data.frame to input text
\end{Description}
%
\begin{Usage}
\begin{verbatim}
make_df_text(data, rownames = FALSE, text = TRUE, col_text = 1)
\end{verbatim}
\end{Usage}
%
\begin{Arguments}
\begin{ldescription}
\item[\code{data}] data.frame

\item[\code{rownames}] rownames=TRUE, FALSE -> nothing rownames

\item[\code{text}] Convert all configurations to char, e.g. '10'

\item[\code{col\_text}] columns to be used as letters, the rest as numeric
\end{ldescription}
\end{Arguments}
%
\begin{Value}
text
\end{Value}
%
\begin{Examples}
\begin{ExampleCode}

## Not run: 

# 데이터 프레임 생성
cho2013_cor <- data.frame(
  필요성 = c(1.000, 0.320, 0.266, 0.506, 0.429, 0.401),
  수업 = c(0.320, 1.000, 0.647, 0.463, 0.548, 0.554),
  교사 = c(0.266, 0.647, 1.000, 0.383, 0.491, 0.545),
  계획 = c(0.506, 0.463, 0.383, 1.000, 0.531, 0.491),
  영향 = c(0.429, 0.548, 0.491, 0.531, 1.000, 0.743),
  만족도 = c(0.401, 0.554, 0.545, 0.491, 0.743, 1.000)
)
cho2013_cor %>% make_df_text()



## End(Not run)

\end{ExampleCode}
\end{Examples}
\inputencoding{utf8}
\HeaderA{make\_label\_list}{Create Label Mapping List for Node Renaming}{make.Rul.label.Rul.list}
%
\begin{Description}
이 함수는 (원래이름, 변경라벨) 쌍으로 구성된 문자 벡터를 입력받아,
각 쌍을 `list(node = ..., to = ...)` 형태의 리스트로 변환한다.
주로 semPlot, qgraph, 사용자 정의 시각화 함수에서
노드 이름 일괄 변경을 위해 사용된다.
\end{Description}
%
\begin{Usage}
\begin{verbatim}
make_label_list(x)
\end{verbatim}
\end{Usage}
%
\begin{Arguments}
\begin{ldescription}
\item[\code{x}] character vector.
반드시 (원래이름, 라벨) 쌍으로 구성되어야 하며,
길이는 짝수여야 한다.
\end{ldescription}
\end{Arguments}
%
\begin{Details}
입력 벡터의 길이가 홀수인 경우,
(원래이름, 라벨) 쌍이 완성되지 않으므로 오류를 발생시킨다.

내부적으로는 `split()`을 사용하여 2개씩 그룹화한 뒤,
`lapply()`를 통해 리스트 형태로 변환한다.
\end{Details}
%
\begin{Value}
A list of lists.
각 원소는 다음 구조를 가진다:
\begin{itemize}

\item{} node: 원래 변수명
\item{} to: 변경할 라벨

\end{itemize}

\end{Value}
%
\begin{Examples}
\begin{ExampleCode}
# Example 1: 기본 사용
labels <- c(
  "REFF", "연구효능감",
  "INTR", "진로흥미",
  "OEXP", "결과기대"
)

make_label_list(labels)

# Example 2: semPlot/qgraph용 라벨 매핑
var_named <- c(
  "efficacy2", "연구효능2",
  "efficacy3", "연구효능3",
  "efficacy4", "연구효능4"
)

label_list <- make_label_list(var_named)

# label_list[[1]]
# $node
# [1] "efficacy2"
# $to
# [1] "연구효능2"

\end{ExampleCode}
\end{Examples}
\inputencoding{utf8}
\HeaderA{make\_layout}{Create a Custom Layout Matrix}{make.Rul.layout}
%
\begin{Description}
Create a Custom Layout Matrix
\end{Description}
%
\begin{Usage}
\begin{verbatim}
make_layout(nc = 10, nr = NULL)
\end{verbatim}
\end{Usage}
%
\begin{Arguments}
\begin{ldescription}
\item[\code{nc}] Integer. Number of columns in the layout matrix. Default is 10.

\item[\code{nr}] Integer or NULL. Number of rows in the layout matrix. If NULL, the number of rows will be automatically determined. Default is NULL.
\end{ldescription}
\end{Arguments}
%
\begin{Value}
A matrix representing the layout, optionally edited interactively using the `edit` function.
\end{Value}
%
\begin{Examples}
\begin{ExampleCode}
## Not run: 
make_layout(nc = 7)
make_layout(nc = 7, nr = 6)
make_layout(3, 6)

## End(Not run)
\end{ExampleCode}
\end{Examples}
\inputencoding{utf8}
\HeaderA{make\_list\_text}{make\_list\_text}{make.Rul.list.Rul.text}
%
\begin{Description}
make\_list\_text
\end{Description}
%
\begin{Usage}
\begin{verbatim}
make_list_text(data)
\end{verbatim}
\end{Usage}
%
\begin{Arguments}
\begin{ldescription}
\item[\code{data}] list data
\end{ldescription}
\end{Arguments}
%
\begin{Value}
text
\end{Value}
%
\begin{Examples}
\begin{ExampleCode}

## Not run: 
#

# Example data (characters)
data1 <- list(
  Speical = c("FieldGoals", "OtherTDs"),
  Rushing = c("YardsRushAtt", "RushYards", "RushFirstDown"),
  Passing = c("YardsPassComp", "PassYards", "PassFirstDown"),
  Offense = c("rush1", "pass1"),
  Scoring = c("PointsGame", "OffensTD", "TDGame")
)
make_list_text(data1)

# Example data (numeric)
data2 <- list(
  Speical = 7:8,
  Rushing = 1:3,
  Passing = 4:6,
  Offense = 18:19,
  Scoring = 9:11
)
make_list_text(data2)



# Example data (numeric)
list(7:8, 1:3, 4:6, 18:19, 9:11)%>%make_list_text()
#
#

## End(Not run)


\end{ExampleCode}
\end{Examples}
\inputencoding{utf8}
\HeaderA{make\_mat\_text}{make\_mat\_text}{make.Rul.mat.Rul.text}
%
\begin{Description}
Converts a matrix into a formatted text representation for layout creation.
\end{Description}
%
\begin{Usage}
\begin{verbatim}
make_mat_text(mat, text = TRUE, byrow = TRUE)
\end{verbatim}
\end{Usage}
%
\begin{Arguments}
\begin{ldescription}
\item[\code{mat}] Matrix data.

\item[\code{text}] Logical. If `TRUE`, the elements are wrapped in quotes. If `FALSE`, numbers or `NA` are returned.

\item[\code{byrow}] Logical. If `TRUE`, the matrix is converted row by row; otherwise, column by column.
\end{ldescription}
\end{Arguments}
%
\begin{Value}
Formatted text representation of the matrix.
\end{Value}
%
\begin{Examples}
\begin{ExampleCode}

## Not run: 
# Example usage
mmm <- matrix(c(1, 2, 3, 4, 5, 6, 7, 8, 9,
10, 11, 12, 13, 14, 15, 16, 17, 18, 19, 20,
21, 22, 23, 24, 25, 26, 27, 28, 29, 30),
 nrow = 5, ncol = 6)

make_mat_text(mmm, text = TRUE, byrow = TRUE)
make_mat_text(mmm, text = TRUE, byrow = FALSE)
make_mat_text(mmm, text = FALSE, byrow = TRUE)
make_mat_text(mmm, text = FALSE, byrow = FALSE)

# Example with NA
mmm_with_na <- matrix(c(1, NA, 3, NA, 5, 6, NA, 8, 9, 10), nrow = 2)
make_mat_text(mmm_with_na, text = FALSE, byrow = TRUE)

data(offense)
offense %>%str()

nfl_path = matrix(
  c(0, 0, 0, 0, 0,
    0, 0, 0, 0, 0,
    0, 0, 0, 0, 0,
    0, 1, 1, 0, 0,
    1, 0, 0, 1, 0),
  nrow = 5, ncol = 5, byrow = TRUE )

rownames(nfl_path) <- c("Speical", "Rushing",
"Passing", "Offense", "Scoring")
colnames(nfl_path) <- c("Speical", "Rushing",
"Passing", "Offense", "Scoring")
nfl_path

nfl_blocks = list(7:8, 1:3, 4:6,1:6, 9:11)

plspm_blocks_match(data=offense, blocks=nfl_blocks,
                   paths=nfl_path)

nfl_mode =c("B","A","A","A","A")

nfl_pls = plspm(Data=offense, nfl_path, nfl_blocks, modes = nfl_mode)
nfl_pls%>%summary()

#########
nfl_pls_boot = plspm_sem(Data=offense, nfl_path,
 nfl_blocks, modes = nfl_mode)

#Layouts extracted from mockups
plspm_semPaths2(nfl_pls_boot,sizeLat = 6)%>%
  plspm_factor_layout()

#Make it input material
plspm_semPaths2(nfl_pls_boot,sizeLat =6)%>%
  plspm_factor_layout() %>%layout_mat(text=T)

#Organizing materials for input
plspm_semPaths2(nfl_pls_boot,sizeLat =6)%>%
  plspm_factor_layout() %>%layout_mat()%>%
  make_mat_text2()

## End(Not run)

\end{ExampleCode}
\end{Examples}
\inputencoding{utf8}
\HeaderA{make\_mat\_text2}{Functions that let you input data in matrix form}{make.Rul.mat.Rul.text2}
%
\begin{Description}
Functions that let you input data in matrix form
\end{Description}
%
\begin{Usage}
\begin{verbatim}
make_mat_text2(data, text = TRUE, byrow = FALSE)
\end{verbatim}
\end{Usage}
%
\begin{Arguments}
\begin{ldescription}
\item[\code{data}] input dat a

\item[\code{text}] text = TRUE is text FALSE numeric

\item[\code{byrow}] byrow=FALSE  or TRUE
\end{ldescription}
\end{Arguments}
%
\begin{Value}
make test
\end{Value}
%
\begin{Examples}
\begin{ExampleCode}

## Not run: 
lay
V1      V2 V3        V4        V5       V6        V7        V8 V9
[1,] "ceo_8" NA "satis_5" "satis_4" ""       "satis_3" "satis_1" NA ""
[2,] ""      NA ""        ""        ""       ""        ""        NA ""
[3,] "ceo_7" NA ""        ""        ""       ""        ""        NA ""
[4,] ""      NA ""        ""        "SATISF" ""        ""        NA ""
[5,] "ceo_6" NA ""        ""        ""       ""        ""        NA "stay_2"
[6,] ""      NA ""        ""        ""       ""        ""        NA ""
[7,] "ceo_4" NA "CEOCOM"  ""        ""       ""        "STAYON"  NA ""
[8,] ""      NA ""        ""        ""       ""        ""        NA ""
[9,] "ceo_3" NA ""        ""        ""       ""        ""        NA "stay_1"
[10,] ""      NA ""        ""        ""       ""        ""        NA ""
[11,] "ceo_2" NA ""        ""        ""       ""        ""        NA ""
[12,] ""      NA ""        ""        ""       ""        ""        NA ""
[13,] "ceo_1" NA ""        ""        ""       ""        ""        NA ""

lay %>% make_mat_text2()


NewMat = matrix(c('ceo_8', NA, 'ceo_7', NA, 'ceo_6',
NA, 'ceo_4', NA, 'ceo_3', NA, 'ceo_2', NA, 'ceo_1',
NA, NA, NA, NA, NA, NA, NA, NA, NA, NA, NA, NA, NA,
'satis_5', NA, NA, NA, NA, NA, 'CEOCOM', NA, NA, NA,
NA, NA, NA, 'satis_4', NA, NA, NA, NA, NA, NA, NA,
 NA, NA, NA, NA, NA, NA, NA, NA, 'SATISF', NA, NA,
  NA, NA, NA, NA, NA, NA, NA, 'satis_3', NA, NA,
  NA, NA, NA, NA, NA, NA, NA, NA, NA, NA, 'satis_1',
  NA, NA, NA, NA, NA, 'STAYON', NA, NA, NA, NA, NA,
  NA, NA, NA, NA, NA, NA, NA, NA, NA, NA, NA, NA,
   NA, NA, NA, NA, NA, NA, 'stay_2', NA, NA, NA,
   'stay_1', NA, NA, NA, NA), nrow = 13, ncol = 9)


colnames(NewMat) <- c("V1", "V2", "V3", "V4", "V5", "V6", "V7", "V8", "V9")

#as.numeric
# NewMat <- apply(NewMat, 2, function(x) as.numeric(gsub("'", "", x)))

V1      V2 V3        V4        V5       V6        V7        V8 V9
[1,] "ceo_8" NA "satis_5" "satis_4" NA       "satis_3" "satis_1" NA NA
[2,] NA      NA NA        NA        NA       NA        NA        NA NA
[3,] "ceo_7" NA NA        NA        NA       NA        NA        NA NA
[4,] NA      NA NA        NA        "SATISF" NA        NA        NA NA
[5,] "ceo_6" NA NA        NA        NA       NA        NA        NA "stay_2"
[6,] NA      NA NA        NA        NA       NA        NA        NA NA
[7,] "ceo_4" NA "CEOCOM"  NA        NA       NA        "STAYON"  NA NA
[8,] NA      NA NA        NA        NA       NA        NA        NA NA
[9,] "ceo_3" NA NA        NA        NA       NA        NA        NA "stay_1"
[10,] NA      NA NA        NA        NA       NA        NA        NA NA
[11,] "ceo_2" NA NA        NA        NA       NA        NA        NA NA
[12,] NA      NA NA        NA        NA       NA        NA        NA NA
[13,] "ceo_1" NA NA        NA        NA       NA        NA        NA NA
#' #'
 lay_str1 %>%make_mat_text3(byrow=T)


 NewMat = matrix(c(NA, NA, 'PEOU', NA, NA, NA, NA,
 NA, NA, NA, NA, NA, NA, NA, 'CUX', NA, NA, NA,
  NA, NA, NA, NA, NA, NA, NA, 'AT', NA, 'USE',
  'SUX', NA, NA, NA, NA, NA, NA, NA, NA, NA, NA,
  NA, NA, NA, NA, NA, 'PU', NA, NA, NA, NA),
  nrow = 7, ncol = 7, byrow = TRUE )

 rownames(NewMat) <- c("row1", "row2", "row3", "row4", "row5", "row6", "row7")
colnames(NewMat) <- c("col1", "col2", "col3", "col4", "col5", "col6", "col7")

NewMat %>%make_mat_text2(byrow=FALSE)
NewMat %>%make_mat_text2(byrow=TRUE)
#'
#'

## End(Not run)


\end{ExampleCode}
\end{Examples}
\inputencoding{utf8}
\HeaderA{make\_mat\_text\_col}{첫번째 열이 존재하는 경우에 사용하는 make\_mat\_text2}{make.Rul.mat.Rul.text.Rul.col}
%
\begin{Description}
첫번째 열이 존재하는 경우에 사용하는 make\_mat\_text2
\end{Description}
%
\begin{Usage}
\begin{verbatim}
make_mat_text_col(data, text = TRUE, firstcol = 1)
\end{verbatim}
\end{Usage}
%
\begin{Arguments}
\begin{ldescription}
\item[\code{data}] data.frame

\item[\code{text}] text TRUE

\item[\code{firstcol}] ncol setting
\end{ldescription}
\end{Arguments}
%
\begin{Value}
text
\end{Value}
%
\begin{Examples}
\begin{ExampleCode}

## Not run: 
#'
data_text = "Tobacco_Dependence (0,5] (5,10] (10,15] (15,20] (20,98]
No_Nicotine_Dependence 130 210 43 114 20
Nicotine_Dependence 119 267 91 254 67"
nt=data_text%>% text2df(header=T)%>%tibble::column_to_rownames("Tobacco_Dependence")
nt %>%tibble()
nt=Char2num(nt,1:5)
nt

## End(Not run)


\end{ExampleCode}
\end{Examples}
\inputencoding{utf8}
\HeaderA{make\_text\_df}{Text to data.frame /It is characterized by the use of blank spaces to separate content.}{make.Rul.text.Rul.df}
%
\begin{Description}
Text to data.frame /It is characterized by the use of blank spaces to separate content.
\end{Description}
%
\begin{Usage}
\begin{verbatim}
make_text_df(text, header = TRUE, type = "res", add_vars = NULL)
\end{verbatim}
\end{Usage}
%
\begin{Arguments}
\begin{ldescription}
\item[\code{text}] text(Row variable names must not contain spaces. )

\item[\code{header}] header = TRUE

\item[\code{type}] res, check

\item[\code{add\_vars}] additinal factor
\end{ldescription}
\end{Arguments}
%
\begin{Value}
table data.frame
\end{Value}
%
\begin{Examples}
\begin{ExampleCode}

## Not run: 
#'
text <- "변인 평균 표준편차 왜도 첨도
성취감 4.24 0.73 -0.82 0.39
인정감 4.19 0.60 -0.36 -0.42
책임감 4.36 0.56 -0.61 -0.29
발전성 3.65 0.89 -0.50 -0.04
감독 3.59 0.88 -0.46 0.13
보수 3.56 0.90 -0.41 -0.18
대인관계 3.75 0.81 -0.14 -0.54
작업환경 3.83 0.75 -0.21 -0.44
안정성 3.28 1.08 -0.28 -0.68
감정이입 3.96 0.72 -0.40 -0.33
직무적합 4.00 0.75 -0.35 -0.78
일체감 4.02 0.75 -0.59 -0.13
현학원재직 3.97 0.79 -0.77 0.37
타학원이직 2.99 1.06 -0.41 -0.77
타직종이직 2.76 1.07 -0.13 -0.70"

text2df(text)

text2df(text, vars=c(rep("동기요인", 4),
              rep("위생요인",5), rep("개인요인",6)  ))

text2df(text, T)%>% make_df_text()

#direct output
"var B se b t
성취감 -.03 .09 -.03 -.33
인정감 .24 .09 .20 2.50
책임감 .30 .08 .23 3.56***
발전성 -.08 .06 -.10 -1.29
감독 .06 .05 .08 1.35
보수 -.14 .07 -.18 -2.01*
대인관계 .06 .06 .06 .89
작업환경 .33 .09 .34 3.82***
안정성 .17 .05 .26 3.42**
성별 -.05 .09 -.03 -.52
연령 -.20 .07 -.22 -3.01**
학력 .05 .07 .04 .70
소득 .00 .04 .01 .12
종사기간 .17 .06 .23 2.71**
기관규모 -.10 .03 -.20 -3.36**" %>%text2df()

reg_text = "var B se b t
성취감 -.03 .09 -.03 -.33
인정감 .24 .09 .20 2.50
책임감 .30 .08 .23 3.56***
발전성 -.08 .06 -.10 -1.29
감독 .06 .05 .08 1.35
보수 -.14 .07 -.18 -2.01*
대인관계 .06 .06 .06 .89
작업환경 .33 .09 .34 3.82***
안정성 .17 .05 .26 3.42**
성별 -.05 .09 -.03 -.52
연령 -.20 .07 -.22 -3.01**
학력 .05 .07 .04 .70
소득 .00 .04 .01 .12
종사기간 .17 .06 .23 2.71**
기관규모 -.10 .03 -.20 -3.36**"

make_text_to_df(reg_text)
#'
make_text_to_df(reg_text) %>% tibble()

#text to numeric
make_text_to_df(reg_text)%>% Char2num() %>% tibble()

make_text_to_df(reg_text)%>% Char2num(iter=2:4) %>% tibble()
#'

text ="d           Estimate  Std.Err  z-value  P(>|z|)
  read~ppsych           -0.275    0.037   -7.385    0.000
    read~motiv             0.461    0.037   12.404    0.000
arith~motiv             0.600    0.036   16.771    0.000
"
make_text_to_df(text)
make_text_to_df(text) %>% tibble::tibble()

text ="             Estimate Std.Error tvalue Pr(>|t|)
(Intercept) -1.336e-07  3.608e-01   0.000        1
ppsych      -2.747e-01  3.730e-02  -7.363  7.51e-13***
motiv        4.613e-01  3.730e-02  12.367  <2e-16***
"
make_text_to_df(text) %>%tibble()
make_text_to_df(text) %>%make_df_text()


text ="                   Estimate  Std.Err  z-value  P(>|z|)

    read~ppsych           -0.275    0.037   -7.385    0.000
    read~motiv             0.461    0.037   12.404    0.000

    arith~ppsych           -0.096    0.037   -2.616    0.009
     arith~motiv             0.576    0.037   15.695    0.000
"
make_text_to_df(text)



text =" 과장 대리 부장 사원 주임 차장 SUM
남자 30(39.5%) 21(38.9%) 34(79.1%) 30(49.2%)  10(37%)     28(71.8%) 153
여자 46(60.5%) 33(61.1%) 9(20.9%) 31(50.8%) 17(63%)    11(28.2%) 147
SUM 76 54 43 61 27 39 300
"
make_text_to_df(text)
#' 
## End(Not run)

\end{ExampleCode}
\end{Examples}
\inputencoding{utf8}
\HeaderA{markdown\_table}{A tool that outputs data to the viewer}{markdown.Rul.table}
%
\begin{Description}
knitr's kable function is the foundation of this package. However, it has many latex/html specific arguments hidden under the ground unless you check its source code. This wrapper function is created to provide better documentation (and auto-complete yay) and at the same time, solve the auto format setting in a better way.
kableExtra uses the built-in bootstrap themes by default in kable\_styling(). Alternatively, you can use a customized table themes for your table. This lightable table style sheet comes with three formats, namely lightable-minimal, lightable-classic, lightable-material and lightable-material-dark with hover and striped options.
\end{Description}
%
\begin{Usage}
\begin{verbatim}
markdown_table(
  data,
  caption = "Table 1. Capiton title input",
  digits = 3,
  font_size = 13,
  full_width = F,
  table = "paper",
  show = "data",
  format = "html",
  row.names = NA,
  col.names = NA,
  align = NULL,
  label = NULL,
  format.args = list(),
  escape = TRUE,
  table.attr = "",
  booktabs = FALSE,
  longtable = FALSE,
  valign = "t",
  position = "center",
  centering = TRUE,
  lightable_options = "basic",
  html_font = "\"Arial Narrow\", arial, helvetica, sans-serif",
  general_title = "Note: ",
  number = NULL,
  alphabet = NULL,
  catout = FALSE,
  general = NULL
)
\end{verbatim}
\end{Usage}
%
\begin{Arguments}
\begin{ldescription}
\item[\code{data}] data.frame or etc data, For , is an R object, which is typically a matrix or data frame. For , a list with each element being a returned value from .kable()xkables()kable()

\item[\code{caption}] table title

\item[\code{digits}] Maximum number of digits for numeric columns, passed to round(). This can also be a vector of length ncol(x), to set the number of digits for individual columns.

\item[\code{font\_size}] table font size

\item[\code{full\_width}] table wide and fit values

\item[\code{table}] tbble type 'paper', 'classic', 'dark', 'minimal'

\item[\code{show}] 'markdown' is output to the viewer, and 'data' is output to the console.

\item[\code{format}] 'markdown' and 'html'. default 'html'. A character string. Possible values are , , (Pandoc's pipe tables), (Pandoc's simple tables), and . The value of this argument will be automatically determined if the function is called within a knitr document. The value can also be set in the global option . If is a function, it must return a character string.latexhtmlpipesimplerstformatknitr.table.formatformat

\item[\code{row.names}] row names in table(Set if necessary), Logical: whether to include row names. By default, row names are included if is neither nor identical to .rownames(x)NULL1:nrow(x)

\item[\code{col.names}] column names in table (Set if necessary)

\item[\code{align}] Table value align, Column alignment: a character vector consisting of 'l' (left), 'c' (center) and/or 'r' (right). By default or if align = NULL, numeric columns are right-aligned, and other columns are left-aligned. If length(align) == 1L, the string will be expanded to a vector of individual letters, e.g. 'clc' becomes c('c', 'l', 'c'), unless the output format is LaTeX.

\item[\code{label}] Table label, he table reference label. By default, the label is obtained from .knitr::opts\_current\$get('label')

\item[\code{format.args}] A list of arguments to be passed to to format table values, e.g. .format()list(big.mark = ',')

\item[\code{escape}] Boolean; whether to escape special characters when producing HTML or LaTeX tables. When , you have to make sure that special characters will not trigger syntax errors in LaTeX or HTML.escape = FALSE

\item[\code{table.attr}] A character string for addition HTML table attributes. This is convenient if you simply want to add a few HTML classes or styles. For example, you can put 'class="table" style="color: red"'.

\item[\code{booktabs}] booktabs - T/F for whether to use the longtable format. If you have a table that will span over two or more pages, you will have to turn this on.

\item[\code{longtable}] longtable

\item[\code{valign}] You probably won't need to adjust this latex option very often. If you are familar with latex tables, this is the optional position for the tabular environment controling the vertical position of the table relative to the baseline of the surrounding text. Possible choices are , and (default)

\item[\code{position}] kable\_styling: A character string determining how to position the table on a page. Possible values include left, center, right, float\_left and float\_right. Please see the package doc site for demonstrations. For a LaTeX table, if ⁠float\_*⁠ is selected, LaTeX package wrapfig will be imported.

\item[\code{centering}] Align cener in Table default TRUE

\item[\code{lightable\_options}] Options to customize lightable. Similar with bootstrap\_options in kable\_styling. Choices include basic, striped and hover.

\item[\code{html\_font}] A string for HTML css font. For example, html\_font = '"Arial Narrow", arial, helvetica, sans-serif'.Everything else you need to specify in kable\_styling.

\item[\code{general\_title}] Section header for general footnotes. Default is "Note: ".

\item[\code{number}] A vector of footnote texts. Footnotes here will be numbered. There is no upper cap for the number of footnotes here

\item[\code{alphabet}] A vector of footnote texts, Footnotes here will be labeled with abc. The vector here should not have more than 26 elements.

\item[\code{catout}] \#p value msg

\item[\code{general}] \#p value msg
\end{ldescription}
\end{Arguments}
%
\begin{Examples}
\begin{ExampleCode}
# example code

  ## Not run: 
  markdown_table(mtcars[1:5,], table="minimal")

  ###*** : p < .001, ** : p < .01, * : p < .05"
  
## End(Not run)
\end{ExampleCode}
\end{Examples}
\inputencoding{utf8}
\HeaderA{mat\_text\_arrange}{mat\_text\_arrange}{mat.Rul.text.Rul.arrange}
%
\begin{Description}
mat\_text\_arrange
\end{Description}
%
\begin{Usage}
\begin{verbatim}
mat_text_arrange(mat, text = TRUE, byrow = FALSE)
\end{verbatim}
\end{Usage}
%
\begin{Arguments}
\begin{ldescription}
\item[\code{mat}] matrix data

\item[\code{text}] TRUE '' FALSE number

\item[\code{byrow}] FALSE, byrow =TRUE fill in row
\end{ldescription}
\end{Arguments}
%
\begin{Value}
text
\end{Value}
%
\begin{Examples}
\begin{ExampleCode}
## Not run: 
data(offense)
offense %>%str()

nfl_path = matrix(
  c(0, 0, 0, 0, 0,
    0, 0, 0, 0, 0,
    0, 0, 0, 0, 0,
    0, 1, 1, 0, 0,
    1, 0, 0, 1, 0),
  nrow = 5, ncol = 5, byrow = TRUE )

rownames(nfl_path) <- c("Speical", "Rushing", "Passing", "Offense", "Scoring")
colnames(nfl_path) <- c("Speical", "Rushing", "Passing", "Offense", "Scoring")
nfl_path

nfl_blocks = list(7:8, 1:3, 4:6,1:6, 9:11)

plspm_blocks_match(data=offense, blocks=nfl_blocks,
                   paths=nfl_path)

nfl_mode =c("B","A","A","A","A")

nfl_pls = plspm(Data=offense, nfl_path, nfl_blocks, modes = nfl_mode)
nfl_pls%>%summary()

#########
nfl_pls_boot = plspm_sem(Data=offense, nfl_path, nfl_blocks, modes = nfl_mode)

#Layouts extracted from mockups
plspm_semPaths2(nfl_pls_boot,sizeLat = 6)%>%
  plspm_factor_layout()

#Make it input material
plspm_semPaths2(nfl_pls_boot,sizeLat =6)%>%
  plspm_factor_layout() %>%layout_mat(text=T)

#Organizing materials for input
plspm_semPaths2(nfl_pls_boot,sizeLat =6)%>%
  plspm_factor_layout() %>%layout_mat()%>%
  mat_text_arrange()

## End(Not run)

\end{ExampleCode}
\end{Examples}
\inputencoding{utf8}
\HeaderA{mcnemar\_test}{mcnemar\_test}{mcnemar.Rul.test}
%
\begin{Description}
mcnemar\_test
\end{Description}
%
\begin{Usage}
\begin{verbatim}
mcnemar_test(mat, correct = TRUE, exact = TRUE)
\end{verbatim}
\end{Usage}
%
\begin{Arguments}
\begin{ldescription}
\item[\code{mat}] matrix

\item[\code{correct}] McNemar's Chi-squared test with continuity correction

\item[\code{exact}] CI check binom.test
\end{ldescription}
\end{Arguments}
%
\begin{Value}
reusult
\end{Value}
%
\begin{Examples}
\begin{ExampleCode}

## Not run: 
#'
Buying_car  = matrix(c(23,28,7,12), nrow=2,
                     dimnames= list(c("시승전_있음","시승전_없음"),
                                    c("시승후_있음","시승후_없음")))
Buying_car
mcnemar_test(Buying_car)
mcnemar_test(Buying_car, correct=FALSE)

gen_table = matrix(c(c(9, 3, 13, 10)),ncol=2,
                   dimnames=list("case"=c("yes", "no"),
                                 "control"=c("yes", "no")   ))
gen_table
#Rbase
mcnemar.test(gen_table, correct = FALSE)
mcnemar_test(gen_table, correct = FALSE)



#모든 종류의 분석을 시행한 경우
bind_cols(
  "TEST"= c("correct and exact", "not Correct"," Not exact", "Not correct, exact"),
  bind_rows(
    mcnemar_test(Buying_car),  # 디폴트ㅡ
    mcnemar_test(Buying_car, correct=FALSE), # 연속보정 중지
    mcnemar_test(Buying_car,exact=FALSE), # 정확도 중지
    mcnemar_test(Buying_car,exact=FALSE, correct=FALSE)# 연속성 보정, 정확도 중지,SAS
  )
)%>%select(-b,-c,-odds_ratio)

## End(Not run)
\end{ExampleCode}
\end{Examples}
\inputencoding{utf8}
\HeaderA{md}{md: Generate Styled Data Tables for Reports}{md}
%
\begin{Description}
This function generates styled tables for reports, supporting different formats and options for enhanced visual presentation.
\end{Description}
%
\begin{Usage}
\begin{verbatim}
md(
  data,
  caption = "Table 1. Capiton title input",
  digits = 3,
  font_size = 13,
  full_width = FALSE,
  table = "paper",
  show = "data",
  format = "html",
  row.names = NA,
  col.names = NA,
  align = NULL,
  label = NULL,
  format.args = list(),
  escape = TRUE,
  table.attr = "",
  booktabs = FALSE,
  longtable = FALSE,
  valign = "t",
  position = "center",
  centering = TRUE,
  lightable_options = "basic",
  html_font = "\"Arial Narrow\", arial, helvetica, sans-serif",
  general_title = "Note: ",
  number = NULL,
  alphabet = NULL,
  catout = FALSE,
  general = NULL,
  pass = FALSE
)
\end{verbatim}
\end{Usage}
%
\begin{Arguments}
\begin{ldescription}
\item[\code{data}] A data frame or tibble to be styled and displayed.

\item[\code{caption}] A character string for the table caption. Default is "Table 1. Capiton title input".

\item[\code{digits}] An integer indicating the number of decimal places. Default is 3.

\item[\code{font\_size}] A numeric value specifying the font size. Default is 13.

\item[\code{full\_width}] Logical, whether the table should use full width. Default is FALSE.

\item[\code{table}] A character string specifying the table style. Options are "paper", "classic", "dark", "minimal", and "basic". Default is "paper".

\item[\code{show}] A character string specifying the type of display. Options are "data", "lm", or "p\_add". Default is "data".

\item[\code{format}] The output format of the table. Default is "html".

\item[\code{row.names}] Logical or character vector for row names. Default is NA.

\item[\code{col.names}] A character vector specifying column names. Default is NA.

\item[\code{align}] A character vector specifying column alignment. Default is NULL.

\item[\code{label}] A character string for table label. Default is NULL.

\item[\code{format.args}] A list of additional arguments passed to `format()`. Default is an empty list.

\item[\code{escape}] Logical, whether to escape special characters. Default is TRUE.

\item[\code{table.attr}] A character string for additional HTML attributes. Default is an empty string.

\item[\code{booktabs}] Logical, whether to use LaTeX booktabs style. Default is FALSE.

\item[\code{longtable}] Logical, whether to use LaTeX longtable style. Default is FALSE.

\item[\code{valign}] A character string specifying vertical alignment. Default is "t".

\item[\code{position}] A character string for table position. Default is "center".

\item[\code{centering}] Logical, whether to center align the table. Default is TRUE.

\item[\code{lightable\_options}] A character string for lightable styling options. Default is "basic".

\item[\code{html\_font}] A character string specifying the font for HTML output. Default is '"Arial Narrow", arial, helvetica, sans-serif'.

\item[\code{general\_title}] A character string specifying the title for footnotes. Default is "Note: ".

\item[\code{number}] A numeric vector specifying footnote numbering. Default is NULL.

\item[\code{alphabet}] A character vector specifying footnote alphabet markers. Default is NULL.

\item[\code{catout}] Logical, whether to output significance markers (e.g., **, *). Default is FALSE.

\item[\code{general}] A character string specifying a general footnote. Default is NULL.

\item[\code{pass}] Logical, whether to process the table styling. If TRUE, returns the original data. Default is FALSE.
\end{ldescription}
\end{Arguments}
%
\begin{Value}
A styled table if `pass = TRUE`, otherwise returns the original data.
\end{Value}
%
\begin{Examples}
\begin{ExampleCode}
## Not run: 
  data <- data.frame(A = c(1.234, 2.345, 3.456), B = c(4.567, 5.678, 6.789))
  md(data, caption = "Example Table", table = "classic")

## End(Not run)
\end{ExampleCode}
\end{Examples}
\inputencoding{utf8}
\HeaderA{meanCentering}{centering each vector}{meanCentering}
%
\begin{Description}
centering each vector
\end{Description}
%
\begin{Usage}
\begin{verbatim}
meanCentering(variable, na.rm = TRUE)
\end{verbatim}
\end{Usage}
%
\begin{Arguments}
\begin{ldescription}
\item[\code{variable}] vector variable in data

\item[\code{na.rm}] missing data chedk
\end{ldescription}
\end{Arguments}
\inputencoding{utf8}
\HeaderA{meanCentering\_df}{meancentering mixed model}{meanCentering.Rul.df}
%
\begin{Description}
meancentering mixed model
\end{Description}
%
\begin{Usage}
\begin{verbatim}
meanCentering_df(data, grp_id, gv, ...)
\end{verbatim}
\end{Usage}
%
\begin{Arguments}
\begin{ldescription}
\item[\code{data}] data.frame

\item[\code{grp\_id}] group id

\item[\code{gv}] group variable, Group Variables That Separate Groups

\item[\code{...}] individual measurement independent variables evaluated at individual level
\end{ldescription}
\end{Arguments}
%
\begin{Examples}
\begin{ExampleCode}
## Not run: 
meanCentering_df(mtcars[,c("cyl","mpg","wt","hp")],"cyl","mpg","wt","hp")|> data.frame()

## End(Not run)



\end{ExampleCode}
\end{Examples}
\inputencoding{utf8}
\HeaderA{mean\_apa}{mean\_apa: Generate descriptive statistics in APA style}{mean.Rul.apa}
%
\begin{Description}
mean\_apa: Generate descriptive statistics in APA style
\end{Description}
%
\begin{Usage}
\begin{verbatim}
mean_apa(df, sd = FALSE)
\end{verbatim}
\end{Usage}
%
\begin{Arguments}
\begin{ldescription}
\item[\code{df}] A data frame containing descriptive statistics, including variables such as N, MEAN, SD, Skew, and Kurt.

\item[\code{sd}] A logical value. If TRUE, includes standard deviation (SD) along with the mean in the output. Default is FALSE.
\end{ldescription}
\end{Arguments}
%
\begin{Value}
A string formatted in APA style, summarizing the descriptive statistics in order of mean values and assessing normality based on skewness and kurtosis.
\end{Value}
%
\begin{Examples}
\begin{ExampleCode}
df <- data.frame(
  var = c("Q33_1", "Q33_2", "Q33_3"),
  N = c(4542, 4542, 4542),
  MEAN = c(1.9489, 1.9337, 1.9243),
  SD = c(0.84679, 0.86568, 0.86960),
  Skew = c(0.72576, 0.63900, 0.70717),
  Kurt = c(0.63369, 0.087763, 0.27470)
)
mean_apa(df, sd = TRUE)

# Example Output:
# "The results of the descriptive statistics analysis are as follows.
# Variables ordered by mean are Q33_1 (M=1.9489, SD=0.84679) > Q33_2 (M=1.9337, SD=0.86568) >
# Q33_3 (M=1.9243, SD=0.86960). Skewness ranged from 0.63900 to 0.72576, and kurtosis ranged from
# 0.087763 to 0.63369. Following Kline (2005), skewness values exceeding an absolute value of 3
# and kurtosis values exceeding an absolute value of 8 or 10 are considered non-normal. Skewness
# and kurtosis met the criteria for normality."
\end{ExampleCode}
\end{Examples}
\inputencoding{utf8}
\HeaderA{mean\_bar}{mean\_bar: Bar plot of means with optional sorting and customizations}{mean.Rul.bar}
%
\begin{Description}
This function generates a bar plot displaying mean values, allowing customization of axis labels, text size, angle, and sorting.
\end{Description}
%
\begin{Usage}
\begin{verbatim}
mean_bar(
  df,
  x.title = "문항",
  sort = FALSE,
  size.text = 4,
  text.angle = 0,
  size.axis = 12
)
\end{verbatim}
\end{Usage}
%
\begin{Arguments}
\begin{ldescription}
\item[\code{df}] A data frame containing 'var' and 'MEAN' columns.

\item[\code{x.title}] A character string for the x-axis title. Default is "문항".

\item[\code{sort}] Logical. If TRUE, sorts the data by mean values in ascending order. Default is FALSE.

\item[\code{size.text}] Numeric. Size of the text labels on bars. Default is 4.

\item[\code{text.angle}] Numeric. Angle of the x-axis text labels. Default is 0.

\item[\code{size.axis}] Numeric. Size of the axis text labels. Default is 12.
\end{ldescription}
\end{Arguments}
%
\begin{Value}
A ggplot object showing a bar plot of means.
\end{Value}
%
\begin{Examples}
\begin{ExampleCode}
## Not run: 
# Example dataset
df <- data.frame(var = c("Q1", "Q2", "Q3"), MEAN = c(3.5, 4.2, 2.8))

# Default bar plot
mean_bar(df)

# Sorted bar plot with customized text size and angle
mean_bar(df, sort = TRUE, size.text = 5, text.angle = 45)

## End(Not run)
\end{ExampleCode}
\end{Examples}
\inputencoding{utf8}
\HeaderA{med\_effect}{DE,IE,TE Effect extraction}{med.Rul.effect}
%
\begin{Description}
DE,IE,TE Effect extraction
\end{Description}
%
\begin{Usage}
\begin{verbatim}
med_effect(x, effect1 = "", caption = "", effect2 = ".", option = 1)
\end{verbatim}
\end{Usage}
%
\begin{Arguments}
\begin{ldescription}
\item[\code{x}] lavaan result

\item[\code{effect1}] first, you must select effect1 = IE, DE, TE

\item[\code{caption}] title

\item[\code{effect2}] '.' 2nd oprion

\item[\code{option}] 1(p\_value) or 2(ci)
\end{ldescription}
\end{Arguments}
%
\begin{Value}
effect result
\end{Value}
\inputencoding{utf8}
\HeaderA{metacorMat}{Creating a meta-structural equation correlation matrix JH, Park}{metacorMat}
%
\begin{Description}
Creating a meta-structural equation correlation matrix JH, Park
\end{Description}
%
\begin{Usage}
\begin{verbatim}
metacorMat(data, startcol, var = "")
\end{verbatim}
\end{Usage}
%
\begin{Arguments}
\begin{ldescription}
\item[\code{data}] data.frame, columns:  study, n, var1\_var2, var1\_var2, ...

\item[\code{startcol}] The column where the correlation coefficient starts, ej, and the column where the correlation coefficient ends are automatically calculated.

\item[\code{var}] Unique variable name with correlation analysis
\end{ldescription}
\end{Arguments}
%
\begin{Examples}
\begin{ExampleCode}
##make matrix
## Not run: 
 ## making correlation matrix ##
cormat1 = metacorMat(rd2, startcol= 3, var=c("po","ne","se","jp"))
 ##variable : "po_ne" "po_se" "po_jp" "ne_se" "ne_jp" "se_jp"

rd2<- data.frame(
  stringsAsFactors = FALSE,
  study = c(1L,
            2L,3L,4L,5L,6L,7L,8L,9L,10L,11L,12L,
            13L,14L,15L,16L,17L,18L,19L,20L,
            21L,22L,23L,24L,25L,26L,27L,28L,29L,
            30L,31L,32L,33L,34L,35L,36L),
  n = c(671L,
        92L,73L,321L,214L,947L,321L,231L,
        531L,567L,213L,928L,269L,593L,78L,983L,
        1504L,247L,79L,875L,68L,201L,891L,
        197L,431L,98L,179L,69L,301L,679L,273L,
        254L,316L,781L,273L,268L),
  po_ne = c(NA,
            NA,-0.33,NA,NA,NA,NA,-0.32,NA,NA,NA,
            NA,NA,NA,-0.21,-0.49,NA,-0.14,-0.29,
            -0.27,NA,-0.37,NA,NA,-0.35,-0.17,NA,
            -0.11,NA,NA,-0.41,NA,NA,NA,NA,-0.39),
  po_se = c(NA,
            0.37,0.29,0.16,0.41,NA,NA,0.21,0.59,
            0.53,0.31,0.38,NA,NA,0.19,NA,NA,NA,
            NA,0.21,0.07,0.57,0.45,NA,0.18,NA,
            0.45,NA,NA,0.09,NA,NA,0.37,0.19,0.23,
            NA),
  po_jp = c(0.16,
            0.39,0.09,0.07,NA,0.04,NA,NA,0.33,
            0.11,NA,0.04,NA,NA,0.21,0.23,0.09,
            0.11,0.31,0.11,-0.07,0.14,NA,0.19,0.13,
            0.25,0.35,0.29,0.35,NA,0.27,0.19,NA,
            NA,0.13,0.19),
  ne_se = c(NA,
            -0.31,NA,NA,NA,NA,NA,-0.19,NA,NA,NA,
            NA,-0.49,NA,-0.45,NA,NA,NA,NA,-0.38,
            NA,-0.39,NA,NA,-0.21,NA,NA,NA,NA,
            -0.23,NA,NA,NA,NA,-0.44,NA),
  ne_jp = c(NA,
            -0.41,NA,NA,NA,-0.18,-0.11,NA,NA,NA,
            NA,NA,-0.21,-0.15,-0.36,-0.29,NA,
            -0.27,-0.29,-0.23,NA,-0.27,NA,NA,-0.19,
            -0.07,NA,-0.31,NA,NA,0.07,NA,NA,NA,
            -0.27,-0.05),
  se_jp = c(NA,
            NA,NA,0.07,NA,NA,NA,NA,0.27,0.07,NA,
            NA,0.48,0.26,0.41,NA,NA,NA,NA,0.31,
            0.19,0.34,NA,NA,0.19,NA,0.39,NA,
            0.57,NA,NA,NA,NA,NA,NA,NA),
  jobyear = c("f",
              "f","s","s","s","s","f","f","s","s",
              "s","s","f","f","f","f","f","f","s",
              "s","s","s","s","f","s","f","f",
              "f","f","f","f","s","s","s","f","f"),
  wage = c(9L,
           1L,1L,1L,4L,9L,2L,3L,3L,3L,2L,7L,
           1L,8L,6L,3L,5L,7L,8L,4L,7L,1L,4L,
           6L,7L,1L,9L,6L,7L,7L,7L,9L,2L,5L,
           7L,3L)
)

## cormat1
## [[1]]
## po ne se   jp
## po 1.00 NA NA 0.16
## ne   NA NA NA   NA
## se   NA NA NA   NA
## jp 0.16 NA NA 1.00
##
## [[2]]
## po    ne    se    jp
## po   NA    NA  0.37  0.39
## ne   NA  1.00 -0.31 -0.41
## se 0.37 -0.31    NA    NA
## jp 0.39 -0.41    NA  1.00
## 출처 meta-analaysis 경로분석모델 1단계 - 상관행렬추정 correlation using metaSEM|자유자재Paper

### 1. How to obtain the correlation matrix using the fixed effects model
library(metaSEM)
fixed1 <- tssem1(Cov=cormat1, n=rd2$n, method="FEM")
summary(fixed1)
coef(fixed1)

### random effects model
random = tssem1(Cov=cormat1, n=rd2$n, method="REM" , RE.type = "Diag")
summary(random)

# 2. How to find correlation matrix using random effects model
remodel = tssem1(Cov=cormat1, n=rd2$n, method="REM" , RE.type = "Diag")
random1 = summary(remodel)
mat1 <- vec2symMat(matrix(random1$coefficients[1:6,1]), diag = FALSE)
colnames(mat1)=c("po","ne","se","jp")
rownames(mat1)=c("po","ne","se","jp")
mat1

##path modeling
### step1
rem = tssem1(Cov=cormat1, n=rd2$n, method="REM" , RE.type = "Diag")
summary(rem)

## path model to A
## se = b31*po + b32*ne +e3
## jp = b41*po + b42*ne + b43*se +e4
A= create.mxMatrix(
c(0,0,0,0,
  0,0,0,0,
  "0.1*b31","0.1*b32",0,0,
  "0.1*b41","0.1*b42","0.1*b43",0),
type="Full", nrow=4, ncol=4, byrow=TRUE, name="A")
A

#exogeous correlation matrix
varnames <- c("po","ne","se","jp")
S= create.mxMatrix(
c(1,
  "0.1*p21",1,
  0,        0,"1*p33",
  0,        0,      0,"1*p44"),
  type="Symm", byrow = TRUE,
  name="S",
  dimnames=list(varnames, varnames)
)
S


### step 2
stage2A = tssem2(rem, Amatrix = A, Smatrix = S, diag.constraints = TRUE)
summary(stage2A)
plot(stage2A, what="est", edge.label.cex = 0.9, edge.label.position=0.7,
    nDigits=3, layout = "spring")


###Alternative model
A2 = create.mxMatrix(
c(0,0,0,0,
  0,0,0,0,
  "0.1*b31","0.1*b32",0,0,
  0,"0.1*b42","0.1*b43",0),
type="Full", nrow=4, ncol=4, byrow=TRUE, name="A")
 A2


stage3A = tssem2(rem, Amatrix = A2, Smatrix = S,
                 diag.constraints = TRUE, intervals.type = "LB")
summary(stage3A)

plot(stage3A, what="est", edge.label.cex = 1, edge.color="gray20",
edge.label.position=0.5, nDigits=4,
layout = "spring", style="lisrel",residScale=16)

# Test the significance of indirect effect
stage4A = tssem2(rem, Amatrix = A2, Smatrix = S, diag.constraints = TRUE,
          intervals.type = "LB",
          mx.algebras = list(
                   IDEpotojp = mxAlgebra(b31*b43, name="IDEpotojp"),
                   IDEnetojp = mxAlgebra(b32*b43, name="IDEnetojp")
                  ))
  summary(stage4A)

## mxAlgebras objects (and their 95% likelihood-based CIs):
## lbound    Estimate      ubound
## IDEpotojp[1,1]  0.05315542  0.08413089  0.11799399
## IDEnetojp[1,1] -0.10185142 -0.06355753 -0.03440116


## End(Not run)

\end{ExampleCode}
\end{Examples}
\inputencoding{utf8}
\HeaderA{min\_max\_scaler}{비율 척도로 변환하는 함수 정의}{min.Rul.max.Rul.scaler}
%
\begin{Description}
비율 척도로 변환하는 함수 정의
\end{Description}
%
\begin{Usage}
\begin{verbatim}
min_max_scaler(data, columns = NULL)
\end{verbatim}
\end{Usage}
%
\begin{Arguments}
\begin{ldescription}
\item[\code{data}] data

\item[\code{columns}] colum all
\end{ldescription}
\end{Arguments}
%
\begin{Value}
df
\end{Value}
%
\begin{Examples}
\begin{ExampleCode}
## Not run: 
#'
 사용 예
data_frame: 데이터프레임 객체
변환할 열 이름을 명시적으로 지정하거나 NULL로 모든 열을 변환
scaled_data <- min_max_scaler(data_frame, columns = c("V1_1", "V2_1", "V3_1"))
print(scaled_data)

## End(Not run)
\end{ExampleCode}
\end{Examples}
\inputencoding{utf8}
\HeaderA{missCheck}{Check and resolve missing data-}{missCheck}
%
\begin{Description}
If there are missing values, print how many there are
\end{Description}
%
\begin{Usage}
\begin{verbatim}
missCheck(data, plot = FALSE)
\end{verbatim}
\end{Usage}
%
\begin{Arguments}
\begin{ldescription}
\item[\code{data}] data.frame

\item[\code{plot}] missing data plot
\end{ldescription}
\end{Arguments}
%
\begin{Value}
Results of checking for missingness
\end{Value}
%
\begin{Examples}
\begin{ExampleCode}
## Not run: 
 sleep %>% str()
sleep %>% missCheck()

## End(Not run)


\end{ExampleCode}
\end{Examples}
\inputencoding{utf8}
\HeaderA{missfillin}{missfillin: Fill Missing Values in a Data Frame}{missfillin}
%
\begin{Description}
This function fills missing values in a data frame using either a specified function (e.g., `mean` or `median`)
or a method such as `pmm` (predictive mean matching) through the `mice` package.
\end{Description}
%
\begin{Usage}
\begin{verbatim}
missfillin(data, fn = mean, method = NULL)
\end{verbatim}
\end{Usage}
%
\begin{Arguments}
\begin{ldescription}
\item[\code{data}] A data frame containing missing values.

\item[\code{fn}] A function to apply for filling missing values (default: `mean`).

\item[\code{method}] The method to use for imputing missing values. Accepts `NULL` (default function-based) or `"pmm"`.
\end{ldescription}
\end{Arguments}
%
\begin{Value}
A data frame with missing values imputed.
\end{Value}
%
\begin{Examples}
\begin{ExampleCode}
## Not run: 
# Example 1: Impute with mean function
data <- data.frame(a = c(1, 2, NA, 4), b = c(NA, 3, 5, 6))
result <- missfillin(data, fn = mean)
print(result)

# Example 2: Impute with median function
data <- data.frame(a = c(NA, 1, 2, 3), b = c(4, NA, NA, 8))
result <- missfillin(data, fn = median)
print(result)

# Example 3: Impute using predictive mean matching (pmm)
data <- data.frame(a = c(1, NA, 3, 4), b = c(NA, 2, NA, 5))
result <- missfillin(data, method = "pmm")
print(result)

# Example 4: Handle data without missing values
data <- data.frame(a = c(1, 2, 3, 4), b = c(4, 5, 6, 7))
result <- missfillin(data)
print(result)

# Example 5: Error when insufficient columns for pmm
data <- data.frame(a = c(1, 2, NA, 4))
result <- tryCatch(missfillin(data, method = "pmm"), error = function(e) e)
print(result)

## End(Not run)
\end{ExampleCode}
\end{Examples}
\inputencoding{utf8}
\HeaderA{missfillinFor}{missfillinFor: Fill Missing Values in Selected Columns of a Data Frame}{missfillinFor}
%
\begin{Description}
This function fills missing values in a specified range of columns within a data frame using either
a specified function (e.g., `mean`, `median`) or a method such as `pmm` (predictive mean matching).
\end{Description}
%
\begin{Usage}
\begin{verbatim}
missfillinFor(data, fn = mean, method = NULL, range = NULL)
\end{verbatim}
\end{Usage}
%
\begin{Arguments}
\begin{ldescription}
\item[\code{data}] A data frame containing missing values.

\item[\code{fn}] A function to apply for filling missing values (default: `mean`).

\item[\code{method}] The method to use for imputing missing values. Accepts `NULL` (default function-based) or `"pmm"`.

\item[\code{range}] A numeric vector indicating column indices or a character string specifying the range of column names (e.g., `"col1:col3"`).
If `NULL`, all columns are considered.
\end{ldescription}
\end{Arguments}
%
\begin{Value}
A data frame with missing values imputed in the specified range of columns.
\end{Value}
%
\begin{Examples}
\begin{ExampleCode}
## Not run: 
# Example 1: Impute all columns with mean function
data <- data.frame(a = c(1, 2, NA, 4), b = c(NA, 3, 5, 6), c = c(NA, NA, 7, 8))
result <- missfillinFor(data, fn = mean)
print(result)

# Example 2: Impute specific column range with median function
data <- data.frame(a = c(1, NA, 3, 4), b = c(NA, 2, NA, 5), c = c(6, 7, 8, 9))
result <- missfillinFor(data, fn = median, range = 1:2)
print(result)

# Example 3: Impute using predictive mean matching (pmm)
data <- data.frame(a = c(1, NA, 3, 4), b = c(NA, 2, NA, 5), c = c(6, 7, 8, 9))
result <- missfillinFor(data, method = "pmm", range = "a:b")
print(result)

# Example 4: Impute using column names
data <- data.frame(a = c(1, 2, NA, 4), b = c(NA, 3, 5, 6), c = c(NA, NA, 7, 8))
result <- missfillinFor(data, fn = mean, range = "a:c")
print(result)

# Example 5: Handle insufficient columns for pmm
data <- data.frame(a = c(1, 2, NA, 4))
result <- tryCatch(missfillinFor(data, method = "pmm", range = 1), error = function(e) e)
print(result)

## End(Not run)
\end{ExampleCode}
\end{Examples}
\inputencoding{utf8}
\HeaderA{mlsem\_fit\_indices}{Extract Model Fit Indices for Multilevel SEM}{mlsem.Rul.fit.Rul.indices}
%
\begin{Description}
This function extracts model fit indices from a fitted
multilevel SEM object (lavaan), with robust fit measures prioritized
when available. The output can be returned in long (APA-style) or
wide format depending on the selected type.
\end{Description}
%
\begin{Usage}
\begin{verbatim}
mlsem_fit_indices(
  fit,
  type = c("data", "long", "wide", "apa"),
  robust = TRUE,
  digits = 3
)
\end{verbatim}
\end{Usage}
%
\begin{Arguments}
\begin{ldescription}
\item[\code{fit}] A fitted lavaan object (e.g., from \code{sem()} with
\code{cluster} specified).

\item[\code{type}] Output type.
One of \code{"data"}, \code{"long"}, \code{"wide"}, or \code{"apa"}.
\begin{itemize}

\item{} \code{"data"}: APA-style long format (Index–Value).
\item{} \code{"long"}: Same as \code{"data"}.
\item{} \code{"wide"}: One-row wide data.frame.
\item{} \code{"apa"}: Same as \code{"wide"} (alias for reporting).

\end{itemize}


\item[\code{robust}] Logical. If TRUE (default), robust fit indices
(e.g., robust CFI, TLI, RMSEA) are used when available.

\item[\code{digits}] Integer. Number of decimal places for rounding.
\end{ldescription}
\end{Arguments}
%
\begin{Details}
This function is designed to be used upstream of
\code{mlsem\_modelfit\_paper()}.
It performs no table formatting and returns numeric values only.
\end{Details}
%
\begin{Value}
A data.frame containing model fit indices.
The structure depends on \code{type}.
\end{Value}
%
\begin{Examples}
\begin{ExampleCode}
## Not run: 
############################################
# Example 1. Typical usage after ML-SEM
############################################

fit_tab_long <- mlsem_fit_indices(fit_mlsem, type = "long")
fit_tab_data <- mlsem_fit_indices(fit_mlsem, type = "data")
fit_tab_wide <- mlsem_fit_indices(fit_mlsem, type = "wide")
fit_tab_apa  <- mlsem_fit_indices(fit_mlsem, type = "apa")

# Pipe-friendly usage
mlsem_fit_indices(fit_mlsem) %>%
  mlsem_modelfit_paper(type = "kable", format = "markdown")


############################################
# Example 2. Simple Multilevel SEM example
############################################

library(lavaan)

# Simulated multilevel data
set.seed(123)
dat <- data.frame(
  university = rep(1:25, each = 40),
  x1 = rnorm(1000),
  x2 = rnorm(1000),
  y1 = rnorm(1000),
  y2 = rnorm(1000)
)

# Multilevel SEM model
model_ml <- '
  level: 1
    X =~ x1 + x2
    Y =~ y1 + y2
    Y ~ X

  level: 2
    Y ~ X
'

fit_mlsem <- sem(
  model_ml,
  data      = dat,
  cluster   = "university",
  estimator = "MLR"
)

# Extract fit indices (APA long format)
fit_tab <- mlsem_fit_indices(fit_mlsem, type = "long")

# Convert to paper-ready table
mlsem_modelfit_paper(fit_tab, type = "kable", format = "markdown")

## End(Not run)
\end{ExampleCode}
\end{Examples}
\inputencoding{utf8}
\HeaderA{mlsem\_fornell\_larcker}{Fornell–Larcker Table with Significance (Multilevel SEM)}{mlsem.Rul.fornell.Rul.larcker}
%
\begin{Description}
Fornell–Larcker Table with Significance (Multilevel SEM)
\end{Description}
%
\begin{Usage}
\begin{verbatim}
mlsem_fornell_larcker(fit, level = c("within", "between"), digits = 3)
\end{verbatim}
\end{Usage}
%
\begin{Arguments}
\begin{ldescription}
\item[\code{fit}] lavaan multilevel SEM object

\item[\code{level}] "within" or "between"

\item[\code{digits}] numeric
\end{ldescription}
\end{Arguments}
\inputencoding{utf8}
\HeaderA{mlsem\_htmt}{HTMT for Multilevel SEM}{mlsem.Rul.htmt}
%
\begin{Description}
This function computes the Heterotrait–Monotrait ratio (HTMT)
for multilevel CFA/SEM models estimated using lavaan.
HTMT is calculated separately for within-level or between-level
measurement models based on observed-variable correlations.
\end{Description}
%
\begin{Usage}
\begin{verbatim}
mlsem_htmt(fit, level = c("within", "between"), digits = 3)
\end{verbatim}
\end{Usage}
%
\begin{Arguments}
\begin{ldescription}
\item[\code{fit}] A lavaan object fitted with a multilevel CFA/SEM model.

\item[\code{level}] Character. One of \code{"within"} or \code{"between"}.

\item[\code{digits}] Integer. Number of digits to round the HTMT values.
\end{ldescription}
\end{Arguments}
%
\begin{Details}
HTMT is defined as the ratio of the mean of heterotrait–heteromethod
correlations to the geometric mean of monotrait–heteromethod
correlations.

Only latent variables with at least two indicators are included.
The function uses level-specific observed-variable correlation
matrices extracted via \code{lavInspect(fit, "cor.ov")}.
\end{Details}
%
\begin{Value}
A data.frame containing the HTMT matrix
(lower triangle only; diagonal and upper triangle are NA).
\end{Value}
%
\begin{References}
Henseler, J., Ringle, C. M., \& Sarstedt, M. (2015).
A new criterion for assessing discriminant validity in
variance-based structural equation modeling.
Journal of the Academy of Marketing Science, 43(1), 115–135.
\end{References}
%
\begin{Examples}
\begin{ExampleCode}
## Not run: 
fit <- lavaan::sem(model, data = Demo.twolevel, cluster = "cluster")

mlsem_htmt(fit, level = "within")
mlsem_htmt(fit, level = "between")

## End(Not run)
\end{ExampleCode}
\end{Examples}
\inputencoding{utf8}
\HeaderA{mlsem\_modelfit\_paper}{Model Fit Table for Multilevel SEM (Paper-ready)}{mlsem.Rul.modelfit.Rul.paper}
%
\begin{Description}
This function formats multilevel SEM model fit indices extracted by
\code{mlsem\_fit\_indices()} into a publication-ready table.
It supports raw data output as well as markdown/html tables via kable.
\end{Description}
%
\begin{Usage}
\begin{verbatim}
mlsem_modelfit_paper(
  x,
  type = c("data", "kable"),
  format = c("raw", "markdown", "html"),
  caption = "Model Fit Indices (Multilevel SEM)"
)
\end{verbatim}
\end{Usage}
%
\begin{Arguments}
\begin{ldescription}
\item[\code{x}] A data.frame in APA long format with columns \code{Index} and \code{Value}.
Typically returned from \code{mlsem\_fit\_indices(fit, type = "long")}.

\item[\code{type}] Output type. One of \code{"data"} or \code{"kable"}.

\item[\code{format}] Output format when \code{type = "kable"}.
One of \code{"raw"}, \code{"markdown"}, or \code{"html"}.

\item[\code{caption}] Table caption used for kable output.
\end{ldescription}
\end{Arguments}
%
\begin{Details}
This function is designed to be used after \code{mlsem\_fit\_indices()}.
All internal formatting helpers are intentionally hidden (prefixed with \code{.})
to keep the public API minimal and stable.
\end{Details}
%
\begin{Value}
A data.frame (when \code{type = "data"} or \code{format = "raw"}),
or a formatted kable object (markdown/html).
\end{Value}
%
\begin{Examples}
\begin{ExampleCode}
## Not run: 
############################################
# Example 1. Typical workflow
############################################

fit_tab <- mlsem_fit_indices(fit_mlsem, type = "long")

mlsem_modelfit_paper(fit_tab)
mlsem_modelfit_paper(fit_tab, type = "data")
mlsem_modelfit_paper(fit_tab, type = "kable")
mlsem_modelfit_paper(fit_tab, type = "kable", format = "raw")
mlsem_modelfit_paper(fit_tab, type = "kable", format = "markdown")
mlsem_modelfit_paper(fit_tab, type = "kable", format = "html")

# Pipe-friendly usage
mlsem_fit_indices(fit_mlsem) %>%
  mlsem_modelfit_paper(type = "kable", format = "markdown")


############################################
# Example 2. Simple Multilevel SEM example
############################################

library(lavaan)

# Simulated data
set.seed(123)
dat <- data.frame(
  univ = rep(1:30, each = 40),
  x1 = rnorm(1200),
  x2 = rnorm(1200),
  y1 = rnorm(1200),
  y2 = rnorm(1200)
)

# Multilevel SEM model
model_ml <- '
  level: 1
    X =~ x1 + x2
    Y =~ y1 + y2
    Y ~ X

  level: 2
    Y ~ X
'

fit_mlsem <- sem(
  model_ml,
  data    = dat,
  cluster = "univ",
  estimator = "MLR"
)

# Extract fit indices
fit_tab <- mlsem_fit_indices(fit_mlsem, type = "long")

# Paper-ready table
mlsem_modelfit_paper(fit_tab, type = "kable", format = "markdown")

## End(Not run)
\end{ExampleCode}
\end{Examples}
\inputencoding{utf8}
\HeaderA{mlsem\_reliability}{Multilevel SEM Reliability (Cronbach's alpha, CR, AVE)}{mlsem.Rul.reliability}
%
\begin{Description}
This function computes reliability indices for multilevel SEM models
fitted with lavaan. Reliability is calculated separately for within-
and between-level measurement models.
\end{Description}
%
\begin{Usage}
\begin{verbatim}
mlsem_reliability(fit, level = c("within", "between"), digits = 3)
\end{verbatim}
\end{Usage}
%
\begin{Arguments}
\begin{ldescription}
\item[\code{fit}] A fitted lavaan multilevel SEM object.

\item[\code{level}] Character. One of \code{"within"} or \code{"between"}.

\item[\code{digits}] Integer. Number of digits to round the results.
\end{ldescription}
\end{Arguments}
%
\begin{Details}
The function returns Cronbach's alpha (model-based, using observed
correlation matrix), Composite Reliability (CR), and Average Variance
Extracted (AVE) based on standardized factor loadings.


\begin{itemize}

\item{} Cronbach's alpha is computed using the model-implied observed
correlation matrix (lavInspect(fit, "cor.ov")).
\item{} CR and AVE are computed from standardized factor loadings.
\item{} Latent variables with fewer than two indicators are excluded.

\end{itemize}

\end{Details}
%
\begin{Value}
A data.frame with latent variables and their reliability indices:
\begin{description}

\item[Latent] Latent variable name
\item[Cronbach] Cronbach's alpha
\item[CR] Composite Reliability
\item[AVE] Average Variance Extracted

\end{description}

\end{Value}
%
\begin{Examples}
\begin{ExampleCode}
## Not run: 
model <- '
level: 1
  fw =~ y1 + y2 + y3

level: 2
  fb =~ y1 + y2 + y3
'

fit <- lavaan::sem(
  model,
  data    = Demo.twolevel,
  cluster = "cluster"
)

mlsem_reliability(fit, level = "within")
mlsem_reliability(fit, level = "between", digits = 4)

## End(Not run)
\end{ExampleCode}
\end{Examples}
\inputencoding{utf8}
\HeaderA{modelfit\_paper}{Model Fit Table for Paper}{modelfit.Rul.paper}
%
\begin{Description}
This function formats model fit indices into a publication-ready table
or returns the underlying data.frame depending on the selected options.
It is designed to be used after extracting model fit information
from a CFA or SEM object (e.g., lavaan output).
\end{Description}
%
\begin{Usage}
\begin{verbatim}
modelfit_paper(
  x,
  type = c("kable", "data"),
  format = c("markdown", "raw", "html"),
  caption = "Model Fit Indices"
)
\end{verbatim}
\end{Usage}
%
\begin{Arguments}
\begin{ldescription}
\item[\code{x}] A data.frame or matrix containing model fit information.
The object must include rows named \code{"criterian"} and \code{"Value"},
and columns such as \code{chisq}, \code{df}, \code{pvalue},
\code{rmsea}, \code{srmr}, \code{cfi}, and \code{tli}.

\item[\code{type}] Character. Output type.
One of \code{"kable"} (default) or \code{"data"}.

\item[\code{format}] Character. Output format when \code{type = "kable"}.
One of \code{"markdown"}, \code{"raw"}, or \code{"html"}.

\item[\code{caption}] Character. Table caption used for kable output.
\end{ldescription}
\end{Arguments}
%
\begin{Details}
The function internally converts the degrees of freedom (df) value
to an integer to avoid outputs such as 97.000 in publication tables.
Formatting using \code{knitr::kable()} is applied only at the final
output stage, ensuring that raw data remain accessible for further
processing when needed.
\end{Details}
%
\begin{Value}
\begin{itemize}

\item{} If \code{type = "data"}: a data.frame.
\item{} If \code{type = "kable"} and \code{format = "raw"}: a data.frame.
\item{} If \code{type = "kable"} and \code{format = "markdown"} or
\code{"html"}: a formatted table object.

\end{itemize}

\end{Value}
%
\begin{Examples}
\begin{ExampleCode}
## Not run: 
# Assuming modelfitdata is returned from cfa2(type = "modelfit")

# Raw data.frame
modelfit_paper(modelfitdata, type = "data")

# Markdown table for papers
modelfit_paper(modelfitdata, type = "kable", format = "markdown")

# HTML table for Viewer or browser
modelfit_paper(modelfitdata, type = "kable", format = "html")

## End(Not run)
\end{ExampleCode}
\end{Examples}
\inputencoding{utf8}
\HeaderA{Modindices}{Modification Indices}{Modindices}
%
\begin{Description}
Given a fitted lavaan object, compute the modification indices (= univariate score tests) for a selected set of fixed-to-zero parameters.
\end{Description}
%
\begin{Usage}
\begin{verbatim}
Modindices(x, mi.level = 10, op1 = "~~")
\end{verbatim}
\end{Usage}
%
\begin{Arguments}
\begin{ldescription}
\item[\code{x}] lavaan result

\item[\code{mi.level}] mi over 10, default

\item[\code{op1}] regression '\textasciitilde{}\textasciitilde{}'
\end{ldescription}
\end{Arguments}
%
\begin{Value}
chisquare
\end{Value}
%
\begin{Examples}
\begin{ExampleCode}
## Not run: 
library(lavaan)
model1 <- '
  # measurement model
    ind60 =~ x1 + x2 + x3
    dem60 =~ y1 + y2 + y3 + y4
    dem65 =~ y5 + y6 + y7 + y8
  # regressions
    dem60 ~ ind60
    dem65 ~ ind60 + dem60
 # residual correlations
    y1 ~~ y5
    y2 ~~ y4 + y6
    y3 ~~ y7
    y4 ~~ y8
    y6 ~~ y8
'

fit1 <- sem(model1, data = PoliticalDemocracy)

 summary(fit1, standardized = TRUE)
cfa2(fit1)

Modindices(fit1)

## End(Not run)



\end{ExampleCode}
\end{Examples}
\inputencoding{utf8}
\HeaderA{ModSEM}{Title}{ModSEM}
%
\begin{Description}
Title
\end{Description}
%
\begin{Usage}
\begin{verbatim}
ModSEM(
  model = NULL,
  data = NULL,
  type = "res",
  ind = TRUE,
  cat = FALSE,
  interact = TRUE,
  auto = TRUE,
  n_name = 1,
  prefix = "a",
  paths_name = TRUE,
  add_model = NULL,
  ordered = NULL,
  sampling.weights = NULL,
  sample.cov = NULL,
  sample.mean = NULL,
  sample.th = NULL,
  sample.nobs = NULL,
  group = NULL,
  cluster = NULL,
  constraints = "",
  WLS.V = NULL,
  NACOV = NULL,
  ov.order = "model",
  ...
)
\end{verbatim}
\end{Usage}
%
\begin{Arguments}
\begin{ldescription}
\item[\code{model}] A description of the user-specified model. Typically, the model is
described using the lavaan model syntax. See model.syntax for more information.
Alternatively, a parameter table (eg. the output of the lavaanify() function)
is also accepted.

\item[\code{data}] An optional data frame containing the observed variables used in the model.
If some variables are declared as ordered factors,
lavaan will treat them as ordinal variables.

\item[\code{type}] res, fit, New\_model, New\_data, imlist

\item[\code{ind}] Automatically Adding Indirect Effects

\item[\code{cat}] output syntax default FALSE

\item[\code{interact}] interaction term extract possinle TRUE

\item[\code{auto}] auto genrate indirect path

\item[\code{n\_name}] indirect name number n\_name =1

\item[\code{prefix}] prefix "a" --"H" change possible

\item[\code{paths\_name}] inde path name default TRUE

\item[\code{add\_model}] add model syntax

\item[\code{ordered}] Character vector. Only used if the data is in a data.frame.
Treat these variables as ordered (ordinal) variables, if they are endogenous
in the model. Importantly, all other variables will be treated as numeric
(unless they are declared as ordered in the data.frame.) Since 0.6-4,
ordered can also be logical. If TRUE, all observed endogenous variables are
\#'    treated as ordered (ordinal). If FALSE, all observed endogenous variables
are considered to be numeric (again, unless they are declared as ordered
in the data.frame.)

\item[\code{sampling.weights}] A variable name in the data frame containing sampling
weight information. Currently only available for non-clustered data.
Depending on the sampling.weights.normalization option, these weights may be
rescaled (or not) so that their sum equals the number of observations
(total or per group).

\item[\code{sample.cov}] Numeric matrix. A sample variance-covariance matrix.
The rownames and/or colnames must contain the observed variable names.
For a multiple group analysis, a list with a variance-covariance matrix
for each group.

\item[\code{sample.mean}] A sample mean vector. For a multiple group analysis,
a list with a mean vector for each group.

\item[\code{sample.th}] Vector of sample-based thresholds.
For a multiple group analysis, a list with a vector of thresholds for each group.

\item[\code{sample.nobs}] Number of observations if the full data frame is missing and
only sample moments are given. For a multiple group analysis, a list or a vector
with the number of observations for each group.

\item[\code{group}] Character. A variable name in the data frame defining
the groups in a multiple group analysis.

\item[\code{cluster}] Character. A (single) variable name in the data frame defining
the clusters in a two-level dataset.

\item[\code{constraints}] Additional (in)equality constraints not yet included
in the model syntax. See model.syntax for more information.

\item[\code{WLS.V}] A user provided weight matrix to be used by estimator "WLS";
if the estimator is "DWLS", only the diagonal of this matrix will be used.
For a multiple group analysis, a list with a weight matrix for each group.
The elements of the weight matrix should be in the following order (if all data is
continuous): first the means (if a meanstructure is involved),
then the lower triangular elements of the covariance matrix including the
diagonal, ordered column by column. In the categorical case: first
the thresholds (including the means for continuous variables),
then the slopes (if any), the variances of continuous variables (if any),
and finally the lower triangular elements of the correlation/covariance
matrix excluding the diagonal, ordered column by column.

\item[\code{NACOV}] A user provided matrix containing the elements of (N times)
the asymptotic variance-covariance matrix of the sample statistics.
For a multiple group analysis, a list with an asymptotic variance-covariance
matrix for each group. See the WLS.V argument for information about
the order of the elements.

\item[\code{ov.order}] Character. If "model" (the default), the order of the observed
variable names (as reflected for example in the output of lavNames())
is determined by the model syntax. If "data", the order is determined
by the data (either the full data.frame or the sample (co)variance matrix).
If the WLS.V and/or NACOV matrices are provided,
this argument is currently set to "data".

\item[\code{...}] Many more additional options can be defined, using 'name = value'. See lavOptions for a complete list.
\end{ldescription}
\end{Arguments}
%
\begin{Value}
data
\end{Value}
%
\begin{Examples}
\begin{ExampleCode}

## Not run: 
#'
## The industrialization and Political Democracy Example
## Bollen (1989), page 332
model <- '
  # latent variable definitions
     ind60 =~ x1 + x2 + x3
     dem60 =~ y1 + a*y2 + b*y3 + c*y4
     dem65 =~ y5 + a*y6 + b*y7 + c*y8

  # regressions
    dem60 ~ ind60
    dem65 ~ ind60 + dem60

  # residual correlations
    y1 ~~ y5
    y2 ~~ y4 + y6
    y3 ~~ y7
    y4 ~~ y8
    y6 ~~ y8
'

fit <- sem(model, data = PoliticalDemocracy)
summary(fit, fit.measures = TRUE)
#'
#'

# # Example usage:
data(MarshWenHau)
str(MarshWenHau)

mwh = MarshWenHau

model2 <- "
f1 =~ x1 + x2 + x3
f1Xf2 =~ f1:f2
f2 =~ x4 + x5 + x6
f3 =~ y1 + y2 + y3

f3 ~ f1 + f2
f3 ~ f1Xf2

f1Xf2 ~~ 0*f1
f1Xf2 ~~ 0*f2
"

# interaction model
intsem =  SEM(model2, data= MarshWenHau)
summary(intsem)
intsem %>%lav_semPaths2(color=list(lat="darkred",man="gold"))


##2nd---------
#base model
lav_extract_mm(model2) %>%cat()
#interaction item list
lav_extract_imlist(model3)
#data
lav_latentProd(mwh, model3)%>%head()

#extracted interaction term
lav_extract_int(model2, lav_extract_imlist(model3))

# or interaction items ; imlist
imlist <- list(
  .f1Xf2 = c("f1", "f2"),
  .i_f1_f4 = c("f1", "f4")
)
# names(imlist)
lav_extract_int(model2, imlist)
#direct
lav_extract_int(model2, imlist= list(.f1Xf2 = c("f1", "f2"),
                                     .i_f1_f4 = c("f1", "f4")))
#estimates model
intsem =  SEM(lav_new_model(model2, lav_extract_imlist(model2)),
              data= lav_latentProd(mwh, model2))

#'

## End(Not run)


\end{ExampleCode}
\end{Examples}
\inputencoding{utf8}
\HeaderA{move\_col}{move\_col}{move.Rul.col}
%
\begin{Description}
move\_col
\end{Description}
%
\begin{Usage}
\begin{verbatim}
move_col(data, from, to, history = FALSE)
\end{verbatim}
\end{Usage}
%
\begin{Arguments}
\begin{ldescription}
\item[\code{data}] df

\item[\code{from}] from

\item[\code{to}] to

\item[\code{history}] histor show text
\end{ldescription}
\end{Arguments}
%
\begin{Value}
df
\end{Value}
%
\begin{Examples}
\begin{ExampleCode}
## Not run: 
#'
# 사용 예시
df <- data.frame(
  A = 1:5,
  B = 6:10,
  C = 11:15,
  D = 16:20
)

# 2번째 열을 4번째 위치로 이동
result_df <- move_col(df, from = 2, to = 4, history = TRUE)
print(result_df)


## End(Not run)

\end{ExampleCode}
\end{Examples}
\inputencoding{utf8}
\HeaderA{move\_cols}{move\_cols replication}{move.Rul.cols}
%
\begin{Description}
move\_cols replication
\end{Description}
%
\begin{Usage}
\begin{verbatim}
move_cols(data, ..., history = TRUE)
\end{verbatim}
\end{Usage}
%
\begin{Arguments}
\begin{ldescription}
\item[\code{data}] data frame

\item[\code{...}] from to, from to

\item[\code{history}] history show text
\end{ldescription}
\end{Arguments}
%
\begin{Value}
df
\end{Value}
%
\begin{Examples}
\begin{ExampleCode}
## Not run: 
#'
# 사용 예시
df <- data.frame(
  A = 1:5,
  B = 6:10,
  C = 11:15,
  D = 16:20
)

# 여러 열을 이동: 2번째 열을 4번째로, 3번째 열을 1번째로 이동
result_df <- move_cols(df, 2, 4, 3, 1, history = TRUE)
print(result_df)


## End(Not run)

\end{ExampleCode}
\end{Examples}
\inputencoding{utf8}
\HeaderA{move\_mat}{move\_mat matrix element coord position}{move.Rul.mat}
%
\begin{Description}
move\_mat matrix element coord position
\end{Description}
%
\begin{Usage}
\begin{verbatim}
move_mat(data, ..., imp = NA, replace = FALSE)
\end{verbatim}
\end{Usage}
%
\begin{Arguments}
\begin{ldescription}
\item[\code{data}] mat

\item[\code{...}] coord

\item[\code{imp}] NA

\item[\code{replace}] FALSE
\end{ldescription}
\end{Arguments}
%
\begin{Value}
mat
\end{Value}
%
\begin{Examples}
\begin{ExampleCode}

## Not run: 
#'
# example
data1 <- matrix(c(NA, NA, NA, NA, NA, NA, NA, NA, NA, "D_P3",
                  NA, NA, NA, NA, NA, NA, "진로태도", NA, NA, "D_P2",
                  NA, NA, NA, NA, NA, NA, NA, NA, NA, "D_P1",
                  "C_S5", NA, NA, NA, NA, NA, NA, NA, NA, "B14",
                  "C_S4", NA, NA, NA, NA, NA, NA, NA, NA, "B10",
                  "C_S3", NA, "자기효능감", NA, NA, NA, "진로준비", NA, NA, "B04",
                  "C_S2", NA, NA, NA, NA, NA, NA, NA, NA, "B02"))

# Test with replace = FALSE
move_mat(data, list(10, 10, 1, 2), list(9, 10, 1, 3), list(8, 10, 6,3), replace = FALSE)
move_mat(data, c(10, 10, 1, 2), c(9, 10, 1, 3), c(8, 10, 6,3), replace = FALSE)


# Test with replace = TRUE
move_mat(data, list(10, 10, 1, 2), list(9, 10, 1, 3), list(8, 10, 6,3), replace = TRUE)
move_mat(data, c(10, 10, 1, 2), c(9, 10, 1, 3), c(8, 10, 6,3), replace = TRUE)

## End(Not run)


\end{ExampleCode}
\end{Examples}
\inputencoding{utf8}
\HeaderA{move\_matpos}{move\_matpos is element change position}{move.Rul.matpos}
%
\begin{Description}
move\_matpos is element change position
\end{Description}
%
\begin{Usage}
\begin{verbatim}
move_matpos(data, ..., imp = NA, replace = FALSE)
\end{verbatim}
\end{Usage}
%
\begin{Arguments}
\begin{ldescription}
\item[\code{data}] mat

\item[\code{...}] position var

\item[\code{imp}] NA

\item[\code{replace}] FALSE
\end{ldescription}
\end{Arguments}
%
\begin{Value}
mat
\end{Value}
%
\begin{Examples}
\begin{ExampleCode}

## Not run: 
#'
# # Original matrix
layout_jut <- matrix(
  c(NA, NA, NA, NA, NA, NA, NA, NA, NA, 'D_P3',
    NA, NA, NA, NA, NA, NA, '진로태도', NA, NA, 'D_P2',
    NA, NA, NA, NA, NA, NA, NA, NA, NA, 'D_P1',
    'C_S5', NA, NA, NA, NA, NA, NA, NA, NA, 'B14',
    'C_S4', NA, NA, NA, NA, NA, NA, NA, NA, 'B10',
    'C_S3', NA, '자기효능감', NA, NA, NA, '진로준비', NA, NA, 'B04',
    'C_S2', NA, NA, NA, NA, NA, NA, NA, NA, 'B02',
    'C_S1', NA, NA, NA, NA, NA, NA, NA, NA, 'A_M3',
    NA, NA, NA, NA, NA, NA, NA, NA, NA, 'A_M2',
    NA, NA, NA, NA, NA, NA, '진로동기', NA, NA, 'A_M1'),
  nrow = 10, ncol = 10, byrow = TRUE)
layout_jut
# Apply the matpos function
layout_jut%>% move_matpos(
  list("진로태도", 3, 5),
  list("진로동기", 8, 5),
  list("B02", 1, 3),
  list("B04", 1, 5),
  list("B10", 1, 7),
  list("B14", 1, 9),
  list("D_P3", 4, 10),
  list("D_P2", 6, 10),
  list("D_P1", 8, 10),
  list("A_M1", 10, 3),
  list("A_M2", 10, 4),
  list("A_M3", 10, 5),
  replace = FALSE
)


matrix(
  c(NA, NA, NA, NA, NA, NA, NA, NA, NA, 'D_P3',
    NA, NA, NA, NA, NA, NA, '진로태도', NA, NA, 'D_P2',
    NA, NA, NA, NA, NA, NA, NA, NA, NA, 'D_P1',
    'C_S5', NA, NA, NA, NA, NA, NA, NA, NA, 'B14',
    'C_S4', NA, NA, NA, NA, NA, NA, NA, NA, 'B10',
    'C_S3', NA, '자기효능감', NA, NA, NA, '진로준비', NA, NA, 'B04',
    'C_S2', NA, NA, NA, NA, NA, NA, NA, NA, 'B02',
    'C_S1', NA, NA, NA, NA, NA, NA, NA, NA, 'A_M3',
    NA, NA, NA, NA, NA, NA, NA, NA, NA, 'A_M2',
    NA, NA, NA, NA, NA, NA, '진로동기', NA, NA, 'A_M1'),
  nrow = 10, ncol = 10, byrow = TRUE)%>% move_matpos(
    list("진로태도", 3, 5),
    list("진로동기", 8, 5),
    list("B02", 1, 3),
    list("B04", 1, 5),
    list("B10", 1, 7),
    list("B14", 1, 9),
    list("D_P3", 4, 10),
    list("D_P2", 6, 10),
    list("D_P1", 8, 10),
    list("A_M1", 10, 3),
    list("A_M2", 10, 4),
    list("A_M3", 10, 5),
    replace = FALSE
  )


## End(Not run)


\end{ExampleCode}
\end{Examples}
\inputencoding{utf8}
\HeaderA{move\_mat\_change}{Functions that replace a single component}{move.Rul.mat.Rul.change}
%
\begin{Description}
Functions that replace a single component
\end{Description}
%
\begin{Usage}
\begin{verbatim}
move_mat_change(mat, pattern, imp)
\end{verbatim}
\end{Usage}
%
\begin{Arguments}
\begin{ldescription}
\item[\code{mat}] matrix

\item[\code{pattern}] pattern

\item[\code{imp}] imputation
\end{ldescription}
\end{Arguments}
%
\begin{Value}
matrix
\end{Value}
%
\begin{Examples}
\begin{ExampleCode}

## Not run: 
#'
plspm_boot_factor_layout(pls_ts)%>%
  move_mat_cut(1,3)%>%
  move_mat_change("int","i_I")%>%make_mat_text()

layout_New = matrix(
  c( NA, 'IMA',  NA, NA, 'i_I',
     NA, 'QUA',NA,  'VAL', 'LOY',
     'EXP', NA, 'SAT',NA, NA),
  nrow = 3, ncol = 5, byrow = TRUE)

## End(Not run)


\end{ExampleCode}
\end{Examples}
\inputencoding{utf8}
\HeaderA{move\_mat\_cut}{move\_mat\_cut}{move.Rul.mat.Rul.cut}
%
\begin{Description}
move\_mat\_cut
\end{Description}
%
\begin{Usage}
\begin{verbatim}
move_mat_cut(mat, start = 1, n = 3)
\end{verbatim}
\end{Usage}
%
\begin{Arguments}
\begin{ldescription}
\item[\code{mat}] matrix

\item[\code{start}] 1

\item[\code{n}] width = 3
\end{ldescription}
\end{Arguments}
%
\begin{Value}
matrix
\end{Value}
%
\begin{Examples}
\begin{ExampleCode}

## Not run: 
#'
# 테스트용 행렬 데이터
mat <- matrix(c(NA, "진로동기", NA, "자기효능감", NA,
    "진로준비", NA, "진로태도", NA), nrow = 3, byrow = TRUE)

# 함수 실행
move_mat_cut(mat, start = 1, n = 3)


## End(Not run)


\end{ExampleCode}
\end{Examples}
\inputencoding{utf8}
\HeaderA{move\_mat\_nc}{move\_mat\_nc}{move.Rul.mat.Rul.nc}
%
\begin{Description}
move\_mat\_nc
\end{Description}
%
\begin{Usage}
\begin{verbatim}
move_mat_nc(mat, row, col, name)
\end{verbatim}
\end{Usage}
%
\begin{Arguments}
\begin{ldescription}
\item[\code{mat}] mat

\item[\code{row}] row

\item[\code{col}] cp;

\item[\code{name}] change name
\end{ldescription}
\end{Arguments}
%
\begin{Value}
mat
\end{Value}
\inputencoding{utf8}
\HeaderA{move\_mat\_ori}{move\_mat\_ori}{move.Rul.mat.Rul.ori}
%
\begin{Description}
move\_mat\_ori
\end{Description}
%
\begin{Usage}
\begin{verbatim}
move_mat_ori(mat, sort = TRUE, type = "res")
\end{verbatim}
\end{Usage}
%
\begin{Arguments}
\begin{ldescription}
\item[\code{mat}] mat

\item[\code{sort}] sort =TRUE

\item[\code{type}] res, coord
\end{ldescription}
\end{Arguments}
%
\begin{Value}
coordnate
\end{Value}
%
\begin{Examples}
\begin{ExampleCode}

## Not run: 
#'
layout_jut <- matrix(
  c(NA, NA, NA, NA, NA, NA, NA, NA, NA, 'D_P3',
    NA, NA, NA, NA, NA, NA, '진로태도', NA, NA, 'D_P2',
    NA, NA, NA, NA, NA, NA, NA, NA, NA, 'D_P1',
    'C_S5', NA, NA, NA, NA, NA, NA, NA, NA, 'B14',
    'C_S4', NA, NA, NA, NA, NA, NA, NA, NA, 'B10',
    'C_S3', NA, '자기효능감', NA, NA, NA, '진로준비', NA, NA, 'B04',
    'C_S2', NA, NA, NA, NA, NA, NA, NA, NA, 'B02',
    'C_S1', NA, NA, NA, NA, NA, NA, NA, NA, 'A_M3',
    NA, NA, NA, NA, NA, NA, NA, NA, NA, 'A_M2',
    NA, NA, NA, NA, NA, NA, '진로동기', NA, NA, 'A_M1'),
  nrow = 10, ncol = 10, byrow = TRUE)
move_mat_ori(layout_jut)

#coord

move_mat_ori(layout_jut)%>%
mutate(Mat = paste0(element," = c(", row_col,")" ))%>%
pull(Mat)%>%
  paste(collapse=", ")%>%
  cat("\n")

##treatment
 move_mat_ori(layout_jut, type="coord")
#'
plspm_make_layout_mat(
  A_M1 = c(10, 10), A_M2 = c(9, 10),
  A_M3 = c(8, 10), B02 = c(7, 10),
  B04 = c(6, 10), B10 = c(5, 10),
   B14 = c(4, 10), C_S1 = c(8, 1),
   C_S2 = c(7, 1), C_S3 = c(6, 1),
   C_S4 = c(5, 1), C_S5 = c(4, 1),
   D_P1 = c(3, 10), D_P2 = c(2, 10),
    D_P3 = c(1, 10), 자기효능감 = c(6, 3),
    진로동기 = c(10, 7), 진로준비 = c(6, 7),
     진로태도 = c(2, 7)
)


## End(Not run)


\end{ExampleCode}
\end{Examples}
\inputencoding{utf8}
\HeaderA{move\_mat\_substr}{move\_mat\_cut}{move.Rul.mat.Rul.substr}
%
\begin{Description}
move\_mat\_cut
\end{Description}
%
\begin{Usage}
\begin{verbatim}
move_mat_substr(mat, start = 1, n = 3)
\end{verbatim}
\end{Usage}
%
\begin{Arguments}
\begin{ldescription}
\item[\code{mat}] matrix

\item[\code{start}] 1

\item[\code{n}] width = 3
\end{ldescription}
\end{Arguments}
%
\begin{Value}
matrix
\end{Value}
%
\begin{Examples}
\begin{ExampleCode}

## Not run: 
#'
# 테스트용 행렬 데이터
mat <- matrix(c(NA, "진로동기", NA, "자기효능감", NA,
    "진로준비", NA, "진로태도", NA), nrow = 3, byrow = TRUE)

# 함수 실행
move_mat_cut(mat, start = 1, n = 3)


## End(Not run)


\end{ExampleCode}
\end{Examples}
\inputencoding{utf8}
\HeaderA{move\_row}{Functions to move rows to a desired position}{move.Rul.row}
%
\begin{Description}
Functions to move rows to a desired position
\end{Description}
%
\begin{Usage}
\begin{verbatim}
move_row(data, from, to, history = FALSE)
\end{verbatim}
\end{Usage}
%
\begin{Arguments}
\begin{ldescription}
\item[\code{data}] data.fram

\item[\code{from}] row from

\item[\code{to}] row to

\item[\code{history}] history
\end{ldescription}
\end{Arguments}
%
\begin{Value}
data
\end{Value}
%
\begin{Examples}
\begin{ExampleCode}

## Not run: 
#' #'
#' # Example data
car = data.frame(
  rownames = c(paste0("Row_",1:6)),
  mpg = c('21', '21', '22.8', '21.4', '18.7', '18.1'),
  cyl = c('6', '6', '4', '6', '8', '6'),
  disp = c('160', '160', '108', '258', '360', '225')
)

car
car%>%move_row(1,6)
car%>%move_row(6,1)
car%>%move_row(6,2)
car%>%move_row(6,3)
car%>%move_row(1,3)
car%>%move_row(6,5)

car
car%>%
  move_row(6,1)%>%
  move_row(1,3)%>%
  move_row(3,2)
#'

## End(Not run)


\end{ExampleCode}
\end{Examples}
\inputencoding{utf8}
\HeaderA{move\_rows}{move\_row rep function}{move.Rul.rows}
%
\begin{Description}
move\_row rep function
\end{Description}
%
\begin{Usage}
\begin{verbatim}
move_rows(data, ..., history = TRUE)
\end{verbatim}
\end{Usage}
%
\begin{Arguments}
\begin{ldescription}
\item[\code{data}] data

\item[\code{...}] movelist c(a,b),c(c,d)...
\end{ldescription}
\end{Arguments}
%
\begin{Value}
data
\end{Value}
%
\begin{Examples}
\begin{ExampleCode}

## Not run: 
#'
data <- data.frame(
  source = c("A05", "A08", "A04", "A06",
  "A09", "A11", "A13", "A02", "A03", "A01",
  "A14", "A15", "A07", "A12", "A10",
  "eigen_value", "Proportion_Var", "Cumulative_Var"),
  MR1 = c("", "", "", "", "", "", "",
  "0.362", "0.449", "0.457", "0.47",
  "0.538", "0.62", "0.792", "0.96",
  "3.363", "0.224", "0.224"),
  MR2 = c("", "0.382", "", "", "0.445",
  "0.724", "0.906", "0.392", "", "",
  "0.487", "", "", "", "", "2.095",
   "0.14", "0.364"),
  MR3 = c("", "", "0.402", "0.963", "", "",
   "", "", "", "", "", "", "", "",
    "", "1.691", "0.113", "0.477")
)
data
#'
move_row(data, 14,1)
move_row(data, 14,1)%>% move_row(1,2)
move_row(data, 14,1)%>% move_row(1,2)%>%
 move_row(2,3)
#'
data %>% move_rows(c(14, 1), c(1, 2))
data %>% move_rows(c(14, 1), c(1, 2), c(2, 3))
#'
car = data.frame(
  rownames = c(paste0("Row_",1:6)),
  mpg = c('21', '21', '22.8', '21.4', '18.7', '18.1'),
  cyl = c('6', '6', '4', '6', '8', '6'),
  disp = c('160', '160', '108', '258', '360', '225'))
cars

#repeated
cars
car %>%
  move_row(6,1)%>%
  move_row(1,3)%>%
  move_row(3,2)

#history
car %>%
  move_row(6,1, history=TRUE)%>%
  move_row(1,3, history=TRUE)%>%
  move_row(3,2, history=TRUE)%>%
  move_row(2,6, history=TRUE)

car%>%move_rows(6,1) #errorcheck
car%>%move_rows(c(6,1), c(1,2),c(2,6))
car%>%move_rows(c(6,1), c(1,2),c(2,6), history=T)

# simple input method
move_rows(car, 6, 1, 1, 2, 2, 6, history=TRUE)
car%>% move_rows( 6, 1, 1, 2, 2, 6, 6,3, history=TRUE)
# Testing the function with direct pair inputs
move_rows(car, 6, 1, 1, 2, 2, 6, history=TRUE)

# Testing the function with list of pairs inputs
move_rows(car, c(6, 1), c(1, 2), c(2, 6), history=TRUE)


# Testing the function with direct pair inputs
move_rows(car, 6, 1,  history=FALSE)
move_rows(car, 6, 1,1,2)
move_rows(car, 6, 1, 1, 2, 2, 6, history=TRUE)
car%>%move_rows(6, 1, 1, 2, 2, 6, history=TRUE)
move_rows(car, 6, 1, 1, 2, 2, 6, 6,3, history=TRUE)
car%>%move_rows(c(6,1), c(1,2),c(2,6),c(6,5), history=T)
car%>%move_rows(6, 1, 1, 2, 2, 6, 6,3, history=T)


# Testing the function with list of pairs inputs
move_rows(car, c(6, 1), history=TRUE)
move_rows(car, c(6, 1),c(1,2), history=TRUE)
move_rows(car, c(6, 1), c(1, 2), c(2, 6),c(6,3), history=TRUE)
#'
#'




## End(Not run)

\end{ExampleCode}
\end{Examples}
\inputencoding{utf8}
\HeaderA{mplus\_est\_df}{mplus\_est\_df}{mplus.Rul.est.Rul.df}
%
\begin{Description}
mplus\_est\_df
\end{Description}
%
\begin{Usage}
\begin{verbatim}
mplus_est_df(text, digits = 3)
\end{verbatim}
\end{Usage}
%
\begin{Arguments}
\begin{ldescription}
\item[\code{text}] text est

\item[\code{digits}] 3
\end{ldescription}
\end{Arguments}
%
\begin{Value}
data.frame
\end{Value}
%
\begin{Examples}
\begin{ExampleCode}


## Not run: 
#'
esttest <- "MODEL RESULTS

                                                    Two-Tailed
                    Estimate       S.E.  Est./S.E.    P-Value

 F1       BY
    X1                 1.000      0.000    999.000    999.000
    X2                 0.771      0.034     22.380      0.000
    X3                 0.683      0.033     20.510      0.000

 F2       BY
    X4                 1.000      0.000    999.000    999.000
    X5                 0.722      0.032     22.400      0.000
    X6                 0.702      0.033     21.320      0.000

 F3       BY
    Y1                 1.000      0.000    999.000    999.000
    Y2                 0.705      0.029     24.023      0.000
    Y3                 0.695      0.029     24.243      0.000

 F2       WITH
    F1                -0.482      0.082     -5.912      0.000

 F3       WITH
    F1                 0.522      0.084      6.203      0.000
    F2                -0.558      0.088     -6.315      0.000

 Variances
    F1                 1.347      0.121     11.095      0.000
    F2                 1.492      0.131     11.382      0.000
    F3                 1.599      0.135     11.868      0.000

 Residual Variances
    X1                 0.364      0.046      7.916      0.000
    X2                 0.245      0.029      8.583      0.000
    X3                 0.320      0.029     11.160      0.000
    X4                 0.347      0.048      7.224      0.000
    X5                 0.279      0.029      9.579      0.000
    X6                 0.342      0.031     10.865      0.000
    Y1                 0.302      0.044      6.931      0.000
    Y2                 0.290      0.028     10.444      0.000
    Y3                 0.269      0.026     10.230      0.000"

mplus_est_df(esttest)
#'

## End(Not run)

\end{ExampleCode}
\end{Examples}
\inputencoding{utf8}
\HeaderA{mplus\_fit\_df}{mplus\_fit\_df: model fit from  Mplus result}{mplus.Rul.fit.Rul.df}
%
\begin{Description}
mplus\_fit\_df: model fit from  Mplus result
\end{Description}
%
\begin{Usage}
\begin{verbatim}
mplus_fit_df(text, digits = 3)
\end{verbatim}
\end{Usage}
%
\begin{Arguments}
\begin{ldescription}
\item[\code{text}] mplut copy data

\item[\code{digits}] 3
\end{ldescription}
\end{Arguments}
%
\begin{Value}
fit data.frame
\end{Value}
%
\begin{Examples}
\begin{ExampleCode}


## Not run: 


text= c("MODEL FIT INFORMATION

Number of Free Parameters                       21

Loglikelihood

          H0 Value                       -4568.339
          H1 Value                       -4551.194

Information Criteria

          Akaike (AIC)                    9178.677
          Bayesian (BIC)                  9263.623
          Sample-Size Adjusted BIC        9196.983
            (n* = (n + 2) / 24)

Chi-Square Test of Model Fit

          Value                             34.290
          Degrees of Freedom                    24
          P-Value                           0.0796

RMSEA (Root Mean Square Error Of Approximation)

          Estimate                           0.032
          90 Percent C.I.                    0.000  0.054
          Probability RMSEA <= .05           0.901

CFI/TLI

          CFI                                0.996
          TLI                                0.994

Chi-Square Test of Model Fit for the Baseline Model

          Value                           2533.266
          Degrees of Freedom                    36
          P-Value                           0.0000

SRMR (Standardized Root Mean Square Residual)

          Value                              0.022")

mplus_fit_df(text)
#' 
## End(Not run)
\end{ExampleCode}
\end{Examples}
\inputencoding{utf8}
\HeaderA{mr}{Functions to move rows to a desired position}{mr}
%
\begin{Description}
Functions to move rows to a desired position
\end{Description}
%
\begin{Usage}
\begin{verbatim}
mr(data, from, to, history = FALSE)
\end{verbatim}
\end{Usage}
%
\begin{Arguments}
\begin{ldescription}
\item[\code{data}] data.fram

\item[\code{from}] row from

\item[\code{to}] row to

\item[\code{history}] history
\end{ldescription}
\end{Arguments}
%
\begin{Value}
data
\end{Value}
%
\begin{Examples}
\begin{ExampleCode}

## Not run: 
#'
#' # Example data
car = data.frame(
  rownames = c(paste0("Row_",1:6)),
  mpg = c('21', '21', '22.8', '21.4', '18.7', '18.1'),
  cyl = c('6', '6', '4', '6', '8', '6'),
  disp = c('160', '160', '108', '258', '360', '225')
)
car
car %>% mr(1,6)
car %>% mr(6,1)
car %>% mr(6,2)
car %>% mr(6,3)
car %>% mr(1,3)
car %>% mr(6,5)

#history
car %>%
  mr(6,1)%>%
  mr(1,3)%>%
  mr(3,2)
#'

## End(Not run)


\end{ExampleCode}
\end{Examples}
\inputencoding{utf8}
\HeaderA{mrs}{move\_row rep function}{mrs}
%
\begin{Description}
move\_row rep function
\end{Description}
%
\begin{Usage}
\begin{verbatim}
mrs(data, ..., history = TRUE)
\end{verbatim}
\end{Usage}
%
\begin{Arguments}
\begin{ldescription}
\item[\code{data}] data

\item[\code{...}] movelist c(a,b),c(c,d)...
\end{ldescription}
\end{Arguments}
%
\begin{Value}
data
\end{Value}
%
\begin{Examples}
\begin{ExampleCode}

## Not run: 
#'
data <- data.frame(
  source = c("A05", "A08", "A04",
  "A06", "A09", "A11", "A13", "A02",
  "A03", "A01", "A14", "A15", "A07",
  "A12", "A10", "eigen_value", "Proportion_Var", "Cumulative_Var"),
  MR1 = c("", "", "", "", "", "", "",
  "0.362", "0.449", "0.457", "0.47",
  "0.538", "0.62", "0.792", "0.96",
  "3.363", "0.224", "0.224"),
  MR2 = c("", "0.382", "", "", "0.445",
  "0.724", "0.906", "0.392", "", "",
  "0.487", "", "", "", "", "2.095", "0.14", "0.364"),
  MR3 = c("", "", "0.402", "0.963", "",
   "", "", "", "", "", "", "", "",
    "", "", "1.691", "0.113", "0.477")
)
data
#'
move_row(data, 14,1)
move_row(data, 14,1)%>% move_row(1,2)
move_row(data, 14,1)%>% move_row(1,2)%>% move_row(2,3)
#'
data %>% move_rows(c(14, 1), c(1, 2))
data %>% move_rows(c(14, 1), c(1, 2), c(2, 3))
#'
car = data.frame(
  rownames = c(paste0("Row_",1:6)),
  mpg = c('21', '21', '22.8', '21.4', '18.7', '18.1'),
  cyl = c('6', '6', '4', '6', '8', '6'),
  disp = c('160', '160', '108', '258', '360', '225'))
cars

#repeated
cars
car %>%
  move_row(6,1)%>%
  move_row(1,3)%>%
  move_row(3,2)

#history
car %>%
  move_row(6,1, history=TRUE)%>%
  move_row(1,3, history=TRUE)%>%
  move_row(3,2, history=TRUE)%>%
  move_row(2,6, history=TRUE)

car%>%move_rows(6,1) #errorcheck
car%>%move_rows(c(6,1), c(1,2),c(2,6))
car%>%move_rows(c(6,1), c(1,2),c(2,6), history=T)

# simple input method
move_rows(car, 6, 1, 1, 2, 2, 6, history=TRUE)
car%>% move_rows( 6, 1, 1, 2, 2, 6, 6,3, history=TRUE)
# Testing the function with direct pair inputs
move_rows(car, 6, 1, 1, 2, 2, 6, history=TRUE)

# Testing the function with list of pairs inputs
move_rows(car, c(6, 1), c(1, 2), c(2, 6), history=TRUE)


# Testing the function with direct pair inputs
move_rows2(car, 6, 1,  history=FALSE)
move_rows2(car, 6, 1,1,2)
move_rows2(car, 6, 1, 1, 2, 2, 6, history=TRUE)
car%>%move_rows2(6, 1, 1, 2, 2, 6, history=TRUE)
move_rows2(car, 6, 1, 1, 2, 2, 6, 6,3, history=TRUE)
car%>%move_rows2(c(6,1), c(1,2),c(2,6),c(6,5), history=T)
car%>%move_rows2(6, 1, 1, 2, 2, 6, 6,3, history=T)


# Testing the function with list of pairs inputs
move_rows2(car, c(6, 1), history=TRUE)
move_rows(car, c(6, 1),c(1,2), history=TRUE)
move_rows2(car, c(6, 1), c(1, 2), c(2, 6),c(6,3), history=TRUE)
#'
#'




## End(Not run)

\end{ExampleCode}
\end{Examples}
\inputencoding{utf8}
\HeaderA{msdn\_contrast}{Contrast test with mean, standard deviation, and number of samples}{msdn.Rul.contrast}
%
\begin{Description}
Contrast test with mean, standard deviation, and number of samples
\end{Description}
%
\begin{Usage}
\begin{verbatim}
msdn_contrast(M, sd, n, ..., digits = 3, type = "res", meanadd = TRUE)
\end{verbatim}
\end{Usage}
%
\begin{Arguments}
\begin{ldescription}
\item[\code{M}] means M = c(26.66, 19.74, 15.10)

\item[\code{sd}] sd  sd = c(4.51, 1.45, 2.56)

\item[\code{n}] sample  n = c(11, 7, 14)

\item[\code{...}] contrast c(1, -1,  0), c(1,  0, -1), c(0,  1, -1) ....

\item[\code{digits}] default 3

\item[\code{type}] res, res\_ci

\item[\code{meanadd}] include means
\end{ldescription}
\end{Arguments}
%
\begin{Value}
conrast result
\end{Value}
%
\begin{Examples}
\begin{ExampleCode}
## Not run: 


msdn_contrast(M = c(26.66, 19.74, 15.10),
              sd = c(4.51, 1.45, 2.56),
              n = c(11, 7, 14),
              c(1, -1, 0) )

contrast_msdn(M = c(26.66, 19.74, 15.10),
              sd = c(4.51, 1.45, 2.56),
              n = c(11, 7, 14),
              c(1, -1, 0) ,
              meanadd = FALSE)

msdn_contrast(M = c(26.66, 19.74, 15.10),
              sd = c(4.51, 1.45, 2.56),
              n = c(11, 7, 14),
              c(1, -1, 0) ,
             meanadd = FALSE,
              type="res_ci")


msdn_contrast(M = c(26.66, 19.74, 15.10),
              sd = c(4.51, 1.45, 2.56),
             n = c(11, 7, 14),
              c(1, -1,  0),
              c(1,  0, -1),
              c(0,  1, -1)
)
msdn_contrast(M = c(26.66, 19.74, 15.10),
              sd = c(4.51, 1.45, 2.56),
              n = c(11, 7, 14),
              contr1 = c(1, -1,  0),
              contr2= c(1,  0, -1),
              contr3= c(0,  1, -1),
              meanadd=F
)
msdn_contrast(M = c(26.66, 19.74, 15.10),
              sd = c(4.51, 1.45, 2.56),
              n = c(11, 7, 14),
             contr1 = c(1, -1,  0),
             contr2= c(1,  0, -1),
              contr3= c(0,  1, -1),
              meanadd=FALSE, type="res_ci"
)
#error
msdn_contrast(M = c(26.66, 19.74, 15.10),
              sd = c(4.51, 1.45, 2.56),
             n = c(11, 7, 14),
              contr = list(c(1, -1,  0),
                            c(1,  0, -1),
                           c(0,  1, -1)
              ))



## End(Not run)
\end{ExampleCode}
\end{Examples}
\inputencoding{utf8}
\HeaderA{msdn\_ttest}{Using ttest with mean, sd, and n}{msdn.Rul.ttest}
%
\begin{Description}
Using ttest with mean, sd, and n
\end{Description}
%
\begin{Usage}
\begin{verbatim}
msdn_ttest(
  n1,
  m1,
  sd1,
  n2,
  m2,
  sd2,
  var1 = NULL,
  var2 = NULL,
  var.equal = FALSE,
  digits = 2
)
\end{verbatim}
\end{Usage}
%
\begin{Arguments}
\begin{ldescription}
\item[\code{n1}] sample1

\item[\code{m1}] mean1

\item[\code{sd1}] sd1

\item[\code{n2}] n2

\item[\code{m2}] m2

\item[\code{sd2}] sd2

\item[\code{var1}] var1 name

\item[\code{var2}] var2 name

\item[\code{var.equal}] TRUE Student FALSE, whech

\item[\code{digits}] 3
\end{ldescription}
\end{Arguments}
%
\begin{Value}
ttest result
\end{Value}
%
\begin{Examples}
\begin{ExampleCode}

## Not run: 
#'
# # Example usage
msdn_ttest(n1 = 603, m1 = 59.6, sd1 = 0.6,
           n2 = 10785, m2 = 46.7, sd2 = 0.3)

#  mysummaryBy(mpg ~ vs, data = mtcars)

# library(jjstat)
msdn_ttest(n1 = 18, m1 = 16.617, sd1 = 3.8607,
           n2 = 14, m2 = 24.557, sd2 = 5.3790,  var.equal=FALSE)
t.test(mpg ~ vs, data = mtcars, var.equal=F)%>%tidy()%>%dall

msdn_ttest(n1 = 18, m1 = 16.617, sd1 = 3.8607,
           n2 = 14, m2 = 24.557, sd2 = 5.3790, var.equal=TRUE)
t.test(mpg ~ vs, data = mtcars, var.equal=T)%>%tidy()

#'
msdn_ttest(n1 = 18, m1 = 16.617, sd1 = 3.8607, var1 = "vs0",
           n2 = 14, m2 = 24.557, sd2 = 5.3790, var2 = "vs1", var.equal = FALSE)
msdn_ttest(n1 = 18, m1 = 16.617, sd1 = 3.8607, var1 = "vs1",
           n2 = 14, m2 = 24.557, sd2 = 5.3790, var2 = "vs2", var.equal = FALSE)



## End(Not run)

\end{ExampleCode}
\end{Examples}
\inputencoding{utf8}
\HeaderA{mutate\_col}{Create, modify, and delete columns col position}{mutate.Rul.col}
%
\begin{Description}
Create, modify, and delete columns col position
\end{Description}
%
\begin{Usage}
\begin{verbatim}
mutate_col(data, ..., col = ncol(data))
\end{verbatim}
\end{Usage}
%
\begin{Arguments}
\begin{ldescription}
\item[\code{data}] data.frame

\item[\code{...}] creat name

\item[\code{col}] position
\end{ldescription}
\end{Arguments}
%
\begin{Value}
data
\end{Value}
%
\begin{Examples}
\begin{ExampleCode}

## Not run: 
#'
fm <- data.frame(
  age = c("20대", "20대", "20대", "30대",
  "30대", "30대", "40대", "40대", "40대"),
  manufacturer = c("google", "apple", "samsung",
  "google", "apple", "samsung", "google",
   "apple", "samsung"),
  values = c(794, 804, 1243, 987, 1055,
  2517, 1528, 1519, 2255)
)
fm

#data
fm %>% mutate_col(gender= rep("a",9))
fm %>% mutate_col(gender= rep("a",9), col = 2)

#position
fm %>% mutate_col(gender = rep("female", nrow(fm)), col = 2)
fm %>% mutate_col(gender = rep("female", nrow(fm)), col = 1)
fm %>% mutate_col(gender = rep("female", nrow(fm)), col = 3)
fm %>% mutate_col(gender = rep("female", nrow(fm)), col = 4)

#Build and adjust with data from within
fm %>% mutate(gender = age)
fm %>% mutate_col(age_int = substring(fm$age,1,2))
fm %>% mutate_col(gender = paste(fm$age,fm$values,sep="//"))

fm %>% mutate_col(gender = paste(.$age,.$values,sep="//"), col=2)
# Double data generation
fm %>%
  mutate_col(insert_2 = rep("female", nrow(fm)), col = 2) %>%
  mutate_col( insert_4 = rep("a",9), col = 4)


# second data
data2 <- data.frame(
  weight = c(4.17, 5.58, 5.18, 6.11, 4.50, 4.61, 5.17, 4.53, 5.33, 5.14,
             4.81, 4.17, 4.41, 3.59, 5.87, 3.83, 6.03, 4.89, 4.32, 4.69,
             6.31, 5.12, 5.54, 5.50, 5.37, 5.29, 4.92, 6.15, 5.80, 5.26),
  ids = c(paste0("C", 1:10), paste0("T",1:10), paste0("S", 1:10))
)
data2 %>%mutate_col(group= .$weight, col=2)


## End(Not run)

\end{ExampleCode}
\end{Examples}
\inputencoding{utf8}
\HeaderA{mydes}{mydes 'n','mean','sd','min','max'}{mydes}
%
\begin{Description}
mydes 'n','mean','sd','min','max'
\end{Description}
%
\begin{Usage}
\begin{verbatim}
mydes(myvariable, var = NULL, digits = 2)
\end{verbatim}
\end{Usage}
%
\begin{Arguments}
\begin{ldescription}
\item[\code{myvariable}] variable vector

\item[\code{var}] input variable name

\item[\code{digits}] digits
\end{ldescription}
\end{Arguments}
%
\begin{Examples}
\begin{ExampleCode}
## Not run: 
## view variable
mydes(mtcars$mpg,"mpg")
##not see variable
mydes(mtcars$mpg)


## End(Not run)

\end{ExampleCode}
\end{Examples}
\inputencoding{utf8}
\HeaderA{mysummary}{my summary descriptive statistics}{mysummary}
%
\begin{Description}
my summary descriptive statistics
\end{Description}
%
\begin{Usage}
\begin{verbatim}
mysummary(myobject, ..., digits = 5, msdn = FALSE, na.rm = FALSE)
\end{verbatim}
\end{Usage}
%
\begin{Arguments}
\begin{ldescription}
\item[\code{myobject}] data.frame, matrix

\item[\code{...}] column variable

\item[\code{digits}] round default 3

\item[\code{msdn}] TRUE only output var, N, mean. sd

\item[\code{na.rm}] FALSE, TRUE is remove NA
\end{ldescription}
\end{Arguments}
%
\begin{Value}
size   MEAN    SD  MIN  MAX
\end{Value}
%
\begin{Examples}
\begin{ExampleCode}
## Not run: 
mysummary(mtcars, "mpg","wt","hp")
## size      MEAN         SD    MIN     MAX
## 1   32  20.09062  6.0269481 10.400  33.900
## 2   32   3.21725  0.9784574  1.513   5.424
## 3   32 146.68750 68.5628685 52.000 335.000

## all variable
 jjstat::mysummary(mtcars, colnames(mtcars))

mysummary(mtcars, colnames(mtcars))

mysummary(mtcars)

mysummary(mtcars, all=FALSE)


## End(Not run)
\end{ExampleCode}
\end{Examples}
\inputencoding{utf8}
\HeaderA{mysummaryBy}{statistics by group}{mysummaryBy}
%
\begin{Description}
statistics by group
\end{Description}
%
\begin{Usage}
\begin{verbatim}
mysummaryBy(
  data,
  formula,
  add_var = NULL,
  stat = FALSE,
  agg = FALSE,
  gm = FALSE,
  fun = NULL,
  msdn = FALSE,
  t = FALSE
)
\end{verbatim}
\end{Usage}
%
\begin{Arguments}
\begin{ldescription}
\item[\code{data}] data.frame

\item[\code{formula}] general formula

\item[\code{add\_var}] one sample name, Used to obtain statistics for one variable

\item[\code{stat}] TRUE t.test and aov(), thes stat ="t.test", or 'aov'

\item[\code{agg}] TRUE, mnay to many variable is showed. This functuion tha makes aggregate() work full

\item[\code{gm}] default FALSE, TRUE Use group data obtained from group descriptive statistics to obtain group-level descriptive statistics (means, standard deviations, etc.) and to analyze a mixed model

\item[\code{msdn}] msdn=TRUE is output mean, sd, N

\item[\code{t}] transpose  when agg=TRUE
\end{ldescription}
\end{Arguments}
%
\begin{Value}
Mea, SD, N, min, max, skew, kurt
\end{Value}
%
\begin{Examples}
\begin{ExampleCode}
## Not run: 
## Used to obtain statistics for one variable
mysummaryBy(mpg ~ 1, data = mtcars)
mysummaryBy(mpg ~ 1, data = mtcars, "mpg")
mysummaryBy(mpg ~ 1, data = mtcars, add_var = "mpg")

##The group variable is 2 or greater
mysummaryBy(mpg ~ vs, data = mtcars)
mysummaryBy(mpg ~ vs+am, data = mtcars)
mysummaryBy(mpg ~ vs+am+cyl, data = mtcars)

##statistic data output
mysummaryBy(mpg ~ 1, data = mtcars, "mpg", stat="t.test")
mysummaryBy(mpg ~ vs, data = mtcars, stat="t.test")
mysummaryBy(mpg ~ vs+am, data = mtcars, stat="aov")
mysummaryBy(mpg ~ vs*cyl, data = mtcars, stat="aov")

##new group mean by group
#' group_by(mtcars, am) %>%
 summarise(mpg = mean(mpg)) %>%
  summarise(mean(mpg),sd(mpg),length(mpg), min(mpg),max(mpg))

### A tibble: 1 × 5
###`mean(mpg)` `sd(mpg)` `length(mpg)` `min(mpg)` `max(mpg)`
###<dbl>     <dbl>         <int>      <dbl>      <dbl>
###  1        20.8      5.12             2       17.1       24.4

mysummaryBy(mpg ~ am , data= mtcars) %>%
    rename(mpg = Mean) %>%   #change name
      mysummaryBy(formula = mpg ~ 1)  # mean of mpg

mysummaryBy(mpg ~ am , data= mtcars) %>%
    rename(mpg = Mean) %>%   #change name
      mysummary("mpg")  # mean of mpg

#multi group transpose table
mysummaryBy(mpg ~ vs+am, data = mtcars, agg=TRUE)

## transpose
mysummaryBy(mpg ~ vs+am, data = mtcars, agg=TRUE, t=TRUE)

##twoway
mysummaryBy(mpg ~ vs+am, data = mtcars, agg=TRUE)


## twoway interaction
mysummaryBy(mpg ~ vs*am, data = mtcars, agg=TRUE, stat="aov")

#add variable
mysummaryBy(mpg ~ vs, data = mtcars, gm=T)
# A tibble: 1 × 9
# grp   dv        N  MEAN    SD   MIN   MAX  Skew  Kurt
# <chr> <chr> <int> <dbl> <dbl> <dbl> <dbl> <dbl> <dbl>
#  vs    mpg       2  20.6  5.61  16.6  24.6     0 -2.75


## End(Not run)



\end{ExampleCode}
\end{Examples}
\inputencoding{utf8}
\HeaderA{mysummary\_apa}{mysummary\_apa: Generate descriptive statistics in APA style (English or Korean)}{mysummary.Rul.apa}
%
\begin{Description}
mysummary\_apa: Generate descriptive statistics in APA style (English or Korean)
\end{Description}
%
\begin{Usage}
\begin{verbatim}
mysummary_apa(df, sd = FALSE, eng = FALSE)
\end{verbatim}
\end{Usage}
%
\begin{Arguments}
\begin{ldescription}
\item[\code{df}] A data frame containing descriptive statistics, including variables such as N, MEAN, SD, Skew, and Kurt.

\item[\code{sd}] A logical value. If TRUE, includes standard deviation (SD) along with the mean in the output. Default is FALSE.

\item[\code{eng}] A logical value. If TRUE, the output is in English. If FALSE, the output is in Korean. Default is TRUE.
\end{ldescription}
\end{Arguments}
%
\begin{Value}
A string formatted in APA style (English or Korean), summarizing the descriptive statistics in order of mean values and assessing normality based on skewness and kurtosis.
\end{Value}
%
\begin{Examples}
\begin{ExampleCode}
## Not run: 
df <- data.frame(
  var = c("Q33_1", "Q33_2", "Q33_3"),
  N = c(4542, 4542, 4542),
  MEAN = c(1.9489, 1.9337, 1.9243),
  SD = c(0.84679, 0.86568, 0.86960),
  Skew = c(0.72576, 0.63900, 0.70717),
  Kurt = c(0.63369, 0.087763, 0.27470)
)
mysummary_apa(df, sd = TRUE, eng = TRUE)
mysummary_apa(df, sd = TRUE, eng = FALSE)

## End(Not run)
\end{ExampleCode}
\end{Examples}
\inputencoding{utf8}
\HeaderA{na2zero}{NA replace to value default 0}{na2zero}
%
\begin{Description}
NA replace to value default 0
\end{Description}
%
\begin{Usage}
\begin{verbatim}
na2zero(data, imp = 0, range = NULL)
\end{verbatim}
\end{Usage}
%
\begin{Arguments}
\begin{ldescription}
\item[\code{data}] data.frame or matrix

\item[\code{imp}] imputation value default 0

\item[\code{range}] range is column
\end{ldescription}
\end{Arguments}
%
\begin{Value}
dat
\end{Value}
%
\begin{Examples}
\begin{ExampleCode}

## Not run: 

data %>% na2zero()
data %>% na2zero(0)
data %>% na2zero(1)

#'

# Example data frame entry
df <- data.frame(a = c(1, NA, 3), b = c(NA, 2, 3))
na2zero(df)

# Vector input examples
vec <- c(1, NA, 3)
print(na2zero(vec))

# Character Vector Input Example
char_vec <- c("a", NA, "b")
print(na2zero(char_vec))
#'


## End(Not run)

\end{ExampleCode}
\end{Examples}
\inputencoding{utf8}
\HeaderA{na\_remover}{na\_remover row unite}{na.Rul.remover}
%
\begin{Description}
na\_remover row unite
\end{Description}
%
\begin{Usage}
\begin{verbatim}
na_remover(df_var, pattern = "", type = "data")
\end{verbatim}
\end{Usage}
%
\begin{Arguments}
\begin{ldescription}
\item[\code{df\_var}] df vector

\item[\code{pattern}] ""

\item[\code{type}] data
\end{ldescription}
\end{Arguments}
%
\begin{Value}
df
\end{Value}
\inputencoding{utf8}
\HeaderA{nice\_aov\_df}{Nice tables for papers in aov\_df results  to  word, html}{nice.Rul.aov.Rul.df}
%
\begin{Description}
Nice tables for papers in aov\_df results  to  word, html
\end{Description}
%
\begin{Usage}
\begin{verbatim}
nice_aov_df(df, show = TRUE, tibble = FALSE, digits = 2)
\end{verbatim}
\end{Usage}
%
\begin{Arguments}
\begin{ldescription}
\item[\code{df}] aov\_Df

\item[\code{show}] TRUE nice, FALSE pass

\item[\code{digits}] 2
\end{ldescription}
\end{Arguments}
%
\begin{Value}
result
\end{Value}
%
\begin{Examples}
\begin{ExampleCode}

## Not run: 
#'
aov_df(data = mtcars,
       dv_var = "mpg",
       iv_vars = c("cyl", "gear", "carb") )

## generate nice table
aov_df(data = mtcars,
       dv_var = "mpg",
       iv_vars = c("cyl", "gear", "carb") )%>%
  nice_aov_df()


aov_df(data = mtcars,
       dv_var = "mpg",
       iv_vars = c("cyl", "gear", "carb") )%>%
  nice_aov_df() %>%web()



## End(Not run)
\end{ExampleCode}
\end{Examples}
\inputencoding{utf8}
\HeaderA{nice\_table}{같은 것이 반복되면 ""처리}{nice.Rul.table}
%
\begin{Description}
같은 것이 반복되면 ""처리
\end{Description}
%
\begin{Usage}
\begin{verbatim}
nice_table(data, exclude = NULL)
\end{verbatim}
\end{Usage}
%
\begin{Arguments}
\begin{ldescription}
\item[\code{data}] data

\item[\code{exclude}] not adapt col
\end{ldescription}
\end{Arguments}
%
\begin{Value}
data ressult
\end{Value}
%
\begin{Examples}
\begin{ExampleCode}

## Not run: 
#'
# Example usage
data <- data.frame(
  Var = c("a", "a", "a", "a", "a", "a", "a", "b",
  "b", "b", "b", "b", "c", "c", "c", "c"),
  alpha_95.CI = c("alpha = 0.855, 95%CI[0.8, 0.9]",
  "alpha = 0.855, 95%CI[0.8, 0.9]",
                  "alpha = 0.855, 95%CI[0.8, 0.9]",
                  "alpha = 0.855, 95%CI[0.8, 0.9]",
                  "alpha = 0.855, 95%CI[0.8, 0.9]",
                  "alpha = 0.855, 95%CI[0.8, 0.9]",
                  "alpha = 0.855, 95%CI[0.8, 0.9]",
                  "alpha = 0.805, 95%CI[0.74, 0.86]",
                  "alpha = 0.805, 95%CI[0.74, 0.86]",
                  "alpha = 0.805, 95%CI[0.74, 0.86]",
                  "alpha = 0.805, 95%CI[0.74, 0.86]",
                  "alpha = 0.805, 95%CI[0.74, 0.86]",
                  "alpha = 0.534, 95%CI[0.36, 0.67]",
                  "alpha = 0.534, 95%CI[0.36, 0.67]",
                  "alpha = 0.534, 95%CI[0.36, 0.67]",
                  "alpha = 0.534, 95%CI[0.36, 0.67]"),
  subfactor = c("B02", "B14", "B10", "B04", "B01", "B15",
   "B03", "B09", "B25", "B07", "B06", "B05", "B17",
   "B18", "B21", "B22"),
  cronbach_alpha = c(0.8175800, 0.8192332,
  0.8375942, 0.8163753, 0.8505183, 0.8488448,
   0.8445658, 0.7607713, 0.7494337,
                     0.7534147, 0.7620378,
                      0.8099670, 0.4376783,
                       0.5111512, 0.4434714, 0.4554295)
)

# Apply the function to the data
nice_table(data)


## End(Not run)


\end{ExampleCode}
\end{Examples}
\inputencoding{utf8}
\HeaderA{nice\_Var}{bind\_alpha\_table문자 정리 함수}{nice.Rul.Var}
%
\begin{Description}
bind\_alpha\_table문자 정리 함수

nice\_Var alpha분석한 경우 첫열의 a1->a
\end{Description}
%
\begin{Usage}
\begin{verbatim}
nice_Var(data)

nice_Var(data)
\end{verbatim}
\end{Usage}
%
\begin{Arguments}
\begin{ldescription}
\item[\code{data}] alpha data
\end{ldescription}
\end{Arguments}
%
\begin{Value}
table data.frame

data.frame
\end{Value}
%
\begin{Examples}
\begin{ExampleCode}

## Not run: 
#'
jutr2 %>% select(B01:B25) %>%
  bind_alpha_table(c(1:5),c(6:10),c(11:15),c(16:20),c(21:25))%>%
  nice_Var()%>%  # <----this
  nice_table()

## End(Not run)
\end{ExampleCode}
\end{Examples}
\inputencoding{utf8}
\HeaderA{Normality\_test}{Multivariate normality tests}{Normality.Rul.test}
%
\begin{Description}
Multivariate normality tests
\end{Description}
%
\begin{Usage}
\begin{verbatim}
Normality_test(multiVariable_data)
\end{verbatim}
\end{Usage}
%
\begin{Arguments}
\begin{ldescription}
\item[\code{multiVariable\_data}] data
\end{ldescription}
\end{Arguments}
%
\begin{Value}
normality
\end{Value}
%
\begin{Examples}
\begin{ExampleCode}

## Not run: 
mtcars%>%Normality_test()

## End(Not run)

\end{ExampleCode}
\end{Examples}
\inputencoding{utf8}
\HeaderA{obs\_exp\_table}{Observation Expectation Table with calculate\_chi\_sig}{obs.Rul.exp.Rul.table}
%
\begin{Description}
Observation Expectation Table with calculate\_chi\_sig
\end{Description}
%
\begin{Usage}
\begin{verbatim}
obs_exp_table(obs_data)
\end{verbatim}
\end{Usage}
%
\begin{Arguments}
\begin{ldescription}
\item[\code{obs\_data}] matrix
\end{ldescription}
\end{Arguments}
%
\begin{Value}
matrix
\end{Value}
\inputencoding{utf8}
\HeaderA{oneway\_aov}{oneway ANVOA calculation}{oneway.Rul.aov}
%
\begin{Description}
oneway ANVOA calculation
\end{Description}
%
\begin{Usage}
\begin{verbatim}
oneway_aov(formula, data, format = "pandoc")
\end{verbatim}
\end{Usage}
%
\begin{Arguments}
\begin{ldescription}
\item[\code{formula}] ex)mpg \textasciitilde{} cyl

\item[\code{data}] data.frame

\item[\code{format}] 'markdwon', 'pandoc', 'html'
\end{ldescription}
\end{Arguments}
%
\begin{Value}
anova result
\end{Value}
%
\begin{Examples}
\begin{ExampleCode}
## Not run: 
oneway_aov(mpg ~ cyl, data = mtcars)

## End(Not run)

\end{ExampleCode}
\end{Examples}
\inputencoding{utf8}
\HeaderA{pall}{Print and Inspect Data with Multiple Output Modes}{pall}
%
\begin{Description}
A unified data inspection function supporting console printing,
RStudio viewing, interactive DataTables (DT), and publication-ready
tables via knitr::kable().
\end{Description}
%
\begin{Usage}
\begin{verbatim}
pall(
  data,
  n = Inf,
  type = c("print", "view", "dt", "kable"),
  dt_browser = TRUE,
  dt_page = 20,
  dt_scroll_y = 650,
  dt_digits = 2,
  dt_highlight_na = TRUE,
  kable_format = c("markdown", "html", "latex", "pipe"),
  kable_digits = 3,
  max_n = 5000,
  summary = TRUE,
  return_data = FALSE
)
\end{verbatim}
\end{Usage}
%
\begin{Arguments}
\begin{ldescription}
\item[\code{data}] A data.frame or tibble.

\item[\code{n}] Number of rows to display. Default is Inf.

\item[\code{type}] Output type: "print", "view", "dt", or "kable".

\item[\code{dt\_browser}] Logical. If TRUE (default), DT output opens in browser.

\item[\code{dt\_page}] Number of rows per DT page.

\item[\code{dt\_scroll\_y}] Vertical scroll height (px) for DT.

\item[\code{dt\_digits}] Decimal digits for numeric variables in DT.

\item[\code{dt\_highlight\_na}] Logical; highlight missing values in DT.

\item[\code{kable\_format}] Output format for kable:
"markdown", "html", "latex", or "pipe".

\item[\code{kable\_digits}] Decimal digits for numeric values in kable output.

\item[\code{max\_n}] Safety limit for maximum rows.

\item[\code{summary}] Logical; print row/column summary.

\item[\code{return\_data}] Logical; return sliced data invisibly.
\end{ldescription}
\end{Arguments}
%
\begin{Details}
In DT mode, the table is rendered in a web browser by default, with
automatic variable-type detection and a summary panel displayed above
the table.

In kable mode, output is routed automatically:
markdown / pipe / latex -> console
html                    -> web browser
\end{Details}
%
\begin{Value}
Invisibly returns the sliced data.
\end{Value}
%
\begin{Examples}
\begin{ExampleCode}
pall(mtcars, n = 10)
pall(iris, type = "dt")
pall(head(mtcars), type = "kable", kable_format = "markdown")
pall(head(mtcars), type = "kable", kable_format = "htmml")

\end{ExampleCode}
\end{Examples}
\inputencoding{utf8}
\HeaderA{path}{Define a Structural Path (DSL)}{path}
%
\begin{Description}
구조방정식에서 잠재변수 간 회귀경로를 정의하기 위한 DSL 함수이다.
\end{Description}
%
\begin{Usage}
\begin{verbatim}
path(from, to)

path(from, to)
\end{verbatim}
\end{Usage}
%
\begin{Arguments}
\begin{ldescription}
\item[\code{from}] var1

\item[\code{to}] var2
\end{ldescription}
\end{Arguments}
%
\begin{Value}
list with elements \code{from}, \code{to}.
\end{Value}
%
\begin{Examples}
\begin{ExampleCode}
path(c("F1","F2"), "F3")
\end{ExampleCode}
\end{Examples}
\inputencoding{utf8}
\HeaderA{patternGraph}{patternGraph}{patternGraph}
%
\begin{Description}
patternGraph
\end{Description}
%
\begin{Usage}
\begin{verbatim}
patternGraph(
  data,
  size = 4,
  strip = 14,
  text = 12,
  axis = 16,
  hjust = -0.4,
  yadd = 10,
  type = "g",
  xlab = "성조형",
  ylab = "빈도",
  show = TRUE,
  tolong = TRUE
)
\end{verbatim}
\end{Usage}
%
\begin{Arguments}
\begin{ldescription}
\item[\code{data}] data.frame

\item[\code{size}] bar text size

\item[\code{strip}] strip size

\item[\code{text}] axis text test size

\item[\code{axis}] aixs size

\item[\code{hjust}] hjust

\item[\code{yadd}] yadd

\item[\code{type}] type "g", "data", "all"

\item[\code{xlab}] xlab

\item[\code{ylab}] ylab

\item[\code{show}] show.legend

\item[\code{tolong}] long format transpose
\end{ldescription}
\end{Arguments}
%
\begin{Value}
graph
\end{Value}
%
\begin{Examples}
\begin{ExampleCode}
## Not run: 
aaa=kge_bind2 %>% filter(지역 =="부산" & N ==1) %>%
  auto_pattern("고유어1") %>%
  kge_chisq_table("a1","성조", "고유어1 부산")

#sAVE
data_long<- aaa$chi_table %>%
  rownames_to_column("syllabic") %>%
  pivot_longer(names_to = "accent", values_to = 'ratio', cols=2:4)
data_long

data_long%>% ggplot(aes(x = accent, y = ratio))+
  geom_bar(stat = "identity", aes( fill = accent),
           position = "dodge", show.legend = FALSE)+
  geom_hline(yintercept = 1, linetype=2, color="gray80")+
  geom_text(aes(label = round(ratio,2)), hjust = -0.1, size = 4)+
  ylim(0,max(data1$ratio)+0.07)+
  coord_flip()+
  theme_bw()+
  theme(axis.text = element_text(size= 12),
        axis.title = element_text(size= 12),
        strip.text = element_text(size= 14)
  )+
  scale_fill_grey(start = 0, end = 0.7) +
  facet_wrap(~ syllabic )

aaa%>% patternGraph()


## End(Not run)
\end{ExampleCode}
\end{Examples}
\inputencoding{utf8}
\HeaderA{patternGraph2}{patternGraph2}{patternGraph2}
%
\begin{Description}
patternGraph2
\end{Description}
%
\begin{Usage}
\begin{verbatim}
patternGraph2(
  data,
  type = "g",
  raw = TRUE,
  Ncol = NULL,
  yadd = 0.45,
  xlab = "음절",
  ylab = "관측기대비율",
  strip_size = 16,
  axis_size = 16,
  text_size = 14,
  Colsnames = NULL,
  size_bartext = 5
)
\end{verbatim}
\end{Usage}
%
\begin{Arguments}
\begin{ldescription}
\item[\code{data}] data,frame

\item[\code{type}] g

\item[\code{raw}] raw data

\item[\code{Ncol}] grid

\item[\code{yadd}] up yadd

\item[\code{xlab}] xlab

\item[\code{ylab}] ylab

\item[\code{strip\_size}] strip

\item[\code{axis\_size}] axis

\item[\code{text\_size}] text

\item[\code{Colsnames}] new select col pivot long

\item[\code{size\_bartext}] bartext
\end{ldescription}
\end{Arguments}
%
\begin{Value}
result fill data
\end{Value}
%
\begin{Examples}
\begin{ExampleCode}
## Not run: 
aaa=kge_bind2 %>% filter(지역 =="부산" & N ==1) %>%
  auto_pattern("고유어1") %>%
  kge_chisq_table("a1","성조", "고유어1 부산")

#sAVE
data_long<- aaa$chi_table %>%
  rownames_to_column("syllabic") %>%
  pivot_longer(names_to = "accent", values_to = 'ratio', cols=2:4)
data_long

data_long%>% ggplot(aes(x = accent, y = ratio))+
  geom_bar(stat = "identity", aes( fill = accent),
           position = "dodge", show.legend = FALSE)+
  geom_hline(yintercept = 1, linetype=2, color="gray80")+
  geom_text(aes(label = round(ratio,2)), hjust = -0.1, size = 4)+
  ylim(0,max(data1$ratio)+0.07)+
  coord_flip()+
  theme_bw()+
  theme(axis.text = element_text(size= 12),
        axis.title = element_text(size= 12),
        strip.text = element_text(size= 14)
  )+
  scale_fill_grey(start = 0, end = 0.7) +
  facet_wrap(~ syllabic )

aaa%>% patternGraph_obs_exp()


## End(Not run)
\end{ExampleCode}
\end{Examples}
\inputencoding{utf8}
\HeaderA{patternGraph\_obs\_exp}{patternGraph\_obs\_exp}{patternGraph.Rul.obs.Rul.exp}
%
\begin{Description}
patternGraph\_obs\_exp
\end{Description}
%
\begin{Usage}
\begin{verbatim}
patternGraph_obs_exp(
  data,
  raw = TRUE,
  Ncol = NULL,
  yadd = 0.3,
  strip_size = 16,
  axis_size = 15,
  text_size = 13,
  size_bartext = 4,
  xlab = "성조형",
  ylab = "관측기대비율",
  Colsnames = NULL,
  type = "g"
)
\end{verbatim}
\end{Usage}
%
\begin{Arguments}
\begin{ldescription}
\item[\code{data}] data,frame

\item[\code{raw}] raw data

\item[\code{Ncol}] grid

\item[\code{yadd}] up yadd

\item[\code{strip\_size}] strip

\item[\code{axis\_size}] axis

\item[\code{text\_size}] text

\item[\code{size\_bartext}] bartext

\item[\code{xlab}] xlab

\item[\code{ylab}] ylab

\item[\code{Colsnames}] new select col pivot long

\item[\code{type}] res
\end{ldescription}
\end{Arguments}
%
\begin{Value}
result fill data
\end{Value}
%
\begin{Examples}
\begin{ExampleCode}
## Not run: 
aaa=kge_bind2 %>% filter(지역 =="부산" & N ==1) %>%
  auto_pattern("고유어1") %>%
  kge_chisq_table("a1","성조", "고유어1 부산")

#sAVE
data_long<- aaa$chi_table %>%
  rownames_to_column("syllabic") %>%
  pivot_longer(names_to = "accent", values_to = 'ratio', cols=2:4)
data_long

data_long%>% ggplot(aes(x = accent, y = ratio))+
  geom_bar(stat = "identity", aes( fill = accent),
           position = "dodge", show.legend = FALSE)+
  geom_hline(yintercept = 1, linetype=2, color="gray80")+
  geom_text(aes(label = round(ratio,2)), hjust = -0.1, size = 4)+
  ylim(0,max(data1$ratio)+0.07)+
  coord_flip()+
  theme_bw()+
  theme(axis.text = element_text(size= 12),
        axis.title = element_text(size= 12),
        strip.text = element_text(size= 14)
  )+
  scale_fill_grey(start = 0, end = 0.7) +
  facet_wrap(~ syllabic )

aaa%>% patternGraph_obs_exp()

#신진영(2013)
shin2013_mat =matrix(c(14,1,23,3,4,0,9,0,13,18,89,8,11,40,46,12,3,4,22,0), ncol= 5)
# colnames(shin2013_mat)=c("지능지수가높다","학업성취도가높다","창의적사고를한다","특정분야에서뛰어난능력을보인다","모든분야에서뛰어난능력을보인다")
colnames(shin2013_mat)=c("지능지수","학업성취도","창의적사고","특정분야","모든분야")
rownames(shin2013_mat)=c("고졸","초대졸","대졸","대학원이상")
shin2013_mat

shin2013_mat%>% data.frame() %>%
  long_df("학력","영재성", cols=2:6) %>%
  add_rows_freq("real") %>%
  chisq_test_kge("학력","영재성", type="data_graph") %>%
  patternGraph_obs_exp_ko()


## End(Not run)
\end{ExampleCode}
\end{Examples}
\inputencoding{utf8}
\HeaderA{patternGraph\_obs\_exp\_ko}{patternGraph\_obs\_exp}{patternGraph.Rul.obs.Rul.exp.Rul.ko}
%
\begin{Description}
patternGraph\_obs\_exp
\end{Description}
%
\begin{Usage}
\begin{verbatim}
patternGraph_obs_exp_ko(
  data,
  raw = TRUE,
  Ncol = NULL,
  yadd = 0.5,
  strip_size = 16,
  axis_size = 16,
  text_size = 13,
  size_bartext = 5,
  xlab = "기준변수",
  ylab = "관측기대비율",
  Colsnames = NULL,
  trans = FALSE,
  type = "g"
)
\end{verbatim}
\end{Usage}
%
\begin{Arguments}
\begin{ldescription}
\item[\code{data}] data,frame

\item[\code{raw}] raw data

\item[\code{Ncol}] grid

\item[\code{yadd}] up yadd

\item[\code{strip\_size}] strip

\item[\code{axis\_size}] axis

\item[\code{text\_size}] text

\item[\code{size\_bartext}] bartext

\item[\code{xlab}] xlab

\item[\code{ylab}] ylab

\item[\code{Colsnames}] new select col pivot long

\item[\code{trans}] trans plot transpose

\item[\code{type}] res
\end{ldescription}
\end{Arguments}
%
\begin{Value}
result fill data
\end{Value}
%
\begin{Examples}
\begin{ExampleCode}
## Not run: 
aaa=kge_bind2 %>% filter(지역 =="부산" & N ==1) %>%
  auto_pattern("고유어1") %>%
  kge_chisq_table("a1","성조", "고유어1 부산")

#sAVE
data_long<- aaa$chi_table %>%
  rownames_to_column("syllabic") %>%
  pivot_longer(names_to = "accent", values_to = 'ratio', cols=2:4)
data_long

data_long%>% ggplot(aes(x = accent, y = ratio))+
  geom_bar(stat = "identity", aes( fill = accent),
           position = "dodge", show.legend = FALSE)+
  geom_hline(yintercept = 1, linetype=2, color="gray80")+
  geom_text(aes(label = round(ratio,2)), hjust = -0.1, size = 4)+
  ylim(0,max(data1$ratio)+0.07)+
  coord_flip()+
  theme_bw()+
  theme(axis.text = element_text(size= 12),
        axis.title = element_text(size= 12),
        strip.text = element_text(size= 14)
  )+
  scale_fill_grey(start = 0, end = 0.7) +
  facet_wrap(~ syllabic )

aaa%>% patternGraph_obs_exp_ko()

#신진영(2013)
shin2013_mat =matrix(c(14,1,23,3,4,0,9,0,13,18,89,8,11,40,46,12,3,4,22,0), ncol= 5)
# colnames(shin2013_mat)=c("지능지수가높다","학업성취도가높다","창의적사고를한다","특정분야에서뛰어난능력을보인다","모든분야에서뛰어난능력을보인다")
colnames(shin2013_mat)=c("지능지수","학업성취도","창의적사고","특정분야","모든분야")
rownames(shin2013_mat)=c("고졸","초대졸","대졸","대학원이상")
shin2013_mat

shin2013_mat%>% data.frame() %>%
  long_df("학력","영재성", cols=2:6) %>%
  add_rows_freq("real") %>%
  chisq_test_kge("학력","영재성", type="data_graph") %>%
  patternGraph_obs_exp_ko()




## End(Not run)

\end{ExampleCode}
\end{Examples}
\inputencoding{utf8}
\HeaderA{pattern\_detector}{pattern\_detector: missing data detect}{pattern.Rul.detector}
%
\begin{Description}
pattern\_detector: missing data detect
\end{Description}
%
\begin{Usage}
\begin{verbatim}
pattern_detector(df, missing = TRUE, data = FALSE, pattern = "")
\end{verbatim}
\end{Usage}
%
\begin{Arguments}
\begin{ldescription}
\item[\code{df}] data

\item[\code{missing}] missing TRUE detecte missing T,F

\item[\code{data}] TRUE data print, FALSE rownames number

\item[\code{pattern}] pattern "" or  NA
\end{ldescription}
\end{Arguments}
%
\begin{Value}
data
\end{Value}
%
\begin{Examples}
\begin{ExampleCode}

## Not run: 

ed_tam %>% pattern_detector()
ed_tam %>% pattern_detector(FALSE,TRUE)%>% head()
# missing data chek
ed_tam %>% pattern_detector(missing=T, data=F)
#missing data output
ed_tam %>% pattern_detector(missing=T, data=T) %>%head(30)

# not missing
ed_tam %>% pattern_detector(data=T,missing=F) %>%head()

# ot missing , not data
ed_tam %>% pattern_detector(data=F,missing=F)



## End(Not run)


\end{ExampleCode}
\end{Examples}
\inputencoding{utf8}
\HeaderA{pcfa}{principal component factor analysis function}{pcfa}
%
\begin{Description}
principal component factor analysis function
\end{Description}
%
\begin{Usage}
\begin{verbatim}
pcfa(
  R = NULL,
  nfactor = NULL,
  rowname = NULL,
  digits = 2,
  xlim = NULL,
  ylim = NULL,
  cex = 1.2,
  vjust = -0.5,
  yadd = 10,
  name = "auto",
  pos = 1
)
\end{verbatim}
\end{Usage}
%
\begin{Arguments}
\begin{ldescription}
\item[\code{R}] correlation matrix

\item[\code{nfactor}] number of factor

\item[\code{rowname}] matrix rowname

\item[\code{digits}] digits = 3

\item[\code{xlim}] NULL

\item[\code{ylim}] NULL

\item[\code{cex}] size =1

\item[\code{yadd}] yadd = 10,

\item[\code{name}] 'auto', 'manual'

\item[\code{pos}] positon  1
\end{ldescription}
\end{Arguments}
%
\begin{Value}
result table, biplot, dimension
\end{Value}
%
\begin{Examples}
\begin{ExampleCode}
## Not run: 
R= matrix(
c(1.00, 0.83, 0.78, 0.70, 0.66, 0.63,
  0.83, 1.00, 0.67, 0.67, 0.65, 0.57,
  0.78, 0.67, 1.00, 0.64, 0.54, 0.51,
  0.70, 0.67, 0.64, 1.00, 0.45, 0.51,
  0.66, 0.65, 0.54, 0.45, 1.00, 0.40,
  0.63, 0.57, 0.51, 0.51, 0.40,1.00),
 byrow=T, nrow=6)
subject=c("classic","france","english","math","pitch","music")

colnames(R)=subject
rownames(R)=subject
R
pcfa(R, nfactor = 2)

library(MVT)
data(examScor)

#if you when input cor
pcfa(cor(examScor), nfactor = 2)

## if you When input data
pcfa( examScor, nfactor = 2)


## End(Not run)

\end{ExampleCode}
\end{Examples}
\inputencoding{utf8}
\HeaderA{pcfa\_porp\_plot}{pcfa\_porp\_plot}{pcfa.Rul.porp.Rul.plot}
%
\begin{Description}
pcfa\_porp\_plot
\end{Description}
%
\begin{Usage}
\begin{verbatim}
pcfa_porp_plot(Data, digits = 2, prop = FALSE, ymax = NA)
\end{verbatim}
\end{Usage}
%
\begin{Arguments}
\begin{ldescription}
\item[\code{Data}] factanal result uniqueness

\item[\code{digits}] digits 3

\item[\code{prop}] FALSE

\item[\code{ymax}] plot y size
\end{ldescription}
\end{Arguments}
%
\begin{Value}
graph
\end{Value}
%
\begin{Examples}
\begin{ExampleCode}
## Not run: 

library(MVT)
library(GPArotation)
data(examScor)
X = examScor
Z = scale(X, scale=T)  # scale
R = cor(X)  # correlation matrix

## factor analysis tools : factanal
mlfa <- factanal(Z, factors = 2, rotation = "varimax")

mlfa|> pcfa_porp_plot()


## End(Not run)
\end{ExampleCode}
\end{Examples}
\inputencoding{utf8}
\HeaderA{pdata}{pdata plspm group change long data format}{pdata}
%
\begin{Description}
pdata plspm group change long data format
\end{Description}
%
\begin{Usage}
\begin{verbatim}
pdata(widedata, rownames_to_column=FALSE)
\end{verbatim}
\end{Usage}
%
\begin{Arguments}
\begin{ldescription}
\item[\code{data}] wide data
\end{ldescription}
\end{Arguments}
%
\begin{Value}
long
\end{Value}
\inputencoding{utf8}
\HeaderA{perform\_chi\_square\_test}{Various output functions after chi-square analysis}{perform.Rul.chi.Rul.square.Rul.test}
%
\begin{Description}
Various output functions after chi-square analysis
\end{Description}
%
\begin{Usage}
\begin{verbatim}
perform_chi_square_test(observed, type = "data", simple = FALSE)
\end{verbatim}
\end{Usage}
%
\begin{Arguments}
\begin{ldescription}
\item[\code{observed}] table data

\item[\code{type}] type data, res, chisq\_test, cramersv, observed\_over\_expected\_ratios, cell\_p\_values, p\_sig

\item[\code{simple}] * or ***
\end{ldescription}
\end{Arguments}
%
\begin{Value}
chisq table obs/exp
\end{Value}
%
\begin{Examples}
\begin{ExampleCode}

## Not run: 

#' #'mat = matrix(c(36, 67, 11,
          0, 35, 44,
         41, 108, 1,
          56, 87, 54),
           nrow = 4, byrow = TRUE.
   dimnames = list(c("공명음", "마찰음",
     "유기음_경음", "평파열음_평파찰음"),  c("H", "H(H)", "L")))

 mat %>% perform_chi_square_test()

 #perform_chi_square_test
mtcars%>%select(am, cyl)%>%table()%>%perform_chi_square_test()
mtcars%>%select(am, cyl)%>%table()%>%perform_chi_square_test("all")

## End(Not run)

\end{ExampleCode}
\end{Examples}
\inputencoding{utf8}
\HeaderA{pivot\_table}{pivot}{pivot.Rul.table}
%
\begin{Description}
pivot
\end{Description}
%
\begin{Usage}
\begin{verbatim}
pivot_table(
  data,
  formula,
  names_from = NULL,
  unite = FALSE,
  col = 1:2,
  remove = TRUE,
  change_col = TRUE
)
\end{verbatim}
\end{Usage}
%
\begin{Arguments}
\begin{ldescription}
\item[\code{data}] data

\item[\code{formula}] formul

\item[\code{names\_from}] column

\item[\code{unite}] 1,2 colum unite

\item[\code{col}] 합쳐질 컬럼

\item[\code{remove}] 기존것 지우기

\item[\code{change\_col}] 열의 위치 바꾸기
\end{ldescription}
\end{Arguments}
%
\begin{Value}
데이터
\end{Value}
%
\begin{Examples}
\begin{ExampleCode}

## Not run: 

#' kge_bind3a %>%
filter(type == "고유어1")%>%
  pivot_table(~ area + age + 성조, unite=TRUE)
#'
kge_bind3a %>%
  filter(type == "고유어1")%>%
  pivot_table(~  age + area+ 성조, unite=FALSE)

#고유어1 공명음  지역별 연령
kge_bind3a %>%
  filter(type == "고유어1" & a1=="공명음")%>%
  pivot_table(~ age + area + 성조, unite=T)  #%>% md()



kge_bind3a %>%
  filter(type == "고유어1")%>%
  pivot_table(~  age +  성조, unite=FALSE)

#고유어1 공명음  지역별 연령
kge_bind3a %>%
  filter(type == "고유어1" & a1=="공명음")%>%
  pivot_table(~  area + 성조, unite=F)  #%>% md()


  
## End(Not run)
#'

\end{ExampleCode}
\end{Examples}
\inputencoding{utf8}
\HeaderA{plot\_lpa}{Latent Profile Analysis Profile Plot}{plot.Rul.lpa}
%
\begin{Description}
Latent Profile Analysis Profile Plot
\end{Description}
%
\begin{Usage}
\begin{verbatim}
plot_lpa(
  data,
  n_profiles = 3,
  model_name = "EEE",
  view = "pair",
  alpha = 0.6,
  gtype = "normal",
  angle = 90,
  size.x = 13,
  size.p = 5,
  ncol = 2,
  fct_reorder = NULL,
  show.legend = FALSE,
  legend.text.size = 16,
  legend_position = "top",
  linewidth = 1,
  title_size = 18,
  seed = 123
)
\end{verbatim}
\end{Usage}
%
\begin{Arguments}
\begin{ldescription}
\item[\code{data}] A data frame containing numeric variables for latent profile analysis.

\item[\code{n\_profiles}] Integer. Number of latent profiles to identify. Default is 3.

\item[\code{model\_name}] Character. Model name to use in Mclust for Gaussian mixture modeling (e.g., "EEE", "EEI", "VVI", "VVV"). Default is "EEE".

\item[\code{view}] Character. Determines the view type of the plots ("pair", "each"). Default is "pair".

\item[\code{alpha}] Numeric. Transparency level for bar plots. Default is 0.6.

\item[\code{gtype}] Character. Graph type to generate ("normal", "trans", "each"). Default is "normal".

\item[\code{angle}] Numeric. Angle for the x-axis text labels. Default is 90.

\item[\code{size.x}] Numeric. Font size for the x-axis text labels. Default is 13.

\item[\code{size.p}] Numeric. Size of points in the plot. Default is 5.

\item[\code{ncol}] Integer. Number of columns for the facet plots in "each" view. Default is 2.

\item[\code{fct\_reorder}] Character vector or NULL. Specifies the order of factor levels for plotting. Default is NULL.

\item[\code{show.legend}] Logical. Whether to show legends for line plots. Default is FALSE.

\item[\code{legend.text.size}] Numeric. Font size for legend text. Default is 16.

\item[\code{legend\_position}] Character. Position of the legend ("top", "bottom", "left", "right"). Default is "top".

\item[\code{linewidth}] Numeric. Line width for plots. Default is 1.

\item[\code{title\_size}] Numeric. Font size for plot titles. Default is 18.

\item[\code{seed}] Integer or NULL. Seed value for reproducibility. Default is 123.
\end{ldescription}
\end{Arguments}
%
\begin{Value}
A list containing:
\begin{ldescription}
\item[\code{std}] Standardized mean data for each profile.
\item[\code{est}] Estimated raw mean data for each profile.
\item[\code{graph}] Combined plots for the profiles.
\item[\code{plot\_data}] Data used for plotting.
\item[\code{Class\_freq}] Frequency table of profile classifications.
\end{ldescription}
\end{Value}
%
\begin{Examples}
\begin{ExampleCode}
## Not run: 
# Example usage with iris dataset
result <- plot_lpa(iris[, -5], n_profiles = 3, model_name = "EEE")
print(result$graph)

## End(Not run)
\end{ExampleCode}
\end{Examples}
\inputencoding{utf8}
\HeaderA{plot\_pie\_chart}{pie chart(using fei)}{plot.Rul.pie.Rul.chart}
%
\begin{Description}
pie chart(using fei)
\end{Description}
%
\begin{Usage}
\begin{verbatim}
plot_pie_chart(
  data,
  title = "",
  size.text = 5,
  title.size = 20,
  vjust = 0.5,
  palette = "Set2"
)
\end{verbatim}
\end{Usage}
%
\begin{Arguments}
\begin{ldescription}
\item[\code{data}] data

\item[\code{title}] title

\item[\code{size.text}] 5

\item[\code{title.size}] 20

\item[\code{vjust}] 0.5

\item[\code{palette}] "Set2", "Set3"
\end{ldescription}
\end{Arguments}
%
\begin{Value}
graph
\end{Value}
%
\begin{Examples}
\begin{ExampleCode}

## Not run: 
#'
# 예시 데이터
data <- data.frame(
  구분 = c("전혀 아니다", "아니다", "중간이다", "그렇다", "매우 그렇다"),
  빈도 = c(52, 90, 304, 257, 114),
  비율 = c(6.36, 11.02, 37.21, 31.46, 13.95)
)

# 함수 호출
plot_pie_chart(data,,"Q1응답",0.5)

## End(Not run)

\end{ExampleCode}
\end{Examples}
\inputencoding{utf8}
\HeaderA{plspm\_add\_blocks}{plspm\_add\_blocks, indicator blocks}{plspm.Rul.add.Rul.blocks}
%
\begin{Description}
plspm\_add\_blocks, indicator blocks
\end{Description}
%
\begin{Usage}
\begin{verbatim}
plspm_add_blocks(newData, blocks, interactionTerm)
\end{verbatim}
\end{Usage}
%
\begin{Arguments}
\begin{ldescription}
\item[\code{newData}] plspm\_make\_prod data

\item[\code{blocks}] clocks

\item[\code{interactionTerm}] interaction Term list data
\end{ldescription}
\end{Arguments}
%
\begin{Value}
list
\end{Value}
%
\begin{Examples}
\begin{ExampleCode}
## Not run: 
#' data(satisfaction)
#'
newData1 <- plspm_make_prod(satisfaction, sat_blocks1,
                            list(intIE = interlist("IMAG",1,2, "EXPE",1,2) ) )


interactionTerm1 = list(list(from = "IMAG*EXPE", to = "LOY"))
plspm_add_blocks(newData1, sat_blocks1, interactionTerm1)


## End(Not run)

\end{ExampleCode}
\end{Examples}
\inputencoding{utf8}
\HeaderA{plspm\_add\_inter\_paths}{plspm\_add\_inter\_paths}{plspm.Rul.add.Rul.inter.Rul.paths}
%
\begin{Description}
plspm\_add\_inter\_paths
\end{Description}
%
\begin{Usage}
\begin{verbatim}
plspm_add_inter_paths(path_mat, paths)
\end{verbatim}
\end{Usage}
%
\begin{Arguments}
\begin{ldescription}
\item[\code{path\_mat}] path matrix data

\item[\code{paths}] add path, interactionTerm
\end{ldescription}
\end{Arguments}
%
\begin{Value}
mat
\end{Value}
%
\begin{Examples}
\begin{ExampleCode}


## Not run: 
#'
data <- data.frame(
  IMAG = c(0, 1, 0, 0, 1, 1),
  EXPE = c(0, 0, 1, 1, 1, 0),
  QUAL = c(0, 0, 0, 1, 1, 0),
  VAL  = c(0, 0, 0, 0, 1, 0),
  SAT  = c(0, 0, 0, 0, 0, 1),
  LOY  = c(0, 0, 0, 0, 0, 0),
  row.names = c("IMAG", "EXPE", "QUAL", "VAL", "SAT", "LOY")
)

# 함수 호출 예시
plspm_add_inter_paths(data, list(
  int1 = list(from = "IMAG*EXPE", to = "LOY"),
  int2 = list(from = "IMAG*QUAL", to = "LOY")
))

plspm_add_inter_paths(data, list(
  list(from = "IMAG*EXPE", to = "LOY"),
  list(from = "IMAG*QUAL", to = "LOY")
))



## End(Not run)


\end{ExampleCode}
\end{Examples}
\inputencoding{utf8}
\HeaderA{plspm\_apa\_boot\_paths}{bootstrap path effect apa analysis}{plspm.Rul.apa.Rul.boot.Rul.paths}
%
\begin{Description}
bootstrap path effect apa analysis
\end{Description}
%
\begin{Usage}
\begin{verbatim}
plspm_apa_boot_paths(data)
\end{verbatim}
\end{Usage}
%
\begin{Arguments}
\begin{ldescription}
\item[\code{data}] plspm\_boot\_paths\_sig
\end{ldescription}
\end{Arguments}
%
\begin{Value}
text
\end{Value}
%
\begin{Examples}
\begin{ExampleCode}

## Not run: 

#'

# 예시 데이터로 함수 호출
nfl_pls_boot2 <- data.frame(
  paths = c("Speical -> Scoring", "Rushing -> Offense", "Passing -> Offense", "Offense -> Scoring"),
  Original = c(0.311, 0.535, 0.789, 0.810),
  Mean.Boot = c(0.098, 0.522, 0.756, 0.762),
  Std.Error = c(0.317, 0.044, 0.252, 0.227),
  t = c("0.310", "11.840***", "2.995**", "3.361***"),
  `95%CI` = I(list(c(-0.484, 0.450), c(0.411, 0.592), c(0.563, 0.948), c(0.587, 0.895)))
)

result <- plspm_apa_boot_paths(nfl_pls_boot2)
cat(result)
#'

## End(Not run)

\end{ExampleCode}
\end{Examples}
\inputencoding{utf8}
\HeaderA{plspm\_apa\_grp}{plspm\_apa\_grp}{plspm.Rul.apa.Rul.grp}
%
\begin{Description}
plspm\_apa\_grp
\end{Description}
%
\begin{Usage}
\begin{verbatim}
plspm_apa_grp(data)
\end{verbatim}
\end{Usage}
%
\begin{Arguments}
\begin{ldescription}
\item[\code{data}] plspm\_grp\_summary
\end{ldescription}
\end{Arguments}
%
\begin{Value}
text
\end{Value}
%
\begin{Examples}
\begin{ExampleCode}

## Not run: 

# Example usage
data14 <- tibble::tribble(
  ~paths, ~global, ~`group.고등학교`, ~`group.중학교`, ~diff, ~t, ~p.value, ~sig,
  "자기효능감->진로동기", 0.88, 0.894, 0.81, 0.085, 0.908, 0.183, "no",
  "자기효능감->진로태도", -0.674, -0.749, 0.438, 1.19, 2.49, 0.007, "yes",
  "자기효능감->진로준비", 0.429, 0.542, 0.595, 0.053, 0.395, 0.347, "no",
  "진로동기->진로준비", 0.446, 0.431, 0.385, 0.046, 0.232, 0.408, "no",
  "진로태도->진로준비", 0.007, 0.098, -0.126, 0.224, 0.626, 0.266, "no"
)

plspm_apa_grp(data14)


## End(Not run)

\end{ExampleCode}
\end{Examples}
\inputencoding{utf8}
\HeaderA{plspm\_apa\_ind}{Function to interpret plspm results}{plspm.Rul.apa.Rul.ind}
%
\begin{Description}
Function to interpret plspm results
\end{Description}
%
\begin{Usage}
\begin{verbatim}
plspm_apa_ind(data)
\end{verbatim}
\end{Usage}
%
\begin{Arguments}
\begin{ldescription}
\item[\code{data}] plspm\_ind\_effect\_test data
\end{ldescription}
\end{Arguments}
%
\begin{Value}
text
\end{Value}
%
\begin{Examples}
\begin{ExampleCode}

## Not run: 
#'
# Example data frame
data11 <- tibble::tibble(
  paths = c("IMAG -> EXPE -> QUAL -> VAL -> SAT -> LOY",
            "IMAG -> EXPE -> QUAL -> SAT -> LOY",
            "IMAG -> EXPE -> VAL -> SAT -> LOY",
            "IMAG -> EXPE -> SAT -> LOY",
            "IMAG -> SAT -> LOY"),
  ind_effect = c(0.0970, 0.0297, 0.0178, -0.000790, 0.0995),
  SE = c(4.44e-4, 4.07e-3, 2.98e-3, 1.91e-2, 3.35e-2),
  Z = c(218.39, 7.31, 5.98, -0.0413, 2.97),
  p.value = c(0, 2.76e-13, 2.26e-9, 9.67e-1, 3.00e-3),
  sig = c("***", "***", "***", "ns", "**")
)

plspm_apa_ind(data11)

## exaample
plspm_ind_effect_test(satpls_boot)%>%plspm_apa_ind()


## End(Not run)
\end{ExampleCode}
\end{Examples}
\inputencoding{utf8}
\HeaderA{plspm\_apa\_paths}{plspm\_apa\_paths}{plspm.Rul.apa.Rul.paths}
%
\begin{Description}
plspm\_apa\_paths
\end{Description}
%
\begin{Usage}
\begin{verbatim}
plspm_apa_paths(data)
\end{verbatim}
\end{Usage}
%
\begin{Arguments}
\begin{ldescription}
\item[\code{data}] plspm\_paths\_sig
\end{ldescription}
\end{Arguments}
%
\begin{Value}
interpretation
\end{Value}
%
\begin{Examples}
\begin{ExampleCode}

## Not run: 
#'
# 예시 데이터 프레임
jutpls_boot1 <- data.frame(
  endogenous = c("진로동기", "진로태도", "진로준비", "진로준비", "진로준비"),
  exogenous = c("자기효능감", "자기효능감", "자기효능감", "진로동기", "진로태도"),
  Est = c(0.880, -0.674, 0.429, 0.446, 0.007),
  SE = c(0.050, 0.078, 0.129, 0.121, 0.078),
  t = c(17.597, -8.656, 3.330, 3.694, 0.087),
  p.value = c("<.001", "<.001", 0.001, "<.001", 0.931),
  sig = c("***", "***", "**", "***", "ns")
)

# 함수 실행
plspm_apa_paths(jutpls_boot)

jutpls_boot %>% plspm_paths_sig()%>% plspm_apa_paths()
#'
jutpls_boot %>% plspm_apa_paths()
jutpls_boot %>% plspm_paths_sig()%>% plspm_apa_paths()
jutpls_boot %>% plspm_paths_sig(unite=T)%>% plspm_apa_paths()



## End(Not run)

\end{ExampleCode}
\end{Examples}
\inputencoding{utf8}
\HeaderA{plspm\_auto\_blocks}{plspm\_auto\_blocks}{plspm.Rul.auto.Rul.blocks}
%
\begin{Description}
plspm\_auto\_blocks
\end{Description}
%
\begin{Usage}
\begin{verbatim}
plspm_auto_blocks(mat, to_lower = TRUE)
\end{verbatim}
\end{Usage}
%
\begin{Arguments}
\begin{ldescription}
\item[\code{mat}] path matrix

\item[\code{to\_lower}] Change to lowercase
\end{ldescription}
\end{Arguments}
%
\begin{Value}
data
\end{Value}
%
\begin{Examples}
\begin{ExampleCode}


## Not run: 

#'
mat <- matrix(c(0, 1, 0, 0, 0, 1, 0, 0,
                0, 0, 1, 1, 1, 0, 0, 0,
                0, 0, 0, 1, 0, 0, 0, 0,
                0, 0, 0, 1, 0, 0, 0, 0,
                0, 0, 0, 0, 0, 1, 0, 0,
                0, 0, 0, 0, 0, 0, 1, 1,
                0, 0, 0, 0, 0, 0, 0, 0,
                0, 0, 0, 0, 0, 0, 0, 0),
              nrow = 8, byrow = TRUE)

rownames(mat) <- c("IMAG", "EXPE", "QUAL", "VAL", "SAT", "LOY", "int_IMEX", "int_IMQU")
colnames(mat) <- c("IMAG", "EXPE", "QUAL", "VAL", "SAT", "LOY", "int_IMEX", "int_IMQU")
mat
# 함수 호출
plspm_auto_blocks(mat)

plspm_add_inter_paths(satpls_boot$model$IDM, list(
  list(from = "IMAG*EXPE", to = "LOY"),
  list(from = "IMAG*QUAL", to = "LOY")
))%>%plspm_auto_blocks()


## End(Not run)

\end{ExampleCode}
\end{Examples}
\inputencoding{utf8}
\HeaderA{plspm\_auto\_genData}{plspm\_auto\_genData}{plspm.Rul.auto.Rul.genData}
%
\begin{Description}
plspm\_auto\_genData
\end{Description}
%
\begin{Usage}
\begin{verbatim}
plspm_auto_genData(data, interactions, to_lower = TRUE)
\end{verbatim}
\end{Usage}
%
\begin{Arguments}
\begin{ldescription}
\item[\code{data}] data.frame

\item[\code{interactions}] interaction term list data

\item[\code{to\_lower}] Change to lowercase
\end{ldescription}
\end{Arguments}
%
\begin{Value}
data
\end{Value}
%
\begin{Examples}
\begin{ExampleCode}

## Not run: 
#'
# 예시 데이터
data1 <- data.frame(
  IMAG = c(-0.06448267, 0.83858784, 0.47542512, 0.42690176, 1.04176442, 0.27894142),
  EXPE = c(0.3442433, 0.7974660, 0.5531210, 0.1281051, 0.4492096, 0.6914954),
  QUAL = c(-0.5254008, 0.3971273, 0.1642778, -0.6892055, 0.7255953, -0.2088824),
  VAL = c(-0.07103969, 0.44321149, 0.56088128, -0.81127296, 1.21587873, -0.29482532),
  SAT = c(-0.20583024, 0.36908558, 0.83019100, -0.05045141, 0.53174538, -0.20212349),
  LOY = c(0.1664431, 0.4924859, 0.8117161, -0.7969487, 0.7348195, 0.1714641)
)

# 함수 실행 예시
plspm_auto_genData(data1, list(
  int1 = list(from = "IMAG*EXPE", to = "LOY"),
  int2 = list(from = "IMAG*QUAL", to = "LOY")
))

# 함수 실행 예시
plspm_auto_genData(satpls_boot$scores , list(
  int1 = list(from = "IMAG*EXPE", to = "LOY"),
  int2 = list(from = "IMAG*QUAL", to = "LOY")
)) %>%head()

# 함수 실행 예시
plspm_auto_genData(satpls_boot$scores , list(
  list(from = "IMAG*EXPE", to = "LOY"),
  list(from = "IMAG*QUAL", to = "LOY")
))%>%head()
#'

## End(Not run)


\end{ExampleCode}
\end{Examples}
\inputencoding{utf8}
\HeaderA{plspm\_blocks}{Function to create blocks------------}{plspm.Rul.blocks}
%
\begin{Description}
Function to create blocks------------
\end{Description}
%
\begin{Usage}
\begin{verbatim}
plspm_blocks(..., dataset = NULL, output = "match")
\end{verbatim}
\end{Usage}
%
\begin{Arguments}
\begin{ldescription}
\item[\code{...}] list data

\item[\code{dataset}] data.frame

\item[\code{output}] 'match', 'match\_colnames' is compare data.frame
\end{ldescription}
\end{Arguments}
%
\begin{Value}
list blocks
\end{Value}
%
\begin{Examples}
\begin{ExampleCode}

## Not run: 

#'
문자열 블록:
  블록의 모든 항목이 문자열이고 output이 "match"일 경우,
   블록을 그대로 유지합니다.

숫자 블록:
  블록의 모든 항목이 숫자이고 output이 "match"일 경우,
  블록을 숫자로 변환합니다.

"match_colnames" 출력 옵션:
  output이 "match_colnames"일 경우, 블록이 숫자인
  경우 해당 인덱스를 데이터셋의 컬럼 이름으로
  변환합니다. 변환된 컬럼 이름 중 데이터셋에
  존재하는 이름만 유지합니다.

기타:
#'   블록의 항목이 혼합되어 있거나 output이
"match_colnames"가 아닌 경우, 블록을 숫자로
변환하고 last_value를 기준으로 연속된 숫자로
변환합니다. 빈 블록의 경우 빈 정수 벡터를
반환합니다.
bl1 <- c("A", "B", "C")
bl2 <- c(1, 2, 3)
dataset <- data.frame(A = 1:5, B = 1:5, C = 1:5, D = 1:5)
dataset
# 함수 호출
plspm_blocks(block1 = bl1, block2 = bl2,
dataset = dataset, output = "match")
plspm_blocks(block1 = bl1, block2 = bl2,output = "match")
plspm_blocks(block1 = bl1, block2 = bl2,
dataset = dataset, output = "match_colnames")
plspm_blocks(block1 = bl1, block2 = bl2,
output = "match_colnames") #character(0)
#'
#'
#'
# selected indicators from satisfaction
satisfaction2 = satisfaction[, c(20:22, 24:26)]
# add dummy variables to satisfaction2
satisfaction2$dummy1 = dummy1
satisfaction2$dummy2 = dummy2
# add product terms to satisfaction2
satisfaction2$sat1m1 = satisfaction2$sat1 * dummy1
satisfaction2$sat2m1 = satisfaction2$sat2 * dummy1
satisfaction2$sat3m1 = satisfaction2$sat3 * dummy1
satisfaction2$sat1m2 = satisfaction2$sat1 * dummy2
satisfaction2$sat2m2 = satisfaction2$sat2 * dummy2
satisfaction2$sat3m2 = satisfaction2$sat3 * dummy2

satisfaction2 %>%str()
#'
#'

plspm_blocks(
  Satis = item(1:3),
  M1= item(7),
  SatisM1 = item(9:11),
  M2 = item(8),
  SatisM2 = item(12:14),
  Loyalty = item(4:6),
  dataset=satisfaction2,
  output="match_colnames"
)


plspm_blocks(
  Satis = 1:3,
  M1 = 7,
  SatisM1 = 9:11,
  M2 = 8,
  SatisM2 = 12:14,
  Loyalty = 4:6,
  dataset = satisfaction2,
  output = "match_colnames"
)

plspm_blocks(
  Satis = 1:3,
  M1 = 7,
  SatisM1 = 9:11,
  M2 = 8,
  SatisM2 = 12:14,
  Loyalty = 4:6,
  dataset = satisfaction2,
  output = "match"
)



#참고
plspm_blocks(
  Satis = item(1:3),
  M1= item(7),
  SatisM1 = item(9:11),
  M2 = item(8),
  SatisM2 = item(12:14),
  Loyalty = item(4:6),
  dataset=satisfaction2,
  output="match"
)



edu_blocks_named2 <- plspm_blocks(
  Support = item(1:4),
  Advising = item(5:8),
  Tutoring= item( 9:12),
  Value = item(13:16),
  Satisfaction = item(17:19),
  Loyalty =item(20:23),
  output = "match_colnames"
edu_blocks_named2

# to lavaan syntax
edu_blocks_named2 %>% plspm_blocks2lav()%>%cat("\n")

# # Example usage
edu_blocks_named1 <- plspm_blocks(
  Support = item("sup.help", "sup.under", "sup.safe", "sup.conc"),
  Advising = item("adv.comp", "adv.acces", "adv.comm", "adv.qual"),
  Tutoring = item("tut.prof", "tut.sched", "tut.stud", "tut.qual"),
  Value = item("val.devel", "val.deci", "val.meet", "val.info"),
  Satisfaction = item("sat.glad", "sat.expe", "sat.over"),
  Loyalty = item("loy.proud", "loy.recom", "loy.asha", "loy.back"),
  dataset = education , # Provide your dataset here
  output = "match_colnames"
)
#'
sat_blocks1 <- plspm_blocks(
  IMAG = item(paste0("imag",1:5)),
  EXPE = item(paste0("expe", 1:5)),
  QUAL = item( paste0("qual", 1:5)),
  VAL = item(paste0("val", 1:4)),
  SAT = item(paste0("sat", 1:4)),
  LOY = item(paste0("loy", 1:4))
)
sat_blocks1
##to lavaan syntax
sat_blocks1 %>% plspm_blocks2lav()%>%cat("\n")




## End(Not run)

\end{ExampleCode}
\end{Examples}
\inputencoding{utf8}
\HeaderA{plspm\_blocks2lav}{blocks to lavaan syntax}{plspm.Rul.blocks2lav}
%
\begin{Description}
blocks to lavaan syntax
\end{Description}
%
\begin{Usage}
\begin{verbatim}
plspm_blocks2lav(listdata, name = FALSE)
\end{verbatim}
\end{Usage}
%
\begin{Arguments}
\begin{ldescription}
\item[\code{listdata}] block list data

\item[\code{name}] defulat FALSE, TRUE generate blocks\_model=
\end{ldescription}
\end{Arguments}
%
\begin{Value}
lavaan syntax
\end{Value}
%
\begin{Examples}
\begin{ExampleCode}

## Not run: 
#'
listdata <- list(
  IMAG = c("imag1", "imag2", "imag3", "imag4", "imag5"),
  EXPE = c("expe1", "expe2", "expe3", "expe4", "expe5"),
  QUAL = c("qual1", "qual2", "qual3", "qual4", "qual5"),
  VAL = c("val1", "val2", "val3", "val4"),
  SAT = c("sat1", "sat2", "sat3", "sat4"),
  LOY = c("loy1", "loy2", "loy3", "loy4")
)
listdata
# Function calls and result output
plspm_blocks2lav(listdata)
plspm_blocks2lav(listdata)%>%cat("\n")

## End(Not run)

\end{ExampleCode}
\end{Examples}
\inputencoding{utf8}
\HeaderA{plspm\_blocks\_match}{Matching plspm blocks names}{plspm.Rul.blocks.Rul.match}
%
\begin{Description}
Matching plspm blocks names
\end{Description}
%
\begin{Usage}
\begin{verbatim}
plspm_blocks_match(data, blocks, paths = NULL, paths_name = NULL)
\end{verbatim}
\end{Usage}
%
\begin{Arguments}
\begin{ldescription}
\item[\code{data}] data

\item[\code{blocks}] plspm blocks

\item[\code{paths}] plspm paths

\item[\code{paths\_name}] path colnames
\end{ldescription}
\end{Arguments}
%
\begin{Value}
block data
\end{Value}
%
\begin{Examples}
\begin{ExampleCode}

## Not run: 

data(satisfaction)


# path matrix
IMAG = c(0,0,0,0,0,0)
EXPE = c(1,0,0,0,0,0)
QUAL = c(0,1,0,0,0,0)
VAL = c(0,1,1,0,0,0)
SAT = c(1,1,1,1,0,0)
LOY = c(1,0,0,0,1,0)
sat_path = rbind(IMAG, EXPE, QUAL, VAL, SAT, LOY)
colnames(sat_path) = rownames(sat_path)
sat_path

# blocks of outer model
sat_blocks = list(1:5, 6:10, 11:15, 16:19, 20:23, 24:27)
sat_blocks


#by hand Very Very tired !!!!
cnames = satisfaction%>%colnames()
latent = colnames(sat_path) #"IMAG", "EXPE", "QUAL", "VAL",  "SAT",  "LOY"
latent

list(IMAGE = cnames[1:5],
     EXPE = cnames[6:10],
     QUAL= cnames[11:15],
     VAL= cnames[16:19],
     SAT= cnames[20:23],
     LOY= cnames[24:27])

#' #Using function : all input
plspm_blocks_match(data=satisfaction, blocks= sat_blocks, paths=sat_path )

# nothing paths --> auto names
plspm_blocks_match(data= satisfaction, blocks=sat_blocks )

# Input paths_name
plspm_blocks_match(data= satisfaction, blocks=sat_blocks,
                   paths_name =c("IMAG", "EXPE", "QUAL", "VAL",  "SAT",  "LOY") )


sat_blocks1 = plspm_blocks_match(satisfaction, paths=sat_path, blocks=sat_blocks )
sat_blocks1

#YOu draw models
paste(plspm_paths2lav(sat_path) ,
      plspm_blocks2lav(sat_blocks1))%>%
  diagram_model(rotation = 1, sizeLat = 6, sizeMan = 4, sizeMan2=3,
                mar=c(2,2,2,2))

## End(Not run)

\end{ExampleCode}
\end{Examples}
\inputencoding{utf8}
\HeaderA{plspm\_boot\_effect}{plspm\_boot\_effect Total effect}{plspm.Rul.boot.Rul.effect}
%
\begin{Description}
plspm\_boot\_effect Total effect
\end{Description}
%
\begin{Usage}
\begin{verbatim}
plspm_boot_effect(plsres, accept = FALSE, type = "res", digits = 3)
\end{verbatim}
\end{Usage}
%
\begin{Arguments}
\begin{ldescription}
\item[\code{plsres}] plsdata

\item[\code{accept}] Hypothesis

\item[\code{type}] res, all

\item[\code{digits}] 3
\end{ldescription}
\end{Arguments}
%
\begin{Value}
data result
\end{Value}
%
\begin{Examples}
\begin{ExampleCode}

## Not run: 

library(plspm)
data(satisfaction)

# path matrix
IMAG = c(0,0,0,0,0,0)
EXPE = c(1,0,0,0,0,0)
QUAL = c(0,1,0,0,0,0)
VAL = c(0,1,1,0,0,0)
SAT = c(1,1,1,1,0,0)
LOY = c(1,0,0,0,1,0)
sat_path = rbind(IMAG, EXPE, QUAL, VAL, SAT, LOY)

# blocks of outer model
sat_blocks = list(1:5, 6:10, 11:15, 16:19, 20:23, 24:27)
sat_mod=c(rep("A", ncol(sat_path)))
# apply plspm
satpls = plspm(satisfaction, sat_path, sat_blocks, modes = sat_mod,
               scaled = FALSE)



  # bootstrap data
pls_boot_effect(satpls_boot)
pls_boot_effect(satpls_boot, accept=T)


## End(Not run)

\end{ExampleCode}
\end{Examples}
\inputencoding{utf8}
\HeaderA{plspm\_boot\_factor\_layout}{plspm\_boot\_factor\_layout}{plspm.Rul.boot.Rul.factor.Rul.layout}
%
\begin{Description}
plspm\_boot\_factor\_layout
\end{Description}
%
\begin{Usage}
\begin{verbatim}
plspm_boot_factor_layout(
  plsres_boot,
  positions = NULL,
  nrow = 4,
  ncol = 5,
  ...
)
\end{verbatim}
\end{Usage}
%
\begin{Arguments}
\begin{ldescription}
\item[\code{plsres\_boot}] bootresult

\item[\code{positions}] position change

\item[\code{nrow}] 4,

\item[\code{ncol}] 5

\item[\code{...}] add c()
\end{ldescription}
\end{Arguments}
%
\begin{Value}
matrix
\end{Value}
%
\begin{Examples}
\begin{ExampleCode}

## Not run: 
#'
# 예시 호출
plsres_boot <- list(model = list(gens = list(
                  lvs_names = c("LV1", "LV2", "LV3", "LV4"))))
result <- plspm_boot_factor_layout(plsres_boot)
print(result)

# 예시 호출
plsres_boot <- list(model =
list(gens = list(
lvs_names = c("LV1", "LV2", "LV3", "LV4", "LV5"))))

plspm_boot_factor_layout(plsres_boot)

# 예시 호출
plsres_boot <- list(model = list(gens =
list(lvs_names = c("LV1", "LV2", "LV3",
"LV4", "LV5", "LV6", "LV7","LV8"))))


plspm_boot_factor_layout(plsres_boot )
###EX)
positions <- list("LV1" = c(2, 3), "LV2" = c(1, 4))


## End(Not run)


\end{ExampleCode}
\end{Examples}
\inputencoding{utf8}
\HeaderA{plspm\_boot\_factor\_point\_to}{plspm\_boot\_factor\_point\_to}{plspm.Rul.boot.Rul.factor.Rul.point.Rul.to}
%
\begin{Description}
plspm\_boot\_factor\_point\_to
\end{Description}
%
\begin{Usage}
\begin{verbatim}
plspm_boot_factor_point_to(
  plsres_boot,
  point = NULL,
  positions = NULL,
  nrow = 4,
  ncol = 5
)
\end{verbatim}
\end{Usage}
%
\begin{Arguments}
\begin{ldescription}
\item[\code{plsres\_boot}] plsres\_boot

\item[\code{point}] point point = list(up = c(2,1))

\item[\code{positions}] positions positions = list(up = c(2,1))

\item[\code{nrow}] 3

\item[\code{ncol}] 3
\end{ldescription}
\end{Arguments}
%
\begin{Value}
data
\end{Value}
%
\begin{Examples}
\begin{ExampleCode}

## Not run: 
#'
plspm_boot_factor_point_to(jutpls_boot)
plspm_boot_factor_point_to(satpls_boot)
# plspm_factor_point_mat(jutpls_boot, point=list())
# point <- list(up = c(1, 2))
plspm_boot_factor_point_to(jutpls_boot, point = list(up = c(2,1)))
plspm_boot_factor_point_to(jutpls_boot, point = list(up = c(2,1), right =c(1,2)))
#'

## End(Not run)



\end{ExampleCode}
\end{Examples}
\inputencoding{utf8}
\HeaderA{plspm\_boot\_ind\_factor}{plspm\_indicator\_factor}{plspm.Rul.boot.Rul.ind.Rul.factor}
%
\begin{Description}
plspm\_indicator\_factor
\end{Description}
%
\begin{Usage}
\begin{verbatim}
plspm_boot_ind_factor(plsres_boot)
\end{verbatim}
\end{Usage}
%
\begin{Arguments}
\begin{ldescription}
\item[\code{plsres\_boot}] plsres\_boot
\end{ldescription}
\end{Arguments}
%
\begin{Value}
plsres\_boot
\end{Value}
%
\begin{Examples}
\begin{ExampleCode}
## Not run: 
plspm_indicator_factor(nfl_pls_boot)

## End(Not run)
\end{ExampleCode}
\end{Examples}
\inputencoding{utf8}
\HeaderA{plspm\_boot\_ind\_order}{plspm\_indicator\_order}{plspm.Rul.boot.Rul.ind.Rul.order}
%
\begin{Description}
plspm\_indicator\_order
\end{Description}
%
\begin{Usage}
\begin{verbatim}
plspm_boot_ind_order(plsres_boot)
\end{verbatim}
\end{Usage}
%
\begin{Arguments}
\begin{ldescription}
\item[\code{plsres\_boot}] plsres\_boot
\end{ldescription}
\end{Arguments}
%
\begin{Value}
plsres\_boot
\end{Value}
%
\begin{Examples}
\begin{ExampleCode}
## Not run: 
#'

## End(Not run)

\end{ExampleCode}
\end{Examples}
\inputencoding{utf8}
\HeaderA{plspm\_boot\_ind\_push}{plspm\_ind\_push}{plspm.Rul.boot.Rul.ind.Rul.push}
%
\begin{Description}
plspm\_ind\_push
\end{Description}
%
\begin{Usage}
\begin{verbatim}
plspm_boot_ind_push(pls_boot, set = 2)
\end{verbatim}
\end{Usage}
%
\begin{Arguments}
\begin{ldescription}
\item[\code{pls\_boot}] pls\_boot

\item[\code{set}] set=1.5 or c(set= c(2,2,1.5,1.5))
\end{ldescription}
\end{Arguments}
%
\begin{Value}
vector
\end{Value}
%
\begin{Examples}
\begin{ExampleCode}

## Not run: 

# 예시 사용법
jutpls_boot <- list(model = list(gens = list(lvs_names = c("자기효능감", "진로동기", "진로태도", "진로준비"))))

# 기본값 사용하는 경우
result_default <- plspm_ind_push(jutpls_boot)
print(result_default)

# 이름과 값이 있는 경우
result_named <- plspm_ind_push(jutpls_boot, set = c(진로동기=2, 자기효능감=3, 진로준비=1.5, 진로태도=2.5))
print(result_named)

# 값만 있는 경우
result_unnamed <- plspm_ind_push(jutpls_boot, set = c(2, 2, 1.5, 1.5))
print(result_unnamed)

# 부분적으로 값이 있는 경우
result_partial <- plspm_ind_push(jutpls_boot, set = c(진로동기=2, 자기효능감=3))
print(result_partial)


## End(Not run)



\end{ExampleCode}
\end{Examples}
\inputencoding{utf8}
\HeaderA{plspm\_boot\_paths\_sig}{boot res sig table}{plspm.Rul.boot.Rul.paths.Rul.sig}
%
\begin{Description}
boot res sig table
\end{Description}
%
\begin{Usage}
\begin{verbatim}
plspm_boot_paths_sig(plsres, type = "df", digits = 3)
\end{verbatim}
\end{Usage}
%
\begin{Arguments}
\begin{ldescription}
\item[\code{plsres}] plspm result

\item[\code{type}] df, vec, vector

\item[\code{digits}] digits=3
\end{ldescription}
\end{Arguments}
%
\begin{Value}
data, vector
\end{Value}
\inputencoding{utf8}
\HeaderA{plspm\_boot\_paths\_sig\_vec}{boot res sig table}{plspm.Rul.boot.Rul.paths.Rul.sig.Rul.vec}
%
\begin{Description}
boot res sig table
\end{Description}
%
\begin{Usage}
\begin{verbatim}
plspm_boot_paths_sig_vec(plsres, type = "vec")
\end{verbatim}
\end{Usage}
%
\begin{Arguments}
\begin{ldescription}
\item[\code{plsres}] plspm result

\item[\code{type}] df, vec, vector
\end{ldescription}
\end{Arguments}
%
\begin{Value}
data, vector
\end{Value}
\inputencoding{utf8}
\HeaderA{plspm\_boot\_summary}{plspm\_boot\_summary}{plspm.Rul.boot.Rul.summary}
%
\begin{Description}
plspm\_boot\_summary
\end{Description}
%
\begin{Usage}
\begin{verbatim}
plspm_boot_summary(pls_boot, type = "res")
\end{verbatim}
\end{Usage}
%
\begin{Arguments}
\begin{ldescription}
\item[\code{pls\_boot}] pls boot data

\item[\code{type}] res, all, weight, loadings, paths, rsq, total.efs
\end{ldescription}
\end{Arguments}
%
\begin{Value}
data
\end{Value}
%
\begin{Examples}
\begin{ExampleCode}

## Not run: 

jutpls_boot%>%plspm_boot_summary()
jutpls_boot$boot%>%plspm_boot_summary()


## End(Not run)


\end{ExampleCode}
\end{Examples}
\inputencoding{utf8}
\HeaderA{plspm\_cfa}{plspm CFA}{plspm.Rul.cfa}
%
\begin{Description}
plspm CFA
\end{Description}
%
\begin{Usage}
\begin{verbatim}
plspm_cfa(plsres_boot, type = "all", axis_x = 1.2)
\end{verbatim}
\end{Usage}
%
\begin{Arguments}
\begin{ldescription}
\item[\code{plsres\_boot}] bootstrap plspm data

\item[\code{type}] all, cfa1, cfa2, item, CRAVE, fl1981, htmt, f2. total\_effect

\item[\code{axis\_x}] 1.2 setting efffecbar text size
\end{ldescription}
\end{Arguments}
%
\begin{Value}
multiple data
\end{Value}
%
\begin{Examples}
\begin{ExampleCode}

## Not run: 

data(satisfaction)
sat_path = plspm_paths(
  row_names = c("IMAG","EXPE","QUAL","VAL","SAT","LOY"),
  relationship = list(
    path(from="IMAG", to=c("EXPE","SAT","LOY")),
    path("EXPE", c("QUAL","VAL","SAT")),
    path("QUAL", c("VAL","SAT")),
    path("VAL",c("SAT")),
    path("SAT","LOY"))
)
sat_path

# blokcs
sat_blocks1 <- plspm_blocks(
  IMAG = item(paste0("imag",1:5)),
  EXPE = item(paste0("expe", 1:5)),
  QUAL = item( paste0("qual", 1:5)),
  VAL = item(paste0("val", 1:4)),
  SAT = item(paste0("sat", 1:4)),
  LOY = item(paste0("loy", 1:4))
)

sat_blocks1

# apply plspm
# satpls = plspm(satisfaction,
#                path_matrix = sat_path,
#                blocks = sat_blocks1,
#                scaled = FALSE)

satpls_boot = plspm_sem(satisfaction, sat_path, sat_blocks1, scaled = FALSE,
                        boot.val =TRUE, br=100)
##
satpls_boot%>%plspm_cfa()
satpls_boot%>%plspm_cfa("res")

satpls_boot%>%plspm_cfa(type="cfa1")
satpls_boot%>%plspm_cfa(type="cfa2")
satpls_boot%>%plspm_cfa(type="cor")
satpls_boot%>%plspm_cfa(type="model")
satpls_boot%>%plspm_cfa(type="inner_model") %>%plspm_inner_model_sig()
satpls_boot%>%plspm_cfa(type="outer_model")
satpls_boot%>%plspm_cfa(type="effectbar")
total

## End(Not run)


\end{ExampleCode}
\end{Examples}
\inputencoding{utf8}
\HeaderA{plspm\_change\_nodeLabels}{plspm\_change\_nodeLabels}{plspm.Rul.change.Rul.nodeLabels}
%
\begin{Description}
plspm\_change\_nodeLabels
\end{Description}
%
\begin{Usage}
\begin{verbatim}
plspm_change_nodeLabels(
  data,
  ...,
  label.cex = 1.2,
  label.scale = TRUE,
  label.prop = 1,
  label.norm = "width",
  plot = TRUE
)
\end{verbatim}
\end{Usage}
%
\begin{Arguments}
\begin{ldescription}
\item[\code{data}] plspm\_semPaths

\item[\code{...}] change node from name , to name

\item[\code{label.cex}] label.cex

\item[\code{label.scale}] label.scale

\item[\code{label.prop}] label.prop

\item[\code{label.norm}] label.norm

\item[\code{plot}] TRUE
\end{ldescription}
\end{Arguments}
%
\begin{Value}
plot
\end{Value}
%
\begin{Examples}
\begin{ExampleCode}

## Not run: 


pathmodel = "
자기효능감 =~ C_S1 +  C_S2 + C_S3 + C_S4 + C_S5
진로태도 =~ B02 + B04 + B10 + B14
진로동기 =~ A_M1 + A_M2 + A_M3
진로준비 =~ D_P1 + D_P2 + D_P3

진로동기 ~ H1*자기효능감
진로태도 ~  H2*자기효능감
진로준비 ~ H5 *진로태도
진로준비 ~  H4 * 진로동기
진로준비 ~  H3*자기효능감
"

plspm_semPaths2(jutpls_boot,   curve=1, layout="tree2",rotation=1,
                edge.label.cex=0.7, mar=c(2,1,2,1),
                edge.label.position=0.6, label.cex=1,
                sizeLat= 10)%>%
  plspm_change_nodeLabels("B02","결정성",
                          "B04", "독립성",
                          "B10", "계획성",
                          "B14", "성향"
  )








## End(Not run)


\end{ExampleCode}
\end{Examples}
\inputencoding{utf8}
\HeaderA{plspm\_col2lav}{plspm\_col2lav, path, or measurement}{plspm.Rul.col2lav}
%
\begin{Description}
plspm\_col2lav, path, or measurement
\end{Description}
%
\begin{Usage}
\begin{verbatim}
plspm_col2lav(data, lhs, rhs, coef, op = "=~")
\end{verbatim}
\end{Usage}
%
\begin{Arguments}
\begin{ldescription}
\item[\code{data}] plspm\$boot\$loadings

\item[\code{lhs}] lhs

\item[\code{rhs}] rhs

\item[\code{coef}] coeff

\item[\code{op}] "=\textasciitilde{}","\textasciitilde{}"
\end{ldescription}
\end{Arguments}
%
\begin{Value}
lavaan lavaan syntac
\end{Value}
%
\begin{Examples}
\begin{ExampleCode}

## Not run: 
#'
satpls_boot <- data.frame(
  block = c("IMAG", "IMAG", "IMAG",
  "IMAG", "IMAG", "EXPE", "EXPE",
  "EXPE", "EXPE", "EXPE", "QUAL",
  "QUAL", "QUAL", "QUAL", "QUAL",
  "VAL", "VAL", "VAL", "VAL", "SAT",
  "SAT", "SAT", "SAT", "LOY", "LOY",
  "LOY", "LOY"),
  name = c("imag1", "imag2", "imag3",
   "imag4", "imag5", "expe1", "expe2",
    "expe3", "expe4", "expe5", "qual1",
     "qual2", "qual3", "qual4", "qual5",
     "val1", "val2", "val3", "val4",
     "sat1", "sat2", "sat3", "sat4",
     "loy1", "loy2", "loy3", "loy4"),
  Original = c(0.709, 0.877, 0.842,
   0.569, 0.778, 0.766, 0.837, 0.76,
   0.718, 0.837, 0.781, 0.882, 0.794,
    0.789, 0.807, 0.865, 0.796, 0.75,
    0.844, 0.92, 0.916, 0.826, 0.818,
     0.907, 0.671, 0.905, 0.682)
)
#data
satpls_boot

#method 1
satpls_boot %>%
plspm_col2lav(lhs = "block", rhs = "name", coef = "Original")

#method 2
satpls_boot %>% plspm_col2lav(1,2,3)

#ethod 3
satpls_boot %>%
plspm_col2lav( lhs = 1, rhs = 2, coef = 3)

# draw model
satpls_boot %>%
plspm_col2lav( lhs = 1, rhs = 2, coef = 3)%>%
  diagram_model(whatLabels = "est", exoCov = F,
  sizeLat=6, sizeMan=6,sizeMan2=2,
                edge.label.cex = 0.6)



## End(Not run)

\end{ExampleCode}
\end{Examples}
\inputencoding{utf8}
\HeaderA{plspm\_cor}{Tplspm result structral cor}{plspm.Rul.cor}
%
\begin{Description}
Tplspm result structral cor
\end{Description}
%
\begin{Usage}
\begin{verbatim}
plspm_cor(plsres)
\end{verbatim}
\end{Usage}
%
\begin{Arguments}
\begin{ldescription}
\item[\code{plsres}] plspm result
\end{ldescription}
\end{Arguments}
%
\begin{Value}
matrix
\end{Value}
\inputencoding{utf8}
\HeaderA{plspm\_CRAVE}{plspm\_CRAVE}{plspm.Rul.CRAVE}
%
\begin{Description}
plspm\_CRAVE
\end{Description}
%
\begin{Usage}
\begin{verbatim}
plspm_CRAVE(plsdata, digits = 3)
\end{verbatim}
\end{Usage}
%
\begin{Arguments}
\begin{ldescription}
\item[\code{plsdata}] plspmdata

\item[\code{digits}] digists 3
\end{ldescription}
\end{Arguments}
%
\begin{Value}
CR and AVE
\end{Value}
%
\begin{Examples}
\begin{ExampleCode}

## Not run: 

data(satisfaction)
#model 1
model1 = "
EXPE ~ IMAG
QUAL ~ EXPE
VAL ~  QUAL+ EXPE
SAT ~ IMAG + EXPE + QUAL + VAL
LOY ~ SAT + IMAG
"
# # Generate the matrix
sat_path = plspm_lav2path(model1,
                          fct_order=c("IMAG", "EXPE","QUAL", "VAL", "SAT", "LOY" ))
# sat_path

# Method 2: paths
sat_path = plspm_paths(
  row_names = c("IMAG","EXPE","QUAL","VAL","SAT","LOY"),
  relationship = list(
    path(from="IMAG", to=c("EXPE","SAT","LOY")),
    path("EXPE", c("QUAL","VAL","SAT")),
    path("QUAL", c("VAL","SAT")),
    path("VAL",c("SAT")),
    path("SAT","LOY")
  )
)
# blokcs
sat_blocks1 <- plspm_blocks(
  IMAG = item(paste0("imag",1:5)),
  EXPE = item(paste0("expe", 1:5)),
  QUAL = item( paste0("qual", 1:5)),
  VAL = item(paste0("val", 1:4)),
  SAT = item(paste0("sat", 1:4)),
  LOY = item(paste0("loy", 1:4))
)
sat_blocks1

# vector of modes (reflective indicators):auto
sat_mod = rep("A", 6)
#############

# apply plspm
satpls = plspm(satisfaction,
                   path_matrix = sat_path,
                   blocks = sat_blocks1,
                   scaled = FALSE)

satpls_boot = plspm_sem(satisfaction, sat_path, sat_blocks1, scaled = FALSE,
                        boot.val =TRUE, br=100)

#CR a AVE
satpls%>% plspm_CRAVE()

satpls_summary%>% plspm_CRAVE()

satpls_boot%>% plspm_CRAVE()

satpls_boot_summary%>% plspm_CRAVE()
#'

## End(Not run)



\end{ExampleCode}
\end{Examples}
\inputencoding{utf8}
\HeaderA{plspm\_crossloadings}{plspm\_crossloadings}{plspm.Rul.crossloadings}
%
\begin{Description}
plspm\_crossloadings
\end{Description}
%
\begin{Usage}
\begin{verbatim}
plspm_crossloadings(plsres, digits = 3)
\end{verbatim}
\end{Usage}
%
\begin{Arguments}
\begin{ldescription}
\item[\code{plsres}] boot result

\item[\code{digits}] reound 3
\end{ldescription}
\end{Arguments}
%
\begin{Value}
data
\end{Value}
\inputencoding{utf8}
\HeaderA{plspm\_crossloadings\_plot}{plspm\_crossloadings\_plot}{plspm.Rul.crossloadings.Rul.plot}
%
\begin{Description}
plspm\_crossloadings\_plot
\end{Description}
%
\begin{Usage}
\begin{verbatim}
plspm_crossloadings_plot(plsres, wide = F)
\end{verbatim}
\end{Usage}
%
\begin{Arguments}
\begin{ldescription}
\item[\code{plsres}] plspm result

\item[\code{wide}] plot type TRUE block, FALSE wide plot
\end{ldescription}
\end{Arguments}
%
\begin{Value}
pplot
\end{Value}
%
\begin{Examples}
\begin{ExampleCode}


## Not run: 

data(satisfaction)
#model 1
model1 = "
EXPE ~ IMAG
QUAL ~ EXPE
VAL ~  QUAL+ EXPE
SAT ~ IMAG + EXPE + QUAL + VAL
LOY ~ SAT + IMAG
"
# # Generate the matrix
sat_path = plspm_lav2path(model1,
                          fct_order=c("IMAG", "EXPE","QUAL", "VAL", "SAT", "LOY" ))
# sat_path

# Method 2: paths
sat_path = plspm_paths(
  row_names = c("IMAG","EXPE","QUAL","VAL","SAT","LOY"),
  relationship = list(
    path(from="IMAG", to=c("EXPE","SAT","LOY")),
    path("EXPE", c("QUAL","VAL","SAT")),
    path("QUAL", c("VAL","SAT")),
    path("VAL",c("SAT")),
    path("SAT","LOY")
  )
)
# blokcs
sat_blocks1 <- plspm_blocks(
  IMAG = item(paste0("imag",1:5)),
  EXPE = item(paste0("expe", 1:5)),
  QUAL = item( paste0("qual", 1:5)),
  VAL = item(paste0("val", 1:4)),
  SAT = item(paste0("sat", 1:4)),
  LOY = item(paste0("loy", 1:4))
)
sat_blocks1

# vector of modes (reflective indicators):auto
sat_mod = rep("A", 6)
#############

# apply plspm
satpls = plspm(satisfaction,
                   path_matrix = sat_path,
                   blocks = sat_blocks1,
                   scaled = FALSE)

satpls_boot = plspm_sem(satisfaction, sat_path, sat_blocks1, scaled = FALSE,
                        boot.val =TRUE, br=100)


#'
satpls %>% plspm_crossloadings_plot()
satpls %>% plspm_crossloadings_plot(T)
satpls$crossloadings %>% plspm_crossloadings_plot()
satpls$crossloadings %>% plspm_crossloadings_plot(T)
#'
#'

## End(Not run)

\end{ExampleCode}
\end{Examples}
\inputencoding{utf8}
\HeaderA{plspm\_curve}{plspm\_curve}{plspm.Rul.curve}
%
\begin{Description}
plspm\_curve
\end{Description}
%
\begin{Usage}
\begin{verbatim}
plspm_curve(plotdata, curve_list = NULL, plot = TRUE)
\end{verbatim}
\end{Usage}
%
\begin{Arguments}
\begin{ldescription}
\item[\code{plotdata}] plot data

\item[\code{curve\_list}] curve list

\item[\code{plot}] TRUE
\end{ldescription}
\end{Arguments}
%
\begin{Value}
plot data
\end{Value}
%
\begin{Examples}
\begin{ExampleCode}

## Not run: 
#'

pathmodel = "
자기효능감 =~ C_S1 +  C_S2 + C_S3 + C_S4 + C_S5
진로태도 =~ B02 + B04 + B10 + B14
진로동기 =~ A_M1 + A_M2 + A_M3
진로준비 =~ D_P1 + D_P2 + D_P3

진로동기 ~ H1*자기효능감
진로태도 ~  H2*자기효능감
진로준비 ~ H5 *진로태도
진로준비 ~  H4 * 진로동기
진로준비 ~  H3*자기효능감
"

plspm_diagram_model(pathmodel, layout="tree2",
                    rotation = 2,residuals = T, residScale = 6,
                    edge.label.position = 0.6)%>%
  plspm_rotate_resid(c(진로동기= 160, 진로태도= 160))

#
plspm_diagram_model(pathmodel, layout="tree2",
                    rotation = 2,residuals = T, residScale = 6,
                    edge.label.position = 0.6)%>%
  plspm_curve(c("진로준비~~ 진로동기" = -2))%>%
  plspm_rotate_resid(c(진로동기= 160, 진로태도= 30))

## End(Not run)

\end{ExampleCode}
\end{Examples}
\inputencoding{utf8}
\HeaderA{plspm\_diagram\_model}{Drawing Structural Equation Models Using the lavaan Syntax (park JH PhD)}{plspm.Rul.diagram.Rul.model}
%
\begin{Description}
Drawing Structural Equation Models Using the lavaan Syntax (park JH PhD)
\end{Description}
%
\begin{Usage}
\begin{verbatim}
plspm_diagram_model(
  model,
  residuals = FALSE,
  whatLabels = "model",
  rotation = 2,
  residScale = 12,
  layout = "tree2",
  edge.label.cex = 1,
  edge.label.position = 0.5,
  edge.color = "black",
  edge.width = 1.5,
  sizeMan = 8,
  sizeMan2 = 3,
  sizeLat = 10,
  sizeLat2 = 4,
  style = "lisrel",
  sig = FALSE,
  exoVar = FALSE,
  exoCov = FALSE,
  curve = 1.5,
  asize = 2,
  mar = c(3, 3, 3, 3),
  nDigits = 3,
  shapeLat = "circle",
  shapeMan = "rectangle",
  shapeInt = "triangle",
  sample.nobs = 100,
  border.width = 2,
  Groups = FALSE,
  group_match = "lat",
  growth = FALSE,
  structural = FALSE,
  edgeLabels = NULL,
  nodeLabels = NULL,
  fixedStyle = 1,
  color = list(lat = "skyblue", man = "yellow"),
  type = "plot"
)
\end{verbatim}
\end{Usage}
%
\begin{Arguments}
\begin{ldescription}
\item[\code{model}] lavaan syntax

\item[\code{residuals}] residauls TRUE

\item[\code{whatLabels}] "model

\item[\code{rotation}] rotation=2

\item[\code{residScale}] residScale=12

\item[\code{layout}] layout="tree"

\item[\code{edge.label.cex}] 1

\item[\code{edge.label.position}] 0.5

\item[\code{edge.color}] edge.color= black

\item[\code{edge.width}] edge.width=1.5

\item[\code{sizeMan}] 8

\item[\code{sizeMan2}] 4

\item[\code{sizeLat}] 10

\item[\code{sizeLat2}] 6

\item[\code{style}] lisrel

\item[\code{sig}] FALSE

\item[\code{exoVar}] FALSE

\item[\code{exoCov}] TRUE

\item[\code{curve}] 1.5

\item[\code{asize}] 1.5

\item[\code{mar}] c(2,8,3,10)

\item[\code{nDigits}] 2

\item[\code{shapeLat}] "circle, ellipse

\item[\code{shapeMan}] retangle

\item[\code{shapeInt}] triangle ,

\item[\code{sample.nobs}] sample.nobs = 100

\item[\code{border.width}] border.width=1

\item[\code{Groups}] FALSE

\item[\code{group\_match}] "FALSE"lat\#'

\item[\code{growth}] growth= FALSE, TRUE  model.type ="grouwth"

\item[\code{structural}] structural = FALSE

\item[\code{edgeLabels}] edgeLabels=NULL

\item[\code{nodeLabels}] nodeLabels=NULL

\item[\code{fixedStyle}] fixedStyle=1 default rbase linetype

\item[\code{color}] lat= skyblue, man yellow

\item[\code{type}] type ="plot, and res is lavaan simdata result
\end{ldescription}
\end{Arguments}
%
\begin{Value}
plot
\end{Value}
%
\begin{Examples}
\begin{ExampleCode}
## Not run: 
#'
model_test <- '
  # latent variable definitions
     ind60 =~ x1 + x2 + x3
     dem60 =~ y1 + a*y2 + b*y3 + c*y4
     dem65 =~ y5 + a*y6 + b*y7 + c*y8

  # regressions
    dem60 ~ ind60
    dem65 ~ ind60 + dem60

  # residual correlations
    # y1 ~~ y5
    # y2 ~~ y4 + y6
    # y3 ~~ y7
    # y4 ~~ y8
    # y6 ~~ y8
'

HSmodel <- ' visual  =~ x1 + x2 + x3
             textual =~ x4 + x5 + x6
             speed   =~ x7 + x8 + x9 '


m1a <- "
y1 =~ x1 + x2 +x5
y2 =~ x3 + x4 +x6
y1 ~~ y2"

diagram_model(model= HSmodel, groups=F,  residuals=F)

diagram_model(model= model_test, style="ram", groups=F)
diagram_model(model= model_test, residuals=T)
diagram_model(model= m1a, residuals=T)


#plspm model
#model 1
model1 = "
EXPE ~ IMAG
QUAL ~ EXPE
VAL ~  QUAL+ EXPE
SAT ~ IMAG + EXPE + QUAL + VAL
LOY ~ SAT + IMAG
"
sat_path = plspm_lav2path(model1,
                          fct_order=c("IMAG", "EXPE","QUAL", "VAL", "SAT", "LOY" ))
sat_path

blocks_model ="
IMAG =~ imag1 + imag2 + imag3 + imag4 + imag5
EXPE =~ expe1 + expe2 + expe3 + expe4 + expe5
QUAL =~ qual1 + qual2 + qual3 + qual4 + qual5
VAL =~ val1 + val2 + val3 + val4
SAT =~ sat1 + sat2 + sat3 + sat4
LOY =~ loy1 + loy2 + loy3 + loy4
"
blocks_model
plspm_lav2blocks(blocks_model)

#model 그리기
paste(model1, blocks_model)%>%

  plspm_diagram_model(rotation = 1, sizeLat = 6, sizeMan = 4, sizeMan2=3,
                mar=c(2,2,2,2))

#'
#'
#'

## End(Not run)
\end{ExampleCode}
\end{Examples}
\inputencoding{utf8}
\HeaderA{plspm\_dULS\_modelfit}{Bollen-Stine Bootstrap Based d\_ULS and d\_G Calculation}{plspm.Rul.dULS.Rul.modelfit}
%
\begin{Description}
This function calculates the d\_ULS and d\_G metrics for a PLS model using the Bollen-Stine bootstrap method.
Confidence intervals are estimated for these indices.
\end{Description}
%
\begin{Usage}
\begin{verbatim}
plspm_dULS_modelfit(pls_model, blocks, bootstrap = 1000, ci = 0.95)
\end{verbatim}
\end{Usage}
%
\begin{Arguments}
\begin{ldescription}
\item[\code{pls\_model}] A PLS model object containing model results.

\item[\code{blocks}] A list defining the measurement model blocks. Each block should contain variable names.

\item[\code{bootstrap}] The number of bootstrap samples to use. Default is 1000.

\item[\code{ci}] The confidence level for the intervals. Default is 0.95.
\end{ldescription}
\end{Arguments}
%
\begin{Value}
A data frame with the calculated d\_ULS and d\_G metrics, along with their confidence intervals and significance indicators.
The columns include:
\begin{itemize}

\item{} \code{Method}: Bootstrap method used (e.g., "Bollen-Stine").
\item{} \code{d\_ULS}: Computed d\_ULS value.
\item{} \code{d\_ULS\_95CI}: Confidence interval for d\_ULS.
\item{} \code{dUsig}: Significance of d\_ULS (whether it falls within the confidence interval).
\item{} \code{d\_G}: Computed d\_G value.
\item{} \code{d\_G\_95CI}: Confidence interval for d\_G.
\item{} \code{dGsig}: Significance of d\_G (whether it falls within the confidence interval).

\end{itemize}

\end{Value}
%
\begin{Examples}
\begin{ExampleCode}
## Not run: 
# Example call with 1000 bootstrap samples
result <- plspm_dULS_modelfit(pls_model = edt_pls, blocks = edblock, bootstrap = 1000, ci = 0.95)
print(result)

## End(Not run)

\end{ExampleCode}
\end{Examples}
\inputencoding{utf8}
\HeaderA{plspm\_edges\_coefs}{plspm의 구조모형의 유의성 자료를 넣는 함수}{plspm.Rul.edges.Rul.coefs}
%
\begin{Description}
plspm의 구조모형의 유의성 자료를 넣는 함수
\end{Description}
%
\begin{Usage}
\begin{verbatim}
plspm_edges_coefs(plsdata, digits = 3, unite = TRUE, type = "df")
\end{verbatim}
\end{Usage}
%
\begin{Arguments}
\begin{ldescription}
\item[\code{plsdata}] plsdata

\item[\code{digits}] round 3

\item[\code{unite}] data combine unite TRUE -> estimat + sig

\item[\code{type}] vec, df
\end{ldescription}
\end{Arguments}
%
\begin{Value}
estimates vectors
\end{Value}
%
\begin{Examples}
\begin{ExampleCode}

## Not run: 

data(satisfaction)
#model 1
model1 = "
EXPE ~ IMAG
QUAL ~ EXPE
VAL ~  QUAL+ EXPE
SAT ~ IMAG + EXPE + QUAL + VAL
LOY ~ SAT + IMAG
"
# # Generate the matrix
sat_path = plspm_lav2path(model1,
                          fct_order=c("IMAG", "EXPE","QUAL", "VAL", "SAT", "LOY" ))
# sat_path

# Method 2: paths
sat_path = plspm_paths(
  row_names = c("IMAG","EXPE","QUAL","VAL","SAT","LOY"),
  relationship = list(
    path(from="IMAG", to=c("EXPE","SAT","LOY")),
    path("EXPE", c("QUAL","VAL","SAT")),
    path("QUAL", c("VAL","SAT")),
    path("VAL",c("SAT")),
    path("SAT","LOY")
  )
)
# blokcs
sat_blocks1 <- plspm_blocks(
  IMAG = item(paste0("imag",1:5)),
  EXPE = item(paste0("expe", 1:5)),
  QUAL = item( paste0("qual", 1:5)),
  VAL = item(paste0("val", 1:4)),
  SAT = item(paste0("sat", 1:4)),
  LOY = item(paste0("loy", 1:4))
)
sat_blocks1

# vector of modes (reflective indicators):auto
sat_mod = rep("A", 6)
#############

# apply plspm
satpls = plspm_sem(satisfaction,
                   path_matrix = sat_path,
                   blocks = sat_blocks1,
                   scaled = FALSE)

satpls_boot = plspm_sem(satisfaction, sat_path, sat_blocks1, scaled = FALSE,
                        boot.val =TRUE, br=100)

plspm_edges_values(satpls)
plspm_edges_values(satpls, unite=F) #no star
plspm_edges_values(satpls$inner_model)
plspm_edges_values(satpls$inner_model, unite=FALSE) #No star
#' #'

lay_p = matrix(c('IMA', NA, NA, NA, NA, NA, NA, NA,
                 NA, NA, NA, NA, NA, NA, NA, NA, NA, NA, NA, 'EXP', NA,
                 NA, NA, NA, NA, NA, NA, NA, NA, NA, NA, NA, NA, NA, 'QUA',
                 NA, NA, NA, NA, NA, NA, NA, NA, NA, NA, NA, NA, NA, NA, NA,
                 NA, NA, NA, NA, NA, NA, NA, NA, NA, 'VAL', NA, NA, NA, NA,
                 NA, NA, NA, NA, NA, NA, NA, NA, NA, NA, 'SAT', NA, NA, NA,
                 NA, NA, NA, NA, NA, NA, NA, NA, NA, NA, NA, NA, NA, 'LOY',
                 NA, NA, NA, NA, NA, NA, NA, NA),
               nrow = 10, ncol = 10, byrow = FALSE)
lay_p
#  Specify a new value for the edgeLabels parameter when calling qgraph
qgraph(qgraph(t(satpls$path_coefs)),
       layout = lay_p,
       posCol = "black",
       fade = FALSE,
       edge.labels = plspm_edges_values(satpls),
       edge.label.cex = 1,
       edge.label.position = 0.4,
       groups = NULL,
       palette = 'pastel'
)
# Specify a new value for the edgeLabels parameter when calling qgraph
qgraph(qgraph(t(satpls$path_coefs)),
       layout = "spring",
       posCol = "black",
       fade = FALSE,
       edge.labels = plspm_edges_values(satpls),
       edge.label.cex = 1,
       edge.label.position = 0.4,
       groups = NULL,
       palette = 'pastel'
)

 Specify a new value for the edgeLabels parameter when calling qgraph
qgraph(qgraph(t(satpls$path_coefs)),
       layout = lay_p,
       posCol = "black",
       fade = FALSE,
       edge.labels = matrix(sprintf("%.4f", t(satpls$path_coefs)),
                            nrow=nrow(t(satpls$path_coefs)),
                            ncol=ncol(t(satpls$path_coefs))),
       edge.label.cex = 1,
       edge.label.position=0.4,
       groups = NULL, palette='pastel' )

 plspm_edges_values(satpls)
satpls %>% plspm_path_coefs_plot()
satpls %>% plspm_path_coefs_plot(node =TRUE)



## End(Not run)
\end{ExampleCode}
\end{Examples}
\inputencoding{utf8}
\HeaderA{plspm\_edges\_loadings}{plspm lodaings significant  vecgtor or df}{plspm.Rul.edges.Rul.loadings}
%
\begin{Description}
plspm lodaings significant  vecgtor or df
\end{Description}
%
\begin{Usage}
\begin{verbatim}
plspm_edges_loadings(plsres, type = "df")
\end{verbatim}
\end{Usage}
%
\begin{Arguments}
\begin{ldescription}
\item[\code{plsres}] plspm result

\item[\code{type}] df, vec, data.frame, vector
\end{ldescription}
\end{Arguments}
%
\begin{Value}
data
\end{Value}
%
\begin{Examples}
\begin{ExampleCode}

## Not run: 

# Method 2: paths
sat_path = plspm_paths(
  row_names = c("IMAG","EXPE","QUAL","VAL","SAT","LOY"),
  relationship = list(
    path(from="IMAG", to=c("EXPE","SAT","LOY")),
    path("EXPE", c("QUAL","VAL","SAT")),
    path("QUAL", c("VAL","SAT")),
    path("VAL",c("SAT")),
    path("SAT","LOY")
  )
)
# blokcs
sat_blocks1 <- plspm_blocks(
  IMAG = item(paste0("imag",1:5)),
  EXPE = item(paste0("expe", 1:5)),
  QUAL = item( paste0("qual", 1:5)),
  VAL = item(paste0("val", 1:4)),
  SAT = item(paste0("sat", 1:4)),
  LOY = item(paste0("loy", 1:4))
)
sat_blocks1

# vector of modes (reflective indicators):auto
sat_mod = rep("A", 6)
#############

# apply plspm
satpls = plspm_sem(satisfaction,
                   path_matrix = sat_path,
                   blocks = sat_blocks1,
                   scaled = FALSE)

satpls_boot = plspm_sem(satisfaction, sat_path, sat_blocks1, scaled = FALSE,
                        boot.val =TRUE, br=100)




# plspm_loadings_values
satpls %>%plspm_edge_loadings()
satpls %>%plspm_edge_loadings("vec")
satpls$boot %>%plspm_edge_loadings()
satpls %>%plspm_edge_loadings()
satpls$boot$loadings%>%plspm_edge_loadings()
satpls$boot$loadings%>%plspm_edge_loadings("vec")


## End(Not run)
\end{ExampleCode}
\end{Examples}
\inputencoding{utf8}
\HeaderA{plspm\_edges\_values}{plspm의 구조모형의 유의성 자료를 넣는 함수}{plspm.Rul.edges.Rul.values}
%
\begin{Description}
plspm의 구조모형의 유의성 자료를 넣는 함수
\end{Description}
%
\begin{Usage}
\begin{verbatim}
plspm_edges_values(plsdata, digits = 3, unite = TRUE, type = "df")
\end{verbatim}
\end{Usage}
%
\begin{Arguments}
\begin{ldescription}
\item[\code{plsdata}] plsdata

\item[\code{digits}] round 3

\item[\code{unite}] data combine unite TRUE -> estimat + sig

\item[\code{type}] vec, df
\end{ldescription}
\end{Arguments}
%
\begin{Value}
estimates vectors
\end{Value}
%
\begin{Examples}
\begin{ExampleCode}

## Not run: 

data(satisfaction)
#model 1
model1 = "
EXPE ~ IMAG
QUAL ~ EXPE
VAL ~  QUAL+ EXPE
SAT ~ IMAG + EXPE + QUAL + VAL
LOY ~ SAT + IMAG
"
# # Generate the matrix
sat_path = plspm_lav2path(model1,
                          fct_order=c("IMAG", "EXPE","QUAL", "VAL", "SAT", "LOY" ))
# sat_path

# Method 2: paths
sat_path = plspm_paths(
  row_names = c("IMAG","EXPE","QUAL","VAL","SAT","LOY"),
  relationship = list(
    path(from="IMAG", to=c("EXPE","SAT","LOY")),
    path("EXPE", c("QUAL","VAL","SAT")),
    path("QUAL", c("VAL","SAT")),
    path("VAL",c("SAT")),
    path("SAT","LOY")
  )
)
# blokcs
sat_blocks1 <- plspm_blocks(
  IMAG = item(paste0("imag",1:5)),
  EXPE = item(paste0("expe", 1:5)),
  QUAL = item( paste0("qual", 1:5)),
  VAL = item(paste0("val", 1:4)),
  SAT = item(paste0("sat", 1:4)),
  LOY = item(paste0("loy", 1:4))
)
sat_blocks1

# vector of modes (reflective indicators):auto
sat_mod = rep("A", 6)
#############

# apply plspm
satpls = plspm_sem(satisfaction,
                   path_matrix = sat_path,
                   blocks = sat_blocks1,
                   scaled = FALSE)

satpls_boot = plspm_sem(satisfaction, sat_path, sat_blocks1, scaled = FALSE,
                        boot.val =TRUE, br=100)

plspm_edges_values(satpls)
plspm_edges_values(satpls, unite=F) #no star
plspm_edges_values(satpls$inner_model)
plspm_edges_values(satpls$inner_model, unite=FALSE) #No star
#' #'

lay_p = matrix(c('IMA', NA, NA, NA, NA, NA, NA, NA,
                 NA, NA, NA, NA, NA, NA, NA, NA, NA, NA, NA, 'EXP', NA,
                 NA, NA, NA, NA, NA, NA, NA, NA, NA, NA, NA, NA, NA, 'QUA',
                 NA, NA, NA, NA, NA, NA, NA, NA, NA, NA, NA, NA, NA, NA, NA,
                 NA, NA, NA, NA, NA, NA, NA, NA, NA, 'VAL', NA, NA, NA, NA,
                 NA, NA, NA, NA, NA, NA, NA, NA, NA, NA, 'SAT', NA, NA, NA,
                 NA, NA, NA, NA, NA, NA, NA, NA, NA, NA, NA, NA, NA, 'LOY',
                 NA, NA, NA, NA, NA, NA, NA, NA),
               nrow = 10, ncol = 10, byrow = FALSE)
lay_p
#  Specify a new value for the edgeLabels parameter when calling qgraph
qgraph(qgraph(t(satpls$path_coefs)),
       layout = lay_p,
       posCol = "black",
       fade = FALSE,
       edge.labels = plspm_edges_values(satpls),
       edge.label.cex = 1,
       edge.label.position = 0.4,
       groups = NULL,
       palette = 'pastel'
)
# Specify a new value for the edgeLabels parameter when calling qgraph
qgraph(qgraph(t(satpls$path_coefs)),
       layout = "spring",
       posCol = "black",
       fade = FALSE,
       edge.labels = plspm_edges_values(satpls),
       edge.label.cex = 1,
       edge.label.position = 0.4,
       groups = NULL,
       palette = 'pastel'
)

 Specify a new value for the edgeLabels parameter when calling qgraph
qgraph(qgraph(t(satpls$path_coefs)),
       layout = lay_p,
       posCol = "black",
       fade = FALSE,
       edge.labels = matrix(sprintf("%.4f", t(satpls$path_coefs)),
                            nrow=nrow(t(satpls$path_coefs)),
                            ncol=ncol(t(satpls$path_coefs))),
       edge.label.cex = 1,
       edge.label.position=0.4,
       groups = NULL, palette='pastel' )

 plspm_edges_values(satpls)
satpls %>% plspm_path_coefs_plot()
satpls %>% plspm_path_coefs_plot(node =TRUE)



## End(Not run)
\end{ExampleCode}
\end{Examples}
\inputencoding{utf8}
\HeaderA{plspm\_edge\_position}{plspm\_edge\_position}{plspm.Rul.edge.Rul.position}
%
\begin{Description}
plspm\_edge\_position
\end{Description}
%
\begin{Usage}
\begin{verbatim}
plspm_edge_position(semPaths_plot, ..., plot = TRUE)
\end{verbatim}
\end{Usage}
%
\begin{Arguments}
\begin{ldescription}
\item[\code{semPaths\_plot}] plot data generated from a PLS-PM model.

\item[\code{...}] triples of arguments specifying `from` node, `to` node, and `new\_position`.

\item[\code{plot}] logical, if `TRUE`, the plot will be displayed with updated edge positions. Default is `TRUE`.
\end{ldescription}
\end{Arguments}
%
\begin{Value}
Modified plot object with updated edge positions.
\end{Value}
%
\begin{Examples}
\begin{ExampleCode}

## Not run: 
# Example 1: Adjusting edge positions
pathmodel <- "
자기효능감 =~ C_S1 +  C_S2 + C_S3 + C_S4 + C_S5
진로태도 =~ B02 + B04 + B10 + B14
진로동기 =~ A_M1 + A_M2 + A_M3
진로준비 =~ D_P1 + D_P2 + D_P3

진로동기 ~ H1*자기효능감
진로태도 ~  H2*자기효능감
진로준비 ~ H5 *진로태도
진로준비 ~  H4 * 진로동기
진로준비 ~  H3*자기효능감
"

semPaths_plot <- plspm_diagram_model(pathmodel, layout="tree2",
                    rotation = 2, residuals = TRUE, residScale = 6,
                    edge.label.position = 0.6)

semPaths_plot %>%
  plspm_edge_position("진로태도", "진로준비", 0.4,
                      "자기효능감", "진로동기", 0.3)

# Example 2: Adjusting multiple edge positions
semPaths_plot %>%
  plspm_edge_position("진로동기", "진로태도", 0.5,
                      "진로태도", "진로준비", 0.7,
                      "자기효능감", "진로준비", 0.6)

## End(Not run)

\end{ExampleCode}
\end{Examples}
\inputencoding{utf8}
\HeaderA{plspm\_edge\_values}{Functions to include significance data for structural models in PLSPM}{plspm.Rul.edge.Rul.values}
%
\begin{Description}
Functions to include significance data for structural models in PLSPM
\end{Description}
%
\begin{Usage}
\begin{verbatim}
plspm_edge_values(plsdata, type = "vec", digits = 3, unite = TRUE)
\end{verbatim}
\end{Usage}
%
\begin{Arguments}
\begin{ldescription}
\item[\code{plsdata}] plsdata

\item[\code{type}] vec, df

\item[\code{digits}] round 3

\item[\code{unite}] data combine unite TRUE -> estimat + sig
\end{ldescription}
\end{Arguments}
%
\begin{Value}
estimates vectors
\end{Value}
\inputencoding{utf8}
\HeaderA{plspm\_effect}{plspm\_effect}{plspm.Rul.effect}
%
\begin{Description}
plspm\_effect
\end{Description}
%
\begin{Usage}
\begin{verbatim}
plspm_effect(pls_boot, digits = 3)
\end{verbatim}
\end{Usage}
%
\begin{Arguments}
\begin{ldescription}
\item[\code{pls\_boot}] pls boot

\item[\code{digits}] round 3
\end{ldescription}
\end{Arguments}
%
\begin{Value}
data
\end{Value}
\inputencoding{utf8}
\HeaderA{plspm\_effectbar}{plspm\_effectbar: Bar graph with direct, indirect, and total effects}{plspm.Rul.effectbar}
%
\begin{Description}
This function creates a bar graph visualizing direct, indirect, and total effects
using plspm result data. Users can input plspm results directly or provide custom data.
\end{Description}
%
\begin{Usage}
\begin{verbatim}
plspm_effectbar(
  x_pls,
  axis_x = 1.3,
  axis_y = 1.3,
  angle = 90,
  hjust = 1,
  axis_x_pos = 1,
  size_text = 4,
  legend.position = "top",
  input = FALSE,
  user_data = NULL,
  ylim = c(0, 1),
  col = c("black", "gray")
)
\end{verbatim}
\end{Usage}
%
\begin{Arguments}
\begin{ldescription}
\item[\code{x\_pls}] A plspm result object or a custom data frame with columns: "relationships", "direct", "indirect".

\item[\code{axis\_x}] Size of x-axis text. Default is 1.3.

\item[\code{axis\_y}] Size of y-axis text. Default is 1.3.

\item[\code{angle}] Angle of x-axis text. Default is 90 degrees.

\item[\code{hjust}] Horizontal alignment of x-axis text. Default is 1.

\item[\code{axis\_x\_pos}] Adjust x-axis position when axis labels are too large. Default is 1.

\item[\code{size\_text}] Size of text displayed on the bars. Default is 4.

\item[\code{legend.position}] Position of the legend. Default is "top". Options include "top", "bottom", "left", "right", or "none".

\item[\code{input}] Boolean indicating whether the user provides a custom dataset. Default is FALSE.

\item[\code{user\_data}] A data frame containing relationships, direct, and indirect columns if `input = TRUE`.

\item[\code{ylim}] A numeric vector of length 2 specifying the y-axis limits. Default is c(0, 1).

\item[\code{col}] A character vector specifying bar colors for "direct" and "indirect" effects. Default is c("black", "gray").
\end{ldescription}
\end{Arguments}
%
\begin{Value}
Returns the processed data frame with rounded values. It also prints a ggplot2 object.
\end{Value}
%
\begin{Examples}
\begin{ExampleCode}
# Basic usage with plspm results
edt_pls %>% plspm_effectbar()

# Directly specifying data with variable renaming
edt_pls$effects %>% select(1:3) %>%
    replace_df_rep(var_named) %>%
    plspm_effectbar(axis_x = 1.6, legend.position = "bottom")
\end{ExampleCode}
\end{Examples}
\inputencoding{utf8}
\HeaderA{plspm\_effect\_bar}{plspm\_compare\_bar}{plspm.Rul.effect.Rul.bar}
%
\begin{Description}
plspm\_compare\_bar
\end{Description}
%
\begin{Usage}
\begin{verbatim}
plspm_effect_bar(
  data,
  x = NULL,
  y = NULL,
  fill = NULL,
  size = 3,
  size_x = 12,
  angle = 90,
  position = "dodge",
  xlab = "",
  ylab = "Values",
  text = TRUE,
  digits = 3
)
\end{verbatim}
\end{Usage}
%
\begin{Arguments}
\begin{ldescription}
\item[\code{data}] arragne data

\item[\code{x}] paths

\item[\code{y}] values

\item[\code{fill}] group

\item[\code{size}] bargrahp text size

\item[\code{size\_x}] textsize

\item[\code{angle}] xangle

\item[\code{position}] dodge, stack

\item[\code{text}] hide bartext
\end{ldescription}
\end{Arguments}
%
\begin{Value}
graph
\end{Value}
%
\begin{Examples}
\begin{ExampleCode}

## Not run: 
#'
# 예시 데이터
REs <- tibble::tribble(
  ~paths, ~grp, ~values,
  "자기효능감->진로동기", "group.남자", 0.902,
  "자기효능감->진로동기", "group.여자", 0.866,
  "자기효능감->진로태도", "group.남자", -0.639,
  "자기효능감->진로태도", "group.여자", -0.672,
  "자기효능감->진로준비", "group.남자", 0.149,
  "자기효능감->진로준비", "group.여자", 0.487,
  "진로동기->진로준비", "group.남자", 0.63,
  "진로동기->진로준비", "group.여자", 0.432,
  "진로태도->진로준비", "group.남자", -0.153,
  "진로태도->진로준비", "group.여자", 0.059
)

# # 함수 호출 예시
REs %>% plspm_effect_bar()
REs %>% plspm_effect_bar(text=FALSE)

## End(Not run)

\end{ExampleCode}
\end{Examples}
\inputencoding{utf8}
\HeaderA{plspm\_extract\_blocks}{olsom boot result model}{plspm.Rul.extract.Rul.blocks}
%
\begin{Description}
olsom boot result model
\end{Description}
%
\begin{Usage}
\begin{verbatim}
plspm_extract_blocks(plsres)
\end{verbatim}
\end{Usage}
%
\begin{Arguments}
\begin{ldescription}
\item[\code{plsres}] result model
\end{ldescription}
\end{Arguments}
%
\begin{Value}
lsit
\end{Value}
%
\begin{Examples}
\begin{ExampleCode}

## Not run: 
# using htmt

## End(Not run)


\end{ExampleCode}
\end{Examples}
\inputencoding{utf8}
\HeaderA{plspm\_f2}{Effect Size Black}{plspm.Rul.f2}
%
\begin{Description}
Effect Size Black
\end{Description}
%
\begin{Usage}
\begin{verbatim}
plspm_f2(plsres, type = "res")
\end{verbatim}
\end{Usage}
%
\begin{Arguments}
\begin{ldescription}
\item[\code{plsres}] pls result

\item[\code{type}] res, f2, R2, ceofs, inner
\end{ldescription}
\end{Arguments}
%
\begin{Value}
data
\end{Value}
%
\begin{Examples}
\begin{ExampleCode}

## Not run: 
#'
#' #'
sat_path = plspm_paths(
  row_names = c("IMAG","EXPE","QUAL","VAL","SAT","LOY"),
  relationship = list(
    path(from="IMAG", to=c("EXPE","SAT","LOY")),
    path("EXPE", c("QUAL","VAL","SAT")),
    path("QUAL", c("VAL","SAT")),
    path("VAL",c("SAT")),
    path("SAT","LOY"))
)
sat_path

# blokcs
sat_blocks1 <- plspm_blocks(
  IMAG = item(paste0("imag",1:5)),
  EXPE = item(paste0("expe", 1:5)),
  QUAL = item( paste0("qual", 1:5)),
  VAL = item(paste0("val", 1:4)),
  SAT = item(paste0("sat", 1:4)),
  LOY = item(paste0("loy", 1:4))
)

sat_blocks1

satpls = plspm(satisfaction,
               path_matrix = sat_path,
               blocks = sat_blocks1,
               scaled = FALSE)

satpls_boot = plspm_sem(satisfaction, sat_path, sat_blocks1, scaled = FALSE,
                        boot.val =TRUE, br=100)
# etc
satpls_boot%>%plspm_f2()
satpls_boot%>%plspm_f2("f2")
satpls_boot%>%plspm_f2("coefs")
satpls_boot%>%plspm_f2("r2")
satpls_boot%>%plspm_f2("inner")

## End(Not run)



\end{ExampleCode}
\end{Examples}
\inputencoding{utf8}
\HeaderA{plspm\_f2\_plot}{Function to plot f2 values}{plspm.Rul.f2.Rul.plot}
%
\begin{Description}
Function to plot f2 values
\end{Description}
%
\begin{Usage}
\begin{verbatim}
plspm_f2_plot(plsres)
\end{verbatim}
\end{Usage}
%
\begin{Arguments}
\begin{ldescription}
\item[\code{plsres}] pls result
\end{ldescription}
\end{Arguments}
%
\begin{Value}
plot
\end{Value}
%
\begin{Examples}
\begin{ExampleCode}
## Not run: 

#'
# Example data for plsres
plsres <- list(
  path_coefs = data.frame(
    자기효능감 = c(0, 0.884, 0.681, 0.373),
    진로동기 = c(0, 0, 0, 0.353),
    진로태도 = c(0, 0, 0, 0.188),
    진로준비 = c(0, 0, 0, 0),
    row.names = c("자기효능감", "진로동기", "진로태도", "진로준비")
  ),
  inner_summary = data.frame(
    R2 = c(0.0000000, 0.7817016, 0.4631848, 0.7195506),
    row.names = c("자기효능감", "진로동기", "진로태도", "진로준비")
  )
)

# Run the function
plspm_f2(plsres, type="res")



plspm_f2_plot(plsres)

## End(Not run)


\end{ExampleCode}
\end{Examples}
\inputencoding{utf8}
\HeaderA{plspm\_factor\_layout}{Extracting the layout matrix of a plspm model}{plspm.Rul.factor.Rul.layout}
%
\begin{Description}
Extracting the layout matrix of a plspm model
\end{Description}
%
\begin{Usage}
\begin{verbatim}
plspm_factor_layout(plspm_model)
\end{verbatim}
\end{Usage}
%
\begin{Arguments}
\begin{ldescription}
\item[\code{plspm\_model}] sempaths
\end{ldescription}
\end{Arguments}
%
\begin{Value}
mat
\end{Value}
%
\begin{Examples}
\begin{ExampleCode}

## Not run: 
#'
data(offense)
offense %>%str()

nfl_path = matrix(
  c(0, 0, 0, 0, 0,
    0, 0, 0, 0, 0,
    0, 0, 0, 0, 0,
    0, 1, 1, 0, 0,
    1, 0, 0, 1, 0),
  nrow = 5, ncol = 5, byrow = TRUE )

rownames(nfl_path) <- c("Speical", "Rushing", "Passing", "Offense", "Scoring")
colnames(nfl_path) <- c("Speical", "Rushing", "Passing", "Offense", "Scoring")
nfl_path

nfl_blocks = list(7:8, 1:3, 4:6,1:6, 9:11)

plspm_blocks_match(data=offense, blocks=nfl_blocks,
                   paths=nfl_path)

nfl_mode =c("B","A","A","A","A")

nfl_pls = plspm(Data=offense, nfl_path, nfl_blocks, modes = nfl_mode)
nfl_pls%>%summary()

#########
nfl_pls_boot = plspm_sem(Data=offense, nfl_path, nfl_blocks, modes = nfl_mode)

#Layouts extracted from mockups
plspm_semPaths2(nfl_pls_boot,sizeLat = 6)%>%
  plspm_factor_layout()

#Make it input material
plspm_semPaths2(nfl_pls_boot,sizeLat =6)%>%
  plspm_factor_layout() %>%layout_mat(text=T)

#Organizing materials for input
plspm_semPaths2(nfl_pls_boot,sizeLat =6)%>%
  plspm_factor_layout() %>%layout_mat()%>%
  mat_text_arrange()

## End(Not run)



\end{ExampleCode}
\end{Examples}
\inputencoding{utf8}
\HeaderA{plspm\_fl}{Fornell and Locker 1981 method, discrimninat validation}{plspm.Rul.fl}
%
\begin{Description}
Fornell and Locker 1981 method, discrimninat validation
\end{Description}
%
\begin{Usage}
\begin{verbatim}
plspm_fl(plsdata)
\end{verbatim}
\end{Usage}
%
\begin{Arguments}
\begin{ldescription}
\item[\code{plsdata}] plsdata or pls summary
\end{ldescription}
\end{Arguments}
%
\begin{Value}
corand AVE
\end{Value}
%
\begin{Examples}
\begin{ExampleCode}

## Not run: 

data(satisfaction)
#model 1
model1 = "
EXPE ~ IMAG
QUAL ~ EXPE
VAL ~  QUAL+ EXPE
SAT ~ IMAG + EXPE + QUAL + VAL
LOY ~ SAT + IMAG
"
# # Generate the matrix
sat_path = plspm_lav2path(model1,
                          fct_order=c("IMAG", "EXPE","QUAL", "VAL", "SAT", "LOY" ))
# sat_path

# Method 2: paths
sat_path = plspm_paths(
  row_names = c("IMAG","EXPE","QUAL","VAL","SAT","LOY"),
  relationship = list(
    path(from="IMAG", to=c("EXPE","SAT","LOY")),
    path("EXPE", c("QUAL","VAL","SAT")),
    path("QUAL", c("VAL","SAT")),
    path("VAL",c("SAT")),
    path("SAT","LOY")
  )
)
# blokcs
sat_blocks1 <- plspm_blocks(
  IMAG = item(paste0("imag",1:5)),
  EXPE = item(paste0("expe", 1:5)),
  QUAL = item( paste0("qual", 1:5)),
  VAL = item(paste0("val", 1:4)),
  SAT = item(paste0("sat", 1:4)),
  LOY = item(paste0("loy", 1:4))
)
sat_blocks1

# vector of modes (reflective indicators):auto
sat_mod = rep("A", 6)
#############

# apply plspm
satpls = plspm_sem(satisfaction,
                   path_matrix = sat_path,
                   blocks = sat_blocks1,
                   scaled = FALSE)

satpls_boot = plspm_sem(satisfaction, sat_path, sat_blocks1, scaled = FALSE,
                        boot.val =TRUE, br=100)

satpls_boot %>% plspm_fl()
satpls_boot_summary %>% plspm_fl()


## End(Not run)
\end{ExampleCode}
\end{Examples}
\inputencoding{utf8}
\HeaderA{plspm\_grp\_comp}{plspm\_grp\_comp:Creating multigroup list data into a data frame}{plspm.Rul.grp.Rul.comp}
%
\begin{Description}
plspm\_grp\_comp:Creating multigroup list data into a data frame
\end{Description}
%
\begin{Usage}
\begin{verbatim}
plspm_grp_comp(lst_data)
\end{verbatim}
\end{Usage}
%
\begin{Arguments}
\begin{ldescription}
\item[\code{lst\_data}] grp data
\end{ldescription}
\end{Arguments}
%
\begin{Value}
data
\end{Value}
%
\begin{Examples}
\begin{ExampleCode}

## Not run: 
#'
plspm_grp_comp(teacher_compare) %>%
  nice_table(exclude=c("df", "sig"))%>%
  replace_df_rep(var_named) %>%Round(2)%>%pall()

## End(Not run)

\end{ExampleCode}
\end{Examples}
\inputencoding{utf8}
\HeaderA{plspm\_grp\_compare\_bar}{plspm\_grp\_compare\_bar}{plspm.Rul.grp.Rul.compare.Rul.bar}
%
\begin{Description}
plspm\_grp\_compare\_bar
\end{Description}
%
\begin{Usage}
\begin{verbatim}
plspm_grp_compare_bar(
  data,
  x = NULL,
  y = NULL,
  fill = NULL,
  size = 3,
  size_x = 12,
  angle = 90,
  position = "dodge",
  xlab = "",
  ylab = "Values",
  text = TRUE
)
\end{verbatim}
\end{Usage}
%
\begin{Arguments}
\begin{ldescription}
\item[\code{data}] arragne data

\item[\code{x}] paths

\item[\code{y}] values

\item[\code{fill}] group

\item[\code{size}] bargrahp text size

\item[\code{size\_x}] textsize

\item[\code{angle}] xangle

\item[\code{position}] dodge, stack

\item[\code{xlab}] default ""

\item[\code{ylab}] values

\item[\code{text}] hide bartext
\end{ldescription}
\end{Arguments}
%
\begin{Value}
graph
\end{Value}
%
\begin{Examples}
\begin{ExampleCode}

## Not run: 
#'
# 예시 데이터
REs <- tibble::tribble(
  ~paths, ~grp, ~values,
  "자기효능감->진로동기", "group.남자", 0.902,
  "자기효능감->진로동기", "group.여자", 0.866,
  "자기효능감->진로태도", "group.남자", -0.639,
  "자기효능감->진로태도", "group.여자", -0.672,
  "자기효능감->진로준비", "group.남자", 0.149,
  "자기효능감->진로준비", "group.여자", 0.487,
  "진로동기->진로준비", "group.남자", 0.63,
  "진로동기->진로준비", "group.여자", 0.432,
  "진로태도->진로준비", "group.남자", -0.153,
  "진로태도->진로준비", "group.여자", 0.059
)

# # 함수 호출 예시
REs %>% compare_bar()
REs %>% compare_bar(text=FALSE)

## End(Not run)
\end{ExampleCode}
\end{Examples}
\inputencoding{utf8}
\HeaderA{plspm\_grp\_comp\_apa}{plspm\_grp\_comp\_apa: Generate APA-style Group Comparison Descriptions}{plspm.Rul.grp.Rul.comp.Rul.apa}
%
\begin{Description}
This function generates APA-style text descriptions for PLS-PM group comparisons.
\end{Description}
%
\begin{Usage}
\begin{verbatim}
plspm_grp_comp_apa(df, edteck = TRUE)
\end{verbatim}
\end{Usage}
%
\begin{Arguments}
\begin{ldescription}
\item[\code{df}] A data frame containing the group comparison results. It should have the columns:
- `comp`: Group comparison identifiers (e.g., "1\_vs\_2").
- `Paths`: The path being analyzed.
- `t`: The t-value for the comparison.
- `p\_value`: The p-value for the comparison.
- `df`: Degrees of freedom for the t-test.
- `sig`: Significance indicator ("yes" if significant, otherwise "no").

\item[\code{edteck}] Logical. If `TRUE`, predefined group names are used for educational technology profiles.
If `FALSE`, group names are auto-generated as "그룹1", "그룹2", etc.
\end{ldescription}
\end{Arguments}
%
\begin{Value}
None. The function prints APA-style descriptions for significant group comparisons.
\end{Value}
%
\begin{Examples}
\begin{ExampleCode}
## Not run: 
df <- data.frame(
  comp = c("1_vs_2", "2_vs_3"),
  Paths = c("Engagement", "Performance"),
  t = c(3.21, 2.45),
  p_value = c(0.0023, 0.0184),
  df = c(45, 50),
  sig = c("yes", "no")
)
plspm_grp_comp_apa(df, edteck = TRUE)

## End(Not run)
\end{ExampleCode}
\end{Examples}
\inputencoding{utf8}
\HeaderA{plspm\_grp\_summary}{plspm\_grp\_summary}{plspm.Rul.grp.Rul.summary}
%
\begin{Description}
plspm\_grp\_summary
\end{Description}
%
\begin{Usage}
\begin{verbatim}
plspm_grp_summary(plsres, digits = 3)
\end{verbatim}
\end{Usage}
%
\begin{Arguments}
\begin{ldescription}
\item[\code{plsres}] plspm boot group

\item[\code{digits}] round 3
\end{ldescription}
\end{Arguments}
%
\begin{Value}
table
\end{Value}
%
\begin{Examples}
\begin{ExampleCode}

## Not run: 
#'
jutpath1 = plspm_paths(
  row_names = c("자기효능감","진로동기","진로태도","진로준비"),
  relationship = list(
    path(from="자기효능감", to=c("진로동기","진로태도","진로준비")),
    path("진로동기", c("진로준비")),
    path("진로태도", c("진로준비"))
  )  )
jutpath1

plspm_paths2lav(jutpath1)%>%cat()
#측정모형
jut_blocks <- plspm_blocks(
  자기효능감 = item("C_S1","C_S2","C_S3","C_S4","C_S5"),
  진로동기 = item("A_M1","A_M2","A_M3"),
  진로태도 = item(  "B02", "B04", "B10","B14" ),#결정성
  진로준비 = item("D_P1","D_P2","D_P3")
)

jutpls_boot = plspm_sem(Data= jut6, path_matrix = jutpath1,
                        blocks = jut_blocks,
                        br = 500, summary= TRUE)

#'
jut_gender = plspm.groups(jutpls_boot,
                          group= factor(jut6$성별),
                          method="bootstrap")
jut_gender_p = plspm.groups(jutpls_boot,
                            group= factor(jut6$성별),
                            method="permutation")

jut_gender%>%plspm_grp_summary()
jut_gender_p%>%plspm_grp_summary()

## End(Not run)

\end{ExampleCode}
\end{Examples}
\inputencoding{utf8}
\HeaderA{plspm\_htmt}{Define a function to calculate HTMT}{plspm.Rul.htmt}
%
\begin{Description}
Define a function to calculate HTMT
\end{Description}
%
\begin{Usage}
\begin{verbatim}
plspm_htmt(
  plsres = NULL,
  data = NULL,
  blocks = NULL,
  htmt2 = TRUE,
  sig = TRUE,
  cut = 0.9,
  digits = 3,
  imp = ""
)
\end{verbatim}
\end{Usage}
%
\begin{Arguments}
\begin{ldescription}
\item[\code{data}] data

\item[\code{blocks}] latent =\textasciitilde{} item, plspl blocks

\item[\code{htmt2}] TRUE Apply geometric mean

\item[\code{sig}] sig= TRUE- > show star

\item[\code{cut}] = 0.9, 0.85. 1

\item[\code{digits}] = 3

\item[\code{imp}] ""
\end{ldescription}
\end{Arguments}
%
\begin{Value}
result
\end{Value}
%
\begin{Examples}
\begin{ExampleCode}
## Not run: 
#'
sat_path = plspm_paths(
row_names = c("IMAG","EXPE","QUAL","VAL","SAT","LOY"),
relationship = list(
  path(from="IMAG", to=c("EXPE","SAT","LOY")),
  path("EXPE", c("QUAL","VAL","SAT")),
  path("QUAL", c("VAL","SAT")),
  path("VAL",c("SAT")),
  path("SAT","LOY"))
)

sat_blocks1 <- plspm_blocks(
  IMAG = item(paste0("imag",1:5)),
  EXPE = item(paste0("expe", 1:5)),
  QUAL = item( paste0("qual", 1:5)),
  VAL = item(paste0("val", 1:4)),
  SAT = item(paste0("sat", 1:4)),
  LOY = item(paste0("loy", 1:4))
)

sat_blocks1
# plspm_blocks2lav(sat_blocks1) %>% cat("\n")
data(satisfaction)
plspm_htmt(satisfaction, sat_blocks1)

#'
#'
plspm_htmt(satisfaction, sat_blocks1)
plspm_htmt(satisfaction, sat_blocks1, sig=F)

plspm_htmt(satisfaction, sat_blocks1,sig=F, htmt2=FALSE)
plspm_htmt(satisfaction, sat_blocks1,sig=T, cut=0.9, htmt2=FALSE)

## Reference
semTools::htmt(plspm_blocks2lav(sat_blocks1),
               satisfaction,
               htmt2 = T)

 satpls_boot = plspm_sem(satisfaction, sat_path, sat_blocks1,
  scaled = FALSE, boot.val =TRUE, br=100)
 #If you've bootstrapped
 plspm_htmt(satpls_boot)
#'


## End(Not run)



\end{ExampleCode}
\end{Examples}
\inputencoding{utf8}
\HeaderA{plspm\_htmt2}{Define a function to calculate HTMT}{plspm.Rul.htmt2}
%
\begin{Description}
Define a function to calculate HTMT
\end{Description}
%
\begin{Usage}
\begin{verbatim}
plspm_htmt2(
  data,
  blocks,
  htmt2 = TRUE,
  sig = TRUE,
  cut = 0.9,
  digits = 3,
  imp = ""
)
\end{verbatim}
\end{Usage}
%
\begin{Arguments}
\begin{ldescription}
\item[\code{data}] data

\item[\code{blocks}] latent =\textasciitilde{} item, plspl blocks

\item[\code{htmt2}] TRUE Apply geometric mean

\item[\code{sig}] sig= TRUE- > show star

\item[\code{cut}] = 0.9, 0.85. 1

\item[\code{digits}] = 3

\item[\code{imp}] imputation
\end{ldescription}
\end{Arguments}
%
\begin{Value}
data
\end{Value}
%
\begin{Examples}
\begin{ExampleCode}

## Not run: 
#'
sat_blocks1 <- plspm_blocks(
  IMAG = item(paste0("imag",1:5)),
  EXPE = item(paste0("expe", 1:5)),
  QUAL = item( paste0("qual", 1:5)),
  VAL = item(paste0("val", 1:4)),
  SAT = item(paste0("sat", 1:4)),
  LOY = item(paste0("loy", 1:4))
)

sat_blocks1
# plspm_blocks2lav(sat_blocks1) %>% cat("\n")
data(satisfaction)
plspm_htmt2(satisfaction, sat_blocks1)

#'
#'
plspm_htmt2(satisfaction, sat_blocks1)
plspm_htmt2(satisfaction, sat_blocks1, sig=F)

plspm_htmt2(satisfaction, sat_blocks1,sig=F, htmt2=FALSE)
plspm_htmt2(satisfaction, sat_blocks1,sig=T, cut=0.9, htmt2=FALSE)

## Reference
semTools::htmt(plspm_blocks2lav(sat_blocks1),
               satisfaction,
               htmt2 = T)
#'


## End(Not run)

\end{ExampleCode}
\end{Examples}
\inputencoding{utf8}
\HeaderA{plspm\_indicator\_factor}{plspm\_indicator\_factor}{plspm.Rul.indicator.Rul.factor}
%
\begin{Description}
plspm\_indicator\_factor
\end{Description}
%
\begin{Usage}
\begin{verbatim}
plspm_indicator_factor(model_syntax)
\end{verbatim}
\end{Usage}
%
\begin{Arguments}
\begin{ldescription}
\item[\code{model\_syntax}] lavaan
\end{ldescription}
\end{Arguments}
%
\begin{Value}
text
\end{Value}
%
\begin{Examples}
\begin{ExampleCode}

## Not run: 
#'
mod <-
  'f1 =~ x01 + x02 + x03 + x06
    f3 =~ x08 + x09 + x10 + x03
    f2 =~ x04 + x05 + x06 + x07
    f4 =~ x11 + x12 + x13 + x14
    f3 ~ f1 + f2
    f4 ~ f3
   '
plspm_indicator_order(mod)

plspm_indicator_factor(mod)

## End(Not run)

\end{ExampleCode}
\end{Examples}
\inputencoding{utf8}
\HeaderA{plspm\_indicator\_order}{plspm\_indicator\_order}{plspm.Rul.indicator.Rul.order}
%
\begin{Description}
plspm\_indicator\_order
\end{Description}
%
\begin{Usage}
\begin{verbatim}
plspm_indicator_order(model_syntax)
\end{verbatim}
\end{Usage}
%
\begin{Arguments}
\begin{ldescription}
\item[\code{model\_syntax}] lavaan
\end{ldescription}
\end{Arguments}
%
\begin{Value}
text
\end{Value}
%
\begin{Examples}
\begin{ExampleCode}
## Not run: 
#'
mod <-
  'f1 =~ x01 + x02 + x03 + x06
    f3 =~ x08 + x09 + x10 + x03
    f2 =~ x04 + x05 + x06 + x07
    f4 =~ x11 + x12 + x13 + x14
    f3 ~ f1 + f2
    f4 ~ f3
   '
plspm_indicator_order(mod)

plspm_indicator_factor(mod)

## End(Not run)

\end{ExampleCode}
\end{Examples}
\inputencoding{utf8}
\HeaderA{plspm\_indicator\_test}{지표의 균일성 평가를 위한 Permutation Test 함수 (병렬 처리 적용)}{plspm.Rul.indicator.Rul.test}
%
\begin{Description}
지표의 균일성 평가를 위한 Permutation Test 함수 (병렬 처리 적용)
\end{Description}
%
\begin{Usage}
\begin{verbatim}
plspm_indicator_test(
  Data,
  path_matrix,
  blocks,
  grp,
  n_perm = 100,
  n_boot = 100,
  n_cores = 12
)
\end{verbatim}
\end{Usage}
%
\begin{Arguments}
\begin{ldescription}
\item[\code{Data}] data

\item[\code{path\_matrix}] path

\item[\code{blocks}] measurement

\item[\code{grp}] grp

\item[\code{n\_perm}] permutation number

\item[\code{n\_boot}] boot n

\item[\code{n\_cores}] cores=8 (recently 12\textasciitilde{}15), exceapt one !!!
\end{ldescription}
\end{Arguments}
%
\begin{Value}
data
\end{Value}
%
\begin{Examples}
\begin{ExampleCode}

## Not run: 
#'
# 예시: 지표의 균일성 평가 실행 (병렬 처리 적용)
result_df2 <- plspm_indicator_test(Data = jut7c, path_matrix = jutpath1, blocks = jut_blocks, grp = "P01")
print(result_df2)

## End(Not run)

\end{ExampleCode}
\end{Examples}
\inputencoding{utf8}
\HeaderA{plspm\_indicator\_test\_multi}{Perform Bootstrap and Permutation Tests for PLS-PM Indicators Across Groups}{plspm.Rul.indicator.Rul.test.Rul.multi}
%
\begin{Description}
This function performs bootstrap and permutation tests for PLS-PM indicators across
selected group pairs or all possible group pairs in the dataset. The results include
group-specific estimates, standard errors, p-values, and invariance checks.
\end{Description}
%
\begin{Usage}
\begin{verbatim}
plspm_indicator_test_multi(
  Data,
  path_matrix,
  blocks,
  grp,
  n_perm = 50,
  n_boot = 50,
  n_cores = 12,
  grp_select = NULL
)
\end{verbatim}
\end{Usage}
%
\begin{Arguments}
\begin{ldescription}
\item[\code{Data}] A data frame containing the dataset to be analyzed.

\item[\code{path\_matrix}] A square matrix defining the path relationships among constructs in the PLS-PM model.

\item[\code{blocks}] A list defining the measurement blocks (i.e., sets of observed variables) for each construct.

\item[\code{grp}] A string specifying the name of the grouping variable in `Data`.

\item[\code{n\_perm}] Number of permutation tests to perform. Default is 50.

\item[\code{n\_boot}] Number of bootstrap resamples to perform. Default is 50.

\item[\code{n\_cores}] Number of cores to use for parallel processing. Default is 12.

\item[\code{grp\_select}] A character vector of length 2 specifying the group pair to compare. Default is NULL (compares all pairs).
\end{ldescription}
\end{Arguments}
%
\begin{Value}
A tibble containing the test results for all indicators, including estimates, standard errors, p-values, and comparisons.
The returned tibble includes the following columns:
- `indicator`: The indicator being tested.
- `g1\_est`: The estimate for the first group.
- `g2\_est`: The estimate for the second group.
- `g1\_se`: The standard error for the first group.
- `g2\_se`: The standard error for the second group.
- `se\_diff`: The difference in standard errors between the groups.
- `p\_value`: The p-value for the test.
- `invariance`: Whether the indicator is invariant between the groups (`"OK"` or `"ns"`).
- `comparison`: The comparison between the two groups (e.g., "group1\_vs\_group2").
\end{Value}
%
\begin{Examples}
\begin{ExampleCode}
## Not run: 
result <- plspm_indicator_test_multi(
  Data = ed_tam2_st,
  path_matrix = edmodel,
  blocks = edblock,
  grp = "Class",
  grp_select = c("1", "3"),
  n_perm = 100,
  n_boot = 50,
  n_cores = 13
)
print(result)

## End(Not run)
\end{ExampleCode}
\end{Examples}
\inputencoding{utf8}
\HeaderA{plspm\_ind\_effect}{plspm\_ind\_effect}{plspm.Rul.ind.Rul.effect}
%
\begin{Description}
plspm\_ind\_effect
\end{Description}
%
\begin{Usage}
\begin{verbatim}
plspm_ind_effect(boot_data, from, through, to, type = "res", digits = 3)
\end{verbatim}
\end{Usage}
%
\begin{Arguments}
\begin{ldescription}
\item[\code{boot\_data}] boot\_data

\item[\code{from}] from

\item[\code{through}] through

\item[\code{to}] to

\item[\code{type}] res, all

\item[\code{digits}] 3
\end{ldescription}
\end{Arguments}
%
\begin{Value}
data
\end{Value}
%
\begin{Examples}
\begin{ExampleCode}

## Not run: 
#'
# Example usage of the function with the provided PLS-PM model
plspm_ind_effect(
  boot_data = jutpls_boot,
  from = "자기효능감",
  through = c("진로동기", "진로태도"),
  to = "진로준비"
)

plspm_ind_effect(
  boot_data = jutpls_boot,
  from = "자기효능감",
  through = c("진로동기", "진로태도"),
  to = "진로준비", type="all"
)

plspm_ind_effect(
  boot_data = jutpls_boot,
  from = "자기효능감",
  through = c("진로동기"),
  to = "진로준비", type="res"
)


## End(Not run)


\end{ExampleCode}
\end{Examples}
\inputencoding{utf8}
\HeaderA{plspm\_ind\_effect\_boot}{plspm ind efffect bootstrap}{plspm.Rul.ind.Rul.effect.Rul.boot}
%
\begin{Description}
plspm ind efffect bootstrap
\end{Description}
%
\begin{Usage}
\begin{verbatim}
plspm_ind_effect_boot(
  plsres = NULL,
  ...,
  data = NULL,
  path_matrix = NULL,
  blocks = NULL,
  modes = NULL,
  B = 100,
  n_cores = NULL
)
\end{verbatim}
\end{Usage}
%
\begin{Arguments}
\begin{ldescription}
\item[\code{plsres}] plspm result boot

\item[\code{...}] paths

\item[\code{data}] data

\item[\code{path\_matrix}] paths

\item[\code{blocks}] blocks

\item[\code{modes}] models

\item[\code{B}] b= 100

\item[\code{n\_cores}] n\_cores =4, CPU core
\end{ldescription}
\end{Arguments}
%
\begin{Value}
resutl paths
\end{Value}
%
\begin{Examples}
\begin{ExampleCode}

## Not run: 

#'
data(satisfaction)
sat_path = plspm_paths(
  row_names = c("IMAG","EXPE","QUAL","VAL","SAT","LOY"),
  relationship = list(
    path(from="IMAG", to=c("EXPE","SAT","LOY")),
    path("EXPE", c("QUAL","VAL","SAT")),
    path("QUAL", c("VAL","SAT")),
    path("VAL",c("SAT")),
    path("SAT","LOY"))
)
sat_path

# blokcs
sat_blocks1 <- plspm_blocks(
  IMAG = item(paste0("imag",1:5)),
  EXPE = item(paste0("expe", 1:5)),
  QUAL = item( paste0("qual", 1:5)),
  VAL = item(paste0("val", 1:4)),
  SAT = item(paste0("sat", 1:4)),
  LOY = item(paste0("loy", 1:4))
)

sat_blocks1

# vector of modes (reflective indicators):auto
sat_mod = rep("A", 6)

# satpls = plspm(satisfaction, sat_path, sat_blocks1, scaled = FALSE)
# satpls
# satpls %>%summary()

satpls_boot = plspm_sem(satisfaction, sat_path, sat_blocks1, scaled = FALSE,
                        boot.val =TRUE, br=100)


long_paths <- c(
  "IMAG -> EXPE -> QUAL -> VAL -> SAT -> LOY",
  "IMAG -> EXPE -> QUAL -> SAT -> LOY",
  "IMAG -> EXPE -> VAL -> SAT -> LOY",
  "IMAG -> EXPE -> SAT -> LOY",
  "IMAG -> SAT -> LOY"
)

result <- plspm_get_boot(satpls_boot, long_paths, B = 100)
result
result %>% Round(3)
result %>% unite_ci()

resultall <- plspm_get_boot(satpls_boot, find_paths(satpls_boot), B = 100)
resultall
#'
plspm_ind_effect_boot(
  plsres = satpls_boot,
  "IMAG -> SAT -> LOY",
  "IMAG -> EXPE -> SAT -> LOY",
  "IMAG -> EXPE -> VAL -> SAT -> LOY",
  "IMAG -> EXPE -> QUAL -> SAT -> LOY",
  "IMAG -> EXPE -> QUAL -> VAL -> SAT -> LOY",
  B = 1000,
  n_cores = 4
)%>%system.time()
#'
#'
plspm_ind_effect_boot(
  plsres = satpls_boot,
  "IMAG -> SAT -> LOY",
  "IMAG -> EXPE -> SAT -> LOY",
  "IMAG -> EXPE -> VAL -> SAT -> LOY",
  "IMAG -> EXPE -> QUAL -> SAT -> LOY",
  "IMAG -> EXPE -> QUAL -> VAL -> SAT -> LOY",
  B = 1000,
  n_cores = 4
)
#'



## End(Not run)




\end{ExampleCode}
\end{Examples}
\inputencoding{utf8}
\HeaderA{plspm\_ind\_effect\_sig}{plspm\_ind\_effect\_sig}{plspm.Rul.ind.Rul.effect.Rul.sig}
%
\begin{Description}
plspm\_ind\_effect\_sig
\end{Description}
%
\begin{Usage}
\begin{verbatim}
plspm_ind_effect_sig(
  boot_data,
  from,
  through,
  to,
  est = "Mean.Boot",
  digits = 3
)
\end{verbatim}
\end{Usage}
%
\begin{Arguments}
\begin{ldescription}
\item[\code{boot\_data}] bootstrap data

\item[\code{from}] from

\item[\code{through}] era

\item[\code{to}] to

\item[\code{est}] Mean.Boot, Original

\item[\code{digits}] 3
\end{ldescription}
\end{Arguments}
%
\begin{Value}
data
\end{Value}
%
\begin{Examples}
\begin{ExampleCode}

## Not run: 
#'
# 함수 호출 예시
plspm_ind_effect_sig(
  boot_data = satpls_boot,
  from = "IMAG",
  through = c("EXPE", "QUAL", "SAT"),
  to = "VAL"
)


## End(Not run)



\end{ExampleCode}
\end{Examples}
\inputencoding{utf8}
\HeaderA{plspm\_ind\_effect\_test}{Test for all indirect effects}{plspm.Rul.ind.Rul.effect.Rul.test}
%
\begin{Description}
Test for all indirect effects
\end{Description}
%
\begin{Usage}
\begin{verbatim}
plspm_ind_effect_test(
  bootdata,
  ...,
  show = TRUE,
  est = "Mean.Boot",
  sobel = "aroian"
)
\end{verbatim}
\end{Usage}
%
\begin{Arguments}
\begin{ldescription}
\item[\code{...}] paths

\item[\code{show}] show paths details

\item[\code{est}] Original or Mean.Boot

\item[\code{sobel}] aroian, sobel, goodman

\item[\code{bootdataboot}] data
\end{ldescription}
\end{Arguments}
%
\begin{Value}
data result
\end{Value}
%
\begin{Examples}
\begin{ExampleCode}

## Not run: 
#'
#'
ple usage:
  satpls_boot <- list(boot = list(paths = data.frame(
    Original = c(0.578959128, 0.200724200, 0.275149576, 0.848344408,
                 0.105477650, -0.002753995, 0.676655529, 0.122144566, 0.589330650, 0.495479322),
    Mean.Boot = c(0.58232765, 0.21209094, 0.27486564, 0.84793899,
                  0.10992051, -0.01008222, 0.67780527, 0.15132543, 0.56312717, 0.50757755),
    Std.Error = c(0.04358293, 0.06056988, 0.07729000, 0.01864573,
                  0.06968826, 0.06507903, 0.07760527, 0.08440236, 0.08723293, 0.07856694),
    perc.025 = c(0.497644719, 0.107561558, 0.140641502, 0.815794494,
                 -0.025566471, -0.147788188, 0.519959683, -0.008629516, 0.392664443, 0.351774130),
    perc.975 = c(0.6646198, 0.3275312, 0.4161131, 0.8828875,
                 0.2545785, 0.1086856, 0.8166486, 0.2991045, 0.7012494, 0.6556333),
    row.names = c("IMAG -> EXPE", "IMAG -> SAT", "IMAG -> LOY",
                  "EXPE -> QUAL", "EXPE -> VAL", "EXPE -> SAT",
                  "QUAL -> VAL", "QUAL -> SAT", "VAL -> SAT", "SAT -> LOY")
  )))
  # bootdata$boot$paths




  plspm_ind_effect_test(satpls_boot,
                      "IMAG -> EXPE -> VAL -> SAT -> LOY",
                      "IMAG -> SAT -> LOY")

  plspm_ind_effect_test(satpls_boot$boot$paths,
                        "IMAG -> EXPE -> VAL -> SAT -> LOY",
                        "IMAG -> SAT -> LOY")

  plspm_ind_effect_test(satpls_boot,
                      "IMAG -> EXPE -> VAL -> SAT -> LOY",
                      "IMAG -> SAT -> LOY", est = "Mean.Boot")

  plspm_ind_effect_test(satpls_boot$boot$paths,
                        "IMAG -> EXPE -> VAL -> SAT -> LOY",
                        "IMAG -> SAT -> LOY", est = "Mean.Boot")
  plspm_ind_effect_test(satpls_boot$boot$paths,
                        "IMAG -> EXPE -> VAL -> SAT -> LOY",
                        "IMAG -> SAT -> LOY",
                        est = "Mean.Boot",
                        sobel="sobel")

  plspm_ind_effect_test(satpls_boot$boot$paths,
                        "IMAG -> EXPE -> VAL -> SAT -> LOY",
                        "IMAG -> SAT -> LOY",
                        est = "Mean.Boot",
                        sobel="goodman")


  # For individual path strings
  plspm_ind_effect_test(satpls_boot)
  plspm_ind_effect_test(satpls_boot, show=TRUE)
  plspm_ind_effect_test(satpls_boot, show=FALSE)

  plspm_ind_effect_test(satpls_boot,
                      "IMAG -> EXPE -> VAL -> SAT -> LOY",
                      "IMAG -> SAT -> LOY")

  # For a list of paths from find_paths(satpls_boot)$paths
  plspm_ind_effect_test(satpls_boot, find_paths(satpls_boot))

  plspm_ind_effect_test(satpls_boot, "IMAG -> EXPE -> SAT -> LOY")
  plspm_ind_effect_test(satpls_boot,  "IMAG -> EXPE -> VAL -> SAT -> LOY",
   "IMAG -> SAT -> LOY")


  find_paths(satpls_boot, type= "ind")
  #default
  find_paths(satpls_boot, type= "paths")

  RRR = find_paths(bootdata, type="paths")
  RRR[1]
  RRR[2]
  RRR[3]
  RRR[4]
  RRR[5]
  RRR[6]

  plspm_ind_effect_test(satpls_boot,  RRR[1],RRR[2],RRR[3],RRR[4],RRR[5] )

  plspm_ind_effect_test(satpls_boot, find_paths(satpls_boot,"ind")%>%unlist())





## End(Not run)

\end{ExampleCode}
\end{Examples}
\inputencoding{utf8}
\HeaderA{plspm\_ind\_factor}{plspm\_indicator\_factor}{plspm.Rul.ind.Rul.factor}
%
\begin{Description}
plspm\_indicator\_factor
\end{Description}
%
\begin{Usage}
\begin{verbatim}
plspm_ind_factor(model_syntax)
\end{verbatim}
\end{Usage}
%
\begin{Arguments}
\begin{ldescription}
\item[\code{model\_syntax}] lavaan
\end{ldescription}
\end{Arguments}
%
\begin{Value}
text
\end{Value}
%
\begin{Examples}
\begin{ExampleCode}

## Not run: 
#'
mod <-
  'f1 =~ x01 + x02 + x03 + x06
    f3 =~ x08 + x09 + x10 + x03
    f2 =~ x04 + x05 + x06 + x07
    f4 =~ x11 + x12 + x13 + x14
    f3 ~ f1 + f2
    f4 ~ f3
   '
plspm_indicator_order(mod)

plspm_indicator_factor(mod)

## End(Not run)

\end{ExampleCode}
\end{Examples}
\inputencoding{utf8}
\HeaderA{plspm\_ind\_order}{plspm\_indicator\_order}{plspm.Rul.ind.Rul.order}
%
\begin{Description}
plspm\_indicator\_order
\end{Description}
%
\begin{Usage}
\begin{verbatim}
plspm_ind_order(model_syntax)
\end{verbatim}
\end{Usage}
%
\begin{Arguments}
\begin{ldescription}
\item[\code{model\_syntax}] lavaan
\end{ldescription}
\end{Arguments}
%
\begin{Value}
text
\end{Value}
%
\begin{Examples}
\begin{ExampleCode}
## Not run: 
#'
mod <-
  'f1 =~ x01 + x02 + x03 + x06
    f3 =~ x08 + x09 + x10 + x03
    f2 =~ x04 + x05 + x06 + x07
    f4 =~ x11 + x12 + x13 + x14
    f3 ~ f1 + f2
    f4 ~ f3
   '
plspm_indicator_order(mod)

plspm_indicator_factor(mod)

## End(Not run)

\end{ExampleCode}
\end{Examples}
\inputencoding{utf8}
\HeaderA{plspm\_ind\_order2}{plspm\_ind\_order2}{plspm.Rul.ind.Rul.order2}
%
\begin{Description}
plspm\_ind\_order2
\end{Description}
%
\begin{Usage}
\begin{verbatim}
plspm_ind_order2(model_syntax)
\end{verbatim}
\end{Usage}
%
\begin{Arguments}
\begin{ldescription}
\item[\code{model\_syntax}] lavaan
\end{ldescription}
\end{Arguments}
%
\begin{Value}
input text
\end{Value}
%
\begin{Examples}
\begin{ExampleCode}

## Not run: 
#'
mod <-
  'f1 =~ x01 + x02 + x03 + x06
    f3 =~ x08 + x09 + x10 + x03
    f2 =~ x04 + x05 + x06 + x07
    f4 =~ x11 + x12 + x13 + x14
    f3 ~ f1 + f2
    f4 ~ f3
   '
plspm_ind_order2(mod)


## End(Not run)

\end{ExampleCode}
\end{Examples}
\inputencoding{utf8}
\HeaderA{plspm\_ind\_push}{plspm\_ind\_push}{plspm.Rul.ind.Rul.push}
%
\begin{Description}
plspm\_ind\_push
\end{Description}
%
\begin{Usage}
\begin{verbatim}
plspm_ind_push(pls_boot, set = 2)
\end{verbatim}
\end{Usage}
%
\begin{Arguments}
\begin{ldescription}
\item[\code{pls\_boot}] pls\_boot

\item[\code{set}] set=1.5 or c(set= c(2,2,1.5,1.5))
\end{ldescription}
\end{Arguments}
%
\begin{Value}
vector
\end{Value}
%
\begin{Examples}
\begin{ExampleCode}

## Not run: 

# 예시 사용법
jutpls_boot <- list(model = list(gens = list(lvs_names = c("자기효능감", "진로동기", "진로태도", "진로준비"))))

# 기본값 사용하는 경우
result_default <- plspm_ind_push(jutpls_boot)
print(result_default)

# 이름과 값이 있는 경우
result_named <- plspm_ind_push(jutpls_boot, set = c(진로동기=2, 자기효능감=3, 진로준비=1.5, 진로태도=2.5))
print(result_named)

# 값만 있는 경우
result_unnamed <- plspm_ind_push(jutpls_boot, set = c(2, 2, 1.5, 1.5))
print(result_unnamed)

# 부분적으로 값이 있는 경우
result_partial <- plspm_ind_push(jutpls_boot, set = c(진로동기=2, 자기효능감=3))
print(result_partial)


## End(Not run)



\end{ExampleCode}
\end{Examples}
\inputencoding{utf8}
\HeaderA{plspm\_ind\_spread}{plspm\_ind\_push}{plspm.Rul.ind.Rul.spread}
%
\begin{Description}
plspm\_ind\_push
\end{Description}
%
\begin{Usage}
\begin{verbatim}
plspm_ind_spread(pls_boot, set = 2)
\end{verbatim}
\end{Usage}
%
\begin{Arguments}
\begin{ldescription}
\item[\code{pls\_boot}] pls\_boot

\item[\code{set}] set=1.5 or c(set= c(2,2,1.5,1.5))
\end{ldescription}
\end{Arguments}
%
\begin{Value}
vector
\end{Value}
%
\begin{Examples}
\begin{ExampleCode}

## Not run: 
#'
# 예시 사용법
jutpls_boot <- list(model = list(gens = list(lvs_names = c("자기효능감", "진로동기", "진로태도", "진로준비"))))

# 기본값 사용하는 경우
result_default <- plspm_ind_spread(jutpls_boot)
print(result_default)

# 이름과 값이 있는 경우
result_named <- plspm_ind_spread(jutpls_boot, set = c(진로동기=2, 자기효능감=3, 진로준비=1.5, 진로태도=2.5))
print(result_named)

# 값만 있는 경우
result_unnamed <- plspm_ind_spread(jutpls_boot, set = c(2, 2, 1.5, 1.5))
print(result_unnamed)

# 부분적으로 값이 있는 경우
result_partial <- plspm_ind_spread(jutpls_boot, set = c(진로동기=2, 자기효능감=3))
print(result_partial)


## End(Not run)



\end{ExampleCode}
\end{Examples}
\inputencoding{utf8}
\HeaderA{plspm\_innerplot}{inner\_plot(Pre-draw a mockup )}{plspm.Rul.innerplot}
%
\begin{Description}
inner\_plot(Pre-draw a mockup )
\end{Description}
%
\begin{Usage}
\begin{verbatim}
plspm_innerplot(
  path_mat,
  edge.label.cex = 2,
  edge.label.position = 0.5,
  edge.labels = TRUE,
  edge.width = 3,
  label.cex = 1.1,
  rotation = 0,
  angle = 90,
  sizeLat = 14,
  groups = NULL,
  layout = "circle",
  negCol = "red",
  fade = FALSE,
  asize = 5,
  borders = TRUE,
  border.width = 4,
  border.color = "gray10",
  labels = NULL,
  label.color = "darkblue",
  posCol = "gray10",
  palette = "pastel",
  curve = 0.2,
  width = 2,
  mar = c(3, 3, 3, 3),
  edge_size = 0.7,
  hypo = TRUE,
  digits = 3
)
\end{verbatim}
\end{Usage}
%
\begin{Arguments}
\begin{ldescription}
\item[\code{path\_mat}] mati

\item[\code{edge.label.cex}] 1.2

\item[\code{edge.label.position}] 0.5

\item[\code{edge.labels}] FALSE

\item[\code{edge.width}] 2

\item[\code{label.cex}] 1.2

\item[\code{rotation}] 1

\item[\code{angle}] rotation angle

\item[\code{groups}] NULL

\item[\code{layout}] spring or layout

\item[\code{negCol}] red

\item[\code{fade}] false

\item[\code{asize}] 4

\item[\code{borders}] TRUE

\item[\code{border.width}] 2

\item[\code{border.color}] gray40

\item[\code{labels}] NULL

\item[\code{label.color}] gray30

\item[\code{posCol}] black

\item[\code{palette}] pastel

\item[\code{curve}] 0.2

\item[\code{width}] 1

\item[\code{mar}] c(3,3,3,3)

\item[\code{edge\_size}] edge.width size control

\item[\code{hypo}] paste H

\item[\code{digits}] 3\#'

\item[\code{vsize}] 8
\end{ldescription}
\end{Arguments}
%
\begin{Value}
graph
\end{Value}
%
\begin{Examples}
\begin{ExampleCode}

## Not run: 

# path matrix
IMAG = c(0,0,0,0,0,0)
EXPE = c(1,0,0,0,0,0)
QUAL = c(0,1,0,0,0,0)
VAL = c(0,1,1,0,0,0)
SAT = c(1,1,1,1,0,0)
LOY = c(1,0,0,0,1,0)
sat_path = rbind(IMAG, EXPE, QUAL, VAL, SAT, LOY)

sat_path%>%plspm_innerplot()
#'
path1 = plspm_path(
  row_names = c("IMAG","EXPE","QUAL","VAL","SAT","LOY"),

  relationship = list(
    path(from="IMAG", to=c("EXPE","SAT","LOY")),
    path("EXPE", c("QUAL","VAL","SAT")),
    path("QUAL", c("VAL","SAT")),
    path("VAL",c("SAT")),
    path("SAT","LOY")
  )
)
path1
path1%>%plspm_innerplot()

#'
model1 = "
EXPE ~ IMAG
QUAL ~ EXPE
VAL ~  QUAL+ EXPE
SAT ~ IMAG + EXPE + QUAL + VAL
LOY ~ SAT + IMAG
"
# Generate the matrix
lavpath = plspm_lav2path(model1)


lavpath%>%plspm_innerplot()

#color
sat_path %>%plspm_innerplot(groups=1:6)
sat_path %>%plspm_innerplot(groups=1:6, hypo= FALSE ) #putput 1


sat_path %>%
 plspm_innerplot(groups=1:6, rotation = 3 )


lay_p = matrix(c('IMA', NA, NA, NA, NA, NA, NA, NA,
                 NA, NA, NA, NA, NA, NA, NA, NA, NA, NA, NA, 'EXP', NA,
                 NA, NA, NA, NA, NA, NA, NA, NA, NA, NA, NA, NA, NA, 'QUA',
                 NA, NA, NA, NA, NA, NA, NA, NA, NA, NA, NA, NA, NA, NA, NA,
                 NA, NA, NA, NA, NA, NA, NA, NA, NA, 'VAL', NA, NA, NA, NA,
                 NA, NA, NA, NA, NA, NA, NA, NA, NA, NA, 'SAT', NA, NA, NA,
                 NA, NA, NA, NA, NA, NA, NA, NA, NA, NA, NA, NA, NA, 'LOY',
                 NA, NA, NA, NA, NA, NA, NA, NA),
               nrow = 10, ncol = 10, byrow = FALSE)
lay_p
#'
#'

## End(Not run)


\end{ExampleCode}
\end{Examples}
\inputencoding{utf8}
\HeaderA{plspm\_inner\_model\_sig}{Showing significance using PLSPM regression}{plspm.Rul.inner.Rul.model.Rul.sig}
%
\begin{Description}
Showing significance using PLSPM regression
\end{Description}
%
\begin{Usage}
\begin{verbatim}
plspm_inner_model_sig(
  data,
  unite = FALSE,
  rhs = "endogenous",
  lhs = "exogenous",
  digits = 3
)
\end{verbatim}
\end{Usage}
%
\begin{Arguments}
\begin{ldescription}
\item[\code{data}] plspm data

\item[\code{unite}] TRUE

\item[\code{rhs}] from

\item[\code{lhs}] to

\item[\code{digits}] round 3
\end{ldescription}
\end{Arguments}
%
\begin{Value}
data result
\end{Value}
%
\begin{Examples}
\begin{ExampleCode}
## Not run: 
library(plspm)
data(satisfaction)

# path matrix
IMAG = c(0,0,0,0,0,0)
EXPE = c(1,0,0,0,0,0)
QUAL = c(0,1,0,0,0,0)
VAL = c(0,1,1,0,0,0)
SAT = c(1,1,1,1,0,0)
LOY = c(1,0,0,0,1,0)
sat_path = rbind(IMAG, EXPE, QUAL, VAL, SAT, LOY)

# blocks of outer model
sat_blocks = list(1:5, 6:10, 11:15, 16:19, 20:23, 24:27)
sat_mod=c(rep("A", ncol(sat_path)))
# apply plspm
satpls = plspm(satisfaction, sat_path, sat_blocks, modes = sat_mod,
               scaled = FALSE)

#default
satpls%>%plspm_paths_sig()
satpls$inner_model %>%plspm_paths_sig()

satpls %>%plspm_paths_sig(unite=FALSE, lhs="row", rhs="col")
satpls$inner_model  %>%innermodel_sig(unite=FALSE, lhs="row", rhs="col")

satpls %>%plspm_paths_sig(unite=FALSE, lhs="row", rhs="col")
satpls$inner_model %>%plspm_paths_sig(unite=FALSE, lhs="row", rhs="col")

#'

## End(Not run)

\end{ExampleCode}
\end{Examples}
\inputencoding{utf8}
\HeaderA{plspm\_interaction}{PLS-PMinteraction effect  Partial Least Squares Path Modeling}{plspm.Rul.interaction}
%
\begin{Description}
PLS-PMinteraction effect  Partial Least Squares Path Modeling
\end{Description}
%
\begin{Usage}
\begin{verbatim}
plspm_interaction(
  Data,
  path_matrix,
  blocks,
  interactionTerm = NULL,
  indicator_prod = NULL,
  inter_modes = NULL,
  method = "two_stage",
  modes = rep("A", ncol(path_matrix)),
  scaling = NULL,
  scheme = "centroid",
  scaled = TRUE,
  tol = 1e-06,
  maxiter = 100,
  plscomp = NULL,
  boot.val = TRUE,
  br = 500,
  seed = NULL,
  dataset = TRUE
)
\end{verbatim}
\end{Usage}
%
\begin{Arguments}
\begin{ldescription}
\item[\code{Data}] matrix or data frame containing the manifest variables.

\item[\code{path\_matrix}] A square (lower triangular) boolean matrix
representing the inner model (i.e. the path relationships between latent variables).

\item[\code{blocks}] list of vectors with column indices
or column names from Data indicating the sets of
manifest variables forming each block (i.e.
which manifest variables correspond to each block).\#'

\item[\code{interactionTerm}] interaction terms

\item[\code{indicator\_prod}] indicator producst

\item[\code{inter\_modes}] new modes auto or direct input

\item[\code{method}] Two\_stage, indicator

\item[\code{modes}] optional argument for runing the non-metric approach;
it is a list of string vectors indicating the type of measurement
scale for each manifest variable specified in blocks. scaling must be
specified when working with non-metric variables. Possible values:
"num" (linear transformation, suitable for numerical variables),
"raw" (no transformation), "nom" (non-monotonic transformation,
suitable for nominal variables), and "ord" (monotonic transformation,
suitable for ordinal variables).

\item[\code{scaling}] character vector indicating the type of measurement f
or each block. Possible values are: "A", "B", "newA", "PLScore",
"PLScow". The length of modes must be equal to the length of blocks.

\item[\code{scheme}] string indicating the type of inner weighting scheme.
Possible values are "centroid", "factorial", or "path".

\item[\code{scaled}] standardized. Only used when scaling = NULL.
When (TRUE, data is scaled to standardized values (mean=0 and variance=1).
The variance is calculated dividing by N instead of N-1).

\item[\code{tol}] decimal value indicating the tolerance criterion
for the iterations (tol=0.000001). Can be specified between 0 and 0.001.

\item[\code{maxiter}] integer indicating the maximum number of iterations
(maxiter=100 by default). The minimum value of maxiter is 100.

\item[\code{plscomp}] optional vector indicating the number of PLS components
(for each block) to be used when handling non-metric data
(only used if scaling is provided)

\item[\code{boot.val}] whether bootstrap validation should be performed.
(FALSE by default).

\item[\code{br}] number bootstrap resamples. Used only when boot.val=TRUE. When boot.val=TRUE, the default number of re-samples is 100.

\item[\code{dataset}] whether the data matrix used in the computations should be retrieved (TRUE by default).
\end{ldescription}
\end{Arguments}
%
\begin{Value}
resutl
\end{Value}
%
\begin{Examples}
\begin{ExampleCode}


## Not run: 
#'

data(satisfaction)
sat_path3 = plspm_paths(
  row_names = c("IMAG","EXPE","QUAL","VAL","SAT","LOY"),
  relationship = list(
    path(from="IMAG", to=c("EXPE","SAT","LOY")),
    path("EXPE", c("QUAL","VAL","SAT")),
    path("QUAL", c("VAL","SAT")),
    path("VAL",c("SAT")),
    path("SAT","LOY"))
)

sat_path2

# blokcs
sat_blocks1 <- plspm_blocks(
  IMAG = item(paste0("imag",1:5)),
  EXPE = item(paste0("expe", 1:5)),
  QUAL = item( paste0("qual", 1:5)),
  VAL = item(paste0("val", 1:4)),
  SAT = item(paste0("sat", 1:4)),
  LOY = item(paste0("loy", 1:4))
)

sat_blocks1

# vector of modes (reflective indicators):auto
sat_mod = rep("A", 6)

resinter = plspm_interaction(Data=satisfaction,
                     sat_path2, blocks = sat_blocks1, modes= sat_mod,
                     interactionTerm = list(list(from = "IMAG*EXPE", to = "LOY"),
                     list(from = "IMAG*QUAL", to = "LOY")))
resinter%>%summary()
x11()
plot(resinter)

# resinter %>%plspm_boot_effect()%>%cut_print(rm_row=TRUE)
resinter %>%plspm_effect()%>%cut_print(rm_row=TRUE)

plspm_paths_sig(resinter, unite=TRUE)
plspm_paths_sig_plot(resinter)
plspm_path_coefs_plot(resinter)

plspm_semPaths2(resinter, structural = T)


#'
## Two stage approach
pls_ts = plspm_interaction(satisfaction, sat_path2, sat_blocks1,
               interactionTerm = list(list(from = "IMAG*EXPE", to = "LOY")) ,
               indicator_prod = list(intIE = interlist("IMAG",1,2, "EXPE",1,2)),
               method="two_stage")
pls_ts %>%summary()


plspm_path_coefs_plot(pls_ts,"circle")
plspm_semPaths2(pls_ts)

# plspm_boot_factor_layout(pls_ts)

## indicator arpproach
pls_int = plspm_interaction(satisfaction, sat_path2, sat_blocks1,
                  interactionTerm = list(list(from = "IMAG*EXPE", to = "LOY")) ,
                indicator_prod = list(intIE = interlist("IMAG", "EXPE")),
                  method="indicator")
pls_int %>%summary()


plspm_path_coefs_plot(pls_int, "circle")
plspm_path_coefs_plot(pls_int)
plspm_semPaths2(pls_int)
#'
#'



## End(Not run)


\end{ExampleCode}
\end{Examples}
\inputencoding{utf8}
\HeaderA{plspm\_lav2blocks}{lavaan syntax to plspm blocks}{plspm.Rul.lav2blocks}
%
\begin{Description}
lavaan syntax to plspm blocks
\end{Description}
%
\begin{Usage}
\begin{verbatim}
plspm_lav2blocks(blocks_model)
\end{verbatim}
\end{Usage}
%
\begin{Arguments}
\begin{ldescription}
\item[\code{blocks\_model}] lavaan syntax
\end{ldescription}
\end{Arguments}
%
\begin{Value}
list
\end{Value}
%
\begin{Examples}
\begin{ExampleCode}
## Not run: 

# Input data
blocks_model <- "
IMAG =~ imag1 + imag2 + imag3 + imag4 + imag5
EXPE =~ expe1 + expe2 + expe3 + expe4 + expe5
QUAL =~ qual1 + qual2 + qual3 + qual4 + qual5
VAL =~ val1 + val2 + val3 + val4
SAT =~ sat1 + sat2 + sat3 + sat4
LOY =~ loy1 + loy2 + loy3 + loy4
"

# Example function calls
plspm_lav2blocks(blocks_model)


## End(Not run)
\end{ExampleCode}
\end{Examples}
\inputencoding{utf8}
\HeaderA{plspm\_lav2paths}{lavaan syntax to plspm path matrix}{plspm.Rul.lav2paths}
%
\begin{Description}
lavaan syntax to plspm path matrix
\end{Description}
%
\begin{Usage}
\begin{verbatim}
plspm_lav2paths(model_syntax, fct_order = NULL)
\end{verbatim}
\end{Usage}
%
\begin{Arguments}
\begin{ldescription}
\item[\code{model\_syntax}] lavaan syntax

\item[\code{fct\_order}] fct\_order is reorder option for rownamnes of matrix
\end{ldescription}
\end{Arguments}
%
\begin{Value}
matrix
\end{Value}
%
\begin{Examples}
\begin{ExampleCode}
## Not run: 
#'
# Example usage
model1 <- "
EXPE ~ IMAG; QUAL ~ EXPE
VAL ~ QUAL + EXPE
SAT ~ IMAG + EXPE + QUAL + VAL
LOY ~ SAT + IMAG
"

 #or
model1 = "
EXPE ~ IMAG
QUAL ~ EXPE
VAL ~  QUAL+ EXPE
SAT ~ IMAG + EXPE + QUAL + VAL
LOY ~ SAT + IMAG
"
# Generate the matrix
plspm_lav2paths(model1)

plspm_lav2paths(model1, fct_order=c("IMAG", "EXPE","QUAL", "VAL", "SAT", "LOY" ))



## End(Not run)

\end{ExampleCode}
\end{Examples}
\inputencoding{utf8}
\HeaderA{plspm\_layout\_mat}{layout\_mat}{plspm.Rul.layout.Rul.mat}
%
\begin{Description}
layout\_mat
\end{Description}
%
\begin{Usage}
\begin{verbatim}
plspm_layout_mat(semplot, text = TRUE, byrow = TRUE, variable = TRUE)
\end{verbatim}
\end{Usage}
%
\begin{Arguments}
\begin{ldescription}
\item[\code{semplot}] semPaths2 data

\item[\code{text}] text True output text matrix

\item[\code{byrow}] byrow TRUE

\item[\code{variable}] variable is rownames adapt
\end{ldescription}
\end{Arguments}
%
\begin{Value}
data
\end{Value}
%
\begin{Examples}
\begin{ExampleCode}


## Not run: 
#'
data(offense)
offense %>%str()

nfl_path = matrix(
  c(0, 0, 0, 0, 0,
    0, 0, 0, 0, 0,
    0, 0, 0, 0, 0,
    0, 1, 1, 0, 0,
    1, 0, 0, 1, 0),
  nrow = 5, ncol = 5, byrow = TRUE )

rownames(nfl_path) <- c("Speical", "Rushing", "Passing", "Offense", "Scoring")
colnames(nfl_path) <- c("Speical", "Rushing", "Passing", "Offense", "Scoring")
nfl_path

nfl_blocks = list(7:8, 1:3, 4:6,1:6, 9:11)

plspm_blocks_match(data=offense, blocks=nfl_blocks,
                   paths=nfl_path)

nfl_mode =c("B","A","A","A","A")

nfl_pls = plspm(Data=offense, nfl_path, nfl_blocks, modes = nfl_mode)
nfl_pls%>%summary()

#########
nfl_pls_boot = plspm_sem(Data=offense, nfl_path, nfl_blocks, modes = nfl_mode)

#Layouts extracted from mockups
plspm_semPaths2(nfl_pls_boot,sizeLat = 6)%>%
  plspm_factor_layout()

#Make it input material
plspm_semPaths2(nfl_pls_boot,sizeLat =6)%>%
  plspm_factor_layout() %>%plspm_layout_mat(text=T)

#Organizing materials for input
plspm_semPaths2(nfl_pls_boot,sizeLat =6)%>%
  plspm_factor_layout() %>%plspm_layout_mat()%>%
  mat_text_arrange()

## End(Not run)

\end{ExampleCode}
\end{Examples}
\inputencoding{utf8}
\HeaderA{plspm\_loadings}{plspm loadings output}{plspm.Rul.loadings}
%
\begin{Description}
plspm loadings output
\end{Description}
%
\begin{Usage}
\begin{verbatim}
plspm_loadings(plsres, digits = 3, unite = TRUE)
\end{verbatim}
\end{Usage}
%
\begin{Arguments}
\begin{ldescription}
\item[\code{plsres}] plspm result

\item[\code{digits}] round =3

\item[\code{unite}] unite t+star
\end{ldescription}
\end{Arguments}
%
\begin{Value}
data table
\end{Value}
%
\begin{Examples}
\begin{ExampleCode}
## Not run: 

plsres %>% plspm_loadings()
plsres$boot %>% plspm_loadings()
plsres$boot$paths %>% plspm_loadings()



## End(Not run)

\end{ExampleCode}
\end{Examples}
\inputencoding{utf8}
\HeaderA{plspm\_loadings\_plot}{plspm\_loadings\_plot}{plspm.Rul.loadings.Rul.plot}
%
\begin{Description}
plspm\_loadings\_plot
\end{Description}
%
\begin{Usage}
\begin{verbatim}
plspm_loadings_plot(
  plsres,
  size.x = 14,
  size.y = 12,
  size.legend = 12,
  size.text = 5,
  show = FALSE
)
\end{verbatim}
\end{Usage}
%
\begin{Arguments}
\begin{ldescription}
\item[\code{plsres}] plspm data

\item[\code{size.x}] axis text 14

\item[\code{size.y}] 12

\item[\code{size.legend}] 12

\item[\code{size.text}] show data
\end{ldescription}
\end{Arguments}
%
\begin{Value}
plot
\end{Value}
%
\begin{Examples}
\begin{ExampleCode}

## Not run: 

data(satisfaction)
#model 1
model1 = "
EXPE ~ IMAG
QUAL ~ EXPE
VAL ~  QUAL+ EXPE
SAT ~ IMAG + EXPE + QUAL + VAL
LOY ~ SAT + IMAG
"
# # Generate the matrix
sat_path = plspm_lav2path(model1,
                          fct_order=c("IMAG", "EXPE","QUAL", "VAL", "SAT", "LOY" ))
# sat_path

# Method 2: paths
sat_path = plspm_paths(
  row_names = c("IMAG","EXPE","QUAL","VAL","SAT","LOY"),
  relationship = list(
    path(from="IMAG", to=c("EXPE","SAT","LOY")),
    path("EXPE", c("QUAL","VAL","SAT")),
    path("QUAL", c("VAL","SAT")),
    path("VAL",c("SAT")),
    path("SAT","LOY")
  )
)
# blokcs
sat_blocks1 <- plspm_blocks(
  IMAG = item(paste0("imag",1:5)),
  EXPE = item(paste0("expe", 1:5)),
  QUAL = item( paste0("qual", 1:5)),
  VAL = item(paste0("val", 1:4)),
  SAT = item(paste0("sat", 1:4)),
  LOY = item(paste0("loy", 1:4))
)
sat_blocks1

# vector of modes (reflective indicators):auto
sat_mod = rep("A", 6)
#############

# apply plspm
satpls = plspm(satisfaction,
                   path_matrix = sat_path,
                   blocks = sat_blocks1,
                   scaled = FALSE)

satpls_boot = plspm_sem(satisfaction, sat_path, sat_blocks1, scaled = FALSE,
                        boot.val =TRUE, br=100)


#plot
satpls %>%plspm_loadings_plot()
satpls$outer_model%>%plspm_loadings_plot()

 
## End(Not run)

\end{ExampleCode}
\end{Examples}
\inputencoding{utf8}
\HeaderA{plspm\_loadings\_values}{plspm lodaings significant  vecgtor or df}{plspm.Rul.loadings.Rul.values}
%
\begin{Description}
plspm lodaings significant  vecgtor or df
\end{Description}
%
\begin{Usage}
\begin{verbatim}
plspm_loadings_values(plsres, type = "df")
\end{verbatim}
\end{Usage}
%
\begin{Arguments}
\begin{ldescription}
\item[\code{plsres}] plspm result

\item[\code{type}] df, vec, data.frame, vector
\end{ldescription}
\end{Arguments}
%
\begin{Value}
data
\end{Value}
%
\begin{Examples}
\begin{ExampleCode}

## Not run: 

# Method 2: paths
sat_path = plspm_paths(
  row_names = c("IMAG","EXPE","QUAL","VAL","SAT","LOY"),
  relationship = list(
    path(from="IMAG", to=c("EXPE","SAT","LOY")),
    path("EXPE", c("QUAL","VAL","SAT")),
    path("QUAL", c("VAL","SAT")),
    path("VAL",c("SAT")),
    path("SAT","LOY")
  )
)
# blokcs
sat_blocks1 <- plspm_blocks(
  IMAG = item(paste0("imag",1:5)),
  EXPE = item(paste0("expe", 1:5)),
  QUAL = item( paste0("qual", 1:5)),
  VAL = item(paste0("val", 1:4)),
  SAT = item(paste0("sat", 1:4)),
  LOY = item(paste0("loy", 1:4))
)
sat_blocks1

# vector of modes (reflective indicators):auto
sat_mod = rep("A", 6)
#############

# apply plspm
satpls = plspm_sem(satisfaction,
                   path_matrix = sat_path,
                   blocks = sat_blocks1,
                   scaled = FALSE)

satpls_boot = plspm_sem(satisfaction, sat_path, sat_blocks1, scaled = FALSE,
                        boot.val =TRUE, br=100)




# plspm_loadings_values
satpls %>%plspm_loadings_values()
satpls %>%plspm_loadings_values("vec")
satpls$boot %>%plspm_loadings_values()
satpls %>%plspm_loadings_values()
satpls$boot$loadings%>%plspm_loadings_values()
satpls$boot$loadings%>%plspm_loadings_values("vec")


## End(Not run)
\end{ExampleCode}
\end{Examples}
\inputencoding{utf8}
\HeaderA{plspm\_loading\_lav}{plspm\_loading to lavaan syntax}{plspm.Rul.loading.Rul.lav}
%
\begin{Description}
plspm\_loading to lavaan syntax
\end{Description}
%
\begin{Usage}
\begin{verbatim}
plspm_loading_lav(input_data, col = "loading")
\end{verbatim}
\end{Usage}
%
\begin{Arguments}
\begin{ldescription}
\item[\code{input\_data}] plspm result

\item[\code{col}] selected col
\end{ldescription}
\end{Arguments}
%
\begin{Value}
lavaan syntax
\end{Value}
%
\begin{Examples}
\begin{ExampleCode}
## Not run: 
#'
data(satisfaction)
sat_path = plspm_paths(
  row_names = c("IMAG","EXPE","QUAL","VAL","SAT","LOY"),
  relationship = list(
    path(from="IMAG", to=c("EXPE","SAT","LOY")),
    path("EXPE", c("QUAL","VAL","SAT")),
    path("QUAL", c("VAL","SAT")),
    path("VAL",c("SAT")),
    path("SAT","LOY"))
)
sat_path

# blokcs
sat_blocks1 <- plspm_blocks(
  IMAG = item(paste0("imag",1:5)),
  EXPE = item(paste0("expe", 1:5)),
  QUAL = item( paste0("qual", 1:5)),
  VAL = item(paste0("val", 1:4)),
  SAT = item(paste0("sat", 1:4)),
  LOY = item(paste0("loy", 1:4))
)
sat_blocks1

# vector of modes (reflective indicators):auto
sat_mod = rep("A", 6)
#############

# apply plspm
satpls = plspm(satisfaction,
               path_matrix = sat_path,
               blocks = sat_blocks1,
               scaled = FALSE)
#lavaan 1
satpls%>%plspm_loading_lav() %>%cat("\n")

#method 2, lavan
satpls$outer_model%>%plspm_loading_lav() %>%cat("\n")

#diagram
x11()
satpls%>%plspm_loading_lav() %>%
  diagram_model(whatLabels = "est",exoCov=F,
                edge.label.cex = 0.6, rotation=4,
                sizeLat=6,
                sizeMan = 6, sizeMan2=1.5, mar=c(1,10,1,10))

## End(Not run)


\end{ExampleCode}
\end{Examples}
\inputencoding{utf8}
\HeaderA{plspm\_make\_layout\_mat}{plspm\_make\_layout\_mat}{plspm.Rul.make.Rul.layout.Rul.mat}
%
\begin{Description}
plspm\_make\_layout\_mat
\end{Description}
%
\begin{Usage}
\begin{verbatim}
plspm_make_layout_mat(...)
\end{verbatim}
\end{Usage}
%
\begin{Arguments}
\begin{ldescription}
\item[\code{...}] variable
\end{ldescription}
\end{Arguments}
%
\begin{Value}
mat
\end{Value}
%
\begin{Examples}
\begin{ExampleCode}

## Not run: 
#'
plspm_make_layout_mat(x1 = c(1, 1),
                      x2 = c(2, 1),
                      x3 = c(2, 2),
                      x4 = c(3, 2),
                      x5 = c(3, 5))

## End(Not run)


\end{ExampleCode}
\end{Examples}
\inputencoding{utf8}
\HeaderA{plspm\_make\_prod}{plspm\_make\_prod , Generate data for indicator interaction multiplied terms}{plspm.Rul.make.Rul.prod}
%
\begin{Description}
plspm\_make\_prod , Generate data for indicator interaction multiplied terms
\end{Description}
%
\begin{Usage}
\begin{verbatim}
plspm_make_prod(data, blocks, indicator, intname = "int", n_name = 2)
\end{verbatim}
\end{Usage}
%
\begin{Arguments}
\begin{ldescription}
\item[\code{data}] ori data

\item[\code{blocks}] ori blocks

\item[\code{indicator}] indicator product

\item[\code{intname}] "int" default

\item[\code{n\_name}] name 2, more 3  int\_AAABBB
\end{ldescription}
\end{Arguments}
%
\begin{Value}
data
\end{Value}
%
\begin{Examples}
\begin{ExampleCode}

## Not run: 
#' data(satisfaction)
#상황이 바뀌는 경우, 데이터를 바꾸어서 생성하는 경우,
plspm_make_prod(satisfaction,
                sat_blocks1,
                indicator= list(intIE = interlist("IMAG",1,2, "EXPE",1,2),
                                intIV = interlist("IMAG",1,2, "QUAL",1)  )
) %>%str()

## End(Not run)

\end{ExampleCode}
\end{Examples}
\inputencoding{utf8}
\HeaderA{plspm\_mat\_lav}{plspm\_mat\_lavaan syntax}{plspm.Rul.mat.Rul.lav}
%
\begin{Description}
plspm\_mat\_lavaan syntax
\end{Description}
%
\begin{Usage}
\begin{verbatim}
plspm_mat_lav(data, free = FALSE, hypo = FALSE)
\end{verbatim}
\end{Usage}
%
\begin{Arguments}
\begin{ldescription}
\item[\code{data}] matrix

\item[\code{free}] TRUE is free parameter

\item[\code{hypo}] hypothesis paste
\end{ldescription}
\end{Arguments}
%
\begin{Value}
syntax lavaan
\end{Value}
%
\begin{Examples}
\begin{ExampleCode}

## Not run: 
#'

# Example data for testing
data_b <- matrix(c(
  0.000, 0.000, 0.000, 0.000, 0.000, 0,
  0.471, 0.000, 0.000, 0.000, 0.000, 0,
  0.000, 0.834, 0.000, 0.000, 0.000, 0,
  0.000, 0.046, 0.701, 0.000, 0.000, 0,
  0.245, -0.017, 0.222, 0.527, 0.000, 0,
  0.182, 0.000, 0.000, 0.000, 0.628, 0
), nrow = 6, byrow = TRUE,
dimnames=list(
  c("IMAG", "EXPE", "QUAL", "VAL", "SAT", "LOY"),
  c("IMAG", "EXPE", "QUAL", "VAL", "SAT", "LOY")
))
plspm_mat_lav(data_b)%>%cat()
plspm_mat_lav(data_b)%>%
  diagram_model(whatLabels = "est", layout="tree2", nDigits=3,
                edge.label.position=0.65, rotation=1)


data_r <- matrix(c(
  0.000, 0.000, 0.000, 0.000, 0.000, 0,
  0.471, 0.000, 0.000, 0.000, 0.000, 0,
  0.000, 0.834, 0.000, 0.000, 0.000, 0,
  0.000, 0.046, 0.701, 0.000, 0.000, 0,
  0.245, -0.017, 0.222, 0.527, 0.000, 0,
  0.182, 0.000, 0.000, 0.000, 0.628, 0
), nrow = 6, byrow = TRUE,
dimnames=list(c("IMAG", "EXPE", "QUAL", "VAL", "SAT", "LOY")))
data_r

plspm_mat_lav(data_r)%>%cat()

plspm_mat_lav(data_r)%>%
  diagram_model(whatLabels = "est", layout="tree2", nDigits=3,
                edge.label.position=0.65, rotation=1)

data_n <- matrix(c(
  0.000, 0.000, 0.000, 0.000, 0.000, 0,
  0.471, 0.000, 0.000, 0.000, 0.000, 0,
  0.000, 0.834, 0.000, 0.000, 0.000, 0,
  0.000, 0.046, 0.701, 0.000, 0.000, 0,
  0.245, -0.017, 0.222, 0.527, 0.000, 0,
  0.182, 0.000, 0.000, 0.000, 0.628, 0
), nrow = 6, byrow = TRUE)
data_n
plspm_mat_lav(data_n)%>%cat()
plspm_mat_lav(data_n)%>%
  diagram_model(whatLabels = "est", layout="tree2", nDigits=3,
                edge.label.position=0.65, rotation=1)



data_null <- matrix(c(
  0.000, 0.000, 0.000, 0.000, 0.000, 0,
  "H1", 0.000, 0.000, 0.000, 0.000, 0,
  0.000, "H2", 0.000, 0.000, 0.000, 0,
  0.000, "H3", "H4", 0.000, 0.000, 0,
  "H5", "H6", "H7", "H8", 0.000, 0,
  "H9", 0.000, 0.000, 0.000, "H10", 0
), nrow = 6, byrow = TRUE,
dimnames=list(
  c("IMAG", "EXPE", "QUAL", "VAL", "SAT", "LOY"),
  c("IMAG", "EXPE", "QUAL", "VAL", "SAT", "LOY")
))

plspm_mat_lav(data_null)%>%cat()

plspm_mat_lav(data_null)%>%
  diagram_model(whatLabels = "model", layout="tree2", nDigits=3,
                edge.label.position=0.65, rotation=1)

#'
lay_p = matrix(c('IMA', NA, NA, NA, NA, NA, NA, NA,
                 NA, NA, NA, NA, NA, NA, NA, NA, NA, NA, NA, 'EXP', NA,
                 NA, NA, NA, NA, NA, NA, NA, NA, NA, NA, NA, NA, NA, 'QUA',
                 NA, NA, NA, NA, NA, NA, NA, NA, NA, NA, NA, NA, NA, NA, NA,
                 NA, NA, NA, NA, NA, NA, NA, NA, NA, 'VAL', NA, NA, NA, NA,
                 NA, NA, NA, NA, NA, NA, NA, NA, NA, NA, 'SAT', NA, NA, NA,
                 NA, NA, NA, NA, NA, NA, NA, NA, NA, NA, NA, NA, NA, 'LOY',
                 NA, NA, NA, NA, NA, NA, NA, NA),
               nrow = 10, ncol = 10, byrow = FALSE)
lay_p

# Using digaram_model

satpls$path_coefs %>% plspm_mat_lav() %>%
  diagram_model(layout="tree2", whatLabels = "est",
                rotation =1, edge.label.position=0.6)

satpls$path_coefs %>% plspm_mat_lav() %>%
  diagram_model(layout="tree2", whatLabels = "est",
                rotation =1, edge.label.position=0.6,
                edgeLabels=satpls_fit$inner_model %>%plspm_edge_values())
#'
#' #'
data <- matrix(c(0, 1, 0, 1, 0, 1, 1, 0, 0), nrow=3, byrow=TRUE)
rownames(data) <- c("자기효능감", "진로동기", "진로태도")
colnames(data) <- c("진로동기", "진로태도", "진로준비")
data
# hypo=TRUE인 경우
cat(plspm_mat_lav(data, hypo=TRUE))
# 진로동기 ~ H1*자기효능감
# 진로태도 ~ H2*자기효능감
# 진로준비 ~ H3*자기효능감
# 진로준비 ~ H4*진로동기
# 진로준비 ~ H5*진로태도

# free=TRUE, hypo=FALSE인 경우
cat(plspm_mat_lav(data, free=TRUE, hypo=FALSE))
# 진로동기 ~ 자기효능감
# 진로태도 ~ 자기효능감
# 진로준비 ~ 자기효능감
# 진로준비 ~ 진로동기
# 진로준비 ~ 진로태도

# free=FALSE인 경우
cat(plspm_mat_lav(data, free=FALSE, hypo=FALSE))
# 진로동기 ~ 1*자기효능감
# 진로태도 ~ 1*자기효능감
# 진로준비 ~ 1*자기효능감
# 진로준비 ~ 1*진로동기
# 진로준비 ~ 1*진로태도


## End(Not run)

\end{ExampleCode}
\end{Examples}
\inputencoding{utf8}
\HeaderA{plspm\_modelfit}{Calculate Model Fit Indices for PLS Models}{plspm.Rul.modelfit}
%
\begin{Description}
This function computes various model fit indices for a Partial Least Squares (PLS) model, including SRMR, NFI, d\_ULS, and d\_G.
Bootstrap confidence intervals for d\_ULS and d\_G are also provided.
\end{Description}
%
\begin{Usage}
\begin{verbatim}
plspm_modelfit(
  pls_model,
  blocks,
  bootstrap = 1000,
  ci = 0.95,
  digits = 5,
  trans = TRUE
)
\end{verbatim}
\end{Usage}
%
\begin{Arguments}
\begin{ldescription}
\item[\code{pls\_model}] A PLS model object containing model results.

\item[\code{blocks}] A list defining the measurement model blocks. Each block should contain variable names.

\item[\code{bootstrap}] The number of bootstrap samples to use. Default is 1000.

\item[\code{ci}] The confidence level for confidence intervals. Default is 0.95.

\item[\code{digits}] The number of decimal places to round results. Default is 5.

\item[\code{trans}] If TRUE, transforms the result into a transposed data frame. Default is TRUE.
\end{ldescription}
\end{Arguments}
%
\begin{Value}
A data frame containing the computed fit indices and their evaluations.
\end{Value}
\inputencoding{utf8}
\HeaderA{plspm\_model\_lav}{plspm\_model\_lav}{plspm.Rul.model.Rul.lav}
%
\begin{Description}
plspm\_model\_lav
\end{Description}
%
\begin{Usage}
\begin{verbatim}
plspm_model_lav(paths, blocks)
\end{verbatim}
\end{Usage}
%
\begin{Arguments}
\begin{ldescription}
\item[\code{paths}] paths

\item[\code{blocks}] blocks
\end{ldescription}
\end{Arguments}
%
\begin{Value}
lavaan syntax
\end{Value}
%
\begin{Examples}
\begin{ExampleCode}

## Not run: 
#'
jutpath1 = plspm_paths(
  row_names = c("자기효능감","진로동기","진로태도","진로준비"),
  relationship = list(
    path(from="자기효능감", to=c("진로동기","진로태도","진로준비")),
    path("진로동기", c("진로준비")),
    path("진로태도", c("진로준비"))
  )  )
jutpath1

#측정모형
jut_blocks <- plspm_blocks(
  자기효능감 = item("C_S1","C_S2","C_S3","C_S4","C_S5"),
  진로동기 = item("A_M1","A_M2","A_M3"),
  진로태도 = item(  "B02", "B04", "B10","B14" ),#결정성
  진로준비 = item("D_P1","D_P2","D_P3")
)

jut_blocks
#'
plspm_model_syntax(jutpath1, jut_blocks) %>%cat()

## End(Not run)

\end{ExampleCode}
\end{Examples}
\inputencoding{utf8}
\HeaderA{plspm\_model\_plot}{Drawing Structural Equation Models Using the lavaan Syntax (park JH PhD)}{plspm.Rul.model.Rul.plot}
%
\begin{Description}
Drawing Structural Equation Models Using the lavaan Syntax (park JH PhD)
\end{Description}
%
\begin{Usage}
\begin{verbatim}
plspm_model_plot(
  model,
  residuals = FALSE,
  whatLabels = "model",
  rotation = 2,
  residScale = 12,
  layout = "tree2",
  edge.label.cex = 1.2,
  edge.label.position = 0.55,
  edge.color = "black",
  edge.width = 1.5,
  sizeMan = 12,
  sizeMan2 = 6,
  sizeLat = 10,
  sizeLat2 = 5,
  label.cex = 1,
  border.width = 2,
  color = list(lat = "gray10", man = "gray90"),
  mar = c(4, 4, 4, 4),
  style = "lisrel",
  shapeLat = "ellipse",
  shapeMan = "circle",
  shapeInt = "triangle",
  exoVar = FALSE,
  exoCov = FALSE,
  curve = 1.5,
  asize = 2,
  nDigits = 3,
  sample.nobs = 100,
  Groups = FALSE,
  group_match = "lat",
  growth = FALSE,
  structural = FALSE,
  edgeLabels = NULL,
  nodeLabels = NULL,
  fixedStyle = 1,
  sig = FALSE,
  type = "plot"
)
\end{verbatim}
\end{Usage}
%
\begin{Arguments}
\begin{ldescription}
\item[\code{residuals}] residauls TRUE

\item[\code{whatLabels}] "model

\item[\code{rotation}] rotation=2

\item[\code{residScale}] residScale=12

\item[\code{layout}] layout="tree"

\item[\code{edge.label.cex}] 1

\item[\code{edge.label.position}] 0.5

\item[\code{edge.color}] edge.color= black

\item[\code{edge.width}] edge.width=1.5

\item[\code{sizeMan}] 8

\item[\code{sizeMan2}] 4

\item[\code{sizeLat}] 10

\item[\code{sizeLat2}] 6

\item[\code{label.cex}] label.cex

\item[\code{border.width}] border.width=1

\item[\code{color}] list(lat="skyblue", man="yellow")

\item[\code{mar}] c(2,8,3,10)

\item[\code{style}] lisrel

\item[\code{shapeLat}] "circle, ellipse

\item[\code{shapeMan}] retangle

\item[\code{shapeInt}] triangle ,

\item[\code{exoVar}] FALSE

\item[\code{exoCov}] TRUE

\item[\code{curve}] 1.5

\item[\code{asize}] 1.5

\item[\code{nDigits}] 2

\item[\code{sample.nobs}] sample.nobs = 100

\item[\code{Groups}] FALSE

\item[\code{group\_match}] "FALSE"lat\#'

\item[\code{growth}] growth= FALSE, TRUE  model.type ="grouwth"

\item[\code{structural}] structural = FALSE

\item[\code{edgeLabels}] edgeLabels=NULL

\item[\code{nodeLabels}] nodeLabels=NULL

\item[\code{fixedStyle}] fixedStyle=1 default rbase linetype

\item[\code{sig}] FALSE

\item[\code{type}] type ="plot, and res is lavaan simdata result

\item[\code{plsres\_boot}] plsres\_boot to model syntax
\end{ldescription}
\end{Arguments}
%
\begin{Value}
plot jutpls\_boot
\end{Value}
%
\begin{Examples}
\begin{ExampleCode}
## Not run: 
#'
model_test <- '
  # latent variable definitions
     ind60 =~ x1 + x2 + x3
     dem60 =~ y1 + a*y2 + b*y3 + c*y4
     dem65 =~ y5 + a*y6 + b*y7 + c*y8

  # regressions
    dem60 ~ ind60
    dem65 ~ ind60 + dem60

  # residual correlations
    # y1 ~~ y5
    # y2 ~~ y4 + y6
    # y3 ~~ y7
    # y4 ~~ y8
    # y6 ~~ y8
'

HSmodel <- ' visual  =~ x1 + x2 + x3
             textual =~ x4 + x5 + x6
             speed   =~ x7 + x8 + x9 '


m1a <- "
y1 =~ x1 + x2 +x5
y2 =~ x3 + x4 +x6
y1 ~~ y2"

diagram_model(model= HSmodel, groups=F,  residuals=F)

diagram_model(model= model_test, style="ram", groups=F)
diagram_model(model= model_test, residuals=T)
diagram_model(model= m1a, residuals=T)


#plspm model
#model 1
model1 = "
EXPE ~ IMAG
QUAL ~ EXPE
VAL ~  QUAL+ EXPE
SAT ~ IMAG + EXPE + QUAL + VAL
LOY ~ SAT + IMAG
"
sat_path = plspm_lav2path(model1,
                          fct_order=c("IMAG", "EXPE","QUAL", "VAL", "SAT", "LOY" ))
sat_path

blocks_model ="
IMAG =~ imag1 + imag2 + imag3 + imag4 + imag5
EXPE =~ expe1 + expe2 + expe3 + expe4 + expe5
QUAL =~ qual1 + qual2 + qual3 + qual4 + qual5
VAL =~ val1 + val2 + val3 + val4
SAT =~ sat1 + sat2 + sat3 + sat4
LOY =~ loy1 + loy2 + loy3 + loy4
"
blocks_model
plspm_lav2blocks(blocks_model)

#model 그리기
paste(model1, blocks_model)%>%
  plspm_model_plot(rotation = 1, sizeLat = 6, sizeMan = 4, sizeMan2=3,
                mar=c(2,2,2,2))

#'
#'
#'

## End(Not run)

\end{ExampleCode}
\end{Examples}
\inputencoding{utf8}
\HeaderA{plspm\_multigroups}{plspm에서 3개 그룹을 비교하는 경우, p보정 방법 포함}{plspm.Rul.multigroups}
%
\begin{Description}
plspm에서 3개 그룹을 비교하는 경우, p보정 방법 포함
\end{Description}
%
\begin{Usage}
\begin{verbatim}
plspm_multigroups(
  plspm_boot,
  group,
  method = "bootstrap",
  type = "all",
  adjust = FALSE,
  adjust.method = "bonferroni"
)
\end{verbatim}
\end{Usage}
%
\begin{Arguments}
\begin{ldescription}
\item[\code{plspm\_boot}] bootstap data

\item[\code{group}] group variable set ="grade"

\item[\code{method}] bootstrap, permutation

\item[\code{type}] all, grp1, grp2, grp3, ...

\item[\code{adjust}] default FALSE

\item[\code{adjust.method}] "bonferroni", "holm", "hochberg", "BH", "BY"
\end{ldescription}
\end{Arguments}
%
\begin{Value}
data
\end{Value}
%
\begin{Examples}
\begin{ExampleCode}
## Not run: 
#'
# 사용 예시

#구조모형
jutpath1 = plspm_paths(
  row_names = c("SelfEff","CarMot","CarAtt","CarAct"),
  relationship = list(
    path(from="SelfEff", to=c("CarMot","CarAtt","CarAct")),
    path("CarMot", c("CarAct")),
    path("CarAtt", c("CarAct"))
  )  )
jutpath1

jutpath1 %>% plspm_paths2lav()%>%cat()


# 측정 모형
jut_blocks <- plspm_blocks(
  SelfEff = item("sef01", "sef02", "sef03", "sef04", "sef05"),
  CarMot = item("mot01","mot02","mot03"),
  #   CarAtt = item( "att01", "att02","att03","att04","att05" ),
  CarAtt = item("att03","att04","att05" ),
  CarAct = item("act01", "act02","act03")
)
jut_blocks

jut_blocks %>%plspm_blocks2lav()%>%cat()



# 데이터 분석 : bootstrap
jutpls_boot = plspm_sem(Data= jut7c,
                        path_matrix = jutpath1,
                        blocks = jut_blocks,seed=19988,
                        br = 5000, summary= TRUE)
# p.adjust() 함수: R에서 제공하는 함수로, 선택된 보정 방법에 맞춰 p-value를 보정한다. p.adjust() 함수는 다양한 다중 비교 보정 방법을 지원한다.
# "bonferroni": 보수적인 Bonferroni 보정.
# "holm": Holm 보정.
# "hochberg": Hochberg 보정.
# "BH": Benjamini-Hochberg 보정.
# "BY": Benjamini-Yekutieli 보정.
# Benjamini-Hochberg 보정을 적용하여 결과 조회
plspm_multigroups(jutpls_boot, group = jut7c$LCA_class,
                  type = "all", adjust = TRUE, adjust.method = "bonferroni")

plspm_multigroups(jutpls_boot, group = jut7c$LCA_class,
                  type = "all", adjust = FALSE, adjust.method = "bonferroni")

plspm_multigroups(jutpls_boot, group = jut7c$LCA_class, method="permutation",
                  type = "all", adjust = FALSE, adjust.method = "bonferroni")

plspm_multigroups(jutpls_boot, group = jut7c$LCA_class, type = "grp1")
plspm_multigroups(jutpls_boot, group = jut7c$LCA_class, type = "grp2")
plspm_multigroups(jutpls_boot, group = jut7c$LCA_class, type = "grp3")

#정한 방법
jut_lca_comparison <- plspm_multigroups(jutpls_boot, group = jut7c$LCA_class)
jut_lca_comparison$Class1_vs_Class2$test
jut_lca_comparison$Class1_vs_Class3$test
jut_lca_comparison$Class2_vs_Class3$test


plspm_multigroups(jutpls_boot, group = jut7c$LCA_class,
                  type = "all", adjust = TRUE, adjust.method = "holm")

plspm_multigroups(jutpls_boot, group = jut7c$LCA_class,
                  type = "all", adjust = FALSE, adjust.method = "holm")

plspm_multigroups(jutpls_boot, group = jut7c$LCA_class,
                  type = "all", adjust = TRUE, adjust.method = "hochberg")
plspm_multigroups(jutpls_boot, group = jut7c$LCA_class,
                  type = "all", adjust = FALSE, adjust.method = "hochberg")

plspm_multigroups(jutpls_boot, group = jut7c$LCA_class,
                  type = "all", adjust = TRUE, adjust.method = "BH")

plspm_multigroups(jutpls_boot, group = jut7c$LCA_class,
                  type = "all", adjust = TRUE, adjust.method = "BY")
#'
#'

## End(Not run)
\end{ExampleCode}
\end{Examples}
\inputencoding{utf8}
\HeaderA{plspm\_paths}{plspm path create}{plspm.Rul.paths}
%
\begin{Description}
plspm path create
\end{Description}
%
\begin{Usage}
\begin{verbatim}
plspm_paths(row_names = NULL, col_names = NULL, relationship = NULL)
\end{verbatim}
\end{Usage}
%
\begin{Arguments}
\begin{ldescription}
\item[\code{relationship}] input path

\item[\code{rownames}] rowname

\item[\code{colnames}] colname
\end{ldescription}
\end{Arguments}
%
\begin{Value}
matrix
\end{Value}
%
\begin{Examples}
\begin{ExampleCode}

## Not run: 
edu_path <- plspm_path(
  rownames = c("Support", "Advising", "Tutoring", "Value", "Satisfaction", "Loyalty"),
  colnames = c("Support", "Advising", "Tutoring", "Value", "Satisfaction", "Loyalty"),
  relationship = list(
    list(from = "Value",
         to = c("Support", "Advising", "Tutoring")),
    list(from = "Satisfaction",
         to = c("Support", "Advising", "Tutoring","Value")),
    list(from = "Loyalty",
         to =c("Satisfaction"))
  )
)
edu_path


edu_path <- plspm_path(
  rownames = c("Support", "Advising", "Tutoring", "Value", "Satisfaction", "Loyalty"),
  colnames = c("Support", "Advising", "Tutoring", "Value", "Satisfaction", "Loyalty"),
  relationship = list(

    path(form = "Support", to ="Advising"),
    path(form = "Tutoring", to ="Advising"),
    path(from = "Value",
         to = c("Support", "Advising", "Tutoring")),
    path(from = "Satisfaction",
         to = c("Support", "Advising", "Tutoring","Value")),
    path (from = "Loyalty",
         to =c("Satisfaction"))
  )
)
edu_path
#'

plspm_paths(
  relationship = list(
    list(from="IMAG", to=c("EXPE","SAT","LOY")),
    list(from="EXPE", to=c("QUAL","VAL","SAT")),
    list(from="QUAL", to=c("VAL","SAT")),
    list(from="VAL", to="SAT"),
    list(from="SAT", to="LOY")
  )
)

# 결과 출력

sat_path1 = plspm_paths(
  row_names = c("IMAG","EXPE","QUAL","VAL","SAT","LOY"),
  relationship = list(
    path(from="IMAG", to=c("EXPE","SAT","LOY")),
    path("EXPE", c("QUAL","VAL","SAT")),
    path("QUAL", c("VAL","SAT")),
    path("VAL",c("SAT")),
    path("SAT","LOY"))
)
sat_path1

#'

sat_path <- plspm_paths(relationship = list(
  path(from="IMAG", to=c("EXPE","SAT","LOY")),
  path(from="EXPE", to=c("QUAL","VAL","SAT")),
  path(from="QUAL", to=c("VAL","SAT")),
  path(from="VAL", to="SAT"),
  path(from="SAT", to="LOY")
) )

sat_path <- plspm_paths(relationship = list(
  path(from="IMAG", to=c("EXPE","SAT","LOY")),
  path(from="EXPE", to=c("QUAL","VAL","SAT")),
  path(from="QUAL", to=c("VAL","SAT")),
  path(from="VAL", to="SAT"),
  path(from="SAT", to="LOY")
),
row_names =  c("IMAG", "EXPE", "QUAL", "VAL", "SAT", "LOY")
)

# 결과 출력
print(sat_path)



## End(Not run)
\end{ExampleCode}
\end{Examples}
\inputencoding{utf8}
\HeaderA{plspm\_paths2lav}{plspm\_paths2lav}{plspm.Rul.paths2lav}
%
\begin{Description}
plspm\_paths2lav
\end{Description}
%
\begin{Usage}
\begin{verbatim}
plspm_paths2lav(path_matrix, name = FALSE, hypo = FALSE)
\end{verbatim}
\end{Usage}
%
\begin{Arguments}
\begin{ldescription}
\item[\code{path\_matrix}] matrix paths

\item[\code{name}] default FALSE, TRUE : paths\_model=
\end{ldescription}
\end{Arguments}
%
\begin{Value}
lavaan syntax
\end{Value}
%
\begin{Examples}
\begin{ExampleCode}
## Not run: 
#'
#' # Example path_matrix
path_matrix <- matrix(c(0, 1, 0, 0, 1, 1,
                        0, 0, 1, 1, 1, 0,
                        0, 0, 0, 1, 1, 0,
                        0, 0, 0, 0, 1, 0,
                        0, 0, 0, 0, 0, 1,
                        0, 0, 0, 0, 0, 0),
                      nrow = 6, ncol = 6, byrow = FALSE)

rownames(path_matrix) <- c("IMAG", "EXPE", "QUAL", "VAL", "SAT", "LOY")
colnames(path_matrix) <- c("IMAG", "EXPE", "QUAL", "VAL", "SAT", "LOY")

cat(plspm_paths2lav(path_matrix))

#Add hypothesis markers
cat(plspm_paths2lav(path_matrix, hypo=TRUE))
#'

## End(Not run)
\end{ExampleCode}
\end{Examples}
\inputencoding{utf8}
\HeaderA{plspm\_paths\_coefs\_plot}{Draw a Model of a Measured PLSPM Structure}{plspm.Rul.paths.Rul.coefs.Rul.plot}
%
\begin{Description}
Draw a Model of a Measured PLSPM Structure
\end{Description}
%
\begin{Usage}
\begin{verbatim}
plspm_paths_coefs_plot(
  plsdata,
  layout = "spring",
  boot = TRUE,
  grp = TRUE,
  groups = NULL,
  digits = 3,
  fade = FALSE,
  node = FALSE,
  border.width = 1.5,
  border.color = "gray90",
  edge.label.cex = 1,
  edge.label.position = 0.6,
  asize = 4,
  vsize = 8,
  esize = 10,
  edge_labels_sig = TRUE,
  posCol = "gray20",
  negCol = "red",
  shape = "circle",
  pastel = "pastel",
  trans = TRUE
)
\end{verbatim}
\end{Usage}
%
\begin{Arguments}
\begin{ldescription}
\item[\code{layout}] default spring

\item[\code{boot}] FALSE inner\_model TRUE boot

\item[\code{grp}] groups on T

\item[\code{groups}] NULL

\item[\code{digits}] 3

\item[\code{fade}] FALSE

\item[\code{node}] nodeNames legeng when node =TRUE

\item[\code{border.width}] 1.5

\item[\code{border.color}] gray30

\item[\code{edge.label.cex}] 1

\item[\code{edge.label.position}] 0.5

\item[\code{asize}] 3 arrow size

\item[\code{vsize}] 8 nodesixe

\item[\code{esize}] 10 edgesize

\item[\code{edge\_labels\_sig}] Tedge\_labels\_sig= TRUE add sig star

\item[\code{posCol}] gray20

\item[\code{negCol}] red

\item[\code{shape}] cicle

\item[\code{pastel}] pastel

\item[\code{trans}] true

\item[\code{data}] plspm data
\end{ldescription}
\end{Arguments}
%
\begin{Value}
plot
\end{Value}
%
\begin{Examples}
\begin{ExampleCode}
## Not run: 

# path matrix
IMAG = c(0,0,0,0,0,0)
EXPE = c(1,0,0,0,0,0)
QUAL = c(0,1,0,0,0,0)
VAL = c(0,1,1,0,0,0)
SAT = c(1,1,1,1,0,0)
LOY = c(1,0,0,0,1,0)
sat_path = rbind(IMAG, EXPE, QUAL, VAL, SAT, LOY)

# blocks of outer model
sat_blocks = list(1:5, 6:10, 11:15, 16:19, 20:23, 24:27)

sat_mod = c(rep("A", 6))
# apply plspm
satpls = plspm(satisfaction, sat_path, sat_blocks, modes = sat_mod,
               scaled = FALSE)


 ##jjstat
data(satisfaction)
#model 1
model1 = "
EXPE ~ IMAG
QUAL ~ EXPE
VAL ~  QUAL+ EXPE
SAT ~ IMAG + EXPE + QUAL + VAL
LOY ~ SAT + IMAG
"
# # Generate the matrix
sat_path = plspm_lav2path(model1,
                          fct_order=c("IMAG", "EXPE","QUAL", "VAL", "SAT", "LOY" ))
# sat_path

# Method 2: paths
sat_path = plspm_paths(
  row_names = c("IMAG","EXPE","QUAL","VAL","SAT","LOY"),
  relationship = list(
    path(from="IMAG", to=c("EXPE","SAT","LOY")),
    path("EXPE", c("QUAL","VAL","SAT")),
    path("QUAL", c("VAL","SAT")),
    path("VAL",c("SAT")),
    path("SAT","LOY")
  )
)
# blokcs
sat_blocks1 <- plspm_blocks(
  IMAG = item(paste0("imag",1:5)),
  EXPE = item(paste0("expe", 1:5)),
  QUAL = item( paste0("qual", 1:5)),
  VAL = item(paste0("val", 1:4)),
  SAT = item(paste0("sat", 1:4)),
  LOY = item(paste0("loy", 1:4))
)
sat_blocks1

# vector of modes (reflective indicators):auto
sat_mod = rep("A", 6)
#############

# apply plspm
satpls = plspm_sem(satisfaction,
                   path_matrix = sat_path,
                   blocks = sat_blocks1,
                   scaled = FALSE)

satpls_boot = plspm_sem(satisfaction, sat_path, sat_blocks1, scaled = FALSE,
                        boot.val =TRUE, br=100)

#default
satpls %>% plspm_path_coefs_plot()



lay_p = matrix(c('IMA', NA, NA, NA, NA, NA, NA,
 NA, NA, NA, NA, NA, NA, NA, NA, NA, NA, NA, NA,
 'EXP', NA, NA, NA, NA, NA, NA, NA, NA, NA, NA,
 NA, NA, NA, NA, 'QUA', NA, NA, NA, NA, NA, NA,
 NA, NA, NA, NA, NA, NA, NA, NA, NA, NA, NA, NA,
 NA, NA, NA, NA, NA, NA, 'VAL', NA, NA, NA, NA,
 NA, NA, NA, NA, NA, NA, NA, NA, NA, NA, 'SAT',
 NA, NA, NA, NA, NA, NA, NA, NA, NA, NA, NA, NA,
 NA, NA, NA, NA, 'LOY', NA, NA, NA, NA, NA, NA, NA, NA),
 nrow = 10, ncol = 10, byrow = FALSE)

satpls %>% plspm_path_coefs_plot(layout=lay_p, edge.label.position = 0.6,
                            border.color="gray99", groups=1:6)


# various type
satpls %>% plspm_path_coefs_plot()
satpls %>% plspm_path_coefs_plot(node =TRUE)

satpls %>% plspm_path_coefs_plot(layout=lay_p, edge.label.position = 0.6,
                                 border.color="gray99", groups=1:6)
satpls %>% plspm_path_coefs_plot(layout=lay_p, edge.label.position = 0.6,
                                 border.color="gray99", groups=1:6, node=TRUE)

satpls %>% plspm_path_coefs_plot(Dshape="square", layout=lay_p,
                                 edge.label.position = 0.6)
#'
satpls_boot%>%plspm_path_coefs_plot()
satpls_boot%>%plspm_path_coefs_plot(boot=F)
#'
satpls %>% plspm_path_coefs_plot()
satpls %>% plspm_path_coefs_plot(node =TRUE)
satpls %>% plspm_path_coefs_plot(grp=FALSE)
satpls %>% plspm_path_coefs_plot(node = FALSE,grp=TRUE)
satpls %>% plspm_path_coefs_plot(node = TRUE,grp=FALSE)
satpls %>% plspm_path_coefs_plot(node = FALSE,grp=FALSE)
satpls %>% plspm_path_coefs_plot(node = TRUE,grp=TRUE)

satpls_boot %>%
plspm_path_coefs_plot(layout =  plspm_boot_factor_layout(satpls_boot)%>%
 move_mat_cut(n=3))

#



## End(Not run)
\end{ExampleCode}
\end{Examples}
\inputencoding{utf8}
\HeaderA{plspm\_paths\_sig}{Showing significance using PLSPM regression}{plspm.Rul.paths.Rul.sig}
%
\begin{Description}
Showing significance using PLSPM regression
\end{Description}
%
\begin{Usage}
\begin{verbatim}
plspm_paths_sig(
  data,
  unite = FALSE,
  rhs = "endogenous",
  lhs = "exogenous",
  digits = 3
)
\end{verbatim}
\end{Usage}
%
\begin{Arguments}
\begin{ldescription}
\item[\code{data}] plspm data

\item[\code{unite}] TRUE

\item[\code{rhs}] from

\item[\code{lhs}] to

\item[\code{digits}] round 3
\end{ldescription}
\end{Arguments}
%
\begin{Value}
data result
\end{Value}
%
\begin{Examples}
\begin{ExampleCode}
## Not run: 
library(plspm)
data(satisfaction)

# path matrix
IMAG = c(0,0,0,0,0,0)
EXPE = c(1,0,0,0,0,0)
QUAL = c(0,1,0,0,0,0)
VAL = c(0,1,1,0,0,0)
SAT = c(1,1,1,1,0,0)
LOY = c(1,0,0,0,1,0)
sat_path = rbind(IMAG, EXPE, QUAL, VAL, SAT, LOY)

# blocks of outer model
sat_blocks = list(1:5, 6:10, 11:15, 16:19, 20:23, 24:27)
sat_mod=c(rep("A", ncol(sat_path)))
# apply plspm
satpls = plspm(satisfaction, sat_path, sat_blocks, modes = sat_mod,
               scaled = FALSE)

#default
satpls%>%plspm_paths_sig()
satpls$inner_model %>% plspm_paths_sig()

satpls %>%plspm_paths_sig(unite=FALSE, lhs="row", rhs="col")
satpls$inner_model  %>%innermodel_sig(unite=FALSE, lhs="row", rhs="col")

satpls %>%plspm_paths_sig(unite=FALSE, lhs="row", rhs="col")
satpls$inner_model %>% plspm_paths_sig(unite=FALSE, lhs="row", rhs="col")

#'

## End(Not run)

\end{ExampleCode}
\end{Examples}
\inputencoding{utf8}
\HeaderA{plspm\_paths\_sig\_plot}{Draw a Model of a Measured PLSPM Structure}{plspm.Rul.paths.Rul.sig.Rul.plot}
%
\begin{Description}
Draw a Model of a Measured PLSPM Structure
\end{Description}
%
\begin{Usage}
\begin{verbatim}
plspm_paths_sig_plot(
  plsdata,
  layout = "spring",
  boot = TRUE,
  grp = TRUE,
  groups = NULL,
  digits = 3,
  fade = FALSE,
  node = FALSE,
  border.width = 1.5,
  border.color = "gray90",
  edge.label.cex = 1,
  edge.label.position = 0.6,
  asize = 4,
  vsize = 8,
  esize = 10,
  edge_labels_sig = TRUE,
  posCol = "gray20",
  negCol = "red",
  shape = "circle",
  pastel = "pastel",
  trans = TRUE
)
\end{verbatim}
\end{Usage}
%
\begin{Arguments}
\begin{ldescription}
\item[\code{layout}] default spring

\item[\code{boot}] FALSE inner\_model TRUE boot

\item[\code{grp}] groups on T

\item[\code{groups}] NULL

\item[\code{digits}] 3

\item[\code{fade}] FALSE

\item[\code{node}] nodeNames legeng when node =TRUE

\item[\code{border.width}] 1.5

\item[\code{border.color}] gray30

\item[\code{edge.label.cex}] 1

\item[\code{edge.label.position}] 0.5

\item[\code{asize}] 3 arrow size

\item[\code{vsize}] 8 nodesixe

\item[\code{esize}] 10 edgesize

\item[\code{edge\_labels\_sig}] Tedge\_labels\_sig= TRUE add sig star

\item[\code{posCol}] gray20

\item[\code{negCol}] red

\item[\code{shape}] cicle

\item[\code{pastel}] pastel

\item[\code{trans}] true

\item[\code{data}] plspm data
\end{ldescription}
\end{Arguments}
%
\begin{Value}
plot
\end{Value}
%
\begin{Examples}
\begin{ExampleCode}
## Not run: 

# path matrix
IMAG = c(0,0,0,0,0,0)
EXPE = c(1,0,0,0,0,0)
QUAL = c(0,1,0,0,0,0)
VAL = c(0,1,1,0,0,0)
SAT = c(1,1,1,1,0,0)
LOY = c(1,0,0,0,1,0)
sat_path = rbind(IMAG, EXPE, QUAL, VAL, SAT, LOY)

# blocks of outer model
sat_blocks = list(1:5, 6:10, 11:15, 16:19, 20:23, 24:27)

sat_mod = c(rep("A", 6))
# apply plspm
satpls = plspm(satisfaction, sat_path, sat_blocks, modes = sat_mod,
               scaled = FALSE)


 ##jjstat
data(satisfaction)
#model 1
model1 = "
EXPE ~ IMAG
QUAL ~ EXPE
VAL ~  QUAL+ EXPE
SAT ~ IMAG + EXPE + QUAL + VAL
LOY ~ SAT + IMAG
"
# # Generate the matrix
sat_path = plspm_lav2path(model1,
                          fct_order=c("IMAG", "EXPE","QUAL", "VAL", "SAT", "LOY" ))
# sat_path

# Method 2: paths
sat_path = plspm_paths(
  row_names = c("IMAG","EXPE","QUAL","VAL","SAT","LOY"),
  relationship = list(
    path(from="IMAG", to=c("EXPE","SAT","LOY")),
    path("EXPE", c("QUAL","VAL","SAT")),
    path("QUAL", c("VAL","SAT")),
    path("VAL",c("SAT")),
    path("SAT","LOY")
  )
)
# blokcs
sat_blocks1 <- plspm_blocks(
  IMAG = item(paste0("imag",1:5)),
  EXPE = item(paste0("expe", 1:5)),
  QUAL = item( paste0("qual", 1:5)),
  VAL = item(paste0("val", 1:4)),
  SAT = item(paste0("sat", 1:4)),
  LOY = item(paste0("loy", 1:4))
)
sat_blocks1

# vector of modes (reflective indicators):auto
sat_mod = rep("A", 6)
#############

# apply plspm
satpls = plspm_sem(satisfaction,
                   path_matrix = sat_path,
                   blocks = sat_blocks1,
                   scaled = FALSE)

satpls_boot = plspm_sem(satisfaction, sat_path, sat_blocks1, scaled = FALSE,
                        boot.val =TRUE, br=100)

#default
satpls %>% plspm_paths_sig_plot()



lay_p = matrix(c('IMA', NA, NA, NA, NA, NA, NA,
 NA, NA, NA, NA, NA, NA, NA, NA, NA, NA, NA, NA,
 'EXP', NA, NA, NA, NA, NA, NA, NA, NA, NA, NA,
 NA, NA, NA, NA, 'QUA', NA, NA, NA, NA, NA, NA,
 NA, NA, NA, NA, NA, NA, NA, NA, NA, NA, NA, NA,
 NA, NA, NA, NA, NA, NA, 'VAL', NA, NA, NA, NA,
 NA, NA, NA, NA, NA, NA, NA, NA, NA, NA, 'SAT',
 NA, NA, NA, NA, NA, NA, NA, NA, NA, NA, NA, NA,
 NA, NA, NA, NA, 'LOY', NA, NA, NA, NA, NA, NA, NA, NA),
 nrow = 10, ncol = 10, byrow = FALSE)

satpls %>% plspm_paths_sig_plot(layout=lay_p, edge.label.position = 0.6,
                            border.color="gray99", groups=1:6)


# various type
satpls %>% plspm_paths_sig_plot()
satpls %>% plspm_paths_sig_plot(node =TRUE)

satpls %>% plspm_paths_sig_plot(layout=lay_p, edge.label.position = 0.6,
                                 border.color="gray99", groups=1:6)
satpls %>% plspm_paths_sig_plot(layout=lay_p, edge.label.position = 0.6,
                                 border.color="gray99", groups=1:6, node=TRUE)

satpls %>% plspm_paths_sig_plot(Dshape="square", layout=lay_p,
                                 edge.label.position = 0.6)
#'
satpls_boot%>%plspm_paths_sig_plot()
satpls_boot%>%plspm_paths_sig_plot(boot=F)
#'
satpls %>% plspm_paths_sig_plot()
satpls %>% plspm_paths_sig_plot(node =TRUE)
satpls %>% plspm_paths_sig_plot(grp=FALSE)
satpls %>% plspm_paths_sig_plot(node = FALSE,grp=TRUE)
satpls %>% plspm_paths_sig_plot(node = TRUE,grp=FALSE)
satpls %>% plspm_paths_sig_plot(node = FALSE,grp=FALSE)
satpls %>% plspm_paths_sig_plot(node = TRUE,grp=TRUE)

satpls_boot %>%
plspm_paths_sig_plot(layout =  plspm_boot_factor_layout(satpls_boot)%>%
 move_mat_cut(n=3))

#



## End(Not run)
\end{ExampleCode}
\end{Examples}
\inputencoding{utf8}
\HeaderA{plspm\_path\_coefs\_plot}{Draw a Model of a Measured PLSPM Structure}{plspm.Rul.path.Rul.coefs.Rul.plot}
%
\begin{Description}
Draw a Model of a Measured PLSPM Structure
\end{Description}
%
\begin{Usage}
\begin{verbatim}
plspm_path_coefs_plot(
  plsdata,
  layout = "spring",
  boot = TRUE,
  grp = TRUE,
  groups = NULL,
  digits = 3,
  fade = FALSE,
  node = FALSE,
  border.width = 1.5,
  border.color = "gray90",
  edge.label.cex = 1,
  edge.label.position = 0.6,
  asize = 4,
  vsize = 8,
  esize = 10,
  edge_labels_sig = TRUE,
  posCol = "gray20",
  negCol = "red",
  shape = "circle",
  pastel = "pastel",
  trans = TRUE
)
\end{verbatim}
\end{Usage}
%
\begin{Arguments}
\begin{ldescription}
\item[\code{layout}] default spring

\item[\code{boot}] FALSE inner\_model TRUE boot

\item[\code{grp}] groups on T

\item[\code{groups}] NULL

\item[\code{digits}] 3

\item[\code{fade}] FALSE

\item[\code{node}] nodeNames legeng when node =TRUE

\item[\code{border.width}] 1.5

\item[\code{border.color}] gray30

\item[\code{edge.label.cex}] 1

\item[\code{edge.label.position}] 0.5

\item[\code{asize}] 3 arrow size

\item[\code{vsize}] 8 nodesixe

\item[\code{esize}] 10 edgesize

\item[\code{edge\_labels\_sig}] Tedge\_labels\_sig= TRUE add sig star

\item[\code{posCol}] gray20

\item[\code{negCol}] red

\item[\code{shape}] cicle

\item[\code{pastel}] pastel

\item[\code{trans}] true

\item[\code{data}] plspm data
\end{ldescription}
\end{Arguments}
%
\begin{Value}
plot
\end{Value}
%
\begin{Examples}
\begin{ExampleCode}
## Not run: 

# path matrix
IMAG = c(0,0,0,0,0,0)
EXPE = c(1,0,0,0,0,0)
QUAL = c(0,1,0,0,0,0)
VAL = c(0,1,1,0,0,0)
SAT = c(1,1,1,1,0,0)
LOY = c(1,0,0,0,1,0)
sat_path = rbind(IMAG, EXPE, QUAL, VAL, SAT, LOY)

# blocks of outer model
sat_blocks = list(1:5, 6:10, 11:15, 16:19, 20:23, 24:27)

sat_mod = c(rep("A", 6))
# apply plspm
satpls = plspm(satisfaction, sat_path, sat_blocks, modes = sat_mod,
               scaled = FALSE)


 ##jjstat
data(satisfaction)
#model 1
model1 = "
EXPE ~ IMAG
QUAL ~ EXPE
VAL ~  QUAL+ EXPE
SAT ~ IMAG + EXPE + QUAL + VAL
LOY ~ SAT + IMAG
"
# # Generate the matrix
sat_path = plspm_lav2path(model1,
                          fct_order=c("IMAG", "EXPE","QUAL", "VAL", "SAT", "LOY" ))
# sat_path

# Method 2: paths
sat_path = plspm_paths(
  row_names = c("IMAG","EXPE","QUAL","VAL","SAT","LOY"),
  relationship = list(
    path(from="IMAG", to=c("EXPE","SAT","LOY")),
    path("EXPE", c("QUAL","VAL","SAT")),
    path("QUAL", c("VAL","SAT")),
    path("VAL",c("SAT")),
    path("SAT","LOY")
  )
)
# blokcs
sat_blocks1 <- plspm_blocks(
  IMAG = item(paste0("imag",1:5)),
  EXPE = item(paste0("expe", 1:5)),
  QUAL = item( paste0("qual", 1:5)),
  VAL = item(paste0("val", 1:4)),
  SAT = item(paste0("sat", 1:4)),
  LOY = item(paste0("loy", 1:4))
)
sat_blocks1

# vector of modes (reflective indicators):auto
sat_mod = rep("A", 6)
#############

# apply plspm
satpls = plspm_sem(satisfaction,
                   path_matrix = sat_path,
                   blocks = sat_blocks1,
                   scaled = FALSE)

satpls_boot = plspm_sem(satisfaction, sat_path, sat_blocks1, scaled = FALSE,
                        boot.val =TRUE, br=100)

#default
satpls %>% plspm_path_coefs_plot()



lay_p = matrix(c('IMA', NA, NA, NA, NA, NA, NA,
 NA, NA, NA, NA, NA, NA, NA, NA, NA, NA, NA, NA,
 'EXP', NA, NA, NA, NA, NA, NA, NA, NA, NA, NA,
 NA, NA, NA, NA, 'QUA', NA, NA, NA, NA, NA, NA,
 NA, NA, NA, NA, NA, NA, NA, NA, NA, NA, NA, NA,
 NA, NA, NA, NA, NA, NA, 'VAL', NA, NA, NA, NA,
 NA, NA, NA, NA, NA, NA, NA, NA, NA, NA, 'SAT',
 NA, NA, NA, NA, NA, NA, NA, NA, NA, NA, NA, NA,
 NA, NA, NA, NA, 'LOY', NA, NA, NA, NA, NA, NA, NA, NA),
 nrow = 10, ncol = 10, byrow = FALSE)

satpls %>% plspm_path_coefs_plot(layout=lay_p, edge.label.position = 0.6,
                            border.color="gray99", groups=1:6)


# various type
satpls %>% plspm_path_coefs_plot()
satpls %>% plspm_path_coefs_plot(node =TRUE)

satpls %>% plspm_path_coefs_plot(layout=lay_p, edge.label.position = 0.6,
                                 border.color="gray99", groups=1:6)
satpls %>% plspm_path_coefs_plot(layout=lay_p, edge.label.position = 0.6,
                                 border.color="gray99", groups=1:6, node=TRUE)

satpls %>% plspm_path_coefs_plot(Dshape="square", layout=lay_p,
                                 edge.label.position = 0.6)
#'
satpls_boot%>%plspm_path_coefs_plot()
satpls_boot%>%plspm_path_coefs_plot(boot=F)
#'
satpls %>% plspm_path_coefs_plot()
satpls %>% plspm_path_coefs_plot(node =TRUE)
satpls %>% plspm_path_coefs_plot(grp=FALSE)
satpls %>% plspm_path_coefs_plot(node = FALSE,grp=TRUE)
satpls %>% plspm_path_coefs_plot(node = TRUE,grp=FALSE)
satpls %>% plspm_path_coefs_plot(node = FALSE,grp=FALSE)
satpls %>% plspm_path_coefs_plot(node = TRUE,grp=TRUE)

satpls_boot %>%
plspm_path_coefs_plot(layout =  plspm_boot_factor_layout(satpls_boot)%>%
 move_mat_cut(n=3))

#



## End(Not run)
\end{ExampleCode}
\end{Examples}
\inputencoding{utf8}
\HeaderA{plspm\_permutation\_test}{잠재변수 평균값과 분산의 동일성을 평가하는 Permutation Test 함수}{plspm.Rul.permutation.Rul.test}
%
\begin{Description}
잠재변수 평균값과 분산의 동일성을 평가하는 Permutation Test 함수
\end{Description}
%
\begin{Usage}
\begin{verbatim}
plspm_permutation_test(Data, path_matrix, blocks, grp, n_perm = 100)
\end{verbatim}
\end{Usage}
%
\begin{Arguments}
\begin{ldescription}
\item[\code{Data}] Data

\item[\code{path\_matrix}] path\_matrix

\item[\code{blocks}] blocks

\item[\code{grp}] grp = "gender"

\item[\code{n\_perm}] default 100
\end{ldescription}
\end{Arguments}
%
\begin{Value}
data
\end{Value}
%
\begin{Examples}
\begin{ExampleCode}
## Not run: 
#'
#'
jutpath1 = plspm_paths(
  row_names = c("SelfEff","CarMot","CarAtt","CarAct"),
  relationship = list(
    path(from="SelfEff", to=c("CarMot","CarAtt","CarAct")),
    path("CarMot", c("CarAct")),
    path("CarAtt", c("CarAct"))
  )  )
jutpath1


jut_blocks <- plspm_blocks(
  SelfEff = list("sef01", "sef02", "sef03", "sef04", "sef05"),
  CarMot = list("mot01","mot02","mot03"),
  #   CarAtt = item( "att01", "att02","att03","att04","att05" ),
  CarAtt = list("att03","att04","att05" ),
  CarAct = list("act01", "act02","act03")
)
jut_blocks


result_df <- plspm_permutation_test(Data = jut7c, path_matrix = jutpath1, blocks = jut_blocks, grp = "P01")


result_df
result_df %>%Round(17)%>%dall()



## End(Not run)

\end{ExampleCode}
\end{Examples}
\inputencoding{utf8}
\HeaderA{plspm\_permutation\_test\_multi}{plspm\_permutation\_test\_multi: Perform pairwise permutation tests for multiple groups}{plspm.Rul.permutation.Rul.test.Rul.multi}
%
\begin{Description}
plspm\_permutation\_test\_multi: Perform pairwise permutation tests for multiple groups
\end{Description}
%
\begin{Usage}
\begin{verbatim}
plspm_permutation_test_multi(Data, path_matrix, blocks, grp, n_perm = 100)
\end{verbatim}
\end{Usage}
%
\begin{Arguments}
\begin{ldescription}
\item[\code{Data}] Data frame containing the data.

\item[\code{path\_matrix}] Path matrix for the PLS-PM model.

\item[\code{blocks}] List of blocks for the PLS-PM model.

\item[\code{grp}] Grouping variable name (as a string).

\item[\code{n\_perm}] Number of permutations to perform (default: 100).
\end{ldescription}
\end{Arguments}
%
\begin{Value}
A list of data frames, where each data frame contains results for a pairwise group comparison.
\end{Value}
%
\begin{Examples}
\begin{ExampleCode}
## Not run: 
# Example usage
data <- iris  # Example dataset
path_matrix <- matrix(c(0, 1, 1, 0), 2, 2)  # Example path matrix
blocks <- list(1:2, 3:4)  # Example blocks
result <- plspm_permutation_test_multi(data, path_matrix, blocks, "Species", n_perm = 100)
print(result)

## End(Not run)
\end{ExampleCode}
\end{Examples}
\inputencoding{utf8}
\HeaderA{plspm\_prod}{Create a multiply variable to create a moderator variable}{plspm.Rul.prod}
%
\begin{Description}
Create a multiply variable to create a moderator variable
\end{Description}
%
\begin{Usage}
\begin{verbatim}
plspm_prod(data, blocks, ..., intname = "int", n_name = 2, id = FALSE)
\end{verbatim}
\end{Usage}
%
\begin{Arguments}
\begin{ldescription}
\item[\code{data}] data

\item[\code{blocks}] latent and measurement blocks

\item[\code{...}] interation terms

\item[\code{id}] discriminat variable
\end{ldescription}
\end{Arguments}
%
\begin{Value}
data
\end{Value}
%
\begin{Examples}
\begin{ExampleCode}

## Not run: 

# Using example data
data(satisfaction)
sat_blocks1 <- list(
  IMAG = paste0("imag", 1:5),
  EXPE = paste0("expe", 1:5),
  QUAL = paste0("qual", 1:5),
  VAL = paste0("val", 1:4),
  SAT = paste0("sat", 1:4),
  LOY = paste0("loy", 1:4)
)
# Partially select to create an interaction term
satisfaction%>% plspm_prod(sat_blocks1,
                           intIS = interlist(c("IMAG", 1, 2, 3), c("SAT", 1, 2, 3))) %>% str()

# Partially select to create an interaction term : repeated
satisfaction%>% plspm_prod(sat_blocks1,
                           intIS = interlist(c("IMAG", 1, 2, 3), c("SAT", 1, 2, 3))
)%>% plspm_prod(sat_blocks1,
                intIS = interlist(c("IMAG", 1, 2, 3), c("EXPE", 1, 2))
) %>%str()


# Create an interaction variable
plspm_prod(satisfaction, sat_blocks1, intIS = interlist("IMAG", "SAT"))    %>%str()
# Create a variable that interacts with multiple variables
plspm_prod(satisfaction, sat_blocks1, intIS = interlist("IMAG", "SAT"),
           intIE = interlist("IMAG", "EXPE"))    %>%str()

satisfaction%>% plspm_prod(sat_blocks1,
                           interlist(c("IMAG", 1, 2, 3), c("SAT", 1, 2, 3))) %>%str()

satisfaction%>% plspm_prod(sat_blocks1,
                           interlist("IMAG", 1, 2, 3, "SAT", 1, 2, 3) ) %>%str()

# plspm_prod(satisfaction, sat_blocks1, list("IMAG", "SAT"),
#                                        list("IMAG", "EXPE")) %>%str()


#  Use vector
plspm_prod(satisfaction, sat_blocks1,
           intIS = list(c("IMAG", 1, 2 ), c("SAT", 3,4)),
           intIE = list(c("IMAG", 1, 2, 3), c("EXPE" ,4,5))
) %>%str()

#interlist use use list
plspm_prod(satisfaction,
           sat_blocks1,
           intIS = interlist(list("IMAG", 1, 2 ,"SAT", 3,4)),
           intIE = interlist(list("IMAG", 1, 2, 3, "EXPE" ,4,5)) ) %>%str()


#interlist use now use list
plspm_prod(satisfaction,
           sat_blocks1,
           intIS = interlist("IMAG", 1, 2 ,"SAT", 3,4),
           intIE = interlist("IMAG", 1, 2, 3, "EXPE" ,4,5)) %>%str()

## End(Not run)



\end{ExampleCode}
\end{Examples}
\inputencoding{utf8}
\HeaderA{plspm\_rotate\_resid}{plspm\_rotate\_resid}{plspm.Rul.rotate.Rul.resid}
%
\begin{Description}
plspm\_rotate\_resid
\end{Description}
%
\begin{Usage}
\begin{verbatim}
plspm_rotate_resid(plotdata, resid_list = NULL, plot = TRUE)
\end{verbatim}
\end{Usage}
%
\begin{Arguments}
\begin{ldescription}
\item[\code{plotdata}] plotdata

\item[\code{resid\_list}] resid list

\item[\code{plot}] TRUE
\end{ldescription}
\end{Arguments}
%
\begin{Value}
plot data
\end{Value}
%
\begin{Examples}
\begin{ExampleCode}

## Not run: 
#'

pathmodel = "
자기효능감 =~ C_S1 +  C_S2 + C_S3 + C_S4 + C_S5
진로태도 =~ B02 + B04 + B10 + B14
진로동기 =~ A_M1 + A_M2 + A_M3
진로준비 =~ D_P1 + D_P2 + D_P3

진로동기 ~ H1*자기효능감
진로태도 ~  H2*자기효능감
진로준비 ~ H5 *진로태도
진로준비 ~  H4 * 진로동기
진로준비 ~  H3*자기효능감
"

plspm_diagram_model(pathmodel, layout="tree2",
                    rotation = 2,residuals = T, residScale = 6,
                    edge.label.position = 0.6)%>%
  plspm_rotate_resid(c(진로동기= 160, 진로태도= 160))

#
plspm_diagram_model(pathmodel, layout="tree2",
                    rotation = 2,residuals = T, residScale = 6,
                    edge.label.position = 0.6)%>%
  plspm_curve(c("진로준비~~ 진로동기" = -2))%>%
  plspm_rotate_resid(c(진로동기= 160, 진로태도= 30))

## End(Not run)


\end{ExampleCode}
\end{Examples}
\inputencoding{utf8}
\HeaderA{plspm\_sem}{PLS-PM: Partial Least Squares Path Modeling}{plspm.Rul.sem}
%
\begin{Description}
PLS-PM: Partial Least Squares Path Modeling
\end{Description}
%
\begin{Usage}
\begin{verbatim}
plspm_sem(
  Data,
  path_matrix,
  blocks,
  interactionTerm = NULL,
  indicator_prod = NULL,
  inter_modes = NULL,
  modes = rep("A", ncol(path_matrix)),
  scaling = NULL,
  scheme = "centroid",
  scaled = TRUE,
  tol = 1e-06,
  maxiter = 100,
  plscomp = NULL,
  boot.val = TRUE,
  br = 500,
  seed = NULL,
  dataset = TRUE,
  summary = TRUE
)
\end{verbatim}
\end{Usage}
%
\begin{Arguments}
\begin{ldescription}
\item[\code{Data}] matrix or data frame containing the manifest variables.

\item[\code{path\_matrix}] A square (lower triangular) boolean matrix
representing the inner model (i.e. the path relationships between latent variables).

\item[\code{blocks}] list of vectors with column indices
or column names from Data indicating the sets of
manifest variables forming each block (i.e.
which manifest variables correspond to each block).

\item[\code{interactionTerm}] interaction effect list

\item[\code{indicator\_prod}] generate indicator product data

\item[\code{inter\_modes}] interaction modes

\item[\code{modes}] optional argument for runing the non-metric approach;
it is a list of string vectors indicating the type of measurement
scale for each manifest variable specified in blocks. scaling must be
specified when working with non-metric variables. Possible values:
"num" (linear transformation, suitable for numerical variables),
"raw" (no transformation), "nom" (non-monotonic transformation,
suitable for nominal variables), and "ord" (monotonic transformation,
suitable for ordinal variables).

\item[\code{scaling}] character vector indicating the type of measurement f
or each block. Possible values are: "A", "B", "newA", "PLScore",
"PLScow". The length of modes must be equal to the length of blocks.

\item[\code{scheme}] string indicating the type of inner weighting scheme.
Possible values are "centroid", "factorial", or "path".

\item[\code{scaled}] standardized. Only used when scaling = NULL.
When (TRUE, data is scaled to standardized values (mean=0 and variance=1).
The variance is calculated dividing by N instead of N-1).

\item[\code{tol}] decimal value indicating the tolerance criterion
for the iterations (tol=0.000001). Can be specified between 0 and 0.001.

\item[\code{maxiter}] integer indicating the maximum number of iterations
(maxiter=100 by default). The minimum value of maxiter is 100.

\item[\code{plscomp}] optional vector indicating the number of PLS components
(for each block) to be used when handling non-metric data
(only used if scaling is provided)

\item[\code{boot.val}] whether bootstrap validation should be performed.
(FALSE by default).

\item[\code{br}] number bootstrap resamples. Used only when boot.val=TRUE. When boot.val=TRUE, the default number of re-samples is 100.

\item[\code{dataset}] whether the data matrix used in the computations should be retrieved (TRUE by default).

\item[\code{summary}] TRUE
\end{ldescription}
\end{Arguments}
%
\begin{Value}
plspm result
\end{Value}
%
\begin{Examples}
\begin{ExampleCode}


## Not run: 
#'
## Not run:
## typical example of PLS-PM in customer satisfaction analysis
## model with six LVs and reflective indicators
# load dataset satisfaction
data(satisfaction)
# path matrix
# IMAG = c(0,0,0,0,0,0)
# EXPE = c(1,0,0,0,0,0)
# QUAL = c(0,1,0,0,0,0)
# VAL = c(0,1,1,0,0,0)
# SAT = c(1,1,1,1,0,0)
# LOY = c(1,0,0,0,1,0)
# sat_path = rbind(IMAG, EXPE, QUAL, VAL, SAT, LOY)

sat_path = plspm_paths(
  row_names = c("IMAG","EXPE","QUAL","VAL","SAT","LOY"),

  relationship = list(
    path(from="IMAG", to=c("EXPE","SAT","LOY")),
    path("EXPE", c("QUAL","VAL","SAT")),
    path("QUAL", c("VAL","SAT")),
    path("VAL",c("SAT")),
    path("SAT","LOY")
  )
)

# plot diagram of path matrix
innerplot(sat_path)
# blocks of outer model
# sat_blocks = list(1:5, 6:10, 11:15, 16:19, 20:23, 24:27)
#' #####################
library(jjstat)
data(satisfaction)
sat_path = plspm_paths(
  row_names = c("IMAG","EXPE","QUAL","VAL","SAT","LOY"),
  relationship = list(
    path(from="IMAG", to=c("EXPE","SAT","LOY")),
    path("EXPE", c("QUAL","VAL","SAT")),
    path("QUAL", c("VAL","SAT")),
    path("VAL",c("SAT")),
    path("SAT","LOY"))
)
sat_path

# blokcs
sat_blocks1 <- plspm_blocks(
  IMAG = item(paste0("imag",1:5)),
  EXPE = item(paste0("expe", 1:5)),
  QUAL = item( paste0("qual", 1:5)),
  VAL = item(paste0("val", 1:4)),
  SAT = item(paste0("sat", 1:4)),
  LOY = item(paste0("loy", 1:4))
)

sat_blocks1

satpls_boot = plspm_sem(satisfaction, sat_path, sat_blocks1, scaled = FALSE,
                        boot.val =TRUE, br=100)


#'
# interction effect two stage method
resinter_sem = plspm_sem(Data=satisfaction,
                         sat_path2, blocks = sat_blocks1, modes= sat_mod,
                         interactionTerm = list(
                           list(from = "IMAG*EXPE", to = "LOY"),
                           list(from = "IMAG*QUAL", to = "LOY")))
resinter_sem
summary(resinter_sem)
plspm_path_coefs_plot(resinter_sem)
#'
#'



#'
#'
## Two stage approach
pls_ts2 = plspm_sem(satisfaction, sat_path2, sat_blocks1,
                    interactionTerm = list(list(from = "IMAG*EXPE", to = "LOY")) )
pls_ts2 %>%summary()
pls_ts2 %>% plspm_path_coefs_plot()

## indicator arpproach
pls_int2 = plspm_sem(satisfaction, sat_path2, sat_blocks1,
                     interactionTerm = list(list(from = "IMAG*EXPE", to = "LOY")) ,
                     indicator_prod = list(intIE = interlist("IMAG",1,2, "EXPE",3,4)))
pls_int2 %>%summary()
pls_int2 %>% plspm_path_coefs_plot()
pls_int2 %>% plspm_semPaths2()


## End(Not run)

\end{ExampleCode}
\end{Examples}
\inputencoding{utf8}
\HeaderA{plspm\_semPaths}{Drawing Structural Equation Models Using the lavaan Syntax (park JH PhD)}{plspm.Rul.semPaths}
%
\begin{Description}
Drawing Structural Equation Models Using the lavaan Syntax (park JH PhD)
\end{Description}
%
\begin{Usage}
\begin{verbatim}
plspm_semPaths(
  plsres_boot,
  residuals = FALSE,
  whatLabels = "model",
  rotation = 2,
  residScale = 12,
  layout = "tree2",
  edge.label.cex = 0.7,
  edge.label.position = 0.55,
  edge.color = "black",
  sizeMan = 6,
  sizeMan2 = 3,
  sizeLat = 7,
  sizeLat2 = 4,
  style = "lisrel",
  sig = FALSE,
  exoVar = FALSE,
  exoCov = FALSE,
  curve = 1.5,
  asize = 2,
  mar = c(2, 4, 2, 4),
  nDigits = 3,
  shapeLat = "circle",
  shapeMan = "rectangle",
  shapeInt = "triangle",
  sample.nobs = 100,
  border.width = 2,
  edge.width = 1.5,
  Groups = FALSE,
  group_match = "lat",
  growth = FALSE,
  structural = FALSE,
  edgeLabels = NULL,
  nodeLabels = NULL,
  fixedStyle = 1,
  color = list(lat = "white", man = "white"),
  label.cex = 1,
  type = "plot"
)
\end{verbatim}
\end{Usage}
%
\begin{Arguments}
\begin{ldescription}
\item[\code{plsres\_boot}] plsres\_boot to model syntax

\item[\code{residuals}] residauls TRUE

\item[\code{whatLabels}] "model

\item[\code{rotation}] rotation=2

\item[\code{residScale}] residScale=12

\item[\code{layout}] layout="tree"

\item[\code{edge.label.cex}] 1

\item[\code{edge.label.position}] 0.5

\item[\code{edge.color}] edge.color= black

\item[\code{sizeMan}] 8

\item[\code{sizeMan2}] 4

\item[\code{sizeLat}] 10

\item[\code{sizeLat2}] 6

\item[\code{style}] lisrel

\item[\code{sig}] FALSE

\item[\code{exoVar}] FALSE

\item[\code{exoCov}] TRUE

\item[\code{curve}] 1.5

\item[\code{asize}] 1.5

\item[\code{mar}] c(2,8,3,10)

\item[\code{nDigits}] 2

\item[\code{shapeLat}] "circle, ellipse

\item[\code{shapeMan}] retangle

\item[\code{shapeInt}] triangle ,

\item[\code{sample.nobs}] sample.nobs = 100

\item[\code{border.width}] border.width=1

\item[\code{edge.width}] edge.width=1.5

\item[\code{Groups}] FALSE

\item[\code{group\_match}] "FALSE"lat\#'

\item[\code{growth}] growth= FALSE, TRUE  model.type ="grouwth"

\item[\code{structural}] structural = FALSE

\item[\code{edgeLabels}] edgeLabels=NULL

\item[\code{nodeLabels}] nodeLabels=NULL

\item[\code{fixedStyle}] fixedStyle=1 default rbase linetype

\item[\code{color}] list(lat="skyblue", man="yellow")

\item[\code{label.cex}] label.cex

\item[\code{type}] type ="plot, and res is lavaan simdata result
\end{ldescription}
\end{Arguments}
%
\begin{Value}
plot jutpls\_boot
\end{Value}
\inputencoding{utf8}
\HeaderA{plspm\_semPaths2}{Drawing Structural Equation Models Using the lavaan Syntax (park JH PhD)}{plspm.Rul.semPaths2}
%
\begin{Description}
Drawing Structural Equation Models Using the lavaan Syntax (park JH PhD)
\end{Description}
%
\begin{Usage}
\begin{verbatim}
plspm_semPaths2(
  plsres_boot,
  residuals = FALSE,
  whatLabels = "model",
  rotation = 2,
  residScale = 12,
  layout = "tree2",
  edge.label.cex = 0.7,
  edge.label.position = 0.55,
  edge.color = "black",
  sizeMan = 6,
  sizeMan2 = 3,
  sizeLat = 7,
  sizeLat2 = 4,
  style = "lisrel",
  sig = FALSE,
  exoVar = FALSE,
  exoCov = FALSE,
  curve = 1.5,
  asize = 2,
  mar = c(2, 4, 2, 4),
  nDigits = 3,
  shapeLat = "circle",
  shapeMan = "rectangle",
  shapeInt = "triangle",
  sample.nobs = 100,
  border.width = 2,
  edge.width = 1.5,
  Groups = FALSE,
  group_match = "lat",
  growth = FALSE,
  structural = FALSE,
  edgeLabels = NULL,
  nodeLabels = NULL,
  fixedStyle = 1,
  color = list(lat = "gray10", man = "gray90"),
  label.cex = 1,
  type = "plot"
)
\end{verbatim}
\end{Usage}
%
\begin{Arguments}
\begin{ldescription}
\item[\code{plsres\_boot}] plsres\_boot to model syntax

\item[\code{residuals}] residauls TRUE

\item[\code{whatLabels}] "model

\item[\code{rotation}] rotation=2

\item[\code{residScale}] residScale=12

\item[\code{layout}] layout="tree"

\item[\code{edge.label.cex}] 1

\item[\code{edge.label.position}] 0.5

\item[\code{edge.color}] edge.color= black

\item[\code{sizeMan}] 8

\item[\code{sizeMan2}] 4

\item[\code{sizeLat}] 10

\item[\code{sizeLat2}] 6

\item[\code{style}] lisrel

\item[\code{sig}] FALSE

\item[\code{exoVar}] FALSE

\item[\code{exoCov}] TRUE

\item[\code{curve}] 1.5

\item[\code{asize}] 1.5

\item[\code{mar}] c(2,8,3,10)

\item[\code{nDigits}] 2

\item[\code{shapeLat}] "circle, ellipse

\item[\code{shapeMan}] retangle

\item[\code{shapeInt}] triangle ,

\item[\code{sample.nobs}] sample.nobs = 100

\item[\code{border.width}] border.width=1

\item[\code{edge.width}] edge.width=1.5

\item[\code{Groups}] FALSE

\item[\code{group\_match}] "FALSE"lat\#'

\item[\code{growth}] growth= FALSE, TRUE  model.type ="grouwth"

\item[\code{structural}] structural = FALSE

\item[\code{edgeLabels}] edgeLabels=NULL

\item[\code{nodeLabels}] nodeLabels=NULL

\item[\code{fixedStyle}] fixedStyle=1 default rbase linetype

\item[\code{color}] list(lat="skyblue", man="yellow")

\item[\code{label.cex}] label.cex

\item[\code{type}] type ="plot, and res is lavaan simdata result
\end{ldescription}
\end{Arguments}
%
\begin{Value}
plot jutpls\_boot
\end{Value}
\inputencoding{utf8}
\HeaderA{plspm\_semPaths\_lav}{Generate lavaan syntax from a model or plspm result}{plspm.Rul.semPaths.Rul.lav}
%
\begin{Description}
Generate lavaan syntax from a model or plspm result
\end{Description}
%
\begin{Usage}
\begin{verbatim}
plspm_semPaths_lav(plsres_boot, type = "lav", paths = NULL, blocks = NULL)
\end{verbatim}
\end{Usage}
%
\begin{Arguments}
\begin{ldescription}
\item[\code{plsres\_boot}] plspm result

\item[\code{type}] lav, ,model, edge, edgeLabels

\item[\code{paths}] paths matrix

\item[\code{blocks}] blocks tree
\end{ldescription}
\end{Arguments}
%
\begin{Value}
lavaan syntax
\end{Value}
%
\begin{Examples}
\begin{ExampleCode}

## Not run: 

library(jjstat)
data(satisfaction)
sat_path = plspm_paths(
  row_names = c("IMAG","EXPE","QUAL","VAL","SAT","LOY"),
  relationship = list(
    path(from="IMAG", to=c("EXPE","SAT","LOY")),
    path("EXPE", c("QUAL","VAL","SAT")),
    path("QUAL", c("VAL","SAT")),
    path("VAL",c("SAT")),
    path("SAT","LOY"))
)
sat_path

# blokcs
sat_blocks1 <- plspm_blocks(
  IMAG = item(paste0("imag",1:5)),
  EXPE = item(paste0("expe", 1:5)),
  QUAL = item( paste0("qual", 1:5)),
  VAL = item(paste0("val", 1:4)),
  SAT = item(paste0("sat", 1:4)),
  LOY = item(paste0("loy", 1:4))
)

sat_blocks1


###lav
satpls_boot%>%
  plspm_semPaths_lav(paths=sat_path, blocks=sat_blocks1)%>%cat()

## model
satpls_boot%>%plspm_semPaths_lav()

## edge
satpls_boot%>%plspm_semPaths_lav("edge")

## draw model
satpls_boot%>%
  plspm_semPaths_lav(paths=sat_path, blocks=sat_blocks1)%>%
  diagram_model(  edgeLabels= satpls_boot%>%plspm_semPaths_lav("edge"),
                  mar=c(2,2,2,2),
                  sizeLat=4, sizeMan=3, sizeMan2= 2, edge.label.cex=0.5, rotation=1)

###
satpls_boot%>%plspm_semPaths_lav()%>%
  diagram_model(edgeLabels=satpls_boot%>%plspm_semPaths_lav("edge"),
                mar=c(2,2,2,2),edge.label.cex=0.5, edge.label.position= 0.7,
                sizeLat=4, sizeMan=3, sizeMan2= 2, rotation=1)


## End(Not run)


\end{ExampleCode}
\end{Examples}
\inputencoding{utf8}
\HeaderA{plspm\_semptools}{plspm\_semptools}{plspm.Rul.semptools}
%
\begin{Description}
plspm\_semptools
\end{Description}
%
\begin{Usage}
\begin{verbatim}
plspm_semptools(
  semPaths_plot,
  pls_boot,
  position = NULL,
  point = NULL,
  set_push = 2.5,
  set_spread = 2.5,
  nrow = 3,
  ncol = 3,
  loading_position = 0.5
)
\end{verbatim}
\end{Usage}
%
\begin{Arguments}
\begin{ldescription}
\item[\code{semPaths\_plot}] plot data

\item[\code{pls\_boot}] boot data

\item[\code{position}] layout  position

\item[\code{point}] point position

\item[\code{set\_push}] each push

\item[\code{set\_spread}] each spread

\item[\code{nrow}] default 3

\item[\code{ncol}] default 3

\item[\code{loading\_position}] 0.5
\end{ldescription}
\end{Arguments}
%
\begin{Value}
plot
\end{Value}
%
\begin{Examples}
\begin{ExampleCode}

## Not run: 
#'

#########################################################
pathmodel = "
자기효능감 =~ C_S1 +  C_S2 + C_S3 + C_S4 + C_S5
진로태도 =~ B02 + B04 + B10 + B14
진로동기 =~ A_M1 + A_M2 + A_M3
진로준비 =~ D_P1 + D_P2 + D_P3

진로동기 ~ H1*자기효능감
진로태도 ~  H2*자기효능감
진로준비 ~ H5*진로태도 + H4*진로동기 + H3*자기효능감
"
pathmodel

plspm_semPaths2(jutpls_boot,   curve=1, layout="tree2",rotation=1,
                sizeMan= 10, color=list(lat="gray10",man="gray98"),
                edge.label.cex=0.7, mar=c(5,5,5,5),
                edge.label.position=0.55,
                sizeLat= 8)%>%
  semptools::set_sem_layout(
    indicator_order =  plspm_indicator_order(pathmodel),
    indicator_factor =  plspm_indicator_factor(pathmodel),
    factor_layout=  matrix(
      c(  NA, '진로태도', NA,
          '자기효능감', NA, '진로준비',
          NA, '진로동기', NA),
      nrow = 3, ncol = 3, byrow = TRUE),
    factor_point_to= matrix(
      c(  NA, 'up', NA,
          'left', NA, 'right',
          NA, 'down', NA),
      nrow = 3, ncol = 3, byrow = TRUE),
    indicator_push <- c(진로태도 = 2, 진로동기 = 2, 자기효능감=1.5, 진로준비 =1.5),
    indicator_spread <- c(자기효능감 = 3, 진로태도 = 2, 진로동기 =2, 진로준비 = 2)
  )%>%  semptools::set_curve(c("진로준비~~ 진로태도" = 0.5))%>%
  semptools::set_curve(c("진로준비~~ 진로동기" = 0))%>%
  plspm_change_nodeLabels("B02","진로결정",
                          "B04", "진로독립",
                          "B10", "진로계획",
                          "B14", "진로성향",
                          "A_M1","진로정체",
                          "A_M2","진로통찰",
                          "A_M3","진로탄력",
                          "C_S1","자기평가",
                          "C_S2","직업정보",
                          "C_S3","계획수립",
                          "C_S4","목표설정",
                          "C_S5","해결방식",
                          "D_P1","진로준비",
                          "D_P2","도구준비",
                          "D_P3","목표달성"
  )
# second ##################################
plspm_indicator_order(pathmodel)
plspm_indicator_factor(pathmodel)


#옵션함수
plspm_boot_ind_order(jutpls_boot)
plspm_boot_ind_factor(jutpls_boot)

plspm_boot_factor_layout(jutpls_boot)
plspm_boot_factor_point_to(jutpls_boot, ncol=3)

plspm_ind_push(jutpls_boot)
plspm_ind_spread(jutpls_boot, set= c(자기효능감=4,진로태도=3, 진로동기=3))

##plot
plspm_semPaths2(jutpls_boot,   curve=1, layout="tree2",rotation=1,
                sizeMan= 10, color=list(lat="gray10",man="gray98"),
                edge.label.cex=0.7, mar=c(5,5,5,5),
                edge.label.position=0.55, edge.color= "tomato4",
                sizeLat= 8)%>%
  semptools::set_sem_layout(
    indicator_order = plspm_boot_ind_order(jutpls_boot),
    indicator_factor = plspm_boot_ind_factor(jutpls_boot),
    factor_layout= plspm_boot_factor_layout(jutpls_boot, ncol=3),
    factor_point_to= plspm_boot_factor_point_to(jutpls_boot, ncol=3),
    indicator_push <- plspm_ind_push(jutpls_boot),
    indicator_spread <-  plspm_ind_spread(jutpls_boot, set= c(자기효능감=4,진로태도=3, 진로동기=3)))%>%
  plspm_curve(c("진로준비~ 진로동기" =0))%>%
  plspm_change_nodeLabels("B02","진로결정",
                          "B04", "진로독립",
                          "B10", "진로계획",
                          "B14", "진로성향",
                          "A_M1","진로정체",
                          "A_M2","진로통찰",
                          "A_M3","진로탄력",
                          "C_S1","자기평가",
                          "C_S2","직업정보",
                          "C_S3","계획수립",
                          "C_S4","목표설정",
                          "C_S5","해결방식",
                          "D_P1","진로준비",
                          "D_P2","도구준비",
                          "D_P3","목표달성"
  )

#'
#' #we want #########################################
plspm_semPaths2(jutpls_boot,   curve=1, layout="tree2",rotation=1,
                sizeMan= 10, color=list(lat="gray10",man="gray98"),
                edge.label.cex=0.7, mar=c(5,5,5,5),
                edge.label.position=0.55, edge.color= "tomato4",
                sizeLat= 8) %>%
  plspm_semptools(jutpls_boot)%>%
  plspm_curve(c("진로준비~ 진로동기" =0))%>%
  plspm_change_nodeLabels("B02","진로결정",
                          "B04", "진로독립",
                          "B10", "진로계획",
                          "B14", "진로성향",
                          "A_M1","진로정체",
                          "A_M2","진로통찰",
                          "A_M3","진로탄력",
                          "C_S1","자기평가",
                          "C_S2","직업정보",
                          "C_S3","계획수립",
                          "C_S4","목표설정",
                          "C_S5","해결방식",
                          "D_P1","진로준비",
                          "D_P2","도구준비",
                          "D_P3","목표달성"
  )
#'



## End(Not run)

\end{ExampleCode}
\end{Examples}
\inputencoding{utf8}
\HeaderA{pls\_modelfit\_apa}{pls\_modelfit\_apa}{pls.Rul.modelfit.Rul.apa}
%
\begin{Description}
pls\_modelfit\_apa
\end{Description}
%
\begin{Usage}
\begin{verbatim}
pls_modelfit_apa(model_fit_data)
\end{verbatim}
\end{Usage}
%
\begin{Arguments}
\begin{ldescription}
\item[\code{model\_fit\_data}] modelfit
\end{ldescription}
\end{Arguments}
%
\begin{Value}
data
\end{Value}
%
\begin{Examples}
\begin{ExampleCode}

## Not run: 
#'
rest= plspm_modelfit(edt_pls, edblock, bootstrap = 1000, ci = 0.95, trans=FALSE)
rest
rest%>%t()%>%rowdata2col(1)%>%data.frame()

pls_modelfit_apa(rest)


## End(Not run)
\end{ExampleCode}
\end{Examples}
\inputencoding{utf8}
\HeaderA{poolcor\_comment}{metasem primary analysis analysis results by JH, Park}{poolcor.Rul.comment}
%
\begin{Description}
metasem primary analysis analysis results by JH, Park
\end{Description}
%
\begin{Usage}
\begin{verbatim}
poolcor_comment(res)
\end{verbatim}
\end{Usage}
%
\begin{Arguments}
\begin{ldescription}
\item[\code{res}] metaSEM result
\end{ldescription}
\end{Arguments}
%
\begin{Examples}
\begin{ExampleCode}
## Not run: 
## making correlation matrix ##
cormat1 = metacorMat(rd2, startcol= 3, var=c("po","ne","se","jp"))
 ##variable : "po_ne" "po_se" "po_jp" "ne_se" "ne_jp" "se_jp"
 varnames <- c("po","ne","se","jp")
rd2<- data.frame(
  stringsAsFactors = FALSE,
  study = c(1L,
            2L,3L,4L,5L,6L,7L,8L,9L,10L,11L,12L,
            13L,14L,15L,16L,17L,18L,19L,20L,
            21L,22L,23L,24L,25L,26L,27L,28L,29L,
            30L,31L,32L,33L,34L,35L,36L),
  n = c(671L,
        92L,73L,321L,214L,947L,321L,231L,
        531L,567L,213L,928L,269L,593L,78L,983L,
        1504L,247L,79L,875L,68L,201L,891L,
        197L,431L,98L,179L,69L,301L,679L,273L,
        254L,316L,781L,273L,268L),
  po_ne = c(NA,
            NA,-0.33,NA,NA,NA,NA,-0.32,NA,NA,NA,
            NA,NA,NA,-0.21,-0.49,NA,-0.14,-0.29,
            -0.27,NA,-0.37,NA,NA,-0.35,-0.17,NA,
            -0.11,NA,NA,-0.41,NA,NA,NA,NA,-0.39),
  po_se = c(NA,
            0.37,0.29,0.16,0.41,NA,NA,0.21,0.59,
            0.53,0.31,0.38,NA,NA,0.19,NA,NA,NA,
            NA,0.21,0.07,0.57,0.45,NA,0.18,NA,
            0.45,NA,NA,0.09,NA,NA,0.37,0.19,0.23,
            NA),
  po_jp = c(0.16,
            0.39,0.09,0.07,NA,0.04,NA,NA,0.33,
            0.11,NA,0.04,NA,NA,0.21,0.23,0.09,
            0.11,0.31,0.11,-0.07,0.14,NA,0.19,0.13,
            0.25,0.35,0.29,0.35,NA,0.27,0.19,NA,
            NA,0.13,0.19),
  ne_se = c(NA,
            -0.31,NA,NA,NA,NA,NA,-0.19,NA,NA,NA,
            NA,-0.49,NA,-0.45,NA,NA,NA,NA,-0.38,
            NA,-0.39,NA,NA,-0.21,NA,NA,NA,NA,
            -0.23,NA,NA,NA,NA,-0.44,NA),
  ne_jp = c(NA,
            -0.41,NA,NA,NA,-0.18,-0.11,NA,NA,NA,
            NA,NA,-0.21,-0.15,-0.36,-0.29,NA,
            -0.27,-0.29,-0.23,NA,-0.27,NA,NA,-0.19,
            -0.07,NA,-0.31,NA,NA,0.07,NA,NA,NA,
            -0.27,-0.05),
  se_jp = c(NA,
            NA,NA,0.07,NA,NA,NA,NA,0.27,0.07,NA,
            NA,0.48,0.26,0.41,NA,NA,NA,NA,0.31,
            0.19,0.34,NA,NA,0.19,NA,0.39,NA,
            0.57,NA,NA,NA,NA,NA,NA,NA),
  jobyear = c("f",
              "f","s","s","s","s","f","f","s","s",
              "s","s","f","f","f","f","f","f","s",
              "s","s","s","s","f","s","f","f",
              "f","f","f","f","s","s","s","f","f"),
  wage = c(9L,
           1L,1L,1L,4L,9L,2L,3L,3L,3L,2L,7L,
           1L,8L,6L,3L,5L,7L,8L,4L,7L,1L,4L,
           6L,7L,1L,9L,6L,7L,7L,7L,9L,2L,5L,
           7L,3L)
)

cormat1


## [[1]]
## po ne se   jp
## po 1.00 NA NA 0.16
## ne   NA NA NA   NA
## se   NA NA NA   NA
## jp 0.16 NA NA 1.00
##
## [[2]]
## po    ne    se    jp
## po   NA    NA  0.37  0.39
## ne   NA  1.00 -0.31 -0.41
## se 0.37 -0.31    NA    NA
## jp 0.39 -0.41    NA  1.00
## 출처 meta-analaysis 경로분석모델 1단계 - 상관행렬추정 correlation using metaSEM|자유자재Paper

### 1. How to obtain the correlation matrix using the fixed effects model
library(metaSEM)
fixed1 <- tssem1(Cov=cormat1, n=rd2$n, method="FEM")
summary(fixed1)
coef(fixed1)

## generate comment
 poolcor_comment(fixed1)

## End(Not run)


\end{ExampleCode}
\end{Examples}
\inputencoding{utf8}
\HeaderA{posthoc\_duncan}{posthoc\_duncan}{posthoc.Rul.duncan}
%
\begin{Description}
posthoc\_duncan
\end{Description}
%
\begin{Usage}
\begin{verbatim}
posthoc_duncan(aov_res, group_var)
\end{verbatim}
\end{Usage}
%
\begin{Arguments}
\begin{ldescription}
\item[\code{aov\_res}] aov result

\item[\code{group\_var}] group
\end{ldescription}
\end{Arguments}
%
\begin{Value}
df
\end{Value}
%
\begin{Examples}
\begin{ExampleCode}
## Not run: 
#'
aov_post_eng <- pwj3_long %>%
  filter(time == "Post" & factor == "eng") %>%
  aov(formula = score ~ Eng_Class)

# Duncan 사후분석 수행
posthoc_duncan(aov_post_eng, "Eng_Class")

## End(Not run)


\end{ExampleCode}
\end{Examples}
\inputencoding{utf8}
\HeaderA{posthoc\_plot}{Post-hoc results and plots}{posthoc.Rul.plot}
%
\begin{Description}
Post-hoc results and plots
\end{Description}
%
\begin{Usage}
\begin{verbatim}
posthoc_plot(
  aov_data,
  type = "all",
  caption = "Table. Multiple comparison",
  console = FALSE,
  group = TRUE,
  adjust = "none"
)
\end{verbatim}
\end{Usage}
%
\begin{Arguments}
\begin{ldescription}
\item[\code{aov\_data}] aov data

\item[\code{type}] type is all, res, g, post hoc

\item[\code{adjust}] node fdr bonferroni
\end{ldescription}
\end{Arguments}
%
\begin{Value}
plot and analysis result
\end{Value}
%
\begin{Examples}
\begin{ExampleCode}

## Not run: 

#Need to replace cyl with factor first. Change to as_trt() first
aov(mpg ~ cyl, data=mtcars %>%as_trt("cyl"))%>%posthoc_plot()

##Error in emmeans::emmeans(., iv) :
aov(mpg ~ factor(cyl), data=mtcars)%>%posthoc_plot()


## End(Not run)

\end{ExampleCode}
\end{Examples}
\inputencoding{utf8}
\HeaderA{PRE}{PRE(proportional reduction in error)}{PRE}
%
\begin{Description}
PRE(proportional reduction in error)
\end{Description}
%
\begin{Usage}
\begin{verbatim}
PRE(model1, model2)
\end{verbatim}
\end{Usage}
%
\begin{Arguments}
\begin{ldescription}
\item[\code{model1}] model 1

\item[\code{model2}] model q mixed model
\end{ldescription}
\end{Arguments}
%
\begin{Value}
random slope, inetcept
\end{Value}
%
\begin{Examples}
\begin{ExampleCode}
## Not run: 
 PRE(gmlm.dich.m1b, gmlm.dich.m2c)

## End(Not run)
\end{ExampleCode}
\end{Examples}
\inputencoding{utf8}
\HeaderA{Pull}{Pull select data.frame col to vector}{Pull}
%
\begin{Description}
Pull select data.frame col to vector
\end{Description}
%
\begin{Usage}
\begin{verbatim}
Pull(df, col = "Freq", name = "term", sep = "")
\end{verbatim}
\end{Usage}
%
\begin{Arguments}
\begin{ldescription}
\item[\code{df}] df

\item[\code{col}] select col, "Freq"

\item[\code{name}] using names from rownames or col data"name"

\item[\code{sep}] paste sep = "\_"
\end{ldescription}
\end{Arguments}
%
\begin{Value}
vector
\end{Value}
%
\begin{Examples}
\begin{ExampleCode}

## Not run: 
#'
# 예시 데이터
df <- tibble::tibble(
  term = c("학교 수업에서", "혼자서 공부할 때", "활용 안함"),
  Freq = c(2058, 1248, 646)
)

# 함수 사용 예시
Pull(df, col="Freq", name="term")

## End(Not run)

\end{ExampleCode}
\end{Examples}
\inputencoding{utf8}
\HeaderA{Pull\_to\_text}{Convert Data Frame Row to Text with Rounded Digits}{Pull.Rul.to.Rul.text}
%
\begin{Description}
Convert Data Frame Row to Text with Rounded Digits
\end{Description}
%
\begin{Usage}
\begin{verbatim}
Pull_to_text(df, digits = 3)
\end{verbatim}
\end{Usage}
%
\begin{Arguments}
\begin{ldescription}
\item[\code{df}] A data frame or tibble with a single row.

\item[\code{digits}] Integer, number of decimal places to round numeric values (default = 3).
\end{ldescription}
\end{Arguments}
%
\begin{Value}
A character string representing the key-value pairs from the data frame.
\end{Value}
%
\begin{Examples}
\begin{ExampleCode}
df <- tibble::tibble(R = 0.205678, adj.R = 0.204123, F = 195.456, p = 3.67e-222, dof = 6, AIC = 10028.1234, BIC = 10080.9876)
Pull_to_text(df, digits = 3)
\end{ExampleCode}
\end{Examples}
\inputencoding{utf8}
\HeaderA{p\_mark\_sig}{p\_mark\_sig is Function that displays significance in data}{p.Rul.mark.Rul.sig}
%
\begin{Description}
p\_mark\_sig is Function that displays significance in data
\end{Description}
%
\begin{Usage}
\begin{verbatim}
p_mark_sig(
  data,
  col = "p.value",
  unite = FALSE,
  unite_col = NULL,
  digits = 3,
  ns = "ns",
  rownames_to_column = FALSE,
  exclude = NULL,
  remove = TRUE
)
\end{verbatim}
\end{Usage}
%
\begin{Arguments}
\begin{ldescription}
\item[\code{data}] input data.frame or matrix

\item[\code{col}] 'col' specifies the name of the column representing the p value in the data. By default, it is set to p.value. For example, if it is written as p, pvalue, p\_value, etc., you can change the name.

\item[\code{unite}] TRUE combine pvalue and sig

\item[\code{unite\_col}] unite\_col is the variable we want to merge with sig

\item[\code{digits}] default 3

\item[\code{ns}] ns="ns"

\item[\code{rownames\_to\_column}] drownames\_to\_column if you need

\item[\code{exclude}] Round exclude varible

\item[\code{remove}] remove var TRUE
\end{ldescription}
\end{Arguments}
%
\begin{Examples}
\begin{ExampleCode}

## Not run: 
## data(ToothGrowth)
library(dplyr)
library(rstatix)

##method 1
 ToothGrowth %>%
  group_by(dose) %>%
  rstatix::t_test(data =., len ~ supp) %>%
  p_mark_sig("p")

  ### NEW
  TukeyHSD(aov(weight ~ feed, chickwts))$feed %>% as.data.frame() %>%
  mutate(sig = add_sig(`p adj`))

  TukeyHSD(aov(weight ~ feed, chickwts))$feed %>% p_mark_sig("p adj")


##method 2
 ToothGrowth %>%
 group_by(dose) %>%
 rstatix::t_test(data =., len ~ supp) %>%
  p_mark_sig("p", unite=TRUE)


 ToothGrowth %>%
     group_by(dose) %>%
    rstatix::t_test(data =., len ~ supp) %>%
    p_mark_sig("p", unite = TRUE, unite_col = "statistic" ,ns="")

 ToothGrowth %>%
   group_by(dose) %>%
   rstatix::t_test(data =., len ~ supp) %>%
   p_mark_sig("p", unite = TRUE, unite_col = "statistic" )

## End(Not run)


\end{ExampleCode}
\end{Examples}
\inputencoding{utf8}
\HeaderA{p\_sig\_cal}{observed expect table with calculate\_chi\_sig}{p.Rul.sig.Rul.cal}
%
\begin{Description}
observed expect table with calculate\_chi\_sig
\end{Description}
%
\begin{Usage}
\begin{verbatim}
p_sig_cal(obs_data)
\end{verbatim}
\end{Usage}
%
\begin{Arguments}
\begin{ldescription}
\item[\code{obs\_data}] matrix
\end{ldescription}
\end{Arguments}
%
\begin{Value}
matrix
\end{Value}
\inputencoding{utf8}
\HeaderA{p\_to\_001}{p value  to <.001}{p.Rul.to.Rul.001}
%
\begin{Description}
p value  to <.001
\end{Description}
%
\begin{Usage}
\begin{verbatim}
p_to_001(df, ..., digits = 3, remove = TRUE)
\end{verbatim}
\end{Usage}
%
\begin{Arguments}
\begin{ldescription}
\item[\code{df}] df data

\item[\code{...}] col names and number

\item[\code{digits}] 3

\item[\code{remove}] default TRUE, col replace
\end{ldescription}
\end{Arguments}
%
\begin{Value}
df
\end{Value}
%
\begin{Examples}
\begin{ExampleCode}
## Not run: 
#'

# 예제 데이터
library(dplyr)

df <- tibble(
  pairwise = c("30분이상_20-30분", "30분이상_10-20분", "30분이상_5-10분", "30분이상_5분이내", "20-30분_10-20분",
               "20-30분_5-10분", "20-30분_5분이내", "10-20분_5-10분", "10-20분_5분이내", "5-10분_5분이내"),
  chisq = c(198, 275, 349, 441, 9.27, 34.2, 87.2, 8.4, 44.6, 16.0),
  df = c(1, 1, 1, 1, 1, 1, 1, 1, 1, 1),
  p.value = c(5.83e-45, 1.09e-61, 8.37e-78, 7.50e-98, 2.33e-3, 4.85e-9, 9.66e-21, 3.76e-3, 2.44e-11, 6.30e-5),
  adj.p = c(1.46e-44, 3


            # 함수 실행 예시
            # 열 번호로 처리
            result <- p_to_001(df, 4, 6, remove = TRUE)

            # 열 이름으로 처리
            result2 <- p_to_001(df, "p.value", "adj.p", remove = TRUE)

            print(result)  # 열 번호로 처리한 결과
            print(result2) # 열 이름으로 처리한 결과
#'

## End(Not run)
\end{ExampleCode}
\end{Examples}
\inputencoding{utf8}
\HeaderA{p\_value\_cal}{Observation Expectation Table with calculate\_chi\_sig}{p.Rul.value.Rul.cal}
%
\begin{Description}
Observation Expectation Table with calculate\_chi\_sig
\end{Description}
%
\begin{Usage}
\begin{verbatim}
p_value_cal(obs_data)
\end{verbatim}
\end{Usage}
%
\begin{Arguments}
\begin{ldescription}
\item[\code{obs\_data}] matrix
\end{ldescription}
\end{Arguments}
%
\begin{Value}
matrix
\end{Value}
\inputencoding{utf8}
\HeaderA{read\_text}{read\_text: Read a Text File with Flexible Handling Using read.delim}{read.Rul.text}
%
\begin{Description}
read\_text: Read a Text File with Flexible Handling Using read.delim
\end{Description}
%
\begin{Usage}
\begin{verbatim}
read_text(file_path, header = TRUE, stringsAsFactors = FALSE, sep = "\t")
\end{verbatim}
\end{Usage}
%
\begin{Arguments}
\begin{ldescription}
\item[\code{file\_path}] The path of the file to read.

\item[\code{header}] Logical. Does the file contain a header? Defaults to TRUE.

\item[\code{stringsAsFactors}] Logical. Should character columns be converted to factors? Defaults to FALSE.
\end{ldescription}
\end{Arguments}
%
\begin{Value}
A data frame containing the file's contents.
\end{Value}
%
\begin{Examples}
\begin{ExampleCode}
## Not run: 
# Save and read a file
teacher_question <- c("Example question 1", "Example question 2")
write_text(teacher_question, file_path = "output.txt")
result <- read_text("output.txt")

## End(Not run)

\end{ExampleCode}
\end{Examples}
\inputencoding{utf8}
\HeaderA{recode\_sep\_vec}{recode\_sep\_vec 벡터데이터를 리코딩}{recode.Rul.sep.Rul.vec}
%
\begin{Description}
recode\_sep\_vec 벡터데이터를 리코딩
\end{Description}
%
\begin{Usage}
\begin{verbatim}
recode_sep_vec(vector, ranges_list)
\end{verbatim}
\end{Usage}
%
\begin{Arguments}
\begin{ldescription}
\item[\code{vector}] vector

\item[\code{ranges\_list}] 변환리스트 list("0\textasciitilde{}5년" = 0:5)
\end{ldescription}
\end{Arguments}
%
\begin{Value}
vector
\end{Value}
%
\begin{Examples}
\begin{ExampleCode}
## Not run: 
#'
# 예시 사용
vec <- c("53", "45", "38", "28", "41", "12", "0", "10", "30", "35")

# 함수를 사용하여 벡터 재코딩
new_vec <- recode_sep_vec(vec, list("0~5년" = 0:5,
                                    "5~10년" = 6:10,
                                    "11~15년" = 11:15,
                                    "16~20년" = 16:20,
                                    "21~25년" = 21:25,
                                    "26~30년" = 26:30,
                                    "30년이상" = 31:60))

# 결과 확인
print(new_vec)

# # # 예시 사용
# vec <- c(53, 45, 45, 38, 36, 28, 38, 28, 41, 55, 39, 29, 31, 35, 37, 34, 34, 35, 58, 46)

# # 함수를 사용하여 벡터 재코딩
new_vec <- recode_sep_vec(vec, list("20대" = 20:29, "30대" = 30:39, "40대" = 40:49, "50대" = 50:59))

# # 결과 확인
# print(new_vec)

## End(Not run)


\end{ExampleCode}
\end{Examples}
\inputencoding{utf8}
\HeaderA{reg\_apa}{regression APA report}{reg.Rul.apa}
%
\begin{Description}
regression APA report
\end{Description}
%
\begin{Usage}
\begin{verbatim}
reg_apa(model, digits = 3, md = FALSE)
\end{verbatim}
\end{Usage}
%
\begin{Arguments}
\begin{ldescription}
\item[\code{model}] lm\_data

\item[\code{digits}] 3

\item[\code{md}] markdown table
\end{ldescription}
\end{Arguments}
%
\begin{Value}
report
\end{Value}
%
\begin{Examples}
\begin{ExampleCode}
## Not run: 
lm(mpg ~ hp + wt + disp + qsec, data = mtcars) %>% reg_apa()

## End(Not run)
reg1
\end{ExampleCode}
\end{Examples}
\inputencoding{utf8}
\HeaderA{remove\_all\_attr}{Remove all ATTRs present in the data}{remove.Rul.all.Rul.attr}
%
\begin{Description}
Remove all ATTRs present in the data
\end{Description}
%
\begin{Usage}
\begin{verbatim}
remove_all_attr(df)
\end{verbatim}
\end{Usage}
%
\begin{Arguments}
\begin{ldescription}
\item[\code{df}] data
\end{ldescription}
\end{Arguments}
%
\begin{Value}
data
\end{Value}
\inputencoding{utf8}
\HeaderA{remove\_attr}{Remove all ATTRs present in the data, deselect}{remove.Rul.attr}
%
\begin{Description}
Remove all ATTRs present in the data, deselect
\end{Description}
%
\begin{Usage}
\begin{verbatim}
remove_attr(df, range = 1:ncol(df))
\end{verbatim}
\end{Usage}
%
\begin{Arguments}
\begin{ldescription}
\item[\code{df}] data

\item[\code{range}] col range
\end{ldescription}
\end{Arguments}
%
\begin{Value}
data.frame
\end{Value}
%
\begin{Examples}
\begin{ExampleCode}

## Not run: 
#'
# Example usage with your data frame jut3
# Remove attributes from all columns
# jut3_cleaned_all <- remove_attr(jut3)

# Remove attributes from columns 1 to 5
# jut3_cleaned_partial <- remove_attr(jut3, range = 1:5)

# Check the structure
# str(jut3_cleaned_all)
# str(jut3_cleaned_partial)

## End(Not run)


\end{ExampleCode}
\end{Examples}
\inputencoding{utf8}
\HeaderA{remove\_blank}{데이터내에서 빈칸 제거}{remove.Rul.blank}
%
\begin{Description}
데이터내에서 빈칸 제거
\end{Description}
%
\begin{Usage}
\begin{verbatim}
remove_blank(
  data,
  col = NULL,
  data_type = c("data.frame", "vector", "character")
)
\end{verbatim}
\end{Usage}
%
\begin{Arguments}
\begin{ldescription}
\item[\code{data}] df

\item[\code{col}] col 숫자와 열이름으로 지정

\item[\code{data\_type}] data.frame, vector
\end{ldescription}
\end{Arguments}
%
\begin{Value}
df
\end{Value}
%
\begin{Examples}
\begin{ExampleCode}
## Not run: 
#'
# 예제 데이터
vector_data <- c("익산 어양 초", "익산옥 야 초")
df_data <- data.frame(
  ID = c(259, 260),
  Name = c("익산 어양 초", "익산옥 야 초"),
  Value = c(1, 1),
  Score = c(0.1237624, 0.1237624)
)

# 벡터에서 공백 제거
clean_vector <- remove_blank(vector_data, data_type = "vector")
print(clean_vector)

# 데이터프레임에서 공백 제거 (열 이름으로)
clean_df <- remove_blank(df_data, col = "Name")
print(clean_df)

# 데이터프레임에서 공백 제거 (열 번호로)
clean_df <- remove_blank(df_data, col = 2)
print(clean_df)
#'


## End(Not run)
\end{ExampleCode}
\end{Examples}
\inputencoding{utf8}
\HeaderA{remove\_cols}{Delete cols in the data}{remove.Rul.cols}
%
\begin{Description}
Delete cols in the data
\end{Description}
%
\begin{Usage}
\begin{verbatim}
remove_cols(data, rm_col = NULL, ...)
\end{verbatim}
\end{Usage}
%
\begin{Arguments}
\begin{ldescription}
\item[\code{data}] data

\item[\code{rm\_col}] first column name or number

\item[\code{...}] add columns
\end{ldescription}
\end{Arguments}
%
\begin{Value}
data
\end{Value}
%
\begin{Examples}
\begin{ExampleCode}
## Not run: 
#'
# 예시 사용법
col_remove(mtcars, 1, "am", 2, "vs", "wt")

# 예시 사용법
col_remove(mtcars, 1, "am", 2, "vs", "wt",4)


# 예시 사용법
col_remove(mtcars, 1, "am", 2, "vs", "wt",3)


# 예시 사용법
col_remove(mtcars, 1, "am", 2, "vs", "wt",3,4)

# 예시 사용법
col_remove(mtcars, rm_col=1:3)
col_remove(mtcars, c(1, 3))
col_remove(mtcars, 1,2, 3)
col_remove(mtcars, 1, "am",2,"vs","wt")
col_remove(mtcars, rm_col="mpg")
col_remove(mtcars, "mpg", "am","vs")


## End(Not run)


\end{ExampleCode}
\end{Examples}
\inputencoding{utf8}
\HeaderA{remove\_filter}{Functions to remove selected rows}{remove.Rul.filter}
%
\begin{Description}
Functions to remove selected rows
\end{Description}
%
\begin{Usage}
\begin{verbatim}
remove_filter(data, ...)
\end{verbatim}
\end{Usage}
%
\begin{Arguments}
\begin{ldescription}
\item[\code{data}] data

\item[\code{...}] condition filter
\end{ldescription}
\end{Arguments}
%
\begin{Value}
data
\end{Value}
%
\begin{Examples}
\begin{ExampleCode}

## Not run: 
#'
mtcars %>% remove_filter(cyl == 6 )
mtcars %>% remove_filter(vs == 1)
mtcars %>% remove_filter(am==0)
mtcars %>% remove_filter(vs == 1)
mtcars %>% remove_filter(cyl == 6 | vs == 1)
mtcars %>% remove_filter(cyl == 6 & vs == 1)
mtcars %>% remove_filter(cyl == 6 | cyl == 4 | am == 0)

## End(Not run)

\end{ExampleCode}
\end{Examples}
\inputencoding{utf8}
\HeaderA{remove\_rows}{remove rows}{remove.Rul.rows}
%
\begin{Description}
remove rows
\end{Description}
%
\begin{Usage}
\begin{verbatim}
remove_rows(data, ..., filter = FALSE)
\end{verbatim}
\end{Usage}
%
\begin{Arguments}
\begin{ldescription}
\item[\code{data}] data

\item[\code{...}] rows

\item[\code{filter}] TRUE possible filter
\end{ldescription}
\end{Arguments}
%
\begin{Value}
data
\end{Value}
%
\begin{Examples}
\begin{ExampleCode}

## Not run: 
#'
# 1~3행 제거
remove_rows(mtcars,  1:30)
remove_rows(mtcars, 1,2,3,6,5,7,8,9,30)
remove_rows(mtcars, rm = c(1,3,4,5,6,7,11:30))
remove_rows(mtcars, cyl == 4, filter=T)

# 1행과 3행 제거
print(remove_rows(mtcars, rm_row = c(1, 3)))

# 필터 조건 적용 (cyl == 6)
remove_rows(mtcars,rm=NULL, cyl == 6) #error
remove_rows(mtcars, cyl == 6 |cyl ==4, filter=T)

# 필터 조건 적용 (cyl == 6 또는 vs == 1)
remove_rows(mtcars, cyl == 6 | vs == 1, filter=T)

# 필터 조건 적용 (cyl == 6 그리고 vs == 1)
remove_rows(mtcars, cyl == 6 | vs == 1, filter=T)

# 필터 조건 적용 후 특정 행 제거 (필터가 우선 적용)
remove_rows(mtcars, 1:3, cyl == 6 | vs == 1)

#'

## End(Not run)


\end{ExampleCode}
\end{Examples}
\inputencoding{utf8}
\HeaderA{rename\_by\_dict}{rename\_by\_dict: Rename both column names and cell values using a dictionary}{rename.Rul.by.Rul.dict}
%
\begin{Description}
This function renames column names and character cell values in a data frame
using a named character vector as dictionary. Especially useful in pipes.
\end{Description}
%
\begin{Usage}
\begin{verbatim}
rename_by_dict(df, dict)
\end{verbatim}
\end{Usage}
%
\begin{Arguments}
\begin{ldescription}
\item[\code{df}] A data.frame or tibble.

\item[\code{dict}] A character vector of key-value pairs (e.g., c("eng1" = "kor1", ...)) or unnamed like c("eng1", "kor1", ...)
\end{ldescription}
\end{Arguments}
%
\begin{Value}
A data.frame with renamed column names and character cell values.
\end{Value}
%
\begin{Examples}
\begin{ExampleCode}
df <- data.frame(latent = c("access", "barrier"), access = 1:2, barrier = 3:4)
var_named <- c("access", "연구접근성", "barrier", "진로장벽")
rename_by_dict(df, var_named)
\end{ExampleCode}
\end{Examples}
\inputencoding{utf8}
\HeaderA{reorder\_row}{Reorder rows in a dataframe with a given row number}{reorder.Rul.row}
%
\begin{Description}
Reorder rows in a dataframe with a given row number
\end{Description}
%
\begin{Usage}
\begin{verbatim}
reorder_row(data, ...)
\end{verbatim}
\end{Usage}
%
\begin{Arguments}
\begin{ldescription}
\item[\code{data}] data.

\item[\code{...}] row order
\end{ldescription}
\end{Arguments}
%
\begin{Value}
Reordered data
\end{Value}
%
\begin{Examples}
\begin{ExampleCode}
## Not run: 
Create a dataframe
df <- data.frame(
  response = c("그렇다", "그렇지않다", "매우그렇다", "보통이다", "전혀그렇지않다"),
  Freq = c(5, 13, 4, 14, 4)
)
df
# Call a function to reorder rows
reorder_row(df, 5, 2, 4, 1,3)  # Put the row numbers in order.

#data process
data_long <- tribble(
  ~response, ~freq,
  "전혀그렇지않다", 1,
  "그렇지않다", 2,
  "보통이다", 3,
  "그렇다", 2,
  "매우그렇다", 0,
  "전혀그렇지않다", 0,
  "그렇지않다", 4,
  "보통이다", 4,
  "그렇다", 0,
  "매우그렇다", 0,
  "전혀그렇지않다", 1,
  "그렇지않다", 3,
  "보통이다", 4,
  "그렇다", 0,
  "매우그렇다", 0,
  "전혀그렇지않다", 1,
  "그렇지않다", 4,
  "보통이다", 3,
  "그렇다", 0,
  "매우그렇다", 0
)
data_long <- data_long %>%as_trt("response")
data_long

data_long %>%add_rows_freq() %>% rowid_to_column("person")%>%
  select(response)%>%table()%>%dall()%>%
  reorder_row(5,2,4,1,3)


## End(Not run)

\end{ExampleCode}
\end{Examples}
\inputencoding{utf8}
\HeaderA{Replace}{Create a function to replace NA with an empty string in an R vector data}{Replace}
%
\begin{Description}
Create a function to replace NA with an empty string in an R vector data
\end{Description}
%
\begin{Usage}
\begin{verbatim}
Replace(vector, pattern = NA, imp = "")
\end{verbatim}
\end{Usage}
%
\begin{Arguments}
\begin{ldescription}
\item[\code{vector}] input data

\item[\code{pattern}] dectect pattern

\item[\code{imp}] imputation
\end{ldescription}
\end{Arguments}
%
\begin{Value}
imputation data
\end{Value}
%
\begin{Examples}
\begin{ExampleCode}

## Not run: 
# Replace(c("A", NA, "B", "", "C", "C"), "C", imp="T")

## End(Not run)

\end{ExampleCode}
\end{Examples}
\inputencoding{utf8}
\HeaderA{replace\_dataframe}{replace\_dataframe}{replace.Rul.dataframe}
%
\begin{Description}
replace\_dataframe
\end{Description}
%
\begin{Usage}
\begin{verbatim}
replace_dataframe(data, pattern, imp)
\end{verbatim}
\end{Usage}
%
\begin{Arguments}
\begin{ldescription}
\item[\code{data}] data

\item[\code{pattern}] patern

\item[\code{imp}] imputation
\end{ldescription}
\end{Arguments}
%
\begin{Value}
data
\end{Value}
\inputencoding{utf8}
\HeaderA{replace\_df}{Search for and fill in empty data}{replace.Rul.df}
%
\begin{Description}
Search for and fill in empty data
\end{Description}
%
\begin{Usage}
\begin{verbatim}
replace_df(df, pattern = NA, imp = "", col = NULL)
\end{verbatim}
\end{Usage}
%
\begin{Arguments}
\begin{ldescription}
\item[\code{df}] data.frame

\item[\code{pattern}] default ""

\item[\code{imp}] iimputation , defalut NA
\end{ldescription}
\end{Arguments}
%
\begin{Value}
data.frame
\end{Value}
%
\begin{Examples}
\begin{ExampleCode}

## Not run: 
#'
New = data.frame(term = c("(Intercept)", "age30대",
 "age40대", "gendermale", "manufacturergoogle",
 "manufacturersamsung",
 "age30대:gendermale", 'age40대:gendermale',
  'age30대:manufacturergoogle', "age40대:manufacturergoogle",
  "age30대:manufacturersamsung", 'age40대:manufacturersamsung'),
estimate = c(6.72127226828228, 0.222516673400011, 0.613427717767302,
 -0.451934840902044, -0.0405765656320795, 0.380612934433137,
  0.00538723055632342, 0.52238521880372, 0.0022025527504637,
  0.0243934423439851, 0.509454602341916, 0.00845576255419509),
std.error = c(0.029551066353128, 0.0393330112847198,
0.0358559955535251, 0.0300846458760622, 0.0387717417205405,
 0.0352086354160786, 0.0383047023682364, 0.0356296478558013,
 0.0519877886919915, 0.0462375166064463,
 0.0454909790784573, 0.0420442735967295),
statistic = c(227.446014568974, 5.65724988074979,
 17.1080933131975, -15.0221093764524, -1.04654998283409,
 10.8102154467288, 0.140641493687488, 14.6615319050554,
  0.0423667327632075, 0.527568176976533,
   11.1990247882612, 0.201115677138321),
p.value = c(0, 1.53817738615173e-08, 1.29158703632758e-65,
5.26066233107499e-51, 0.295307180741282, 3.0794597517271e-27,
 0.888153166786026, 1.13678932467476e-48,
 0.966206347935429, 0.59779908524448,
 4.12246475726497e-29, 0.840608123856931)
)%>%tibble()

New

dfs= New%>%p_mark_sig(ns="", unite=TRUE, unite_col="estimate")

dfs%>%
  replace_df("0.000","< .001", col="p.value")

dfs%>%
  replace_df("0.000","< .001", col=c("p.value", "std.error"))



## End(Not run)


\end{ExampleCode}
\end{Examples}
\inputencoding{utf8}
\HeaderA{replace\_df2}{Search for and fill in NA data}{replace.Rul.df2}
%
\begin{Description}
Search for and fill in NA data
\end{Description}
%
\begin{Usage}
\begin{verbatim}
replace_df2(df, pattern = NA, imp = "", cat = TRUE)
\end{verbatim}
\end{Usage}
%
\begin{Arguments}
\begin{ldescription}
\item[\code{df}] data.frame

\item[\code{pattern}] default ""

\item[\code{imp}] iimputation , defalut NA
\end{ldescription}
\end{Arguments}
%
\begin{Value}
data.frame
\end{Value}
\inputencoding{utf8}
\HeaderA{replace\_df\_rep}{replace\_df\_repeated adation}{replace.Rul.df.Rul.rep}
%
\begin{Description}
replace\_df\_repeated adation
\end{Description}
%
\begin{Usage}
\begin{verbatim}
replace_df_rep(data, ..., detect = TRUE, chr = TRUE)
\end{verbatim}
\end{Usage}
%
\begin{Arguments}
\begin{ldescription}
\item[\code{data}] data

\item[\code{...}] pattern and imputation

\item[\code{detect}] detect TRUE in data, FALSE data

\item[\code{chr}] transformation fac  to chr
\end{ldescription}
\end{Arguments}
%
\begin{Value}
data
\end{Value}
%
\begin{Examples}
\begin{ExampleCode}

## Not run: 

# Use case 1: Passing pattern-replacement string
# pairs directly as variable arguments
data <- data.frame(
  col1 = c("aa", "bb", "cc"),
  col2 = c("aa", "bb", "dd")
)

print("Before Replacement:")
print(data)

result1 <- replace_df_rep(data, "aa", "a1", "bb", "b1")
print("After Replacement (direct arguments):")
print(result1)

# Use case 2: Passing as a single vector
changes <- c("aa", "a1", "bb", "b1")
result2 <- replace_df_rep(data, changes)
print("After Replacement (vector):")
print(result2)
#'

#'
# 예시 데이터
data <- data.frame(
  Question = c("나는 나의 세 가지 가장 큰 약점을 열거할 수 있다.",
   "나의 행동은 나의 핵심 가치를 반영한다.",
   "나는 스스로 결정하기 전에 다른 사람들의 의견을 구한다."),
  factor = c("자기", "자기", "자기"),
  stringsAsFactors = FALSE
)

data
replace_df_rep(data, "약점", "강점","행동","수행",
"스스로","다른 사람이",  detect = TRUE)
replace_df_rep(data, "약점",
"강점","행동","수행",  detect = FALSE)



## End(Not run)



\end{ExampleCode}
\end{Examples}
\inputencoding{utf8}
\HeaderA{replace\_in\_dataframe}{replace\_in\_dataframe}{replace.Rul.in.Rul.dataframe}
%
\begin{Description}
replace\_in\_dataframe
\end{Description}
%
\begin{Usage}
\begin{verbatim}
replace_in_dataframe(data, changes)
\end{verbatim}
\end{Usage}
%
\begin{Arguments}
\begin{ldescription}
\item[\code{data}] data

\item[\code{changes}] change
\end{ldescription}
\end{Arguments}
%
\begin{Value}
data
\end{Value}
\inputencoding{utf8}
\HeaderA{replace\_in\_df}{replace\_in\_df}{replace.Rul.in.Rul.df}
%
\begin{Description}
replace\_in\_df
\end{Description}
%
\begin{Usage}
\begin{verbatim}
replace_in_df(data, pattern, imp)
\end{verbatim}
\end{Usage}
%
\begin{Arguments}
\begin{ldescription}
\item[\code{data}] data

\item[\code{pattern}] pattern

\item[\code{imp}] ""
\end{ldescription}
\end{Arguments}
%
\begin{Value}
data
\end{Value}
%
\begin{Examples}
\begin{ExampleCode}
## Not run: 
#'
data <- data.frame(
  Question = c("나는 나의 세 가지 가장 큰
  약점을 열거할 수 있다.", "나의 행동은 나의
  핵심 가치를 반영한다.", "나는 스스로 결정하기
  전에 다른 사람들의 의견을 구한다."),
  factor = c("자기", "자기", "자기"),
  stringsAsFactors = FALSE
)

data
data %>% replace_in_df("약점", "강점")

## End(Not run)
\end{ExampleCode}
\end{Examples}
\inputencoding{utf8}
\HeaderA{Replace\_na}{Create a function to replace NA with an empty string in an R vector data}{Replace.Rul.na}
%
\begin{Description}
Create a function to replace NA with an empty string in an R vector data
\end{Description}
%
\begin{Usage}
\begin{verbatim}
Replace_na(vector, pattern = NA, imp = "")
\end{verbatim}
\end{Usage}
%
\begin{Arguments}
\begin{ldescription}
\item[\code{vector}] input data

\item[\code{pattern}] dectect pattern

\item[\code{imp}] imputation
\end{ldescription}
\end{Arguments}
%
\begin{Value}
imputation data
\end{Value}
%
\begin{Examples}
\begin{ExampleCode}

## Not run: 
# Replace(c("A", NA, "B", "", "C", "C"), "C", imp="T")

## End(Not run)

\end{ExampleCode}
\end{Examples}
\inputencoding{utf8}
\HeaderA{report\_aov}{anova report ,report\_aov}{report.Rul.aov}
%
\begin{Description}
anova report ,report\_aov
\end{Description}
%
\begin{Usage}
\begin{verbatim}
report_aov(data, digits = 2, trans = FALSE)
\end{verbatim}
\end{Usage}
%
\begin{Arguments}
\begin{ldescription}
\item[\code{data}] aov data

\item[\code{digits}] default 3

\item[\code{trans}] TRUE english to korea
\end{ldescription}
\end{Arguments}
%
\begin{Value}
report
\end{Value}
%
\begin{Examples}
\begin{ExampleCode}
## Not run: 
aov(len ~ supp*factor(dose), data= ToothGrowth) %>% report_aov()
aov(len ~ supp*factor(dose), data= ToothGrowth) %>% report_aov(trans=TRUE)



## End(Not run)


\end{ExampleCode}
\end{Examples}
\inputencoding{utf8}
\HeaderA{report\_aov\_ko}{repot anova}{report.Rul.aov.Rul.ko}
%
\begin{Description}
repot anova
\end{Description}
%
\begin{Usage}
\begin{verbatim}
report_aov_ko(data, digits = 2, print = TRUE)
\end{verbatim}
\end{Usage}
%
\begin{Arguments}
\begin{ldescription}
\item[\code{data}] aov data

\item[\code{digits}] 3

\item[\code{print}] table show
\end{ldescription}
\end{Arguments}
%
\begin{Value}
table
\end{Value}
%
\begin{Examples}
\begin{ExampleCode}

## Not run: 
aov(mpg ~ cyl* gear +carb, data=mtcars) %>%report_aov_ko()

aov(mpg ~ cyl* gear +carb, data=mtcars) %>%report_aov_ko(print=T)
aov(mpg ~ cyl* gear +carb, data=mtcars) %>%report_aov_ko(print=F)


## End(Not run)

\end{ExampleCode}
\end{Examples}
\inputencoding{utf8}
\HeaderA{report\_lme}{linear mixed model report}{report.Rul.lme}
%
\begin{Description}
Mixed model summary
\end{Description}
%
\begin{Usage}
\begin{verbatim}
report_lme(
  lmedata,
  type = "basic",
  form = "lmer",
  apa = FALSE,
  fit_more = FALSE,
  ranef_sig = FALSE,
  show.effect = FALSE,
  show.ci = FALSE
)
\end{verbatim}
\end{Usage}
%
\begin{Arguments}
\begin{ldescription}
\item[\code{lmedata}] lmedata is lmer() function result

\item[\code{type}] = 'all' is res, 'Fixed\_effect','Random\_effect','ICC', 'ConfidenceInterval\_95','Satterthwaite\_method','FIT','APA'

\item[\code{form}] 'lmer','glmer' = 'logit'='dich', 'poisson' ='pois'

\item[\code{apa}] = default FALSE, if you  REML is FALSE

\item[\code{fit\_more}] = default FALSE, TRUE detailed report

\item[\code{ranef\_sig}] ranef\_sig =TRUE  random effect test

\item[\code{show.effect}] if you show.effect = TRUE -> show ranef, fixef, coef

\item[\code{show.ci}] TRUE calculation 95percent CI
\end{ldescription}
\end{Arguments}
%
\begin{Examples}
\begin{ExampleCode}
## Not run: 
data(sleepstudy)

sleepstudy |> str()

lmer(Reaction ~ Days + (Days | Subject), sleepstudy) |> lme_report()

lmer(Reaction ~ Days + (Days | Subject), sleepstudy) |> lme_report(apa=TRUE)


book.url <- "https://stat.ethz.ch/~meier/teaching/book-anova"
quality <- readRDS(url(file.path(book.url, "data/quality.rds")))
str(quality)
fit.quality <- lmer(score ~ (1 | employee) + (1 | batch) +
                     (1 | employee:batch), data = quality)
summary(fit.quality)
fit.quality|> lme_report()
fit.quality|> lme_report(apa= TRUE) #error the app is FALSE
fit.quality|> lme_report(apa=FALSE) #error the app is FALSE
fit.quality|> ICC()

lmer(Reaction ~ Days + (Days | Subject), sleepstudy) |>
SEM212::lme_report(apa=TRUE, fit_more = T)
lmer(Reaction ~ Days + (Days | Subject), sleepstudy) |>
SEM212::lme_report(apa=TRUE, fit_more = F)  #default FALSE, it is show that AIC and BIC


## End(Not run)

\end{ExampleCode}
\end{Examples}
\inputencoding{utf8}
\HeaderA{report\_manova}{Functions to report MANOVA results}{report.Rul.manova}
%
\begin{Description}
Functions to report MANOVA results
\end{Description}
%
\begin{Usage}
\begin{verbatim}
report_manova(model_manova, type = "res", digits = 2)
\end{verbatim}
\end{Usage}
%
\begin{Arguments}
\begin{ldescription}
\item[\code{model\_manova}] manova result

\item[\code{type}] default res, and all, summary\_aov, Pillai, Wilks, Hotelling\_Lawley, Roy, Multivariate\_Tests

\item[\code{digits}] default 2
\end{ldescription}
\end{Arguments}
%
\begin{Value}
report MANOVA table
\end{Value}
%
\begin{Examples}
\begin{ExampleCode}
## Not run: 

Skulls = heplots::data(Skulls)

# execute
manova(cbind(mb,bh,bl,nh)~ epoch, data=Skulls) %>%report_manova()

manova(cbind(mb,bh,bl,nh)~ epoch, data=Skulls) %>%report_manova(type="all")

manova(cbind(mb,bh,bl,nh)~ epoch, data=Skulls) %>%report_manova(type="summmary_aov")

manova(cbind(mb,bh,bl,nh)~ epoch, data=Skulls) %>%report_manova(type="Pillai")

manova(cbind(mb,bh,bl,nh)~ epoch, data=Skulls) %>%report_manova(type="Wilks")

manova(cbind(mb,bh,bl,nh)~ epoch, data=Skulls) %>%report_manova(type="Hotelling_Lawley")

manova(cbind(mb,bh,bl,nh)~ epoch, data=Skulls) %>%report_manova(type="Roy")

manova(cbind(mb,bh,bl,nh)~ epoch, data=Skulls) %>%report_manova(type="Multivariate_Tests")



Y <- cbind.data.frame(tear=c(6.5, 6.2, 5.8, 6.5, 6.5, 6.9, 7.2, 6.9, 6.1, 6.3,
                             6.7, 6.6, 7.2, 7.1, 6.8, 7.1, 7.0, 7.2, 7.5, 7.6),
                      gloss=c(9.5, 9.9, 9.6, 9.6, 9.2, 9.1, 10.0, 9.9, 9.5, 9.4,
          9.1, 9.3, 8.3, 8.4, 8.5, 9.2, 8.8, 9.7, 10.1, 9.2),
          opacity=c(4.4, 6.4, 3.0, 4.1, 0.8, 5.7, 2.0, 3.9, 1.9, 5.7,
                    2.8, 4.1, 3.8, 1.6, 3.4, 8.4, 5.2, 6.9, 2.7, 1.9),
          rate = gl(2,10, labels = c("Low", "High")),
          additive = gl(2, 5, length = 20, labels = c("Low", "High")) )

manova(cbind(tear, gloss, opacity) ~ rate * additive, data=Y)%>%report_manova()

manova(cbind(tear, gloss, opacity) ~ rate * additive, data=Y)%>%report_manova(type="all")
manova(cbind(tear, gloss, opacity) ~ rate * additive, data=Y)%>%
                           report_manova(type="summmary_aov")
manova(cbind(tear, gloss, opacity) ~ rate * additive, data=Y) %>%
  report_manova("Pillai")
manova(cbind(tear, gloss, opacity) ~ rate * additive, data=Y)%>%
  report_manova("Wilks")
manova(cbind(tear, gloss, opacity) ~ rate * additive, data=Y)%>%
  report_manova("Hotelling_Lawley")
manova(cbind(tear, gloss, opacity) ~ rate * additive, data=Y)%>%
  report_manova("Roy")
manova(cbind(tear, gloss, opacity) ~ rate * additive, data=Y)%>%
  report_manova("Multivariate_Tests")


## End(Not run)


\end{ExampleCode}
\end{Examples}
\inputencoding{utf8}
\HeaderA{report\_stat}{report nomality , t.test , var.test}{report.Rul.stat}
%
\begin{Description}
report nomality , t.test , var.test
\end{Description}
%
\begin{Usage}
\begin{verbatim}
report_stat(data, type = "normality", digits = 3, trans = FALSE, show = "data")
\end{verbatim}
\end{Usage}
%
\begin{Arguments}
\begin{ldescription}
\item[\code{data}] test result

\item[\code{type}] normality, t.test, var.test(leven)

\item[\code{digits}] default 3

\item[\code{trans}] translate TRUE korean language

\item[\code{show}] normal is source and output, data is output
\end{ldescription}
\end{Arguments}
%
\begin{Value}
report data
\end{Value}
%
\begin{Examples}
\begin{ExampleCode}
## Not run: 

shapiro.test(final) %>% report_stat()

var.test(mpg ~ vs, data = mtcars)%>%
report_stat("var.test")

car::leveneTest(mpg~ factor(vs), mtcars)%>%
 report_stat("leven.test")

car::leveneTest(mpg~ factor(vs), mtcars)%>%
report_stat("leven.test", trans=TRUE)


fligner.test(mpg ~ vs,data= mtcars) %>%
report_stat("t.test")

t.test(mpg ~ vs, data= mtcars, var.equal = FALSE)%>%
report_stat("t.test")

t.test(len ~ supp, data = ToothGrowth, var.equal=TRUE) %>%
report_stat("t.test")

shapiro.test(final) %>% report_stat()%>% jjstat::k("en","ko")
shapiro.test(final) %>% report_stat(trans =TRUE)

car::leveneTest(mpg~ factor(vs), mtcars)%>%
report_stat("leven.test", trans=TRUE)



## End(Not run)

\end{ExampleCode}
\end{Examples}
\inputencoding{utf8}
\HeaderA{report\_ttest}{report\_ttest}{report.Rul.ttest}
%
\begin{Description}
report\_ttest
\end{Description}
%
\begin{Usage}
\begin{verbatim}
report_ttest(ttest_result, trans = FALSE, type = "report")
\end{verbatim}
\end{Usage}
%
\begin{Arguments}
\begin{ldescription}
\item[\code{ttest\_result}] t.test result

\item[\code{trans}] true english to korean

\item[\code{type}] default 'report', 'all' is t.table and APA
\end{ldescription}
\end{Arguments}
%
\begin{Value}
APA report
\end{Value}
%
\begin{Examples}
\begin{ExampleCode}
## Not run: 
## english report
t.test(mpg ~ vs, mtcars, var.equal=TRUE) %>%
        report_ttest()

## translation
t.test(mpg ~ vs, mtcars, var.equal=TRUE) %>%
report_ttest(trans=T)

t.test(mpg ~ vs, mtcars) %>% report_ttest("all")

## End(Not run)
\end{ExampleCode}
\end{Examples}
\inputencoding{utf8}
\HeaderA{report\_vartest}{var.test report}{report.Rul.vartest}
%
\begin{Description}
var.test report
\end{Description}
%
\begin{Usage}
\begin{verbatim}
report_vartest(var_test, trans = FALSE)
\end{verbatim}
\end{Usage}
%
\begin{Arguments}
\begin{ldescription}
\item[\code{var\_test}] var.test result

\item[\code{trans}] tans true english to korean
\end{ldescription}
\end{Arguments}
%
\begin{Value}
explain result apa
\end{Value}
%
\begin{Examples}
\begin{ExampleCode}
## Not run: 
var.test(mpg ~ vs, data = mtcars)%>%
     report_vartest()

var.test(mpg ~ vs, data = mtcars)%>%
     report_vartest() %>%
     jjstat::kakaoi()

var.test(mpg ~ vs, data = mtcars)%>%
report_vartest(trans=TRUE)


## End(Not run)
\end{ExampleCode}
\end{Examples}
\inputencoding{utf8}
\HeaderA{revcode}{Custom function to perform reverse coding for multiple variables with custom n values}{revcode}
%
\begin{Description}
Custom function to perform reverse coding for multiple variables with custom n values
\end{Description}
%
\begin{Usage}
\begin{verbatim}
revcode(data, n = 8, ..., postfix = "_r", replace = FALSE)
\end{verbatim}
\end{Usage}
%
\begin{Arguments}
\begin{ldescription}
\item[\code{data}] Data frame containing the variables

\item[\code{n}] Vector of maximum scale values (one for each vadata:image/png;base64,iVBORw0KGgoAAAANSUhEUgAAABIAAAASCAYAAABWzo5XAAAAbElEQVR4Xs2RQQrAMAgEfZgf7W9LAguybljJpR3wEse5JOL3ZObDb4x1loDhHbBOFU6i2Ddnw2KNiXcdAXygJlwE8OFVBHDgKrLgSInN4WMe9iXiqIVsTMjH7z/GhNTEibOxQswcYIWYOR/zAjBJfiXh3jZ6AAAAAElFTkSuQmCCriable), c("v1","v2",...)

\item[\code{...}] Variable names to be reverse coded

\item[\code{postfix}] postfix is reverse variable name default '\_r'

\item[\code{replace}] How to enter and replace the reverse coding results
\end{ldescription}
\end{Arguments}
%
\begin{Value}
add reverse data in data.frame
\end{Value}
%
\begin{Examples}
\begin{ExampleCode}
## Not run: 

### How to decode all variables by applying only one n
revcode(stat_onl, 8, "upgrade", "satisfy", "fatigue" )|>
select(upgrade, satisfy, fatigue, upgrade_r, satisfy_r, fatigue_r)

### the method for the reverse coding variable being applied separately
### to each variable and decoding
revcode(stat_onl, n= c(8,8,8), "upgrade", "satisfy", "fatigue" )|>
select(upgrade, satisfy, fatigue, upgrade_r, satisfy_r, fatigue_r)

## mtcars data
## Since the inverse coding value of the cyl variable is 4,6,8,
## n enters 4 + 8 = 12; the minimum + maximum = n value.
revcode(mtcars[,c("cyl", "gear","carb")], 12, "cyl")


revcode(mtcars[,c("cyl", "gear","carb")],
            n = c(12, 6, 9),"cyl","gear","carb")
#'
df <- data.frame(
  A = c(1, 2, 3, 4, 5),
  B = c(5, 4, 3, 2, 1)
)

# Passing variable names as symbols
df_new1 <- revcode(df, n = 6, A, B)
print(df_new1)
str(df_new1)
print(attr(df_new1$A_r, "reverse")) # Output properties for verification
print(attr(df_new1$B_r, "reverse"))

# pass the variable name as a string
df_new2 <- revcode(df, n = 6, "A", "B")
print(df_new2)
str(df_new2)
print(attr(df_new2$A_r, "reverse"))  # Output properties for verification
print(attr(df_new2$B_r, "reverse"))

# Reverse-code to replace original column
df_replaced <- revcode(df, n = 6, A, B, replace = TRUE)
print(df_replaced)
str(df_replaced)
print(attr(df_replaced$A, "reverse"))  # Output properties for verification
print(attr(df_replaced$B, "reverse"))

#'


## End(Not run)
\end{ExampleCode}
\end{Examples}
\inputencoding{utf8}
\HeaderA{rev\_coding}{Reverse Coding for Likert Scale Items}{rev.Rul.coding}
%
\begin{Description}
This function performs reverse coding on a numeric vector, commonly used in Likert scale data.
The function assumes the Likert scale ranges from 1 to (n - 1), where `n` is typically 6 for a 5-point scale.
\end{Description}
%
\begin{Usage}
\begin{verbatim}
rev_coding(vec, n = 6, check = FALSE)
\end{verbatim}
\end{Usage}
%
\begin{Arguments}
\begin{ldescription}
\item[\code{vec}] Numeric vector to be reverse-coded.

\item[\code{n}] Integer value indicating the maximum Likert scale value plus 1 (e.g., use `n = 6` for 1 to 5 scale).

\item[\code{check}] Logical; if TRUE, returns a data frame showing the original and reverse-coded values side-by-side.
\end{ldescription}
\end{Arguments}
%
\begin{Value}
A reverse-coded vector or a data frame if `check = TRUE`. Returns an error message if input is not a vector
or values exceed the expected maximum.
\end{Value}
%
\begin{Examples}
\begin{ExampleCode}
# Example 1: Basic reverse coding (Likert 1~5, n=6)
x <- c(1, 2, 3, 4, 5)
rev_coding(x)

# Example 2: View original and reverse-coded values
rev_coding(x, check = TRUE)

# Example 3: 7-point Likert scale (1~7), use n = 8
x7 <- c(1, 3, 5, 7)
rev_coding(x7, n = 8)

# Example 4: Using data from mtcars - reverse code 'gear' variable
rev_coding(mtcars$gear, n = 6)

# Example 5: Check = TRUE for inspection
rev_coding(mtcars$gear, check = TRUE)

# Example 6: Using 'cyl' column from mtcars (values 4, 6, 8), Likert scale mismatch
rev_coding(mtcars$cyl, n = 6)

# Example 7: Using 'Sepal.Length' from iris (continuous values), inappropriate for Likert
rev_coding(iris$Sepal.Length, n = 6)

# Example 8: Using ToothGrowth dataset - 'dose' column
rev_coding(ToothGrowth$dose, n = 6)

# Example 9: Inputting a data frame (error case)
rev_coding(mtcars)

# Example 10: Reverse coding with missing values
x_miss <- c(1, 2, NA, 4, 5)
rev_coding(x_miss, n = 6)

# Example 11: Reverse coding with missing values and check = TRUE
rev_coding(x_miss, n = 6, check = TRUE)
\end{ExampleCode}
\end{Examples}
\inputencoding{utf8}
\HeaderA{rev\_coding\_extract\_only}{Extract only reverse-coded variables from a data.frame}{rev.Rul.coding.Rul.extract.Rul.only}
%
\begin{Description}
Extract only reverse-coded variables from a data.frame
\end{Description}
%
\begin{Usage}
\begin{verbatim}
rev_coding_extract_only(data, drop_attr = FALSE)
\end{verbatim}
\end{Usage}
%
\begin{Arguments}
\begin{ldescription}
\item[\code{data}] A data.frame of class 'revcoded\_df'

\item[\code{drop\_attr}] Logical. If TRUE, removes attributes and classes (default = FALSE)
\end{ldescription}
\end{Arguments}
%
\begin{Value}
A data.frame containing only reverse-coded variables
\end{Value}
%
\begin{Examples}
\begin{ExampleCode}
df <- data.frame(A = 1:5, B = 5:1)
res <- rev_coding_rep(df, 6, A, B)
rev_coding_extract_only(res)
\end{ExampleCode}
\end{Examples}
\inputencoding{utf8}
\HeaderA{rev\_coding\_map}{Create a mapping table between original and reverse-coded variables}{rev.Rul.coding.Rul.map}
%
\begin{Description}
Create a mapping table between original and reverse-coded variables
\end{Description}
%
\begin{Usage}
\begin{verbatim}
rev_coding_map(data)
\end{verbatim}
\end{Usage}
%
\begin{Arguments}
\begin{ldescription}
\item[\code{data}] A data.frame of class 'revcoded\_df'
\end{ldescription}
\end{Arguments}
%
\begin{Value}
A data.frame with original and reverse-coded variable name pairs
\end{Value}
%
\begin{Examples}
\begin{ExampleCode}
df <- data.frame(A = 1:5, B = 5:1)
res <- rev_coding_rep(df, 6, A, B)
rev_coding_map(res)
\end{ExampleCode}
\end{Examples}
\inputencoding{utf8}
\HeaderA{rev\_coding\_plot\_var}{Plot reverse-coded variable against original values}{rev.Rul.coding.Rul.plot.Rul.var}
%
\begin{Description}
Displays a comparison between the original and reverse-coded values.
Works only when original variable is still available in parent environment or in the given data.
\end{Description}
%
\begin{Usage}
\begin{verbatim}
rev_coding_plot_var(x, data = NULL, ...)
\end{verbatim}
\end{Usage}
%
\begin{Arguments}
\begin{ldescription}
\item[\code{x}] A reverse-coded vector with class 'reverse\_var'

\item[\code{data}] Optional: original data.frame containing the original variable

\item[\code{...}] Additional graphical arguments passed to plot
\end{ldescription}
\end{Arguments}
%
\begin{Value}
A scatter plot of original vs. reverse-coded values
\end{Value}
%
\begin{Examples}
\begin{ExampleCode}
## Not run: 
df <- data.frame(A = 1:5)
res <- rev_coding_rep(df, 6, A)
rev_coding_plot_var(res$A_r, data = df)

## End(Not run)
\end{ExampleCode}
\end{Examples}
\inputencoding{utf8}
\HeaderA{rev\_coding\_rep}{Reverse Coding for Multiple Variables with S3 Methods}{rev.Rul.coding.Rul.rep}
%
\begin{Description}
Reverse-codes specified variables in a data frame by subtracting values
from a given maximum scale value. Supports symbol and character input
for variable names. Adds S3 classes for customized print and summary behavior.
\end{Description}
%
\begin{Usage}
\begin{verbatim}
rev_coding_rep(data, n = 8, ..., postfix = "_r", replace = FALSE)
\end{verbatim}
\end{Usage}
%
\begin{Arguments}
\begin{ldescription}
\item[\code{data}] A data.frame containing variables to be reverse coded.

\item[\code{n}] A numeric scalar or numeric vector specifying maximum scale values.
If length 1, the value is recycled for all variables.

\item[\code{...}] Variables to be reverse coded. Can be symbols or character strings.

\item[\code{postfix}] Character string appended to new variable names.
Default is "\_r".

\item[\code{replace}] Logical. If TRUE, replaces original variables.
\end{ldescription}
\end{Arguments}
%
\begin{Value}
A data.frame of class 'revcoded\_df'.
Reverse-coded variables have class 'reverse\_var'.
\end{Value}
%
\begin{Examples}
\begin{ExampleCode}
## Not run: 
############################################
# 0. 기본 데이터 생성
############################################
df <- data.frame(
  A = c(1, 2, 3, 4, 5),
  B = c(5, 4, 3, 2, 1),
  C = c(2, 3, 3, 4, 5)
)

############################################
# 1. 기본 사용: 하나의 n 값으로 여러 변수 역코딩
############################################
res1 <- rev_coding_rep(df, n = 6, A, B)
print(res1)
summary(res1$A_r)

############################################
# 2. 문자열로 변수 이름 전달
############################################
res2 <- rev_coding_rep(df, n = 6, "A", "B")
print(res2)

############################################
# 3. replace = TRUE : 기존 변수 덮어쓰기
############################################
res3 <- rev_coding_rep(df, n = 6, A, B, replace = TRUE)
print(res3)
summary(res3$A)

############################################
# 4. 변수별로 다른 n 값 적용
############################################
res4 <- rev_coding_rep(df, n = c(6, 5, 7), A, B, C)
print(res4)

############################################
# 5. postfix 변경
############################################
res5 <- rev_coding_rep(df, n = 6, A, B, postfix = "_inv")
names(res5)

############################################
# 6. NA 값이 포함된 데이터
############################################
df_na <- data.frame(
  Q1 = c(1, NA, 3, 5),
  Q2 = c(5, 4, NA, 2)
)
res6 <- rev_coding_rep(df_na, n = 6, Q1, Q2)
print(res6)
summary(res6$Q1_r)

############################################
# 7. mtcars 예제 (범주형처럼 쓰이는 수치형 변수)
############################################
# cyl = {4,6,8} → min + max = 12
res7 <- rev_coding_rep(
  mtcars[, c("cyl", "gear", "carb")],
  n = c(12, 6, 9),
  "cyl", "gear", "carb"
)
head(res7)

############################################
# 8. tidyverse 파이프라인에서 사용
############################################
library(dplyr)
df %>%
  rev_coding_rep(n = 6, A, B) %>%
  select(A_r, B_r) %>%
  summary()

############################################
# 9. reverse-coded 변수만 자동 추출
############################################
get_reverse_vars <- function(x) {
  names(x)[sapply(x, function(col) inherits(col, "reverse_var"))]
}
res9 <- rev_coding_rep(df, n = 6, A, B)
get_reverse_vars(res9)

############################################
# 10. 심리척도 예제 (Likert scale)
############################################
likert <- data.frame(
  stress1 = c(1,2,3,4,5),
  stress2 = c(2,3,4,4,5),
  stress3 = c(5,4,3,2,1)
)
res10 <- rev_coding_rep(likert, n = 6, stress3)
print(res10)

############################################
# 11. 실수 방지 예제 (factor 주의)
############################################
df_factor <- data.frame(Q1 = factor(c(1,2,3)))
# rev_coding_rep(df_factor, n = 6, Q1)  # 권장하지 않음

## End(Not run)
\end{ExampleCode}
\end{Examples}
\inputencoding{utf8}
\HeaderA{rev\_coding\_restore\_original}{Remove reverse-coded variables, keeping only original variables}{rev.Rul.coding.Rul.restore.Rul.original}
%
\begin{Description}
Remove reverse-coded variables, keeping only original variables
\end{Description}
%
\begin{Usage}
\begin{verbatim}
rev_coding_restore_original(data)
\end{verbatim}
\end{Usage}
%
\begin{Arguments}
\begin{ldescription}
\item[\code{data}] A data.frame of class 'revcoded\_df'
\end{ldescription}
\end{Arguments}
%
\begin{Value}
A data.frame without reverse-coded variables
\end{Value}
%
\begin{Examples}
\begin{ExampleCode}
df <- data.frame(A = 1:5, B = 5:1)
res <- rev_coding_rep(df, 6, A, B)
rev_coding_restore_original(res)
\end{ExampleCode}
\end{Examples}
\inputencoding{utf8}
\HeaderA{rev\_coding\_validate\_scale}{Validate reverse coding scale (n): scale 유효성 검사}{rev.Rul.coding.Rul.validate.Rul.scale}
%
\begin{Description}
Checks if the provided max scale (n) is valid for each variable.
A valid n is usually: min(x) + max(x)
\end{Description}
%
\begin{Usage}
\begin{verbatim}
rev_coding_validate_scale(data, vars, n, strict = FALSE)
\end{verbatim}
\end{Usage}
%
\begin{Arguments}
\begin{ldescription}
\item[\code{data}] A data.frame

\item[\code{vars}] Character vector or symbols: variable names

\item[\code{n}] Vector of expected scale maxima (same length as vars or length 1)

\item[\code{strict}] Logical. If TRUE, gives warning when mismatch found
\end{ldescription}
\end{Arguments}
%
\begin{Value}
A data.frame with columns: var, min, max, expected\_n, provided\_n, match
\end{Value}
%
\begin{Examples}
\begin{ExampleCode}
df <- data.frame(A = 1:5, B = 5:1)
rev_coding_validate_scale(df, c("A", "B"), n = 6)
rev_coding_validate_scale(df, c("A", "B"), n = c(6, 6))
\end{ExampleCode}
\end{Examples}
\inputencoding{utf8}
\HeaderA{Rload}{Rload}{Rload}
%
\begin{Description}
Rload
\end{Description}
%
\begin{Usage}
\begin{verbatim}
Rload(file, verbose = TRUE)
\end{verbatim}
\end{Usage}
%
\begin{Arguments}
\begin{ldescription}
\item[\code{file}] position

\item[\code{verbose}] default TRUR
\end{ldescription}
\end{Arguments}
%
\begin{Value}
datafile
\end{Value}
%
\begin{Examples}
\begin{ExampleCode}

## Not run: 
#'
#'
#'
Rload("pwj2.RData", verbose = TRUE)
pwj2 %>%head()
Rload(file="F:/Rwork/02_OutAna/Parkwonjun_yonsei_2024/pwj2.RData")

## End(Not run)


\end{ExampleCode}
\end{Examples}
\inputencoding{utf8}
\HeaderA{Round}{Find and round numeric variables only}{Round}
%
\begin{Description}
Find and round numeric variables only
\end{Description}
%
\begin{Usage}
\begin{verbatim}
Round(data, digits = 3, type = "tibble", exclude = NULL)
\end{verbatim}
\end{Usage}
%
\begin{Arguments}
\begin{ldescription}
\item[\code{data}] data.frame

\item[\code{digits}] default 2

\item[\code{type}] data type tibble or data.frame , matrix

\item[\code{exclude}] exclude variable not adapted
\end{ldescription}
\end{Arguments}
%
\begin{Value}
rounding data
\end{Value}
%
\begin{Examples}
\begin{ExampleCode}
## Not run: 

# vector
Round(12.345654)

##data.frame
mtcars%>% Round(1)

 ##
dff = data.frame(A = c(1.23456, 2.34567, 3.45678),
                 B = c(4.56789, 5.67890, 6.78901),
                C = c("a", "b", "c"))
 ## rounding 1
 dff%>%Round(1)

#exclue variable B
 dff%>%Round(1, exclude ="B")


## End(Not run)
\end{ExampleCode}
\end{Examples}
\inputencoding{utf8}
\HeaderA{round2}{numeric round}{round2}
%
\begin{Description}
numeric round
\end{Description}
%
\begin{Usage}
\begin{verbatim}
round2(data, digits = 2)
\end{verbatim}
\end{Usage}
%
\begin{Arguments}
\begin{ldescription}
\item[\code{data}] data

\item[\code{digits}] 2
\end{ldescription}
\end{Arguments}
%
\begin{Value}
data.frame
\end{Value}
%
\begin{Examples}
\begin{ExampleCode}

## Not run: 

dataset = data.frame(
var = c("mpg", "cyl", "disp", "hp",
"drat", "wt", "qsec", "vs", "am",
"gear", "carb"),
N = c(32, 32, 32, 32, 32, 32, 32,
 32, 32, 32, 32),
MEAN = c(20.090625, 6.1875, 230.721875,
 146.6875, 3.5965625, 3.21725, 17.84875,
 0.4375, 0.40625, 3.6875, 2.8125),
SD = c(6.0269480520891, 1.78592164694654,
 123.938693831382, 68.5628684893206,
  0.534678736070971, 0.978457442989697,
  1.78694323609684, 0.504016128774185,
   0.498990917235846, 0.737804065256947,
    1.61519997763185),
MIN = c(10.4, 4, 71.1, 52, 2.76,
1.513, 14.5, 0, 0, 3, 1),
MAX = c(33.9, 8, 472, 335, 4.93,
5.424, 22.9, 1, 1, 5, 8),
Skew = c(0.610655017573288, -0.174611913930169,
 0.381657034147599, 0.726023656361273,
 0.265903904603492, 0.423146464177225,
 0.369045278509436, 0.240257688369044,
  0.364015894329692, 0.52885446231085, 1.05087376449851),
Kurt = c(-0.372766029820891, -1.76211977398177,
 -1.2072119466699, -0.135551121097421,
 -0.714700615734489, -0.0227107528393127,
 0.335114219571393, -2.00193762400794,
 -1.92474142839069, -1.06975067515432,
  1.25704307347743))
dataset
dataset%>% round2()



## End(Not run)


\end{ExampleCode}
\end{Examples}
\inputencoding{utf8}
\HeaderA{round\_col}{round\_col, 한개의 열만 round처리}{round.Rul.col}
%
\begin{Description}
round\_col, 한개의 열만 round처리
\end{Description}
%
\begin{Usage}
\begin{verbatim}
round_col(df, col = ncol(df), digits = 3, scientific = FALSE)
\end{verbatim}
\end{Usage}
%
\begin{Arguments}
\begin{ldescription}
\item[\code{df}] df

\item[\code{col}] select col , default last column

\item[\code{digits}] default 3
\end{ldescription}
\end{Arguments}
%
\begin{Value}
data
\end{Value}
\inputencoding{utf8}
\HeaderA{row2col}{rownames to column}{row2col}
%
\begin{Description}
rownames to column
\end{Description}
%
\begin{Usage}
\begin{verbatim}
row2col(data, colname = "rownames")
\end{verbatim}
\end{Usage}
%
\begin{Arguments}
\begin{ldescription}
\item[\code{data}] data.fram

\item[\code{colname}] new column default rownames
\end{ldescription}
\end{Arguments}
%
\begin{Value}
data
\end{Value}
%
\begin{Examples}
\begin{ExampleCode}

## Not run: 
mtcars %>%row2col("car")

## End(Not run)

\end{ExampleCode}
\end{Examples}
\inputencoding{utf8}
\HeaderA{row2cols}{rownames to columns}{row2cols}
%
\begin{Description}
rownames to columns
\end{Description}
%
\begin{Usage}
\begin{verbatim}
row2cols(data, colname = "rows")
\end{verbatim}
\end{Usage}
%
\begin{Arguments}
\begin{ldescription}
\item[\code{data}] data.fram

\item[\code{colname}] new column default rownames
\end{ldescription}
\end{Arguments}
%
\begin{Value}
data
\end{Value}
%
\begin{Examples}
\begin{ExampleCode}

## Not run: 
mtcars %>%row2cols("car")

## End(Not run)

\end{ExampleCode}
\end{Examples}
\inputencoding{utf8}
\HeaderA{rowdata2col}{Select from ROW to make it a column name}{rowdata2col}
%
\begin{Description}
Select from ROW to make it a column name
\end{Description}
%
\begin{Usage}
\begin{verbatim}
rowdata2col(data, row)
\end{verbatim}
\end{Usage}
%
\begin{Arguments}
\begin{ldescription}
\item[\code{data}] data

\item[\code{row}] select row
\end{ldescription}
\end{Arguments}
%
\begin{Value}
data
\end{Value}
%
\begin{Examples}
\begin{ExampleCode}


## Not run: 

#'
# 예제 데이터 생성
data <- data.frame(
  X1 = c(1.000000, 13.437500, 5.464659, 16.000000),
  X2 = c(2.000000, 11.800000, 4.813671, 15.000000),
  X3 = c(3.000000, 12.473684, 4.032761, 19.000000)
)
rownames(data) <- c("var", "Mean", "SD", "N")
# 함수 사용 예시
result <- rowdata2col(data, "var")
print(result)



## End(Not run)
\end{ExampleCode}
\end{Examples}
\inputencoding{utf8}
\HeaderA{rownames2colnames}{rownames2colnames}{rownames2colnames}
%
\begin{Description}
rownames2colnames
\end{Description}
%
\begin{Usage}
\begin{verbatim}
rownames2colnames(df, reverse = FALSE)
\end{verbatim}
\end{Usage}
%
\begin{Arguments}
\begin{ldescription}
\item[\code{df}] df

\item[\code{reverse}] reverse = FALSE
\end{ldescription}
\end{Arguments}
%
\begin{Value}
df
\end{Value}
%
\begin{Examples}
\begin{ExampleCode}
## Not run: 
#'

# 예시 데이터프레임
df <- data.frame(a = 1:3, b = 4:6)
rownames(df) <- c("row1", "row2", "row3")

# rownames를 colnames로 변환
new_df <- rownames2colnames(df)
print(new_df)

# colnames를 rownames로 변환
df_with_colnames_as_rownames <- rownames2colnames(new_df, reverse = TRUE)
print(df_with_colnames_as_rownames)

## End(Not run)
\end{ExampleCode}
\end{Examples}
\inputencoding{utf8}
\HeaderA{rows\_to\_colnames}{Convert a Row into Column Names of a Data Frame}{rows.Rul.to.Rul.colnames}
%
\begin{Description}
This function sets the values of a specified row in a data frame as the new column names.
The row used for column names is then removed from the data.
\end{Description}
%
\begin{Usage}
\begin{verbatim}
rows_to_colnames(df, row = 1)
\end{verbatim}
\end{Usage}
%
\begin{Arguments}
\begin{ldescription}
\item[\code{df}] A data frame. The input data where one row will be promoted to column names.

\item[\code{row}] An integer indicating the row number to use as column names. Defaults to 1.
\end{ldescription}
\end{Arguments}
%
\begin{Value}
A data frame with updated column names and the specified row removed.
\end{Value}
%
\begin{Examples}
\begin{ExampleCode}
# Create a simple data frame
df <- data.frame(
  V1 = c("Name", "Alice", "Bob"),
  V2 = c("Age", 25, 30),
  V3 = c("Gender", "F", "M"),
  stringsAsFactors = FALSE
)

# Convert the first row to column names
new_df <- rows_to_colnames(df)
print(new_df)

# Use the second row instead
df2 <- data.frame(
  A = c("Header1", "X", "Y"),
  B = c("Header2", 1, 2),
  stringsAsFactors = FALSE
)
rows_to_colnames(df2, row = 2)
\end{ExampleCode}
\end{Examples}
\inputencoding{utf8}
\HeaderA{row\_filter}{select rows filter function}{row.Rul.filter}
%
\begin{Description}
select rows filter function
\end{Description}
%
\begin{Usage}
\begin{verbatim}
row_filter(data, ...)
\end{verbatim}
\end{Usage}
%
\begin{Arguments}
\begin{ldescription}
\item[\code{data}] data

\item[\code{...}] row condition
\end{ldescription}
\end{Arguments}
%
\begin{Value}
data
\end{Value}
%
\begin{Examples}
\begin{ExampleCode}

## Not run: 
#'
mtcars %>% row_filter(cyl == 6 )
mtcars %>% row_filter(vs == 1)
mtcars %>% row_filter(cyl == 6 | vs == 1)
mtcars %>% row_filter(cyl == 6 & vs == 1)

## End(Not run)



\end{ExampleCode}
\end{Examples}
\inputencoding{utf8}
\HeaderA{Rsave}{RData fast save}{Rsave}
%
\begin{Description}
RData fast save
\end{Description}
%
\begin{Usage}
\begin{verbatim}
Rsave(df, file = "")
\end{verbatim}
\end{Usage}
%
\begin{Arguments}
\begin{ldescription}
\item[\code{df}] data to save

\item[\code{file}] filename
\end{ldescription}
\end{Arguments}
%
\begin{Value}
data
\end{Value}
%
\begin{Examples}
\begin{ExampleCode}
## Not run: 

Rsave(Kge_jbs)           # 기본 파일 이름: Kge_jbs.RData
Rsave(Kge_jbs, file = "my_custom_name")  # 사용자 지정 파일 이름: my_custom_name.RData

## End(Not run)
\end{ExampleCode}
\end{Examples}
\inputencoding{utf8}
\HeaderA{Rsave\_xlsx}{Save Data to Excel File}{Rsave.Rul.xlsx}
%
\begin{Description}
This function saves a given data frame to an Excel file in the current working directory.
The file is named after the variable name of the input data.
\end{Description}
%
\begin{Usage}
\begin{verbatim}
Rsave_xlsx(data)
\end{verbatim}
\end{Usage}
%
\begin{Arguments}
\begin{ldescription}
\item[\code{data}] A data frame to be saved as an Excel file.
\end{ldescription}
\end{Arguments}
%
\begin{Value}
None. The function saves the data to a file and prints the save path and a reload command.
\end{Value}
%
\begin{Examples}
\begin{ExampleCode}
## Not run: 
# Example: Save a data frame to an Excel file
df <- data.frame(
  Name = c("Alice", "Bob", "Charlie"),
  Age = c(25, 30, 35)
)
Rsave_xlsx(df)

# The file "df.xlsx" will be saved in the current working directory.
# Reload the saved file with:
# df <- openxlsx::read.xlsx(file = "df.xlsx")

## End(Not run)
\end{ExampleCode}
\end{Examples}
\inputencoding{utf8}
\HeaderA{save\_xlsx}{Save Data to Excel File}{save.Rul.xlsx}
%
\begin{Description}
This function saves a given data frame to an Excel file in the current working directory.
The file is named after the variable name of the input data.
\end{Description}
%
\begin{Usage}
\begin{verbatim}
save_xlsx(data)
\end{verbatim}
\end{Usage}
%
\begin{Arguments}
\begin{ldescription}
\item[\code{data}] A data frame to be saved as an Excel file.
\end{ldescription}
\end{Arguments}
%
\begin{Value}
None. The function saves the data to a file and prints the save path and a reload command.
\end{Value}
%
\begin{Examples}
\begin{ExampleCode}
## Not run: 
# Example: Save a data frame to an Excel file
df <- data.frame(
  Name = c("Alice", "Bob", "Charlie"),
  Age = c(25, 30, 35)
)
Rsave_xlsx(df)

# The file "df.xlsx" will be saved in the current working directory.
# Reload the saved file with:
# df <- openxlsx::read.xlsx(file = "df.xlsx")

## End(Not run)
\end{ExampleCode}
\end{Examples}
\inputencoding{utf8}
\HeaderA{select\_format}{select\_format}{select.Rul.format}
%
\begin{Description}
select\_format
\end{Description}
%
\begin{Usage}
\begin{verbatim}
select_format(data, ..., nsmall = 3, digits = 3)
\end{verbatim}
\end{Usage}
%
\begin{Arguments}
\begin{ldescription}
\item[\code{data}] data

\item[\code{...}] col selection

\item[\code{nsmall}] default 3

\item[\code{digits}] round 3
\end{ldescription}
\end{Arguments}
%
\begin{Value}
data
\end{Value}
%
\begin{Examples}
\begin{ExampleCode}


## Not run: 


# Usage example
data2 <- data.frame(a = rnorm(5), b = rnorm(5), c = rnorm(5), d = rnorm(5), e = rnorm(5))
select_format(data2, 1, 2)
select_format(data2, 1:5)
select_format(data2, a)
select_format(data2, a:e)
select_format(data2, "a")

## End(Not run)
\end{ExampleCode}
\end{Examples}
\inputencoding{utf8}
\HeaderA{SEM}{Title}{SEM}
%
\begin{Description}
Title
\end{Description}
%
\begin{Usage}
\begin{verbatim}
SEM(
  model = NULL,
  data = NULL,
  type = "res",
  ind = TRUE,
  cat = FALSE,
  interact = TRUE,
  auto = TRUE,
  n_name = 1,
  prefix = "a",
  paths_name = TRUE,
  add_model = NULL,
  ordered = NULL,
  sampling.weights = NULL,
  sample.cov = NULL,
  sample.mean = NULL,
  sample.th = NULL,
  sample.nobs = NULL,
  group = NULL,
  cluster = NULL,
  constraints = "",
  WLS.V = NULL,
  NACOV = NULL,
  ov.order = "model",
  ...
)
\end{verbatim}
\end{Usage}
%
\begin{Arguments}
\begin{ldescription}
\item[\code{model}] A description of the user-specified model. Typically, the model is
described using the lavaan model syntax. See model.syntax for more information.
Alternatively, a parameter table (eg. the output of the lavaanify() function)
is also accepted.

\item[\code{data}] An optional data frame containing the observed variables used in the model.
If some variables are declared as ordered factors,
lavaan will treat them as ordinal variables.

\item[\code{type}] res, fit, New\_model, New\_data, imlist

\item[\code{ind}] Automatically Adding Indirect Effects

\item[\code{cat}] output syntax default FALSE

\item[\code{interact}] interaction term extract possinle TRUE

\item[\code{auto}] auto genrate indirect path

\item[\code{n\_name}] indirect name number n\_name =1

\item[\code{prefix}] prefix "a" --"H" change possible

\item[\code{paths\_name}] inde path name default TRUE

\item[\code{add\_model}] add model syntax

\item[\code{ordered}] Character vector. Only used if the data is in a data.frame.
Treat these variables as ordered (ordinal) variables, if they are endogenous
in the model. Importantly, all other variables will be treated as numeric
(unless they are declared as ordered in the data.frame.) Since 0.6-4,
ordered can also be logical. If TRUE, all observed endogenous variables are
\#'    treated as ordered (ordinal). If FALSE, all observed endogenous variables
are considered to be numeric (again, unless they are declared as ordered
in the data.frame.)

\item[\code{sampling.weights}] A variable name in the data frame containing sampling
weight information. Currently only available for non-clustered data.
Depending on the sampling.weights.normalization option, these weights may be
rescaled (or not) so that their sum equals the number of observations
(total or per group).

\item[\code{sample.cov}] Numeric matrix. A sample variance-covariance matrix.
The rownames and/or colnames must contain the observed variable names.
For a multiple group analysis, a list with a variance-covariance matrix
for each group.

\item[\code{sample.mean}] A sample mean vector. For a multiple group analysis,
a list with a mean vector for each group.

\item[\code{sample.th}] Vector of sample-based thresholds.
For a multiple group analysis, a list with a vector of thresholds for each group.

\item[\code{sample.nobs}] Number of observations if the full data frame is missing and
only sample moments are given. For a multiple group analysis, a list or a vector
with the number of observations for each group.

\item[\code{group}] Character. A variable name in the data frame defining
the groups in a multiple group analysis.

\item[\code{cluster}] Character. A (single) variable name in the data frame defining
the clusters in a two-level dataset.

\item[\code{constraints}] Additional (in)equality constraints not yet included
in the model syntax. See model.syntax for more information.

\item[\code{WLS.V}] A user provided weight matrix to be used by estimator "WLS";
if the estimator is "DWLS", only the diagonal of this matrix will be used.
For a multiple group analysis, a list with a weight matrix for each group.
The elements of the weight matrix should be in the following order (if all data is
continuous): first the means (if a meanstructure is involved),
then the lower triangular elements of the covariance matrix including the
diagonal, ordered column by column. In the categorical case: first
the thresholds (including the means for continuous variables),
then the slopes (if any), the variances of continuous variables (if any),
and finally the lower triangular elements of the correlation/covariance
matrix excluding the diagonal, ordered column by column.

\item[\code{NACOV}] A user provided matrix containing the elements of (N times)
the asymptotic variance-covariance matrix of the sample statistics.
For a multiple group analysis, a list with an asymptotic variance-covariance
matrix for each group. See the WLS.V argument for information about
the order of the elements.

\item[\code{ov.order}] Character. If "model" (the default), the order of the observed
variable names (as reflected for example in the output of lavNames())
is determined by the model syntax. If "data", the order is determined
by the data (either the full data.frame or the sample (co)variance matrix).
If the WLS.V and/or NACOV matrices are provided,
this argument is currently set to "data".

\item[\code{...}] Many more additional options can be defined, using 'name = value'. See lavOptions for a complete list.
\end{ldescription}
\end{Arguments}
%
\begin{Value}
data
\end{Value}
%
\begin{Examples}
\begin{ExampleCode}

## Not run: 
#'
## The industrialization and Political Democracy Example
## Bollen (1989), page 332
model <- '
  # latent variable definitions
     ind60 =~ x1 + x2 + x3
     dem60 =~ y1 + a*y2 + b*y3 + c*y4
     dem65 =~ y5 + a*y6 + b*y7 + c*y8

  # regressions
    dem60 ~ ind60
    dem65 ~ ind60 + dem60

  # residual correlations
    y1 ~~ y5
    y2 ~~ y4 + y6
    y3 ~~ y7
    y4 ~~ y8
    y6 ~~ y8
'

fit <- sem(model, data = PoliticalDemocracy)
summary(fit, fit.measures = TRUE)
#'
#'

# # Example usage:
data(MarshWenHau)
str(MarshWenHau)

mwh = MarshWenHau

model2 <- "
f1 =~ x1 + x2 + x3
f1Xf2 =~ f1:f2
f2 =~ x4 + x5 + x6
f3 =~ y1 + y2 + y3

f3 ~ f1 + f2
f3 ~ f1Xf2

f1Xf2 ~~ 0*f1
f1Xf2 ~~ 0*f2
"

# interaction model
intsem =  SEM(model2, data= MarshWenHau)
summary(intsem)
intsem %>%lav_semPaths2(color=list(lat="darkred",man="gold"))


##2nd---------
#base model
lav_extract_mm(model2) %>%cat()
#interaction item list
lav_extract_imlist(model3)
#data
lav_latentProd(mwh, model3)%>%head()

#extracted interaction term
lav_extract_int(model2, lav_extract_imlist(model3))

# or interaction items ; imlist
imlist <- list(
  .f1Xf2 = c("f1", "f2"),
  .i_f1_f4 = c("f1", "f4")
)
# names(imlist)
lav_extract_int(model2, imlist)
#direct
lav_extract_int(model2, imlist= list(.f1Xf2 = c("f1", "f2"),
                                     .i_f1_f4 = c("f1", "f4")))
#estimates model
intsem =  SEM(lav_new_model(model2, lav_extract_imlist(model2)),
              data= lav_latentProd(mwh, model2))

#'

## End(Not run)


\end{ExampleCode}
\end{Examples}
\inputencoding{utf8}
\HeaderA{semPaths2}{semPlot::semPaths wrapper with customization for grouped or non-grouped path diagrams}{semPaths2}
%
\begin{Description}
A convenient wrapper around `semPlot::semPaths` with additional options for layout, labels,
grouping, and significance marking using `semptools::mark\_sig`.
\end{Description}
%
\begin{Usage}
\begin{verbatim}
semPaths2(
  lav_obj,
  whatLabels = "est",
  rotation = 2,
  nDigits = 2,
  residScale = 18,
  mar = c(2, 8, 3, 8),
  edge.label.cex = 1,
  edge.label.position = 0.5,
  sizeMan = 6,
  sizeMan2 = 4,
  sizeLat = 8,
  sizeLat2 = 6,
  border.width = 2,
  edge.width = 1.5,
  style = "lisrel",
  sig = TRUE,
  asize = 1.5,
  shapeLat = "circle",
  residuals = FALSE,
  curve = 1,
  layout = "tree",
  exoVar = TRUE,
  exoCov = TRUE,
  structual = FALSE,
  nodeLabels = NULL,
  groups = FALSE
)
\end{verbatim}
\end{Usage}
%
\begin{Arguments}
\begin{ldescription}
\item[\code{lav\_obj}] A fitted lavaan object.

\item[\code{whatLabels}] Type of labels to display on edges (e.g., `"est"` for estimates).

\item[\code{rotation}] Rotation angle for layout.

\item[\code{nDigits}] Number of digits to round estimates.

\item[\code{residScale}] Scaling factor for residuals.

\item[\code{mar}] Margins of the plot.

\item[\code{edge.label.cex}] Character expansion factor for edge labels.

\item[\code{edge.label.position}] Position of edge labels along the edge.

\item[\code{sizeMan}] Size of manifest variables.

\item[\code{sizeMan2}] Size of exogenous manifest variables.

\item[\code{sizeLat}] Size of latent variables.

\item[\code{sizeLat2}] Size of exogenous latent variables.

\item[\code{border.width}] Border width of nodes.

\item[\code{edge.width}] Line width of edges.

\item[\code{style}] Style of path diagram (e.g., `"lisrel"`).

\item[\code{sig}] Logical. If TRUE, significant paths are marked.

\item[\code{asize}] Arrow size.

\item[\code{shapeLat}] Shape of latent variables (e.g., `"circle"`).

\item[\code{residuals}] Logical. If TRUE, show residuals.

\item[\code{curve}] Curvature of edges.

\item[\code{layout}] Layout of the graph (e.g., `"tree"`, `"circle"`).

\item[\code{exoVar}] Logical. If TRUE, include exogenous variables.

\item[\code{exoCov}] Logical. If TRUE, include covariances of exogenous variables.

\item[\code{structual}] Logical. If TRUE, plot structural model only.

\item[\code{nodeLabels}] Optional custom node labels.

\item[\code{groups}] Logical. If TRUE, color nodes by latent groupings.
\end{ldescription}
\end{Arguments}
%
\begin{Value}
A path diagram object, optionally plotted if `sig = TRUE`.
\end{Value}
%
\begin{Examples}
\begin{ExampleCode}
## Not run: 
library(lavaan)
library(semPlot)
library(semptools)

model <- '
  visual  =~ x1 + x2 + x3
  textual =~ x4 + x5 + x6
  speed   =~ x7 + x8 + x9
'
fit <- cfa(model, data = HolzingerSwineford1939)

semPaths2(fit)
semPaths2(fit, sig = FALSE)

## End(Not run)
\end{ExampleCode}
\end{Examples}
\inputencoding{utf8}
\HeaderA{sem\_apa}{SEM direct effect interpretation}{sem.Rul.apa}
%
\begin{Description}
SEM direct effect interpretation
\end{Description}
%
\begin{Usage}
\begin{verbatim}
sem_apa(data, md = FALSE, caption = "Table: ", print = TRUE)
\end{verbatim}
\end{Usage}
%
\begin{Arguments}
\begin{ldescription}
\item[\code{data}] DE\_effect obj

\item[\code{md}] markdown\_table output

\item[\code{caption}] caption= title

\item[\code{print}] print TRUE
\end{ldescription}
\end{Arguments}
%
\begin{Value}
report apa
\end{Value}
%
\begin{Examples}
\begin{ExampleCode}

## Not run: 

 #library(lavaan) # only needed once per session
 models <- '
   # measurement model
     ind60 =~ x1 + x2 + x3
     dem60 =~ y1 + y2 + y3 + y4
     dem65 =~ y5 + y6 + y7 + y8
 # regressions
     dem60 ~ ind60
     dem65 ~ ind60 + dem60
   # residual correlations
     y1 ~~ y5
     y2 ~~ y4 + y6
     y3 ~~ y7
     y4 ~~ y8
    y6 ~~ y8
 '
 fits <- sem(models, data=PoliticalDemocracy)
 summary(fits, standardized=TRUE)

 ##Direct effect
 DE_effect(fits)
 DE_effect(fits)%>% sem_apa()

example(sem)
DE_effect(fit)
DE_effect(fit) %>% sem_apa()


## End(Not run)

\end{ExampleCode}
\end{Examples}
\inputencoding{utf8}
\HeaderA{sem\_apa\_ie}{sem\_apa\_ie}{sem.Rul.apa.Rul.ie}
%
\begin{Description}
sem\_apa\_ie
\end{Description}
%
\begin{Usage}
\begin{verbatim}
sem_apa_ie(data, msg = "IE", md = FALSE, est = 3, std = 4, sig = 7, p = 6)
\end{verbatim}
\end{Usage}
%
\begin{Arguments}
\begin{ldescription}
\item[\code{data}] data

\item[\code{msg}] msg="IE"

\item[\code{md}] md = FALSE

\item[\code{est}] est=3

\item[\code{std}] std=4

\item[\code{sig}] sig = 7

\item[\code{p}] p=6
\end{ldescription}
\end{Arguments}
%
\begin{Value}
ie report
\end{Value}
%
\begin{Examples}
\begin{ExampleCode}

## Not run: 
example(sem)
IE_effect(fit)
IE_effect(fit) %>% sem_apa_ie()
sem_effect_ci(satis_sem1, ci=FALSE)
sem_effect_ci(satis_sem1, ci=FALSE)%>%sem_apa_ie()



## End(Not run)
\end{ExampleCode}
\end{Examples}
\inputencoding{utf8}
\HeaderA{sem\_apa\_ie\_CI}{sem\_apa\_ie\_CI}{sem.Rul.apa.Rul.ie.Rul.CI}
%
\begin{Description}
sem\_apa\_ie\_CI
\end{Description}
%
\begin{Usage}
\begin{verbatim}
sem_apa_ie_CI(
  data,
  caption = "매개효과 분석 결과 ",
  md = FALSE,
  est = 2,
  std = 3,
  sig = 6,
  CI = 7,
  Z = FALSE,
  tot_filter = "총",
  ind_filter = "IE"
)
\end{verbatim}
\end{Usage}
%
\begin{Arguments}
\begin{ldescription}
\item[\code{data}] lavaan object

\item[\code{caption}] "매개효과 분석 결과",

\item[\code{md}] markdown

\item[\code{est}] est sel 2

\item[\code{std}] std sel 3

\item[\code{sig}] sig sel 6

\item[\code{CI}] CI seledct 7

\item[\code{Z}] Z=FALSE

\item[\code{tot\_filter}] tot\_filter="총",

\item[\code{ind\_filter}] ind\_filter="IE"
\end{ldescription}
\end{Arguments}
%
\begin{Value}
table
\end{Value}
%
\begin{Examples}
\begin{ExampleCode}

## Not run: 

sem_effect_ci(satis_sem1)%>%sem_apa_ie_CI()


## End(Not run)
\end{ExampleCode}
\end{Examples}
\inputencoding{utf8}
\HeaderA{sem\_APA\_path\_table}{Create an APA-style path coefficient table from a lavaan SEM object}{sem.Rul.APA.Rul.path.Rul.table}
%
\begin{Description}
This function extracts structural path coefficients from a fitted
\code{lavaan} SEM object and returns an APA-style summary table
including unstandardized estimates, standard errors, standardized
coefficients, p-values, and significance symbols.
\end{Description}
%
\begin{Usage}
\begin{verbatim}
sem_APA_path_table(fit, only_sig = FALSE)
\end{verbatim}
\end{Usage}
%
\begin{Arguments}
\begin{ldescription}
\item[\code{fit}] A fitted \code{lavaan} SEM object (e.g., from \code{sem()} or \code{cfa()}).

\item[\code{only\_sig}] Logical. If \code{TRUE}, only statistically significant
paths (p < .05) are returned. Default is \code{FALSE}.
\end{ldescription}
\end{Arguments}
%
\begin{Details}
Only regression paths (operator \code{"\textasciitilde{}"}) are included.
Covariances, variances, and intercepts are excluded.
\end{Details}
%
\begin{Value}
A data.frame containing:
\begin{description}

\item[Path] Structural path in the form "Outcome <- Predictor"
\item[B] Unstandardized regression coefficient
\item[SE] Standard error of the coefficient
\item[Beta] Standardized regression coefficient (std.all)
\item[p] Formatted p-value
\item[sig] Significance symbol (* p < .05, ** p < .01, *** p < .001)

\end{description}

\end{Value}
%
\begin{Examples}
\begin{ExampleCode}
## Not run: 
#--------------------------------------------------
# Example 1: Simple mediation model
#--------------------------------------------------
library(lavaan)

model_med <- '
  M ~ a*X
  Y ~ b*M + c*X
'

fit_med <- sem(model_med, data = mydata)

# All structural paths
sem_APA_path_table(fit_med)

# Only significant paths
sem_APA_path_table(fit_med, only_sig = TRUE)

#--------------------------------------------------
# Example 2: Multiple predictors
#--------------------------------------------------
model_multi <- '
  Y ~ X1 + X2 + X3
'

fit_multi <- sem(model_multi, data = mydata)

sem_APA_path_table(fit_multi)

#--------------------------------------------------
# Example 3: Full SEM with latent variables
#--------------------------------------------------
model_sem <- '
  F1 =~ x1 + x2 + x3
  F2 =~ y1 + y2 + y3

  F2 ~ F1
'

fit_sem_full <- sem(model_sem, data = mydata)

sem_APA_path_table(fit_sem_full)
sem_APA_path_table(fit_sem_full, only_sig = TRUE)

## End(Not run)
\end{ExampleCode}
\end{Examples}
\inputencoding{utf8}
\HeaderA{sem\_count\_df}{Functions for calculating degrees of freedom in SEM}{sem.Rul.count.Rul.df}
%
\begin{Description}
Functions for calculating degrees of freedom in SEM
\end{Description}
%
\begin{Usage}
\begin{verbatim}
sem_count_df(
  model,
  type = "fit",
  growth = FALSE,
  sample.nobs = 100,
  whatLabels = "model",
  sig = FALSE,
  plot = FALSE,
  rotation = 2
)
\end{verbatim}
\end{Usage}
%
\begin{Arguments}
\begin{ldescription}
\item[\code{model}] lavaan syntax

\item[\code{type}] fit, all

\item[\code{growth}] grouwth =FALSE

\item[\code{sample.nobs}] 100(default)

\item[\code{whatLabels}] model, If it is an actual measurement, est

\item[\code{sig}] FALSE not outpur significant

\item[\code{plot}] plot FALSE

\item[\code{rotation}] rotation 4
\end{ldescription}
\end{Arguments}
%
\begin{Value}
table and plot
\end{Value}
%
\begin{Examples}
\begin{ExampleCode}

## Not run: 
#'
m1 <- "
y1 =~ x1 + x2 +x5
y2 =~ x3 + x4 +x6
y1 ~~ y2"
sem_count_df(m1)
sem_count_df(m1,"fit")

#그래프를 보고 싶을 때
diagram_model(m1)
#이함수에서 보려면
sem_count_df(m1, plot=TRUE)
#전체를 보여주는 함수
sem_count_df(m1,"all")
sem_count_df(m1,"all", plot=TRUE )
#'
#'
m1a <- "
y1 =~ x1 + x2 +x5
y2 =~ x3 + x4 +x6
y1 ~~ 0*y2"
sem_count_df(m1a)
sem_count_df(m1a,"all")

#이함수에서 보려면
sem_count_df(m1, plot=TRUE)

#simulation data lavaan
sem_count_df(m1,"lav")


## End(Not run)
\end{ExampleCode}
\end{Examples}
\inputencoding{utf8}
\HeaderA{sem\_effect}{SEM direct or indirect  effect result}{sem.Rul.effect}
%
\begin{Description}
SEM direct or indirect  effect result
\end{Description}
%
\begin{Usage}
\begin{verbatim}
sem_effect(
  model_sem,
  effect = "~",
  effect2 = NULL,
  ci = FALSE,
  lower = 8,
  upper = 9,
  unite = FALSE,
  unite_col = "z"
)
\end{verbatim}
\end{Usage}
%
\begin{Arguments}
\begin{ldescription}
\item[\code{model\_sem}] lavaan object

\item[\code{effect}] select op

\item[\code{effect2}] select op 2nd option

\item[\code{ci}] FALSE

\item[\code{lower}] 8th col select

\item[\code{upper}] 9th col select

\item[\code{unite}] TRUE unite p +sig

\item[\code{unite\_col}] select "z"
\end{ldescription}
\end{Arguments}
%
\begin{Value}
repot
\end{Value}
%
\begin{Examples}
\begin{ExampleCode}
## Not run: 

 #library(lavaan) # only needed once per session

models1 <- '
  # measurement model
    ind60 =~ x1 + x2 + x3
    dem60 =~ y1 + y2 + y3 + y4
    dem65 =~ y5 + y6 + y7 + y8
    # regressions
    dem60 ~ ind60
    dem65 ~ ind60 + dem60
      # residual correlations
    y1 ~~ y5
    y2 ~~ y4 + y6
    y3 ~~ y7
    y4 ~~ y8
    y6 ~~ y8
'
models2 ='
  # measurement model
    ind60 =~ x1 + x2 + x3
    dem60 =~ y1 + y2 + y3 + y4
    dem65 =~ y5 + y6 + y7 + y8
     # regressions
    dem60 ~ a1*ind60
    dem65 ~ a2*ind60 + a3*dem60
 #NEw
   ie := a1*a3
   te := a2 +a1*a3
'
fits1 <- sem(models1, data=PoliticalDemocracy)
fits2 <- sem(models2, data=PoliticalDemocracy)
summary(fits1, standardized=TRUE)
summary(fits2, standardized=TRUE)

parameterEstimates(fits1, ci=F, stand=T)
parameterEstimates(fits2, ci=F, stand=T)
sem_effect(fits1)
sem_effect(fits1, effect="=~")
sem_effect(fits1, effect="=~", effect2="~")
sem_effect(fits1, effect="~~")
sem_effect(fits1, effect=":=")

sem_effect(fits2)
sem_effect(fits2, effect="=~")
sem_effect(fits2, effect="~~")
sem_effect(fits2, effect=":=")
sem_effect(fits2, effect=":=", effect2="~")
 sem_effect(fits)%>% interpretation_de()
 # same method
 sem_effect(fits)%>% report_sem_de()


## End(Not run)


\end{ExampleCode}
\end{Examples}
\inputencoding{utf8}
\HeaderA{sem\_effect\_ci}{SEM indirect effect}{sem.Rul.effect.Rul.ci}
%
\begin{Description}
SEM indirect effect
\end{Description}
%
\begin{Usage}
\begin{verbatim}
sem_effect_ci(
  model_sem,
  effect1 = "DE",
  effect2 = "IE",
  effect3 = "TE",
  label = FALSE,
  ci = TRUE
)
\end{verbatim}
\end{Usage}
%
\begin{Arguments}
\begin{ldescription}
\item[\code{model\_sem}] lavaan objgect

\item[\code{effect1}] "DE"

\item[\code{effect2}] "IE"

\item[\code{effect3}] "TE"

\item[\code{label}] label

\item[\code{ci}] confidence interval
\end{ldescription}
\end{Arguments}
%
\begin{Value}
parameterEstimates()
\end{Value}
%
\begin{Examples}
\begin{ExampleCode}

## Not run: 
#library(lavaan) # only needed once per session

models1 <- '
  # measurement model
    ind60 =~ x1 + x2 + x3
    dem60 =~ y1 + y2 + y3 + y4
    dem65 =~ y5 + y6 + y7 + y8
    # regressions
    dem60 ~ ind60
    dem65 ~ ind60 + dem60
      # residual correlations
    y1 ~~ y5
    y2 ~~ y4 + y6
    y3 ~~ y7
    y4 ~~ y8
    y6 ~~ y8

fits1 <- sem(models1, data=PoliticalDemocracy)


models2 ='
  # measurement model
    ind60 =~ x1 + x2 + x3
    dem60 =~ y1 + y2 + y3 + y4
    dem65 =~ y5 + y6 + y7 + y8
     # regressions
    dem60 ~ a1*ind60
    dem65 ~ a2*ind60 + a3*dem60
 #NEw
   ie := a1*a3
   te := a2 +a1*a3
   diff := a1 - a3

fits2 <- sem(models2, data=PoliticalDemocracy)

summary(fits1, standardized=TRUE)
summary(fits2, standardized=TRUE)
fits2 %>% sem_effect_ie()

## End(Not run)

\end{ExampleCode}
\end{Examples}
\inputencoding{utf8}
\HeaderA{sem\_effect\_diff}{Diff effect}{sem.Rul.effect.Rul.diff}
%
\begin{Description}
Diff effect
\end{Description}
%
\begin{Usage}
\begin{verbatim}
sem_effect_diff(
  model_sem,
  effect1 = "Diff",
  effect2 = "diff",
  label = FALSE,
  ci = TRUE
)
\end{verbatim}
\end{Usage}
%
\begin{Arguments}
\begin{ldescription}
\item[\code{model\_sem}] lavaan objgect

\item[\code{effect1}] "Diff"

\item[\code{effect2}] "diff"

\item[\code{label}] label

\item[\code{ci}] confidence interval
\end{ldescription}
\end{Arguments}
%
\begin{Value}
parameterEstimates()
\end{Value}
%
\begin{Examples}
\begin{ExampleCode}

## Not run: 
#library(lavaan) # only needed once per session

models1 <- '
  # measurement model
    ind60 =~ x1 + x2 + x3
    dem60 =~ y1 + y2 + y3 + y4
    dem65 =~ y5 + y6 + y7 + y8
    # regressions
    dem60 ~ ind60
    dem65 ~ ind60 + dem60
      # residual correlations
    y1 ~~ y5
    y2 ~~ y4 + y6
    y3 ~~ y7
    y4 ~~ y8
    y6 ~~ y8

models2 ='
  # measurement model
    ind60 =~ x1 + x2 + x3
    dem60 =~ y1 + y2 + y3 + y4
    dem65 =~ y5 + y6 + y7 + y8
     # regressions
    dem60 ~ a1*ind60
    dem65 ~ a2*ind60 + a3*dem60
 #NEw
   ie := a1*a3
   te := a2 +a1*a3
   Diff := a1 - a3




fits1 <- sem(models1, data=PoliticalDemocracy)
fits2 <- sem(models2, data=PoliticalDemocracy)
summary(fits1, standardized=TRUE)
summary(fits2, standardized=TRUE)
fits2 %>% sem_effect_diff()

## End(Not run)
\end{ExampleCode}
\end{Examples}
\inputencoding{utf8}
\HeaderA{sem\_effect\_ie}{SEM indirect effect}{sem.Rul.effect.Rul.ie}
%
\begin{Description}
SEM indirect effect
\end{Description}
%
\begin{Usage}
\begin{verbatim}
sem_effect_ie(
  model_sem,
  effect1 = "IE",
  effect2 = "IE",
  label = FALSE,
  ci = TRUE
)
\end{verbatim}
\end{Usage}
%
\begin{Arguments}
\begin{ldescription}
\item[\code{model\_sem}] lavaan objgect

\item[\code{effect1}] "IE"

\item[\code{effect2}] "IE"

\item[\code{label}] label

\item[\code{ci}] confidence interval
\end{ldescription}
\end{Arguments}
%
\begin{Value}
parameterEstimates()
\end{Value}
%
\begin{Examples}
\begin{ExampleCode}

## Not run: 
#library(lavaan) # only needed once per session

models1 <- '
  # measurement model
    ind60 =~ x1 + x2 + x3
    dem60 =~ y1 + y2 + y3 + y4
    dem65 =~ y5 + y6 + y7 + y8
    # regressions
    dem60 ~ ind60
    dem65 ~ ind60 + dem60
      # residual correlations
    y1 ~~ y5
    y2 ~~ y4 + y6
    y3 ~~ y7
    y4 ~~ y8
    y6 ~~ y8

models2 ='
  # measurement model
    ind60 =~ x1 + x2 + x3
    dem60 =~ y1 + y2 + y3 + y4
    dem65 =~ y5 + y6 + y7 + y8
     # regressions
    dem60 ~ a1*ind60
    dem65 ~ a2*ind60 + a3*dem60
 #NEw
   ie := a1*a3
   te := a2 +a1*a3
   diff := a1 - a3




fits1 <- sem(models1, data=PoliticalDemocracy)
fits2 <- sem(models2, data=PoliticalDemocracy)
summary(fits1, standardized=TRUE)
summary(fits2, standardized=TRUE)
fits2 %>% sem_effect_ie()

## End(Not run)

\end{ExampleCode}
\end{Examples}
\inputencoding{utf8}
\HeaderA{sem\_effect\_ind}{SEM indirect effect}{sem.Rul.effect.Rul.ind}
%
\begin{Description}
SEM indirect effect
\end{Description}
%
\begin{Usage}
\begin{verbatim}
sem_effect_ind(
  model_sem,
  effect1 = "DE",
  effect2 = "IE",
  digits = 3,
  label = FALSE,
  ci = TRUE
)
\end{verbatim}
\end{Usage}
%
\begin{Arguments}
\begin{ldescription}
\item[\code{model\_sem}] lavaan objgect

\item[\code{effect1}] "IE"

\item[\code{effect2}] "IE"

\item[\code{digits}] digits=3

\item[\code{label}] label

\item[\code{ci}] confidence interval
\end{ldescription}
\end{Arguments}
%
\begin{Value}
parameterEstimates()
\end{Value}
%
\begin{Examples}
\begin{ExampleCode}

## Not run: 
#library(lavaan) # only needed once per session

models1 <- '
  # measurement model
    ind60 =~ x1 + x2 + x3
    dem60 =~ y1 + y2 + y3 + y4
    dem65 =~ y5 + y6 + y7 + y8
    # regressions
    dem60 ~ ind60
    dem65 ~ ind60 + dem60
      # residual correlations
    y1 ~~ y5
    y2 ~~ y4 + y6
    y3 ~~ y7
    y4 ~~ y8
    y6 ~~ y8

fits1 <- sem(models1, data=PoliticalDemocracy)


models2 ='
  # measurement model
    ind60 =~ x1 + x2 + x3
    dem60 =~ y1 + y2 + y3 + y4
    dem65 =~ y5 + y6 + y7 + y8
     # regressions
    dem60 ~ a1*ind60
    dem65 ~ a2*ind60 + a3*dem60
 #NEw
   ie := a1*a3
   te := a2 +a1*a3
   diff := a1 - a3

fits2 <- sem(models2, data=PoliticalDemocracy)

summary(fits1, standardized=TRUE)
summary(fits2, standardized=TRUE)
fits2 %>% sem_effect_ie()

## End(Not run)

\end{ExampleCode}
\end{Examples}
\inputencoding{utf8}
\HeaderA{sem\_effect\_te}{Total effect}{sem.Rul.effect.Rul.te}
%
\begin{Description}
Total effect
\end{Description}
%
\begin{Usage}
\begin{verbatim}
sem_effect_te(
  model_sem,
  effect1 = "TE",
  effect2 = "te",
  label = FALSE,
  ci = TRUE
)
\end{verbatim}
\end{Usage}
%
\begin{Arguments}
\begin{ldescription}
\item[\code{model\_sem}] lavaan objgect

\item[\code{effect1}] "IE"

\item[\code{effect2}] "IE"

\item[\code{label}] label

\item[\code{ci}] confidence interval
\end{ldescription}
\end{Arguments}
%
\begin{Value}
parameterEstimates()
\end{Value}
%
\begin{Examples}
\begin{ExampleCode}

## Not run: 
#library(lavaan) # only needed once per session

models1 <- '
  # measurement model
    ind60 =~ x1 + x2 + x3
    dem60 =~ y1 + y2 + y3 + y4
    dem65 =~ y5 + y6 + y7 + y8
    # regressions
    dem60 ~ ind60
    dem65 ~ ind60 + dem60
      # residual correlations
    y1 ~~ y5
    y2 ~~ y4 + y6
    y3 ~~ y7
    y4 ~~ y8
    y6 ~~ y8

models2 ='
  # measurement model
    ind60 =~ x1 + x2 + x3
    dem60 =~ y1 + y2 + y3 + y4
    dem65 =~ y5 + y6 + y7 + y8
     # regressions
    dem60 ~ a1*ind60
    dem65 ~ a2*ind60 + a3*dem60
 #NEw
   ie := a1*a3
   te := a2 +a1*a3
   diff := a1 - a3




fits1 <- sem(models1, data=PoliticalDemocracy)
fits2 <- sem(models2, data=PoliticalDemocracy)
summary(fits1, standardized=TRUE)
summary(fits2, standardized=TRUE)
fits2 %>% sem_effect_te()

## End(Not run)

\end{ExampleCode}
\end{Examples}
\inputencoding{utf8}
\HeaderA{sem\_fit}{Structural Equation Model fit}{sem.Rul.fit}
%
\begin{Description}
Structural Equation Model fit
\end{Description}
%
\begin{Usage}
\begin{verbatim}
sem_fit(..., nice_table = FALSE)
\end{verbatim}
\end{Usage}
%
\begin{Arguments}
\begin{ldescription}
\item[\code{...}] lavaan model

\item[\code{nice\_table}] markdown table TRUE
\end{ldescription}
\end{Arguments}
%
\begin{Value}
table
\end{Value}
%
\begin{Examples}
\begin{ExampleCode}
## Not run: 
#'
model_test1 <- '
  # latent variable definitions
     ind60 =~ x1 + x2 + x3
     dem60 =~ y1 + a*y2 + b*y3 + c*y4
     dem65 =~ y5 + a*y6 + b*y7 + c*y8

  # regressions
    dem60 ~ ind60
    dem65 ~ ind60 + dem60

  # residual correlations
    y1 ~~ y5
    y2 ~~ y4 + y6
    y3 ~~ y7
    y4 ~~ y8
    y6 ~~ y8
'

model_test2 <- '
  # latent variable definitions
     ind60 =~ x1 + x2 + x3
     dem60 =~ y1 + a*y2 + b*y3 + c*y4
     dem65 =~ y5 + a*y6 + b*y7 + c*y8

  # regressions
    dem60 ~ ind60
    dem65 ~ ind60 + dem60

'

fit1 <- sem(model_test1, data = PoliticalDemocracy)
fit2 <- sem(model_test2, data = PoliticalDemocracy)

sem_fit(fit1, fit2)
sem_fit("model-1" = fit1,"model-2" = fit2)
#'

## End(Not run)
\end{ExampleCode}
\end{Examples}
\inputencoding{utf8}
\HeaderA{sem\_lav\_hypo}{Hypothesis generation functions}{sem.Rul.lav.Rul.hypo}
%
\begin{Description}
Hypothesis generation functions
\end{Description}
%
\begin{Usage}
\begin{verbatim}
sem_lav_hypo(model)
\end{verbatim}
\end{Usage}
%
\begin{Arguments}
\begin{ldescription}
\item[\code{model}] lavaan syntax
\end{ldescription}
\end{Arguments}
%
\begin{Value}
text
\end{Value}
%
\begin{Examples}
\begin{ExampleCode}

## Not run: 
#'
# 예제 모델
model1a <- "
EXPE ~ H1*IMAG
QUAL ~ H2*EXPE
VAL ~  H3*QUAL + H4*EXPE
SAT ~ H5*IMAG + H6*EXPE + H7*QUAL + H8*VAL
LOY ~ H9*SAT + H10*IMAG
"

# 함수 실행
sem_lav_hypo(model1a)
#'

# Example usage with the provided model
model10 <- "
IMAG =~ imag1 + imag2 + imag3 + imag4 + imag5
EXPE =~ expe1 + expe2 + expe3 + expe4 + expe5
QUAL =~ qual1 + qual2 + qual3 + qual4 + qual5
VAL =~ val1 + val2 + val3 + val4
SAT =~ sat1 + sat2 + sat3 + sat4
LOY =~ loy1 + loy2 + loy3 + loy4

EXPE ~ IMAG
QUAL ~ EXPE
VAL ~  QUAL+ EXPE
SAT ~ IMAG + EXPE + QUAL + VAL
LOY ~ SAT + IMAG
"
lav_extract_sm(model10)%>%sem_lav_hypo()
lav_extract_sm(model10, prefix="H")%>%sem_lav_hypo()



#'

## End(Not run)

\end{ExampleCode}
\end{Examples}
\inputencoding{utf8}
\HeaderA{sem\_lav\_ind\_hypo}{Function that takes path data as input and outputs an indirect effect hypothesis.}{sem.Rul.lav.Rul.ind.Rul.hypo}
%
\begin{Description}
Function that takes path data as input and outputs an indirect effect hypothesis.
\end{Description}
%
\begin{Usage}
\begin{verbatim}
sem_lav_ind_hypo(paths)
\end{verbatim}
\end{Usage}
%
\begin{Arguments}
\begin{ldescription}
\item[\code{paths}] indpath
\end{ldescription}
\end{Arguments}
%
\begin{Value}
text
\end{Value}
%
\begin{Examples}
\begin{ExampleCode}

## Not run: 
#'
# 모델 문자열
model1a <- "
EXPE ~ H1*IMAG
QUAL ~ H2*EXPE
VAL ~ H3*QUAL + H4*EXPE
SAT ~ H5*IMAG + H6*EXPE + H7*QUAL + H8*VAL
LOY ~ H9*SAT + H10*IMAG
"

model1 <- "
EXPE ~ IMAG
QUAL ~ EXPE
VAL ~ QUAL + EXPE
SAT ~ IMAG + EXPE + QUAL + VAL
LOY ~ SAT + IMAG
"
lav_extract_path(model1a)%>%
  lav_generate_ind_paths() %>%
  lav_ind_hypo()


## End(Not run)

\end{ExampleCode}
\end{Examples}
\inputencoding{utf8}
\HeaderA{sem\_model}{Drawing Structural Equation Models Using the lavaan Syntax (park JH PhD)}{sem.Rul.model}
%
\begin{Description}
Drawing Structural Equation Models Using the lavaan Syntax (park JH PhD)
\end{Description}
%
\begin{Usage}
\begin{verbatim}
sem_model(
  model,
  residuals = FALSE,
  whatLabels = "model",
  rotation = 2,
  residScale = 12,
  layout = "tree2",
  edge.label.cex = 1,
  edge.label.position = 0.5,
  edge.color = "black",
  sizeMan = 6,
  sizeMan2 = 3,
  sizeLat = 8,
  sizeLat2 = 4,
  style = "lisrel",
  sig = FALSE,
  exoVar = FALSE,
  exoCov = TRUE,
  curve = 1.5,
  asize = 2,
  mar = c(3, 3, 3, 3),
  nDigits = 3,
  shapeLat = "circle",
  shapeMan = "rectangle",
  shapeInt = "triangle",
  sample.nobs = 100,
  border.width = 2,
  edge.width = 1.5,
  Groups = FALSE,
  group_match = "lat",
  growth = FALSE,
  structural = FALSE,
  edgeLabels = NULL,
  nodeLabels = NULL,
  fixedStyle = 1,
  color = list(lat = "gray10", man = "gold"),
  type = "plot"
)
\end{verbatim}
\end{Usage}
%
\begin{Arguments}
\begin{ldescription}
\item[\code{model}] lavaan syntax

\item[\code{residuals}] residauls TRUE

\item[\code{whatLabels}] "model

\item[\code{rotation}] rotation=2

\item[\code{residScale}] residScale=12

\item[\code{layout}] layout="tree"

\item[\code{edge.label.cex}] 1

\item[\code{edge.label.position}] 0.5

\item[\code{edge.color}] edge.color= black

\item[\code{sizeMan}] 8

\item[\code{sizeMan2}] 4

\item[\code{sizeLat}] 10

\item[\code{sizeLat2}] 6

\item[\code{style}] lisrel

\item[\code{sig}] FALSE

\item[\code{exoVar}] FALSE

\item[\code{exoCov}] TRUE

\item[\code{curve}] 1.5

\item[\code{asize}] 1.5

\item[\code{mar}] c(2,8,3,10)

\item[\code{nDigits}] 2

\item[\code{shapeLat}] "circle, ellipse

\item[\code{shapeMan}] retangle

\item[\code{shapeInt}] triangle ,

\item[\code{sample.nobs}] sample.nobs = 100

\item[\code{border.width}] border.width=1

\item[\code{edge.width}] edge.width=1.5

\item[\code{Groups}] FALSE

\item[\code{group\_match}] group\_match lat

\item[\code{growth}] growth= FALSE, TRUE  model.type ="grouwth"

\item[\code{structural}] structural = FALSE

\item[\code{edgeLabels}] edgeLabels=NULL

\item[\code{nodeLabels}] nodeLabels=NULL

\item[\code{fixedStyle}] fixedStyle=1 default rbase linetype

\item[\code{type}] type ="plot, and res is lavaan simdata result
\end{ldescription}
\end{Arguments}
%
\begin{Value}
plot
\end{Value}
%
\begin{Examples}
\begin{ExampleCode}
## Not run: 
#'
model_test <- '
  # latent variable definitions
     ind60 =~ x1 + x2 + x3
     dem60 =~ y1 + a*y2 + b*y3 + c*y4
     dem65 =~ y5 + a*y6 + b*y7 + c*y8

  # regressions
    dem60 ~ ind60
    dem65 ~ ind60 + dem60

  # residual correlations
    # y1 ~~ y5
    # y2 ~~ y4 + y6
    # y3 ~~ y7
    # y4 ~~ y8
    # y6 ~~ y8
'

HSmodel <- ' visual  =~ x1 + x2 + x3
             textual =~ x4 + x5 + x6
             speed   =~ x7 + x8 + x9 '


m1a <- "
y1 =~ x1 + x2 +x5
y2 =~ x3 + x4 +x6
y1 ~~ y2"

sem_model(model= HSmodel, groups=F,  residuals=F)

sem_model(model= model_test, style="ram", groups=F)
sem_model(model= model_test, residuals=T)
sem_model(model= m1a, residuals=T)

#'

## End(Not run)
\end{ExampleCode}
\end{Examples}
\inputencoding{utf8}
\HeaderA{sem\_plot}{Fast sem plot draw}{sem.Rul.plot}
%
\begin{Description}
Fast sem plot draw
\end{Description}
%
\begin{Usage}
\begin{verbatim}
sem_plot(
  data,
  whatLabels = "std",
  type = "res1",
  opt = 1,
  layout = "tree3",
  curve = 2,
  rotation = 2,
  sizeLat = 10,
  sizeLat2 = 6,
  sizeMan = 8,
  sizeMan2 = 4,
  edge.label.cex = 0.9,
  edge.label.position = 0.5,
  mar = c(1, 5, 1, 5),
  groups = "lat",
  pastel = TRUE,
  fade = FALSE,
  se = FALSE,
  exoVar = FALSE,
  residuals = FALSE,
  residScale = 12,
  intercepts = FALSE,
  semptools = TRUE
)
\end{verbatim}
\end{Usage}
%
\begin{Arguments}
\begin{ldescription}
\item[\code{data}] lavaaan object

\item[\code{whatLabels}] ste

\item[\code{type}] type res 1

\item[\code{opt}] 1, 2, 3,

\item[\code{layout}] layY(opt =2, you make layout), or tree3

\item[\code{curve}] -2

\item[\code{rotation}] rotation = 2

\item[\code{sizeLat}] 8

\item[\code{sizeLat2}] 6

\item[\code{sizeMan}] 6

\item[\code{sizeMan2}] 3

\item[\code{edge.label.cex}] 0.7

\item[\code{edge.label.position}] 0.6

\item[\code{mar}] c(1,5,1,5)

\item[\code{groups}] groups

\item[\code{pastel}] pastel

\item[\code{fade}] default FALSE when opt =3

\item[\code{se}] se false

\item[\code{exoVar}] exoVar=FLASE

\item[\code{residuals}] residuals F

\item[\code{residScale}] 12

\item[\code{intercepts}] intercepts =FALSE becouse semptools

\item[\code{semptools}] using semptools TRUE
\end{ldescription}
\end{Arguments}
%
\begin{Value}
plot
\end{Value}
%
\begin{Examples}
\begin{ExampleCode}

## Not run: 

# library(lavaan)
#import data
example(sem)

sem_plot(fit)

example(cfa)
sem_plot(fit)
#'
HS.model <- '
    visual  =~ x1 + b1*x2 + x3
    textual =~ x4 + b2*x5 + x6
    speed   =~ x7 + b3*x8 + x9
speed ~ visual + textual

'
fit1 <- sem(HS.model, data = HolzingerSwineford1939)

fit1 %>%sem_plot(residuals=T, curve=2)

## End(Not run)
\end{ExampleCode}
\end{Examples}
\inputencoding{utf8}
\HeaderA{shapiro\_test\_all}{Custom function to perform Shapiro-Wilk test on all columns of a data frame}{shapiro.Rul.test.Rul.all}
%
\begin{Description}
Custom function to perform Shapiro-Wilk test on all columns of a data frame
\end{Description}
%
\begin{Usage}
\begin{verbatim}
shapiro_test_all(data, digits = 3)
\end{verbatim}
\end{Usage}
%
\begin{Arguments}
\begin{ldescription}
\item[\code{data}] data.frame

\item[\code{digits}] default 3
\end{ldescription}
\end{Arguments}
%
\begin{Value}
nomalitytest result
\end{Value}
%
\begin{Examples}
\begin{ExampleCode}
## Not run: 
data <- mtcars[,c("mpg","disp","hp","drat","wt","qsec")]

## Shapiro-Wilke Test Results for All Columns
shapiro_test_all(data)


## End(Not run)

\end{ExampleCode}
\end{Examples}
\inputencoding{utf8}
\HeaderA{shapiro\_test\_df}{Calculate Shapiro-Wilke test results for all columns}{shapiro.Rul.test.Rul.df}
%
\begin{Description}
Calculate Shapiro-Wilke test results for all columns
\end{Description}
%
\begin{Usage}
\begin{verbatim}
shapiro_test_df(data, digits = 3)
\end{verbatim}
\end{Usage}
%
\begin{Arguments}
\begin{ldescription}
\item[\code{data}] data.frame

\item[\code{digits}] default 3
\end{ldescription}
\end{Arguments}
%
\begin{Value}
nomalitytest result
\end{Value}
%
\begin{Examples}
\begin{ExampleCode}
## Not run: 
data <- mtcars[,c("mpg","disp","hp","drat","wt","qsec")]

## Shapiro-Wilke Test Results for All Columns
shapiro_test_df(data)


## End(Not run)
\end{ExampleCode}
\end{Examples}
\inputencoding{utf8}
\HeaderA{simdata\_cov}{Regression with covariance, generate structural equation data restoration}{simdata.Rul.cov}
%
\begin{Description}
Regression with covariance, generate structural equation data restoration
\end{Description}
%
\begin{Usage}
\begin{verbatim}
simdata_cov(
  covmat = NULL,
  n = 300,
  dv_mean = 0,
  iv_means = rep(0, ncol(covmat) - 1),
  means = NULL,
  cormat = NULL,
  sds = rep(1, ncol(cormat)),
  seed = NULL
)
\end{verbatim}
\end{Usage}
%
\begin{Arguments}
\begin{ldescription}
\item[\code{covmat}] cov matrix

\item[\code{n}] sample size default = 300

\item[\code{dv\_mean}] dv mean

\item[\code{iv\_means}] iv means

\item[\code{means}] total means default null

\item[\code{cormat}] correlation matrix

\item[\code{sds}] sds is need cor to cov

\item[\code{seed}] set seed today Sys.Data
\end{ldescription}
\end{Arguments}
%
\begin{Value}
sim data
\end{Value}
%
\begin{Examples}
\begin{ExampleCode}

## Not run: 
#'
#simdata mean setting 0
#set.seed(todaty)
cov(mtcars[,c("mpg","hp","wt")]) %>%cov_simdata(n=32)
cov(mtcars[,c("mpg","hp","wt")]) %>%cov_simdata(n=32, means=c(0,0,0))

cov(mtcars[,c("mpg","hp","wt")]) %>%
  cov_simdata(n=32, means=c(20.09, 146.69 ,3.22))

cov(mtcars[,c("mpg","hp","wt")]) %>%
  cov_simdata(n=32, dv_mean=20.09, iv_means=c(146.69 ,3.22))

#simdata
cov(mtcars[,c("mpg","hp","wt")]) %>%
  cov_simdata(n=32) %>%
  lm(formula = mpg ~ hp + wt)%>%summary()

#revised mean data Closer to real-world data
cov(mtcars[,c("mpg","hp","wt")]) %>%
  cov_simdata(n=32, dv_mean=20.09, iv_means=c(146.69 ,3.22)) %>%
  lm(formula = mpg ~ hp + wt)%>%summary()


#real data
lm(formula = mpg ~ hp +wt,data=mtcars)%>%summary()



# define sample statistics
x <- MASS::Cars93[,c("EngineSize", "Horsepower", "RPM")]
covx <- cov(x)
covx <- matrix(c(1.076122, 39.776999, -339.163745,
                 39.776999, 2743.079, 1146.634,
                 -339.163745, 1146.634, 356088.7097), nrow = 3,
               dimnames = list(c("EngineSize", "Horsepower", "RPM"),
                               c("EngineSize", "Horsepower", "RPM")))
covx
# 함수 실행
sim_data <- simdata_cov(covmat = covx, n = 100)
head(sim_data)
# check
sim_data <- simdata_cov(covmat = covx, n = 100)
lm( EngineSize~ Horsepower+RPM, data = sim_data) %>% summary()
#real check
lm( EngineSize ~ Horsepower +RPM , data=x)%>% summary()



#'

## End(Not run)
\end{ExampleCode}
\end{Examples}
\inputencoding{utf8}
\HeaderA{simdata\_sem}{Structural equation data generation using the lavaan syntax}{simdata.Rul.sem}
%
\begin{Description}
Structural equation data generation using the lavaan syntax
\end{Description}
%
\begin{Usage}
\begin{verbatim}
simdata_sem(model, N = 100, datatype = "data.frame", empirical = TRUE)
\end{verbatim}
\end{Usage}
%
\begin{Arguments}
\begin{ldescription}
\item[\code{model}] lavaan model

\item[\code{N}] sample size N

\item[\code{datatype}] "data.frame", "cor", "matrix"

\item[\code{empirical}] empirical
\end{ldescription}
\end{Arguments}
%
\begin{Value}
data
\end{Value}
%
\begin{Examples}
\begin{ExampleCode}

## Not run: 
#'
# devtools::install_github("M-E-Rademaker/cSEM")
# devtools::install_github("M-E-Rademaker/cSEM.DGP")
#
# library(cSEM.DGP)
# library(cSEM)

# sample data
HS.model <- ' visual  =~ x1 + x2 + x3
              textual =~ x4 + x5 + x6
              speed   =~ x7 + x8 + x9
              '

fit <- cfa(HS.model, data = HolzingerSwineford1939)
summary(fit)

#lavaan
HS_model <- ' visual  =~ 0.77*x1 + 0.42*x2 + 0.58*x3
              textual =~ 0.85*x4 + 0.86*x5 + 0.84*x6
              speed   =~ 0.57*x7 + 0.72*x8 + 0.67*x9
              visual ~~  0.46*textual
              visual ~~  0.47*speed
              textual ~~  0.28*speed
          '
#cSEM, cSEM.DGP
# HS_simdata = generateData(.model = HS_model,
#                           .return_type = "data.frame",
#                           .empirical = TRUE, .N= 301)

HS_simdata1 = simdata_sem(model = HS_model,
                          datatype = "data.frame",
                          empirical = TRUE, N= 301)
fit1 <- cfa(HS.model, data = HS_simdata1)
summary(fit1)
fit1 %>% diagram2("std",rotation =2, mar=c(3,12,3,8), groups=T)
#'
#'
#'


## End(Not run)
\end{ExampleCode}
\end{Examples}
\inputencoding{utf8}
\HeaderA{simple\_bar}{Tsimple\_bar}{simple.Rul.bar}
%
\begin{Description}
Tsimple\_bar
\end{Description}
%
\begin{Usage}
\begin{verbatim}
simple_bar(
  data,
  add_y = 10,
  color = 2:17,
  vjust = 1,
  cex = 1,
  line = FALSE,
  main = NULL,
  xlab = "x",
  ylab = "y"
)
\end{verbatim}
\end{Usage}
%
\begin{Arguments}
\begin{ldescription}
\item[\code{data}] table

\item[\code{add\_y}] max 30

\item[\code{color}] color=1:16

\item[\code{vjust}] vjust =10

\item[\code{cex}] cex=1

\item[\code{line}] T/F

\item[\code{main}] main sistitle

\item[\code{xlab}] xlab

\item[\code{ylab}] ylab
\end{ldescription}
\end{Arguments}
%
\begin{Value}
bar graph
\end{Value}
%
\begin{Examples}
\begin{ExampleCode}
## Not run: 
table(mtcars$cyl) %>% simple_bar(cex=1.5)

data(bioChemists, package = "pscl")
bioChemists <- bioChemists %>%
rename(
  Num_Article = art, #articles in last 3 years of PhD
   Sex = fem, #coded 1 if female
   Married = mar, #coded 1 if married
   Num_Kid5 = kid5, #number of childeren under age 6
   PhD_Quality = phd, #prestige of PhD program
   Num_MentArticle = ment #articles by mentor in last 3 years
   )
  table(bioChemists$Num_Article) %>%
      simple_bar(cex=1.5, add_y= 30, vjust=10)


## End(Not run)


\end{ExampleCode}
\end{Examples}
\inputencoding{utf8}
\HeaderA{SKEW}{skew계산 함수}{SKEW}
%
\begin{Description}
skew계산 함수
\end{Description}
%
\begin{Usage}
\begin{verbatim}
SKEW(x)
\end{verbatim}
\end{Usage}
%
\begin{Arguments}
\begin{ldescription}
\item[\code{x}] data  vector
\end{ldescription}
\end{Arguments}
%
\begin{Value}
skewness
\end{Value}
\inputencoding{utf8}
\HeaderA{Sobel\_test}{sobel test}{Sobel.Rul.test}
%
\begin{Description}
sobel test
\end{Description}
%
\begin{Usage}
\begin{verbatim}
Sobel_test(
  a,
  b,
  sa,
  sb,
  roburst = FALSE,
  n = NULL,
  k1 = 2,
  k2 = 2,
  digits = 5,
  conf_level = 0.95
)
\end{verbatim}
\end{Usage}
%
\begin{Arguments}
\begin{ldescription}
\item[\code{a}] est1

\item[\code{b}] est2

\item[\code{sa}] se1

\item[\code{sb}] se2

\item[\code{roburst}] roburst T whide correction

\item[\code{n}] sample size

\item[\code{k1}] number of regressin coefs, default 2

\item[\code{k2}] number of regressin coefs, default 2

\item[\code{digits}] 5

\item[\code{type}] Sobel, Aroian, Goodman
\end{ldescription}
\end{Arguments}
%
\begin{Value}
test result
\end{Value}
%
\begin{Examples}
\begin{ExampleCode}

## Not run: 
#'
# Example usage
a <- 0.5  # Example coefficient a
b <- 0.3  # Example coefficient b
sa <- 0.1 # Standard error of a
sb <- 0.05 # Standard error of b

Sobel_test(a, b, sa, sb)
Sobel_test(a, b, sa, sb, "all")
Sobel_test(a, b, sa, sb, "Sobel")

Sobel_test(a, b, sa, sb,"Aroian")

Sobel_test(a, b, sa, sb,"Goodman")


# roburst se correction : White(1980)
Sobel_test(0.2, 0.3,0.05, 0.04, TRUE, 100)


## End(Not run)

\end{ExampleCode}
\end{Examples}
\inputencoding{utf8}
\HeaderA{sobel\_test\_extend}{sobel\_test\_extend}{sobel.Rul.test.Rul.extend}
%
\begin{Description}
sobel\_test\_extend
\end{Description}
%
\begin{Usage}
\begin{verbatim}
sobel_test_extend(
  coefficients,
  se_values,
  digits = 5,
  show = TRUE,
  sobel = "aroian"
)
\end{verbatim}
\end{Usage}
%
\begin{Arguments}
\begin{ldescription}
\item[\code{coefficients}] coefficients c()

\item[\code{se\_values}] se

\item[\code{digits}] 5

\item[\code{show}] TRUE

\item[\code{sobel}] aroian(default), sobel, goodman
\end{ldescription}
\end{Arguments}
%
\begin{Value}
seobel resutl
\end{Value}
%
\begin{Examples}
\begin{ExampleCode}

## Not run: 
#'

# Example usage
sdata <- data.frame(
  paths = c("IMAG -> EXPE", "IMAG -> SAT",
              "IMAG -> LOY", "EXPE -> QUAL", "EXPE -> VAL", "EXPE -> SAT",
            "QUAL -> VAL", "QUAL -> SAT", "VAL -> SAT", "SAT -> LOY"),
  Original = c(0.578959128, 0.200724200, 0.275149576,
              0.848344408, 0.105477650, -0.002753995,
               0.676655529, 0.122144566, 0.589330650, 0.495479322),
  Mean.Boot = c(0.58232765, 0.21209094, 0.27486564,
                0.84793899, 0.10992051, -0.01008222,
                0.67780527, 0.15132543, 0.56312717, 0.50757755),
  Std.Error = c(0.04358293, 0.06056988, 0.07729000,
                 0.01864573, 0.06968826, 0.06507903,
                0.07760527, 0.08440236, 0.08723293, 0.07856694),
  perc.025 = c(0.497644719, 0.107561558, 0.140641502,
                0.815794494, -0.025566471, -0.147788188,
               0.519959683, -0.008629516, 0.392664443, 0.351774130),
  perc.975 = c(0.6646198, 0.3275312, 0.4161131,
                0.8828875, 0.2545785, 0.1086856,
               0.8166486, 0.2991045, 0.7012494, 0.6556333)
)



sobel_test_extend(coefficients= c("IMAG -> EXPE" = 0.578959128,
                                  "EXPE -> SAT" = -0.002753995,
                                  "SAT -> LOY" = 0.495479322),
                  se_values = c(se1 = 0.04358293,
                                se2 = 0.06507903,
                                se3 = 0.0785669))
# Example usage
coefficients1 <- c("IMAG -> EXPE" = 0.578959128,
                  "EXPE -> SAT" = -0.002753995, "SAT -> LOY" = 0.495479322)
se_values1 <- c(se1 = 0.04358293, se2 = 0.06507903, se3 = 0.0785669)

sobel_test_extend(coefficients1, se_values1)


# Test with unnamed values
coefficients2 <- c(0.578959128, -0.002753995, 0.495479322)
se_values2 <- c(0.04358293, 0.06507903, 0.0785669)

sobel_test_extend(coefficients2, se_values2)


# Test with est and se named values
coefficients3 <- c(imag2expe = 0.578959128,
                  expe2sat = -0.002753995, sat2loy = 0.495479322)
se_values3 <- c(se1 = 0.04358293, se2 = 0.06507903, se3 = 0.0785669)

sobel_test_extend(coefficients3, se_values3)




## End(Not run)


\end{ExampleCode}
\end{Examples}
\inputencoding{utf8}
\HeaderA{sobel\_test\_general}{Generalized Sobel Test function}{sobel.Rul.test.Rul.general}
%
\begin{Description}
Generalized Sobel Test function
\end{Description}
%
\begin{Usage}
\begin{verbatim}
sobel_test_general(
  a,
  SE,
  roburst = FALSE,
  n = NULL,
  k = 2,
  digits = 5,
  conf = 0.95
)
\end{verbatim}
\end{Usage}
%
\begin{Arguments}
\begin{ldescription}
\item[\code{a}] est

\item[\code{SE}] se

\item[\code{roburst}] roburst T whide correction

\item[\code{n}] sample size

\item[\code{k}] number of regressin coefs, default 2

\item[\code{digits}] 5

\item[\code{conf}] 0.90, 0.95, 0.99, 0.999
\end{ldescription}
\end{Arguments}
%
\begin{Value}
result
\end{Value}
%
\begin{Examples}
\begin{ExampleCode}

## Not run: 
#'

general_Sobel_test(a = c(0.5, 0.3), SE = c(0.1, 0.05))

Sobel_test_2(a=0.5, b=0.3, sa=0.1, sb =0.05)

# Input data
a <- c(0.2, 0.3, 0.4, 0.5)
SE <- c(0.05, 0.04, 0.03, 0.02)

# Calculate results
sobel_test_general(a = c(0.2, 0.3, 0.4, 0.5), SE = c(0.05, 0.04, 0.03, 0.02))
aroian_sobel_test_general(a = c(0.2, 0.3, 0.4, 0.5), SE = c(0.05, 0.04, 0.03, 0.02))
Goodman_sobel_test_general(a = c(0.2, 0.3, 0.4, 0.5), SE = c(0.05, 0.04, 0.03, 0.02))
general_Sobel_test(a = c(0.2, 0.3, 0.4, 0.5), SE = c(0.05, 0.04, 0.03, 0.02))
#'
# 예제 데이터
a <- c(0.5, 0.3, 0.4)
SE <- c(0.1, 0.1, 0.1)

# 함수 실행 (기본 95% 신뢰구간)
sobel_test_general(a, SE)

# 함수 실행 (90% 신뢰구간)
result_90 <- sobel_test_general(a, SE, conf=0.90)
print(result_90)

# 함수 실행 (99% 신뢰구간)
result_99 <- sobel_test_general(a, SE, conf=0.99)
print(result_99)

# 함수 실행 (99.9% 신뢰구간)
result_999 <- sobel_test_general(a, SE, conf=0.999)
print(result_999)


## End(Not run)


\end{ExampleCode}
\end{Examples}
\inputencoding{utf8}
\HeaderA{sobel\_test\_serial}{sobel\_test\_extend}{sobel.Rul.test.Rul.serial}
%
\begin{Description}
sobel\_test\_extend
\end{Description}
%
\begin{Usage}
\begin{verbatim}
sobel_test_serial(
  coefficients,
  se_values,
  digits = 5,
  show = TRUE,
  sobel = "aroian"
)
\end{verbatim}
\end{Usage}
%
\begin{Arguments}
\begin{ldescription}
\item[\code{coefficients}] coefficients c()

\item[\code{se\_values}] se

\item[\code{digits}] 5

\item[\code{show}] TRUE

\item[\code{sobel}] aroina, sobel, goodman
\end{ldescription}
\end{Arguments}
%
\begin{Value}
seobel resutl
\end{Value}
%
\begin{Examples}
\begin{ExampleCode}

## Not run: 
#'

# Example usage
sdata <- data.frame(
  paths = c("IMAG -> EXPE", "IMAG -> SAT",
              "IMAG -> LOY", "EXPE -> QUAL", "EXPE -> VAL", "EXPE -> SAT",
            "QUAL -> VAL", "QUAL -> SAT", "VAL -> SAT", "SAT -> LOY"),
  Original = c(0.578959128, 0.200724200, 0.275149576,
              0.848344408, 0.105477650, -0.002753995,
               0.676655529, 0.122144566, 0.589330650, 0.495479322),
  Mean.Boot = c(0.58232765, 0.21209094, 0.27486564,
                0.84793899, 0.10992051, -0.01008222,
                0.67780527, 0.15132543, 0.56312717, 0.50757755),
  Std.Error = c(0.04358293, 0.06056988, 0.07729000,
                 0.01864573, 0.06968826, 0.06507903,
                0.07760527, 0.08440236, 0.08723293, 0.07856694),
  perc.025 = c(0.497644719, 0.107561558, 0.140641502,
                0.815794494, -0.025566471, -0.147788188,
               0.519959683, -0.008629516, 0.392664443, 0.351774130),
  perc.975 = c(0.6646198, 0.3275312, 0.4161131,
                0.8828875, 0.2545785, 0.1086856,
               0.8166486, 0.2991045, 0.7012494, 0.6556333)
)



sobel_test_serial(coefficients= c("IMAG -> EXPE" = 0.578959128,
                                  "EXPE -> SAT" = -0.002753995,
                                  "SAT -> LOY" = 0.495479322),
                  se_values = c(se1 = 0.04358293,
                                se2 = 0.06507903,
                                se3 = 0.0785669))
# Example usage
coefficients1 <- c("IMAG -> EXPE" = 0.578959128,
                  "EXPE -> SAT" = -0.002753995, "SAT -> LOY" = 0.495479322)
se_values1 <- c(se1 = 0.04358293, se2 = 0.06507903, se3 = 0.0785669)

sobel_test_serial(coefficients1, se_values1)


# Test with unnamed values
coefficients2 <- c(0.578959128, -0.002753995, 0.495479322)
se_values2 <- c(0.04358293, 0.06507903, 0.0785669)

sobel_test_serial(coefficients2, se_values2)


# Test with est and se named values
coefficients3 <- c(imag2expe = 0.578959128,
                  expe2sat = -0.002753995, sat2loy = 0.495479322)
se_values3 <- c(se1 = 0.04358293, se2 = 0.06507903, se3 = 0.0785669)

sobel_test_serial(coefficients3, se_values3)




## End(Not run)


\end{ExampleCode}
\end{Examples}
\inputencoding{utf8}
\HeaderA{split\_korean\_word}{split\_korean\_word- Separating vowels and consonants in Korean}{split.Rul.korean.Rul.word}
%
\begin{Description}
split\_korean\_word- Separating vowels and consonants in Korean
\end{Description}
%
\begin{Usage}
\begin{verbatim}
split_korean_word(word, paste = TRUE)
\end{verbatim}
\end{Usage}
%
\begin{Arguments}
\begin{ldescription}
\item[\code{word}] input word, one or more

\item[\code{paste}] Whether you want the outputs to be separate or combined into one
\end{ldescription}
\end{Arguments}
%
\begin{Value}
text data
\end{Value}
%
\begin{Examples}
\begin{ExampleCode}
## Not run: 
split_korean_word("책")
## [1] "ㅊ,ㅐ,ㄱ"
split_korean_word("채")
## [1] "ㅊ,ㅐ,"
split_korean_word("고양이")
#' ##[1] "ㄱ,ㅗ,,ㅇ,ㅑ,ㅇ,ㅇ,ㅣ,"
split_korean_word("곱창밥")
## [1] "ㄱ,ㅗ,ㅂ,ㅊ,ㅏ,ㅇ,ㅂ,ㅏ,ㅂ"

## End(Not run)

\end{ExampleCode}
\end{Examples}
\inputencoding{utf8}
\HeaderA{split\_korean\_word\_check}{Tsplit\_korean\_word\_check is datacheck}{split.Rul.korean.Rul.word.Rul.check}
%
\begin{Description}
Tsplit\_korean\_word\_check is datacheck
\end{Description}
%
\begin{Usage}
\begin{verbatim}
split_korean_word_check(worddata, type = "all")
\end{verbatim}
\end{Usage}
%
\begin{Arguments}
\begin{ldescription}
\item[\code{worddata}] worddata

\item[\code{type}] type defaul all, text, df , all is input and output, text is result , df is data.frame
\end{ldescription}
\end{Arguments}
%
\begin{Value}
split word
\end{Value}
%
\begin{Examples}
\begin{ExampleCode}
## Not run: 
split_korean_word_check("호랑이")

## End(Not run)

\end{ExampleCode}
\end{Examples}
\inputencoding{utf8}
\HeaderA{split\_kw\_df}{split\_kw\_df Create a data column by separating the initial, neutral, and final Korean characters.}{split.Rul.kw.Rul.df}
%
\begin{Description}
split\_kw\_df Create a data column by separating the initial, neutral, and final Korean characters.
\end{Description}
%
\begin{Usage}
\begin{verbatim}
split_kw_df(df, sel_col = "word", remove = TRUE)
\end{verbatim}
\end{Usage}
%
\begin{Arguments}
\begin{ldescription}
\item[\code{df}] data.frame

\item[\code{sel\_col}] select word column

\item[\code{remove}] remove False is overlab, TRUE is generate new data
\end{ldescription}
\end{Arguments}
%
\begin{Value}
data.frame
\end{Value}
%
\begin{Examples}
\begin{ExampleCode}
## Not run: 
wd <- data.frame(
id =  paste0("BY", 1:5),
word = c("토끼", "휘파람", "뻐스", "우리냥","우리낙"),
tone = c("HH","LHH","HL","LHH","LHL" ),
gender =c("F","F","M","M","M"),
age = c(20,60,20,60, 30),
area= c("n","s","n","s","s"),
mix= c(FALSE,FALSE,FALSE,TRUE,TRUE),
stringsAsFactors = FALSE)

wd%>%split_kw_df("word")
wd%>%split_kw_df()


## End(Not run)

\end{ExampleCode}
\end{Examples}
\inputencoding{utf8}
\HeaderA{split\_kw\_df\_match}{split\_kw\_df\_match, using Goeun Kim, PhD (Seoul National University 2024)}{split.Rul.kw.Rul.df.Rul.match}
%
\begin{Description}
split\_kw\_df\_match, using Goeun Kim, PhD (Seoul National University 2024)
\end{Description}
%
\begin{Usage}
\begin{verbatim}
split_kw_df_match(df, sel_col = "word", remove = TRUE, type = "res")
\end{verbatim}
\end{Usage}
%
\begin{Arguments}
\begin{ldescription}
\item[\code{df}] data.frame

\item[\code{sel\_col}] word column

\item[\code{remove}] TRUE is newdata add

\item[\code{type}] type is res and all
\end{ldescription}
\end{Arguments}
%
\begin{Value}
data.frame
\end{Value}
%
\begin{Examples}
\begin{ExampleCode}

## Not run: 

#' wd <- data.frame(
id =  paste0("BY", 1:5),
word = c("토끼", "휘파람", "뻐스", "우리냥","우리낙"),
tone = c("HH","LHH","HL","LHH","LHL" ),
gender =c("F","F","M","M","M"),
age = c(20,60,20,60, 30),
area= c("n","s","n","s","s"),
mix= c(FALSE,FALSE,FALSE,TRUE,TRUE),
stringsAsFactors = FALSE)

# Create new data
wd%>%split_kw_df_match("word", type="res")

wd%>%split_kw_df_match("word", type="all")

wd2 <- data.frame(id =  paste0("BY", 1:5),
word = c("토끼", "휘파람", "뻐스", "우리냥","우리낙망"))

wd2%>%split_kw_df_match("word", type="res")

wd2%>%split_kw_df_match("word", type="all")



## End(Not run)



\end{ExampleCode}
\end{Examples}
\inputencoding{utf8}
\HeaderA{split\_word}{Separate words into syllables}{split.Rul.word}
%
\begin{Description}
Separate words into syllables
\end{Description}
%
\begin{Usage}
\begin{verbatim}
split_word(word)
\end{verbatim}
\end{Usage}
%
\begin{Arguments}
\begin{ldescription}
\item[\code{word}] Words being typed
\end{ldescription}
\end{Arguments}
%
\begin{Value}
Syllable-separated data
\end{Value}
%
\begin{Examples}
\begin{ExampleCode}

## Not run: 
split_word("HLL")
split_word("호랑이")


## End(Not run)

\end{ExampleCode}
\end{Examples}
\inputencoding{utf8}
\HeaderA{str\_cut}{str\_cut function}{str.Rul.cut}
%
\begin{Description}
A function that modifies the values of a specified column in a data frame by cutting the string based on a specific separator (sep).
\end{Description}
%
\begin{Usage}
\begin{verbatim}
str_cut(df, col = "term", sep = "\\(")
\end{verbatim}
\end{Usage}
%
\begin{Arguments}
\begin{ldescription}
\item[\code{df}] A data frame.

\item[\code{col}] The name of the column to modify in the data frame. Default is "term".

\item[\code{sep}] The separator used to cut the string. Default is "\bsl{}(".
\end{ldescription}
\end{Arguments}
%
\begin{Value}
A modified data frame.
\end{Value}
%
\begin{Examples}
\begin{ExampleCode}
## Not run: 
# Example usage
data <- data.frame(
  term = c("AI 학습 진단 및 문제은행 (예: AI 펭톡, EBS 단추)",
           "기타(직접입력)",
           "메타버스/게임화 도구 (예: Zep, 띵커벨, 퀴즈앤, Kahoot! 등)")
)

# Using the function
result <- str_cut(data, "term")
print(result)

## End(Not run)
\end{ExampleCode}
\end{Examples}
\inputencoding{utf8}
\HeaderA{summary.jjstat\_sem}{Summary Method for jjstat\_sem Objects}{summary.jjstat.Rul.sem}
%
\begin{Description}
lavaan 스타일의 요약 출력을 제공하는
\code{jjstat\_sem} 클래스의 summary 메서드이다.
\end{Description}
%
\begin{Usage}
\begin{verbatim}
## S3 method for class 'jjstat_sem'
summary(object, fit.measures = TRUE, standardized = TRUE, digits = 3, ...)
\end{verbatim}
\end{Usage}
%
\begin{Arguments}
\begin{ldescription}
\item[\code{object}] object of class \code{jjstat\_sem}.

\item[\code{fit.measures}] logical.
모형 적합도 지수 출력 여부.

\item[\code{standardized}] logical.
표준화 계수(std.lv, std.all) 출력 여부.

\item[\code{digits}] integer.
출력 시 반올림 자릿수.

\item[\code{...}] additional arguments (unused).
\end{ldescription}
\end{Arguments}
%
\begin{Details}
본 함수는 lavaan의 \code{summary()} 출력 형식을 참고하여,
교육용 SEM 엔진의 결과를 사람이 읽기 쉬운 형태로 요약한다.
표준오차는 bootstrap 기반으로 계산된다.
\end{Details}
%
\begin{Value}
Invisibly returns the original \code{jjstat\_sem} object.
\end{Value}
%
\begin{Examples}
\begin{ExampleCode}
## Not run: 
res <- jjstat_SEM(
  data = HolzingerSwineford1939[,1:6],
  latent_list = list(
    F1 = c("x1","x2","x3"),
    F2 = c("x4","x5","x6")
  ),
  structural = list(path("F1","F2")),
  nboot = 100
)

summary(res)

## End(Not run)
\end{ExampleCode}
\end{Examples}
\inputencoding{utf8}
\HeaderA{summary\_bartlett}{Factorization Bartlett test results}{summary.Rul.bartlett}
%
\begin{Description}
Factorization Bartlett test results
\end{Description}
%
\begin{Usage}
\begin{verbatim}
summary_bartlett(data, test = "df", cat = TRUE)
\end{verbatim}
\end{Usage}
%
\begin{Arguments}
\begin{ldescription}
\item[\code{data}] data. cor data. result data

\item[\code{test}] perform test df, cor, none

\item[\code{cat}] show method
\end{ldescription}
\end{Arguments}
%
\begin{Value}
result
\end{Value}
%
\begin{Examples}
\begin{ExampleCode}

## Not run: 
#'
# Henry Kaiser (1970) introduced an Measure of Sampling
# Adequacy  (MSA) of factor analytic data matrices.
# Kaiser and Rice (1974) then modified it.
#'

New = data.frame(
  A01 = c(4, 4, 5, 4, 5, 5, 5, 4, 4, 2, 4, 4, 5, 4, 4, 4, 4, 5, 1, 3),
  A02 = c(2, 3, 3, 5, 5, 4, 5, 5, 4, 4, 4, 4, 4, 3, 3, 5, 4, 4, 3, 1),
  A03 = c(3, 5, 4, 5, 4, 4, 4, 4, 4, 4, 5, 4, 5, 4, 4, 4, 4, 5, 3, 3),
  A04 = c(4, 4, 4, 4, 5, 4, 5, 4, 4, 3, 5, 4, 5, 5, 5, 4, 5, 5, 1, 5),
  A05 = c(3, 2, 3, 5, 4, 4, 4, 4, 4, 4, 4, 3, 5, 2, 3, 4, 5, 5, 4, 2)
)
New %>%summary_bartlett("df")
New %>%cor()%>% summary_bartlett("cor")
New %>%cor()%>%psych::cortest.bartlett()%>% summary_bartlett()


## End(Not run)
\end{ExampleCode}
\end{Examples}
\inputencoding{utf8}
\HeaderA{summary\_fa}{FA,PCA result summary}{summary.Rul.fa}
%
\begin{Description}
FA,PCA result summary
\end{Description}
%
\begin{Usage}
\begin{verbatim}
summary_fa(fa_data, digits = 3)
\end{verbatim}
\end{Usage}
%
\begin{Arguments}
\begin{ldescription}
\item[\code{fa\_data}] fa, factanal, pricipla result

\item[\code{digits}] roundign defaylt 4
\end{ldescription}
\end{Arguments}
%
\begin{Value}
table
\end{Value}
%
\begin{Examples}
\begin{ExampleCode}

library(MVT)

#  Mardia, K.V., Kent, J.T., and Bibby, J.M. (1979).  _Multivariate
#  Analysis_. Academic Press, London.
data(examScor)
factanal(scale(examScor), factors = 2, rotation = "varimax")%>% summary_fa()

#
psych::principal(cor(examScor), nfactor = 2) %>%summary_fa()

psych::prcomp(cor(examScor), nfactor = 2) %>%summary_fa()

factanal(scale(examScor), factors = 2) %>%summary_fa()

\end{ExampleCode}
\end{Examples}
\inputencoding{utf8}
\HeaderA{summary\_kmo}{Kaiser-Meyer-Olkin factor adequacy (KMO)}{summary.Rul.kmo}
%
\begin{Description}
Kaiser-Meyer-Olkin factor adequacy (KMO)
\end{Description}
%
\begin{Usage}
\begin{verbatim}
summary_kmo(data, test = "df", cat = TRUE)
\end{verbatim}
\end{Usage}
%
\begin{Arguments}
\begin{ldescription}
\item[\code{data}] data, cor data, result data

\item[\code{test}] none, cor, df

\item[\code{cat}] TRUE
\end{ldescription}
\end{Arguments}
%
\begin{Value}
result
\end{Value}
%
\begin{Examples}
\begin{ExampleCode}
## Not run: 
#'
## Example usage
New = data.frame(
  A01 = c(4, 4, 5, 4, 5, 5, 5, 4, 4, 2, 4, 4, 5, 4, 4, 4, 4, 5, 1, 3),
  A02 = c(2, 3, 3, 5, 5, 4, 5, 5, 4, 4, 4, 4, 4, 3, 3, 5, 4, 4, 3, 1),
  A03 = c(3, 5, 4, 5, 4, 4, 4, 4, 4, 4, 5, 4, 5, 4, 4, 4, 4, 5, 3, 3),
  A04 = c(4, 4, 4, 4, 5, 4, 5, 4, 4, 3, 5, 4, 5, 5, 5, 4, 5, 5, 1, 5),
  A05 = c(3, 2, 3, 5, 4, 4, 4, 4, 4, 4, 4, 3, 5, 2, 3, 4, 5, 5, 4, 2)
)
New %>%summary_kmo("df")
New %>%summary_kmo("df", cat=FALSE)
New %>%cor()%>% summary_kmo("cor")
New %>%cor()%>%psych::KMO()%>% summary_kmo()

## End(Not run)

\end{ExampleCode}
\end{Examples}
\inputencoding{utf8}
\HeaderA{summary\_sem}{summary sem result and apa report}{summary.Rul.sem}
%
\begin{Description}
summary sem result and apa report
\end{Description}
%
\begin{Usage}
\begin{verbatim}
summary_sem(
  lav_obj,
  whatLabels = "est",
  caption = "Result of path Coefficient",
  effect = "~",
  effect2 = NULL,
  md = TRUE,
  type = "all",
  opt = 1,
  sizeMan = 8,
  sizeMan2 = 5,
  layout = "tree2",
  mar = c(1, 5, 1, 5),
  intercepts = FALSE
)
\end{verbatim}
\end{Usage}
%
\begin{Arguments}
\begin{ldescription}
\item[\code{lav\_obj}] lavaan object

\item[\code{whatLabels}] whatLabels est

\item[\code{caption}] caption possible

\item[\code{effect}] select "\textasciitilde{}"

\item[\code{effect2}] select 2nd

\item[\code{md}] markdown default TRUEz

\item[\code{type}] all, effect, apa, plot

\item[\code{opt}] opt=1\textasciitilde{}4

\item[\code{sizeMan}] sizeMan = 10

\item[\code{sizeMan2}] sizeMan2 =6

\item[\code{layout}] layout ="tree2"

\item[\code{mar}] mar = c(1,5,1,5)

\item[\code{intercepts}] intercept=FALSE semptools
\end{ldescription}
\end{Arguments}
%
\begin{Value}
result and md
\end{Value}
%
\begin{Examples}
\begin{ExampleCode}

## Not run: 
example(sem)
fit %>% summary_sem()


## End(Not run)
\end{ExampleCode}
\end{Examples}
\inputencoding{utf8}
\HeaderA{table\_apa}{Report frequency analysis result}{table.Rul.apa}
%
\begin{Description}
Report frequency analysis result
\end{Description}
%
\begin{Usage}
\begin{verbatim}
table_apa(data, print = TRUE, md = FALSE, unit = "명")
\end{verbatim}
\end{Usage}
%
\begin{Arguments}
\begin{ldescription}
\item[\code{data}] table\_df objecte

\item[\code{print}] consolt outpur

\item[\code{md}] markdown\_table output

\item[\code{unit}] unit ="명" if you want chage unit
\end{ldescription}
\end{Arguments}
%
\begin{Value}
report data apa
\end{Value}
%
\begin{Examples}
\begin{ExampleCode}

## Not run: 

sats %>%table_df(digits=3) %>% table_apa()
sats %>%table_df(digits=3) %>% table_apa(print=F)
sats %>%table_df(digits=3) %>% table_apa(print=T, md= T)

## End(Not run)

\end{ExampleCode}
\end{Examples}
\inputencoding{utf8}
\HeaderA{table\_df}{Frequency Analysis Functions}{table.Rul.df}
%
\begin{Description}
Frequency Analysis Functions
\end{Description}
%
\begin{Usage}
\begin{verbatim}
table_df(data, exclude = NULL, digits = 4, cat = TRUE)
\end{verbatim}
\end{Usage}
%
\begin{Arguments}
\begin{ldescription}
\item[\code{data}] data.frame

\item[\code{exclude}] Select variables to exclude from analysis

\item[\code{digits}] round

\item[\code{cat}] TRUE Data frequency analysis results message output
\end{ldescription}
\end{Arguments}
%
\begin{Value}
Publication frequency table
\end{Value}
%
\begin{Examples}
\begin{ExampleCode}

## Not run: 
## Identifying variables
 mtcars %>% as_trt("cyl","am","vs","gear")%>% tibble()

 mtcars %>% as_trt("cyl","am","vs","gear")%>% table_df()

 iris %>% as_trt("Species")%>% table_df()
 iris %>% as_trt("Species")%>% table_df() %>%table_apa()


 mtcars %>% as_trt("cyl","am","vs","gear")%>% table_df(exclude=c("cyl", "am"))
 sats %>% table_df()



## End(Not run)


\end{ExampleCode}
\end{Examples}
\inputencoding{utf8}
\HeaderA{table\_df\_bar}{Generate a Bar Plot with Automatic Column Renaming}{table.Rul.df.Rul.bar}
%
\begin{Description}
Generate a Bar Plot with Automatic Column Renaming
\end{Description}
%
\begin{Usage}
\begin{verbatim}
table_df_bar(
  df,
  size.x = 20,
  size.text = 6,
  add_range = 100,
  x.title = "Items",
  flip = TRUE,
  hjust = -0.1,
  vjust = -0.5,
  x11_width = 8,
  x11_height = 5,
  linetype = "dashed",
  linecolor = "gray50",
  show_line = FALSE,
  levels = c("전혀 아니다", "아니다", "중간이다", "그렇다",
    "매우 그렇다"),
  sort = FALSE
)
\end{verbatim}
\end{Usage}
%
\begin{Arguments}
\begin{ldescription}
\item[\code{df}] A data frame with at least three columns: the first column for terms, the second for frequencies, and the third for proportions.

\item[\code{size.x}] Numeric. Font size for axis text. Default is 20.

\item[\code{size.text}] Numeric. Font size for text labels on bars. Default is 6.

\item[\code{add\_range}] Numeric. Additional range added to the maximum frequency for y-axis limits. Default is 100.

\item[\code{x.title}] Character. Title for the x-axis. Default is "Items".

\item[\code{flip}] Logical. Whether to flip the plot to horizontal orientation. Default is TRUE.

\item[\code{hjust}] Numeric. Horizontal adjustment for text labels on bars. Default is -0.1.

\item[\code{vjust}] Numeric. Vertical adjustment for text labels on bars. Default is -0.5.

\item[\code{x11\_width}] Numeric. Width of the x11 window. Default is 8.

\item[\code{x11\_height}] Numeric. Height of the x11 window. Default is 5.

\item[\code{linetype}] Character. Line type for the optional line. Default is "dashed".

\item[\code{linecolor}] Character. Color of the optional line. Default is "gray50".

\item[\code{show\_line}] Logical. Whether to display a line connecting the bars. Default is FALSE.

\item[\code{levels}] Character vector. Custom order for categories on the x-axis. Default is c("전혀 아니다", "아니다", "중간이다", "그렇다", "매우 그렇다").

\item[\code{sort}] Logical. Whether to sort the bars by frequency in descending order. Default is FALSE.
\end{ldescription}
\end{Arguments}
%
\begin{Value}
A bar plot displayed in an x11 window.
\end{Value}
%
\begin{Examples}
\begin{ExampleCode}
## Not run: 
data <- data.frame(term = c("1", "2", "3", "4"), Freq = c(250, 66, 215, 364),
                   `prop(%)` = c(27.9, 7.37, 24.0, 40.7))
table_df_bar(data)

## End(Not run)
\end{ExampleCode}
\end{Examples}
\inputencoding{utf8}
\HeaderA{table\_df\_bar2}{Bar and Line Plot for Frequency Data}{table.Rul.df.Rul.bar2}
%
\begin{Description}
This function creates a bar and line plot using term frequency and proportion data.
\end{Description}
%
\begin{Usage}
\begin{verbatim}
table_df_bar2(
  df,
  size.x = 16,
  size.text = 4,
  add_range = 100,
  x.title = "Items",
  hjust = -0.1,
  vjust = -0.5,
  x11_width = 8,
  x11_height = 5,
  linetype = "dashed",
  linecolor = "gray50",
  show_line = TRUE,
  remove_bar = FALSE,
  sort = FALSE,
  text_angle = 45
)
\end{verbatim}
\end{Usage}
%
\begin{Arguments}
\begin{ldescription}
\item[\code{df}] A data frame with three columns: term (character), Freq (numeric)

\item[\code{size.x}] Numeric. Font size for x-axis labels. Default is 16.

\item[\code{size.text}] Numeric. Font size for text labels inside the plot. Default is 4.

\item[\code{add\_range}] Numeric. Extra space added to the top of the y-axis. Default is 100.

\item[\code{x.title}] Character. Title for the x-axis. Default is "Items".

\item[\code{hjust}] Numeric. Horizontal justification for text labels. Default is -0.1.

\item[\code{vjust}] Numeric. Vertical justification for text labels. Default is -0.5.

\item[\code{x11\_width}] Numeric. Width of the X11 plot window. Default is 8.

\item[\code{x11\_height}] Numeric. Height of the X11 plot window. Default is 5.

\item[\code{linetype}] Character. Line type used for connecting points. Default is "dashed".

\item[\code{linecolor}] Character. Line color used for connecting points. Default is "gray50".

\item[\code{show\_line}] Logical. Whether to show the connecting line. Default is TRUE.

\item[\code{remove\_bar}] Logical. Whether to remove the bar plot. Default is FALSE.

\item[\code{sort}] Logical. Whether to sort the data in descending order of frequency. Default is FALSE.

\item[\code{text\_angle}] Numeric. Angle for the x-axis text. Default is 45.
\end{ldescription}
\end{Arguments}
%
\begin{Value}
None. This function opens a new graphics window and displays a plot.
\end{Value}
%
\begin{Examples}
\begin{ExampleCode}
## Not run: 
df <- data.frame(term = c("A", "B", "C"),
                 Freq = c(50, 30, 40),
                 `prop(%)` = c(50, 30, 40))
table_df_bar2(df, sort = TRUE)

## End(Not run)

\end{ExampleCode}
\end{Examples}
\inputencoding{utf8}
\HeaderA{table\_df\_colnames}{table\_df\_colnames}{table.Rul.df.Rul.colnames}
%
\begin{Description}
table\_df\_colnames
\end{Description}
%
\begin{Usage}
\begin{verbatim}
table_df_colnames(df, input = "구분")
\end{verbatim}
\end{Usage}
%
\begin{Arguments}
\begin{ldescription}
\item[\code{df}] df

\item[\code{input}] "구분
\end{ldescription}
\end{Arguments}
%
\begin{Value}
df
\end{Value}
\inputencoding{utf8}
\HeaderA{table\_df\_multi}{Table and Summary Function for Data Frames}{table.Rul.df.Rul.multi}
%
\begin{Description}
This function processes a data frame to clean missing rows, calculates the frequency
and proportion of terms, and provides different outputs based on the selected type.
\end{Description}
%
\begin{Usage}
\begin{verbatim}
table_df_multi(df, type = "res", digits = 4, msg = TRUE)
\end{verbatim}
\end{Usage}
%
\begin{Arguments}
\begin{ldescription}
\item[\code{df}] A data frame to process.

\item[\code{type}] A character string indicating the type of result to return. Options are:
\begin{itemize}

\item{} \code{"res"}: Returns the frequency table (default).
\item{} \code{"freq\_table"}: Returns the frequency table.
\item{} \code{"missing"}: Returns the list of missing rows.
\item{} \code{"missing\_list"}: Returns the list of missing rows.
\item{} \code{"summary"}: Returns a summary of the frequency table.
\item{} \code{"data"}: Returns the cleaned data frame.

\end{itemize}


\item[\code{digits}] An integer specifying the number of decimal places to round proportions. Default is 4.

\item[\code{msg}] A logical value. If \code{TRUE}, displays messages about the process. Default is \code{TRUE}.
\end{ldescription}
\end{Arguments}
%
\begin{Value}
Depending on the \code{type} parameter:
\begin{itemize}

\item{} A frequency table (\code{"res"}, \code{"freq\_table"}, or \code{"summary"}).
\item{} A list of indices of missing rows (\code{"missing"} or \code{"missing\_list"}).
\item{} The cleaned data frame (\code{"data"}).

\end{itemize}

\end{Value}
%
\begin{Examples}
\begin{ExampleCode}
## Not run: 
# Example: Survey Data with Repeated Questions
# Survey data includes repeated questions with different phrasing
survey_data <- data.frame(
  question1 = c("Yes", "No", "Yes", "", "No"),
  question2 = c("No", "Yes", "Yes", "", ""),
  question3 = c("Yes", "", "No", "Yes", "No")
)

# Process the survey data to identify duplicate responses
result <- table_df_multi(survey_data, type = "res", digits = 2, msg = TRUE)

# Output the frequency table
print(result)

# Additional example to find missing rows
missing_info <- table_df_multi(survey_data, type = "missing", msg = FALSE)
print(missing_info)

## End(Not run)
\end{ExampleCode}
\end{Examples}
\inputencoding{utf8}
\HeaderA{table\_freq}{Frequency analysis table}{table.Rul.freq}
%
\begin{Description}
Frequency analysis table
\end{Description}
%
\begin{Usage}
\begin{verbatim}
table_freq(
  data,
  ...,
  prop = TRUE,
  plot = FALSE,
  angle = 0,
  size_text = 4,
  size_axis = 12,
  legend.position = "top",
  reorder = FALSE,
  type = "res"
)
\end{verbatim}
\end{Usage}
%
\begin{Arguments}
\begin{ldescription}
\item[\code{data}] data.frame

\item[\code{...}] vars If you input varabale names "var1", "var2" ,...

\item[\code{prop}] proportion value, default TRUE

\item[\code{plot}] plot

\item[\code{angle}] axis text angle

\item[\code{size\_text}] default 3, bargraph text

\item[\code{size\_axis}] default 12, axis test

\item[\code{legend.position}] legend.position="" none. top, bottom, left, right

\item[\code{reorder}] TRUE is plot x-variable reoder

\item[\code{type}] all, res=data, g=plot=graph
\end{ldescription}
\end{Arguments}
%
\begin{Value}
Freqency table
\end{Value}
%
\begin{Examples}
\begin{ExampleCode}
## Not run: 
Mtcars = mtcars
Mtcars$am = factor(Mtcars$am, levels=c(0,1), labels= c("automatic","manual" ))
Mtcars$vs  = factor(Mtcars$vs, levels=c(0,1), labels= c("V-shaped","straight" ))
Mtcars$cyl  = factor(Mtcars$cyl, levels=c(4,6,8), labels= c("cyl-4","cyl-6","cyl-8"))

Mtcars %>% table_freq("am")
Mtcars %>% table_freq("am", prop=F)
Mtcars %>% table_freq("vs", plot=FALSE)

Mtcars %>% table_freq("vs")
Mtcars %>% table_freq("vs","am")

Mtcars %>% table_freq("vs","am","cyl")

Mtcars %>% table_freq("vs","am","cyl", angle=90)
Mtcars %>% table_freq("vs","am","cyl", angle=90, reorder = TRUE)

Mtcars %>% table_freq("vs","am","cyl", angle=90)

##error we need option  type = "g"
Mtcars %>% table_freq("vs","am","cyl", angle=90, reorder = TRUE)+ylim(0,14)
Mtcars %>% table_freq("vs","am","cyl", angle=90, reorder = TRUE, type="g")+ylim(0,14)

Mtcars %>% table_freq("vs","am") %>% arrange(am)
Mtcars %>% table_freq("vs","am") %>% arrange(vs)
Mtcars %>% table_freq("vs","am") %>% arrange(Freq)

#ONly data : $result  $g -> NULL
 mtcars %>% table_freq("vs","am", type = "all", plot = F)

# this NULL
mtcars %>% table_freq("vs","am", plot = F, type="g")

## End(Not run)



\end{ExampleCode}
\end{Examples}
\inputencoding{utf8}
\HeaderA{table\_reorder}{Reorder rows of data}{table.Rul.reorder}
%
\begin{Description}
Reorder rows of data
\end{Description}
%
\begin{Usage}
\begin{verbatim}
table_reorder(data, reorder = 1:nrow(data))
\end{verbatim}
\end{Usage}
%
\begin{Arguments}
\begin{ldescription}
\item[\code{data}] data.fram

\item[\code{reorder}] reorder=c(select row order)
\end{ldescription}
\end{Arguments}
%
\begin{Value}
data
\end{Value}
%
\begin{Examples}
\begin{ExampleCode}

## Not run: 
#'
# Example data
data0 <- data.frame(
  source = c("A05", "A08", "A04", "A06",
  "A09", "A11", "A13", "A02", "A03", "A01",
  "A14", "A15", "A07", "A12", "A10",
  "eigen_value", "Proportion_Var", "Cumulative_Var"),
  MR1 = c("", "", "", "", "", "", "", "0.362",
  "0.449", "0.457", "0.47", "0.538", "0.62",
  "0.792", "0.96", "3.363", "0.224", "0.224"),
  MR2 = c("", "0.382", "", "", "0.445", "0.724",
  "0.906", "0.392", "", "", "0.487", "", "", "",
   "", "2.095", "0.14", "0.364"),
  MR3 = c("", "", "0.402", "0.963", "", "", "",
  "", "", "", "", "", "", "", "", "1.691",
  "0.113", "0.477")
)

# 함수 테스트
new_order <- c(2, 3, 1, 4:16)
# 1번 columns backwards, and the rest in order
table_reorder(data0, new_order)

## End(Not run)

\end{ExampleCode}
\end{Examples}
\inputencoding{utf8}
\HeaderA{table\_sum}{Functions to sum parts of a table}{table.Rul.sum}
%
\begin{Description}
Functions to sum parts of a table
\end{Description}
%
\begin{Usage}
\begin{verbatim}
table_sum(data, ..., trans = FALSE, sep = "_")
\end{verbatim}
\end{Usage}
%
\begin{Arguments}
\begin{ldescription}
\item[\code{data}] table,data.frame

\item[\code{...}] contrast excon= c(1, 1, -1, -1)

\item[\code{trans}] data transpose, con must be entered to match the row in the data

\item[\code{sep}] combine sep ="/'
\end{ldescription}
\end{Arguments}
%
\begin{Value}
table
\end{Value}
%
\begin{Examples}
\begin{ExampleCode}

## Not run: 
#'
#'
data <- data.frame(T = c(30, 40, 60, 120),
                   S = c(50, 70, 60, 80))
rownames(data) <- LETTERS[1:4]

data = as.matrix(data)


table_sum(data, con= c(1, 1, -1, -1))
table_sum(data, con= c(1, 1, 2, -1))
table_sum(data, con= c(1, 1, 2, 3))
table_sum(data, con= c(0, 1, 2, 1))
table_sum(data, con= c(0, 1, 0, 2))

New = data.frame(
  row.names= c('No_Nicotine_Dependence', 'Nicotine_Dependence'),

  "A" = c(130, 119),
  "B" = c(210, 267),
  "C" = c(43, 91),
  "D" = c(114, 254),
  "E" = c(20, 67)
);New

t(New)%>%table_sum(c(1,1,3,2,2))

New %>%table_sum(c(1,1))
t(New)
t(New)%>%table_sum(c(1,1,1,2,2))

New%>%table_sum(c(1,1,3,2,2), trans=TRUE)
New%>%table_sum(1,1,3,2,2, trans=TRUE)
New%>%table_sum(c(1,1), trans=FALSE)


## End(Not run)


\end{ExampleCode}
\end{Examples}
\inputencoding{utf8}
\HeaderA{text2df}{Text to data.frame /It is characterized by the use of blank spaces to separate content.}{text2df}
%
\begin{Description}
Text to data.frame /It is characterized by the use of blank spaces to separate content.
\end{Description}
%
\begin{Usage}
\begin{verbatim}
text2df(text, header = TRUE, type = "res", add_vars = NULL)
\end{verbatim}
\end{Usage}
%
\begin{Arguments}
\begin{ldescription}
\item[\code{text}] text(Row variable names must not contain spaces. )

\item[\code{header}] header = TRUE

\item[\code{type}] res, check

\item[\code{add\_vars}] additinal factor
\end{ldescription}
\end{Arguments}
%
\begin{Value}
table data.frame
\end{Value}
%
\begin{Examples}
\begin{ExampleCode}

## Not run: 
#'
text <- "변인 평균 표준편차 왜도 첨도
성취감 4.24 0.73 -0.82 0.39
인정감 4.19 0.60 -0.36 -0.42
책임감 4.36 0.56 -0.61 -0.29
발전성 3.65 0.89 -0.50 -0.04
감독 3.59 0.88 -0.46 0.13
보수 3.56 0.90 -0.41 -0.18
대인관계 3.75 0.81 -0.14 -0.54
작업환경 3.83 0.75 -0.21 -0.44
안정성 3.28 1.08 -0.28 -0.68
감정이입 3.96 0.72 -0.40 -0.33
직무적합 4.00 0.75 -0.35 -0.78
일체감 4.02 0.75 -0.59 -0.13
현학원재직 3.97 0.79 -0.77 0.37
타학원이직 2.99 1.06 -0.41 -0.77
타직종이직 2.76 1.07 -0.13 -0.70"

text2df(text)

text2df(text, vars=c(rep("동기요인", 4), rep("위생요인",5), rep("개인요인",6)  ))

text2df(text, T)%>% make_df_text()

#direct output
"var B se b t
성취감 -.03 .09 -.03 -.33
인정감 .24 .09 .20 2.50
책임감 .30 .08 .23 3.56***
발전성 -.08 .06 -.10 -1.29
감독 .06 .05 .08 1.35
보수 -.14 .07 -.18 -2.01*
대인관계 .06 .06 .06 .89
작업환경 .33 .09 .34 3.82***
안정성 .17 .05 .26 3.42**
성별 -.05 .09 -.03 -.52
연령 -.20 .07 -.22 -3.01**
학력 .05 .07 .04 .70
소득 .00 .04 .01 .12
종사기간 .17 .06 .23 2.71**
기관규모 -.10 .03 -.20 -3.36**" %>%text2df()

reg_text = "var B se b t
성취감 -.03 .09 -.03 -.33
인정감 .24 .09 .20 2.50
책임감 .30 .08 .23 3.56***
발전성 -.08 .06 -.10 -1.29
감독 .06 .05 .08 1.35
보수 -.14 .07 -.18 -2.01*
대인관계 .06 .06 .06 .89
작업환경 .33 .09 .34 3.82***
안정성 .17 .05 .26 3.42**
성별 -.05 .09 -.03 -.52
연령 -.20 .07 -.22 -3.01**
학력 .05 .07 .04 .70
소득 .00 .04 .01 .12
종사기간 .17 .06 .23 2.71**
기관규모 -.10 .03 -.20 -3.36**"

text2df(reg_text)
#'
text2df(reg_text) %>% tibble()

#text to numeric
text2df(reg_text)%>% Char2num() %>% tibble()

text2df(reg_text)%>% Char2num(iter=2:4) %>% tibble()
#'

text ="d           Estimate  Std.Err  z-value  P(>|z|)
  read~ppsych           -0.275    0.037   -7.385    0.000
    read~motiv             0.461    0.037   12.404    0.000
arith~motiv             0.600    0.036   16.771    0.000
"
text2df(text)
text2df(text) %>% tibble::tibble()

text ="             Estimate Std.Error tvalue Pr(>|t|)
(Intercept) -1.336e-07  3.608e-01   0.000        1
ppsych      -2.747e-01  3.730e-02  -7.363  7.51e-13***
motiv        4.613e-01  3.730e-02  12.367  <2e-16***
"
text2df(text) %>%tibble()
text2df(text) %>%make_df_text()


text ="                   Estimate  Std.Err  z-value  P(>|z|)

    read~ppsych           -0.275    0.037   -7.385    0.000
    read~motiv             0.461    0.037   12.404    0.000

    arith~ppsych           -0.096    0.037   -2.616    0.009
     arith~motiv             0.576    0.037   15.695    0.000
"
text2df(text)



text =" 과장 대리 부장 사원 주임 차장 SUM
남자 30(39.5%) 21(38.9%) 34(79.1%) 30(49.2%)  10(37%)     28(71.8%) 153
여자 46(60.5%) 33(61.1%) 9(20.9%) 31(50.8%) 17(63%)    11(28.2%) 147
SUM 76 54 43 61 27 39 300
"
text2df(text)
#' 
## End(Not run)

\end{ExampleCode}
\end{Examples}
\inputencoding{utf8}
\HeaderA{TE\_effect}{Total effect}{TE.Rul.effect}
%
\begin{Description}
Total effect
\end{Description}
%
\begin{Usage}
\begin{verbatim}
TE_effect(model_sem, effect1 = "TE", effect2 = "te", label = FALSE, ci = F)
\end{verbatim}
\end{Usage}
%
\begin{Arguments}
\begin{ldescription}
\item[\code{model\_sem}] lavaan objgect

\item[\code{effect1}] "IE"

\item[\code{effect2}] "IE"

\item[\code{label}] label

\item[\code{ci}] confidence interval
\end{ldescription}
\end{Arguments}
%
\begin{Value}
parameterEstimates()
\end{Value}
%
\begin{Examples}
\begin{ExampleCode}

## Not run: 
#library(lavaan) # only needed once per session

models1 <- '
  # measurement model
    ind60 =~ x1 + x2 + x3
    dem60 =~ y1 + y2 + y3 + y4
    dem65 =~ y5 + y6 + y7 + y8
    # regressions
    dem60 ~ ind60
    dem65 ~ ind60 + dem60
      # residual correlations
    y1 ~~ y5
    y2 ~~ y4 + y6
    y3 ~~ y7
    y4 ~~ y8
    y6 ~~ y8

models2 ='
  # measurement model
    ind60 =~ x1 + x2 + x3
    dem60 =~ y1 + y2 + y3 + y4
    dem65 =~ y5 + y6 + y7 + y8
     # regressions
    dem60 ~ a1*ind60
    dem65 ~ a2*ind60 + a3*dem60
 #NEw
   ie := a1*a3
   te := a2 +a1*a3
   diff := a1 - a3




fits1 <- sem(models1, data=PoliticalDemocracy)
fits2 <- sem(models2, data=PoliticalDemocracy)
summary(fits1, standardized=TRUE)
summary(fits2, standardized=TRUE)
fits2 %>% TE_effect()

## End(Not run)

\end{ExampleCode}
\end{Examples}
\inputencoding{utf8}
\HeaderA{to\_long}{longdata transfomation}{to.Rul.long}
%
\begin{Description}
longdata transfomation
\end{Description}
%
\begin{Usage}
\begin{verbatim}
to_long(
  data,
  names_to = "names",
  values_to = "values",
  cols = NULL,
  fix = NULL,
  rownames_to_column = FALSE,
  rowname = "rows"
)
\end{verbatim}
\end{Usage}
%
\begin{Arguments}
\begin{ldescription}
\item[\code{data}] widedata

\item[\code{names\_to}] name

\item[\code{values\_to}] value

\item[\code{cols}] colrange

\item[\code{fix}] column fixed (thereofore not using cols )

\item[\code{rownames\_to\_column}] rownames\_to\_column=FALSE

\item[\code{rowname}] defai;t acce
\end{ldescription}
\end{Arguments}
%
\begin{Value}
longdata
\end{Value}
%
\begin{Examples}
\begin{ExampleCode}
## Not run: 

mtcars[,1:4] %>%to_long("car", "value")
mtcars[,1:4] %>%to_long("car", "value", fix="cyl")
mtcars[,1:4] %>%to_long("car", "value", fix= 1:2)
mtcars[,1:4] %>%to_long("car", "value", fix= c(1))
mtcars[,1:4] %>%to_long("car", "value", fix= c(2))
mtcars[,1:4] %>%to_long("car", "value", fix=c("cyl","mpg"))
mtcars[,1:6] %>%long_df("car", "value", rownames_to_column = F, cols = c(1,3:6))







## End(Not run)
\end{ExampleCode}
\end{Examples}
\inputencoding{utf8}
\HeaderA{to\_table}{data.frame to table}{to.Rul.table}
%
\begin{Description}
data.frame to table
\end{Description}
%
\begin{Usage}
\begin{verbatim}
to_table(data, sel = ncol(data) - 1, value = ncol(data), type = "mat")
\end{verbatim}
\end{Usage}
%
\begin{Arguments}
\begin{ldescription}
\item[\code{data}] data.frame

\item[\code{sel}] select column (names\_to)

\item[\code{value}] values\_to

\item[\code{type}] mat(matrix), df(data.frame)
\end{ldescription}
\end{Arguments}
%
\begin{Value}
table
\end{Value}
%
\begin{Examples}
\begin{ExampleCode}
## Not run: 
data <- data.frame(am = c(0, 1, 0, 1),
                   vs = c(0, 0, 1, 1),
                   Freq = c(12, 6, 7, 7))
data
#make  table
to_table(data)

mtcars %>%select(am, vs) %>% table()

mtcars %>%select(am, vs) %>% table() %>%data.frame()

mtcars %>%select(am, vs) %>% table() %>%data.frame() %>%to_table()

mtcars %>%select(am, vs) %>% table() %>%data.frame() %>%to_table()%>%addmargins()

mtcars %>%select(am, vs) %>% table() %>%data.frame() %>%to_table()%>%add_ratio()

mtcars %>%select(am, vs) %>% table() %>%data.frame() %>%
to_table()%>%add_ratio()%>%as.matrix()%>%addmargins()

mtcars %>%select(am, vs) %>% table() %>%data.frame() %>%to_table()%>%add_ratio_df()

mtcars %>%select(am, vs) %>% table() %>%data.frame() %>%to_table()%>%add_sum()

mtcars %>%select(am, vs) %>% table() %>%data.frame() %>%to_table(sel="vs")

mtcars %>%select(am, cyl) %>% table()

mtcars %>%select(am, cyl) %>% table() %>%data.frame()

mtcars %>%select(am, cyl) %>% table() %>%data.frame() %>%to_table()

mtcars %>%select(am, cyl) %>% table() %>%data.frame() %>%to_table()%>%add_ratio()

mtcars %>%select(am, cyl) %>% table() %>%data.frame() %>%to_table()%>%add_ratio_df()

mtcars %>%select(cyl, am) %>% table() %>%data.frame() %>%to_table() %>%add_sum()

mat = matrix(c(10, 50, 10, 28, 22, 16, 68,12, 14, 124, 116, 120),
ncol = 4,
dimnames = list(c("20s","30s","40s"),
                c("A","B","C","D")))
mat

mat  %>% long_df("age","brand","count") %>%to_table()
# mean data to table
aggregate(len ~ ., data = ToothGrowth, mean) %>% xtabs(formula = len ~.)
#this is jjstat method
aggregate(len ~ ., data = ToothGrowth, mean) %>% to_table()

aggregate(len ~ ., data = ToothGrowth, mean) %>% to_table("dose","len")

#same method
aggregate(len ~ ., data = ToothGrowth, mean) %>% xtabs(len ~ supp + dose, data =.)


## End(Not run)
\end{ExampleCode}
\end{Examples}
\inputencoding{utf8}
\HeaderA{to\_wide}{pivot\_wider}{to.Rul.wide}
%
\begin{Description}
pivot\_wider
\end{Description}
%
\begin{Usage}
\begin{verbatim}
to_wide(
  data,
  names_from = "names",
  values_from = "values",
  names_prefix = "",
  names_sep = "_",
  names_sort = FALSE,
  rm_cols = NULL,
  add_id = NULL
)
\end{verbatim}
\end{Usage}
%
\begin{Arguments}
\begin{ldescription}
\item[\code{data}] long format data

\item[\code{names\_from}] names\_from

\item[\code{values\_from}] values\_from

\item[\code{names\_prefix}] names\_prefix +variable

\item[\code{names\_sep}] sep

\item[\code{names\_sort}] colnames sort

\item[\code{rm\_cols}] Remove variables that interfere with wide format by creating individual IDs

\item[\code{add\_id}] Select Repeating Groups to create IDs, line in column names or column numbers
\end{ldescription}
\end{Arguments}
%
\begin{Value}
wide data
\end{Value}
%
\begin{Examples}
\begin{ExampleCode}


## Not run: 
mtcars[1:3, 1:7]
#this is error
mtcars[1:4, 1:5] %>% to_long() %>% wide_df()

mtcars[1:3, 1:7] %>% add_id()%>% to_wide()
mtcars[1:3, 1:7] %>% to_wide(add_id=TRUE)
mtcars[1:3, 1:7] %>% to_wide(cols= 3:7)
mtcars[1:3, 1:7] %>% to_wide(fix=1:2)
mtcars[1:3, 1:7] %>% to_wide(fix= 1:3, rownames_to_column=F)

mtcars[1:3, 1:7] %>% to_wide() %>% mutate(ik = "bc", io = "a")
mtcars[1:3, 1:7] %>% to_wide() %>% mutate(ik = "bc", io = "a")%>% to_wide()

mtcars[1:3, 1:7] %>% to_wide() %>%
  mutate(ik = "bc", io = "a")%>%
  to_wide(names_from = c("names","ik"))

mtcars[1:3, 1:7] %>% to_wide() %>%
  mutate(ik = rownames, io = "a")%>%
  to_wide(names_from = c("ik"))%>%pall()

mtcars[1:3, 1:7] %>% to_wide() %>%
  mutate(ik = "bc", io = "a")%>%
  to_wide(names_from = 2)



#일반적인 데이터에서  롱포맷으로 바꾸기
data <- data.frame(
  weight = c(4.17, 5.58, 5.18, 6.11, 4.50, 4.61, 5.17, 4.53, 5.33, 5.14,
             4.81, 4.17, 4.41, 3.59, 5.87, 3.83, 6.03, 4.89, 4.32, 4.69,
             6.31, 5.12, 5.54, 5.50, 5.37, 5.29, 4.92, 6.15, 5.80, 5.26),
  group = c(rep("ctrl", 10), rep("trt1", 10), rep("trt2", 10))
)
data
data%>%add_rep_id("group")

data%>%add_rep_id("group")%>%jjstat::as_trt("id") %>%wide_df("group", "weight")

data  %>%wide_df("group", "weight", add_id="group")


#개별적인 그룹을 만들어진 경우
data2 <- data.frame(
  weight = c(4.17, 5.58, 5.18, 6.11, 4.50, 4.61, 5.17, 4.53, 5.33, 5.14,
             4.81, 4.17, 4.41, 3.59, 5.87, 3.83, 6.03, 4.89, 4.32, 4.69,
             6.31, 5.12, 5.54, 5.50, 5.37, 5.29, 4.92, 6.15, 5.80, 5.26),
  ids = c(paste0("C", 1:10), paste0("T",1:10), paste0("S", 1:10))
)

data2

data2 %>%jjstat::mutate_col(group= c(rep("ctrl", 10),
                 rep("trt1", 10), rep("trt2", 10)), col=2)

# 그룹을 나누어주는 아이디 생성
data2 %>%
jjstat::mutate_col(group= c(rep("ctrl", 10),
        rep("trt1", 10), rep("trt2", 10)), col=2)%>%
        add_rep_id("group")

# 개별적인 그룹을 만드는 열이 존재하는 경우 --> 문제발생
data2 %>%
    jjstat::mutate_col(group= c(rep("ctrl", 10),
          rep("trt1", 10), rep("trt2", 10)), col=2)%>%
          wide_df("group", "weight", add_id="group")

#개별적인 그룹을 만드는 열을 제거하는 경우 --> 문제해결
data2 %>%
  jjstat::mutate_col(group= c(rep("ctrl", 10),
  rep("trt1", 10), rep("trt2", 10)),  col=2)%>%
  wide_df("group", "weight", add_id="group", rm_cols="ids") %>%print(n=Inf)
#'
data3= data.frame(
  weight=c(1:9),
  group=c(rep("a",3),rep("b",3),rep("c",3))
)
data3
data3 %>% to_wide("group", "weight", add_id=2)


## End(Not run)
\end{ExampleCode}
\end{Examples}
\inputencoding{utf8}
\HeaderA{ttest\_df}{Repeat t.test to create a table}{ttest.Rul.df}
%
\begin{Description}
Repeat t.test to create a table
\end{Description}
%
\begin{Usage}
\begin{verbatim}
ttest_df(
  data,
  dv_var = NULL,
  iv_vars = NULL,
  grp_mean = TRUE,
  unite_t = FALSE,
  unite_p = FALSE,
  digits = 2,
  var.equal = TRUE,
  paired = FALSE,
  sig = FALSE
)
\end{verbatim}
\end{Usage}
%
\begin{Arguments}
\begin{ldescription}
\item[\code{data}] data. fram

\item[\code{dv\_var}] ddpendent variable

\item[\code{iv\_vars}] Independent variables

\item[\code{grp\_mean}] mean each level

\item[\code{unite\_t}] tvalue+sig

\item[\code{unite\_p}] p\_vakye+sig

\item[\code{digits}] digits 2

\item[\code{var.equal}] variance equal -student ttest

\item[\code{paired}] paired t.test (TRUE)

\item[\code{sig}] star
\end{ldescription}
\end{Arguments}
%
\begin{Value}
t.test table
\end{Value}
%
\begin{Examples}
\begin{ExampleCode}

## Not run: 

ttest_df(data = mtcars, dv_var = "mpg", iv_vars = c("am", "vs"))

ttest_df(data = mtcars, dv_var = "mpg", iv_vars = c("am", "vs"), sig=T)


ttest_df(data = mtcars, dv_var = "mpg", iv_vars = c("am", "vs"), sig=T, unite_t=TRUE)

ttest_df(data = mtcars, dv_var = "mpg", iv_vars = c("am", "vs"), sig=T, unite_p=TRUE)

ttest_df(data = mtcars, dv_var = "mpg", iv_vars = c("am", "vs"), grp_mean=FALSE)



## End(Not run)


\end{ExampleCode}
\end{Examples}
\inputencoding{utf8}
\HeaderA{t\_sep}{t\_sep separate t value and star}{t.Rul.sep}
%
\begin{Description}
t\_sep separate t value and star
\end{Description}
%
\begin{Usage}
\begin{verbatim}
t_sep(data)
\end{verbatim}
\end{Usage}
%
\begin{Arguments}
\begin{ldescription}
\item[\code{data}] data
\end{ldescription}
\end{Arguments}
%
\begin{Value}
add column
\end{Value}
%
\begin{Examples}
\begin{ExampleCode}

## Not run: 
#'
reg_text = "var B se b t
성취감 -.03 .09 -.03 -.33
인정감 .24 .09 .20 2.50
책임감 .30 .08 .23 3.56***
발전성 -.08 .06 -.10 -1.29
감독 .06 .05 .08 1.35
보수 -.14 .07 -.18 -2.01*
대인관계 .06 .06 .06 .89
작업환경 .33 .09 .34 3.82***
안정성 .17 .05 .26 3.42**
성별 -.05 .09 -.03 -.52
연령 -.20 .07 -.22 -3.01**
학력 .05 .07 .04 .70
소득 .00 .04 .01 .12
종사기간 .17 .06 .23 2.71**
기관규모 -.10 .03 -.20 -3.36**"


text2df(reg_text)

text2df(reg_text) %>%t_sep()


## End(Not run)

\end{ExampleCode}
\end{Examples}
\inputencoding{utf8}
\HeaderA{t\_test\_report}{Advanceed t test report function}{t.Rul.test.Rul.report}
%
\begin{Description}
Advanceed t test report function
\end{Description}
%
\begin{Usage}
\begin{verbatim}
t_test_report(
  data,
  iv = NULL,
  dv = NULL,
  type = "all",
  xlab = "independent Variable",
  ylab = "dependent Variable",
  translate = FALSE,
  mu = 0,
  paired = FALSE,
  conf.level = 0.95
)
\end{verbatim}
\end{Usage}
%
\begin{Arguments}
\begin{ldescription}
\item[\code{data}] data.frame

\item[\code{iv}] independent variable, if you want to change the name, you can change the setting

\item[\code{dv}] dependent variable, if you want to change the name, you can change the setting

\item[\code{type}] type option default is 'all'. and  'var.test', 'var.test.report', 't.test','t.test.report', 'boxplot', 'var.test.full', 't.test.full', 'descriptive' output each result

\item[\code{xlab}] graph name dependent variable

\item[\code{ylab}] graph name dependent variable

\item[\code{translate}] report translate defualt FALSE, TRUE is korean

\item[\code{mu}] a number indicating the true value of the mean (or difference in means if you are performing a two sample test).

\item[\code{paired}] a logical indicating whether you want a paired t-test.
var.equal a logical variable indicating whether to treat the two variances as being equal. If TRUE then the pooled variance is used to estimate the variance otherwise the Welch (or Satterthwaite) approximation to the degrees of freedom is used.

\item[\code{conf.level}] confidence level of the interval.
\end{ldescription}
\end{Arguments}
%
\begin{Examples}
\begin{ExampleCode}
# mtcars data
## Not run: 
##default analysis
t_test_report(mtcars, iv="vs",dv="mpg")
##$var_test_report

##
##$t_test_result_tibble
## A tibble: 1 × 7
##IV    DV       df t.value   p.value sig   test_type
##<chr> <chr> <dbl>   <dbl>     <dbl> <chr> <chr>
##  1 vs    mpg      30   -4.86 0.0000342 ***   student t.test
##
##$t_test_report
## New namaing xlab and ylab
t_test_report(mtcars, iv="vs",dv="mpg",xlab="transmission", ylab= "males per galon")

## End(Not run)
\end{ExampleCode}
\end{Examples}
\inputencoding{utf8}
\HeaderA{unCombine}{Functions to uncombind combind\_data again}{unCombine}
%
\begin{Description}
Functions to uncombind combind\_data again
\end{Description}
%
\begin{Usage}
\begin{verbatim}
unCombine(
  data,
  into = c("value", "sig"),
  sep = " ",
  type = "all",
  pattern = NULL,
  imp = ""
)
\end{verbatim}
\end{Usage}
%
\begin{Arguments}
\begin{ldescription}
\item[\code{data}] conbind data

\item[\code{into}] column name

\item[\code{sep}] separate sep= " "

\item[\code{type}] "all", res1 = m1, res2=m2

\item[\code{pattern}] remove tail

\item[\code{imp}] imputation tail
\end{ldescription}
\end{Arguments}
%
\begin{Value}
results  two data
\end{Value}
%
\begin{Examples}
\begin{ExampleCode}
## Not run: 

Generate example data2
combined_data2 <- data.frame(
  "H" = c("1.0587162 ***", "1.1269248 ***", "0.9163776 ***", "0.9530237 ***"),
  "H(H)" = c("1.1477347 ***", "0.5743712 ***", "1.4060606 ***", "0.8624314 ***"),
  "L" = c("0.50877193 ***", "1.94957983 *", "0.03515152 ***", "1.44531611 **")
)
# Generate example data2
combined_data <- data.frame(
  "H" = c("1.0587162(***)", "1.1269248(***)", "0.9163776(***)", "0.9530237(***)"),
  "H(H)" = c("1.1477347(***)", "0.5743712(***)", "1.4060606(***)", "0.8624314(***)"),
  "L" = c("0.50877193(***)", "1.94957983(*)", "0.03515152(***)", "1.44531611(**)")
)

combined_data2
unCombind(combined_data2,sep=" ")

combined_data
unCombind(combined_data, sep="\\(", pattern=")")

#iris cor
iris_cor = iris %>% select(-Species) %>%psych::corr.test()
iris_cor$r
iris_cor$p

combine_data(iris_cor$r%>%round(2),
             iris_cor$p%>%round(5),
             left="(", right=")")
 sig = add_significance_symbols(iris_cor$p)
 sig

iris_cor_combine =combine_data( round(iris_cor$r, 3) , sig)
iris_cor_combine

 combine_data( round(iris_cor$r, 3) , sig) %>%
 lowerMat(diag = 1, fill="")







## End(Not run)
\end{ExampleCode}
\end{Examples}
\inputencoding{utf8}
\HeaderA{uncombine\_data}{Functions to uncombind combind\_data again}{uncombine.Rul.data}
%
\begin{Description}
Functions to uncombind combind\_data again
\end{Description}
%
\begin{Usage}
\begin{verbatim}
uncombine_data(
  data,
  into = c("value", "sig"),
  sep = " ",
  type = "res",
  pattern = NULL,
  imp = ""
)
\end{verbatim}
\end{Usage}
%
\begin{Arguments}
\begin{ldescription}
\item[\code{data}] conbind data

\item[\code{into}] column name

\item[\code{sep}] separate sep= " "

\item[\code{type}] "all", res1 = m1, res2=m2

\item[\code{pattern}] remove tail

\item[\code{imp}] imputation tail
\end{ldescription}
\end{Arguments}
%
\begin{Value}
results  two data
\end{Value}
%
\begin{Examples}
\begin{ExampleCode}
## Not run: 

Generate example data2
combined_data2 <- data.frame(
  "H" = c("1.0587162 ***", "1.1269248 ***", "0.9163776 ***", "0.9530237 ***"),
  "H(H)" = c("1.1477347 ***", "0.5743712 ***", "1.4060606 ***", "0.8624314 ***"),
  "L" = c("0.50877193 ***", "1.94957983 *", "0.03515152 ***", "1.44531611 **")
)
# Generate example data2
combined_data <- data.frame(
  "H" = c("1.0587162(***)", "1.1269248(***)", "0.9163776(***)", "0.9530237(***)"),
  "H(H)" = c("1.1477347(***)", "0.5743712(***)", "1.4060606(***)", "0.8624314(***)"),
  "L" = c("0.50877193(***)", "1.94957983(*)", "0.03515152(***)", "1.44531611(**)")
)

combined_data2
unCombind(combined_data2,sep=" ")

combined_data
unCombind(combined_data, sep="\\(", pattern=")")



## End(Not run)
\end{ExampleCode}
\end{Examples}
\inputencoding{utf8}
\HeaderA{unCount}{Functions that generate rows with a frequency}{unCount}
%
\begin{Description}
Functions that generate rows with a frequency
\end{Description}
%
\begin{Usage}
\begin{verbatim}
unCount(data, type = "real", sel = ncol(data), fix = NULL)
\end{verbatim}
\end{Usage}
%
\begin{Arguments}
\begin{ldescription}
\item[\code{data}] data

\item[\code{type}] res, all

\item[\code{sel}] frequency column

\item[\code{fix}] final select column
\end{ldescription}
\end{Arguments}
%
\begin{Value}
data
\end{Value}
%
\begin{Examples}
\begin{ExampleCode}

## Not run: 
dff = data.frame(
  trt = c("r","s","z"),
  con = c("f","m","l"),
  num = c(10,0,2),
  fre = c(4,0,5) )
dff
dff %>%unCount()
dff %>%unCount()
dff %>%unCount(fix=1)
dff %>%unCount(type= "res", fix=1:2)
dff %>%unCount(type= "res", fix="con")
dff %>%unCount(type= "all")
dff %>%unCount(type= "real")
dff %>%unCount(type= "zero")
dff %>%unCount(type= "zero", fix="con")

# Making frequency data into real data
dff %>%unCount()%>% add_rows_freq()
 dff_real =  dff %>%unCount()%>% add_rows_freq()
 dff_real

 freq_data =  mtcars %>% select(am, vs) %>% table()
 freq_data%>% data.frame()
 freq_data%>% data.frame() %>% unCount(type ="zero")
 freq_data%>% data.frame() %>% unCount(type ="zero") %>%table()
rowid_to_column

## End(Not run)

\end{ExampleCode}
\end{Examples}
\inputencoding{utf8}
\HeaderA{Unite}{UInte data}{Unite}
%
\begin{Description}
UInte data
\end{Description}
%
\begin{Usage}
\begin{verbatim}
Unite(data, col1, col2, Col = NULL, sep = "", remove = TRUE)
\end{verbatim}
\end{Usage}
%
\begin{Arguments}
\begin{ldescription}
\item[\code{data}] data

\item[\code{col1}] select col1

\item[\code{col2}] select col2

\item[\code{Col}] col names

\item[\code{sep}] ""

\item[\code{remove}] true

\item[\code{left}] " or "["

\item[\code{right}] ""  or  "]"
\end{ldescription}
\end{Arguments}
%
\begin{Value}
data
\end{Value}
%
\begin{Examples}
\begin{ExampleCode}
## Not run: 
#'
# 예시 데이터 생성
data <- data.frame(
  block = c("IMAG", "IMAG", "IMAG", "IMAG"),
  name = c("imag1", "imag2", "imag3", "imag4"),
  Original = c("0.7093265 ***", "0.8773078 ***", "0.8417485 ***", "0.5691934 ***"),
  sig = c("***", "***", "***", "***"),
  Mean.Boot = c(0.7088215, 0.8750787, 0.8401757, 0.5647333),
  Std.Error = c(0.05949724, 0.02402834, 0.02970117, 0.07686419),
  perc.025 = c(0.5774529, 0.8231642, 0.7784983, 0.3876417),
  perc.975 = c(0.8105249, 0.9146114, 0.8863100, 0.6961572)
)

# 예시 호출

# Unite using column numbers
result1 <- Unite(data, 3, 4)
print(result1)

# Unite using column numbers with a new column name
result2 <- Unite(data, 3, 4, "New")
print(result2)

# Unite using column names
result3 <- Unite(data, "Original", "sig")
print(result3)

# Unite using column names with a new column name
result4 <- Unite(data, "Original", "sig", "New")
print(result4)



## End(Not run)
#
\end{ExampleCode}
\end{Examples}
\inputencoding{utf8}
\HeaderA{unite\_ci}{unite\_ci if you add ci}{unite.Rul.ci}
%
\begin{Description}
unite\_ci if you add ci
\end{Description}
%
\begin{Usage}
\begin{verbatim}
unite_ci(
  data,
  col1 = ncol(data) - 1,
  col2 = ncol(data),
  left = "[",
  right = "]",
  digits = 3,
  colname = "95%CI",
  remove = TRUE
)
\end{verbatim}
\end{Usage}
%
\begin{Arguments}
\begin{ldescription}
\item[\code{data}] A data frame containing the data

\item[\code{col1}] lower

\item[\code{col2}] upper

\item[\code{left}] lefe (

\item[\code{right}] right )

\item[\code{digits}] round3

\item[\code{colname}] "95CI"

\item[\code{remove}] orinal column
\end{ldescription}
\end{Arguments}
%
\begin{Value}
data
\end{Value}
%
\begin{Examples}
\begin{ExampleCode}

## Not run: 


# Sample data to simulate the scenario
data <- data.frame(
  paths = c('자기효능감-C_S1', '자기효능감-C_S2', '자기효능감-C_S3'),
  Original = c(0.883, 0.826, 0.794),
  Mean.Boot = c(0.885, 0.833, 0.8),
  Std.Error = c(0.022, 0.041, 0.039),
  t = c('40.771***', '20.22***', '20.612***'),
  perc.025 = c(0.837, 0.733, 0.725),
  perc.975 = c(0.923, 0.893, 0.872)
)

# Apply the function to the dataframe with remove = TRUE
unite_ci(data, col1 = 6, col2 = 7, remove = TRUE)

# Apply the function to the dataframe with remove = FALSE
unite_ci(data, col1 = 'perc.025', col2 = 'perc.975', remove = FALSE)


## End(Not run)

\end{ExampleCode}
\end{Examples}
\inputencoding{utf8}
\HeaderA{Unite\_col}{Unite Columns in a Data Frame}{Unite.Rul.col}
%
\begin{Description}
Unite Columns in a Data Frame
\end{Description}
%
\begin{Usage}
\begin{verbatim}
Unite_col(
  data,
  col1,
  col2,
  Col = NULL,
  sep = ",",
  remove = TRUE,
  left = "[",
  right = "]"
)
\end{verbatim}
\end{Usage}
%
\begin{Arguments}
\begin{ldescription}
\item[\code{data}] A data frame.

\item[\code{col1}] The first column to unite. Can be numeric (column index) or character (column name).

\item[\code{col2}] The second column to unite. Can be numeric (column index) or character (column name).

\item[\code{Col}] Optional, name for the new column (default = NULL).

\item[\code{sep}] Separator string to use between the combined values (default = ",").

\item[\code{remove}] Logical, whether to remove the original columns after uniting (default = TRUE).

\item[\code{left}] String to add to the left side of the combined values (default = "[").

\item[\code{right}] String to add to the right side of the combined values (default = "]").
\end{ldescription}
\end{Arguments}
%
\begin{Value}
A data frame with the combined column.
\end{Value}
%
\begin{Examples}
\begin{ExampleCode}
## Not run: 
# Example data
data <- data.frame(
  block = c("IMAG", "IMAG", "IMAG", "IMAG"),
  name = c("imag1", "imag2", "imag3", "imag4"),
  Original = c("0.7093265 ***", "0.8773078 ***", "0.8417485 ***", "0.5691934 ***"),
  sig = c("***", "***", "***", "***"),
  Mean.Boot = c(0.7088215, 0.8750787, 0.8401757, 0.5647333),
  Std.Error = c(0.05949724, 0.02402834, 0.02970117, 0.07686419),
  perc.025 = c(0.5774529, 0.8231642, 0.7784983, 0.3876417),
  perc.975 = c(0.8105249, 0.9146114, 0.8863100, 0.6961572)
)

# Unite using column numbers
result1 <- Unite_col(data, 3, 4)
print(result1)

# Unite using column numbers with a new column name
result2 <- Unite_col(data, 3, 4, "New")
print(result2)

# Unite using column names
result3 <- Unite_col(data, "Original", "sig")
print(result3)

# Unite using column names with a new column name
result4 <- Unite_col(data, "Original", "sig", "New")
print(result4)

## End(Not run)
\end{ExampleCode}
\end{Examples}
\inputencoding{utf8}
\HeaderA{Unite\_confint}{Unite Columns Confint in a Data Frame}{Unite.Rul.confint}
%
\begin{Description}
Unite Columns Confint in a Data Frame
\end{Description}
%
\begin{Usage}
\begin{verbatim}
Unite_confint(
  data,
  col1,
  col2,
  Col = NULL,
  sep = ", ",
  remove = TRUE,
  left = "[",
  right = "]"
)
\end{verbatim}
\end{Usage}
%
\begin{Arguments}
\begin{ldescription}
\item[\code{data}] A data frame.

\item[\code{col1}] The first column to unite. Can be numeric (column index) or character (column name).

\item[\code{col2}] The second column to unite. Can be numeric (column index) or character (column name).

\item[\code{Col}] Optional, name for the new column (default = NULL).

\item[\code{sep}] Separator string to use between the combined values (default = ",").

\item[\code{remove}] Logical, whether to remove the original columns after uniting (default = TRUE).

\item[\code{left}] String to add to the left side of the combined values (default = "[").

\item[\code{right}] String to add to the right side of the combined values (default = "]").
\end{ldescription}
\end{Arguments}
%
\begin{Value}
A data frame with the combined column.
\end{Value}
%
\begin{Examples}
\begin{ExampleCode}
## Not run: 
# Example data
data <- data.frame(
  block = c("IMAG", "IMAG", "IMAG", "IMAG"),
  name = c("imag1", "imag2", "imag3", "imag4"),
  Original = c("0.7093265 ***", "0.8773078 ***", "0.8417485 ***", "0.5691934 ***"),
  sig = c("***", "***", "***", "***"),
  Mean.Boot = c(0.7088215, 0.8750787, 0.8401757, 0.5647333),
  Std.Error = c(0.05949724, 0.02402834, 0.02970117, 0.07686419),
  perc.025 = c(0.5774529, 0.8231642, 0.7784983, 0.3876417),
  perc.975 = c(0.8105249, 0.9146114, 0.8863100, 0.6961572)
)

# Unite using column numbers
result1 <- Unite_confint(data, 3, 4)
print(result1)

# Unite using column numbers with a new column name
result2 <- Unite_confint(data, 3, 4, "New")
print(result2)

# Unite using column names
result3 <- Unite_confint(data, "Original", "sig")
print(result3)

# Unite using column names with a new column name
result4 <- Unite_confint(data, "Original", "sig", "New")
print(result4)

## End(Not run)
\end{ExampleCode}
\end{Examples}
\inputencoding{utf8}
\HeaderA{unite\_rows}{행을 결합해주는 함수}{unite.Rul.rows}
%
\begin{Description}
행을 결합해주는 함수
\end{Description}
%
\begin{Usage}
\begin{verbatim}
unite_rows(data, row1, row2, sep = ", ", left = "[", right = "]")
\end{verbatim}
\end{Usage}
%
\begin{Arguments}
\begin{ldescription}
\item[\code{data}] data

\item[\code{row1}] row1 합쳐질 곳

\item[\code{row2}] row2 가져올 곳

\item[\code{sep}] sep "," default

\item[\code{left}] left "["

\item[\code{right}] right"]"
\end{ldescription}
\end{Arguments}
%
\begin{Value}
data
\end{Value}
%
\begin{Examples}
\begin{ExampleCode}

## Not run: 
#'
data17 <- data.frame(
  criterian = c("Chisq", "df", "p >0.05", "RMSEA <0.05",
                "90%CI.lower", "90%CI.upper", "p <= 0.05",
                "SRMR <0.08", "GFI >0.95", "CFI >0.95", "TLI >0.90", "lower", "lower"),
  Value = c(40.179, 38.000, 0.374, 0.028, 0.000, 0.087,
            0.665, 0.056, 0.920, 0.997, 0.995, 3153.636, 3218.526)
)
data17
# 함수 실행 예시
unite_rows(data17, row1 = 5, row2 = 6)
unite_rows(data17, row1 = 5, row2 = 6)%>%move_row(6,5)

mtcars[1:5,]
mtcars[1:5,] %>%unite_rows(1,3)

## End(Not run)



\end{ExampleCode}
\end{Examples}
\inputencoding{utf8}
\HeaderA{validate\_residual\_cov}{Validate Residual Covariance Specification}{validate.Rul.residual.Rul.cov}
%
\begin{Description}
관측변수 간 잔차공분산 설정의 유효성을 검사한다.
\end{Description}
%
\begin{Usage}
\begin{verbatim}
validate_residual_cov(residual_cov, obs_vars)
\end{verbatim}
\end{Usage}
%
\begin{Arguments}
\begin{ldescription}
\item[\code{residual\_cov}] list. c("y1","y2") 또는 c("y1","y2","0").

\item[\code{obs\_vars}] character. 관측변수 이름 벡터.
\end{ldescription}
\end{Arguments}
%
\begin{Value}
list with elements \code{pairs}, \code{fixed}.
\end{Value}
\inputencoding{utf8}
\HeaderA{validate\_variance}{Validate Variance Constraints}{validate.Rul.variance}
%
\begin{Description}
관측·잠재변수 분산 고정 또는 자유 설정을 검증한다.
\end{Description}
%
\begin{Usage}
\begin{verbatim}
validate_variance(variance, obs_vars, latent_vars)
\end{verbatim}
\end{Usage}
%
\begin{Arguments}
\begin{ldescription}
\item[\code{variance}] list. c("x","x") 또는 c("f","f","1").

\item[\code{obs\_vars}] character. 관측변수 이름.

\item[\code{latent\_vars}] character. 잠재변수 이름.
\end{ldescription}
\end{Arguments}
%
\begin{Value}
list with elements \code{free}, \code{fixed}.
\end{Value}
\inputencoding{utf8}
\HeaderA{wald\_test}{wald\_test}{wald.Rul.test}
%
\begin{Description}
wald\_test
\end{Description}
%
\begin{Usage}
\begin{verbatim}
wald_test(
  b1,
  se1,
  b2,
  se2,
  est1 = NULL,
  est2 = NULL,
  digits = 3,
  type = "res2"
)
\end{verbatim}
\end{Usage}
%
\begin{Arguments}
\begin{ldescription}
\item[\code{b1}] fires coef

\item[\code{se1}] firtse se

\item[\code{b2}] secoond coef

\item[\code{se2}] second se

\item[\code{est1}] est1 name

\item[\code{est2}] est2 name

\item[\code{digits}] round

\item[\code{type}] res1 res2,
\end{ldescription}
\end{Arguments}
%
\begin{Value}
test result
\end{Value}
%
\begin{Examples}
\begin{ExampleCode}
## Not run: 
##Clogg, C. C., Petkova, E., & Haritou, A. (1995).
Statistical methods for comparing regression coefficients between models.
American Journal of Sociology, 100(5), 1261-1293.
##https://blog.naver.com/shoutjoy/223294372137

 #' wald_test(-0.795,-0.795/-1.138, -0.795+0.149,  (-0.795+0.149)/1.344)

wald_test(-0.795,-0.795/-1.138,-0.795+0.149,
          (-0.795+0.149)/1.344,
           est1 = "남성",
           est2 = "여성")


## End(Not run)


wald_test <- function(b1, se1, b2, se2, est1=NULL, est2=NULL){
  library(broom)
  t1 = b1/se1
  t2 = b2/se2

  z = (b1 - b2)/ sqrt(se1^2 + se2^2)
  p = 2*(1-pnorm(abs(z)))

  var1 = paste0(est1,"_est1 = ", b1,", se = ", round(se1, 3), ", t = ", t1,"." )
  var2 = paste0(est2,"_est2 = ", b2,", se = ", round(se2, 3), ", t = ", t2,"." )

  res = broom::tidy(c(z.value = z,  pvalue = p))
  colnames(res)=c("statistics","value")

  res = res%>% dplry::mutate(coef=c(var1, var2))
  res
  # list(res, var1, var2)
}


\end{ExampleCode}
\end{Examples}
\inputencoding{utf8}
\HeaderA{wald\_test2}{wald\_test2 input t value}{wald.Rul.test2}
%
\begin{Description}
wald\_test2 input t value
\end{Description}
%
\begin{Usage}
\begin{verbatim}
wald_test2(
  b1,
  t1,
  b2,
  t2,
  est1 = NULL,
  est2 = NULL,
  gender = FALSE,
  digits = 3,
  type = "res1"
)
\end{verbatim}
\end{Usage}
%
\begin{Arguments}
\begin{ldescription}
\item[\code{b1}] fires coef

\item[\code{t1}] firtse t

\item[\code{b2}] secoond coef

\item[\code{t2}] second t

\item[\code{est1}] est1 name

\item[\code{est2}] est2 name

\item[\code{gender}] For dichotomous variables with intercepts, the Gender

\item[\code{digits}] round

\item[\code{type}] res=res1. res2 = all
\end{ldescription}
\end{Arguments}
%
\begin{Value}
ztest
\end{Value}
%
\begin{Examples}
\begin{ExampleCode}
## Not run: 

wald_test2(-0.795,8.245, -0.795+0.149,  (-0.795+0.149)/1.344)

wald_test2(-0.795,-0.795/-1.138,-0.795+0.149,
           (-0.795+0.149)/1.344,est1="남성",est2="여성")


##gender variable
##first method
wald_test2(1.311, 1.628, 1.311-0.366,  -2.732)
##second method
wald_test2(1.311, 1.628, -0.366,  -2.732, gender=T)
wald_test2(1.311, 1.628, -0.366,  -2.732,"male","female", gender=T)


## End(Not run)



\end{ExampleCode}
\end{Examples}
\inputencoding{utf8}
\HeaderA{wald\_test\_df}{wald\_test\_df: Wald test for linear model coefficients using a data frame or tibble}{wald.Rul.test.Rul.df}
%
\begin{Description}
This function performs a Wald test to compare two coefficients from a linear model.
It calculates the difference, z-statistic, and p-value between the selected coefficients.
\end{Description}
%
\begin{Usage}
\begin{verbatim}
wald_test_df(data, row1 = 2, row2 = 3, rname1 = "", rname2 = "", sel = "row")
\end{verbatim}
\end{Usage}
%
\begin{Arguments}
\begin{ldescription}
\item[\code{data}] An object of class `lm`, representing a fitted linear model.

\item[\code{row1}] Integer. The row number of the first coefficient to compare (default = 2).

\item[\code{row2}] Integer. The row number of the second coefficient to compare (default = 3).

\item[\code{rname1}] Character. Name of the first term to select when `sel = "term"`. Defaults to "".

\item[\code{rname2}] Character. Name of the second term to select when `sel = "term"`. Defaults to "".

\item[\code{sel}] Character. Method for selecting coefficients: "row" (default) for row numbers,
or "term" for term names.
\end{ldescription}
\end{Arguments}
%
\begin{Value}
A tibble containing the estimated coefficients, standard errors, difference, z-statistic,
and p-value.
\end{Value}
%
\begin{Examples}
\begin{ExampleCode}
## Not run: 
# Using default row selection
lm(mpg ~ hp + wt + drat, mtcars) %>% wald_test_df()

# Selecting specific rows by number
lm(mpg ~ hp + wt + drat, mtcars) %>% wald_test_df(row1 = 2, row2 = 4)

# Selecting specific terms by name
lm(mpg ~ hp + wt + drat, mtcars) %>% wald_test_df(sel = "term", rname1 = "hp", rname2 = "wt")

## End(Not run)
\end{ExampleCode}
\end{Examples}
\inputencoding{utf8}
\HeaderA{wald\_test\_gender}{wald\_test2 input t value gender}{wald.Rul.test.Rul.gender}
%
\begin{Description}
wald\_test2 input t value gender
\end{Description}
%
\begin{Usage}
\begin{verbatim}
wald_test_gender(
  b1,
  t1,
  b2,
  t2,
  est1 = NULL,
  est2 = NULL,
  gender = TRUE,
  digits = 3
)
\end{verbatim}
\end{Usage}
%
\begin{Arguments}
\begin{ldescription}
\item[\code{b1}] fires coef

\item[\code{t1}] firtse t

\item[\code{b2}] secoond coef

\item[\code{t2}] second t

\item[\code{est1}] est1 name

\item[\code{est2}] est2 name

\item[\code{gender}] For dichotomous variables with intercepts, the Gender

\item[\code{digits}] round
\end{ldescription}
\end{Arguments}
%
\begin{Value}
ztest
\end{Value}
%
\begin{Examples}
\begin{ExampleCode}
## Not run: 

wald_test2(-0.795,8.245, -0.795+0.149,  (-0.795+0.149)/1.344)

#Regression: gender variable should consider intercept.
wald_test2(-0.795,-0.795/-1.138,-0.795+0.149,
           (-0.795+0.149)/1.344,est1="남성",est2="여성")


##gender variable
##first method
wald_test_gender(1.311, 1.628, 1.311-0.366,  -2.732)
##second method
wald_test_gender(1.311, 1.628, -0.366,  -2.732, gender=T)
wald_test_gender(1.311, 1.628, -0.366,  -2.732,"male","female", gender=T)


## End(Not run)



\end{ExampleCode}
\end{Examples}
\inputencoding{utf8}
\HeaderA{wald\_test\_general}{General Wald Test for Linear Hypotheses}{wald.Rul.test.Rul.general}
%
\begin{Description}
This function performs a Wald test for general linear hypotheses on a fitted model.
The hypothesis can be specified as a string (e.g., `"x1 - x2 + x3 = 0"`) or as
a numeric vector/matrix (e.g., `c(1, -1, 1)`). The function provides an option
for detailed output using the `full` parameter.
\end{Description}
%
\begin{Usage}
\begin{verbatim}
wald_test_general(
  model,
  hypothesis_matrix,
  rhs = 0,
  test = c("F", "Chisq"),
  vcov_func = NULL,
  full = FALSE
)
\end{verbatim}
\end{Usage}
%
\begin{Arguments}
\begin{ldescription}
\item[\code{model}] A fitted model object, such as an object from `lm` or `glm`.

\item[\code{hypothesis\_matrix}] A hypothesis specified as either a character string (e.g., `"x1 - x2 + x3 = 0"`)
or a numeric vector/matrix (e.g., `c(1, -1, 1)`).

\item[\code{rhs}] Right-hand side of the hypothesis equation. Default is `0`.

\item[\code{test}] A character string specifying the type of test to perform. Options are `"F"` (default) or `"Chisq"`.

\item[\code{vcov\_func}] An optional function for specifying a custom covariance matrix. Default is `NULL`.

\item[\code{full}] Logical. If `TRUE`, returns a tibble with full results. If `FALSE`, prints summary information
and returns a simplified tibble. Default is `FALSE`.
\end{ldescription}
\end{Arguments}
%
\begin{Value}
A tibble with results. If `full = TRUE`, the tibble includes:
- `var\_comparison`: The hypothesis expressed as a standardized string.
- `W`: The Wald statistic.
- `df`: Degrees of freedom for the test.
- `p\_value`: The p-value associated with the test.
- `method`: The type of test performed (`"F"` or `"Chisq"`).
- `mat`: A vectorized representation of the hypothesis matrix.
If `full = FALSE`, the tibble includes only:
- `var\_comparison`: The hypothesis expressed as a standardized string.
- `W`: The Wald statistic.
- `df`: Degrees of freedom for the test.
- `p\_value`: The p-value associated with the test.
\end{Value}
%
\begin{Examples}
\begin{ExampleCode}
# Example: Linear model with Wald test
set.seed(123)
data <- data.frame(x1 = rnorm(100), x2 = rnorm(100), x3 = rnorm(100), y = rnorm(100))
model <- lm(y ~ x1 + x2 + x3, data = data)

# Perform Wald test with text-based hypothesis
result1 <- wald_test_general(model, "x1 - x2 + x3 = 0", test = "F", full = TRUE)
print(result1)

# Perform Wald test with matrix-based hypothesis
result2 <- wald_test_general(model, c(1, -1, 1), test = "F", full = TRUE)
print(result2)

# Perform Wald test with custom covariance matrix
custom_vcov <- function(model) {
  vcov_matrix <- vcov(model)
  # Example: Inflate variance for demonstration
  vcov_matrix * 1.5
}
result_custom <- wald_test_general(model, c(1, -1, 1), test = "F", vcov_func = custom_vcov, full = TRUE)
print(result_custom)

\end{ExampleCode}
\end{Examples}
\inputencoding{utf8}
\HeaderA{wald\_test\_general\_df}{Perform Wald Tests for Multiple Hypotheses}{wald.Rul.test.Rul.general.Rul.df}
%
\begin{Description}
This function allows multiple hypotheses to be tested on a fitted model. It internally
uses `wald\_test\_general` to perform each test and combines the results into a single
tibble.
\end{Description}
%
\begin{Usage}
\begin{verbatim}
wald_test_general_df(
  model,
  ...,
  rhs = 0,
  test = c("F", "Chisq"),
  vcov_func = NULL,
  full = FALSE
)
\end{verbatim}
\end{Usage}
%
\begin{Arguments}
\begin{ldescription}
\item[\code{model}] A fitted model object, such as an object from `lm` or `glm`.

\item[\code{...}] Hypotheses to test. Each hypothesis can be specified as a character string
(e.g., `"x1 - x2 + x3 = 0"`) or a numeric vector/matrix (e.g., `c(1, -1, 1)`).

\item[\code{rhs}] Right-hand side of the hypothesis equations. Default is `0`.

\item[\code{test}] A character string specifying the type of test to perform. Options are `"F"`
(default) or `"Chisq"`.

\item[\code{vcov\_func}] An optional function for specifying a custom covariance matrix. Default is `NULL`.

\item[\code{full}] Logical. If `TRUE`, returns a tibble with full results. Default is `FALSE`.
\end{ldescription}
\end{Arguments}
%
\begin{Value}
A tibble with results for all tested hypotheses. Each row corresponds to one hypothesis.
\end{Value}
%
\begin{Examples}
\begin{ExampleCode}
# Example: Multiple hypotheses with a linear model
set.seed(123)
data <- data.frame(x1 = rnorm(100), x2 = rnorm(100), x3 = rnorm(100), y = rnorm(100))
model <- lm(y ~ x1 + x2 + x3, data = data)

# Perform multiple Wald tests
result <- wald_test_general_df(model, "x1 - x2 + x3 = 0", "x1 - x2 = 0", c(1, -1, 1), test = "F", full = TRUE)
print(result)

\end{ExampleCode}
\end{Examples}
\inputencoding{utf8}
\HeaderA{wald\_test\_general\_mat}{Wald Test for Hypotheses Specified as Matrices}{wald.Rul.test.Rul.general.Rul.mat}
%
\begin{Description}
This function performs a Wald test for general linear hypotheses on a fitted model.
The hypothesis is specified as a numeric vector or matrix. The function computes
the Wald statistic and associated p-value based on the input hypothesis matrix.
\end{Description}
%
\begin{Usage}
\begin{verbatim}
wald_test_general_mat(
  model,
  hypothesis_matrix,
  rhs = 0,
  test = c("F", "Chisq"),
  vcov_func = NULL
)
\end{verbatim}
\end{Usage}
%
\begin{Arguments}
\begin{ldescription}
\item[\code{model}] A fitted model object, such as an object from `lm` or `glm`.

\item[\code{hypothesis\_matrix}] A numeric vector or matrix representing the hypothesis to test.
For example, `c(1, -1, 1)` represents the hypothesis `x1 - x2 + x3 = 0`.

\item[\code{rhs}] A numeric value for the right-hand side of the hypothesis equation. Default is `0`.

\item[\code{test}] A character string specifying the type of test to perform. Options are `"F"`
(default) or `"Chisq"`.

\item[\code{vcov\_func}] An optional function for specifying a custom covariance matrix. Default is `NULL`.
\end{ldescription}
\end{Arguments}
%
\begin{Value}
A tibble with the following columns:
- `var\_comparison`: A standardized version of the hypothesis in text form.
- `W`: The Wald statistic.
- `df`: Degrees of freedom for the test.
- `p\_value`: The p-value associated with the test.
- `method`: The type of test performed (`"F"` or `"Chisq"`).
- `mat`: A vectorized representation of the hypothesis matrix.
\end{Value}
%
\begin{Examples}
\begin{ExampleCode}
# Example: Linear model with Wald test
set.seed(123)
data <- data.frame(x1 = rnorm(100), x2 = rnorm(100), x3 = rnorm(100), y = rnorm(100))
model <- lm(y ~ x1 + x2 + x3, data = data)

# Perform Wald test with matrix-based hypothesis
result <- wald_test_general_mat(model, c(1, -1, 1), test = "F")
print(result)

# Perform Wald test with a custom covariance matrix
custom_vcov <- function(model) {
  vcov_matrix <- vcov(model)
  # Example: Inflate variance for demonstration
  vcov_matrix * 1.5
}
result_custom <- wald_test_general_mat(model, c(1, -1, 1), test = "F", vcov_func = custom_vcov)
print(result_custom)

\end{ExampleCode}
\end{Examples}
\inputencoding{utf8}
\HeaderA{wald\_test\_general\_str}{Wald Test for Hypotheses Specified as Strings}{wald.Rul.test.Rul.general.Rul.str}
%
\begin{Description}
This function performs a Wald test for general linear hypotheses on a fitted model.
The hypothesis is specified as a string (e.g., `"x1 - x2 + x3 = 0"`) and is converted
into a matrix internally for computation. This allows for flexible specification of
hypotheses in an intuitive text-based format.
\end{Description}
%
\begin{Usage}
\begin{verbatim}
wald_test_general_str(
  model,
  hypothesis,
  rhs = 0,
  test = c("F", "Chisq"),
  vcov_func = NULL
)
\end{verbatim}
\end{Usage}
%
\begin{Arguments}
\begin{ldescription}
\item[\code{model}] A fitted model object, such as an object from `lm` or `glm`.

\item[\code{hypothesis}] A character string representing the hypothesis to test.
For example, `"x1 - x2 + x3 = 0"`. The right-hand side (`rhs`) is
automatically incorporated into the hypothesis.

\item[\code{rhs}] A numeric value for the right-hand side of the hypothesis equation.
Default is `0`.

\item[\code{test}] A character string specifying the type of test to perform. Options
are `"F"` (default) or `"Chisq"`.

\item[\code{vcov\_func}] An optional function for specifying a custom covariance matrix.
Default is `NULL`, in which case the model's default covariance matrix is used.
\end{ldescription}
\end{Arguments}
%
\begin{Value}
A tibble with the following columns:
- `var\_comparison`: A standardized version of the hypothesis in text form.
- `W`: The Wald statistic.
- `df`: Degrees of freedom for the test.
- `p\_value`: The p-value associated with the test.
- `method`: The type of test performed (`"F"` or `"Chisq"`).
- `mat`: A vectorized representation of the hypothesis matrix.
\end{Value}
%
\begin{Examples}
\begin{ExampleCode}
# Example: Linear model with Wald test
set.seed(123)
data <- data.frame(x1 = rnorm(100), x2 = rnorm(100), x3 = rnorm(100), y = rnorm(100))
model <- lm(y ~ x1 + x2 + x3, data = data)

# Perform Wald test with text-based hypothesis
result <- wald_test_general_str(model, "x1 - x2 + x3 = 0", test = "F")
print(result)

\end{ExampleCode}
\end{Examples}
\inputencoding{utf8}
\HeaderA{wald\_test\_lm}{Pairwise Wald Test for Linear Model Coefficients}{wald.Rul.test.Rul.lm}
%
\begin{Description}
This function performs pairwise Wald tests for coefficients of a linear model.
\end{Description}
%
\begin{Usage}
\begin{verbatim}
wald_test_lm(input_data, raw = FALSE)
\end{verbatim}
\end{Usage}
%
\begin{Arguments}
\begin{ldescription}
\item[\code{input\_data}] An object of type `lm`, `summary` of an `lm` object, or a tidy dataframe
containing `term`, `estimate`, and `std.error` columns.

\item[\code{raw}] Logical. If `TRUE`, returns raw results. If `FALSE`, returns formatted results with significance markings.
\end{ldescription}
\end{Arguments}
%
\begin{Value}
A dataframe containing pairwise Wald test results. The output includes:
- `var\_comparison`: Pairwise variable comparison.
- `est1`, `se1`, `est2`, `se2`: Estimates and standard errors of the variables being compared.
- `Z`: Wald statistic.
- `p\_value`: p-value for the Wald test.
\end{Value}
%
\begin{Examples}
\begin{ExampleCode}

## Not run: 
# Example with an lm object
model <- lm(mpg ~ hp + wt, data = mtcars)
wald_test_lm(model, raw = TRUE)
wald_test_lm(model, raw = FALSE)

# Example with a summary of lm object
model_summary <- summary(model)
wald_test_lm(model_summary, raw = TRUE)

# Example with a tidy dataframe
tidy_data <- broom::tidy(model)
wald_test_lm(tidy_data, raw = FALSE)

## End(Not run)
\end{ExampleCode}
\end{Examples}
\inputencoding{utf8}
\HeaderA{web}{Functions that output to a web browser}{web}
%
\begin{Description}
Functions that output to a web browser
\end{Description}
%
\begin{Usage}
\begin{verbatim}
web(
  data,
  out = TRUE,
  tibble = TRUE,
  digits = 3,
  font_size = 14,
  file = "output.html"
)
\end{verbatim}
\end{Usage}
%
\begin{Arguments}
\begin{ldescription}
\item[\code{data}] data

\item[\code{out}] TRUE browser

\item[\code{tibble}] show data

\item[\code{digits}] default 3

\item[\code{font\_size}] default 14

\item[\code{file}] file name default "output.html
\end{ldescription}
\end{Arguments}
%
\begin{Value}
web output
\end{Value}
%
\begin{Examples}
\begin{ExampleCode}

## Not run: 
mtcars%>%web()


## End(Not run)

\end{ExampleCode}
\end{Examples}
\inputencoding{utf8}
\HeaderA{webr}{data to web}{webr}
%
\begin{Description}
data to web
\end{Description}
%
\begin{Usage}
\begin{verbatim}
webr(
  data,
  show = TRUE,
  digits = 3,
  font_size = 18,
  border_width = 1,
  border_color = "#ccc",
  file = "output.html"
)
\end{verbatim}
\end{Usage}
%
\begin{Arguments}
\begin{ldescription}
\item[\code{data}] data.frame

\item[\code{show}] TRUE: web FALSE" console

\item[\code{digits}] round 3

\item[\code{font\_size}] 14

\item[\code{border\_width}] 1

\item[\code{border\_color}] "gray"

\item[\code{file}] file = "output.html"
\end{ldescription}
\end{Arguments}
%
\begin{Value}
web
\end{Value}
%
\begin{Examples}
\begin{ExampleCode}

## Not run: 
#'
# Example usage
webr(mtcars)
webr(mtcars, font_size=20)

## End(Not run)

\end{ExampleCode}
\end{Examples}
\inputencoding{utf8}
\HeaderA{white\_se}{White의 수정된 표준오차를 계산하는 함수 정의}{white.Rul.se}
%
\begin{Description}
White의 수정된 표준오차를 계산하는 함수 정의
\end{Description}
%
\begin{Usage}
\begin{verbatim}
white_se(original_se, n = NULL, k = 2, type = "white")
\end{verbatim}
\end{Usage}
%
\begin{Arguments}
\begin{ldescription}
\item[\code{original\_se}] original\_se

\item[\code{n}] sample size

\item[\code{k}] number of regression coefs

\item[\code{type}] all, white, roburst
\end{ldescription}
\end{Arguments}
%
\begin{Value}
se data
\end{Value}
%
\begin{Examples}
\begin{ExampleCode}
## Not run: 
#'
# 함수 사용 예시
original_se <- 0.009619  # 예시 표준오차
n <- 100  # 예시 표본 크기
k <- 2  # 예시 회귀 계수 개수 (절편 포함)

# 함수 호출 및 결과 출력
white_se(original_se= 0.009619, n= 100, k=2)

## End(Not run)

\end{ExampleCode}
\end{Examples}
\inputencoding{utf8}
\HeaderA{wide\_df}{pivot\_wider}{wide.Rul.df}
%
\begin{Description}
pivot\_wider
\end{Description}
%
\begin{Usage}
\begin{verbatim}
wide_df(
  data,
  names_from = "names",
  values_from = "values",
  names_prefix = "",
  names_sep = "_",
  names_sort = FALSE,
  rm_cols = NULL,
  add_id = NULL
)
\end{verbatim}
\end{Usage}
%
\begin{Arguments}
\begin{ldescription}
\item[\code{data}] long format data

\item[\code{names\_from}] names\_from

\item[\code{values\_from}] values\_from

\item[\code{names\_prefix}] names\_prefix +variable

\item[\code{names\_sep}] sep

\item[\code{names\_sort}] colnames sort

\item[\code{rm\_cols}] Remove variables that interfere with wide format by creating individual IDs

\item[\code{add\_id}] Select Repeating Groups to create IDs, line in column names or column numbers
\end{ldescription}
\end{Arguments}
%
\begin{Value}
wide data
\end{Value}
%
\begin{Examples}
\begin{ExampleCode}

## Not run: 
mtcars[1:3, 1:7]
mtcars[1:3, 1:7] %>% long_df()
mtcars[1:3, 1:7] %>% long_df(cols= 3:7)
mtcars[1:3, 1:7] %>% long_df(fix=1:2)
mtcars[1:3, 1:7] %>% long_df(fix= 1:3, rownames_to_column=F)

mtcars[1:3, 1:7] %>% long_df() %>% mutate(ik = "bc", io = "a")

#wide_df
mtcars[1:3, 1:7] %>% long_df() %>% mutate(ik = "bc", io = "a")%>% wide_df()

mtcars[1:3, 1:7] %>% long_df() %>%
  mutate(ik = "bc", io = "a")%>%
  wide_df(names_from = c("names","ik"))


mtcars[1:3, 1:7] %>% long_df() %>%
  mutate(ik = rownames, io = "a")%>%
  wide_df(names_from = c("ik"))%>%pall()

dfa = mtcars[1:3, 1:7] %>% long_df()

dfa%>%
  mutate(ik = "bc", io = "a")%>%
#'   wide_df(names_from = 2)


#일반적인 데이터에서  롱포맷으로 바꾸기
data <- data.frame(
  weight = c(4.17, 5.58, 5.18, 6.11, 4.50, 4.61, 5.17, 4.53, 5.33, 5.14,
             4.81, 4.17, 4.41, 3.59, 5.87, 3.83, 6.03, 4.89, 4.32, 4.69,
             6.31, 5.12, 5.54, 5.50, 5.37, 5.29, 4.92, 6.15, 5.80, 5.26),
  group = c(rep("ctrl", 10), rep("trt1", 10), rep("trt2", 10))
)
data
data%>%add_rep_id("group")

data%>%add_rep_id("group")%>%jjstat::as_trt("id") %>%wide_df("group", "weight")

data  %>%wide_df("group", "weight", add_id="group")


#개별적인 그룹을 만들어진 경우
data2 <- data.frame(
  weight = c(4.17, 5.58, 5.18, 6.11, 4.50, 4.61, 5.17, 4.53, 5.33, 5.14,
             4.81, 4.17, 4.41, 3.59, 5.87, 3.83, 6.03, 4.89, 4.32, 4.69,
             6.31, 5.12, 5.54, 5.50, 5.37, 5.29, 4.92, 6.15, 5.80, 5.26),
  ids = c(paste0("C", 1:10), paste0("T",1:10), paste0("S", 1:10))
)

data2

data2 %>%jjstat::mutate_col(group= c(rep("ctrl", 10),
                 rep("trt1", 10), rep("trt2", 10)), col=2)

# 그룹을 나누어주는 아이디 생성
data2 %>%
jjstat::mutate_col(group= c(rep("ctrl", 10),
        rep("trt1", 10), rep("trt2", 10)), col=2)%>%
        add_rep_id("group")

# 개별적인 그룹을 만드는 열이 존재하는 경우 --> 문제발생
data2 %>%
    jjstat::mutate_col(group= c(rep("ctrl", 10),
          rep("trt1", 10), rep("trt2", 10)), col=2)%>%
          wide_df("group", "weight", add_id="group")

#개별적인 그룹을 만드는 열을 제거하는 경우 --> 문제해결
data2 %>%
  jjstat::mutate_col(group= c(rep("ctrl", 10),
  rep("trt1", 10), rep("trt2", 10)),  col=2)%>%
  wide_df("group", "weight", add_id="group", rm_cols="ids") %>%print(n=Inf)
#'
data3= data.frame(
  weight=c(1:9),
  group=c(rep("a",3),rep("b",3),rep("c",3))
)
data3
data3 %>% wide_df("group", "weight", add_id=2)

## End(Not run)

\end{ExampleCode}
\end{Examples}
\inputencoding{utf8}
\HeaderA{write\_text}{write\_text: Save R Object as Text File and Provide Read Instructions}{write.Rul.text}
%
\begin{Description}
write\_text: Save R Object as Text File and Provide Read Instructions
\end{Description}
%
\begin{Usage}
\begin{verbatim}
write_text(data, file_path = NULL, folder_open = TRUE)
\end{verbatim}
\end{Usage}
%
\begin{Arguments}
\begin{ldescription}
\item[\code{data}] Any R object to be saved as a text file.

\item[\code{file\_path}] Optional. A string specifying the file path where the text file will be saved.
If not provided, the file name is derived from the object's name with ".text" extension.

\item[\code{folder\_open}] Logical indicating whether to open the folder containing the saved file. Defaults to TRUE.
\end{ldescription}
\end{Arguments}
%
\begin{Value}
NULL. The function saves a text file and optionally opens the folder where it is saved.
\end{Value}
%
\begin{Examples}
\begin{ExampleCode}
## Not run: 
# Example: Save a summary of a linear model
lmres <- lm(mpg ~ hp + wt, data = mtcars) %>% summary()
write_text(lmres)
#'
#'
lmres = lm(mpg~ hp+wt, data=mtcars) %>% summary()
write_text(lmres)

# resultlm <- read_text("F:/Rwork/07 futureDuck/2024FEI/lmres.text")
# resultlm
resultlm <- readLines("F:/Rwork/07 futureDuck/2024FEI/lmres.text")



## End(Not run)
\end{ExampleCode}
\end{Examples}
\inputencoding{utf8}
\HeaderA{ztest}{Z test compare coefficient}{ztest}
%
\begin{Description}
Z test compare coefficient
\end{Description}
%
\begin{Usage}
\begin{verbatim}
ztest(b1, se1, b2, se2, type = "res", digits = 3)
\end{verbatim}
\end{Usage}
%
\begin{Arguments}
\begin{ldescription}
\item[\code{b1}] b1 coeff

\item[\code{se1}] se1 standard error

\item[\code{b2}] b2 coeff

\item[\code{se2}] se2 se

\item[\code{type}] result type res1, res2

\item[\code{digits}] rounding defualt 3
\end{ldescription}
\end{Arguments}
%
\begin{Value}
z value and p value
\end{Value}
%
\begin{Examples}
\begin{ExampleCode}
## Not run: 
# hp, wt
ztest(b1= -0.032230,se1= 0.00892,   b2=-3.23 , se2= 0.796 )

## End(Not run)


\end{ExampleCode}
\end{Examples}
\inputencoding{utf8}
\HeaderA{ztest\_df}{Z-test using data.frame(or tibble)}{ztest.Rul.df}
%
\begin{Description}
Z-test using data.frame(or tibble)
\end{Description}
%
\begin{Usage}
\begin{verbatim}
ztest_df(data, row1 = 2, row2 = 3, rname1 = "", rname2 = "", sel = "row")
\end{verbatim}
\end{Usage}
%
\begin{Arguments}
\begin{ldescription}
\item[\code{data}] lm() and data.frame

\item[\code{row1}] row1 = 2

\item[\code{row2}] row2 = 3

\item[\code{rname1}] select first row name

\item[\code{rname2}] select second row name

\item[\code{sel}] row is number, term rowname
\end{ldescription}
\end{Arguments}
%
\begin{Value}
diff, z, p value
\end{Value}
%
\begin{Examples}
\begin{ExampleCode}
## Not run: 
 lm(mpg ~ hp + wt + drat, mtcars) %>%
 ztest_df()

# #selecting variable
lm(mpg ~ hp + wt + drat, mtcars) %>%
   ztest_df(row1=2, row2=3)
#
# # When other variables are selected and analyzed
 lm(mpg ~ hp + wt + drat, mtcars) %>%
   ztest_df(row1=2, row2=4)
#
# # When variable-names are selected
 lm(mpg ~ hp + wt + drat, mtcars) %>%
   ztest_df(sel="term", rname1="hp", rname2="wt")
#
 lm(mpg ~ hp + wt + drat, mtcars) %>%
   ztest_df(sel="term", rname1="hp", rname2="drat")
 
## End(Not run)


\end{ExampleCode}
\end{Examples}
\end{document}
